% !TeX root = ../Book1.tex
\BookPart{分布、Sobolev 与有界变差函数}{b1:part:04}
% 本编导读（章级标题，不计入编号）
\BookPartGuide

对一列光滑函数取极限，函数值可以收敛，导数却未必逐点收敛。若一开始就把可微性作为不可放松的要求，许多由近似过程产生的对象会被排除在外。本编从分部积分出发，把导数理解为对测试函数的作用，再用可积性与范数控制挑出适于分析的函数空间。

\chapref{b1:ch:20}建立分布和弱导数，说明哪些导数仍能由函数表示。\chapref{b1:ch:21}把函数与弱导数同时放进 \(L^p\)，研究所得空间的完备性、基本运算和代表。随后\chapref{b1:ch:22}处理逼近与延拓，\chapref{b1:ch:23}处理边界值。这样，积分意义下定义的对象才逐渐获得可供计算的光滑近似，以及在边界上可以辨认的值。

另一组问题关心控制能够带来什么。\chapref{b1:ch:24}的嵌入与不等式把导数的可积性转化为函数的额外控制，\chapref{b1:ch:25}说明何时还能得到紧性。\chapref{b1:ch:26}的容量则细化例外集的大小。最后两章允许导数成为测度，进入有界变差函数、有限周长和边界结构；前一编的低维面积与可整流性将在这里重新出现。

编内主线及其分支可概括为
\[
\begin{tikzcd}[column sep=1.8em,row sep=1.5em]
  \text{第20章}\arrow[r]
    & \text{第21章}\arrow[r]\arrow[d]\arrow[dd]
    & \text{第22章}\arrow[r]\arrow[d]\arrow[dr]
    & \text{第23章} & \\
  & \text{第24章}\arrow[r]\arrow[d]
    & \text{第25章}\arrow[r]
    & \text{第27章}\arrow[r]
    & \text{第28章} \\
  & \text{第26章} & & &
\end{tikzcd}
\]
第一行给出弱导数、Sobolev 空间、逼近延拓与迹的函数空间主线；\chapref{b1:ch:24}、\chapref{b1:ch:25}把估计推进到嵌入与紧性，\chapref{b1:ch:26}则沿容量形成精细代表支线。\chapref{b1:ch:27}同时回到分布导数、逼近和紧性，随后由\chapref{b1:ch:28}进入有界变差函数的精细结构。

初读可以先沿“弱导数—Sobolev 空间—光滑逼近—基本嵌入与紧性”的路线前进。这条函数空间主线大部分只需要前两编的积分与泛函分析，以及\chapref{b1:ch:13}的卷积逼近，不必先掌握完整的几何测度论。但一维绝对连续代表会用到\secref{b1:sec:1604}，近似微分会调用\secref{b1:sec:1603}；迹与后续有界变差函数的几何结构还会使用第三编的覆盖、面积与可整流性工具。这里应区分“定义弱导数所需的起点”和“精细代表与几何结论所需的后续输入”，也不必把\chapref{b1:ch:19}的全部深层密度判据预先读完。

Brézis 的《Functional Analysis, Sobolev Spaces and Partial Differential Equations》\cite{Brezis2011Functional}适合把本编与前面的泛函分析连读；Leoni 的《A First Course in Sobolev Spaces》\cite{Leoni2017First}提供围绕函数、逼近及代表的另一条教材路线。需要集中查阅嵌入、边界与更细的函数空间时，可再使用 Adams 与 Fournier 的《Sobolev Spaces》\cite{Adams2003Sobolev}。这些伴读各有侧重，阅读时应先比较定义域、指数端点和边界假设，而不是只对照定理名称。

% !TeX root = ../../Book1.tex
\chapter{分布与弱导数}
\label{b1:ch:20}
\BookChapterNumber{b1:ch:20}
% 本章摘要（不计入编号）
\begin{chapterabstract}
分部积分把求导转移到测试函数上，使不可微函数乃至点质量也能拥有导数。本章构造紧支集光滑测试空间，定义分布及其基本操作，证明分布的局部光滑化，再从中识别能够由局部可积函数表示的弱导数。尖角说明少量经典不可微点并不妨碍弱导数存在；跳跃函数与 Cantor 函数则说明，仅有经典导数几乎处处存在仍不足以保证存在局部可积弱导数。
\end{chapterabstract}

% 本章目标（不计入编号）
\begin{learninggoals}
\begin{itemize}
\item\label{b1:item:20-goal-tests}能使用光滑截断与单位分解，解释测试函数收敛为何同时要求共同紧支集和各阶导数控制。
\item\label{b1:item:20-goal-distributions}能从连续线性泛函的定义计算分布导数、限制、支集和光滑乘法，区分点质量与普通函数。
\item\label{b1:item:20-goal-smoothing}能说明分布与光滑核卷积的定义域、微分公式和收敛意义，避免把局部光滑化误作任意延拓。
\item\label{b1:item:20-goal-weak}能检验弱导数的存在性与唯一性，连接经典微分、局部 Lipschitz 函数和一维绝对连续代表。
\end{itemize}
\end{learninggoals}

从\secref{b1:sec:2001}到\secref{b1:sec:2004}，测试空间、其连续对偶以及对偶中仍由函数表示的部分依次出现。测试空间的局部凸拓扑只展开到本章需要的有限阶估计；读者第一次阅读时可以先掌握共同紧支集的收敛判据，再回看其拓扑构造。主要前置工具是\chapref{b1:ch:13}的卷积与 \(L^p\) 平移连续性、\chapref{b1:ch:14}的 Lebesgue 点定理，以及\chapref{b1:ch:16}的 Rademacher 定理和一维绝对连续函数理论。

Hunter 与 Nachtergaele 的《Applied Analysis》\cite[Chapter 11 and Section 12.9]{Hunter2001Applied}适合对照对偶配对、导数和典型例子，其中 Chapter 11 主要面向温和分布，本章则在一般开集上使用紧支集测试函数，不施加无穷远增长条件。Schwartz 的《广义函数论》\cite[第一至三章]{Schwartz2010DistributionsChinese}把测试空间、局部化和导数系统地联系起来，可用作进一步阅读。Sobolev 空间的应用路线可接 Evans 的《Partial Differential Equations》\cite[Chapter 5]{Evans2010Partial}；本章先完成这条路线所需的弱导数语言。

% !TeX root = ../../Book1.tex
\section{测试函数}
\label{b1:sec:2001}

对光滑函数作分部积分，可以把一个函数的导数移到另一个函数上。若第二个函数在边界附近为零，边界项也随之消失。于是，即使第一个函数不可微，积分的另一边仍可能有意义。为了沿着这条思路推广导数，先要选择一类既能任意次求导、又能把计算限制在开集内部的函数。

\ProseParagraphBreak

本节固定非空开集 \(\Omega\subset\mathbb R^d\)，标量域为 \(\mathbb K=\mathbb R\) 或 \(\mathbb C\)。记号 \(K\Subset\Omega\) 表示 \(K\) 是包含于 \(\Omega\) 的紧集。本节的拓扑构造只使用半范数、紧集穷竭和光滑截断，不预借分布或 Sobolev 空间的性质。
% 来源边界：Hunter2001Applied第11章主要使用Schwartz空间；作者公开稿印刷294页明确区分D'与S'。
% 本节DK完备性、单位分解和D的局部凸拓扑由下文完整构造，不把该书温和分布版本冒称为这里的证明。

\subsection{紧支集光滑函数}
\label{b1:subsec:200101}

对多重指标 \(\alpha=(\alpha_1,\ldots,\alpha_d)\in\mathbb N_0^d\)，约定
\[
 |\alpha|\coloneqq\sum_{i=1}^d\alpha_i,\quad
 \alpha!\coloneqq\prod_{i=1}^d\alpha_i!,\quad
 D^\alpha\coloneqq\partial_1^{\alpha_1}\cdots\partial_d^{\alpha_d}.
\]
不等式 \(\beta\leq\alpha\) 指每个坐标分别满足不等式，且
\(\binom\alpha\beta=\alpha!/[\beta!(\alpha-\beta)!]\)。这里 \(D^0\) 为恒等操作。

\begin{definition}\label{b1:def:test-function}
设 \(\varphi:\Omega\rightarrow\mathbb K\) 光滑。如果其支集
\[
 \operatorname{supp}\varphi
 \coloneqq\overline{\{\,x\in\Omega:\varphi(x)\ne0\,\}}^{\,\Omega}
\]
为 \(\Omega\) 的紧子集，则称 \(\varphi\) 为 \(\Omega\) 上的\term[ceshi-hanshu]{测试函数}[test function]。全体测试函数构成的线性空间记作
\[
 \mathcal D(\Omega)\coloneqq C_c^\infty(\Omega).
\]
\end{definition}
\symidx{\(\mathcal D(\Omega)\)：紧支集光滑测试函数空间}

测试函数也有“试验函数”这一译名；《分析学》第2版中译本第6.2节采用后者\footnote{该译名及节次可核对\href{https://academic.hep.com.cn/br/CN/chapter/978-7-04-017381-9/chapter06}{高等教育出版社公布的第六章目录}。这里仅借目录说明名称，不以目录作为下文拓扑构造的依据。}\termidx[shiyan-hanshu]{试验函数}。

\ProseParagraphBreak

支集的紧性和它位于开集内部同样重要。例如 \(\Omega=(0,1)\) 上的常函数 \(1\) 不是测试函数，尽管它处处光滑且有界。测试函数在靠近 \(\partial\Omega\) 的某个区域内为零，因此在 \(\Omega\) 外补零后仍是 \(\mathbb R^d\) 上的光滑函数。以后对测试函数取全空间上确界或积分时，均使用这个零延拓。

\ProseParagraphBreak

这种函数并非只有零函数。令
\[
 \vartheta(t)=
 \begin{cases}e^{-1/t},&t>0,\\0,&t\leq0.\end{cases}
\]
在正半轴上，每阶导数均是 \(e^{-1/t}\) 与 \(1/t\) 的某个多项式的乘积。由 \(s^Ne^{-s}\to0\)，这些导数在 \(t\downarrow0\) 时都趋于零；逐阶检查零点的差商，得到 \(\vartheta\in C^\infty(\mathbb R)\)。于是 \(x\mapsto\vartheta(1-|x|^2)\) 是支集为闭单位球、在开单位球内严格为正的光滑函数。平移与缩放便能在 \(\Omega\) 内任一球上放置这样的函数。

\subsection{基本操作}
\label{b1:subsec:200102}

若 \(\varphi\in\mathcal D(\Omega)\)，则 \(D^\alpha\varphi\in\mathcal D(\Omega)\)，且微分不会扩大支集。若 \(a\in C^\infty(\Omega)\)，则 \(a\varphi\) 也是测试函数；即使 \(a\) 无界，其在 \(\operatorname{supp}\varphi\) 附近的每个固定阶导数仍有界。Leibniz 公式为
\begin{equation}\label{b1:eq:test-leibniz}
 D^\alpha(a\varphi)
 =\sum_{\beta\leq\alpha}\binom\alpha\beta
 D^\beta a\,D^{\alpha-\beta}\varphi.
\end{equation}
这由一阶乘积法则对 \(|\alpha|\) 归纳得到，二项系数记录各次微分落在两个因子上的分配方式。

\ProseParagraphBreak

对 \(K\Subset\Omega\) 记
\[
 \mathcal D_K=\{\,\varphi\in\mathcal D(\Omega):\operatorname{supp}\varphi\subset K\,\},
 \quad p_{K,m}(\varphi)=\max_{|\alpha|\leq m}\sup_{x\in K}|D^\alpha\varphi(x)|.
\]
上确界可等价地取遍 \(\Omega\)，空紧集上的半范数取零。因为导数在 \(K\) 外为零，\eqref{b1:eq:test-leibniz}给出
\[
 p_{K,m}(a\varphi)\leq C_{a,K,m}\cdot p_{K,m}(\varphi),
 \quad p_{K,m}(D^\alpha\varphi)\leq p_{K,m+|\alpha|}(\varphi).
\]
常数只涉及 \(a\) 在 \(K\) 上不超过 \(m\) 阶的导数。

\ProseParagraphBreak

再记 \((\tau_h\varphi)(x)=\varphi(x-h)\)。当 \(h\to0\) 且充分小时，这些平移的支集落在 \(\Omega\) 的同一个紧子集中。各阶导数的一致连续性给出 \(D^\alpha(\tau_h\varphi-\varphi)\to0\) 一致成立。沿坐标方向还可用一维积分公式写出
\[
 \frac{\varphi(x-he_i)-\varphi(x)}h
 =-\int_0^1\partial_i\varphi(x-the_i)\,\mathrm dt.
\]
对任意阶导数重复这一等式，便知差商以相同方式趋于 \(-\partial_i\varphi\)。这比只知道每个点的差商收敛更强：后面连续泛函正需要这种各阶一致控制。

\begin{lemma}\label{b1:lem:smooth-cutoff}
设 \(K\Subset\Omega\)，\(U\subset\Omega\) 为包含 \(K\) 的开集。存在 \(\chi\in\mathcal D(U)\)，满足 \(0\leq\chi\leq1\)，且 \(\chi\) 在 \(K\) 的一个邻域内恒等于 \(1\)。
\end{lemma}
\begin{proof}
由上面的 \(\vartheta\) 构造光滑阶跃
\[
 b(t)=\frac{\vartheta(2-t)}{\vartheta(2-t)+\vartheta(t-1)}.
\]
分母处处为正，且 \(b=1\) 于 \((-\infty,1]\)，\(b=0\) 于 \([2,\infty)\)。选有限个球 \(B(x_j,r_j)\) 覆盖 \(K\)，使 \(\overline B(x_j,\sqrt2r_j)\subset U\)。令 \(b_j(x)=b(|x-x_j|^2/r_j^2)\)，并取
\[
 \chi=1-\prod_j(1-b_j).
\]
每个 \(b_j\) 在其小球上等于 \(1\)，在大球外为零。因此 \(\chi\) 具有所需值域、支集和邻域恒等性质。空紧集情形取 \(\chi=0\)。
\end{proof}

光滑截断把一个局部对象变成紧支集对象，但不会在截断恒等于 \(1\) 的区域内改变它。这个局部性质比截断函数的具体公式更常用。

\subsection{单位分解}
\label{b1:subsec:200103}

\begin{definition}\label{b1:def:smooth-partition-unity}
设 \((U_\lambda)_{\lambda\in\Lambda}\) 为 \(\Omega\) 的开覆盖。如果一族非负光滑函数 \((\psi_j)\) 的支集在 \(\Omega\) 中局部有限，每个支集包含于某个 \(U_{\lambda(j)}\)，并且
\[
 \sum_j\psi_j(x)=1\quad(x\in\Omega),
\]
则称它为从属于该覆盖的\term[guanghua-danwei-fenjie]{光滑单位分解}[smooth partition of unity]。局部有限指每个点有一个邻域，只与有限个支集相交。
\end{definition}

\begin{theorem}\label{b1:thm:smooth-partition-unity}
每个欧氏开集的开覆盖都有一族可数、从属于该覆盖的光滑单位分解，而且每个分解函数均有紧支集。
\end{theorem}
\begin{proof}
取紧集穷竭
\[
 K_j=\{\,x\in\Omega:|x|\leq j,\ \operatorname{dist}(x,\mathbb R^d\setminus\Omega)\geq1/j\,\},
 \quad j\geq1,
\]
空补集的距离约定为 \(+\infty\)，并令 \(K_0=K_{-1}=\varnothing\)。这些紧集满足 \(K_j\subset\operatorname{int}K_{j+1}\)，且其内部之并为 \(\Omega\)。紧集
\(A_j=K_j\setminus\operatorname{int}K_{j-1}\) 包含在开集 \(\operatorname{int}K_{j+1}\setminus K_{j-2}\) 内。

对每个 \(A_j\)，选有限个小球覆盖它，使各小球的稍大闭球仍包含在 \(\operatorname{int}K_{j+1}\setminus K_{j-2}\) 与某个 \(U_\lambda\) 的交中。由\cref{b1:lem:smooth-cutoff}，可取在小球闭包上等于 \(1\)、支集在相应大球内的非负光滑函数 \(\eta_{j,l}\)。这族函数的支集局部有限：给定点先取其邻域包含在某个 \(\operatorname{int}K_N\) 中，所有 \(j\geq N+2\) 的支集均避开该邻域，其余层只有有限个函数。

各层小球覆盖 \(\Omega\)，故局部有限和 \(s=\sum_{j,l}\eta_{j,l}\) 光滑且处处为正。令 \(\psi_{j,l}=\eta_{j,l}/s\)，再把双指标排成一个序列。除以正光滑函数不扩大支集，所需各性质随即成立。
\end{proof}

\begin{corollary}\label{b1:cor:precompact-smooth-partition}
若 \(\Omega\subseteq\mathbb R^d\) 为非空开集，则存在可数的光滑单位分解 \((\psi_i)\) 及局部有限开覆盖 \((V_i)\)，使
\[
 \psi_i\in C_c^\infty(V_i),\qquad V_i\Subset\Omega.
\]
特别地，每个紧集只与有限多个 \(V_i\) 相交。
\end{corollary}
\begin{proof}
在上一证明中保留承载各个 \(\eta_{j,l}\) 的稍大开球，并与相应的 \(\psi_{j,l}\) 一同枚举。这些开球的闭包均包含在 \(\Omega\) 中；每层只有有限个，而足够远的层避开任意给定紧集，故它们构成局部有限开覆盖。归一化不扩大支集，所以相应的 \(\psi_i\) 属于 \(C_c^\infty(V_i)\)。
\end{proof}

若 \(\varphi\) 是测试函数，局部有限性和 \(\operatorname{supp}\varphi\) 的紧性使
\[
 \varphi=\sum_j\psi_j\varphi
\]
实际上只有有限个非零项。于是可以先在每个覆盖开集上验证一个关于测试函数的线性恒等式，再把有限项相加恢复全局恒等式。局部有限的要求不能换成“可数”二字；一般可数和既未必光滑，也未必允许逐项操作。

\subsection{测试函数空间拓扑}
\label{b1:subsec:200104}

先在 \(\mathcal D_K\) 上使用半范数 \(p_{K,m}\) 所给的拓扑。其收敛就是各阶导数一致收敛，可由度量
\[
 \mathrm{d}_K(\varphi,\psi)=\sum_{m=0}^\infty2^{-m-1}\min\{1,p_{K,m}(\varphi-\psi)\}
\]
描述。这个空间完备。事实上，Cauchy 列的每个导数都有一致极限 \(g_\alpha\)。在包含于 \(\Omega\) 的小球中，将普通积分恒等式
\[
 D^\alpha\varphi_j(x+he_i)-D^\alpha\varphi_j(x)
 =\int_0^h D^{\alpha+e_i}\varphi_j(x+te_i)\,\mathrm dt
\]
传到极限，得到 \(\partial_i g_\alpha=g_{\alpha+e_i}\)。故 \(g_0\) 光滑，其支集仍在 \(K\)，且原列在每个半范数中收敛到它。

\ProseParagraphBreak

整个 \(\mathcal D(\Omega)\) 是这些固定支集空间的并，却不能直接赋予 \(C^\infty(\Omega)\) 的局部一致拓扑。例如，在 \(\Omega=\mathbb R^d\) 中，一个固定小凸起向无穷远平移，在每个紧集上终将为零，但其积分始终不变。若这样的列在测试空间中趋于零，积分泛函就会失去连续性。

\ProseParagraphBreak

以下从半范数族构造完整局部凸拓扑的部分属于技术性补充；第一次阅读可以先接受随后\cref{b1:prop:test-sequence-convergence}与\cref{b1:prop:test-functional-continuity}给出的序列判据和有限阶估计，再回读这一构造。为避免把一个收敛列的判据误当成对全部开集的定义，设 \(\mathscr P\) 为 \(\mathcal D(\Omega)\) 上所有这样的半范数 \(q\)：对每个 \(K\Subset\Omega\)，存在 \(C>0\) 与整数 \(m\geq0\)，使
\[
 q(\varphi)\leq C\cdot p_{K,m}(\varphi)\quad(\varphi\in\mathcal D_K).
\]
以有限个 \(q\in\mathscr P\) 的小于正数的集合之交作为零点邻域基，定义 \(\mathcal D(\Omega)\) 的拓扑。半范数的不等式保证加法、数乘连续，而且每个 \(\mathcal D_K\) 的包含映射连续。由于这些邻域都是凸的，这是一种局部凸拓扑；它也是使各包含映射连续的最强局部凸拓扑。这最后一句只须注意：任何另一个由半范数给出的此类拓扑，其定义半范数在每个 \(\mathcal D_K\) 上连续，因而都已收入 \(\mathscr P\)。

\begin{proposition}\label{b1:prop:test-sequence-convergence}
序列 \(\varphi_j\) 在 \(\mathcal D(\Omega)\) 中趋于零，当且仅当存在同一个 \(K\Subset\Omega\) 包含所有 \(\operatorname{supp}\varphi_j\)，并且对每个 \(m\geq0\)，有 \(p_{K,m}(\varphi_j)\to0\)。
\end{proposition}
\begin{proof}
若具有右侧性质，则每个 \(q\in\mathscr P\) 的固定紧集估计给出 \(q(\varphi_j)\to0\)，所以在所定义拓扑中收敛。

反之，假定 \(\varphi_j\to0\)。全空间上各阶导数的上确界半范数属于 \(\mathscr P\)，故这些量趋于零。还须证明支集可统一限制。若不能，对上面紧集穷竭可选一个子列 \(\varphi_{j_l}\) 及点 \(x_l\notin K_l\)，使 \(\varphi_{j_l}(x_l)\ne0\)。每个有限的前缀有共同紧支集，因此指标可递增选取。点族 \((x_l)\) 在 \(\Omega\) 中局部有限。于是
\[
 q(\varphi)=\sum_l\frac{|\varphi(x_l)|}{|\varphi_{j_l}(x_l)|}
\]
是良定义的半范数：每个紧支集只遇到有限个 \(x_l\)，在固定 \(\mathcal D_K\) 上有 \(q\leq C_K\cdot p_{K,0}\)。故 \(q\in\mathscr P\)，但 \(q(\varphi_{j_l})\geq1\)，与收敛矛盾。这证明共同紧支集存在。
\end{proof}

\begin{proposition}\label{b1:prop:test-functional-continuity}
设 \(T:\mathcal D(\Omega)\rightarrow\mathbb K\) 线性。下列三个条件等价：
\begin{equivalences}
\item\label{b1:item:test-functional-topological}泛函 \(T\) 在上述拓扑中连续；
\item\label{b1:item:test-functional-sequential}每当 \(\varphi_j\to0\) 于 \(\mathcal D(\Omega)\)，都有 \(T\varphi_j\to0\)；
\item\label{b1:item:test-functional-order}对每个 \(K\Subset\Omega\)，存在 \(C_K>0\)、整数 \(m_K\geq0\)，使
\[
 |T\varphi|\leq C_K\cdot p_{K,m_K}(\varphi)\quad(\varphi\in\mathcal D_K).
\]
\end{equivalences}
\end{proposition}
\begin{proof}
连续性蕴涵序列连续性。若第三项对某个 \(K\) 失败，则对每个 \(j\geq1\) 可取 \(u_j\in\mathcal D_K\)，满足 \(|Tu_j|>j\cdot p_{K,j}(u_j)\)。令 \(v_j=u_j/|Tu_j|\)，则对每个固定的 \(m\)，在 \(j\geq m\) 时有 \(p_{K,m}(v_j)<1/j\)，但 \(|Tv_j|=1\)。这与第二项矛盾。最后，第三项说明半范数 \(\varphi\mapsto|T\varphi|\) 属于 \(\mathscr P\)，从而 \(T\) 连续。
\end{proof}

这个有限阶估计是后面使用测试函数拓扑的主要接口。微分和乘以固定光滑函数都保持测试列收敛；在线性泛函上，这足以证明转移过去的相应操作仍连续。紧集改变时允许估计中的阶数和常数一起改变，没有理由要求一个无界开集上的所有位置受同一有限阶控制。

% !TeX root = ../../Book1.tex
\section{分布}
\label{b1:sec:2002}

一个函数可以通过与测试函数的积分来辨认；分部积分又能把求导转移到测试函数上。沿着这个方向，研究对象不必先在每一点有函数值，只需规定它怎样作用于测试函数。这样得到的空间既包含局部可积函数，也容纳点质量及其任意阶导数。

% 实读 Hunter--Nachtergaele 作者官方 ch11.pdf，印刷291--292、294--297页：
% 11.2 的对偶配对、例11.5--11.6，11.3 的光滑乘法、定义11.10与例11.14。
% 原书这一段主要讨论 Schwartz 测试空间及温和分布；294页明确另有 D'。
% 本节改用任意开集上的 D，并自足展开局部有限阶、限制、支集和局部唯一性；
% 不移用原书的无穷远多项式增长条件，不依赖尚未建立的20.3光滑化。
% 中文异名据 Schwartz2010DistributionsChinese 高等教育出版社目录第一章核定；
% 本轮未读该中译本正文，目录只作为名称及伴读位置的证据，不承担数学证明。

\subsection{分布的定义}
\label{b1:subsec:200201}

本节固定开集 \(\Omega\subseteq\mathbb R^d\)，默认测试函数取复值，沿用 \(\mathcal D(\Omega)=C_c^\infty(\Omega)\)。光滑函数的偏导数记作 \(\partial^\alpha\)，与上一节的 \(D^\alpha\) 相同。

\begin{definition}\label{b1:def:distribution}
设 \(T:\mathcal D(\Omega)\rightarrow\mathbb C\) 为复线性泛函。若每当 \(\varphi_j\rightarrow0\) 为测试函数收敛列时都有 \(T(\varphi_j)\rightarrow0\)，则称 \(T\) 为 \(\Omega\) 上的\term[fenbu]{分布}[distribution]。所有这样的泛函组成向量空间 \(\mathcal D'(\Omega)\)，配对记作
\[
 \langle T,\varphi\rangle=T(\varphi).
\]
配对对两个变量均线性，不对测试函数取复共轭。
\end{definition}
\symidx{\(\mathcal D'(\Omega)\)：分布空间}

分布也称广义函数；这一中文名称见 Schwartz 的《广义函数论》\cite[第一章]{Schwartz2010DistributionsChinese}\termidx[guangyi-hanshu]{广义函数}\foreignidx{generalized function}。这里的“广义”不是给每个点补一个广义数值，而是改变描述对象的方式。实值理论采用实测试函数和实线性泛函，以下证明不变。

由\cref{b1:prop:test-functional-continuity}，上述连续性等价于：对每个紧集 \(K\subseteq\Omega\)，存在常数 \(C_K\) 和非负整数 \(m_K\)，使
\begin{equation}\label{b1:eq:distribution-local-order}
 |\langle T,\varphi\rangle|
 \leq C_K p_{K,m_K}(\varphi)
 \quad(\varphi\in\mathcal D_K).
\end{equation}
这就是分布的局部有限阶估计。阶数允许随 \(K\) 改变；若能对所有紧集选用同一个整数 \(m\)，则称这个\term[fenbu-de-jie]{分布的阶}[order of a distribution]不超过 \(m\)。局部估计并不要求 \(\Omega\) 上存在一个统一的阶数或常数。

\begin{definition}\label{b1:def:distribution-convergence}
设 \(T_j,T\in\mathcal D'(\Omega)\)。若对每个固定的 \(\varphi\in\mathcal D(\Omega)\) 都有
\[
 \langle T_j,\varphi\rangle\longrightarrow\langle T,\varphi\rangle,
\]
则称 \(T_j\) 在分布意义下收敛于 \(T\)，记作 \(T_j\rightarrow T\) 于 \(\mathcal D'(\Omega)\)。
\end{definition}

这里变化的是分布而不是测试函数；它与定义分布时要求的测试函数列连续性是两个不同的量词安排。

\begin{definition}\label{b1:def:distribution-support}
对开集 \(U\subseteq\Omega\)，将 \(\varphi\in\mathcal D(U)\) 在 \(\Omega\setminus U\) 上延为零，所得测试函数记为 \(\widetilde\varphi\)。规定
\[
 \langle T|_U,\varphi\rangle=\langle T,\widetilde\varphi\rangle,
\]
称 \(T|_U\) 为 \(T\) 在 \(U\) 上的限制。若 \(T|_U=0\)，则称 \(T\) 在 \(U\) 上为零。定义\term[fenbu-zhiji]{分布的支集}[support of a distribution]为
\[
 \operatorname{supp}T
 =\Omega\setminus\bigcup\{\,U\subseteq\Omega:U\text{ open},\ T|_U=0\,\}.
\]
\end{definition}
\symidx{\(\operatorname{supp}T\)：分布的支集}

\begin{proposition}\label{b1:prop:distribution-locality}
限制 \(T|_U\) 是 \(U\) 上的分布，并与进一步限制相容。支集是 \(\Omega\) 中的相对闭集；若测试函数的支集避开 \(\operatorname{supp}T\)，则 \(T\) 在该测试函数上的值为零。特别地，两个分布在某个开覆盖的每个成员上相等，就在 \(\Omega\) 上相等。
\end{proposition}
\begin{proof}
紧支集测试函数在 \(U\) 的边界附近已经为零，因而延零后仍光滑，且延零保持共同紧支集及所有阶导数的一致收敛。这证明限制的连续性；复线性和限制的相容性直接由定义得到。

记定义中零开集的并为 \(V\)。若 \(\varphi\in\mathcal D(V)\)，其紧支集被有限个这样的开集 \(U_1,\ldots,U_N\) 覆盖。由\cref{b1:thm:smooth-partition-unity}，可将它写成有限和 \(\varphi=\sum_{\ell=1}^N\varphi_\ell\)，其中 \(\varphi_\ell\in\mathcal D(U_\ell)\)。每项配对为零，所以 \(T|_V=0\)。支集外测试的结论随即成立。对两个分布之差作同一论证，就得到开覆盖上的唯一性。
\end{proof}

支集的意义还在于局部化：如果两个测试函数在 \(\operatorname{supp}T\) 的某个邻域相同，它们之差的支集便避开 \(\operatorname{supp}T\)，所以配对相同。对紧支集分布，这使我们能够扩大测试函数的范围。

\begin{corollary}\label{b1:cor:compact-distribution-smooth-action}
若 \(S=\operatorname{supp}T\) 是包含于 \(\Omega\) 的紧集，则 \(T\) 有有限阶。选取在 \(S\) 的邻域为 \(1\) 的 \(\chi\in\mathcal D(\Omega)\)，可以对任意 \(f\in C^\infty(\Omega)\) 规定
\[
 \langle T,f\rangle=\langle T,\chi f\rangle.
\]
这一规定不依赖 \(\chi\)，并与原来在测试函数上的作用一致。存在固定紧集 \(L\subseteq\Omega\)、整数 \(m\) 和常数 \(C\)，使
\[
 |\langle T,f\rangle|
 \leq C\max_{|\alpha|\leq m}\sup_{x\in L}|\partial^\alpha f(x)|.
\]
\end{corollary}
\begin{proof}
截断函数由\cref{b1:lem:smooth-cutoff}给出。若 \(\chi_1,\chi_2\) 都满足要求，则 \((\chi_1-\chi_2)f\) 为紧支集光滑函数，并在 \(S\) 的邻域为零，由上一命题，其配对为零。对原本紧支集的 \(f\)，同样将 \((1-\chi)f\) 配对为零，得到一致性。

在 \(L=\operatorname{supp}\chi\) 上应用\eqref{b1:eq:distribution-local-order}，再展开 \(\chi f\) 的有限阶导数，就得到所述估计。对任意测试函数 \(\varphi\)，\(L\) 上的导数上确界不超过 \(\operatorname{supp}\varphi\) 上相同阶导数的上确界，因为所有这些导数在该支集之外为零。因此这里的同一个 \(m\) 控制所有测试函数，证明 \(T\) 有有限阶。
\end{proof}

\subsection{正则分布}
\label{b1:subsec:200202}

\begin{definition}\label{b1:def:regular-distribution}
给定 \(f\in L^1_{\mathrm{loc}}(\Omega)\)，规定
\[
 \langle T_f,\varphi\rangle
 =\int_\Omega f(x)\varphi(x)\,\mathrm dx.
\]
由这种局部可积函数产生的分布称为\term[zhengze-fenbu]{正则分布}[regular distribution]。
\end{definition}

这个规定确实给出分布。若 \(\operatorname{supp}\varphi\subseteq K\)，则
\[
 |\langle T_f,\varphi\rangle|
 \leq\left(\int_K|f|\,\mathrm dx\right)p_{K,0}(\varphi).
\]
因此 \(T_f\) 为零阶分布。积分只看到几乎处处等价类；下面证明它也没有丢失等价类以外的信息。

\begin{proposition}\label{b1:prop:regular-distribution-injective}
映射 \(f\mapsto T_f\) 是从 \(L^1_{\mathrm{loc}}(\Omega)\) 到 \(\mathcal D'(\Omega)\) 的线性单射。若 \(f_j\rightarrow f\) 于 \(L^1_{\mathrm{loc}}\)，则 \(T_{f_j}\rightarrow T_f\) 于分布意义。
\end{proposition}
\begin{proof}
只需证明 \(T_f=0\) 蕴含 \(f=0\) 几乎处处。选取非负的 \(\rho\in C_c^\infty\bigl(B(0,1)\bigr)\)，使 \(\int\rho=1\)。这样的函数由上一节的光滑截断归一化得到。对 \(f\) 的一个 Lebesgue 点 \(x\in\Omega\)，当 \(r>0\) 足够小时，
\[
 \varphi_{x,r}(y)=r^{-d}\rho\left(\frac{y-x}r\right)
\]
属于 \(\mathcal D(\Omega)\)，所以 \(\int f\varphi_{x,r}=0\)。利用其积分为 \(1\)，有
\[
 \begin{aligned}
 |f(x)|
 &=\left|\int_\Omega\bigl(f(x)-f(y)\bigr)\varphi_{x,r}(y)\,\mathrm dy\right|\\
 &\leq\frac{\|\rho\|_\infty}{r^d}
        \int_{B(x,r)}|f(y)-f(x)|\,\mathrm dy
 \longrightarrow0.
 \end{aligned}
\]
由\cref{b1:thm:lebesgue-point}，几乎每一点都是这样的点，故单射性成立。最后，对固定测试函数及 \(K=\operatorname{supp}\varphi\)，
\[
 |\langle T_{f_j}-T_f,\varphi\rangle|
 \leq\|\varphi\|_\infty\cdot\|f_j-f\|_{L^1(K)}
 \longrightarrow0.
\]
\end{proof}

因此以后可以把 \(f\) 与 \(T_f\) 视为同一对象，但不能据此为一般分布写下点值 \(T(x)\)。同一命题在每个开子集上适用，故 \(T_f\) 在一个开集上为零，当且仅当 \(f\) 在那里几乎处处为零。这也说明分布支集对应的是函数的本质支集，而不是任选代表的非零点集合之闭包。

\subsection{Dirac 分布}
\label{b1:subsec:200203}

\begin{definition}\label{b1:def:dirac-distribution}
给定 \(a\in\Omega\)，以
\[
 \langle\delta_a,\varphi\rangle=\varphi(a)
\]
定义 \(a\) 处的\term[dirac-fenbu]{Dirac 分布}[Dirac distribution]。
\end{definition}

估计 \(|\varphi(a)|\leq p_{K,0}(\varphi)\) 对 \(a\in K\) 成立；若 \(a\notin K\)，配对为零。因此 \(\delta_a\) 是零阶分布。它与点质量测度作用于测试函数所得的泛函一致，但不是局部可积函数。

\begin{proposition}\label{b1:prop:dirac-not-regular}
Dirac 分布不是正则分布。更一般地，对任意多重指标 \(\alpha\)，泛函
\[
 \langle S_{a,\alpha},\varphi\rangle
 =(-1)^{|\alpha|}\partial^\alpha\varphi(a)
\]
是非正则分布，且其支集恰为 \(\{a\}\)。
\end{proposition}
\begin{proof}
点上导数由 \(p_{K,|\alpha|}\) 控制，所以 \(S_{a,\alpha}\) 连续。选取在原点附近为 \(1\) 的光滑截断 \(\chi\)，支集包含于 \(B(0,1)\)，令
\[
 \psi(z)=\frac{z^\alpha}{\alpha!}\chi(z),\quad
 \psi_r(x)=r^{|\alpha|}\psi\left(\frac{x-a}r\right).
\]
当 \(r\) 足够小时，\(\psi_r\in\mathcal D(\Omega)\)，且 \(\partial^\alpha\psi_r(a)=1\)。若 \(S_{a,\alpha}=T_g\)，则一方面配对的绝对值恒为 \(1\)，另一方面
\[
 |\langle T_g,\psi_r\rangle|
 \leq r^{|\alpha|}\|\psi\|_\infty
            \int_{B(a,r)}|g(x)|\,\mathrm dx
 \longrightarrow0,
\]
其中使用 \(g\) 的局部可积性及积分的绝对连续性，产生矛盾。取 \(\alpha=0\) 即包含 Dirac 分布的情形。

支集避开 \(a\) 的测试函数在 \(a\) 的邻域恒为零，其所有点上导数也为零；反过来，任意包含 \(a\) 的开集都容纳上述某个 \(\psi_r\)，其配对非零。由支集定义得到结论。
\end{proof}

Dirac 分布揭示了“零阶”与“正则”之间的区别：只需要函数值估计的泛函仍可能集中在一个 Lebesgue 零测集上。Hunter 与 Nachtergaele 的《Applied Analysis》\cite[Sections 11.2--11.3, pp. 291--296]{Hunter2001Applied}给出了点质量、正则分布和导数的同类构造；该书此处主要使用温和分布，本章的紧支集测试则不要求无穷远处的增长限制。

\subsection{分布导数}
\label{b1:subsec:200204}

若 \(f\) 足够光滑，反复分部积分会把它的偏导数转移到测试函数上，而紧支集消去边界项。这提示我们把右侧作为一般分布的求导规则。

\begin{definition}\label{b1:def:distributional-derivative}
设 \(T\in\mathcal D'(\Omega)\)，\(\alpha\) 为多重指标。定义其\term[fenbu-daoshu]{分布导数}[distributional derivative]为
\begin{equation}\label{b1:eq:distributional-derivative}
 \langle\partial^\alpha T,\varphi\rangle
 =(-1)^{|\alpha|}\langle T,\partial^\alpha\varphi\rangle.
\end{equation}
对 \(a\in C^\infty(\Omega)\)，定义光滑函数与分布的乘法为
\[
 \langle aT,\varphi\rangle=\langle T,a\varphi\rangle.
\]
\end{definition}
\symidx{\(\partial^\alpha T\)：分布导数}

\begin{proposition}\label{b1:prop:distribution-derivative-calculus}
上述导数与乘积都是分布。分布偏导数可交换，并满足
\[
 \partial^\alpha\partial^\beta T=\partial^{\alpha+\beta}T.
\]
求导与乘以固定光滑函数都保持分布收敛。若 \(f\in C^{|\alpha|}(\Omega)\)，则
\[
 \partial^\alpha T_f=T_{\partial^\alpha f}.
\]
此外，限制与求导、光滑乘法相容，且
\[
 \operatorname{supp}(\partial^\alpha T)\subseteq\operatorname{supp}T,\quad
 \operatorname{supp}(aT)\subseteq\operatorname{supp}a\cap\operatorname{supp}T.
\]
\end{proposition}
\begin{proof}
固定 \(K\)，若 \(T\) 满足阶数为 \(m\) 的局部估计，则
\[
 |\langle\partial^\alpha T,\varphi\rangle|
 \leq C_K p_{K,m+|\alpha|}(\varphi).
\]
光滑乘积的通常求导公式给出
\[
 p_{K,m}(a\varphi)
 \leq C_{K,m,a}p_{K,m}(\varphi),
\]
因为其中只出现 \(a\) 在 \(K\) 上有限多个有界导数。代入 \(T\) 的估计，就证明了 \(aT\) 的连续性。

对同一测试函数连续使用定义，得到
\[
 \langle\partial^\alpha\partial^\beta T,\varphi\rangle
 =(-1)^{|\alpha|+|\beta|}
       \langle T,\partial^{\alpha+\beta}\varphi\rangle,
\]
所以混合偏导数可交换。若 \(T_j\rightarrow T\)，分别用固定测试函数 \(\partial^\alpha\varphi\) 和 \(a\varphi\) 测试，即得对应导数与乘积的收敛。

对经典可微的 \(f\)，先将测试函数用光滑划分分为支集包含在 \(\Omega\) 内小矩形中的有限项。对每项使用 Fubini 定理和一维分部积分，反复 \(|\alpha|\) 次可得
\[
 \int_\Omega(\partial^\alpha f)\varphi\,\mathrm dx
 =(-1)^{|\alpha|}\int_\Omega f\,\partial^\alpha\varphi\,\mathrm dx.
\]
测试函数在每个矩形边界附近为零，故没有边界项。求和证明与经典导数一致。

延零与测试函数的求导、乘法相容，所以限制也相容。若 \(T\) 在某开集上为零，则其导数和 \(aT\) 在该开集上为零；若 \(a\) 在开集上恒为零，则 \(aT\) 也为零。支集的包含关系由此得到。
\end{proof}

\begin{proposition}[Leibniz 公式]\label{b1:prop:distribution-leibniz}
对 \(a\in C^\infty(\Omega)\) 与任意多重指标 \(\alpha\)，有
\[
 \partial^\alpha(aT)
 =\sum_{\beta\leq\alpha}\binom{\alpha}{\beta}
      (\partial^\beta a)\,\partial^{\alpha-\beta}T.
\]
\end{proposition}
\begin{proof}
先求一阶导数。对任意测试函数，
\[
 \begin{aligned}
 \langle\partial_i(aT),\varphi\rangle
 &=-\langle T,a\,\partial_i\varphi\rangle\\
 &=-\langle T,\partial_i(a\varphi)\rangle
       +\langle T,(\partial_i a)\varphi\rangle\\
 &=\langle a\,\partial_iT+(\partial_i a)T,\varphi\rangle.
 \end{aligned}
\]
再对 \(|\alpha|\) 归纳。给归纳式施加 \(\partial_i\)，每一项分别对光滑系数和分布求导；同一项的两个系数按
\(\binom{\alpha}{\beta}+\binom{\alpha}{\beta-e_i}
=\binom{\alpha+e_i}{\beta}\)
合并，越界系数规定为零，便得到下一阶公式。
\end{proof}

由定义，上一小节的点导数泛函正是 \(\partial^\alpha\delta_a\)，所以每个分布都能任意阶求导，但导数不必是正则分布。例如在一维有 \(x\delta_0=0\)，一阶 Leibniz 公式给出
\[
 x\delta_0'=-\delta_0.
\]
这些运算不依赖点值，却仍能给出精确的代数恒等式。对于一般两个分布，上述定义没有赋予它们乘法，因而不能把 \(\delta_0^2\) 等符号作为已经存在的对象使用。

局部有限阶也不等于全局有限阶：把越来越高阶的点导数放在趋于无穷远的位置，可以使每个紧集只遇到有限多项，却找不到统一的阶数界。具体构造留在习题\ref{b1:ex:20-locally-finite-unbounded-order}。

分布空间容纳这种随位置变化的阶数，同时又允许统一求导。本章的下一步是研究如何用光滑函数逼近分布；再把分布导数要求为局部可积函数，就导向\secref{b1:sec:2004}的弱导数。

% !TeX root = ../../Book1.tex
\section{卷积与光滑化}
\label{b1:sec:2003}

\subsecref{b1:subsec:130403}已经用可积核构造函数的逼近恒等。本节把平均推广到分布：把核的平移看成测试函数，由分布作用于它。所得对象虽然是普通光滑函数，其导数却能保留原分布的微分信息。整个构造先在远离边界的区域进行，不预设分布可以延拓到开集之外。

% 实读 Hunter2001Applied 作者官方 ch11.pdf，11.3--11.4，印刷 pp.297--301。
% 原书讨论 Schwartz 测试函数与缓增分布，11.23 只有稠密性证明概要；
% 本节一般开集 D' 的有限阶差商、共同紧支集 Riemann 和及局部化证明另行完整组织。
% 名称依据：Oleinik2008PDEChinese官方目录1.2用“磨光函数”；只核名称，未读该节定义。
% “磨光子”另见北京大学公开《实变函数》前言，来源记录限定到名称，不据此声称核对其归一化。

\subsection{磨光函数}
\label{b1:subsec:200301}

取非负偶函数 \(\rho\in C_c^\infty(\mathbb R^d)\)，使
\[
 \operatorname{supp}\rho\subset\overline{B(0,1)},
 \quad \int_{\mathbb R^d}\rho=1,
 \quad \rho_\varepsilon(x)\coloneqq\varepsilon^{-d}\cdot \rho(x/\varepsilon).
\]
满足上述条件的核称为本章的\term[moguang-hanshu]{磨光函数}[mollifier]。名称可对照 Oleinik 的《偏微分方程讲义》\cite[第1.2节]{Oleinik2008PDEChinese}；中文资料也使用“磨光子”这一名称。\termidx[moguangzi]{磨光子}\footnote{北京大学数学学院公开的《实变函数》前言，在介绍\chapref{b1:ch:04}的 \(L^p\) 与 Sobolev 空间工具时采用“磨光子”。该用词说明见\href{https://math101.pku.edu.cn/docs/2025-04/b74312fec0b54f0c90ff66270f2bded8.pdf}{学校公布的前言}；这里仅提示名称，不据前言推断其核函数的全部约定。}

\ProseParagraphBreak

例如可将 \(|x|<1\) 上的 \(\exp\bigl(-1/(1-|x|^2)\bigr)\) 在球外补零，再除以其正积分；光滑性已在\secref{b1:sec:2001}的截断构造中说明。这是\cref{b1:def:approximate-identity}中的一种特殊核，光滑性和紧支集是本节新增的要求，偶性则用来简化反射记号。

\ProseParagraphBreak

对开集 \(\Omega\subset\mathbb R^d\)，记
\[
 \Omega_\varepsilon
 \coloneqq\{\,x\in\Omega\mid\operatorname{dist}(x,\Omega^c)>\varepsilon\,\}.
\]
约定到空集的距离为无穷，故 \(\Omega=\mathbb R^d\) 时 \(\Omega_\varepsilon=\mathbb R^d\)。若 \(x\in\Omega_\varepsilon\)，则 \(y\mapsto\rho_\varepsilon(x-y)\) 的支集紧含于 \(\Omega\)；这才保证一般分布可以作用于它。

\subsection{分布卷积}
\label{b1:subsec:200302}

\begin{definition}\label{b1:def:distribution-test-convolution}
设 \(T\in\mathcal D'(\Omega)\)、\(\eta\in C_c^\infty(\mathbb R^d)\)，且 \(\operatorname{supp}\eta\subset\overline{B(0,a)}\)，其中 \(a>0\)。在 \(\Omega_a\) 上定义 \(T\) 与测试函数 \(\eta\) 的\term[fenbu-juanji]{分布卷积}[convolution of a distribution with a test function]为
\[
 (T*\eta)(x)\coloneqq\langle T_y,\eta(x-y)\rangle.
\]
下标 \(y\) 指明分布所作用的变量，不表示对 \(T\) 作逐点取值。
\end{definition}
\symidx{\(T*\eta\)：分布与测试函数的卷积}

若 \(T=T_f\) 来自 \(f\in L^1_{\mathrm{loc}}(\Omega)\)，该定义成为
\[
 (T_f*\eta)(x)=\int_\Omega f(y)\cdot\eta(x-y)\,\mathrm dy.
\]
被积函数支集紧含于 \(\Omega\)，所以每个 \(x\in\Omega_a\) 处都绝对可积，并与前文函数卷积一致。这里没有定义两个任意分布的卷积；即使两个常函数 \(1\) 的通常卷积，在全空间也会发散。紧支集等附加条件可以支持更广的分布卷积，但不参与本节证明。

\begin{proposition}\label{b1:prop:distribution-convolution-smooth}
上述 \(T*\eta\) 属于 \(C^\infty(\Omega_a)\)。对任意多重指标 \(\alpha\)，在 \(\Omega_a\) 上有
\[
 D^\alpha(T*\eta)=T*(D^\alpha\eta)=(D^\alpha T)*\eta.
\]
\end{proposition}
\begin{proof}
固定 \(K\Subset\Omega_a\)，取 \(K\) 的稍大紧邻域 \(K'\Subset\Omega_a\)。当 \(x\in K'\) 时，所有 \(\eta(x-\cdot)\) 的支集都在紧集 \(H=K'+\overline{B(0,a)}\Subset\Omega\) 中。由\eqref{b1:eq:distribution-local-order}，存在 \(C,m\) 使
\[
 |\langle T,\psi\rangle|\le C\cdot p_{H,m}(\psi)
 \quad(\psi\in\mathcal D_H).
\]
对每个坐标方向 \(e_j\)，Taylor 公式及 \(\eta\) 各阶导数的有界性给出
\[
 p_{H,m}\left(
 \frac{\eta(x+he_j-\cdot)-\eta(x-\cdot)}h
 -(\partial_j\eta)(x-\cdot)\right)\le C_\eta|h|
\]
对 \(x\in K\) 和充分小的 \(h\) 一致成立。将 \(T\) 作用于差商，便证明第一阶导数公式；同一估计也证明所得导数连续。把 \(\eta\) 换成其各阶导数，迭代即得全部光滑性与第一项等式。

分布导数的定义又给出
\[
 \langle D^\alpha T_y,\eta(x-y)\rangle
 =(-1)^{|\alpha|}\langle T_y,D_y^\alpha\eta(x-y)\rangle
 =\langle T_y,(D^\alpha\eta)(x-y)\rangle,
\]
因为对 \(y\) 微分还产生一个 \((-1)^{|\alpha|}\)。这证明第二项等式。
\end{proof}

例如在全空间，\(\delta_b*\eta=\eta(\cdot-b)\)，而 \((D^\alpha\delta_b)*\eta=D^\alpha\eta(\cdot-b)\)。点质量及其导数被变成普通光滑函数，但导数阶数没有丢失。

\subsection{正则化}
\label{b1:subsec:200303}

为了讨论收敛，还须把光滑函数 \(T*\eta\) 再与测试函数配对。记反射核 \(\check\eta(z)=\eta(-z)\)。反射不可由卷积记号中擅自省去；当 \(\eta=\rho_\varepsilon\) 时，偶性才使它等于原核。

\begin{proposition}\label{b1:prop:distribution-convolution-pairing}
若 \(\varphi\in\mathcal D(\Omega_a)\)，将其在外部补零，则
\[
 \langle T*\eta,\varphi\rangle
 =\langle T,\check\eta*\varphi\rangle.
\]
右端的卷积属于 \(\mathcal D(\Omega)\)。
\end{proposition}
\begin{proof}
令 \(K=\operatorname{supp}\varphi\)，则
\[
 q(y)=\int_{\mathbb R^d}\varphi(x)\cdot\eta(x-y)\,\mathrm dx
 =(\check\eta*\varphi)(y)
\]
光滑，且支集包含在 \(H=K+\overline{B(0,a)}\Subset\Omega\) 中。以下证明 \(T\) 与这个积分确能交换。

取包含 \(K\) 的有界闭矩形 \(Q\)，并将其分成网格趋于零的小矩形 \(Q_i\)，各选一点 \(x_i\)。函数值 Riemann 和
\[
 q_N(y)=\sum_i m_d(Q_i)\cdot\varphi(x_i)\cdot\eta(x_i-y)
\]
都属于 \(\mathcal D_H\)。对任意固定阶数 \(j\)，\(\varphi(x)D_y^\beta\eta(x-y)\) 在共同紧集上一致连续，\(|\beta|\le j\)，所以 Riemann 和误差关于 \(y\) 一致趋零，即 \(p_{H,j}(q_N-q)\to0\)。有限阶估计因此给出 \(\langle T,q_N\rangle\to\langle T,q\rangle\)。另一方面，由线性性，每个配对都是连续紧支集函数 \(x\mapsto\varphi(x)(T*\eta)(x)\) 的 Riemann 和；在 \(\Omega_a\) 外该乘积补零。它们趋于 \(\int\varphi(x)(T*\eta)(x)\,\mathrm dx\)，等式得证。
\end{proof}

\begin{theorem}[分布的局部光滑化]\label{b1:thm:distribution-regularization}
设 \(T\in\mathcal D'(\Omega)\)，令 \(T_\varepsilon=T*\rho_\varepsilon\in C^\infty(\Omega_\varepsilon)\)。则对每个固定的 \(\varphi\in\mathcal D(\Omega)\)，当 \(\varepsilon\) 充分小时配对有定义，且
\[
 \langle T_\varepsilon,\varphi\rangle\longrightarrow\langle T,\varphi\rangle.
\]
每个 \(D^\alpha T_\varepsilon\) 也以同样的局部意义趋于 \(D^\alpha T\)。
\end{theorem}
\begin{proof}
设 \(K=\operatorname{supp}\varphi\)，取 \(a>0\) 使 \(H=K+\overline{B(0,a)}\Subset\Omega\)。当 \(0<\varepsilon<a\) 时，\(\check\rho_\varepsilon*\varphi\) 和 \(\varphi\) 的支集均在 \(H\) 中，而且
\[
 D^\beta(\check\rho_\varepsilon*\varphi-\varphi)(y)
 =\int\rho(z)\cdot\bigl(D^\beta\varphi(y+\varepsilon z)-D^\beta\varphi(y)\bigr)\,\mathrm dz.
\]
因 \(D^\beta\varphi\) 一致连续、\(\rho\) 只在单位球中有质量，每个阶数的上确界范数都趋于零。前一命题与 \(T\) 的连续性给出所求配对极限。再将 \(T\) 换成 \(D^\alpha T\)，并用卷积微分公式，即得导数结论。
\end{proof}

这里 \(T_\varepsilon\) 的定义域随 \(\varepsilon\) 改变；“局部”指每个固定测试函数最终都能使用，不是把 \(T_\varepsilon\) 未经说明地当成整个 \(\Omega\) 上的函数。也不要求 \(T_\varepsilon\) 本身有逐点极限：例如 \(\delta_0*\rho_\varepsilon=\rho_\varepsilon\)，其峰值通常随 \(\varepsilon^{-d}\) 增长。

\subsection{逼近恒等}
\label{b1:subsec:200304}

\begin{proposition}\label{b1:prop:local-smoothing-approximation}
设 \(K\Subset\Omega\)。若 \(f\in L^p_{\mathrm{loc}}(\Omega)\)、\(1\le p<\infty\)，则
\[
 \|f*\rho_\varepsilon-f\|_{L^p(K)}\longrightarrow0.
\]
若 \(f\in C^k(\Omega)\)，则对所有 \(|\alpha|\le k\)，\(D^\alpha(f*\rho_\varepsilon)\to D^\alpha f\) 在 \(K\) 上一致成立。
\end{proposition}
\begin{proof}
由\secref{b1:sec:2001}的光滑截断，取 \(\chi\in C_c^\infty(\Omega)\)，使它在 \(K\) 的一个邻域上等于 \(1\)。将 \(\chi f\) 补零到全空间；当 \(\varepsilon\) 充分小时，它与 \(f\) 的卷积在 \(K\) 上相同。有限 \(p\) 的结论因此直接归结为\cref{b1:prop:approximate-identity-lp}。

若 \(f\in C^k\)，则 \(\chi f\) 的各阶导数都是连续紧支集函数。对每个 \(|\alpha|\le k\)，卷积微分公式把所需差转为 \((D^\alpha(\chi f))*\rho_\varepsilon-D^\alpha(\chi f)\)，再用同一命题的一致收敛部分即可。
\end{proof}

不能把有限 \(p\) 的结论照搬到 \(p=\infty\)。例如 \(f=\mathbf1_{(0,\infty)}\) 在一维原点附近有跳跃，偶核卷积在原点取值 \(1/2\)，并在附近连续，故其与 \(f\) 的本质上确界距离至少为 \(1/2\)。局部分布收敛、有限 \(p\) 范数收敛和光滑函数的逐阶一致收敛是不同层次，后续使用时须说明所需的是哪一种。

卷积平均怎样消除跳跃并保持局部信息，可见\cref{b1:fig:visual-f15}。
% F15：本书原创矢量示意；依据所在节的定义、证明或明确模型。
\begin{figure}[htbp]
\centering
\begin{tikzpicture}[x=1cm,y=1cm,>=stealth,every node/.style={font=\small},line width=.65pt]
\node[] at (6,3.4) {单位积分光滑核的卷积逼近};
\draw[->,gray] (0.88,0)--(11.8,0) node[right] {\(x\)};\draw[->,gray] (1,-0.12)--(1,2.7);
\draw[gray,dashed] plot coordinates {(3.0000,0.0000) (3.0000,2.2000) (7.0000,2.2000) (7.0000,0.0000)};
\draw[AnalysisBlue,thick] plot coordinates {(0.8000,0.0000) (0.8384,0.0000) (0.8767,0.0000) (0.9151,0.0000) (0.9534,0.0000) (0.9918,0.0000) (1.0301,0.0000) (1.0685,0.0000) (1.1068,0.0000) (1.1452,0.0000) (1.1836,0.0000) (1.2219,0.0000) (1.2603,0.0000) (1.2986,0.0000) (1.3370,0.0000) (1.3753,0.0000) (1.4137,0.0000) (1.4521,0.0000) (1.4904,0.0000) (1.5288,0.0001) (1.5671,0.0005) (1.6055,0.0018) (1.6438,0.0044) (1.6822,0.0087) (1.7205,0.0150) (1.7589,0.0236) (1.7973,0.0344) (1.8356,0.0474) (1.8740,0.0627) (1.9123,0.0801) (1.9507,0.0997) (1.9890,0.1212) (2.0274,0.1446) (2.0658,0.1698) (2.1041,0.1967) (2.1425,0.2251) (2.1808,0.2550) (2.2192,0.2863) (2.2575,0.3189) (2.2959,0.3527) (2.3342,0.3877) (2.3726,0.4236) (2.4110,0.4605) (2.4493,0.4983) (2.4877,0.5369) (2.5260,0.5762) (2.5644,0.6163) (2.6027,0.6569) (2.6411,0.6981) (2.6795,0.7398) (2.7178,0.7819) (2.7562,0.8244) (2.7945,0.8672) (2.8329,0.9103) (2.8712,0.9536) (2.9096,0.9971) (2.9479,1.0407) (2.9863,1.0844) (3.0247,1.1281) (3.0630,1.1718) (3.1014,1.2153) (3.1397,1.2588) (3.1781,1.3020) (3.2164,1.3451) (3.2548,1.3878) (3.2932,1.4302) (3.3315,1.4722) (3.3699,1.5138) (3.4082,1.5548) (3.4466,1.5952) (3.4849,1.6351) (3.5233,1.6742) (3.5616,1.7126) (3.6000,1.7501) (3.6384,1.7868) (3.6767,1.8224) (3.7151,1.8570) (3.7534,1.8905) (3.7918,1.9227) (3.8301,1.9537) (3.8685,1.9832) (3.9068,2.0112) (3.9452,2.0376) (3.9836,2.0623) (4.0219,2.0852) (4.0603,2.1061) (4.0986,2.1251) (4.1370,2.1419) (4.1753,2.1566) (4.2137,2.1690) (4.2521,2.1791) (4.2904,2.1870) (4.3288,2.1927) (4.3671,2.1965) (4.4055,2.1987) (4.4438,2.1997) (4.4822,2.2000) (4.5205,2.2000) (4.5589,2.2000) (4.5973,2.2000) (4.6356,2.2000) (4.6740,2.2000) (4.7123,2.2000) (4.7507,2.2000) (4.7890,2.2000) (4.8274,2.2000) (4.8658,2.2000) (4.9041,2.2000) (4.9425,2.2000) (4.9808,2.2000) (5.0192,2.2000) (5.0575,2.2000) (5.0959,2.2000) (5.1342,2.2000) (5.1726,2.2000) (5.2110,2.2000) (5.2493,2.2000) (5.2877,2.2000) (5.3260,2.2000) (5.3644,2.2000) (5.4027,2.2000) (5.4411,2.2000) (5.4795,2.2000) (5.5178,2.2000) (5.5562,2.1997) (5.5945,2.1987) (5.6329,2.1965) (5.6712,2.1927) (5.7096,2.1870) (5.7479,2.1791) (5.7863,2.1690) (5.8247,2.1566) (5.8630,2.1419) (5.9014,2.1251) (5.9397,2.1061) (5.9781,2.0852) (6.0164,2.0623) (6.0548,2.0376) (6.0932,2.0112) (6.1315,1.9832) (6.1699,1.9537) (6.2082,1.9227) (6.2466,1.8905) (6.2849,1.8570) (6.3233,1.8224) (6.3616,1.7868) (6.4000,1.7501) (6.4384,1.7126) (6.4767,1.6742) (6.5151,1.6351) (6.5534,1.5952) (6.5918,1.5548) (6.6301,1.5138) (6.6685,1.4722) (6.7068,1.4302) (6.7452,1.3878) (6.7836,1.3451) (6.8219,1.3020) (6.8603,1.2588) (6.8986,1.2153) (6.9370,1.1718) (6.9753,1.1281) (7.0137,1.0844) (7.0521,1.0407) (7.0904,0.9971) (7.1288,0.9536) (7.1671,0.9103) (7.2055,0.8672) (7.2438,0.8244) (7.2822,0.7819) (7.3205,0.7398) (7.3589,0.6981) (7.3973,0.6569) (7.4356,0.6163) (7.4740,0.5762) (7.5123,0.5369) (7.5507,0.4983) (7.5890,0.4605) (7.6274,0.4236) (7.6658,0.3877) (7.7041,0.3527) (7.7425,0.3189) (7.7808,0.2863) (7.8192,0.2550) (7.8575,0.2251) (7.8959,0.1967) (7.9342,0.1698) (7.9726,0.1446) (8.0110,0.1212) (8.0493,0.0997) (8.0877,0.0801) (8.1260,0.0627) (8.1644,0.0474) (8.2027,0.0344) (8.2411,0.0236) (8.2795,0.0150) (8.3178,0.0087) (8.3562,0.0044) (8.3945,0.0018) (8.4329,0.0005) (8.4712,0.0001) (8.5096,0.0000) (8.5479,0.0000) (8.5863,-0.0000) (8.6247,0.0000) (8.6630,0.0000) (8.7014,0.0000) (8.7397,0.0000) (8.7781,0.0000) (8.8164,0.0000) (8.8548,0.0000) (8.8932,0.0000) (8.9315,0.0000) (8.9699,0.0000) (9.0082,0.0000) (9.0466,0.0000) (9.0849,0.0000) (9.1233,0.0000) (9.1616,0.0000) (9.2000,0.0000)};
\node[AnalysisBlue] at (10,2.4) {\(\varepsilon=0.8\)};
\draw[orange!80!black,thick] plot coordinates {(0.8000,0.0000) (0.8384,0.0000) (0.8767,0.0000) (0.9151,0.0000) (0.9534,0.0000) (0.9918,0.0000) (1.0301,0.0000) (1.0685,0.0000) (1.1068,0.0000) (1.1452,0.0000) (1.1836,0.0000) (1.2219,0.0000) (1.2603,0.0000) (1.2986,0.0000) (1.3370,0.0000) (1.3753,0.0000) (1.4137,0.0000) (1.4521,0.0000) (1.4904,0.0000) (1.5288,0.0000) (1.5671,0.0000) (1.6055,0.0000) (1.6438,0.0000) (1.6822,0.0000) (1.7205,0.0000) (1.7589,0.0000) (1.7973,0.0000) (1.8356,0.0000) (1.8740,0.0000) (1.9123,0.0000) (1.9507,0.0000) (1.9890,0.0000) (2.0274,0.0000) (2.0658,0.0000) (2.1041,0.0000) (2.1425,0.0000) (2.1808,0.0000) (2.2192,0.0000) (2.2575,0.0000) (2.2959,0.0012) (2.3342,0.0069) (2.3726,0.0203) (2.4110,0.0425) (2.4493,0.0737) (2.4877,0.1133) (2.5260,0.1606) (2.5644,0.2148) (2.6027,0.2750) (2.6411,0.3405) (2.6795,0.4106) (2.7178,0.4847) (2.7562,0.5621) (2.7945,0.6423) (2.8329,0.7248) (2.8712,0.8091) (2.9096,0.8949) (2.9479,0.9816) (2.9863,1.0688) (3.0247,1.1562) (3.0630,1.2433) (3.1014,1.3297) (3.1397,1.4151) (3.1781,1.4990) (3.2164,1.5809) (3.2548,1.6603) (3.2932,1.7368) (3.3315,1.8098) (3.3699,1.8787) (3.4082,1.9428) (3.4466,2.0014) (3.4849,2.0537) (3.5233,2.0989) (3.5616,2.1361) (3.6000,2.1648) (3.6384,2.1844) (3.6767,2.1954) (3.7151,2.1994) (3.7534,2.2000) (3.7918,2.2000) (3.8301,2.2000) (3.8685,2.2000) (3.9068,2.2000) (3.9452,2.2000) (3.9836,2.2000) (4.0219,2.2000) (4.0603,2.2000) (4.0986,2.2000) (4.1370,2.2000) (4.1753,2.2000) (4.2137,2.2000) (4.2521,2.2000) (4.2904,2.2000) (4.3288,2.2000) (4.3671,2.2000) (4.4055,2.2000) (4.4438,2.2000) (4.4822,2.2000) (4.5205,2.2000) (4.5589,2.2000) (4.5973,2.2000) (4.6356,2.2000) (4.6740,2.2000) (4.7123,2.2000) (4.7507,2.2000) (4.7890,2.2000) (4.8274,2.2000) (4.8658,2.2000) (4.9041,2.2000) (4.9425,2.2000) (4.9808,2.2000) (5.0192,2.2000) (5.0575,2.2000) (5.0959,2.2000) (5.1342,2.2000) (5.1726,2.2000) (5.2110,2.2000) (5.2493,2.2000) (5.2877,2.2000) (5.3260,2.2000) (5.3644,2.2000) (5.4027,2.2000) (5.4411,2.2000) (5.4795,2.2000) (5.5178,2.2000) (5.5562,2.2000) (5.5945,2.2000) (5.6329,2.2000) (5.6712,2.2000) (5.7096,2.2000) (5.7479,2.2000) (5.7863,2.2000) (5.8247,2.2000) (5.8630,2.2000) (5.9014,2.2000) (5.9397,2.2000) (5.9781,2.2000) (6.0164,2.2000) (6.0548,2.2000) (6.0932,2.2000) (6.1315,2.2000) (6.1699,2.2000) (6.2082,2.2000) (6.2466,2.2000) (6.2849,2.1994) (6.3233,2.1954) (6.3616,2.1844) (6.4000,2.1648) (6.4384,2.1361) (6.4767,2.0989) (6.5151,2.0537) (6.5534,2.0014) (6.5918,1.9428) (6.6301,1.8787) (6.6685,1.8098) (6.7068,1.7368) (6.7452,1.6603) (6.7836,1.5809) (6.8219,1.4990) (6.8603,1.4151) (6.8986,1.3297) (6.9370,1.2433) (6.9753,1.1562) (7.0137,1.0688) (7.0521,0.9816) (7.0904,0.8949) (7.1288,0.8091) (7.1671,0.7248) (7.2055,0.6423) (7.2438,0.5621) (7.2822,0.4847) (7.3205,0.4106) (7.3589,0.3405) (7.3973,0.2750) (7.4356,0.2148) (7.4740,0.1606) (7.5123,0.1133) (7.5507,0.0737) (7.5890,0.0425) (7.6274,0.0203) (7.6658,0.0069) (7.7041,0.0012) (7.7425,0.0000) (7.7808,0.0000) (7.8192,0.0000) (7.8575,0.0000) (7.8959,0.0000) (7.9342,0.0000) (7.9726,0.0000) (8.0110,0.0000) (8.0493,0.0000) (8.0877,0.0000) (8.1260,0.0000) (8.1644,0.0000) (8.2027,0.0000) (8.2411,0.0000) (8.2795,0.0000) (8.3178,0.0000) (8.3562,0.0000) (8.3945,0.0000) (8.4329,0.0000) (8.4712,0.0000) (8.5096,0.0000) (8.5479,0.0000) (8.5863,0.0000) (8.6247,0.0000) (8.6630,0.0000) (8.7014,0.0000) (8.7397,0.0000) (8.7781,0.0000) (8.8164,0.0000) (8.8548,0.0000) (8.8932,0.0000) (8.9315,0.0000) (8.9699,0.0000) (9.0082,0.0000) (9.0466,0.0000) (9.0849,0.0000) (9.1233,0.0000) (9.1616,0.0000) (9.2000,0.0000)};
\node[orange!80!black] at (10,1.7999999999999998) {\(\varepsilon=0.4\)};
\draw[black,thick] plot coordinates {(0.8000,0.0000) (0.8384,0.0000) (0.8767,0.0000) (0.9151,0.0000) (0.9534,0.0000) (0.9918,0.0000) (1.0301,0.0000) (1.0685,0.0000) (1.1068,0.0000) (1.1452,0.0000) (1.1836,0.0000) (1.2219,0.0000) (1.2603,0.0000) (1.2986,0.0000) (1.3370,0.0000) (1.3753,0.0000) (1.4137,0.0000) (1.4521,0.0000) (1.4904,0.0000) (1.5288,0.0000) (1.5671,0.0000) (1.6055,0.0000) (1.6438,0.0000) (1.6822,0.0000) (1.7205,0.0000) (1.7589,0.0000) (1.7973,0.0000) (1.8356,0.0000) (1.8740,0.0000) (1.9123,0.0000) (1.9507,0.0000) (1.9890,0.0000) (2.0274,0.0000) (2.0658,0.0000) (2.1041,0.0000) (2.1425,0.0000) (2.1808,0.0000) (2.2192,0.0000) (2.2575,0.0000) (2.2959,0.0000) (2.3342,0.0000) (2.3726,0.0000) (2.4110,0.0000) (2.4493,0.0000) (2.4877,0.0000) (2.5260,0.0000) (2.5644,0.0000) (2.6027,0.0000) (2.6411,0.0000) (2.6795,0.0000) (2.7178,0.0000) (2.7562,0.0113) (2.7945,0.0762) (2.8329,0.2000) (2.8712,0.3684) (2.9096,0.5677) (2.9479,0.7869) (2.9863,1.0168) (3.0247,1.2495) (3.0630,1.4772) (3.1014,1.6917) (3.1397,1.8834) (3.1781,2.0406) (3.2164,2.1488) (3.2548,2.1959) (3.2932,2.2000) (3.3315,2.2000) (3.3699,2.2000) (3.4082,2.2000) (3.4466,2.2000) (3.4849,2.2000) (3.5233,2.2000) (3.5616,2.2000) (3.6000,2.2000) (3.6384,2.2000) (3.6767,2.2000) (3.7151,2.2000) (3.7534,2.2000) (3.7918,2.2000) (3.8301,2.2000) (3.8685,2.2000) (3.9068,2.2000) (3.9452,2.2000) (3.9836,2.2000) (4.0219,2.2000) (4.0603,2.2000) (4.0986,2.2000) (4.1370,2.2000) (4.1753,2.2000) (4.2137,2.2000) (4.2521,2.2000) (4.2904,2.2000) (4.3288,2.2000) (4.3671,2.2000) (4.4055,2.2000) (4.4438,2.2000) (4.4822,2.2000) (4.5205,2.2000) (4.5589,2.2000) (4.5973,2.2000) (4.6356,2.2000) (4.6740,2.2000) (4.7123,2.2000) (4.7507,2.2000) (4.7890,2.2000) (4.8274,2.2000) (4.8658,2.2000) (4.9041,2.2000) (4.9425,2.2000) (4.9808,2.2000) (5.0192,2.2000) (5.0575,2.2000) (5.0959,2.2000) (5.1342,2.2000) (5.1726,2.2000) (5.2110,2.2000) (5.2493,2.2000) (5.2877,2.2000) (5.3260,2.2000) (5.3644,2.2000) (5.4027,2.2000) (5.4411,2.2000) (5.4795,2.2000) (5.5178,2.2000) (5.5562,2.2000) (5.5945,2.2000) (5.6329,2.2000) (5.6712,2.2000) (5.7096,2.2000) (5.7479,2.2000) (5.7863,2.2000) (5.8247,2.2000) (5.8630,2.2000) (5.9014,2.2000) (5.9397,2.2000) (5.9781,2.2000) (6.0164,2.2000) (6.0548,2.2000) (6.0932,2.2000) (6.1315,2.2000) (6.1699,2.2000) (6.2082,2.2000) (6.2466,2.2000) (6.2849,2.2000) (6.3233,2.2000) (6.3616,2.2000) (6.4000,2.2000) (6.4384,2.2000) (6.4767,2.2000) (6.5151,2.2000) (6.5534,2.2000) (6.5918,2.2000) (6.6301,2.2000) (6.6685,2.2000) (6.7068,2.2000) (6.7452,2.1959) (6.7836,2.1488) (6.8219,2.0406) (6.8603,1.8834) (6.8986,1.6917) (6.9370,1.4772) (6.9753,1.2495) (7.0137,1.0168) (7.0521,0.7869) (7.0904,0.5677) (7.1288,0.3684) (7.1671,0.2000) (7.2055,0.0762) (7.2438,0.0113) (7.2822,0.0000) (7.3205,0.0000) (7.3589,0.0000) (7.3973,0.0000) (7.4356,0.0000) (7.4740,0.0000) (7.5123,0.0000) (7.5507,0.0000) (7.5890,0.0000) (7.6274,0.0000) (7.6658,0.0000) (7.7041,0.0000) (7.7425,0.0000) (7.7808,0.0000) (7.8192,0.0000) (7.8575,0.0000) (7.8959,0.0000) (7.9342,0.0000) (7.9726,0.0000) (8.0110,0.0000) (8.0493,0.0000) (8.0877,0.0000) (8.1260,0.0000) (8.1644,0.0000) (8.2027,0.0000) (8.2411,0.0000) (8.2795,0.0000) (8.3178,0.0000) (8.3562,0.0000) (8.3945,0.0000) (8.4329,0.0000) (8.4712,0.0000) (8.5096,0.0000) (8.5479,0.0000) (8.5863,0.0000) (8.6247,0.0000) (8.6630,0.0000) (8.7014,0.0000) (8.7397,0.0000) (8.7781,0.0000) (8.8164,0.0000) (8.8548,0.0000) (8.8932,0.0000) (8.9315,0.0000) (8.9699,0.0000) (9.0082,0.0000) (9.0466,0.0000) (9.0849,0.0000) (9.1233,0.0000) (9.1616,0.0000) (9.2000,0.0000)};
\node[black] at (10,1.2) {\(\varepsilon=0.15\)};

\end{tikzpicture}
\caption[卷积如何移动、平均与光滑化]{卷积如何移动、平均与光滑化。虚线为区间示性函数；实线为它与尺度化光滑核的卷积。核取 \(C\exp(-1/(1-x^2))\) 于 \(|x|<1\)，区间外为零，其中常数使积分为一。}
\label{b1:fig:visual-f15}
\end{figure}


% !TeX root = ../../Book1.tex
\section{弱导数}
\label{b1:sec:2004}

分布总可以求导，但导数未必还能由函数表示。弱导数研究的正是后一个问题：给定局部可积函数，把它作为正则分布求导后，结果是否仍是一个局部可积函数？这个要求比几乎处处存在经典导数更强，因为分部积分还会检测跳跃以及集中在零测集上的变化。

% 实读：Hunter2001Applied，作者官方ch12.pdf，§12.9印刷362--364页，
% 定义12.63--12.68；本节主要对照12.64的多指标弱导数与L1loc范围。
% 原PDF文本层乱码，实际读作者公开原页图像，不声称这些页含本节全部证明。
% 唯一性承接2002正则嵌入；局部Lipschitz证明用16章Rademacher及差商平移；
% 一维代表用零积分测试函数原函数、Fubini及1604绝对连续微积分基本定理自足证明。

\subsection{定义}
\label{b1:subsec:200401}

设 \(\Omega\subseteq\mathbb R^d\) 为开集，函数可取实值或复值。沿用多指标记号 \(\alpha=(\alpha_1,\ldots,\alpha_d)\in\mathbb N_0^d\)，其中 \(|\alpha|=\alpha_1+\cdots+\alpha_d\)，\(\partial^\alpha=\partial_1^{\alpha_1}\cdots\partial_d^{\alpha_d}\) 在光滑函数上表示通常的偏导数。

\begin{definition}
\label{b1:def:weak-derivative}
给定 \(u,g\in L^1_{\mathrm{loc}}(\Omega)\) 和多指标 \(\alpha\)，若对每个 \(\varphi\in C_c^\infty(\Omega)\) 都有
\begin{equation}
\label{b1:eq:weak-derivative-test}
 \int_\Omega g\varphi\,\mathrm dx
 =(-1)^{|\alpha|}\int_\Omega u\,\partial^\alpha\varphi\,\mathrm dx,
\end{equation}
则称 \(g\) 为 \(u\) 的 \(\alpha\) 阶\term[ruo-daoshu]{弱导数}[weak derivative]，记为 \(\partial^\alpha u=g\)。当 \(\alpha=e_i\) 时，记为 \(\partial_i u\)。若全部一阶弱偏导数存在，便把它们合记为 \(\nabla u\coloneqq(\partial_1u,\ldots,\partial_du)\)。
\end{definition}
\symidx{\(\partial^\alpha u\)：局部可积函数的多指标弱导数}

所有积分只涉及测试函数的紧支集，所以局部可积性已经足够，不要求 \(u\) 或 \(g\) 在整个 \(\Omega\) 上可积。定义不含复共轭，与前面的复线性分布配对一致。当 \(|\alpha|=0\) 时，它只是函数本身；正阶时，定义并没有先要求经典差商存在。

用正则分布记号，\eqref{b1:eq:weak-derivative-test}就是 \(\partial^\alpha T_u=T_g\)。因此，本书把“\(u\) 有弱导数”限定为其分布导数能够由 \(L^1_{\mathrm{loc}}\) 函数表示；不能因为分布导数总存在，就认为弱导数也总存在。Hunter 与 Nachtergaele 的《Applied Analysis》\cite[Section 12.9, Definition 12.64]{Hunter2001Applied}采用这一局部可积函数的定义，并由此引入下一章的函数空间。

\subsection{唯一性}
\label{b1:subsec:200402}

\begin{proposition}
\label{b1:prop:weak-derivative-unique}
一个给定阶数的弱导数若存在，则在几乎处处相等的意义下唯一。把 \(u\) 或其弱导数在零测集上改值，不改变弱导数关系。
\end{proposition}
\begin{proof}
若 \(g_1,g_2\) 都满足定义，则 \(\int_\Omega(g_1-g_2)\varphi=0\) 对所有测试函数成立。由\cref{b1:prop:regular-distribution-injective}，\(g_1=g_2\) 几乎处处。零测集上的改值不影响定义中的积分，因此也不影响这个等价类。
\end{proof}

唯一性并不指定每一点的导数值。例如 \(u(x)=|x|\) 的一阶弱导数将在下文确定为 \(\operatorname{sgn}x\)，它在原点取哪个有限值都无关紧要。反过来，若要求导数具有额外的连续代表或指定的点值，就需要另外证明这种代表存在。

\begin{proposition}
\label{b1:prop:weak-derivative-locality}
设 \(u,g\in L^1_{\mathrm{loc}}(\Omega)\)，且 \(\{U_j\}\) 是 \(\Omega\) 的一个开覆盖。则 \(\partial^\alpha u=g\) 在 \(\Omega\) 上成立，当且仅当它在每个 \(U_j\) 上成立。
\end{proposition}
\begin{proof}
整体关系限制到开子集时仍成立，因为子集内的测试函数补零后仍是 \(\Omega\) 中的测试函数。反过来，分布 \(S=\partial^\alpha T_u-T_g\) 在每个 \(U_j\) 上为零。由\cref{b1:prop:distribution-locality}，\(S=0\) 在整个 \(\Omega\) 上成立，正是所需关系。
\end{proof}

弱导数关系还便于取极限。若 \(u_j\rightarrow u\)、\(g_j\rightarrow g\) 于 \(L^1_{\mathrm{loc}}(\Omega)\)，且 \(\partial^\alpha u_j=g_j\)，则 \(\partial^\alpha u=g\)。事实上，对固定测试函数，\(\varphi\) 与 \(\partial^\alpha\varphi\) 在同一个紧集上有界；两个积分的误差分别受相应的局部 \(L^1\) 误差乘这个界控制，因而可在\eqref{b1:eq:weak-derivative-test}中直接取极限。

这里需要同时控制函数和导数，而不是从 \(u_j\) 的收敛自动推断导数收敛。例如 \(u_j(x)=j^{-1}\sin(j^2x)\) 一致趋于零，但其经典导数 \(j\cos(j^2x)\) 在原点甚至趋于无穷。分布导数仍由前节的连续性趋于零；要求导数在某个函数空间中收敛，则是另外一个有实质内容的条件。

\subsection{分部积分}
\label{b1:subsec:200403}

\begin{proposition}[Leibniz]
\label{b1:prop:weak-leibniz}
若 \(u\) 的全部 \(|\beta|\leq r\) 阶弱导数局部可积，且 \(a\in C^\infty(\Omega)\)，则对 \(|\alpha|\leq r\)，
\[
 \partial^\alpha(au)
 =\sum_{\beta\leq\alpha}\binom\alpha\beta
    (\partial^{\alpha-\beta}a)(\partial^\beta u),
 \quad
 \binom\alpha\beta=\prod_{i=1}^d\binom{\alpha_i}{\beta_i}.
\]
其中 \(\beta\leq\alpha\) 指逐分量不等式，等式按弱导数解释。
\end{proposition}
\begin{proof}
先设 \(r=1\)，记 \(g_i=\partial_i u\)。将 \(a\varphi\) 代入定义，得到
\[
 \begin{aligned}
 -\int_\Omega au\,\partial_i\varphi
 &=-\int_\Omega u\,\partial_i(a\varphi)
    +\int_\Omega u(\partial_i a)\varphi\\
 &=\int_\Omega\bigl(ag_i+u\partial_i a\bigr)\varphi.
 \end{aligned}
\]
右侧函数局部可积，故这就证明一阶公式。高阶时，分布导数可交换，所以已存在的 \(\partial^\beta u\) 的弱偏导数为 \(\partial^{\beta+e_i}u\)。反复应用一阶公式；对每个坐标，\(\alpha_i\) 次求导中有 \(\beta_i\) 次落在 \(u\) 上的选择数为 \(\binom{\alpha_i}{\beta_i}\)，各坐标的选择数相乘，便得到所列系数及公式。
\end{proof}

上述计算是内部的分部积分，不要求 \(\Omega\) 的边界光滑。它也允许把测试函数类扩大到 \(C_c^1(\Omega)\)：若 \(\partial_i u=g_i\in L^1_{\mathrm{loc}}\)，则
\begin{equation}
\label{b1:eq:weak-integration-by-parts}
 \int_\Omega u\,\partial_i v\,\mathrm dx
 =-\int_\Omega g_i v\,\mathrm dx
 \quad(v\in C_c^1(\Omega)).
\end{equation}
为证明这一点，把 \(v\) 在域外补零，并用\secref{b1:sec:2003}的磨光函数取 \(v_\varepsilon=v*\rho_\varepsilon\)。当 \(\varepsilon\) 足够小时，这些函数的支集包含在同一个 \(K\Subset\Omega\) 中，且 \(v_\varepsilon\rightarrow v\)、\(\partial_i v_\varepsilon\rightarrow\partial_i v\) 一致收敛。这两个一致收敛也可直接由 \(v\) 及其一阶导数的一致连续性和 \(\int\rho_\varepsilon=1\) 得到。对 \(v_\varepsilon\) 使用定义，再由 \(u,g_i\in L^1(K)\) 取极限，得到\eqref{b1:eq:weak-integration-by-parts}。

紧支集条件在这里有实际作用：测试函数在边界附近为零，所以不出现边界项。对一直延伸到边界的函数使用分部积分，需要边界值及迹的概念，不能从上式省略紧支集条件后直接得到。

与此相应，在开子集上限制弱导数是安全的，在域外补零却需要额外检查。例如常函数 \(u=1\) 在 \((0,1)\) 上的弱导数为零，但其零延拓 \(\overline u=\mathbf1_{(0,1)}\) 满足
\[
 \langle\partial T_{\overline u},\varphi\rangle
 =-\int_0^1\varphi'(x)\,\mathrm dx
 =\varphi(0)-\varphi(1).
\]
因此零延拓产生了 \(\delta_0-\delta_1\) 两个端点项，并不是把内部的零导数也补零。若先乘以 \(\chi\in C_c^\infty(\Omega)\)，则 \(\chi u\) 在边界附近已经为零，Leibniz 公式给出的导数便可以与函数一起补零。这是局部化与直接截断定义域的区别。

\subsection{经典导数与局部 Lipschitz 函数}
\label{b1:subsec:200404}

若 \(u\in C^1(\Omega)\)，则经典偏导数也是弱偏导数。事实上，\(u\varphi\) 在域外补零后属于 \(C_c^1(\mathbb R^d)\)。沿第 \(i\) 个坐标积分，由一维微积分基本定理和 Fubini 定理，\(\int\partial_i(u\varphi)=0\)，展开乘积即得定义中的等式。重复这个论证，也能处理 \(C^r\) 函数的 \(r\) 阶导数。

\begin{theorem}
\label{b1:thm:lipschitz-weak-derivative}
局部 Lipschitz 函数 \(u:\Omega\rightarrow\mathbb R\) 或 \(\mathbb C\) 的各一阶弱导数存在，属于 \(L^\infty_{\mathrm{loc}}(\Omega)\)，并几乎处处等于经典偏导数。在 \(u\) 为 \(L\)-Lipschitz 的开集上，\(|\partial_i u|\leq L\) 几乎处处。
\end{theorem}
\begin{proof}
先在 \(u\) 为 \(L\)-Lipschitz 的有界开球 \(B\Subset\Omega\) 内证明。由\cref{b1:thm:rademacher}，经典微分几乎处处存在；再由\cref{b1:prop:lipschitz-differentiability-borel}，把不可微点处的微分矩阵补零后得到 Borel 矩阵函数。令 \(g_i\) 为其第 \(i\) 个偏导分量，则 \(g_i\) 可测且 \(|g_i|\leq L\)。将 \(u|_B\) 在球外补零为 \(u_0\)。对 \(\varphi\in C_c^\infty(B)\)，平移变量给出
\[
 \int_{\mathbb R^d}\dfrac{u_0(x+he_i)-u_0(x)}h\varphi(x)\,\mathrm dx
 =\int_{\mathbb R^d}u_0(x)
   \dfrac{\varphi(x-he_i)-\varphi(x)}h\,\mathrm dx.
\]
当 \(|h|\) 足够小时，左侧非零测试点所涉及的两点都在球内，差商被 \(L\) 控制，并几乎处处趋于 \(g_i\)。右侧测试差商支集留在固定的紧子集内，被 \(\|\partial_i\varphi\|_\infty\) 控制，且趋于 \(-\partial_i\varphi\)。两边应用支配收敛定理，得 \(\int_Bg_i\varphi=-\int_Bu\partial_i\varphi\)。最后用这样的球覆盖 \(\Omega\)，由弱导数的局部性和唯一性拼合。
\end{proof}

因此 \(|x|\) 的弱导数为 \(\operatorname{sgn}x\)，尖点不妨碍一阶弱导数存在。一维情形还有一个更精确的刻画，它同时说明绝对连续性为什么会出现。

\subsection{一维弱导数与绝对连续代表}
\label{b1:subsec:200405}

\begin{theorem}[一维弱导数与绝对连续代表]
\label{b1:thm:one-dimensional-weak-ac}
设 \(I\subseteq\mathbb R\) 为非空开区间，\(u,g\in L^1_{\mathrm{loc}}(I)\)。则 \(u'=g\) 作为弱导数成立，当且仅当 \(u\) 有一个代表 \(\widetilde u\)，在每个紧子区间上绝对连续，且 \(\widetilde u'=g\) 几乎处处。这个连续代表唯一，并满足
\[
 \widetilde u(t)-\widetilde u(s)=\int_s^t g(r)\,\mathrm dr
 \quad(s,t\in I).
\]
\end{theorem}
\begin{proof}
先证实值情形。复值情形分别对 \(u,g\) 的实部和虚部应用实值结论，再把所得两个绝对连续代表组合起来；因而以下构造中的函数均取实值。

先证明：若 \(w\in L^1_{\mathrm{loc}}(I)\) 的弱导数为零，则 \(w\) 几乎处处为常数。对积分为零的 \(\psi\in C_c^\infty(I)\)，先在 \(I\) 外补零，其原函数 \(\Psi(x)=\int_{-\infty}^x\psi(t)\,\mathrm dt\) 仍属于 \(C_c^\infty(I)\)，故 \(\int_Iw\psi=\int_Iw\Psi'=0\)。取固定的 \(\eta\in C_c^\infty(I)\)，使 \(\int_I\eta=1\)。对任意测试函数 \(\varphi\)，将上述结论用于 \(\psi=\varphi-(\int_I\varphi)\eta\)，便得
\[
 \int_I w\varphi=c\int_I\varphi,\quad c=\int_Iw\eta.
\]
正则分布嵌入的单射性说明 \(w=c\) 几乎处处。

现固定 \(x_0\in I\)，定义 \(G(x)=\int_{x_0}^xg(t)\,\mathrm dt\)。由\cref{b1:thm:ac-fundamental-calculus}，\(G\) 在每个紧子区间上绝对连续，且 \(G'=g\) 几乎处处。它也有弱导数 \(g\)：给定测试函数，取包含其支集的 \([a,b]\subset I\)，写 \(G(x)=G(a)+\int_a^xg(t)\,\mathrm dt\)，由 Fubini 定理，
\[
 \int_a^bG(x)\varphi'(x)\,\mathrm dx
 =\int_a^b g(t)\left(\int_t^b\varphi'(x)\,\mathrm dx\right)\mathrm dt
 =-\int_a^b g(t)\varphi(t)\,\mathrm dt.
\]
若 \(u'=g\)，则 \(u-G\) 的弱导数为零，所以 \(u=c+G\) 几乎处处，所需代表为 \(\widetilde u=c+G\)。反向则由绝对连续函数的积分表示及同一个 Fubini 计算得到。两个连续代表若几乎处处相等，则处处相等，否则连续性会给出一个它们不相等的开子区间，故代表唯一。
\end{proof}

这个刻画给出了超出 Lipschitz 类的具体函数。若 \(0<\beta<1\)，则 \(u(x)=|x|^\beta\) 在原点附近不是 Lipschitz 函数，但
\[
 |x|^\beta=\int_0^x\beta\operatorname{sgn}(t)|t|^{\beta-1}\,\mathrm dt.
\]
被积函数局部可积，因而它就是 \(u\) 的弱导数。进一步，对 \(1\leq p<\infty\)，该导数在原点附近属于 \(L^p\)，当且仅当 \(p(1-\beta)<1\)；这可由积分 \(\int_0^r t^{-p(1-\beta)}\,\mathrm dt\) 直接判断。在临界等号处，积分按对数发散。导数存在与导数具有何种可积性，是两个不同层次的问题。

最后看两个反例。令 \(H=\mathbf1_{(0,\infty)}\)。它的经典导数在原点外为零，但
\[
 \langle\partial T_H,\varphi\rangle
 =-\int_0^\infty\varphi'(x)\,\mathrm dx
 =\varphi(0),
\]
所以分布导数是 \(\delta_0\)，而不是零。它不能由局部可积函数表示：若能，则相应函数在不含原点的开集上由正则嵌入的单射性为零，于是几乎处处为零，却无法对满足 \(\varphi(0)=1\) 的测试函数给出值 \(1\)。因此 \(H\) 没有本节意义下的一阶弱导数。

\cref{b1:fig:jump-distributional-derivative}把这一差别画成最简单的模型：经典导数在跳跃点外为零，分布导数却把全部跳跃质量集中到该点。
\begin{figure}[htbp]
  \centering
  \begin{tikzpicture}[>=stealth]
    \begin{scope}
      \draw[->] (-1.8,0)--(1.9,0);
      \draw[->] (0,-0.25)--(0,1.45);
      \draw[thick] (-1.55,0)--(0,0);
      \draw[thick] (0,1)--(1.55,1);
      \draw[dashed] (0,0)--(0,1);
      \fill[white] (0,0) circle (2pt);
      \draw (0,0) circle (2pt);
      \fill (0,1) circle (2pt);
      \node[above right] at (0,1) {$H$};
    \end{scope}
    \draw[->,thick] (2.25,0.62)--(3.65,0.62) node[midway,above] {$\partial$};
    \begin{scope}[shift={(5.6,0)}]
      \draw[->] (-1.5,0)--(1.55,0);
      \draw[->] (0,-0.25)--(0,1.45);
      \draw[->,very thick] (0,0)--(0,1.15);
      \node[above right] at (0,1.05) {$\delta_0$};
    \end{scope}
  \end{tikzpicture}
  \caption[跳跃的分布导数]{跳跃的分布导数。Heaviside 函数的变化集中在原点，故 \(\partial T_H=\delta_0\)。}
  \label{b1:fig:jump-distributional-derivative}
\end{figure}

连续性也不能排除这种问题。\cref{b1:ex:14-cantor-function}构造的 Cantor 函数 \(F_C\) 连续、非恒定，且经典导数几乎处处为零。如果它在 \((0,1)\) 上有局部可积的弱导数，由刚证的定理，便有局部绝对连续代表；该代表与连续的 \(F_C\) 必须处处相等。其弱导数于是等于经典导数，几乎处处为零，迫使 \(F_C\) 为常数，矛盾。跳跃函数和 Cantor 函数分别说明，点上的跳跃与零测集上的连续变化，都可能被经典的几乎处处导数遗漏。

% 本章小结（不计入编号）
\begin{chaptersummary}
紧支集光滑函数允许任意次分部积分，又把计算限制在开集内部。分布的连续性在每个固定紧集上化成一个有限阶估计，因此可以把测试函数上的微分、光滑乘法和局部化转移到分布上。一般分布无需有点值，也无需由可积密度给出。

光滑化使分布暂时成为普通光滑函数，并保留微分关系；其收敛首先是逐个测试函数的配对收敛，不能自动理解为一致收敛或逐点收敛。弱导数则进一步要求所得分布由局部可积函数表示。在一维，这与局部绝对连续代表相联系；跳跃产生点质量，Cantor 函数提醒我们不能只凭几乎处处的经典导数判断整个变化。下一章将对函数和弱导数同时施加 \(L^p\) 控制，建立 Sobolev 空间。
\end{chaptersummary}


\BookExercises

% \chapref{b1:ch:20}习题临时稿；不另设 BookExercises 入口。
% 十一题均为依据本章已建立工具自拟的变式，不冒认教材原题编号。
% 训练分工：测试拓扑/局部阶；乘法核/主值构造；跳跃/分布ODE；磨光误差/弱导数端点。
% 前置：20.1测试列判据，20.2有限阶与分布运算，20.3卷积，20.4弱导数；无21章依赖。

\begin{exercise}\label{b1:ex:20-basic-distribution-pairings}
在 \(\mathcal D'(\mathbb R)\) 中按定义计算下列对象对测试函数 \(\varphi\) 的作用，并辨认所得分布：
\[
 \delta_a',\qquad x\delta_0',\qquad (T_x)'.
\]
其中 \(T_x\) 表示函数 \(x\mapsto x\) 产生的正则分布。
\end{exercise}
% 来源：自拟；先以三项按定义配对训练点导数、光滑乘法与正则分布求导。
\begin{solution}
分布导数的定义直接给出
\[
 \langle\delta_a',\varphi\rangle=-\varphi'(a).
\]
再由光滑函数与分布的乘法，
\[
 \langle x\delta_0',\varphi\rangle
 =\langle\delta_0',x\varphi\rangle
 =-(x\varphi)'(0)=-\varphi(0),
\]
所以 \(x\delta_0'=-\delta_0\)。最后作一次普通分部积分，测试函数的边界项为零，故
\[
 \langle(T_x)',\varphi\rangle
 =-\int_{\mathbb R}x\varphi'(x)\,\mathrm dx
 =\int_{\mathbb R}\varphi(x)\,\mathrm dx
 =\langle T_1,\varphi\rangle.
\]
因此 \((T_x)'=T_1\)。
\end{solution}

\begin{exercise}\label{b1:ex:20-basic-weak-derivatives}
在 \(\mathbb R\) 上分别考虑 \(1\)、\(x\)、\(|x|\) 和 \(H=\mathbf1_{(0,\infty)}\)。
\begin{enumerate}
\item 求前三个函数的一阶弱导数。
\item 证明 \(H\) 的分布导数为 \(\delta_0\)，因而它没有局部可积的一阶弱导数。
\item 求 \(|x|\) 的二阶分布导数，并判断它是否存在局部可积的二阶弱导数。
\end{enumerate}
\end{exercise}
% 来源：自拟；用最基本模型区分经典不可微点、弱导数和集中型分布导数。
\begin{solution}
常函数与恒等函数的经典导数分别为 \(0\) 和 \(1\)，由经典导数与弱导数的一致性，它们的一阶弱导数也是这两个函数。函数 \(|x|\) 局部 Lipschitz，所以由\cref{b1:thm:lipschitz-weak-derivative}，
\[
 (|x|)'=\operatorname{sgn}x
 \quad\text{几乎处处}.
\]
原点处给 \(\operatorname{sgn}x\) 赋哪个有限值都不影响弱导数。

对任意测试函数，
\[
 \langle(T_H)',\varphi\rangle
 =-\int_0^\infty\varphi'(x)\,\mathrm dx
 =\varphi(0)=\langle\delta_0,\varphi\rangle.
\]
若 \(H\) 有局部可积弱导数，该函数产生的正则分布就应等于 \(\delta_0\)，与\cref{b1:prop:dirac-not-regular}矛盾。

同理，分段积分给出
\[
 \langle(\operatorname{sgn}x)',\varphi\rangle=2\varphi(0),
\]
所以 \((T_{|x|})''=2\delta_0\)。它不是正则分布，故 \(|x|\) 没有本章意义下的二阶局部可积弱导数。尖点容许一阶弱导数，却在再求一次导数时产生点质量。
\end{solution}

\begin{exercise}\label{b1:ex:20-dirac-mollification}
设 \(a\in\mathbb R^d\)，\(\rho_\varepsilon\) 为本章的磨光函数。证明
\[
 \delta_a*\rho_\varepsilon=\rho_\varepsilon(\,\cdot-a),
 \qquad
 (\partial^\alpha\delta_a)*\rho_\varepsilon
 =\partial^\alpha\rho_\varepsilon(\,\cdot-a).
\]
再直接验证这两个光滑函数产生的正则分布分别趋于 \(\delta_a\) 和 \(\partial^\alpha\delta_a\)。
\end{exercise}
% 来源：自拟；直接练习点质量的卷积、导数交换与分布收敛。
\begin{solution}
由分布卷积的定义，
\[
 (\delta_a*\rho_\varepsilon)(x)
 =\langle(\delta_a)_y,\rho_\varepsilon(x-y)\rangle
 =\rho_\varepsilon(x-a).
\]
卷积微分公式随即给出第二个等式。对任意 \(\varphi\in\mathcal D(\mathbb R^d)\)，变量代换 \(x=a+\varepsilon z\) 得
\[
 \int_{\mathbb R^d}\rho_\varepsilon(x-a)\varphi(x)\,\mathrm dx
 =\int_{\mathbb R^d}\rho(z)\varphi(a+\varepsilon z)\,\mathrm dz
 \longrightarrow\varphi(a),
\]
其中使用 \(\rho\) 的紧支集、单位积分和 \(\varphi\) 的连续性。因此相应正则分布趋于 \(\delta_a\)。对导数再作分部积分，
\[
 \int\partial^\alpha\rho_\varepsilon(x-a)\varphi(x)\,\mathrm dx
 =(-1)^{|\alpha|}\int\rho_\varepsilon(x-a)\partial^\alpha\varphi(x)\,\mathrm dx
 \longrightarrow(-1)^{|\alpha|}\partial^\alpha\varphi(a),
\]
右端正是 \(\langle\partial^\alpha\delta_a,\varphi\rangle\)。
\end{solution}

\begin{exercise}\label{b1:ex:20-test-convergence-boundaries}
\begin{enumerate}
\item\label{b1:item:20-escaping-flat-amplitude}
取 \(0\ne\psi\in C_c^\infty\bigl((-1/4,1/4)\bigr)\)，在 \(\Omega=(0,1)\) 上令
\[
 \varphi_n(x)=e^{-n^2}\psi\bigl(n^2(x-1/n)\bigr),
 \quad n\ge4.
\]
证明它的每个固定阶导数在整个 \(\Omega\) 上一致趋零，但 \(\varphi_n\) 不在 \(\mathcal D(\Omega)\) 中趋于零。
\item\label{b1:item:20-oscillatory-regular-distributions}
取 \(\chi\in\mathcal D\bigl((-2,2)\bigr)\)，使 \(\chi=1\) 于 \([-1,1]\)，令 \(u_n(x)=n^{-1}\chi(x)\sin(nx)\)。证明 \(T_{u_n}\to0\) 于 \(\mathcal D'\bigl((-2,2)\bigr)\)，但 \(u_n\not\to0\) 于测试函数空间。
\end{enumerate}
\end{exercise}
% 来源：自拟；逃逸例增加超快衰减，使所有全域导数上确界都趋零；高频例比较测试对象与正则分布两种收敛。
\begin{solution}
第一问中，\(\varphi_n\) 的支集落在 \([1/n-1/(4n^2),1/n+1/(4n^2)]\subset(0,1)\)。对每个固定整数 \(j\ge0\)，
\[
 \|\varphi_n^{(j)}\|_\infty
 =e^{-n^2}n^{2j}\|\psi^{(j)}\|_\infty\longrightarrow0.
\]
然而取 \(z_0\) 使 \(\psi(z_0)\ne0\)，则 \(x_n=1/n+z_0/n^2\) 为非零取值点，且 \(x_n\to0\)。而 \((0,1)\) 的任一紧子集都与端点保持正距离，所以不存在共同紧集包含全部支集。由\cref{b1:prop:test-sequence-convergence}，测试函数收敛不成立。振幅再小，也不能消除支集逃向边界的问题。

第二问的全部支集都在 \(\operatorname{supp}\chi\) 中。对任意固定测试函数 \(\eta\)，
\[
 |\langle T_{u_n},\eta\rangle|
 \le\frac{\|\chi\|_\infty}{n}\int_{-2}^2|\eta(x)|\,\mathrm dx\longrightarrow0.
\]
但在 \([-1,1]\) 上有 \(u_n'(x)=\cos(nx)\)，特别是 \(u_n'(0)=1\)。因此一阶导数不一致趋零，仍不满足测试列判据。分布收敛只逐个测试其积分，不能据此反推函数及其导数的测试空间收敛。
\end{solution}

\begin{optionalexercise}\label{b1:ex:20-locally-finite-unbounded-order}
在 \(\mathbb R\) 上定义线性泛函
\[
 \langle T,\varphi\rangle
 =\sum_{n=1}^\infty(-1)^n\varphi^{(n)}(n),
 \quad \varphi\in\mathcal D(\mathbb R).
\]
证明它是支集恰为 \(\{1,2,\ldots\}\) 的分布，但不存在一个固定整数 \(m\)，使其在每个紧集上都满足阶数不超过 \(m\) 的估计。这里允许估计常数随紧集改变。
\end{optionalexercise}
% 来源：自拟；将局部有限的点支集和逐渐增加的导数阶数组合，具体实现正文关于局部阶的警告。
\begin{solution}
测试函数的紧支集只包含有限个正整数，其余点处的函数及全部导数为零，所以这个和实际上有限。对固定紧集 \(K\)，若其中最大的正整数为 \(N\)，则
\[
 |\langle T,\varphi\rangle|\le Np_{K,N}(\varphi)
 \quad(\varphi\in\mathcal D_K).
\]
若 \(K\) 不含正整数，左端恒为零。局部有限阶估计证明 \(T\in\mathcal D'(\mathbb R)\)。

在不含正整数的开集上，\(T\) 为零；在正整数 \(N\) 附近，它与 \(\delta_N^{(N)}\) 相同。取支集足够小、在 \(N\) 处的 \(N\) 阶导数非零的测试函数，便知它在 \(N\) 的任一邻域不恒为零。因此支集恰为所列离散集合。

假设存在全局通用阶数 \(m\)。选整数 \(N>m\)，固定 \(K=[N-1/4,N+1/4]\)，再取 \(\psi\in C_c^\infty\bigl((-1/4,1/4)\bigr)\) 使 \(\psi^{(N)}(0)=1\)。对 \(0<\varepsilon<1\)，令
\[
 \psi_\varepsilon(x)=\varepsilon^m\psi\left(\frac{x-N}{\varepsilon}\right).
\]
其支集在 \(K\) 中，且所有 \(j\le m\) 阶导数的上确界均有独立于 \(\varepsilon\) 的界。但
\[
 |\langle T,\psi_\varepsilon\rangle|=\varepsilon^{m-N}\longrightarrow\infty,
\]
与 \(|\langle T,\varphi\rangle|\le C_Kp_{K,m}(\varphi)\) 矛盾。这不是常数在远处增大的问题，而是在同一个选定紧集上，第 \(m\) 阶控制已经不够。
\end{solution}

\begin{optionalexercise}\label{b1:ex:20-polynomial-annihilator}
设整数 \(m\ge1\)，\(T\in\mathcal D'(\mathbb R)\)。证明
\[
 x^mT=0
 \quad\Longleftrightarrow\quad
 T=\sum_{j=0}^{m-1}c_j\delta_0^{(j)}
\]
其中系数 \(c_j\in\mathbb C\) 唯一。解答须说明所用的除以 \(x^m\) 的商为何光滑且有紧支集，不预借单点支集分布的一般结构定理。并求 \(x\delta_0^{(j)}\)。
\end{optionalexercise}
% 来源：自拟；用有限阶 Taylor 分解求光滑乘法算子的核，实质扩展正文的一阶点质量恒等式。
\begin{solution}
固定 \(\chi\in\mathcal D(\mathbb R)\)，使 \(\chi=1\) 于原点附近。对任意测试函数 \(\varphi\)，令
\[
 r(x)=\varphi(x)-\chi(x)\sum_{j=0}^{m-1}\frac{\varphi^{(j)}(0)}{j!}x^j.
\]
函数 \(r\) 有紧支集，且原点处低于 \(m\) 阶的导数全部为零。在原点附近，带积分余项的 Taylor 公式给出
\[
 \frac{r(x)}{x^m}
 =\frac1{(m-1)!}\int_0^1(1-t)^{m-1}\varphi^{(m)}(tx)\,\mathrm dt.
\]
右端在 \(x=0\) 也光滑，任意阶微分可在紧区间积分下进行。因此这个商在原点光滑延拓；在原点外按通常除法定义，且因 \(r\) 紧支集而仍有紧支集。记所得测试函数为 \(\psi\)。

若 \(x^mT=0\)，则 \(\langle T,x^m\psi\rangle=0\)，所以
\[
 \langle T,\varphi\rangle
 =\sum_{j=0}^{m-1}\frac{\langle T,x^j\chi\rangle}{j!}\varphi^{(j)}(0).
\]
与 \(\langle\delta_0^{(j)},\varphi\rangle=(-1)^j\varphi^{(j)}(0)\) 比较，即得
\[
 c_j=\frac{(-1)^j}{j!}\langle T,x^j\chi\rangle.
\]
函数 \(x^j\chi/j!\) 在原点的各阶导数可分别取出相应系数，故系数唯一。

对于 \(j\ge1\)，乘积求导给出
\[
 \langle x\delta_0^{(j)},\varphi\rangle
 =(-1)^j(x\varphi)^{(j)}(0)
 =(-1)^jj\varphi^{(j-1)}(0),
\]
所以 \(x\delta_0^{(j)}=-j\delta_0^{(j-1)}\)，而 \(x\delta_0=0\)。迭代可知 \(x^m\delta_0^{(j)}=0\) 对 \(j<m\) 成立，这证明了反向蕴涵。
\end{solution}

\begin{optionalexercise}\label{b1:ex:20-principal-value-logarithm}
对 \(\varphi\in\mathcal D(\mathbb R)\)，考察
\[
 \langle P,\varphi\rangle
 =\lim_{\varepsilon\downarrow0}
 \int_{|x|>\varepsilon}\frac{\varphi(x)}x\,\mathrm dx.
\]
\begin{enumerate}
\item\label{b1:item:20-principal-value-existence}
证明极限存在并定义一个分布，且 \(P=(T_{\log|x|})'\)。
\item\label{b1:item:20-principal-value-nonregular}
证明 \(P\) 不是正则分布，\(xP=T_1\)，并求方程 \(xT=T_1\) 的全部分布解。
\end{enumerate}
\end{optionalexercise}
% 来源：自拟；对称截断构造及其边界项均在此独立证明，不预借 Hilbert 变换。
\begin{solution}
固定 \(R>0\)，使 \(\operatorname{supp}\varphi\subset(-R,R)\)。对称截断可写为
\[
 \int_{\varepsilon<|x|<R}\frac{\varphi(x)}x\,\mathrm dx
 =\int_\varepsilon^R\frac{\varphi(x)-\varphi(-x)}x\,\mathrm dx.
\]
由微积分基本定理，分子绝对值不超过 \(2x\|\varphi'\|_\infty\)，故极限存在，而且
\[
 |\langle P,\varphi\rangle|\le2R\|\varphi'\|_\infty.
\]
对每个固定紧支集可统一选 \(R\)，这证明分布的连续性。

函数 \(\log|x|\) 在零点附近可积。在 \((-R,-\varepsilon)\) 与 \((\varepsilon,R)\) 上分别分部积分，得到
\[
 -\int_{\varepsilon<|x|<R}\log|x|\,\varphi'(x)\,\mathrm dx
 =\int_{\varepsilon<|x|<R}\frac{\varphi(x)}x\,\mathrm dx
 +\log\varepsilon\bigl(\varphi(\varepsilon)-\varphi(-\varepsilon)\bigr).
\]
最后一项的绝对值至多为 \(2\varepsilon|\log\varepsilon|\cdot\|\varphi'\|_\infty\)，趋于零；左端因局部可积性趋于完整积分。这就证明 \((T_{\log|x|})'=P\)。

若 \(P=T_g\) 且 \(g\in L^1_{\mathrm{loc}}\)，则在不含原点的开区间上，测试积分给出 \(g=1/x\) 几乎处处。用可数个这样的区间覆盖 \(\mathbb R\setminus\{0\}\)，便得 \(g=1/x\) 几乎处处，与零点附近的绝对可积性矛盾。因此 \(P\) 非正则。

直接测试可得
\[
 \langle xP,\varphi\rangle
 =\lim_{\varepsilon\downarrow0}\int_{|x|>\varepsilon}\varphi(x)\,\mathrm dx
 =\langle T_1,\varphi\rangle.
\]
所以任意解与 \(P\) 的差满足 \(x(T-P)=0\)。由\cref{b1:ex:20-polynomial-annihilator}，全部解恰为 \(P+c\delta_0\)，其中 \(c\in\mathbb C\)。
\end{solution}

\begin{exercise}\label{b1:ex:20-piecewise-jump-derivatives}
设 \(a_1<\cdots<a_N\)，函数 \(f\) 在这些点分出的各开区间内为 \(C^2\) 函数，且在每个有限端点有有限的 \(f,f',f''\) 单侧连续延拓。在断点任意赋值。用 \(f'_{\mathrm{cl}}\)、\(f''_{\mathrm{cl}}\) 表示各段经典导数，并记
\[
 [f]_{a_j}=f(a_j+)-f(a_j-),
 \quad [f']_{a_j}=f'(a_j+)-f'(a_j-).
\]
证明
\[
 (T_f)'=T_{f'_{\mathrm{cl}}}+\sum_{j=1}^N[f]_{a_j}\delta_{a_j},
\]
\[
 (T_f)''=T_{f''_{\mathrm{cl}}}
 +\sum_{j=1}^N[f']_{a_j}\delta_{a_j}
 +\sum_{j=1}^N[f]_{a_j}\delta_{a_j}'.
\]
据此给出 \(f\) 存在一阶、二阶局部可积弱导数的充要条件。
\end{exercise}
% 来源：自拟；将单跳跃检测扩展为有限多界面并区分函数跳跃和一阶导数跳跃的分布阶数。
\begin{solution}
对固定测试函数，在其支集外选左右端点，再在所有 \(a_j\) 处分段。在一个相邻断点之间的区间 \((a,b)\) 上，分部积分给出
\[
 -\int_a^bf\varphi'
 =\int_a^bf'_{\mathrm{cl}}\varphi
 +f(a+)\varphi(a)-f(b-)\varphi(b).
\]
求和后，外端点项因测试函数在那里为零而消失；每个内断点的两个贡献合成 \([f]_{a_j}\varphi(a_j)\)，所以一阶公式成立。将相同计算用于分段函数 \(f'_{\mathrm{cl}}\)，并对点质量项求导，就得到二阶公式。

各段的经典导数均局部可积。有限点支集的正则分布只能为零：在这些点的补集上它为零，正则嵌入单射性迫使其代表几乎处处为零。另一方面，可在每个断点附近独立选择测试函数的值与一阶导数，从而分别取出 \(\delta_{a_j}\) 与 \(\delta_{a_j}'\) 的系数。因此一阶分布导数为正则分布，当且仅当所有 \([f]_{a_j}=0\)；二阶分布导数为正则分布，当且仅当所有 \([f]_{a_j}\) 和 \([f']_{a_j}\) 都为零。

故一阶弱导数存在等价于各段函数值能连续拼接；二阶弱导数存在等价于函数值及一阶导数都能连续拼接。断点上预先填写的单点值不改变这些结论，因为弱导数只依赖几乎处处等价类。
\end{solution}

\begin{optionalexercise}\label{b1:ex:20-distribution-first-order-ode}
设 \(I\subset\mathbb R\) 为非空开区间，\(a\in\mathbb C\)。证明 \(\mathcal D'(I)\) 中的方程
\[
 T'+aT=0
\]
的每个解都是光滑函数 \(ce^{-ax}\) 所产生的正则分布。不得预先假设 \(T\) 有函数代表。若 \(I\) 改成不连通开集，常数应如何选择？
\end{optionalexercise}
% 来源：自拟；把积分因子方法推广到一般分布，补证零导数分布为常数，而非预借正则代表。
\begin{solution}
先设 \(S\in\mathcal D'(I)\) 且 \(S'=0\)。若 \(\psi\in\mathcal D(I)\) 且 \(\int_I\psi=0\)，将其补零后令
\[
 \Psi(x)=\int_{-\infty}^x\psi(t)\,\mathrm dt.
\]
在 \(\operatorname{supp}\psi\) 左右两侧它都为零，且其支集位于 \(I\) 内的紧区间，所以 \(\Psi\in\mathcal D(I)\)。于是
\[
 \langle S,\psi\rangle=\langle S,\Psi'\rangle=-\langle S',\Psi\rangle=0.
\]
取固定 \(\eta\in\mathcal D(I)\)，使 \(\int\eta=1\)。对任意 \(\varphi\)，函数 \(\varphi-(\int\varphi)\eta\) 积分为零，故
\[
 \langle S,\varphi\rangle=c\int_I\varphi,
 \quad c=\langle S,\eta\rangle.
\]
这证明一般分布 \(S\) 已是常函数 \(c\) 的正则分布。

对原方程令 \(S=e^{ax}T\)。由分布的光滑乘法及 Leibniz 公式，
\[
 S'=e^{ax}(T'+aT)=0.
\]
所以 \(e^{ax}T=T_c\)，乘以处处非零的 \(e^{-ax}\) 得 \(T=T_{ce^{-ax}}\)。反过来，这些光滑函数的经典导数满足原方程，因而也是分布解。

若底空间为不连通开集，每个连通分支都是开区间，可在各分支上独立选择一个常数。这些解能拼成局部光滑函数：一个紧子集可被有限个包含于开集内的区间覆盖，只会遇到有限个分支，故不存在分支端点处的内部跳跃条件。
\end{solution}

\begin{optionalexercise}\label{b1:ex:20-smoothing-variable-coefficient}
设 \(\Omega\subset\mathbb R^d\) 为开集，\(U\) 为开集且 \(K\Subset U\)、\(\overline U\Subset\Omega\)，\(a\in C^\infty(\Omega)\)，\(f\in L^p_{\mathrm{loc}}(\Omega)\)，\(1\le p\le\infty\)。沿用本章非负单位积分光滑核 \(\rho_\varepsilon\)，令
\[
 C_\varepsilon(x)=a(x)(f*\rho_\varepsilon)(x)-\bigl((af)*\rho_\varepsilon\bigr)(x).
\]
证明当 \(\varepsilon\) 充分小时，
\[
 \|C_\varepsilon\|_{L^p(K)}
 \le\varepsilon\sup_U|\nabla a|\cdot\|f\|_{L^p(U)}
       \int_{\mathbb R^d}|z|\rho(z)\,\mathrm dz.
\]
最后在一维取 \(T=\delta_0\)、\(a(x)=x\)，证明相应表达式
\[
 x(T*\rho_\varepsilon)-(xT)*\rho_\varepsilon
\]
虽在分布意义下趋于零，却不在原点的任一固定邻域上依 \(L^\infty\) 范数趋于零。
\end{optionalexercise}
% 来源：自拟；量化光滑乘法与卷积不交换，有限p和无穷p只用于误差估计，不宣称一般L∞磨光范数收敛。
\begin{solution}
选择 \(\varepsilon<\operatorname{dist}(K,U^c)\)。若 \(x\in K\)、\(|z|\le\varepsilon\)，则连接 \(x\) 与 \(x-z\) 的整条线段都在 \(U\) 中。沿线段积分给出
\[
 |a(x)-a(x-z)|\le |z|\sup_U|\nabla a|.
\]
把 \(f\mathbf1_U\) 在外部补零，记作 \(f_0\)。卷积的差可写为
\[
 C_\varepsilon(x)=\int\bigl(a(x)-a(x-z)\bigr)f_0(x-z)\rho_\varepsilon(z)\,\mathrm dz.
\]
因 \(\rho\ge0\)，其绝对值受 \(\sup_U|\nabla a|\) 乘以 \(k_\varepsilon*|f_0|\) 控制，其中 \(k_\varepsilon(z)=|z|\rho_\varepsilon(z)\)。由\cref{b1:lem:convolution-lp-bound}，
\[
 \|C_\varepsilon\|_{L^p(K)}
 \le\sup_U|\nabla a|\cdot\|k_\varepsilon\|_1\cdot\|f_0\|_p.
\]
而 \(\|k_\varepsilon\|_1=\varepsilon\int|z|\rho(z)\,\mathrm dz\)，这就给出所求估计，包括 \(p=\infty\)。

对点质量，\(x\delta_0=0\)、\(\delta_0*\rho_\varepsilon=\rho_\varepsilon\)，所以差等于 \(x\rho_\varepsilon(x)\)。在一维令 \(x=\varepsilon z\)，可得
\[
 \sup_x|x\rho_\varepsilon(x)|=\sup_z|z|\rho(z)>0.
\]
核非负且积分为 \(1\)，不能只集中在单点，故右端确实为正。支集随 \(\varepsilon\downarrow0\) 缩进任一固定原点邻域，因而局部 \(L^\infty\) 收敛失败。另一方面，
\[
 \|x\rho_\varepsilon(x)\|_{L^1(\mathbb R)}
 =\varepsilon\int|z|\rho(z)\,\mathrm dz\longrightarrow0,
\]
所以它在分布意义下趋于零。函数误差的范数结论不能直接推广给没有局部可积代表的一般分布。
\end{solution}

\begin{exercise}\label{b1:ex:20-power-weak-derivative-thresholds}
对实数 \(a>-1\)，在 \((-1,1)\) 上令 \(u_a(x)=|x|^a\)，原点处任取有限值；\(a=0\) 时按几乎处处恒等于 \(1\) 解释。确定哪些 \(a\) 使它分别有一阶、二阶局部可积弱导数，并在存在时给出公式。特别说明 \(a=1\) 时二阶分布导数是什么。
\end{exercise}
% 来源：自拟；把分段分部积分扩展到幂次奇性，端点a=1单独处理，不将经典导数代入零点来猜弱导数。
\begin{solution}
条件 \(a>-1\) 保证 \(u_a\) 在零点附近可积。若 \(a>0\)，候选一阶弱导数为
\[
 g_a(x)=a\operatorname{sgn}(x)|x|^{a-1}\quad(x\ne0),
\]
它也局部可积。在原点两侧截去 \((-\varepsilon,\varepsilon)\) 后分部积分；测试函数在外端点为零，所得内边界项是
\[
 \varepsilon^a\bigl(\varphi(\varepsilon)-\varphi(-\varepsilon)\bigr),
\]
其绝对值不超过 \(2\varepsilon^{a+1}\|\varphi'\|_\infty\)，趋于零。两侧积分由绝对可积性传到极限，证明 \(u_a'=g_a\) 为弱导数。当 \(a=0\) 时弱导数为零。

若 \(-1<a<0\) 且有局部可积弱导数 \(g\)，则在原点外的各开区间上，经典与弱导数一致性及唯一性迫使 \(g=g_a\) 几乎处处。但 \(\int_0^r|g_a|=\infty\)，与 \(g\in L^1_{\mathrm{loc}}\) 矛盾。因此一阶弱导数存在的充要条件是 \(a\ge0\)。

若 \(a>1\)，对 \(g_a\) 重复两侧分部积分，内边界项为
\[
 a\varepsilon^{a-1}\bigl(\varphi(\varepsilon)+\varphi(-\varepsilon)\bigr)\longrightarrow0.
\]
其经典导数 \(a(a-1)|x|^{a-2}\) 局部可积，所以
\[
 u_a''=a(a-1)|x|^{a-2}
\]
为二阶弱导数。当 \(a=0\) 时二阶导数仍为零。

若 \(a\ne0,1\) 且 \(-1<a\le1\)，任何局部可积的二阶弱导数在原点外都必须等于 \(a(a-1)|x|^{a-2}\)，但该函数在零点附近不可积，所以不存在。剩下 \(a=1\)：一阶弱导数为 \(\operatorname{sgn}x\)，上式中的边界项此时趋于 \(2\varphi(0)\)，两侧经典二阶导数为零，故
\[
 (T_{|x|})''=2\delta_0.
\]
它不是正则分布，因此 \(|x|\) 没有二阶局部可积弱导数。综上，二阶弱导数存在当且仅当 \(a=0\) 或 \(a>1\)。一阶导数在尖点处的跳跃，恰好成为二阶导数中的点质量。
\end{solution}

% !TeX root = ../../Book1.tex
\chapter{Sobolev 空间}
\label{b1:ch:21}
\BookChapterNumber{b1:ch:21}
% 本章摘要（不计入编号）
\begin{chapterabstract}
本章把弱导数的积分恒等式与 \(L^p\) 范数结合，建立 Sobolev 空间。先用有限组导数分量证明完备性，并在相应指数范围内证明可分性与自反性，再通过局部光滑逼近推导乘积、复合、截断及坐标变换的规则。最后选取适当的函数代表，说明一维绝对连续性怎样转成多维的沿坐标线绝对连续性，以及弱梯度怎样决定几乎处处的近似微分。空间的范数结构与代表的点态性质是两条相互联系、却不能混同的主线。
\end{chapterabstract}

% 本章目标（不计入编号）
\begin{learninggoals}
\begin{enumerate}
\item\label{b1:goal:21-definition}掌握 \(W^{k,p}\)、\(H^k\) 与局部空间的定义，区分弱导数存在、全域可积性和具体代表的正则性。
\item\label{b1:goal:21-structure}把 Sobolev 空间实现为有限乘积 \(L^p\) 的闭子空间，并据此判断完备性、可分性、自反性及分量弱收敛。
\item\label{b1:goal:21-calculus}证明乘积、链式、截断和坐标变换公式，明确局部有界性、零阶可积性以及指数端点承担的作用。
\item\label{b1:goal:21-representatives}构造一个同时适用于全部坐标方向的绝对连续代表，区分经典偏导、通常微分和近似微分，并识别后者与弱梯度的关系。
\end{enumerate}
\end{learninggoals}

本章的起点是\chapref{b1:ch:20}的弱导数及局部卷积。\secref{b1:sec:2102}使用第八至十章的 Banach 空间、Hahn--Banach 定理与 \(L^p\) 对偶。\secref{b1:sec:2103}前三类局部运算沿用第二十章的卷积；坐标变换和通常可微性还调用\cref{b1:thm:rademacher}，坐标变换另需\cref{b1:thm:weighted-area-formula}。代表理论使用\secref{b1:sec:1604}的一维绝对连续微积分与\chapref{b1:ch:14}的 Lebesgue 点；近似可微性再接上\chapref{b1:ch:13}的 Hardy--Littlewood 极大函数和\secref{b1:sec:1603}的 Lipschitz 分片。

\ProseParagraphBreak

因此，定义与空间结构只依赖前两编及卷积工具，点态代表和坐标变换还使用第三编。这些调用都发生在开集内部，既不要求边界光滑，也不预先使用下一章的全域光滑逼近。

\ProseParagraphBreak

初读时宜先贯通定义、完备性、局部运算和一维代表，再细读共同直线代表的选取与近似微分的证明。后两步中的量词尤其重要：沿几乎每条坐标线具有绝对连续性，不等于对每个代表都能逐点求导；弱梯度有界与仅属于某个有限 \(L^p\) 空间，也给出不同强度的点态结论。

\ProseParagraphBreak

Hunter 与 Nachtergaele 的《Applied Analysis》\cite[Sections 12.9--12.10]{Hunter2001Applied}提供定义和后续性质的概览，其中若干性质只陈述而不证明。需要系统伴读时，可参阅 Brézis 的《Functional Analysis, Sobolev Spaces and Partial Differential Equations》\cite{Brezis2011Functional}，或 Adams 与 Fournier 的《Sobolev Spaces》\cite{Adams2003Sobolev}。局部均值比较所用的证明路线，另在相应位置注明来源。

% !TeX root = ../../Book1.tex
\section{定义}
\label{b1:sec:2101}

上一章把导数关系写成测试积分恒等式。现在给函数及其弱导数同时加上可积性要求，就能用一个范数度量它们。这种做法适合处理取极限的问题：函数可以不再光滑，但只要各阶弱导数仍有适当的积分控制，极限就有希望留在同一个空间中。

\ProseParagraphBreak

以下固定非空开集 \(\Omega\subset\mathbb R^d\)，使用 Lebesgue 测度，标量域为 \(\mathbb K=\mathbb R\) 或 \(\mathbb C\)。不预设 \(\Omega\) 有界、连通或边界光滑。弱导数始终按\cref{b1:def:weak-derivative}理解；函数及其导数均按几乎处处相等取等价类。
% 实读来源：Hunter2001Applied作者公开第12章印刷362--364页，定义12.66及其后的范数、内积。
% 此处完整规定本书多指标范数；对局部空间、边界可积性及高阶计数另作说明。

\subsection{\texorpdfstring{\(W^{1,p}\)}{W1,p}}
\label{b1:subsec:210101}

\begin{definition}\label{b1:def:sobolev-first-order}
给定 \(1\le p\le\infty\)，若 \(u\in L^p(\Omega)\)，且它的每个一阶弱偏导数 \(\partial_i u\) 均存在并属于 \(L^p(\Omega)\)，则称 \(u\) 属于一阶\term[sobolev-kongjian]{Sobolev 空间}[Sobolev space]，记为 \(u\in W^{1,p}(\Omega)\)。规定
\[
 \|u\|_{W^{1,p}(\Omega)}
 \coloneqq\left(\|u\|_{L^p(\Omega)}^p+
         \sum_{i=1}^d\|\partial_i u\|_{L^p(\Omega)}^p\right)^{1/p}
 \quad(1\le p<\infty),
\]
\[
 \|u\|_{W^{1,\infty}(\Omega)}
 \coloneqq\max\bigl\{\|u\|_{L^\infty(\Omega)},
                \|\partial_1u\|_{L^\infty(\Omega)},\ldots,
                \|\partial_du\|_{L^\infty(\Omega)}\bigr\}.
\]
\end{definition}
\symidx{\(W^{1,p}(\Omega)\)：函数及一阶弱导数属于\(L^p\)的空间}

这些弱导数首先确实可以按上一章定义。因为对每个紧集 \(K\Subset\Omega\)，Hölder 不等式给出
\[
 \int_K|u|\,\mathrm dx
 \le m_d(K)^{1-1/p}\|u\|_{L^p(\Omega)}
 \quad(1\le p<\infty),
\]
而 \(p=\infty\) 时右端换为 \(m_d(K)\|u\|_\infty\)。因此 \(L^p(\Omega)\subset L^1_{\mathrm{loc}}(\Omega)\)，即使 \(\Omega\) 的总体积无限也成立。\secref{b1:sec:2004}的唯一性又保证范数不依赖弱导数代表的选择。

\ProseParagraphBreak

定义中既控制导数，也控制函数本身。在有界开集上，非零常函数的梯度为零，函数本身却不是零等价类；因此只用 \(\|\nabla u\|_p\) 不能定义这里的范数。涉及边界条件或把常数取商的空间，需要以后另作定义。

\ProseParagraphBreak

在 \((-1,1)\) 上，\(u(x)=|x|\) 属于每个 \(W^{1,p}\)，包括 \(p=\infty\)，尽管它在原点不可微。相反，一个函数即使在开集内处处光滑，也不保证满足全域积分条件。例如 \(u(x)=\sqrt{x}\) 在 \((0,1)\) 内光滑，且函数本身有界，但
\[
 \int_0^1|u'(x)|^p\,\mathrm dx
 =2^{-p}\int_0^1x^{-p/2}\,\mathrm dx
\]
只在 \(p<2\) 时有限。因此它属于 \(W^{1,p}(0,1)\) 当且仅当 \(1\le p<2\)。这里的障碍发生在趋近边界时，不在任何内部点。

\subsection{\texorpdfstring{\(W^{k,p}\)}{Wk,p}}
\label{b1:subsec:210102}

令 \(k\in\mathbb N_0\)，记
\[
 A_k\coloneqq\{\alpha\in\mathbb N_0^d:|\alpha|\le k\},
 \qquad N_k\coloneqq\#A_k=\binom{d+k}{k}.
\]
约定 \(\partial^0u=u\)。对高阶空间，需要每一个不超过 \(k\) 阶的弱导数，而不只是某几个最高阶组合。

\begin{definition}\label{b1:def:sobolev-higher-order}
若 \(u\in L^p(\Omega)\)，且对每个 \(\alpha\in A_k\)，弱导数 \(\partial^\alpha u\) 都属于 \(L^p(\Omega)\)，则称 \(u\in W^{k,p}(\Omega)\)。其范数规定为
\begin{equation}\label{b1:eq:sobolev-norm}
 \|u\|_{W^{k,p}(\Omega)}
 \coloneqq\left(\sum_{\alpha\in A_k}
              \|\partial^\alpha u\|_{L^p(\Omega)}^p\right)^{1/p}
 \quad(p<\infty),
\end{equation}
而 \(p=\infty\) 时取这些导数的 \(L^\infty\) 范数的最大值。特别地，\(W^{0,p}(\Omega)=L^p(\Omega)\)。
\end{definition}
\symidx{\(W^{k,p}(\Omega)\)：函数及所有不超过\(k\)阶的弱导数均属于\(L^p\)的空间}

对整数 \(m\geq1\)，正式约定
\[
 W^{k,p}(\Omega;\mathbb R^m)
 \coloneqq\bigl\{u=(u^1,\ldots,u^m):u^a\in W^{k,p}(\Omega),\ 1\leq a\leq m\bigr\}.
\]
这里
\[
 \partial^\alpha u
 \coloneqq\bigl(\partial^\alpha u^1,\ldots,\partial^\alpha u^m\bigr),
\]
并固定用 \(\mathbb R^m\) 的欧氏范数计算每个 \(\|\partial^\alpha u\|_{L^p(\Omega;\mathbb R^m)}\)。向量值 Sobolev 范数仍按\eqref{b1:eq:sobolev-norm}对 \(\alpha\in A_k\) 作 \(p\) 次和；当 \(p=\infty\) 时，取这些向量值 \(L^\infty\) 范数的最大值。这一约定固定了后文定量估计中的有限维范数，而不只依赖有限维范数等价性。
\symidx{\(W^{k,p}(\Omega;\mathbb R^m)\)：各分量属于\(W^{k,p}(\Omega)\)的有限维向量值 Sobolev 空间}

\ProseParagraphBreak

若 \(|\alpha|+|\beta|\le k\)，则 \(\partial^\alpha(\partial^\beta u)=\partial^{\alpha+\beta}u\)。这不是另加一条经典可微假设：等式先由分布导数可交换得到，再由局部可积代表的唯一性成为弱导数的等式。因此也可以递归地检查 \(u\) 及其低阶导数是否具有下一阶弱导数。

\ProseParagraphBreak

同一空间常配有不同但等价的范数。例如令
\[
 S(u)=\sum_{\alpha\in A_k}\|\partial^\alpha u\|_p.
\]
对 \(p<\infty\)，有限维 Hölder 不等式给出
\[
 \|u\|_{W^{k,p}}\le S(u)
 \le N_k^{1-1/p}\|u\|_{W^{k,p}},
\]
对 \(p=\infty\) 则有 \(\|u\|_{W^{k,\infty}}\le S(u)\le N_k\|u\|_{W^{k,\infty}}\)。所以空间及收敛概念相同，但后续若要追踪精确常数，须回到本书的\eqref{b1:eq:sobolev-norm}，不能把等价范数当作数值相等。

\ProseParagraphBreak

当 \(k\ge1\) 时，记最高阶部分为
\[
 |u|_{W^{k,p}(\Omega)}
 \coloneqq\left(\sum_{|\alpha|=k}\|\partial^\alpha u\|_p^p\right)^{1/p}
 \quad(p<\infty),
\]
无穷端点仍取最大值。这一般只是半范数。比如在有界域上，次数低于 \(k\) 的多项式具有零的最高阶半范数，却未必为零函数。

\ProseParagraphBreak

提高 \(k\) 会缩小空间：\(W^{k+1,p}(\Omega)\subset W^{k,p}(\Omega)\)，且本书范数下包含映射的范数不超过 \(1\)。提高 \(p\) 则还要看域的测度。若 \(m_d(\Omega)<\infty\) 且 \(1\le p\le q\le\infty\)，逐项使用 Hölder 可得
\[
 \|u\|_{W^{k,p}(\Omega)}
 \le N_k^{1/p-1/q}\cdot m_d(\Omega)^{1/p-1/q}
       \|u\|_{W^{k,q}(\Omega)}.
\]
在无限测度域上没有这条无条件包含关系；常函数 \(1\) 在 \(\mathbb R^d\) 上属于 \(W^{k,\infty}\)，却不属于任何有限 \(p\) 的 \(W^{k,p}\)。

\subsection{\texorpdfstring{\(H^k=W^{k,2}\)}{Hk=Wk,2}}
\label{b1:subsec:210103}

\begin{definition}\label{b1:def:sobolev-hilbert}
将 \(W^{k,2}(\Omega)\) 记作 \(H^k(\Omega)\)，并定义
\begin{equation}\label{b1:eq:sobolev-inner-product}
 \langle u,v\rangle_{H^k(\Omega)}
 \coloneqq\sum_{\alpha\in A_k}\int_\Omega
       \partial^\alpha u\,\overline{\partial^\alpha v}\,\mathrm dx.
\end{equation}
实值情形省略复共轭。
\end{definition}
\symidx{\(H^k(\Omega)\)：\(W^{k,2}(\Omega)\)及其导数内积}

内积对第一变量线性，与\chapref{b1:ch:08}约定一致。零阶项保证正定性，并且 \(\langle u,u\rangle_{H^k}=\|u\|_{W^{k,2}}^2\)。因此 \(H^k\) 不仅表示另一种记号，还突出了二次能量与正交方法可以使用的结构；要称它为 Hilbert 空间，仍须在下一节证明完备性。

\ProseParagraphBreak

这里每个多重指标只计一次。例如二维 \(H^2\) 范数含有一个 \(\|\partial_1\partial_2u\|_2^2\) 项；若把 Hessian 矩阵每个位置都逐项平方求和，混合导数会计两次，得到的是另一个等价范数。两种约定都常用，本文需要数值恒等式时始终以\eqref{b1:eq:sobolev-inner-product}为准。

\ProseParagraphBreak

Hunter 与 Nachtergaele 的《Applied Analysis》\cite[Section 12.9, Definition 12.66 and the following discussion]{Hunter2001Applied}把这些空间放在弱导数之后引入。接下来我们将把导数作为 \(L^p\) 中的有限组分量，直接从\chapref{b1:ch:10}的完备性推导其基本结构。

\subsection{局部 Sobolev 空间}
\label{b1:subsec:210104}

全域可积性会同时检测函数在边界附近及无穷远处的行为。研究开集内部的运算时，往往只需要把这些位置暂时排除。

\begin{definition}\label{b1:def:sobolev-local}
若 \(u\in L^1_{\mathrm{loc}}(\Omega)\)，且每个 \(|\alpha|\le k\) 阶弱导数都存在并属于 \(L^p_{\mathrm{loc}}(\Omega)\)，则称 \(u\) 属于\term[jubu-sobolev-kongjian]{局部 Sobolev 空间}[local Sobolev space]，记作 \(W^{k,p}_{\mathrm{loc}}(\Omega)\)。这里 \(L^p_{\mathrm{loc}}\) 表示在每个紧子集上 \(p\) 次幂可积，或在 \(p=\infty\) 时本性有界。
\end{definition}
\symidx{\(W^{k,p}_{\mathrm{loc}}(\Omega)\)：局部Sobolev空间}

向量值局部空间也逐分量定义：
\[
 W^{k,p}_{\mathrm{loc}}(\Omega;\mathbb R^m)
 \coloneqq\bigl\{u=(u^1,\ldots,u^m):u^a\in W^{k,p}_{\mathrm{loc}}(\Omega),\ 1\leq a\leq m\bigr\}.
\]
因此后文的向量值链式法则不再把“逐分量理解”作为未写出的定义。

\ProseParagraphBreak

等价地，对每个开集 \(V\) 满足 \(\overline V\Subset\Omega\)，都有 \(u|_V\in W^{k,p}(V)\)。一个方向由限制立即得到。反过来，用可数个这类 \(V\) 覆盖 \(\Omega\)；重叠部分的弱导数由唯一性相等，故能拼接为同一个局部 \(L^p\) 函数，再用\cref{b1:prop:weak-derivative-locality}得到整个 \(\Omega\) 上的弱导数关系。

\ProseParagraphBreak

例如 \(x\mapsto1/x\) 属于 \((0,1)\) 上的每个 \(W^{k,p}_{\mathrm{loc}}\)，但并不属于该区间上的 \(L^1\)。另一个方向，\(W^{k,p}(\Omega)\) 总包含于 \(W^{k,p}_{\mathrm{loc}}(\Omega)\)，而局部成员只有在所有导数的全域范数也有限时才回到全域空间。

\ProseParagraphBreak

局部收敛同样按内部区域理解：\(u_j\to u\) 于 \(W^{k,p}_{\mathrm{loc}}(\Omega)\)，是指对上述每个 \(V\)，都有 \(\|u_j-u\|_{W^{k,p}(V)}\to0\)。只要选一列内部开集穷竭，使每个紧子集最终落在其中，检查这可数列就足够。局部化并不规定边界值，更不意味着把函数在域外补零后仍属于同阶 Sobolev 空间；上一章的端点点质量已经说明零延拓可能产生额外的分布导数。

% !TeX root = ../../Book1.tex
\section{基本结构}
\label{b1:sec:2102}

弱导数的定义把函数与一组积分恒等式联系起来，Sobolev 范数则同时控制函数和这些导数。两者配合，才能回答极限是否仍留在空间内。本节把所有导数放入一个有限乘积空间：范数、完备性和弱收敛都可以在这些分量上研究，而分量之间必须满足的相容关系由测试函数保存。

\ProseParagraphBreak

Hunter 与 Nachtergaele 的《Applied Analysis》\cite[Section 12.9, p. 363]{Hunter2001Applied}在 Sobolev 空间定义之后列出 Banach 与 Hilbert 结构。本节从\chapref{b1:ch:10}的 \(L^p\) 理论出发证明这些结论，并展开连续泛函如何由导数分量表示。以下固定非空开集 \(\Omega\subseteq\mathbb R^d\)、\(k\in\mathbb N_0\) 及标量域 \(\mathbb K=\mathbb R\) 或 \(\mathbb C\)，不要求 \(\Omega\) 有界或具有规则边界。

% 实读来源：Hunter2001Applied，作者官方 ch12.pdf，印刷 pp.362--367；
% 定义12.66及Banach/Hilbert结构的陈述在p.363，未附证明。
% pp.365--367为作者明确不附证明的性质综述，不作为本节技术证明来源。
% 下文有限乘积、闭性、可分性及对偶证明以本书第8--10章已证结果展开，
% 不使用第22章的全局光滑稠密性或边界延拓。

\subsection{导数分量与完备性}
\label{b1:subsec:210201}

沿用前节的有限指标集 \(A_k=\{\,\alpha\in\mathbb N_0^d\mid|\alpha|\leq k\,\}\)，并置
\[
 Y_p=\prod_{\alpha\in A_k}L^p(\Omega;\mathbb K).
\]
对 \(v=(v_\alpha)_{\alpha\in A_k}\)，赋予范数
\[
 \|v\|_{Y_p}=
 \begin{cases}
  \left(\displaystyle\sum_{\alpha\in A_k}\|v_\alpha\|_p^p\right)^{1/p},
       &1\leq p<\infty,\\[2mm]
  \displaystyle\max_{\alpha\in A_k}\|v_\alpha\|_\infty,
       &p=\infty.
 \end{cases}
\]
这里各分量相互独立；任取这样的函数族，通常不能把它们全部看成同一个函数的导数。

\begin{proposition}\label{b1:prop:sobolev-derivative-embedding}
映射
\[
 \mathcal J:W^{k,p}(\Omega)\longrightarrow Y_p,
 \quad u\longmapsto(\partial^\alpha u)_{\alpha\in A_k}
\]
是线性等距嵌入。特别地，前一节规定的 \(\|\cdot\|_{W^{k,p}}\) 确为范数，每个导数映射
\(u\mapsto\partial^\alpha u\) 都是范数不超过 \(1\) 的有界线性映射。
\end{proposition}
\begin{proof}
弱导数恒等式对函数线性，唯一性保证这些运算在几乎处处等价类上有定义。有限 \(p\) 的三角不等式来自各分量中的 Minkowski 不等式及有限序列的 Minkowski 不等式；后者是对有限计数测度应用\cref{b1:thm:lp-minkowski}。无穷端点直接取各分量范数的最大值。范数等式
\(\|\mathcal Ju\|_{Y_p}=\|u\|_{W^{k,p}}\) 来自定义。由于指标集含零指标，\(\mathcal Ju=0\) 特别蕴含 \(u=0\) 几乎处处，所以映射单射，范数也正定。最后，每个分量的范数不超过整个乘积范数，得到导数映射的估计。
\end{proof}

\begin{theorem}[Sobolev 空间的完备性]\label{b1:thm:sobolev-complete}
对每个 \(1\leq p\leq\infty\)，像空间 \(\mathcal J\bigl(W^{k,p}(\Omega)\bigr)\) 是 \(Y_p\) 的闭线性子空间，因而 \(W^{k,p}(\Omega)\) 是 Banach 空间。
\end{theorem}
\begin{proof}
先看 \(Y_p\)。其中的 Cauchy 列在每个分量上都是 \(L^p\) 中的 Cauchy 列。由\cref{b1:thm:lp-complete}，各分量都有极限；分量只有有限个，所以这些极限组成的向量也是原序列在 \(Y_p\) 中的极限。因此 \(Y_p\) 完备。

设 \(u_j\in W^{k,p}(\Omega)\)，并且 \(\mathcal Ju_j\to(v_\alpha)_\alpha\) 于 \(Y_p\)。取零阶分量 \(u=v_0\)。我们需要证明 \(v_\alpha=\partial^\alpha u\)，不能只把它们命名为导数。给定 \(\varphi\in C_c^\infty(\Omega)\)，记 \(p'\) 为共轭指数，端点取 \(1'=\infty\)、\(\infty'=1\)。函数 \(\varphi\) 及其各阶导数均有紧支集，故属于 \(L^{p'}(\Omega)\)。Hölder 不等式给出
\[
 \left|\int_\Omega(\partial^\alpha u_j-v_\alpha)\varphi\,\mathrm dx\right|
 \leq\|\partial^\alpha u_j-v_\alpha\|_p\cdot\|\varphi\|_{p'}
 \longrightarrow0,
\]
也给出 \(\int_\Omega u_j\partial^\alpha\varphi\to\int_\Omega u\partial^\alpha\varphi\)。于是对每个 \(\alpha\in A_k\)，
\[
 \int_\Omega v_\alpha\varphi\,\mathrm dx
 =\lim_{j\to\infty}\int_\Omega\partial^\alpha u_j\,\varphi\,\mathrm dx
 =(-1)^{|\alpha|}\int_\Omega u\partial^\alpha\varphi\,\mathrm dx.
\]
所有分量均局部可积，所以上式正是弱导数的定义。因而 \(u\in W^{k,p}(\Omega)\) 且 \(\mathcal Ju=(v_\alpha)_\alpha\)，像空间闭。最后应用\cref{b1:prop:closed-subspace-complete}，再通过等距映射 \(\mathcal J\) 把完备性传回原空间。
\end{proof}

证明中没有在 \(\partial\Omega\) 上做任何运算。测试函数的支集留在域内，使极限所需的积分始终受控；这也是任意开集都适用的原因。另一方面，不能在乘积空间中随意指定导数。例如在 \((0,1)\) 上，\((0,1)\in L^p\times L^p\) 不可能等于 \((u,u')\)：零阶分量要求 \(u=0\)，其弱导数便也只能为零。

\subsection{Hilbert 结构}
\label{b1:subsec:210202}

\begin{corollary}\label{b1:cor:sobolev-hilbert}
空间 \(H^k(\Omega)=W^{k,2}(\Omega)\) 关于内积
\[
 \langle u,v\rangle_{H^k}
 =\sum_{\alpha\in A_k}
   \int_\Omega\partial^\alpha u\,\overline{\partial^\alpha v}\,\mathrm dx
\]
是 Hilbert 空间，其内积范数恰为前一节的 \(W^{k,2}\) 范数。
\end{corollary}
\begin{proof}
每一项由 \(L^2\) 中的 Cauchy--Schwarz 不等式保证有限。积分的线性性与复共轭规则给出第一变量线性和共轭对称性；且
\[
 \langle u,u\rangle_{H^k}
 =\sum_{\alpha\in A_k}\|\partial^\alpha u\|_2^2
 =\|u\|_{W^{k,2}}^2.
\]
零阶项保证正定性。前一定理已证明该范数完备，因此它是 Hilbert 内积。
\end{proof}

这使\chapref{b1:ch:08}的正交投影与表示定理可以用于 Sobolev 函数。具体地，由\cref{b1:thm:hilbert-riesz}，每个 \(\Lambda\in H^k(\Omega)^*\) 都有唯一 \(w\in H^k(\Omega)\)，使
\[
 \Lambda(u)=\sum_{\alpha\in A_k}
       \int_\Omega\partial^\alpha u\,\overline{\partial^\alpha w}\,\mathrm dx,
 \quad \|\Lambda\|=\|w\|_{H^k}.
\]
复情形中，\(w\mapsto\Lambda\) 是共轭线性的，不能把上式中的共轭静默去掉。

\ProseParagraphBreak

例如，给定这个泛函，考虑
\[
 \mathcal E(u)=\dfrac12\|u\|_{H^k}^2-\operatorname{Re}\Lambda(u).
\]
把 \(\Lambda(u)=\langle u,w\rangle_{H^k}\) 代入并展开内积，得到
\[
 \mathcal E(u)=\dfrac12\|u-w\|_{H^k}^2-\dfrac12\|w\|_{H^k}^2.
\]
因此最小值在 \(u=w\) 处唯一达到。这里同时控制了零阶及高阶导数；若删去零阶项，常数函数可能有零能量，便不能不加条件地沿用同一范数和唯一性论证。

\subsection{可分性与自反性}
\label{b1:subsec:210203}

\begin{proposition}\label{b1:prop:sobolev-separable-reflexive}
若 \(1\leq p<\infty\)，则 \(W^{k,p}(\Omega)\) 可分；若 \(1<p<\infty\)，则它还自反。
\end{proposition}
\begin{proof}
先给出 \(L^p(\Omega)\) 的可数稠密集。由\cref{b1:prop:lp-continuous-dense}，只须在 \(\mathbb R^d\) 上逼近连续紧支集函数。把空间划成边长 \(2^{-j}\) 的半开立方体，在有限多个这样的立方体上取有理数值，其余处取零；复情形取实部、虚部均有理的数。所有这些阶梯函数组成可数集。

给定 \(f\in C_c(\mathbb R^d)\)，固定一个内部包含其支集的大立方体。对充分细的网格，仅在与支集相交的网格立方体上取样，再把样本值逼近为上述有理数值。函数 \(f\) 一致连续，故所得阶梯函数与 \(f\) 的一致误差趋于零；误差支集又全落在一个固定有界集内，因此 \(L^p\) 误差也趋于零。任意 \(L^p(\Omega)\) 函数在域外补零后属于 \(L^p(\mathbb R^d)\)，将上述逼近限制回 \(\Omega\)，便得到所需可数稠密集。这里的补零只用于 \(L^p\) 函数，没有断言它保持弱导数。

有限乘积 \(Y_p\) 因而可分。可分度量空间的子空间仍可分：用一个可数稠密集中的点作中心、正有理数作半径，得到可数拓扑基；对其中每个与子空间相交的基元素，选取一个交点，这些交点便在子空间中稠密。将此事实用于 \(\mathcal J\bigl(W^{k,p}(\Omega)\bigr)\)，可分性得证。

为证明自反性，把有限乘积看成一个 \(L^p\) 空间。令
\[
 Z=\Omega\times A_k,
 \quad \nu(A)=\sum_{\alpha\in A_k}
       m_d\bigl(\{\,x\in\Omega\mid(x,\alpha)\in A\,\}\bigr),
\]
其中可测集是各截面都为 Lebesgue 可测的集合。有限求和定义了一个 \(\sigma\)-有限测度，映射 \((v_\alpha)_\alpha\mapsto[(x,\alpha)\mapsto v_\alpha(x)]\) 将 \(Y_p\) 线性等距地识别为 \(L^p(\nu)\)。若 \(1<p<\infty\)，由\cref{b1:thm:lp-reflexive}，\(Y_p\) 自反；其闭线性子空间由\cref{b1:cor:closed-subspace-reflexive}也自反。通过 \(\mathcal J\) 即得结论。
\end{proof}

有限 \(p\) 的可分性并不来自“所有 Sobolev 函数都能由光滑函数逼近”这一尚未证明的命题。我们只使用了 \(L^p\) 的可分性及度量子空间的性质；由此选出的稠密族未必光滑。

\ProseParagraphBreak

无穷端点不能作同样结论，甚至 \(W^{1,\infty}(0,1)\) 也不可分。函数本身可以连续且一致有界，其导数中的跳跃位置仍能在无穷范数下彼此分离；本章习题将给出这一边界的具体构造。

\subsection{弱收敛}
\label{b1:subsec:210204}

范数收敛等价于有限多个导数分量的 \(L^p\) 范数收敛。弱收敛也有相应的分量刻画，但不能仅由范数定义直接宣布这一点：还需说明原空间上的每个连续泛函都能由分量上的泛函组合而成。

\begin{proposition}\label{b1:prop:sobolev-functional-representation}
设 \(1\leq p<\infty\)，\(q=p'\)。每个 \(\Lambda\in W^{k,p}(\Omega)^*\) 都可写为
\begin{equation}\label{b1:eq:sobolev-functional-representation}
 \Lambda(u)=\sum_{\alpha\in A_k}
   \int_\Omega g_\alpha\,\partial^\alpha u\,\mathrm dx,
 \quad g_\alpha\in L^q(\Omega).
\end{equation}
反过来，每个这样的有限函数族都定义连续线性泛函，且
\[
 \|\Lambda\|\leq\|(g_\alpha)_\alpha\|_{Y_q}.
\]
表示族通常不唯一，但可选得此处等号成立；换言之，\(\|\Lambda\|\) 是所有表示族的 \(Y_q\) 范数的最小值。
\end{proposition}
\begin{proof}
在 \(\mathcal J\bigl(W^{k,p}(\Omega)\bigr)\) 上定义 \(F(\mathcal Ju)=\Lambda(u)\)。等距性保证 \(F\) 良定义且范数等于 \(\|\Lambda\|\)。由\cref{b1:thm:hahn-banach}，它有同范数延拓 \(\widetilde F\in Y_p^*\)。使用上一证明中的 \(\sigma\)-有限测度空间 \((Z,\nu)\)，由\cref{b1:thm:lp-dual}及 \(p=1\) 时的\cref{b1:thm:l1-dual-sigma-finite}，存在 \(g\in L^q(\nu)\)，满足
\[
 \widetilde F(v)=\int_Zgv\,\mathrm d\nu,
 \quad\|g\|_{L^q(\nu)}=\|\widetilde F\|=\|\Lambda\|.
\]
写成截面 \(g_\alpha(x)=g(x,\alpha)\)，再限制到 \(v=\mathcal Ju\)，即得所需表示和等范数。反方向在 \((Z,\nu)\) 上应用 Hölder 不等式，得到
\[
 |\Lambda(u)|
 \leq\|(g_\alpha)_\alpha\|_{Y_q}\cdot\|\mathcal Ju\|_{Y_p}
 =\|(g_\alpha)_\alpha\|_{Y_q}\cdot\|u\|_{W^{k,p}}.
\]
于是每个表示族的范数均不小于 \(\|\Lambda\|\)，先前构造的族达到此下界。
\end{proof}

这个表示沿用 \(L^p\) 对偶的双线性积分约定，\(g_\alpha\) 上不取共轭。它与 Hilbert 表示的区别还在于各分量未被要求来自同一函数。非唯一性可以直接看见：若 \(k\geq1\)、\(\psi\in C_c^\infty(\Omega)\)，取 \(g_0=\partial_1\psi\)、\(g_{e_1}=\psi\)，其余分量取零，则弱导数恒等式给出
\[
 \int_\Omega u\partial_1\psi\,\mathrm dx
 +\int_\Omega\psi\partial_1u\,\mathrm dx=0.
\]
所以可以把这一族加到任何表示族上而不改变泛函。它并不抵触 Hilbert 表示向量 \(w\) 的唯一性。

\begin{proposition}\label{b1:prop:sobolev-weak-components}
设 \(1\leq p\leq\infty\)，\(u_j,u\in W^{k,p}(\Omega)\)。则
\[
 u_j\rightharpoonup u\text{ 于 }W^{k,p}(\Omega)
 \quad\Longleftrightarrow\quad
 \partial^\alpha u_j\rightharpoonup\partial^\alpha u
 \text{ 于 }L^p(\Omega)\quad(\alpha\in A_k).
\]
当 \(p<\infty\) 时，右边又等价于对每个 \(\alpha\in A_k\) 和 \(g\in L^{p'}(\Omega)\)，
\[
 \int_\Omega g\partial^\alpha u_j\,\mathrm dx
 \longrightarrow\int_\Omega g\partial^\alpha u\,\mathrm dx.
\]
\end{proposition}
\begin{proof}
若在 Sobolev 空间中弱收敛，将 \(L^p\) 上任意连续线性泛函与有界导数映射 \(u\mapsto\partial^\alpha u\) 复合，便得到各分量的弱收敛。

反过来，设每个分量弱收敛。对任意 \(\Lambda\in W^{k,p}(\Omega)^*\)，再次用 Hahn--Banach 定理将其通过 \(\mathcal J\) 延拓为 \(\widetilde F\in Y_p^*\)。记 \(\iota_\alpha:L^p(\Omega)\to Y_p\) 为只在第 \(\alpha\) 个分量放入输入函数的映射，令 \(F_\alpha=\widetilde F\circ\iota_\alpha\)。这些都是 \(L^p\) 上的连续线性泛函，且有限分量分解给出
\[
 \Lambda(u_j)=\sum_{\alpha\in A_k}F_\alpha(\partial^\alpha u_j)
 \longrightarrow\sum_{\alpha\in A_k}F_\alpha(\partial^\alpha u)
 =\Lambda(u).
\]
此论证也适用于 \(p=\infty\)。有限 \(p\) 时再用 \(L^p\) 对偶表示，便得到积分形式的刻画。
\end{proof}

在无穷端点，只用 \(L^1\) 函数检验各分量，表达的是各分量的\WeakStar{}收敛，不能据此改称 \(W^{k,\infty}\) 的弱收敛。上一个命题的端点版本使用整个 \((L^\infty)^*\)。

\begin{corollary}\label{b1:cor:sobolev-weak-subsequence}
若 \(1<p<\infty\) 且 \((u_j)\) 在 \(W^{k,p}(\Omega)\) 中有界，则存在子列 \((u_{j_\ell})\) 及 \(u\in W^{k,p}(\Omega)\)，使 \(u_{j_\ell}\rightharpoonup u\)。并且
\[
 \|u\|_{W^{k,p}}\leq\liminf_{\ell\to\infty}\|u_{j_\ell}\|_{W^{k,p}}.
\]
\end{corollary}
\begin{proof}
由自反性及\cref{b1:thm:reflexive-weak-subsequence}取得弱收敛子列；最后的不等式来自\cref{b1:prop:norm-weak-lsc}。
\end{proof}

也可以先对各阶导数应用 \(L^p\) 的弱子列结论。分量有限，逐次抽取后有共同子列；把测试函数代入各分量的弱极限，如完备性证明那样即可辨认零阶极限的弱导数。前一命题保证最后得到的确实是 Sobolev 空间中的弱收敛。这种分量做法常用于逐阶估计近似解。

\ProseParagraphBreak

当 \(p=1\) 时，有界性本身不再保证上述抽列结论，导数质量可能集中在越来越小的区域内。阻止这种集中的条件属于\secref{b1:sec:1005}的一致可积性理论，不由 Sobolev 范数的有界性自动产生。章末将用一维有界序列说明这个障碍。

% !TeX root = ../../Book1.tex
\section{运算规则}
\label{b1:sec:2103}

经典微积分的运算公式能否保留，要同时回答两个问题：等式在分布意义下是否成立，右端是否具有所要求的可积性。下面先在开集内部把函数光滑化，再用局部积分收敛传递公式。整个过程只使用第20章的卷积，不依赖下一章的全域光滑稠密性或延拓定理。

\begin{lemma}\label{b1:lem:sobolev-local-smoothing}
设 \(u\in W^{k,p}_{\mathrm{loc}}(\Omega)\)，\(U,V\) 为开集，\(\overline U\Subset V\)、\(\overline V\Subset\Omega\)。取第20章的非负单位积分磨光函数，令 \(u_\varepsilon=u*\rho_\varepsilon\)。当 \(\varepsilon\) 充分小时，\(u_\varepsilon\) 在 \(U\) 的邻域光滑，且
\[
 \partial^\alpha u_\varepsilon
   =(\partial^\alpha u)*\rho_\varepsilon
 \quad\text{在 }U\text{ 上},\qquad |\alpha|\le k.
\]
若 \(p<\infty\)，则 \(u_\varepsilon\to u\) 于 \(W^{k,p}(U)\)。若 \(p=\infty\)，则对每个有限 \(q\ge1\) 有 \(W^{k,q}(U)\) 收敛，并有
\[
 \|u_\varepsilon\|_{W^{k,\infty}(U)}
 \le\|u\|_{W^{k,\infty}(V)}.
\]
\end{lemma}
\begin{proof}
选择 \(\varepsilon<\operatorname{dist}(\overline U,V^c)\)。卷积在 \(U\) 上只读取 \(V\) 中的值，导数交换由\cref{b1:prop:distribution-convolution-smooth}成立。对有限个导数分别应用\cref{b1:prop:local-smoothing-approximation}，得到有限 \(p\) 的范数收敛。若 \(p=\infty\)，这些导数在 \(V\) 上也属于每个有限 \(L^q\)，故局部收敛仍成立；非负核与单位积分则给出每个导数的本性界，取最大值即得所列估计。
\end{proof}

后面的证明还会使用这样的取极限方式：若 \(u_j\to u\)、\(g_{i,j}\to g_i\) 于 \(L^1_{\mathrm{loc}}\)，且 \(\partial_i u_j=g_{i,j}\)，则 \(\partial_i u=g_i\)。只需对固定测试函数在分部积分恒等式两端取极限，正如第\ref{b1:sec:2004}节已说明的那样。

无穷端点只提供有限 \(q\) 的局部逼近和统一界，不能改写为 \(W^{k,\infty}\) 范数逼近。例如 \(|x|\) 的弱导数为符号函数，光滑函数不可能在原点附近依本性上确界范数逼近这个跳跃。

\subsection{乘积法则}
\label{b1:subsec:210301}

先考虑局部有界的因子，这使乘积右端仍处于原来的 \(L^p\) 尺度。

\begin{proposition}[Leibniz]\label{b1:prop:sobolev-product}
若 \(u,v\in W^{1,p}_{\mathrm{loc}}(\Omega)\cap L^\infty_{\mathrm{loc}}(\Omega)\)，则 \(uv\in W^{1,p}_{\mathrm{loc}}(\Omega)\)，并且
\begin{equation}\label{b1:eq:sobolev-product}
 \partial_i(uv)=u\partial_i v+v\partial_i u
 \quad\text{几乎处处},\qquad 1\le i\le d.
\end{equation}
若两个函数均属于全域 \(W^{1,p}(\Omega)\cap L^\infty(\Omega)\)，则乘积也属于该空间，且
\[
 \|uv\|_{W^{1,p}}
 \le\|u\|_\infty\cdot\|v\|_{W^{1,p}}
      +\|v\|_\infty\cdot\|u\|_{W^{1,p}}.
\]
\end{proposition}
\begin{proof}
先设 \(p<\infty\)，固定 \(U,V\) 如前一引理。局部卷积给出光滑 \(u_j,v_j\)，它们在 \(U\) 中依 \(W^{1,p}\) 收敛，并可取一个子列，使两列函数几乎处处收敛。非负核保证 \(|u_j|\)、\(|v_j|\) 具有由原函数在 \(V\) 上的本性界给出的公共界。

因此 \(u_jv_j\to uv\) 于 \(L^p(U)\)。对一个导数项，分解
\[
 u_j\partial_i v_j-u\partial_i v
 =u_j(\partial_i v_j-\partial_i v)
   +(u_j-u)\partial_i v.
\]
第一项由统一有界性和导数的 \(L^p\) 收敛趋零；第二项由几乎处处收敛及以 \(C|\partial_i v|^p\) 为控制函数的控制收敛定理趋零。另一导数项同理。对经典乘积公式取极限，再遍历内部开集，就得到\eqref{b1:eq:sobolev-product}。

若 \(p=\infty\)，先把相同函数视为局部 \(W^{1,1}\) 成员，用刚证明的公式；其右端由假设局部本性有界，所以确实给出 \(W^{1,\infty}_{\mathrm{loc}}\) 关系。全域情形直接估计右端各分量，应用有限分量的 Minkowski 不等式，即得范数界；乘积的本性有界性也显然成立。
\end{proof}

只知道 \(u,v\in W^{1,p}\) 时，不能无条件断言乘积仍属于 \(W^{1,p}\)。弱导数公式并不自动提高 \(u\partial_i v\) 的可积性；第\ref{b1:ch:24}章的嵌入定理以后会提供另外一些充分条件。习题\ref{b1:ex:21-product-integrability-loss}则直接构造乘积失去原有可积性的例子。

若其中一个因子为光滑函数，不需要假设另一因子局部有界。第20章的弱 Leibniz 公式直接给出
\[
 \partial_i(au)=a\partial_i u+(\partial_i a)u.
\]
特别地，\(a\in C_c^\infty(\Omega)\) 时可以截出一个内部部分；这里产生的额外项 \(u\partial_i a\) 不能省略。

\subsection{链式法则}
\label{b1:subsec:210302}

本小节先讨论实值函数。对于向量值映射，把每个分量分别放在 Sobolev 空间中，证明完全相同；复值函数可以按两个实分量处理，而不是擅自要求复合映射全纯。

\begin{theorem}[链式法则]\label{b1:thm:sobolev-chain}
设 \(u\in W^{1,p}_{\mathrm{loc}}(\Omega;\mathbb R^m)\)，\(F\in C^1(\mathbb R^m;\mathbb R^\ell)\)，且 \(DF\) 有界。则
\[
 F\circ u\in W^{1,p}_{\mathrm{loc}}(\Omega;\mathbb R^\ell),
 \qquad
 \partial_i(F\circ u)=DF(u)\partial_i u
 \quad\text{几乎处处}.
\]
若 \(u\) 局部本性有界，则可去掉 \(DF\) 全域有界的要求。
\end{theorem}
\begin{proof}
有界导数使 \(F\) 为 Lipschitz 映射，因而
\[
 |F(u)|\le |F(0)|+\|DF\|_\infty\cdot|u|.
\]
右端局部属于 \(L^p\)，包括无穷端点。先取 \(p<\infty\)，并在固定内部区域用前面的卷积引理选取 \(u_j\to u\) 于 \(W^{1,p}\)，同时几乎处处收敛。Lipschitz 界给出 \(F(u_j)\to F(u)\) 于 \(L^p\)。导数分解为
\[
 DF(u_j)\partial_i u_j-DF(u)\partial_i u
 =DF(u_j)(\partial_i u_j-\partial_i u)
   +\bigl(DF(u_j)-DF(u)\bigr)\partial_i u.
\]
第一项的 \(L^p\) 范数趋零。第二项由于 \(DF\) 连续而几乎处处趋零，并受 \(2\|DF\|_\infty\cdot|\partial_i u|\) 控制，所以也依 \(L^p\) 趋零。经典链式法则在极限下给出弱导数公式。

\(p=\infty\) 时先在局部使用 \(p=1\) 的结论，再用所得公式的本性界。若 \(u\) 仅局部本性有界，对固定区域取包含其值域和卷积值域的紧球，用光滑截断把 \(F\) 改成该球邻域上与其相等、导数全域有界的映射，即归结为已证情形。
\end{proof}

全域结论还需要零阶项可积。若 \(u\in W^{1,p}(\Omega)\)，\(DF\) 有界且 \(F(0)=0\)，则 \(F(u)\in W^{1,p}(\Omega)\)。若 \(p<\infty\) 且 \(m_d(\Omega)<\infty\)，允许 \(F(0)\ne0\) 也可以；但在无限测度域上，\(u=0\)、\(F(t)=1+t\) 已经说明不能遗漏这个条件。\(p=\infty\) 时，全域本性有界的 \(u\) 经连续 \(F\) 映射后仍本性有界。

稍后要把光滑的 \(F\) 换成折线函数。折点处不能任意填写一个经典导数便结束论证，所需事实是：Sobolev 函数在一个水平集上没有非零弱梯度。

\begin{proposition}\label{b1:prop:sobolev-gradient-on-level-sets}
设实值 \(u\in W^{1,p}_{\mathrm{loc}}(\Omega)\)。对每个 \(c\in\mathbb R\)，有
\[
 \nabla u=0\quad\text{几乎处处于 }\{u=c\}.
\]
这里的水平集可以用任一可测代表定义，结论不受零测修改影响。
\end{proposition}
\begin{proof}
取 \(0\le\theta\le1\)、\(\theta\in C_c^\infty((-1,1))\)、\(\theta(0)=1\)，令
\[
 F_\varepsilon(t)=\int_c^t\theta\left(\frac{s-c}{\varepsilon}\right)\,\mathrm ds.
\]
有 \(|F_\varepsilon|\le2\varepsilon\)，而链式法则给出
\[
 \partial_i F_\varepsilon(u)
 =\theta\left(\frac{u-c}{\varepsilon}\right)\partial_i u.
\]
函数在每个内部有界区域趋于零，右端则由控制收敛在局部 \(L^1\) 中趋于 \(\mathbf1_{\{u=c\}}\partial_i u\)。弱导数关系的闭性说明后一个极限是零函数的弱导数，因此等于零。这个证明也包括 \(p=\infty\)，并没有使用尚待证明的沿直线代表或近似可微性。
\end{proof}

\subsection{截断}
\label{b1:subsec:210303}

对实数 \(M>0\)，定义截断函数
\[
 T_M(t)=\max\{-M,\min\{t,M\}\}.
\]
它保留 \([-M,M]\) 内的值，并把其余值压到两个端点。与把定义域截成一个子集不同，这是一种作用于函数值的操作。

\begin{proposition}\label{b1:prop:sobolev-truncation}
若实值 \(u\in W^{1,p}_{\mathrm{loc}}(\Omega)\)，则
\[
 T_M(u)\in W^{1,p}_{\mathrm{loc}}(\Omega),\qquad
 \nabla T_M(u)=\mathbf1_{\{|u|<M\}}\nabla u
 \quad\text{几乎处处}.
\]
正部 \(u^+=\max\{u,0\}\)、负部 \(u^-=\max\{-u,0\}\) 和绝对值也属于同一局部空间，且
\[
 \nabla u^+=\mathbf1_{\{u>0\}}\nabla u,\quad
 \nabla u^-=-\mathbf1_{\{u<0\}}\nabla u,\quad
 \nabla|u|=\operatorname{sgn}(u)\nabla u.
\]
约定 \(\operatorname{sgn}(0)=0\)。若 \(u\in W^{1,p}(\Omega)\)，上述函数也属于全域空间。
\end{proposition}
\begin{proof}
在实轴上用非负偶光滑核 \(\eta_\varepsilon\) 卷积 \(T_M\)，记为 \(F_\varepsilon\)。由于 \(T_M\) 奇且为 \(1\)-Lipschitz，\(F_\varepsilon(0)=0\)、\(|F_\varepsilon'|\le1\)，并且 \(F_\varepsilon\to T_M\) 一致成立。分段求导或直接积分可得
\[
 F_\varepsilon'=\mathbf1_{(-M,M)}*\eta_\varepsilon.
\]
因此当 \(t\ne\pm M\) 时，导数趋于 \(\mathbf1_{(-M,M)}(t)\)。在 \(\{u=\pm M\}\) 上，前一命题保证 \(\nabla u=0\)，所以无论近似导数在两个端点的极限为何，都有
\[
 F_\varepsilon'(u)\nabla u
 \longrightarrow\mathbf1_{\{|u|<M\}}\nabla u
 \quad\text{于 }L^1_{\mathrm{loc}}.
\]
由链式法则和弱导数的闭性得到所求关系；右端的局部 \(L^p\) 控制给出空间归属。

对 \(t^+\) 作同样卷积，并减去卷积在零点的值，使近似函数仍在零点为零。其导数在 \(t\ne0\) 处趋于 \(\mathbf1_{(0,\infty)}(t)\)，在零水平集上仍由梯度为零消除折点的不确定性。于是正部公式成立。负部及绝对值由 \((-u)^+\) 和 \(u^++u^-\) 得到。

最后，\(|T_M(u)|\)、\(u^+\)、\(u^-\) 和 \(|u|\) 均不超过 \(|u|\)，各导数的绝对值也由原导数控制，故全域情形不会丢失零阶或一阶可积性。
\end{proof}

若 \(1\le p<\infty\) 且 \(u\in W^{1,p}(\Omega)\)，则
\begin{equation}\label{b1:eq:sobolev-truncation-convergence}
 T_M(u)\longrightarrow u\quad\text{于 }W^{1,p}(\Omega)
 \quad(M\to\infty).
\end{equation}
事实上，函数误差由 \(|u|\mathbf1_{\{|u|>M\}}\) 控制，导数误差由 \(|\partial_i u|\mathbf1_{\{|u|\ge M\}}\) 控制，对它们的 \(p\) 次幂分别用控制收敛即可。在全域 \(W^{1,\infty}\) 中，对一个固定 \(u\)，只要 \(M\ge\|u\|_\infty\)，截断已经与原函数几乎处处相等；这不意味着正部等非线性运算对 \(W^{1,\infty}\) 范数总是连续，习题会区分这两个问题。

\subsection{坐标变换}
\label{b1:subsec:210304}

改变坐标还会改变积分域和零测集。若一个坐标映射把零测集拉回成正测度集，\(u\circ\Phi\) 就可能依赖 \(u\) 的代表；因此先明确映射的几何控制。

\begin{theorem}\label{b1:thm:sobolev-bilipschitz-change}
设 \(U,V\subset\mathbb R^d\) 为开集，\(\Phi:U\to V\) 是双 Lipschitz 同胚，且
\[
 |\Phi(x)-\Phi(y)|\le L|x-y|,\qquad
 |\Phi^{-1}(z)-\Phi^{-1}(w)|\le M|z-w|.
\]
若 \(u\in W^{1,p}(V)\)，则 \(u\circ\Phi\in W^{1,p}(U)\)，且
\begin{equation}\label{b1:eq:sobolev-coordinate-chain}
 \nabla(u\circ\Phi)(x)
   =D\Phi(x)^{\mathsf T}(\nabla u)(\Phi(x))
 \quad\text{对几乎处处的 }x\in U.
\end{equation}
包含零阶项的范数满足
\[
 \|u\circ\Phi\|_{W^{1,p}(U)}
 \le C(d,p,L,M)\|u\|_{W^{1,p}(V)}.
\]
\end{theorem}
\begin{proof}
由 Lipschitz 映射的零测集性质，\(\Phi^{-1}\) 把 \(V\) 中的 Lebesgue 零测集映成 \(U\) 中的零测集，所以复合不依赖代表。对非负可测 \(g\)，第17章的加权面积公式用于单射映射 \(\Phi^{-1}\)，给出
\begin{equation}\label{b1:eq:sobolev-change-measure-bound}
 \int_U g(\Phi(x))\,\mathrm dx
 =\int_V g(y)|\det D\Phi^{-1}(y)|\,\mathrm dy
 \le M^d\int_Vg(y)\,\mathrm dy.
\end{equation}
这里 Rademacher 定理保证导数几乎处处存在，算子范数不超过 \(M\)，从而行列式绝对值不超过 \(M^d\)。

先设 \(p<\infty\)，固定 \(U_0\) 满足 \(\overline{U_0}\Subset U\)。其像的闭包紧含于 \(V\)，因而可以在像的一个内部邻域对 \(u\) 作局部 \(W^{1,p}\) 光滑逼近 \(u_j\)。函数 \(u_j\circ\Phi\) 局部 Lipschitz；经典链式法则在 \(\Phi\) 的可微点给出
\[
 \nabla(u_j\circ\Phi)
 =D\Phi^{\mathsf T}(\nabla u_j)\circ\Phi.
\]
由\cref{b1:thm:lipschitz-weak-derivative}，这也是弱导数等式。估计\eqref{b1:eq:sobolev-change-measure-bound}使复合后的函数差及导数差在 \(U_0\) 中依 \(L^p\) 趋零；\(D\Phi\) 的算子范数又不超过 \(L\)，故右端也在 \(L^p\) 中收敛。弱导数关系的闭性于是给出\eqref{b1:eq:sobolev-coordinate-chain}。由于 \(U_0\) 任意，这一关系在整个 \(U\) 上成立。

若 \(p=\infty\)，同一局部证明先使用有限指数 \(1\)，得到公式后再取本性界。最后，对有限 \(p\) 用\eqref{b1:eq:sobolev-change-measure-bound}，对无穷端点用零测集拉回性质，并结合有限维向量范数等价性，得到包含函数与全部一阶偏导数的全域范数估计。
\end{proof}

对 \(\Phi^{-1}\) 再用一次定理可知，这个复合映射实际给出两个 \(W^{1,p}\) 空间之间的有界线性同构。它并不保持本书范数的数值，也不由单向 Lipschitz 控制推出逆方向的有界性。

高阶空间需要更强的坐标条件。若 \(\Phi\) 还是 \(C^k\) 微分同胚，并且 \(1\) 至 \(k\) 阶的各导数有界，则相同方法给出 \(u\circ\Phi\in W^{k,p}(U)\) 及相应有界估计：对光滑 \(u_j\)，逐次求导后，每一项都是某个 \((\partial^\beta u_j)\circ\Phi\) 乘以有限个 \(\Phi\) 的导数，且 \(|\beta|\le k\)；系数有统一界，故逐项用换元估计取极限即可。若还要得到反向的 \(W^{k,p}\) 有界映射，同样要控制 \(\Phi^{-1}\) 的高阶导数。仅有双 Lipschitz 性只控制一阶，不能代替这些条件。

% !TeX root = ../../Book1.tex
\section{Sobolev 函数的代表}
\label{b1:sec:2104}

Sobolev 空间按几乎处处相等取商，而连续性、沿直线的变化和差商都涉及点值。讨论这些性质时，必须先说明选择什么代表。本节先把一维积分表示推广到几乎所有坐标直线，再区分通常微分与近似微分。以下以实值函数为主；有限维向量值函数的相应结论可逐分量应用。

% 实读 Hajlasz2003SobolevMetric，作者公开正式版，Contemp. Math. 338 (2003)，
% §2 pp.175--177：ACL 定义、定理2.3及式(2.5)；§7定理7.13 pp.192--193；
% §9定理9.4及完整证明 pp.204--205。
% 原文7.13经上梯度并使用全域光滑稠密性；本节不移用该前置，
% ACL 改用第20章局部磨光及可数内矩形上的公共快速收敛子列。
% 双点估计改编9.4的均值望远镜路线，立方体平均振荡先由ACL自足证明。

\subsection{一维绝对连续代表}
\label{b1:subsec:210401}

一维情形中，弱导数不只决定积分配对，也决定代表在任意两点之间的增量。

\begin{proposition}\label{b1:prop:sobolev-interval-representative}
设 \(I\) 为开区间、\(1\leq p\leq\infty\)。函数 \(u\) 属于 \(W^{1,p}_{\mathrm{loc}}(I)\)，当且仅当它有局部绝对连续代表 \(\widetilde u\)，且 \(\widetilde u,\widetilde u'\in L^p_{\mathrm{loc}}(I)\)。这个代表唯一，经典导数与弱导数几乎处处相等。

若 \(I=(a,b)\) 有界且 \(u\in W^{1,p}(I)\)，则该代表还唯一延伸为 \([a,b]\) 上的绝对连续函数，并满足
\begin{equation}\label{b1:eq:interval-sobolev-sup}
 \|\widetilde u\|_{L^\infty([a,b])}
 \leq\dfrac{\|u\|_{L^1(I)}}{b-a}+\|u'\|_{L^1(I)}.
\end{equation}
\end{proposition}
\begin{proof}
局部刻画及唯一性就是定理\ref{b1:thm:one-dimensional-weak-ac}。在有界区间上，由 Hölder 不等式，\(u,u'\in L^1(I)\)。固定内点 \(x_0\)，积分表示
\[
 \widetilde u(x)=\widetilde u(x_0)+\int_{x_0}^x u'(t)\,\mathrm dt
\]
一直延伸到两端，并在整个闭区间上绝对连续。对任意 \(x,y\in[a,b]\)，有
\[
 |\widetilde u(x)|\leq|\widetilde u(y)|+\int_a^b|u'(t)|\,\mathrm dt.
\]
对 \(y\) 取平均，再对 \(x\) 取上确界，即得\eqref{b1:eq:interval-sobolev-sup}。
\end{proof}

当 \(1<p<\infty\) 时，同一积分表示给出
\[
 |\widetilde u(t)-\widetilde u(s)|
 \leq\|u'\|_{L^p(I)}\cdot|t-s|^{1-1/p}
 \quad(s,t\in[a,b]).
\]
当 \(p=\infty\) 时，右侧改为 \(\|u'\|_\infty\cdot|t-s|\)；当 \(p=1\) 时，对 \(s<t\) 的增量由 \(\int_s^t|u'|\) 控制，但不能从 \(L^1\) 范数单独读出一个正的 Hölder 指数。估计\eqref{b1:eq:interval-sobolev-sup}还有另一个用途：函数及其一阶导数的 \(L^1\) 收敛，会使对应的绝对连续代表一致收敛。下一小节将在几乎每条直线上使用这一点。

\subsection{ACL 性质}
\label{b1:subsec:210402}

在多维开集内，平行于某个坐标轴的直线可能与区域相交为多个区间，所以这里的绝对连续性应逐个紧子区间理解。先设 \(d\geq2\)；\(d=1\) 时，下面的条件就是局部绝对连续。

\begin{definition}\label{b1:def:acl}
设 \(v:\Omega\rightarrow\mathbb R\) 为有限值 Borel 函数。若对每个坐标方向 \(e_i\)，对其正交补中几乎所有 \(z\)，限制
\[
 t\longmapsto v(z+te_i)
\]
在开集 \(\{\,t\in\mathbb R\mid z+te_i\in\Omega\,\}\) 的每个紧子区间上绝对连续，则称 \(v\) \term[yan-zuobiaoxian-jueduilianxu]{沿坐标线绝对连续}[absolutely continuous on lines]，简称具有 ACL 性质。这里的“几乎所有”使用横向变量 \(z\) 上的 \((d-1)\) 维 Lebesgue 测度。
\end{definition}

\begin{theorem}\label{b1:thm:sobolev-acl}
设 \(1\leq p\leq\infty\)，且 \(u,g_1,\ldots,g_d\in L^p_{\mathrm{loc}}(\Omega)\)。则 \(\partial_i u=g_i\) 对所有 \(i\) 成立，当且仅当 \(u\) 有一个具有 ACL 性质的代表 \(\widetilde u\)，其经典坐标偏导数几乎处处存在，且等于 \(g_i\)。
\end{theorem}
\begin{proof}
先设弱导数关系成立。列出所有闭包紧含于 \(\Omega\) 的有理端点开矩形 \(Q_m\)。由\cref{b1:lem:sobolev-local-smoothing}，可以选 \(\varepsilon_j\downarrow0\)，使 \(u_j=u*\rho_{\varepsilon_j}\) 在 \(Q_1,\ldots,Q_j\) 的邻域上都有定义，且
\[
 \|u_j-u\|_{L^1(Q_m)}
 +\sum_{i=1}^d\|\partial_i u_j-g_i\|_{L^1(Q_m)}<2^{-j}
 \quad(m\leq j).
\]
这里仅用了内矩形上的局部 \(L^1\) 逼近。将每个 \(u_j\) 在其定义域外补零，只为得到 \(\Omega\) 上的 Borel 函数；不宣称补零后的函数在全域光滑或保留弱导数。

固定 \(Q_m\) 和方向 \(i\)，写 \(Q_m=Q_m'\times I\)，把第 \(i\) 个坐标放在最后。对上述可求和误差应用 Tonelli 定理，可知对几乎每个 \(z\in Q_m'\)，
\[
 \sum_{j\geq m}\left(
  \|u_j(z,\cdot)-u(z,\cdot)\|_{L^1(I)}
  +\|\partial_i u_j(z,\cdot)-g_i(z,\cdot)\|_{L^1(I)}
 \right)<\infty.
\]
因此这条直线上的 \(u_j\) 在一维 \(W^{1,1}(I)\) 中为 Cauchy 列。将\eqref{b1:eq:interval-sobolev-sup}用于差 \(u_j-u_k\)，得到它们在 \(I\) 上一致收敛到某个函数 \(v_z\)。又因导数在 \(L^1(I)\) 中趋于 \(g_i(z,\cdot)\)，在光滑函数的积分公式中取极限，得
\[
 v_z(t)-v_z(s)=\int_s^t g_i(z,r)\,\mathrm dr
 \quad(s,t\in I).
\]
所以 \(v_z\) 绝对连续，导数几乎处处为该截面的 \(g_i\)。它也与 \(u(z,\cdot)\) 几乎处处相等。

现在在 \(\Omega\) 中统一定义
\[
 \widetilde u(x)=
 \begin{cases}
  \displaystyle\lim_{j\to\infty}u_j(x),&\text{若极限存在且有限},\\
  0,&\text{其他情形}.
 \end{cases}
\]
这是一个 Borel 函数，并由各矩形上的可求和 \(L^1\) 误差知 \(\widetilde u=u\) 几乎处处。在前述每条良好直线上，它在矩形内的每一点都等于一致极限 \(v_z\)。对可数个矩形及有限个方向合并横向零测例外后，仍只排除了各方向中的一个横向零测集。任意紧线段可用有限个这样的矩形覆盖，局部积分表示在交叠处来自同一个函数 \(\widetilde u\)，故可拼接。这证明了同一个代表同时具有全部坐标方向的 ACL 性质。

反过来，在内矩形中取测试函数，沿第 \(i\) 个方向在几乎每条良好直线上作一维分部积分，再由 Fubini 定理得
\[
 \int_{Q_m}\widetilde u\,\partial_i\varphi
 =-\int_{Q_m}g_i\varphi.
\]
各积分因局部可积性而有定义；测试函数在相应区间端点附近为零。弱导数的局部性随后把这些关系拼成整个 \(\Omega\) 上的关系。
\end{proof}

沿几乎所有直线成立，与沿每一条直线成立不同。维数至少为二时，甚至把一个 ACL 代表在单点改值，也只会破坏有限条坐标直线上的绝对连续性，而不改变 ACL 性质。因此 ACL 代表不像一维连续代表那样逐点唯一。定理所需的是存在一个共同代表，并不是给等价类的每一点都指定不可更改的值。

\subsection{\texorpdfstring{\(W^{1,\infty}\)}{W1,infinity} 与局部 Lipschitz 代表}
\label{b1:subsec:210403}

由 ACL 性质，Sobolev 函数的合适代表有几乎处处存在的坐标偏导数，且它们组成弱梯度。但坐标偏导存在不等于通常的全微分存在。本小节先给出能够恢复通常微分的有界梯度情形。

\begin{theorem}\label{b1:thm:w1infty-lipschitz}
函数 \(u\) 属于 \(W^{1,\infty}_{\mathrm{loc}}(\Omega)\)，当且仅当它有局部 Lipschitz 代表。该连续代表唯一，并几乎处处通常可微，其经典微分为弱梯度所确定的线性映射。

更具体地，若 \(B\) 为开球、\(\overline B\Subset\Omega\)，且 \(|\nabla u|\leq L\) 几乎处处于 \(B\)，则此代表在 \(B\) 上为 \(L\)-Lipschitz 函数。
\end{theorem}

这里 \(|\nabla u|\) 使用 \(\mathbb R^d\) 的欧氏范数。它与本章按偏导分量取最大值的范数满足
\[
 \bigl\||\nabla u|\bigr\|_{L^\infty(B)}
 \leq \sqrt d\,\max_{1\leq i\leq d}\|\partial_i u\|_{L^\infty(B)}
 \leq \sqrt d\,\|u\|_{W^{1,\infty}(B)}.
\]
所以若只从本章的 \(W^{1,\infty}\) 范数读取梯度控制，可以直接得到 \(\sqrt d\) 倍的 Lipschitz 常数；定理中的 \(L\) 则是欧氏梯度模的本性上界，不应与前一常数混同。

\begin{proof}
先设 \(u\in W^{1,\infty}_{\mathrm{loc}}\)。固定所述开球，令 \(u_\varepsilon=u*\rho_\varepsilon\)。局部卷积微分公式给出
\[
 \nabla u_\varepsilon=(\nabla u)*\rho_\varepsilon,
 \quad |\nabla u_\varepsilon|\leq L
\]
于卷积只涉及 \(B\) 内点的区域。取 \(B\) 中两个 Lebesgue 点 \(x,y\)。线段 \([x,y]\) 紧含于 \(B\)，故当 \(\varepsilon\) 足够小时，在这条线段上积分光滑函数的导数可得
\[
 |u_\varepsilon(x)-u_\varepsilon(y)|\leq L|x-y|.
\]
由\cref{b1:prop:approx-identity-lebesgue-point}，两端趋于 \(u\) 在相应点的 Lebesgue 值。于是同一 Lipschitz 不等式在 \(B\) 的 Lebesgue 点集上成立。

该点集满测度，因而在 \(B\) 中稠密。对任意点取其中的逼近点列，Lipschitz 估计使函数值为 Cauchy 列；估计也保证极限不依赖所选点列，并在整个 \(B\) 上保留常数 \(L\)。所得连续函数与 \(u\) 几乎处处相等。在交叠球上，两个连续代表几乎处处相等，所以处处相等，由此拼成局部 Lipschitz 代表。唯一性同样来自连续性与几乎处处相等。

反向由定理\ref{b1:thm:lipschitz-weak-derivative}得到。最后对局部 Lipschitz 代表使用 Rademacher 定理\ref{b1:thm:rademacher}；经典偏导与弱偏导几乎处处一致，故全微分的系数正是弱梯度。
\end{proof}

这里的“局部”不能无条件去掉。例如在 \(\Omega=(-1,0)\cup(0,1)\) 上，令 \(u\) 左侧为零、右侧为一，则 \(u\in W^{1,\infty}(\Omega)\)，弱导数为零；但跨越缺失的原点取两侧点，差商可以任意大，故它不是 \(\Omega\) 上的 Lipschitz 函数。

即使函数属于 \(W^{1,\infty}\)，通常可微的断言也针对所选的局部 Lipschitz 代表，而非任意代表。在有界开集上，\(\mathbf1_{\mathbb Q^d}\) 代表零元素，却处处不连续。一般有限 \(p\) 时，上述有界梯度证明不适用；不能把本小节标题理解成“所有 Sobolev 代表都几乎处处通常可微”。对全部 \(1\leq p\leq\infty\) 统一成立的结论是下一小节的近似可微性。

\subsection{近似可微性}
\label{b1:subsec:210404}

近似微分允许忽略密度为零的异常点，但仍要求一阶误差相对于距离趋于零。本小节的证明链为
\[
 \begin{gathered}
  \text{ACL}\Longrightarrow\text{立方体平均振荡估计}\\
  \Longrightarrow\text{极大函数双点估计}\\
  \Longrightarrow\text{可数 Lipschitz 分片}\\
  \Longrightarrow\operatorname{ap}Du=\nabla u.
 \end{gathered}
\]
其中层集 \(\{Mg\leq n\}\) 的作用是在每一层上得到 Lipschitz 分片，并不是把原函数提升为整个邻域上的 Lipschitz 函数。为从弱导数得到上述控制，先把沿坐标线的积分公式转成局部平均振荡的估计。对有限正体积的集合 \(Q\)，简记
\[
 u_Q=\dfrac1{m_d(Q)}\int_Q u.
\]

\begin{lemma}\label{b1:lem:sobolev-cube-oscillation}
设 \(u\in W^{1,1}_{\mathrm{loc}}(\Omega)\)，\(Q\) 为边长 \(\ell\) 的坐标开立方体，且 \(\overline Q\Subset\Omega\)。则
\begin{equation}\label{b1:eq:sobolev-cube-oscillation}
 \int_Q|u-u_Q|
 \leq\ell\sum_{i=1}^d\int_Q|\partial_i u|.
\end{equation}
\end{lemma}
\begin{proof}
选取 ACL 代表。对几乎所有 \(x,y\in Q\)，依次把 \(x\) 的各坐标换成 \(y\) 的坐标，用 \(d\) 条坐标线段连接它们。所有这些线段都属于良好直线的例外集之外；此断言对几乎所有点对成立，正是对每个方向的横向零测集应用 Fubini 定理的结果。沿线段积分，再用三角不等式，得到
\[
 |u(x)-u(y)|
 \leq\sum_{i=1}^d\int_{[x_i,y_i]}
 |\partial_i u(y_1,\ldots,y_{i-1},t,x_{i+1},\ldots,x_d)|\,\mathrm dt,
\]
其中 \([x_i,y_i]\) 表示两端点之间的无向区间。将此式对 \(x,y\) 积分。第 \(i\) 项中未出现的 \(d-1\) 个坐标积分产生因子 \(\ell^{d-1}\)；对固定 \(t\)，使区间 \([x_i,y_i]\) 包含 \(t\) 的点对面积至多为 \(\ell^2\)。由 Tonelli 定理，第 \(i\) 项的二重积分不超过 \(\ell^{d+1}\int_Q|\partial_i u|\)。因此
\[
 \int_Q|u(x)-u_Q|\,\mathrm dx
 \leq\dfrac1{m_d(Q)}\int_Q\int_Q|u(x)-u(y)|\,\mathrm dy\,\mathrm dx
 \leq\ell\sum_{i=1}^d\int_Q|\partial_i u|,
\]
因为 \(m_d(Q)=\ell^d\)。
\end{proof}

下一步比较两个点附近的一列局部均值。Hajłasz 的《Sobolev spaces on metric-measure spaces》\cite[Section 9, Theorem 9.4]{Hajlasz2003SobolevMetric}用这种逐半径的望远镜求和，从平均振荡得到双点控制；以下在欧氏立方体上给出所需证明。

\begin{proposition}\label{b1:prop:sobolev-maximal-difference}
设 \(w\in W^{1,1}(\mathbb R^d)\)，令 \(g=\sum_{i=1}^d|\partial_iw|\)。存在只依赖于 \(d\) 的常数 \(C_d\)，使在 \(w\) 的任意两个 Lebesgue 点 \(x,y\) 上，其 Lebesgue 值满足
\begin{equation}\label{b1:eq:sobolev-maximal-difference}
 |w(x)-w(y)|\leq C_d|x-y|\bigl(Mg(x)+Mg(y)\bigr).
\end{equation}
这里 \(M\) 是\secref{b1:sec:1303}的中心极大算子。
\end{proposition}
\begin{proof}
只需考虑 \(x\ne y\) 且两处极大函数都有限的情形。记 \(r=|x-y|\)，\(Q(x,t)=x+(-t,t)^d\)，并置 \(Q_j=Q(x,2^{-j}r)\)。由于立方体包含在可比半径的球内，Lebesgue 点性质给出 \(w_{Q_j}\rightarrow w(x)\)。由\eqref{b1:eq:sobolev-cube-oscillation}，
\[
 \begin{aligned}
 |w_{Q_{j+1}}-w_{Q_j}|
 &\leq\dfrac1{m_d(Q_{j+1})}\int_{Q_j}|w-w_{Q_j}|\\
 &\leq C_d\,2^{-j}r\,\dfrac1{m_d(Q_j)}\int_{Q_j}g
 \leq C_d\,2^{-j}r\,Mg(x).
 \end{aligned}
\]
最后一步使用 \(Q(x,t)\subseteq B(x,\sqrt d\,t)\) 及两者体积比只依赖维数。求和得到
\[
 |w(x)-w_{Q(x,r)}|\leq C_drMg(x).
\]
在 \(y\) 处也有同样的估计。

还需比较半径为 \(r\) 的两个均值。因为 \(Q(x,r)\) 和 \(Q(y,r)\) 都包含在 \(R=Q(x,2r)\) 中，体积比有统一界，故
\[
 \begin{aligned}
 |w_{Q(x,r)}-w_{Q(y,r)}|
 &\leq|w_{Q(x,r)}-w_R|+|w_{Q(y,r)}-w_R|\\
 &\leq\dfrac{C_d}{m_d(R)}\int_R|w-w_R|
 \leq C_drMg(x).
 \end{aligned}
\]
合并三项便得到\eqref{b1:eq:sobolev-maximal-difference}。
\end{proof}

极大函数在某个层集上有界时，上式就成为该层集上的 Lipschitz 估计。这里不需要 \(M\) 在 \(L^1\) 中强有界，只需要它对可积函数几乎处处有限，因此端点 \(p=1\) 也能纳入论证。

\begin{theorem}\label{b1:thm:sobolev-approximate-differentiability}
设 \(1\leq p\leq\infty\)，且 \(u\in W^{1,p}_{\mathrm{loc}}(\Omega)\)。则任一有限值可测代表都几乎处处近似可微，并且
\[
 \operatorname{ap}Du(x)h=\nabla u(x)\cdot h
 \quad(h\in\mathbb R^d)
\]
在几乎所有 \(x\in\Omega\) 上成立。
\end{theorem}

这里“任一代表”的量词是：每选定一个有限值可测代表，都存在一个可以依赖该代表的零测例外集，使结论在其外成立。不同代表只在零测集上不同，因而“几乎处处近似可微”作为等价类性质保持不变；这不表示所有代表在同一个共同点集上都近似可微。

\begin{proof}
局部 \(L^p\) 包含于局部 \(L^1\)，所以先把 \(u\) 看成 \(W^{1,1}_{\mathrm{loc}}\) 函数。固定 \(K\Subset\Omega\)，选 \(\chi\in C_c^\infty(\Omega)\)，使 \(\chi=1\) 于 \(K\) 的某个邻域。由光滑乘法公式，将 \(w=\chi u\) 及其弱导数补零到 \(\mathbb R^d\)，得到 \(w\in W^{1,1}(\mathbb R^d)\)。这一步的补零没有边界项，因为 \(\chi\) 紧支撑于 \(\Omega\)。

令 \(g=\sum_i|\partial_iw|\in L^1(\mathbb R^d)\)。由 Hardy--Littlewood 弱型不等式\ref{b1:thm:hl-weak-11}，\(Mg\) 几乎处处有限。从 \(K\) 删去零测例外，使剩下的点都是 \(w\) 的 Lebesgue 点，且所选的 \(u\) 值等于该点的 Lebesgue 值。再按 \(Mg(x)\leq n\) 分层，得可测集 \(E_n\)，其并覆盖 \(K\) 中几乎所有点。由\eqref{b1:eq:sobolev-maximal-difference}，\(u|_{E_n}\) 为 \(2C_dn\)-Lipschitz 函数。

每个 \(E_n\) 体积有限。由紧内正则性，可在其中选择可数个紧集，除去一个零测集后覆盖 \(E_n\)。再用可数个 \(K\Subset\Omega\) 覆盖 \(\Omega\)，就得到几乎处处覆盖定义域的紧集族 \((K_j)\)，每个 \(u|_{K_j}\) 都是 Lipschitz 的。由\cref{b1:thm:approximate-lipschitz-pieces}，\(u\) 几乎处处近似可微。

还要识别近似微分，而不只是证明它存在。把 \(u|_{K_j}\) 延拓为全空间 Lipschitz 函数 \(F_j\)。由定理\ref{b1:thm:lipschitz-weak-derivative}，\(F_j\) 的弱梯度等于其几乎处处经典梯度。又因 \(u-F_j=0\) 于 \(K_j\)，命题\ref{b1:prop:sobolev-gradient-on-level-sets}给出
\[
 \nabla u=\nabla F_j=DF_j
 \quad\text{几乎处处于 }K_j.
\]
其中最后一项按行向量与线性泛函的对应理解。在 \(K_j\) 的密度点且 \(F_j\) 可微的点上，沿 \(K_j\) 两函数相等，所以密度一集合判据给出 \(\operatorname{ap}Du=DF_j\)。合并这些等式，再合并可数个零测例外，便得所需识别。
\end{proof}

这个结论把弱导数与点态的一阶几何联系起来，但所得 Lipschitz 常数依赖于所选片，不能据此推出整个邻域上的 Lipschitz 性质。它也不声称 Sobolev 映射把每个零测集映成零测集：近似微分忽略定义域中的零测改值，像的测度却可能受到这些改值影响。以后将微分性质用于面积、迹或嵌入问题时，仍须分别检查那些问题需要的代表与额外条件。

% 本章小结（不计入编号）
\begin{chaptersummary}
Sobolev 空间的基本结构来自两个事实：各阶导数可以作为有限个 \(L^p\) 分量，而弱导数关系在这些分量的适当极限下保持。因此函数与导数不能独立指定，却可以一起使用 Banach 空间和弱收敛的方法。

\ProseParagraphBreak

运算规则首先是局部的。光滑逼近让经典公式进入弱导数框架，可积性条件再决定结果属于哪个全域空间。水平集上的梯度为零消除了截断函数折点处的不确定性；双 Lipschitz 坐标变换既控制导数，也保证零测集与积分估计正确转移。

\ProseParagraphBreak

对函数代表，一维绝对连续性、多维沿坐标线绝对连续性和局部 Lipschitz 性是不同层次。一般 Sobolev 函数几乎处处的近似微分等于弱梯度，但这不提供无条件的全域连续代表或边界值。下一章将处理全域光滑逼近、零边界空间与延拓，为以后讨论迹、嵌入和紧性准备工具。
\end{chaptersummary}


\BookExercises

% \chapref{b1:ch:21}习题临时稿；由章聚合文件已有的 BookExercises 入口统一接入。
% 十一题均为依据已读本章纲要、实际前章正文与已确认正文接口自拟的工具变式。
% 不冒认教材原题编号；不用第22章全局稠密性、第24章嵌入或尚未建立的迹理论。

\begin{exercise}\label{b1:ex:21-basic-definition-norm}
令 \(Q=(0,1)^2\)，\(u(x,y)=xy\)。对每个 \(1\leq p\leq\infty\)，直接从弱导数定义验证 \(u\in W^{1,p}(Q)\)，求出两个弱偏导数，并计算本章约定的 \(\|u\|_{W^{1,p}(Q)}\)。
\end{exercise}
% 来源：自拟；用于训练弱导数定义与本章逐分量范数，不调用后续嵌入或正则性。
\begin{solution}
取任意 \(\varphi\in C_c^\infty(Q)\)。由 Fubini 定理和一维分部积分，
\[
 \int_Qxy\,\partial_x\varphi(x,y)\,\mathrm dx\,\mathrm dy
 =\int_0^1y\left(\int_0^1x\,\partial_x\varphi(x,y)\,\mathrm dx\right)\mathrm dy
 =-\int_Qy\varphi(x,y)\,\mathrm dx\,\mathrm dy.
\]
测试函数在端点附近为零，所以没有边界项。交换两个变量得到
\[
 \partial_xu=y,\qquad \partial_yu=x
\]
的弱导数等式。三者在有界方形上都属于每个 \(L^p\)，故 \(u\in W^{1,p}(Q)\)。

若 \(1\leq p<\infty\)，则
\[
 \|u\|_p^p=\frac1{(p+1)^2},\qquad
 \|\partial_xu\|_p^p=\|\partial_yu\|_p^p=\frac1{p+1}.
\]
因此
\[
 \|u\|_{W^{1,p}(Q)}
 =\left(\frac1{(p+1)^2}+\frac2{p+1}\right)^{1/p}.
\]
当 \(p=\infty\) 时，\(u\)、\(x\) 与 \(y\) 的本性上确界都为 \(1\)，而本章范数取三者的最大值，故 \(\|u\|_{W^{1,\infty}(Q)}=1\)。
\end{solution}

\begin{exercise}\label{b1:ex:21-basic-vector-chain}
设 \(\Omega\subset\mathbb R^d\) 为开集、\(1\leq p\leq\infty\)，且
\[
 u=(u^1,u^2)\in W^{1,p}_{\mathrm{loc}}(\Omega;\mathbb R^2).
\]
令 \(F(s,t)=\sin s+\cos t\)。证明 \(F\circ u\in W^{1,p}_{\mathrm{loc}}(\Omega)\)，并对每个 \(1\leq i\leq d\) 求出 \(\partial_i(F\circ u)\)。
\end{exercise}
% 来源：自拟；直接训练向量值记号及正文链式法则的基础计算。
\begin{solution}
映射 \(F\in C^1(\mathbb R^2)\)，且
\[
 DF(s,t)=(\cos s,-\sin t)
\]
一致有界。由\cref{b1:thm:sobolev-chain}，复合函数属于所述局部 Sobolev 空间，并且
\[
 \partial_i(F\circ u)
 =\cos(u^1)\,\partial_i u^1
  -\sin(u^2)\,\partial_i u^2
 \quad\text{几乎处处于 }\Omega.
\]
右端属于 \(L^p_{\mathrm{loc}}\)，因为两个三角函数因子的绝对值不超过 \(1\)。局部陈述不要求 \(F\circ u\) 在无界域上具有全域 \(L^p\) 范数。
\end{solution}

\begin{exercise}\label{b1:ex:21-basic-derivative-embedding}
在 \(I=(0,1)\) 上令
\[
 u_n(x)=x+n^{-2}\sin(nx),\qquad u(x)=x.
\]
证明对每个 \(1\leq p\leq\infty\)，都有 \(u_n\to u\) 于 \(W^{1,p}(I)\)。再用\cref{b1:prop:sobolev-derivative-embedding}中的映射 \(\mathcal J\)，解释为什么这里的 Sobolev 收敛恰好等价于函数分量与一阶导数分量同时在 \(L^p(I)\) 中收敛。
\end{exercise}
% 来源：自拟；用于训练有限乘积嵌入和Sobolev范数收敛的直接判定。
\begin{solution}
两函数光滑，且
\[
 u_n-u=n^{-2}\sin(nx),\qquad
 u_n'-u'=n^{-1}\cos(nx).
\]
由于 \(m_1(I)=1\)，对每个有限 \(p\) 有
\[
 \|u_n-u\|_p\leq n^{-2},\qquad
 \|u_n'-u'\|_p\leq n^{-1},
\]
所以
\[
 \|u_n-u\|_{W^{1,p}(I)}
 \leq\bigl(n^{-2p}+n^{-p}\bigr)^{1/p}\longrightarrow0.
\]
在 \(p=\infty\) 时，本章范数取两个分量范数的最大值，同样有
\[
 \|u_n-u\|_{W^{1,\infty}(I)}\leq n^{-1}\longrightarrow0.
\]
这里 \(\mathcal Jv=(v,v')\)，而 \(Y_p\) 的乘积范数正按有限 \(p\) 的 \(p\) 次和或无穷端点的最大值定义。因此 \(\|\mathcal J(u_n-u)\|_{Y_p}=\|u_n-u\|_{W^{1,p}}\)，所述两个分量同时收敛与 Sobolev 范数收敛完全等价。
\end{solution}

\begin{optionalexercise}\label{b1:ex:21-scaling-concentration}
设整数 \(d\ge2\)、\(1\le p<d\)，记 \(B=B(0,1)\)，取非零 \(\psi\in C_c^\infty(B)\)，并令
\[
 u_r(x)=r^{1-d/p}\psi(x/r),\quad 0<r<1.
\]
计算任意多指标 \(\alpha\) 对应的 \(\|\partial^\alpha u_r\|_{L^p(B)}\)。证明 \(u_r\to0\) 于 \(L^p(B)\)，其一阶导数范数保持不变，而 \(\|u_r\|_\infty\to\infty\)。据此证明，不存在独立于 \(u\) 的常数 \(C\)，使
\[
 \|u\|_{L^\infty(B)}
 \le C\|u\|_{W^{1,p}(B)}
 \quad\bigl(u\in C_c^\infty(B)\bigr).
\]
\end{optionalexercise}
% 来源：自拟；由弱导数与线性换元的尺度计算构造集中列，只排除p<d的所列估计，不预用嵌入定理。
\begin{solution}
函数 \(u_r\) 的支集包含在 \(B(0,r)\) 中，经典求导与换元 \(x=ry\) 给出
\[
 \begin{aligned}
 \|\partial^\alpha u_r\|_p^p
 &=r^{p-d-p|\alpha|}\int_{B(0,r)}
       |(\partial^\alpha\psi)(x/r)|^p\,\mathrm dx\\
 &=r^{p(1-|\alpha|)}\|\partial^\alpha\psi\|_p^p.
 \end{aligned}
\]
因此
\[
 \|\partial^\alpha u_r\|_p
 =r^{1-|\alpha|}\|\partial^\alpha\psi\|_p.
\]
当 \(|\alpha|=0\) 时右端趋零；当 \(|\alpha|=1\) 时它不随 \(r\) 改变。由于非零紧支集光滑函数不可能在连通的全空间上为常数，至少一个一阶导数的范数为正。故这列函数的 \(W^{1,p}\) 范数一致有界，而一阶导数并没有一同趋零。

另一方面，\(\psi\) 连续，故其上确界等于本性上确界，且
\[
 \|u_r\|_\infty=r^{1-d/p}\|\psi\|_\infty\longrightarrow\infty,
\]
因为 \(p<d\)。若题中统一估计成立，左端便会被一致有界的 Sobolev 范数控制，矛盾。这个例子只处理严格不等式 \(p<d\)；当 \(p=d\) 时，同一伸缩的振幅不再增长，不能从本题推出临界指数的结论。
\end{solution}

\begin{optionalexercise}\label{b1:ex:21-endpoint-structure}
\begin{enumerate}
\item\label{b1:item:21-w11-no-weak-subsequence}
在 \(I=(0,1)\) 上令 \(v_n(x)=\min(nx,1)\)。证明 \((v_n)\) 在 \(W^{1,1}(I)\) 中有界，但没有弱收敛子列。
\item\label{b1:item:21-w1infty-nonseparable}
对 \(t\in(0,1)\)，令 \(w_t(x)=|x-t|\)。证明这些元素在 \(W^{1,\infty}(I)\) 中两两距离至少为 \(2\)，并据此证明该空间不可分。
\end{enumerate}
\end{optionalexercise}
% 来源：自拟；第10章端点弱测试与本章连续导数映射的具体应用，完整折点构造未在正文展开。
\begin{solution}
第一问中，\(v_n\) 为 Lipschitz 函数，弱导数几乎处处等于经典导数，所以
\[
 v_n'=n\mathbf1_{(0,1/n)},\quad
 \|v_n\|_1\le1,\quad \|v_n'\|_1=1.
\]
这证明了 Sobolev 范数有界。

假设某个子列 \((v_{n_j})\) 在 \(W^{1,1}(I)\) 中弱收敛。由\cref{b1:prop:sobolev-derivative-embedding}，导数映射从 \(W^{1,1}(I)\) 到 \(L^1(I)\) 有界线性，相应导数必在 \(L^1(I)\) 中弱收敛到某个 \(g\)。固定 \(a\in(0,1)\)，以及任意在 \((0,a)\) 上几乎处处为零的 \(h\in L^\infty(I)\)。当 \(n_j>1/a\) 时，\(\int_Iv_{n_j}'h=0\)，故 \(\int_Igh=0\)。这些有界测试说明 \(g=0\) 几乎处处于 \((a,1)\)。取 \(a=1/k\)、\(k\ge2\)，便得 \(g=0\) 几乎处处于 \(I\)。但常函数 \(h=1\) 也是允许的对偶测试，且
\[
 \int_Ig=\lim_j\int_Iv_{n_j}'=1,
\]
矛盾。因此不存在所假设的弱收敛子列。事实上，\(v_n\to1\) 于 \(L^1(I)\)，失去弱紧性的是向端点集中的导数。

第二问中，\(w_t'=\operatorname{sgn}(x-t)\) 几乎处处。若 \(s<t\)，则在正测度区间 \((s,t)\) 上有
\[
 w_s'=1,\quad w_t'=-1.
\]
所以 \(\|w_s'-w_t'\|_\infty=2\)，从而 \(\|w_s-w_t\|_{W^{1,\infty}}\ge2\)。

若空间有可数稠密子集 \(\{z_j\}\)，则每个开球 \(B(w_t,1/2)\) 至少包含一个 \(z_j\)。这些球两两不交，否则三角不等式将使两个中心的距离小于 \(1\)。这样便给不可数多个球分别配上不同的可数集元素，矛盾。故 \(W^{1,\infty}(I)\) 不可分。
\end{solution}

\begin{exercise}\label{b1:ex:21-oscillatory-gradient-energy}
在 \(I=(0,2\pi)\) 上令 \(u_n(x)=\sin(nx)/n\)。证明 \(u_n\rightharpoonup0\) 于 \(H^1(I)\)，但不依 \(H^1\) 范数趋零。进一步证明，对任意 \(g\in C([0,2\pi])\)，
\[
 \lim_{n\to\infty}\int_0^{2\pi}g(x)|u_n'(x)|^2\,\mathrm dx
 =\frac12\int_0^{2\pi}g(x)\,\mathrm dx.
\]
解答不预借 Fourier 理论。
\end{exercise}
% 来源：自拟；与第20章分布收敛例相比，新增H1弱极限识别及导数能量的非零极限。
\begin{solution}
先证明 \(\cos(nx)\rightharpoonup0\) 于 \(L^2(I)\)。对一个子区间 \((a,b)\subset I\)，
\[
 \left|\int_a^b\cos(nx)\,\mathrm dx\right|
 =\frac{|\sin(nb)-\sin(na)|}{n}\le\frac2n.
\]
因此与任意有限区间阶梯函数相乘后的积分都趋于零。这些阶梯函数在 \(L^2(I)\) 中稠密：先用\secref{b1:sec:1002}的连续紧支集函数逼近，再在有限区间上利用一致连续性作阶梯逼近。又因 \(\|\cos(nx)\|_2=\sqrt\pi\)，对 \(h\in L^2(I)\) 及一个阶梯函数 \(s\)，有
\[
 \left|\int_I\cos(nx)h\right|
 \le\left|\int_I\cos(nx)s\right|
     +\sqrt\pi\|h-s\|_2.
\]
先选 \(s\)，再令 \(n\to\infty\)，即得所需弱收敛。

现在 \(\|u_n\|_2=\sqrt\pi/n\to0\)，而 \(u_n'=\cos(nx)\rightharpoonup0\) 于 \(L^2\)。由\cref{b1:prop:sobolev-weak-components}，\(u_n\rightharpoonup0\) 于 \(H^1\)。也可以直接对任意 \(v\in H^1(I)\) 测试内积：零阶积分由 \(\|u_n\|_2\to0\) 控制，一阶积分由刚证的弱收敛控制。不过
\[
 \|u_n'\|_2^2=\pi
\]
始终不变，所以不可能强收敛到零。

最后使用恒等式 \(\cos^2(nx)=\bigl(1+\cos(2nx)\bigr)/2\)。函数 \(g\) 属于 \(L^2(I)\)，上面的弱收敛结论同样适用于 \(\cos(2nx)\)，于是
\[
 \int_Ig|u_n'|^2
 =\frac12\int_Ig+\frac12\int_Ig\cos(2nx)
 \longrightarrow\frac12\int_Ig.
\]
函数值的强收敛与导数的弱收敛都不能使导数平方直接趋向极限导数的平方。
\end{solution}

\begin{exercise}\label{b1:ex:21-product-integrability-loss}
在 \(B=B(0,1)\subset\mathbb R^3\) 上取 \(u(x)=|x|^{-1/3}\)，原点处任意赋一个有限值。证明
\[
 u\in H^1(B),\quad
 u^2\in L^2(B)\cap W^{1,1}(B),\quad
 u^2\notin H^1(B).
\]
说明这与有界 Sobolev 函数的乘积法则并不矛盾。
\end{exercise}
% 来源：自拟；根正文只说明W1p一般非代数，本题独立计算高维奇性，用ACL反向判据确认弱导数。
\begin{solution}
在原点外，经典梯度为
\[
 \nabla u(x)=-\frac13|x|^{-7/3}x,\quad
 \nabla(u^2)(x)=-\frac23|x|^{-8/3}x.
\]
所以其大小分别为 \(\tfrac13|x|^{-4/3}\) 与 \(\tfrac23|x|^{-5/3}\)。由三维径向积分公式，\(|x|^{-q}\) 在单位球上可积，当且仅当 \(q<3\)，因为相应径向积分是常数 \(4\pi\) 乘以 \(\int_0^1r^{2-q}\,\mathrm dr\)。因此
\[
 u\in L^2(B),\quad |\nabla u|\in L^2(B),\quad
 u^2\in L^2(B),\quad |\nabla(u^2)|\in L^1(B).
\]
这里所用的奇性指数依次为 \(2/3\)、\(8/3\)、\(4/3\)、\(5/3\)，都小于 \(3\)。

还须确认经典梯度的候选确实是弱梯度，不能只凭原点外的计算略去这一点。对任一坐标方向，除了横向参数为零、因而穿过原点的那一条直线外，函数 \(u\) 和 \(u^2\) 在相应球内截线上都光滑，并在每个紧子区间上绝对连续。异常直线的参数集在 \(\mathbb R^2\) 中为零测集。因此两函数都具有 ACL 性质。结合已算出的函数和候选偏导数的可积性，\cref{b1:thm:sobolev-acl}的反向判据给出 \(u\in H^1(B)\) 与 \(u^2\in W^{1,1}(B)\)。

然而
\[
 \int_B|\nabla(u^2)|^2\,\mathrm dx
 =\frac{16\pi}{9}\int_0^1r^{-4/3}\,\mathrm dr=\infty.
\]
若 \(u^2\) 另有一个 \(L^2\) 弱梯度，由弱导数唯一性，它必须与已经得到的 \(L^1\) 弱梯度几乎处处相同，这与上式矛盾。因此 \(u^2\notin H^1(B)\)。

函数 \(u\) 并不本性有界。乘积的弱导数依旧满足 \(\nabla(u^2)=2u\nabla u\)，但两个 \(L^2\) 因子的乘积由 Hölder 不等式只能统一保证属于 \(L^1\)，不能保证仍属于 \(L^2\)。失去的是原有可积指数，而不是乘积求导公式本身。
\end{solution}

\begin{exercise}\label{b1:ex:21-positive-part-continuity}
设 \(\Omega\subset\mathbb R^d\) 为开集、\(1\le p<\infty\)，实值函数
\(u_n\to u\)、\(v_n\to v\) 于 \(W^{1,p}(\Omega)\)。利用正文的正部强连续性证明
\[
 \max\{u_n,v_n\}\longrightarrow\max\{u,v\},\qquad
 \min\{u_n,v_n\}\longrightarrow\min\{u,v\}
\]
于同一 Sobolev 空间。再在 \((-1,1)\) 上用 \(u_n(x)=x-1/n\)、\(u(x)=x\)，说明正部映射在 \(W^{1,\infty}\) 中不连续。
\end{exercise}
% 来源：自拟；由正文正部强连续性推出上下格运算的强连续性，并保留无穷端点反例。
\begin{solution}
差函数满足 \(u_n-v_n\to u-v\) 于 \(W^{1,p}(\Omega)\)。由\cref{b1:prop:sobolev-lattice-strong-continuity}及恒等式
\[
 \max\{a,b\}=b+(a-b)_+,\qquad
 \min\{a,b\}=a-(a-b)_+,
\]
分别把右端的加法、减法和正部运算通过极限，便得到两个结论。

无穷指数时，给定的 \(u_n-u=-1/n\)，且二者一阶导数都为 \(1\)，故 \(u_n\to u\) 于 \(W^{1,\infty}\bigl((-1,1)\bigr)\)。但在 \((0,1/n)\) 上，\(((u_n)_+)'=0\)、\((u_+)'=1\)，所以
\[
 \|((u_n)_+)'-(u_+)'\|_\infty=1.
\]
区间虽不断缩小，本性上确界仍能检测其中的完整跳差。这正是有限 \(p\) 的控制收敛论证不能照搬的地方。
\end{solution}

\begin{optionalexercise}\label{b1:ex:21-coordinate-change-boundaries}
\begin{enumerate}
\item\label{b1:item:21-inverse-lipschitz-required}
映射 \(\Phi(x)=x^3\) 是 \(I=(-1,1)\) 到自身的光滑 Lipschitz 双射。取坐标函数 \(u(x)=x\)，证明 \(u\in H^1(I)\)，但 \(u\circ\Phi^{-1}\notin H^1(I)\)。
\item\label{b1:item:21-bilipschitz-not-second-order}
令 \(\Psi(x)=x+\tfrac12|x|\)。证明它是 \(I=(-1,1)\) 到 \(J=(-1/2,3/2)\) 的双 Lipschitz 映射，却存在一个函数 \(v\)，它属于所有 \(W^{2,p}(J)\)、\(1\le p\le\infty\)，但 \(v\circ\Psi\) 不属于任何 \(W^{2,p}(I)\)。
\end{enumerate}
\end{optionalexercise}
% 来源：自拟；用退化坐标与折线坐标分别说明逆Lipschitz条件和高阶坐标正则性不能省略。
\begin{solution}
第一问中，\(|\Phi'(x)|=3x^2\le3\)，所以 \(\Phi\) 为 Lipschitz 映射；严格单调性及端点极限说明它是题中双射。其逆映射在零点附近并非 Lipschitz，因为
\[
 \frac{|\Phi^{-1}(h)-\Phi^{-1}(0)|}{|h|}=|h|^{-2/3}\longrightarrow\infty.
\]
坐标函数 \(u\) 及其一阶导数都平方可积。复合函数则为
\[
 (u\circ\Phi^{-1})(x)=\operatorname{sgn}(x)|x|^{1/3}.
\]
它在原点外的经典导数是 \(\tfrac13|x|^{-2/3}\)。若有平方可积的弱导数，该弱导数在两个不含原点的开区间上必须与经典导数几乎处处相同，但
\[
 \int_0^1\left|\frac13x^{-2/3}\right|^2\,\mathrm dx=\infty,
\]
所以不存在这样的弱导数。正向映射光滑且 Lipschitz，并不足以控制经过其逆映射复合后的 Sobolev 范数。

第二问中，\(\Psi\) 在负半轴的斜率为 \(1/2\)，在正半轴的斜率为 \(3/2\)。对任意 \(x<y\)，分段积分给出
\[
 \frac12(y-x)\le\Psi(y)-\Psi(x)\le\frac32(y-x).
\]
故它及其逆映射都 Lipschitz，端点极限给出像恰为 \(J\)。

取 \(v(y)=y\)。在有界区间 \(J\) 上，\(v\)、\(v'=1\)、\(v''=0\) 属于所有所列 \(L^p\) 空间，因此 \(v\in W^{2,p}(J)\)。复合 \(v\circ\Psi=\Psi\) 的弱一阶导数为
\[
 \Psi'=\frac12\mathbf1_{(-1,0)}+\frac32\mathbf1_{(0,1)}.
\]
这个导数在原点的跳跃为 \(1\)。由\cref{b1:ex:20-piecewise-jump-derivatives}的跳跃公式，或直接在原点两侧分部积分，
\[
 \partial^2T_\Psi=\delta_0.
\]
点质量不是正则分布，因此 \(\Psi\) 没有局部可积的二阶弱导数，更不属于任何所列 \(W^{2,p}(I)\)。双 Lipschitz 条件适用于本章的一阶坐标变换，不能自动提升为任意阶的结论。
\end{solution}

\begin{exercise}\label{b1:ex:21-acl-representative-choices}
在 \(\Omega=(-1,1)^2\) 上考虑零函数以及
\[
 v(x,y)=\mathbf1_{\mathbb Q}(x),\quad
 w(x,y)=\mathbf1_{\{(0,0)\}}(x,y).
\]
证明三者属于同一个 \(W^{1,p}(\Omega)\) 等价类，\(1\le p\le\infty\)，但 \(v\) 在每条水平线上都不是绝对连续函数，\(w\) 却仍具有同时针对两个坐标方向的 ACL 性质。解释为何这既不违背共同 ACL 代表的存在性，也不违背一维绝对连续代表的唯一性。最后说明两次零集改值不改变几乎处处的近似微分。
\end{exercise}
% 来源：自拟；从零集改值推进到ACL的逐方向量词及高维代表非逐点唯一性，不重复正文普通可微反例。
\begin{solution}
集合 \(\mathbb Q\cap(-1,1)\) 为一维零测集，由 Fubini 定理，\(v\) 的非零集合在 \(\Omega\) 中为二维零测集；\(w\) 的非零集合也为零测集。因此两者都几乎处处等于零，定义中的函数积分与弱导数积分均不变。它们代表的是零 Sobolev 元素，所有弱偏导数也都是零等价类。

固定任意 \(y\in(-1,1)\)，函数 \(x\mapsto v(x,y)\) 是有理点的指示函数。每个区间都同时含有有理点与无理点，故该截面处处不连续，特别不绝对连续。因此 \(v\) 连“几乎每条水平线”上的绝对连续性都不满足。它在竖直线上为常数，并不能补救另一个坐标方向的失败。

对 \(w\)，若 \(y\ne0\)，则整个水平截面恒为零；若 \(x\ne0\)，则整个竖直截面也恒为零。每个方向只须排除参数为零的那一条直线。因此 \(w\) 是一个同时满足两个方向 ACL 条件的代表，却在原点不等于处处为零的代表。共同 ACL 代表的存在性并没有说每个任意改值的代表都可以使用，也没有断言高维 ACL 代表处处唯一。

一维绝对连续代表的唯一性没有失效。在原点所在的那两条坐标线上，\(w\) 的截面并不连续，因而它们正是 ACL 条件允许排除的异常截面，不能拿来与零截面应用一维唯一性。在其余合格截线上，两者确实处处相同。

最后设 \(E\) 为某一次改值的二维零测集，取 \(z\notin E\)。改值后的函数在 \(z\) 取值为零，在 \(E\) 外也处处为零。对任意 \(\eta>0\)，
\[
 \{\,h:0<|h|<r,\ |f(z+h)-f(z)|>\eta|h|\,\}
 \subseteq E-z.
\]
右侧在每个球中都为零测集，所以左侧的相对测度为零。这就说明近似微分在 \(z\) 存在且为零。仅有改值集上的基点可能不满足这个论证，而它本身是零测集，故几乎处处近似微分保持不变。
\end{solution}

\begin{optionalexercise}\label{b1:ex:21-highest-derivative-kernel}
设 \(\Omega\subset\mathbb R^d\) 为非空连通开集，整数 \(k\ge1\)，\(1\le p\le\infty\)，且 \(u\in W^{k,p}_{\mathrm{loc}}(\Omega)\)。证明以下两条等价：
\begin{equivalences}
\item\label{b1:item:21-highest-derivatives-zero}
所有 \(|\alpha|=k\) 阶弱导数 \(\partial^\alpha u\) 都为零。
\item\label{b1:item:21-polynomial-representative}
函数 \(u\) 几乎处处等于一个总次数不超过 \(k-1\) 的多项式。
\end{equivalences}
去掉连通性后，应怎样修改第二条？
\end{optionalexercise}
% 来源：自拟；局部磨光、有限维闭性与连通拼合的综合应用，不预借Poincare不等式或全域稠密定理。
\begin{solution}
多项式的经典导数也是弱导数，所以第二条直接蕴涵第一条。现假设第一条成立。

固定开球 \(B=B(x_0,r)\)，使 \(\overline{B(x_0,2r)}\subset\Omega\)。由分布卷积的微分公式，当 \(\varepsilon>0\) 足够小时，局部磨光 \(u_\varepsilon=u*\rho_\varepsilon\) 在 \(B\) 的邻域光滑，并满足
\[
 \partial^\alpha u_\varepsilon=0\quad(|\alpha|=k).
\]
对 \(x\in B\)，沿连接 \(x_0\) 与 \(x\) 的线段令 \(F(t)=u_\varepsilon\bigl(x_0+t(x-x_0)\bigr)\)。链式求导表明 \(F^{(k)}=0\)，因为展开后的每项都含一个 \(k\) 阶偏导数。由一维 Taylor 公式，
\[
 u_\varepsilon(x)
 =\sum_{|\alpha|<k}
     \frac{\partial^\alpha u_\varepsilon(x_0)}{\alpha!}(x-x_0)^\alpha.
\]
因此每个 \(u_\varepsilon|_B\) 都是总次数不超过 \(k-1\) 的多项式的限制。

函数 \(u\) 在 \(\overline B\) 的邻域局部可积，即使 \(p=\infty\) 也如此。由\cref{b1:prop:local-smoothing-approximation}的 \(L^1\) 版本，\(u_\varepsilon\to u\) 于 \(L^1(B)\)。固定次数多项式在 \(B\) 上的限制构成一个有限维线性子空间；其 \(L^1(B)\) 范数确实是范数，因为连续多项式若几乎处处为零，便在开球上处处为零，从而是零多项式。由有限维子空间的闭性，\(u|_B\) 几乎处处等于某个次数不超过 \(k-1\) 的多项式 \(P_B\)。这里所用的逼近只有局部 \(L^1\) 收敛，不需要无穷范数的磨光收敛。

若两个这样的球 \(B_1,B_2\) 相交，则 \(P_{B_1}-P_{B_2}\) 在非空开集 \(B_1\cap B_2\) 上几乎处处为零。连续性使其在那里处处为零，多项式恒等性进而给出 \(P_{B_1}=P_{B_2}\) 作为全空间多项式。为看出后一事实，可在交集内取一个小球；多项式在那里为零，其在球心的全部偏导数也为零，而有限 Taylor 展开恰是这个多项式本身。

上述小球覆盖 \(\Omega\)。固定其中一个球，取所有能够通过有限条相交小球链与它连接的球的并集。这个并集开，补集也是其余链类的球的并集，因而也开。由 \(\Omega\) 连通，只能有一个链类，所以全部 \(P_B\) 都是同一个多项式 \(P\)。再由开覆盖的可数子覆盖性，球内各自的零测例外可以并成一个零测集，从而 \(u=P\) 几乎处处于整个 \(\Omega\)。

若 \(\Omega\) 不连通，则在每个连通分支上分别得到一个总次数不超过 \(k-1\) 的多项式，不同分支间不要求系数相同。反过来，按分支给出的这类函数具有所需的局部正则性。确实，若 \(K\Subset\Omega\)，则 \(\delta=\operatorname{dist}(K,\Omega^c)>0\)。对每个 \(x\in K\)，连通球 \(B(x,\delta/2)\) 包含于 \(\Omega\)，因而完全落在 \(x\) 所属的连通分支中；由紧性可取有限个这样的球覆盖 \(K\)，所以 \(K\) 只遇到有限个分支。分支边界不在 \(\Omega\) 内，也不会产生内部跳跃项。
\end{solution}

% !TeX root = ../../Book1.tex
\chapter{Sobolev 函数的逼近、局部化与延拓}
\label{b1:ch:22}
\BookChapterNumber{b1:ch:22}
\begin{chapterabstract}
本章研究如何把 Sobolev 函数换成便于计算的函数，同时保留弱导数与范数控制。内部紧支集允许补零，并在有限指数下允许磨光；逐层分解使任意开集上的局部磨光成为全域光滑逼近。无穷指数下卷积仍有范数界，却一般不在 Sobolev 范数中收敛。要求逼近函数具有紧支集，则产生零边界空间，并使边界几何进入问题。最后通过半空间反射、Lipschitz 图表及单位分解构造有界延拓，再用截断推导一个基本的局部能量估计。
\end{chapterabstract}

\begin{learninggoals}
\begin{enumerate}
\item\label{b1:goal:22-approximation}在 \(1\le p<\infty\) 的范围内证明内部紧支集与全空间的光滑逼近，掌握 Meyers--Serrin 定理中磨光尺度、支集和可和误差的配合，并辨认 \(p=\infty\) 时强逼近失效的原因。
\item\label{b1:goal:22-zero-boundary}理解 \(W_0^{1,p}\) 的闭包定义，区分补零、零端点值和内部光滑逼近，并辨认无穷指数的额外障碍。
\item\label{b1:goal:22-extension}通过沿线绝对连续性证明半空间反射，再把反射放入 Lipschitz 边界图表，构造固定的有界线性延拓。
\item\label{b1:goal:22-localization}追踪截断导数对估计常数的影响，用紧支集逼近合法化 Sobolev 测试函数，并推导局部梯度能量界。
\end{enumerate}
\end{learninggoals}

\chapref{b1:ch:21}的完备性、弱乘积公式、沿坐标线绝对连续性与双 Lipschitz 坐标变换，是本章的直接前置。光滑单位分解与磨光函数来自\chapref{b1:ch:20}。证明安排先处理开集内部，再处理边界；这样可以看清哪些结论对任意开集成立，哪些结论必须使用边界图表。

阅读时应把三种要求分开：在域内光滑、具有内部紧支集、能够延拓到边界之外。常数函数在有界区间上已经光滑，却不能由内部紧支集函数依一阶 Sobolev 范数逼近；有界 Lipschitz 区域上的函数可以受控延拓，却仍可能在 \(W^{1,\infty}\) 范数下无法磨光逼近。这些差别并不是符号细节，而会决定后续边界条件的含义。

Meyers 与 Serrin 的《\(H=W\)》\cite{MeyersSerrin1964HW}是全域光滑稠密性的原始文献，篇幅很短，适合在读完\secref{b1:sec:2201}后对照。需要把本章放回 Sobolev 空间的系统理论时，可伴读 Adams 与 Fournier 的《Sobolev Spaces》\cite{Adams2003Sobolev}及 Brézis 的《Functional Analysis, Sobolev Spaces and Partial Differential Equations》\cite{Brezis2011Functional}。延拓中的几何组织另在\secref{b1:sec:2203}说明具体参考范围。

% !TeX root = ../../Book1.tex
\section{光滑逼近}
\label{b1:sec:2201}

\chapref{b1:ch:21}已用局部卷积建立 Sobolev 函数的运算法则。现在要把这种局部逼近接成一个在全域范数下有效的光滑函数。困难不在卷积本身，而在于一个固定半径的核可能越过边界，许多个局部误差又未必能够相加。本节为不同位置选择不同的磨光尺度，同时控制支集和误差。

Meyers 与 Serrin 的短文《\(H=W\)》\cite[pp. 1055--1056]{MeyersSerrin1964HW}证明了任意开集上的光滑稠密性。原文的 \(H^{k,p}\) 指域内光滑函数在 \(W^{k,p}\) 中的闭包，不是本书专用于 \(W^{k,2}\) 的记号 \(H^k\)。以下直接陈述稠密性，不另外引入同名空间。

% 实读来源：MeyersSerrin1964HW，PNAS 51(6) (1964), pp.1055--1056，
% DOI 10.1073/pnas.51.6.1055，PMCID PMC300210；Europe PMC正式扫描副本两页均已阅读。
% p.1055：任意开集、复值、k>=0、1<=p<infty，原文用各导数Lp范数之和；
% p.1056：局部Sobolev函数的全局小误差逼近、环层单位分解及逐层卷积证明。
% 本节用已证完备性处理可和误差级数，补明紧支集补零、支集膨胀与局部有限性；
% 原文末段关于光滑到边界的区别仅作范围说明，不改成本书的紧支集稠密性结论。

\subsection{局部卷积}
\label{b1:subsec:220101}

设 \(\Omega\subseteq\mathbb R^d\) 为非空开集，\(k\in\mathbb N_0\)。仍使用\secref{b1:sec:2003}的非负偶磨光函数 \(\rho\)，其支集包含于 \(\overline B(0,1)\)，积分等于 \(1\)，并记 \(\rho_\varepsilon(x)=\varepsilon^{-d}\rho(x/\varepsilon)\)。对 \(u\in W^{k,p}_{\mathrm{loc}}(\Omega)\)，卷积 \(u*\rho_\varepsilon\) 只先定义在
\[
 \Omega_\varepsilon
 =\{\,x\in\Omega\mid\operatorname{dist}(x,\Omega^c)>\varepsilon\,\}.
\]
由\cref{b1:lem:sobolev-local-smoothing}，在任意固定内部相对紧开集上，这些卷积在有限 \(p\) 的 Sobolev 范数中趋于 \(u\)，并与不超过 \(k\) 阶的弱导数交换。不能在尚未定义 \(u\) 的地方替它指定零值，再假定这些导数关系仍成立。

若函数的非零部分已经留在内部紧集上，补零却没有这个问题。下面将“具有内部紧支集”明确解释为：存在 \(K\Subset\Omega\)，使函数在 \(\Omega\setminus K\) 上几乎处处为零。

\begin{lemma}\label{b1:lem:compact-sobolev-zero-extension}
设 \(1\leq p\leq\infty\)、\(u\in W^{k,p}(\Omega)\)，且 \(u\) 具有内部紧支集 \(K\)。在 \(\Omega\) 外将 \(u\) 补为零，所得函数记为 \(\widetilde u\)。则 \(\widetilde u\in W^{k,p}(\mathbb R^d)\)，并且
\begin{equation}\label{b1:eq:compact-sobolev-zero-extension}
 \partial^\alpha\widetilde u=\widetilde{\partial^\alpha u}
 \quad(|\alpha|\leq k),
 \quad
 \|\widetilde u\|_{W^{k,p}(\mathbb R^d)}
 =\|u\|_{W^{k,p}(\Omega)}.
\end{equation}
\end{lemma}
\begin{proof}
由弱导数的局部性及唯一性，\(\partial^\alpha u=0\) 几乎处处于 \(\Omega\setminus K\)。选取 \(\chi\in C_c^\infty(\Omega)\)，使其在 \(K\) 的一个邻域内等于 \(1\)；这样的截断由\cref{b1:lem:smooth-cutoff}给出。

任取 \(\varphi\in C_c^\infty(\mathbb R^d)\)。函数 \(\chi\varphi\) 是 \(\Omega\) 中的测试函数，而在 \(K\) 附近有 \(\partial^\alpha(\chi\varphi)=\partial^\alpha\varphi\)。因此
\[
 \begin{aligned}
 \int_{\mathbb R^d}\widetilde u\,\partial^\alpha\varphi\,\mathrm dx
 &=\int_\Omega u\partial^\alpha(\chi\varphi)\,\mathrm dx\\
 &=(-1)^{|\alpha|}\int_\Omega\partial^\alpha u\,\chi\varphi\,\mathrm dx\\
 &=(-1)^{|\alpha|}\int_{\mathbb R^d}
        \widetilde{\partial^\alpha u}\,\varphi\,\mathrm dx.
 \end{aligned}
\]
候选导数属于 \(L^p(\mathbb R^d)\)，于是得到弱导数等式。补零不改变任何分量的 \(L^p\) 范数，对有限 \(p\) 求 \(p\) 次和、对无穷端点取最大值，即得范数等式。
\end{proof}

特别地，若只有 \(u\in W^{k,p}_{\mathrm{loc}}(\Omega)\)，而 \(\psi\in C_c^\infty(\Omega)\)，则 \(\psi u\) 仍可使用这个引理。由\cref{b1:prop:weak-leibniz}，
\[
 \partial^\alpha(\psi u)
 =\sum_{\beta\leq\alpha}\binom\alpha\beta
       (\partial^{\alpha-\beta}\psi)(\partial^\beta u).
\]
每个光滑系数都在同一个内部紧集上有界，其支集外为零，所以每项全局属于 \(L^p(\Omega)\)。由此 \(\psi u\in W^{k,p}(\Omega)\)，且具有内部紧支集。这是随后使用单位分解的基本步骤。

\subsection{Friedrichs 磨光函数}
\label{b1:subsec:220102}

Friedrichs 磨光使用上述光滑核把函数换成邻域内的加权平均。\secref{b1:sec:2003}已经解释了这一操作为何使分布光滑；这里还要保留全部 Sobolev 范数估计。

\begin{proposition}\label{b1:prop:sobolev-friedrichs-smoothing}
若 \(u\in W^{k,p}(\mathbb R^d)\)、\(1\leq p\leq\infty\)，则
\(u_\varepsilon=\rho_\varepsilon*u\) 属于 \(C^\infty(\mathbb R^d)\cap W^{k,p}(\mathbb R^d)\)，且
\[
 \partial^\alpha u_\varepsilon
 =\rho_\varepsilon*\partial^\alpha u,
 \quad
 \|u_\varepsilon\|_{W^{k,p}(\mathbb R^d)}
 \leq\|u\|_{W^{k,p}(\mathbb R^d)}.
\]
若 \(p<\infty\)，还有 \(u_\varepsilon\to u\) 于 \(W^{k,p}(\mathbb R^d)\)。
\end{proposition}
\begin{proof}
光滑性和导数交换由\cref{b1:prop:distribution-convolution-smooth}成立。对 \(f\in L^p(\mathbb R^d)\)、\(p<\infty\)，在概率测度 \(\rho_\varepsilon(y)\,\mathrm dy\) 上使用 Hölder 不等式，得到
\[
 \bigl|\rho_\varepsilon*f(x)\bigr|^p
 \leq\int_{\mathbb R^d}\rho_\varepsilon(y)|f(x-y)|^p\,\mathrm dy.
\]
对 \(x\) 积分，Tonelli 定理与平移不变性给出 \(\|\rho_\varepsilon*f\|_p\leq\|f\|_p\)。当 \(p=\infty\) 时，非负核与单位积分直接给出同一估计。将它应用于有限多个导数，便得 Sobolev 范数界。

若 \(p<\infty\)，由\cref{b1:prop:approximate-identity-lp}，每个 \(\rho_\varepsilon*\partial^\alpha u\) 都在 \(L^p\) 中趋于 \(\partial^\alpha u\)。这些误差的 \(p\) 次方只有有限项，其和也趋零，得到所需范数收敛。
\end{proof}

无穷端点的结论只有范数不增，不能把最后一句也扩展到 \(p=\infty\)。\secref{b1:sec:2103}已用 \(|x|\) 的导数跳跃说明：即使在一个有界区间内部，光滑逼近也未必在 \(W^{1,\infty}\) 范数中成立。

\begin{proposition}\label{b1:prop:compact-sobolev-smooth-approximation}
设 \(1\leq p<\infty\)、\(u\in W^{k,p}(\Omega)\)，且 \(u\) 在 \(\Omega\setminus K\) 上几乎处处为零，其中 \(K\Subset\Omega\)。对任意包含 \(K\) 的开集 \(V\subseteq\Omega\)，当 \(\varepsilon\) 充分小时，
\[
 u_\varepsilon=\rho_\varepsilon*\widetilde u\in C_c^\infty(V),
 \quad
 \|u_\varepsilon-u\|_{W^{k,p}(\Omega)}\longrightarrow0.
\]
更精确地，\(\operatorname{supp}u_\varepsilon\subseteq K+\overline B(0,\varepsilon)\)，其中右侧表示两集合中向量的所有和。
\end{proposition}
\begin{proof}
由内部紧支集补零引理，\(\widetilde u\in W^{k,p}(\mathbb R^d)\)，故前一命题给出全空间范数收敛，限制到 \(\Omega\) 后仍收敛。若 \(\operatorname{dist}(x,K)>\varepsilon\)，卷积所读取的点都在 \(K\) 外，因此卷积为零。由其连续性，支集包含于紧集 \(K+\overline B(0,\varepsilon)\)。当 \(K\ne\varnothing\) 时，取
\(\varepsilon<\operatorname{dist}(K,V^c)\)，便使这个紧集包含于 \(V\)；空补集的距离仍约定为无穷。若 \(K=\varnothing\)，原函数及其卷积均为零，结论直接成立。
\end{proof}

这个命题既控制逼近误差，也允许预先指定一个支集必须留在其中的开集。仅有误差趋零还不足以拼接局部逼近：支集若在磨光后聚集到同一点，原来局部有限的函数族就可能失去局部有限性。

\subsection{Meyers–Serrin 定理}
\label{b1:subsec:220103}

\term[meyers-serrin-dingli]{Meyers--Serrin 定理}[Meyers--Serrin theorem]把内部磨光推广到任意开集上的全域范数逼近。它要求逼近函数在 \(\Omega\) 内光滑，不要求这些函数可以光滑延伸到边界外，也不要求它们具有紧支集。

\begin{theorem}[Meyers--Serrin]\label{b1:thm:meyers-serrin}
设 \(\Omega\subseteq\mathbb R^d\) 为非空开集，\(k\in\mathbb N_0\)、\(1\leq p<\infty\)。则 \(C^\infty(\Omega)\cap W^{k,p}(\Omega)\) 在实或复空间 \(W^{k,p}(\Omega)\) 中稠密。

更一般地，给定 \(u\in W^{k,p}_{\mathrm{loc}}(\Omega)\) 及 \(\eta>0\)，存在 \(v\in C^\infty(\Omega)\)，使 \(v-u\in W^{k,p}(\Omega)\) 且
\[
 \|v-u\|_{W^{k,p}(\Omega)}<\eta.
\]
后一陈述不要求 \(u\) 或 \(v\) 各自属于全局 \(W^{k,p}(\Omega)\)。
\end{theorem}
\begin{proof}
证明较一般的陈述。沿用\cref{b1:thm:smooth-partition-unity}证明中的紧集穷竭，记
\[
 K_j=\{\,x\in\Omega\mid |x|\leq j,
                 \operatorname{dist}(x,\Omega^c)\geq1/j\,\},
 \quad K_0=K_{-1}=\varnothing,
\]
并置 \(V_j=\operatorname{int}K_{j+1}\setminus K_{j-2}\)。紧层 \(K_j\setminus\operatorname{int}K_{j-1}\) 覆盖 \(\Omega\)，且各自包含于 \(V_j\)，故这些开集也覆盖 \(\Omega\)。它们局部有限：任意一点有一个包含在某个 \(\operatorname{int}K_N\) 中的邻域，所有 \(j\geq N+2\) 的 \(V_j\) 都避开它。

该单位分解证明在每一层构造有限多个函数，其支集紧含于 \(V_j\)。把这些函数排成 \((\psi_i)_{i\geq1}\)，略去零函数，记第 \(i\) 个函数所属层为 \(j(i)\)。于是
\[
 \psi_i\geq0,
 \quad\sum_i\psi_i=1,
 \quad\operatorname{supp}\psi_i\Subset V_{j(i)},
\]
且每层只出现有限次。令 \(w_i=\psi_i u\)，并在 \(\Omega\) 外补零。前面的乘积讨论及补零引理保证 \(\widetilde w_i\in W^{k,p}(\mathbb R^d)\)。

对每个 \(i\)，使用\cref{b1:prop:compact-sobolev-smooth-approximation}，选取半径 \(\varepsilon_i>0\)，使
\[
 v_i=\rho_{\varepsilon_i}*\widetilde w_i\in C_c^\infty\bigl(V_{j(i)}\bigr),
 \quad
 \|v_i-w_i\|_{W^{k,p}(\Omega)}<\eta2^{-i-1}.
\]
这里先要求 \(\varepsilon_i<\operatorname{dist}\bigl(\operatorname{supp}\psi_i,V_{j(i)}^c\bigr)\)，再进一步缩小半径以达到误差条件。两个要求可以同时实现。

函数族 \((v_i)\) 仍局部有限。事实上，每个紧集只遇到有限层 \(V_j\)，而每层只有有限项。因此
\[
 v=\sum_{i=1}^\infty v_i
\]
在每个点的某个邻域内都是有限个光滑函数的和，故 \(v\in C^\infty(\Omega)\)。

剩下的全域估计应对误差求和。记 \(e_i=v_i-w_i\)。因为
\[
 \sum_i\|e_i\|_{W^{k,p}(\Omega)}\leq\eta/2,
\]
由\cref{b1:thm:sobolev-complete}及\ref{b1:prop:banach-series-criterion}，级数 \(\sum_i e_i\) 在 \(W^{k,p}(\Omega)\) 中收敛到某个 \(e\)，且 \(\|e\|_{W^{k,p}}\leq\eta/2\)。在任意内部相对紧开集上，\(v_i\) 和 \(w_i\) 都只有有限个非零项；足够长的部分和因而几乎处处等于
\[
 \sum_i(v_i-w_i)=v-u\sum_i\psi_i=v-u.
\]
另一方面，部分和在全域 \(L^p\) 中趋于 \(e\)，从而也在这个开集上趋于 \(e\)。故 \(e=v-u\) 几乎处处于该开集。再取可数个这样的开集穷竭 \(\Omega\)，得到全域等式与所需误差界。

若原来的 \(u\in W^{k,p}(\Omega)\)，则 \(v=u+e\) 也属于该空间。对每个正整数 \(n\) 取 \(\eta=1/n\)，便得到定理中的稠密性。
\end{proof}

证明没有宣称 \(\sum_{i=1}^N\psi_i u\) 在全域 Sobolev 范数中趋于 \(u\)。单位分解的导数可能随位置增大；局部有限只保证局部运算合法，并不自动给出全域范数收敛。逐项磨光后的误差另有可和估计，完备性用在这组误差上。

\subsection{\texorpdfstring{\(C_c^\infty\)}{Cc∞} 的稠密性问题}
\label{b1:subsec:220104}

Meyers--Serrin 定理产生的光滑函数可能一直延伸到边界附近或无穷远处。紧支集是额外要求，需要单独判断。

\begin{theorem}\label{b1:thm:whole-space-sobolev-density}
对 \(k\in\mathbb N_0\)、\(1\leq p<\infty\)，空间 \(C_c^\infty(\mathbb R^d)\) 在 \(W^{k,p}(\mathbb R^d)\) 中稠密。
\end{theorem}
\begin{proof}
由光滑截断引理，取 \(0\leq\chi\leq1\)，使 \(\chi=1\) 于 \(\overline B(0,1)\) 的一个邻域，且 \(\chi\in C_c^\infty\bigl(B(0,2)\bigr)\)。对 \(R\geq1\)，置 \(\chi_R(x)=\chi(x/R)\)。任取 \(u\in W^{k,p}(\mathbb R^d)\)。由 Leibniz 公式，对 \(|\alpha|\leq k\)，
\[
 \begin{aligned}
 \partial^\alpha(\chi_Ru-u)
 &=(\chi_R-1)\partial^\alpha u\\
 &\quad+\sum_{\substack{\beta\leq\alpha\\|\beta|>0}}
   \binom\alpha\beta R^{-|\beta|}
   (\partial^\beta\chi)(x/R)\partial^{\alpha-\beta}u.
 \end{aligned}
\]
第一项的 \(L^p\) 范数不超过
\(\|\mathbf1_{\mathbb R^d\setminus B(0,R)}\partial^\alpha u\|_p\)，由有限 \(p\) 的可积尾部趋于零。余下每项的范数不超过
\[
 \binom\alpha\beta R^{-|\beta|}
 \|\partial^\beta\chi\|_\infty\cdot\|\partial^{\alpha-\beta}u\|_p,
\]
也趋于零。有限个多重指标合起来给出 \(\chi_Ru\to u\) 于 \(W^{k,p}(\mathbb R^d)\)。

给定 \(\eta>0\)，先取 \(R\) 使截断误差小于 \(\eta/2\)。函数 \(\chi_Ru\) 具有紧支集，再由\cref{b1:prop:compact-sobolev-smooth-approximation}取 \(v\in C_c^\infty(\mathbb R^d)\)，使 \(\|v-\chi_Ru\|_{W^{k,p}}<\eta/2\)。三角不等式完成逼近。
\end{proof}

这里的截断只切掉无穷远处的尾部，并且 \(\chi_R\) 的非零阶导数带有衰减因子 \(R^{-|\beta|}\)。若试图在一般域的边界附近作越来越陡的截断，其导数反而可能增大，不能原样复制这个估计。

例如在 \(I=(0,1)\) 上取 \(u=1\)。它属于所有有限 \(p\) 的 \(W^{1,p}(I)\)。对任意 \(\varphi\in C_c^\infty(I)\)，函数 \(1-\varphi\) 的连续代表在两个端点都等于 \(1\)，因而由\eqref{b1:eq:interval-sobolev-sup}，
\[
 \begin{aligned}
 1&\leq\|1-\varphi\|_{L^1(I)}+\|\varphi'\|_{L^1(I)}\\
  &\leq\|1-\varphi\|_{L^p(I)}+\|\varphi'\|_{L^p(I)}\\
  &\leq2^{1-1/p}\|1-\varphi\|_{W^{1,p}(I)}.
 \end{aligned}
\]
最后一步是两项有限序列的 Hölder 不等式，也包括 \(p=1\)。因此这些紧支集函数到常数 \(1\) 的距离至少为 \(2^{1/p-1}>0\)。同一估计还排除了任意 \(k\geq1\) 时在 \(W^{k,p}(I)\) 中的逼近，因为高阶范数也控制其中的零阶和一阶项。问题来自边界附近强制为零的要求，不是域内光滑性不足；常数 \(1\) 本身已经光滑。

零阶情形应另作区分。对任意开集 \(\Omega\) 和有限 \(p\)，\(C_c^\infty(\Omega)\) 在 \(L^p(\Omega)\) 中仍稠密。事实上，取内部紧集穷竭 \(K_j\) 及 \(0\leq\chi_j\leq1\)、\(\chi_j\in C_c^\infty(\Omega)\)，使 \(\chi_j=1\) 于 \(K_j\)。则 \(\chi_j u\to u\) 于 \(L^p\)，这是控制收敛定理的直接应用；再对每个 \(\chi_j u\) 使用 \(k=0\) 的内部紧支集磨光命题即可。只控制函数值时，边界层的 \(L^p\) 误差可以消失；同时控制导数后，这个边界层便可能留下不可忽略的代价。下一节将把能够由紧支集光滑函数逼近的 Sobolev 函数单独组成一个闭子空间。

% !TeX root = ../../Book1.tex
\section{\texorpdfstring{\(W_0^{1,p}\)}{W₀(1,p)}}
\label{b1:sec:2202}

域内光滑逼近不要求近似函数在边界附近消失。若近似函数还须有紧支集，上一节的一维反例已经表明，得到的闭包可能是真子空间。本节就以这个闭包定义零边界空间，再说明它与补零、截断以及经典端点值之间的关系。对一般开集，定义本身不预设迹算子存在。

\subsection{定义}
\label{b1:subsec:220201}

\begin{definition}\label{b1:def:sobolev-zero-boundary}
设 \(\Omega\subset\mathbb R^d\) 为非空开集、\(1\le p\le\infty\)。定义
\[
 W_0^{1,p}(\Omega)
 =\overline{C_c^\infty(\Omega)}^{\,W^{1,p}(\Omega)}.
\]
也就是说，\(u\in W_0^{1,p}(\Omega)\) 当且仅当存在 \(\varphi_j\in C_c^\infty(\Omega)\)，使 \(\|u-\varphi_j\|_{W^{1,p}(\Omega)}\to0\)。另记 \(H_0^1(\Omega)=W_0^{1,2}(\Omega)\)。
\end{definition}
\symidx{\(W_0^{1,p}(\Omega)\)：紧支集光滑函数在\(W^{1,p}\)中的闭包}
\symidx{\(H_0^1(\Omega)\)：\(W_0^{1,2}(\Omega)\)}

它是 \(W^{1,p}\) 的闭线性子空间，所以关于原范数完备；\(H_0^1\) 关于原 \(H^1\) 内积为 Hilbert 空间。若 \(1<p<\infty\)，闭子空间的自反性又给出 \(W_0^{1,p}\) 自反。这些结论都来自定义中的闭包，不需要对边界作几何假设。

下标 \(0\) 通常用来表达零边界条件。不过在这个定义中，“零”首先落实为可用内部紧支集函数逼近，而不是给每个边界点指定一个数。一般 Sobolev 元素仅是几乎处处等价类，边界点甚至不属于定义域；若没有额外的代表定理，写下逐点边界值就还没有确定一个运算。

\(p=\infty\) 时，本书也严格使用上述范数闭包。它通常比“有零边界值的 Lipschitz 函数”更小，不能把有限 \(p\) 的稠密性结论直接搬到这个端点。需要时，应明确其他文献所用 \(W_0^{1,\infty}\) 是否采用同一定义。

\subsection{零边界条件}
\label{b1:subsec:220202}

先把不涉及边界点取值、却总是有效的补零结论证明出来。对 \(u\) 在 \(\Omega\) 上的可测代表，令
\[
 E_0u(x)=
 \begin{cases}
 u(x),&x\in\Omega,\\
 0,&x\notin\Omega.
 \end{cases}
\]
这个操作在 \(L^p\) 等价类上始终有定义并保持范数，困难只在于补零后的导数。

\begin{proposition}\label{b1:prop:w0-zero-extension}
若 \(u\in W_0^{1,p}(\Omega)\)，则 \(E_0u\in W^{1,p}(\mathbb R^d)\)，并且
\[
 \partial_i(E_0u)=E_0(\partial_i u),\qquad
 \|E_0u\|_{W^{1,p}(\mathbb R^d)}=\|u\|_{W^{1,p}(\Omega)}.
\]
因此 \(E_0\) 是一个线性等距嵌入，结论包括 \(p=\infty\)。
\end{proposition}
\begin{proof}
取定义中的 \(\varphi_j\)。每个 \(E_0\varphi_j\) 都属于 \(C_c^\infty(\mathbb R^d)\)，而且经典导数在域内外直接给出
\[
 \partial_i(E_0\varphi_j)=E_0(\partial_i\varphi_j).
\]
因此 \(\|E_0\varphi_j-E_0\varphi_k\|_{W^{1,p}(\mathbb R^d)}
=\|\varphi_j-\varphi_k\|_{W^{1,p}(\Omega)}\)。由完备性，它们在全空间的 \(W^{1,p}\) 中收敛。另一方面，\(L^p\) 补零等距性表明其零阶极限为 \(E_0u\)，各一阶分量的极限为 \(E_0(\partial_i u)\)。这就同时证明了导数公式及范数等式。
\end{proof}

反向在任意开集上不成立。令
\[
 \Omega=(-1,0)\cup(0,1),\qquad u(x)=1-|x|\quad(x\in\Omega).
\]
补零后所得等价类由全实轴上的帐篷函数 \((1-|x|)^+\) 表示；仅在原点处填的值不同，不影响等价类，所以 \(E_0u\in W^{1,p}(\mathbb R)\) 对全部 \(p\) 成立。然而若 \(\varphi_j\in C_c^\infty(\Omega)\) 依 \(W^{1,p}(\Omega)\) 逼近 \(u\)，限制到 \((0,1)\)，由\secref{b1:sec:2104}的一维上确界估计，绝对连续代表必须在 \([0,1]\) 上一致收敛。近似函数在 \(0\) 处都为零，极限代表在 \(0\) 处却为一，矛盾。

这里补零只改变一个内部缺失点，积分看不见这个点；但沿两侧区间逼近时，端点行为仍受到导数范数控制。因此“补零没有产生奇异导数”与“能够在域内由紧支集函数逼近”是两条不同的条件。

无穷端点还有另一种障碍。在 \((0,1)\) 上，\(u(x)=x(1-x)\) 光滑且两个端点值都为零，但它不属于 \(W_0^{1,\infty}(0,1)\)。事实上，每个 \(\varphi\in C_c^\infty(0,1)\) 在零点附近恒为零，故 \(\varphi'=0\) 在该邻域；而 \(u'(x)=1-2x\) 趋于一。因此
\[
 \|u'-\varphi'\|_{L^\infty(0,1)}\ge1.
\]
这不是端点函数值的问题，而是强无穷范数还要求导数在靠近端点时也能被一致控制。

\subsection{截断逼近}
\label{b1:subsec:220203}

以后在积分恒等式中使用 Sobolev 函数作测试时，常需先控制它的幅值。有限 \(p\) 下可以在零边界空间内完成这一步。

\begin{proposition}\label{b1:prop:w0-truncation}
设 \(1\le p<\infty\)，实值 \(u\in W_0^{1,p}(\Omega)\)。则对每个 \(M>0\)，截断 \(T_M(u)\) 仍属于 \(W_0^{1,p}(\Omega)\)，且
\[
 T_M(u)\longrightarrow u\quad\text{于 }W^{1,p}(\Omega)
 \quad(M\to\infty).
\]
正部、负部和绝对值也保持 \(W_0^{1,p}\)。若另外 \(|u|\le M\) 几乎处处，则可以选取 \(\psi_j\in C_c^\infty(\Omega)\)，使 \(|\psi_j|\le M\) 且 \(\psi_j\to u\) 于 \(W^{1,p}\)。
\end{proposition}
\begin{proof}
取 \(\varphi_j\in C_c^\infty(\Omega)\) 逼近 \(u\)。固定 \(M\)，先说明
\[
 T_M(\varphi_j)\longrightarrow T_M(u)\quad\text{于 }W^{1,p}(\Omega).
\]
零阶收敛由 \(T_M\) 的 \(1\)-Lipschitz 性得到。一阶导数使用上一章的截断公式，分解为
\[
 \begin{aligned}
 \partial_iT_M(\varphi_j)-\partial_iT_M(u)
 &=\mathbf1_{\{|\varphi_j|<M\}}(\partial_i\varphi_j-\partial_i u)\\
 &\quad+
  \bigl(\mathbf1_{\{|\varphi_j|<M\}}-\mathbf1_{\{|u|<M\}}\bigr)\partial_i u.
 \end{aligned}
\]
第一项在 \(L^p\) 中趋零。沿任意一条 \(\varphi_j\to u\) 几乎处处的子列，第二项在 \(|u|\ne M\) 处几乎处处趋零；在 \(|u|=M\) 上，水平集梯度为零。控制收敛于是给出该项的 \(L^p\) 收敛。对任意子列再抽取这样的子列，可把结论传回全列，正如习题\ref{b1:ex:21-positive-part-continuity}中的论证。

每个 \(T_M(\varphi_j)\) 都是内部紧支集的 \(W^{1,p}\) 函数。由命题\ref{b1:prop:compact-sobolev-smooth-approximation}，它属于 \(W_0^{1,p}\)。再用闭性，得到 \(T_M(u)\in W_0^{1,p}\)。当 \(M\to\infty\) 时的强收敛已由\eqref{b1:eq:sobolev-truncation-convergence}证明。正部、负部和绝对值用相同的紧支集逼近，并使用上一章相应的强连续性，即可得到。

若 \(|u|\le M\)，则 \(T_M(u)=u\)。对每个 \(T_M(\varphi_j)\) 选择足够小的非负单位积分卷积，使支集仍留在 \(\Omega\)，且 \(W^{1,p}\) 误差小于 \(1/j\)。卷积不增加幅值上界，故得到所需的 \(\psi_j\)。
\end{proof}

最后一段同时说明了一个常用事实：内部紧支集的 \(W^{1,p}\) 函数在有限 \(p\) 下属于 \(W_0^{1,p}\)。无穷端点不能以同一理由推出结论，因为磨光并不保证 \(W^{1,\infty}\) 范数收敛；一个内部尖点已经可能构成障碍。这里对截断保持性的断言也只在有限 \(p\) 范围内使用。

\subsection{与迹的关系}
\label{b1:subsec:220204}

在有界区间 \(I=(a,b)\) 上，上一章已经给每个 \(W^{1,p}(I)\) 元素选出了唯一的闭区间绝对连续代表。于是端点取值
\[
 \gamma_Iu=(\widetilde u(a),\widetilde u(b))
\]
是良定义的线性映射；由\eqref{b1:eq:interval-sobolev-sup}及 Hölder 不等式，它对 \(W^{1,p}\) 范数连续。这是一般迹算子的一维原型。

\begin{theorem}\label{b1:thm:interval-zero-boundary}
若 \(I=(a,b)\) 有界、\(1\le p<\infty\)，则
\[
 W_0^{1,p}(I)
 =\{\,u\in W^{1,p}(I):\widetilde u(a)=\widetilde u(b)=0\,\}
 =\ker\gamma_I.
\]
\end{theorem}
\begin{proof}
一个方向由 \(\gamma_I\) 连续且在 \(C_c^\infty(I)\) 上为零立即得到。反过来，设两个端点值均为零。对充分小的 \(\varepsilon>0\)，选 \(0\le\eta_\varepsilon\le1\)，使其在距端点不超过 \(\varepsilon\) 处为零、在 \([a+2\varepsilon,b-2\varepsilon]\) 上为一，且
\[
 \eta_\varepsilon\in C_c^\infty(I),\qquad
 |\eta_\varepsilon'|\le C/\varepsilon.
\]
乘积 \(u_\varepsilon=\eta_\varepsilon u\) 为内部紧支集的 Sobolev 函数，因此属于 \(W_0^{1,p}(I)\)。函数差 \((1-\eta_\varepsilon)u\) 与导数项 \((1-\eta_\varepsilon)u'\) 的 \(L^p\) 范数都因边界层缩小而趋零。需要单独估计的是 \(\eta_\varepsilon'u\)。

由 \(u(a)=0\) 和绝对连续性，对 \(a<x<a+2\varepsilon\)，Hölder 不等式给出
\[
 |u(x)|^p
 \le (x-a)^{p-1}\int_a^x|u'(t)|^p\,\mathrm dt.
\]
\(p=1\) 时这就是积分三角不等式。再次积分并用 Tonelli 定理，得到
\[
 \varepsilon^{-p}\int_a^{a+2\varepsilon}|u(x)|^p\,\mathrm dx
 \le 2^p\int_a^{a+2\varepsilon}|u'(t)|^p\,\mathrm dt
 \longrightarrow0.
\]
在右端点使用 \(u(b)=0\) 得到同样估计。因此 \(\eta_\varepsilon'u\to0\) 于 \(L^p\)，所以 \(u_\varepsilon\to u\) 于 \(W^{1,p}(I)\)。利用 \(W_0^{1,p}\) 的闭性完成证明。
\end{proof}

这个证明显示了零端点值的具体作用：它把 \(u\) 在薄边界层中的大小变成 \(u'\) 在同一层上的积分，从而抵消截断导数中的 \(\varepsilon^{-1}\)。仅仅知道 \(u\) 有界，不足以完成这一步。前面的 \(x(1-x)\) 例子则说明，定理的有限 \(p\) 范围不可直接放宽到无穷端点。

对有界开集，还有一个不要求边界规则的充分条件。若实值 \(u\in W^{1,p}(\Omega)\cap C(\overline\Omega)\)、\(p<\infty\)，并且 \(u=0\) 在 \(\partial\Omega\) 上，则 \(u\in W_0^{1,p}(\Omega)\)。为此取
\[
 v_\varepsilon=u-T_\varepsilon(u).
\]
由连续性与零边界值，\(v_\varepsilon\) 的支集包含于紧集 \(\{x\in\overline\Omega:|u(x)|\ge\varepsilon\}\Subset\Omega\)，所以它属于 \(W_0^{1,p}\)。而 \(T_\varepsilon(u)\to0\) 于 \(L^p\)，其弱梯度为 \(\mathbf1_{\{|u|<\varepsilon\}}\nabla u\)，由水平集梯度为零及控制收敛也趋于零。故 \(v_\varepsilon\to u\) 于 \(W^{1,p}\)。复值情形对实部、虚部分别处理即可。

一般 Sobolev 元素未必具有这样的连续代表；一般边界也未必允许只用逐点限制描述其行为。下一章将对适当的区域构造有界迹算子，并在明确条件下辨认其核。在此之前，\(W_0^{1,p}\) 始终以紧支集光滑闭包为定义，而不是以尚未定义的边界取值替换它。

% !TeX root = ../../Book1.tex
\section{延拓算子}
\label{b1:sec:2203}

在开集内部磨光时，可以把卷积半径选得足够小；接近边界以后，这种选择不再给出同一个全域函数。延拓提供了另一条路线：先在开集外补出函数，同时控制弱导数，再在整个空间中使用卷积。本节只处理一阶空间，且保留两个端点 \(p=1\) 与 \(p=\infty\)。

有限图表、反射和单位分解的几何组织，可对照 Di Nezza、Palatucci 与 Valdinoci 的《Hitchhiker's guide to the fractional Sobolev spaces》\cite[Section 5]{DiNezza2012Hitchhikers}。该文讨论分数阶空间；下面的一阶证明则使用前章的沿线绝对连续代表与双 Lipschitz 换元，不由分数阶结论直接推出。
% 实读来源：DiNezza2012Hitchhikers，作者公开副本 arXiv:1104.4345，
% 第5节 Lemmas 5.1--5.3 与 Theorem 5.4 的完整证明。
% 原文范围为 0<s<1、1<=p<infty；这里只借鉴图表、截断和反射的组织，
% 一阶及 p=infty 的论证由本书21.3--21.4接口独立完成。

\subsection{半空间反射}
\label{b1:subsec:220301}

记上半空间为 \(H_+=\mathbb R^{d-1}\times(0,\infty)\)，点写成 \(x=(x',t)\)。反射时保留函数值，法向导数则改变符号。关键不是这个形式计算，而是证明两个半空间的弱导数能跨过中间的超平面拼在一起。

\begin{theorem}[半空间偶反射]\label{b1:thm:half-space-sobolev-extension}
设 \(1\le p\le\infty\)。对 \(u\in W^{1,p}(H_+)\)，在 \(t\ne0\) 处规定
\[
 (E_+u)(x',t)=u(x',|t|).
\]
该公式定义一个从 \(W^{1,p}(H_+)\) 到 \(W^{1,p}(\mathbb R^d)\) 的线性映射，其弱导数为
\begin{equation}\label{b1:eq:half-space-reflection-derivatives}
 \partial_i(E_+u)(x',t)=
 \begin{cases}
  (\partial_i u)(x',|t|),&1\le i<d,\\
  \operatorname{sgn}(t)(\partial_d u)(x',|t|),&i=d.
 \end{cases}
\end{equation}
此外，按本书的 Sobolev 范数，
\[
 \|E_+u\|_{W^{1,p}(\mathbb R^d)}
 =2^{1/p}\|u\|_{W^{1,p}(H_+)}\quad(1\le p<\infty),
\]
而 \(p=\infty\) 时两侧范数相等。
\end{theorem}
\begin{proof}
先取\cref{b1:thm:sobolev-acl}给出的共同 Borel 代表。由 Fubini 定理及有界盒上的 Hölder 不等式，对几乎每个 \(x'\)，函数 \(u(x',\cdot)\) 与 \(\partial_d u(x',\cdot)\) 在每个 \((0,R)\) 上均可积。只需先取正整数 \(R\)，便可使这里的例外集与 \(R\) 无关。沿线的绝对连续性与\cref{b1:prop:sobolev-interval-representative}于是给出端点值 \(a(x')\)，并有
\[
 u(x',t)=a(x')+\int_0^t\partial_d u(x',r)\,\mathrm dr
 \quad(0\le t\le R).
\]
两侧反射使用同一个 \(a(x')\)，故所得函数在 \([-R,R]\) 上绝对连续，导数为\eqref{b1:eq:half-space-reflection-derivatives}的法向分量。端点值可取为 \(u(x',1/k)\) 的极限；在有限极限不存在处取零，从而仍得到 Borel 代表。

对其他坐标方向，\(t\ne0\) 的直线只是原来水平直线的反射，其绝对连续性及切向导数公式均得保留。超平面 \(t=0\) 所在的整层直线，在这些方向的横向参数中为零测集，可以一并排除。因此新函数有同一个适用于全部方向的沿线绝对连续代表。

反射保持 Lebesgue 测度，故新函数及候选弱导数都在 \(L^p(\mathbb R^d)\) 中。由 ACL 的反向判据，它们确实满足弱导数关系。对有限 \(p\)，函数和每个一阶导数的 \(p\) 次积分都恰好加倍；对 \(p=\infty\)，本性上确界不变，这就给出范数等式。最后，反射把零测集映为零测集，公式不依赖 \(u\) 的代表；超平面上的赋值不改变等价类。因而算子在等价类层面是线性的。
\end{proof}

如果把 \(u(t)=e^{-t}\) 从 \((0,\infty)\) 直接补零到负半轴，原点两侧的值不吻合，分布导数便含有点质量。偶反射得到 \(e^{-|t|}\)，两侧连续接合，其一阶弱导数只有普通函数部分。这里需要的是函数值接合，而不是要求法向导数也在边界两侧相等。

\subsection{Lipschitz 区域}
\label{b1:subsec:220302}

半空间模型可以在边界附近经过坐标变换使用。为此，边界不必具有连续法向，但须能在足够小的坐标柱中写成一张 Lipschitz 图像。

\begin{definition}\label{b1:def:lipschitz-domain}
设 \(\Omega\subset\mathbb R^d\) 为开集。若每个 \(a\in\partial\Omega\) 都有刚性坐标，使 \(a\) 的坐标为零，并有 \(R,H>0\) 及 Lipschitz 函数
\(\varphi:(-R,R)^{d-1}\rightarrow(-H,H)\)，满足 \(\varphi(0)=0\) 和
\[
 \Omega\cap C=\{\,(z,s)\in C\mid s>\varphi(z)\,\},
 \quad C=(-R,R)^{d-1}\times(-H,H),
\]
则称 \(\Omega\) 为\term[lipschitz-quyu]{Lipschitz 区域}[Lipschitz domain]。本章不要求它连通；需要有界性时另行说明。维数为一时，横向坐标省略，条件就是边界附近为区间的一端。
\end{definition}

例如平面第一象限在角点旋转坐标后可写成 \(s>|z|\)，因此角点本身不妨碍这一条件。有限个互不接触的闭球的内部之并，也符合本章不要求连通的约定。这里允许折角，但不允许在同一局部图柱中把开集放在图像的两侧。

下面固定图表尺寸，以便在反射以后仍有空间补零。设图函数的 Lipschitz 常数为 \(L\)，选取 \(\rho>0\)，使
\[
 2\rho<R,\quad 2L\sqrt{d-1}\rho<H/4,
 \quad h=H/4.
\]
若 \(L=0\) 或 \(d=1\)，第二个限制自动满足。取对称盒
\[
 \begin{split}
 D&=(-2\rho,2\rho)^{d-1}\times(-2h,2h),\\
 D_1&=(-\rho,\rho)^{d-1}\times(-h,h),
 \end{split}
\]
并在上述刚性坐标中规定
\[
 F(z,t)=(z,\varphi(z)+t),\quad
 F^{-1}(z,s)=(z,s-\varphi(z)).
\]
由于 \(|\varphi(z)|<H/4\) 且 \(|t|<H/2\)，整个 \(F(D)\) 都留在原图柱内。映射 \(F\) 与 \(F^{-1}\) 的 Lipschitz 常数至多为 \(1+L\)。回到原坐标以后，刚性变换不破坏这一控制。

记 \(U=F(D)\)、\(V=F(D_1)\) 及 \(D^+=D\cap H_+\)。于是 \(\overline V\subset U\)，且
\[
 \Omega\cap U=F(D^+).
\]
两个盒都关于 \(t=0\) 对称；内盒中的集合反射后仍在内盒中。后面将截断函数放在 \(V\) 内，却在较大的 \(U\) 中完成延拓，内外两层之间的余量用来避免侧面和顶面上的跳跃。

\subsection{延拓算子}
\label{b1:subsec:220303}

\begin{definition}\label{b1:def:sobolev-extension-operator}
给定开集 \(\Omega\) 及 \(1\le p\le\infty\)，若线性映射
\[
 E:W^{1,p}(\Omega)\longrightarrow W^{1,p}(\mathbb R^d)
\]
满足 \((Eu)|_\Omega=u\) 几乎处处，则称它为该空间的\term[yantuo-suanzi]{延拓算子}[extension operator]。若还有与 \(u\) 无关的常数 \(C\)，使
\[
 \|Eu\|_{W^{1,p}(\mathbb R^d)}
 \le C\|u\|_{W^{1,p}(\Omega)},
\]
则称这个延拓算子有界。
\end{definition}

先处理落在一个边界图表内的部分。此时截断函数在物理空间中光滑；拉回后只保证 Lipschitz，所以先作光滑乘法、再作坐标变换，可以直接使用已经建立的两条规则。

\begin{lemma}\label{b1:lem:localized-graph-extension}
固定上一小节满足 \(\overline V\Subset U\) 的内外图表及 \(\eta\in C_c^\infty(V)\)。存在与 \(p\) 无关的线性构造 \(E_\eta\)，使每个 \(u\in W^{1,p}(\Omega)\) 都满足
\[
 E_\eta u\in W^{1,p}(\mathbb R^d),\quad
 (E_\eta u)|_\Omega=\eta u,
\]
并有
\[
 \|E_\eta u\|_{W^{1,p}(\mathbb R^d)}
 \le C_{\eta,F,p}\|u\|_{W^{1,p}(\Omega\cap U)}.
\]
延拓的支集包含在一个只依赖于 \(\eta,F\) 的固定紧集内，该紧集包含于 \(U\)。
\end{lemma}
\begin{proof}
光滑乘法的弱 Leibniz 公式给出
\[
 \partial_i(\eta u)=\eta\partial_i u+(\partial_i\eta)u.
\]
因此乘法在这里的 Sobolev 范数下有界，不要求 \(u\) 本性有界。由\cref{b1:thm:sobolev-bilipschitz-change}，函数
\[
 v=(\eta u)\circ F\in W^{1,p}(D^+)
\]
也有相应范数控制。它在 \(D^+\setminus\overline{D_1}\) 上为零，其弱导数在这一开集上同样为零。

将 \(v\) 在 \(H_+\setminus D^+\) 上补零，记为 \(\widehat v\)。这里的支集可能贴着 \(t=0\)，不能把它当作紧含于 \(D^+\) 的内部函数。不过，开集
\[
 D^+,\quad H_+\setminus\overline{D_1}
\]
覆盖整个 \(H_+\)，且在交集上，\(v\) 与零函数及其弱导数都相同。任取 \(H_+\) 上的测试函数，利用\cref{b1:thm:smooth-partition-unity}将它分解为有限个测试函数之和，每个分量支集落在上述某个开集中。在各分量上分别使用 \(v\) 或零函数的弱导数恒等式，再求和，便得到 \(\widehat v\) 的弱导数。候选导数正是原导数的零延拓，因此
\[
 \|\widehat v\|_{W^{1,p}(H_+)}=\|v\|_{W^{1,p}(D^+)}.
\]

现在作半空间偶反射 \(w=E_+\widehat v\)。令
\[
 S=F^{-1}(\operatorname{supp}\eta)\cap\{\,t\ge0\,\},
 \quad \mathcal R(z,t)=(z,-t).
\]
集合 \(S\) 是 \(D_1\) 中的紧集，\(w\) 在 \(S\cup\mathcal R(S)\) 之外几乎处处为零。因此反射后的支集仍紧含于 \(D_1\)。再在 \(U\) 上把 \(w\) 推回物理空间：
\[
 e=w\circ F^{-1}.
\]
双 Lipschitz 换元保证 \(e\in W^{1,p}(U)\)，且它的支集包含于固定紧集
\(F\bigl(S\cup\mathcal R(S)\bigr)\Subset U\)。此时可用\cref{b1:lem:compact-sobolev-zero-extension}把 \(e\) 补零到整个空间，所得函数就是 \(E_\eta u\)。

在 \(\Omega\cap U\) 上，拉回坐标的最后一分量为正，所以反射并未改变原函数，故 \(E_\eta u=\eta u\)；在 \(\Omega\setminus U\) 上两者都为零。各次乘法、复合、反射与补零都是固定的线性操作。依次使用它们的范数界，就得到引理中的估计。
\end{proof}

\subsection{有界延拓}
\label{b1:subsec:220304}

有界性使边界能够由有限个内图表覆盖，从而把局部延拓合成一个全域算子。

\begin{theorem}[Lipschitz 区域的一阶延拓]\label{b1:thm:lipschitz-domain-extension}
设 \(\Omega\subset\mathbb R^d\) 为有界 Lipschitz 区域。可以固定同一套线性延拓公式，使它对每个 \(1\le p\le\infty\) 都定义有界算子
\[
 E:W^{1,p}(\Omega)\longrightarrow W^{1,p}(\mathbb R^d),
 \quad (Eu)|_\Omega=u,
\]
并满足
\begin{equation}\label{b1:eq:lipschitz-domain-extension-bound}
 \|Eu\|_{W^{1,p}(\mathbb R^d)}
 \le C_{\Omega,p}\|u\|_{W^{1,p}(\Omega)}.
\end{equation}
所有 \(Eu\) 均可取支集包含于同一个固定紧集。
\end{theorem}
\begin{proof}
由边界的紧性，选有限个内图表 \(V_1,\ldots,V_N\) 覆盖 \(\partial\Omega\)，对应外图表记为 \(U_j\)。集合
\[
 K_0=\overline\Omega\setminus\bigcup_{j=1}^N V_j
\]
是 \(\Omega\) 内的紧集。选开集 \(V_0\) 包含 \(K_0\)，并使 \(\overline{V_0}\Subset\Omega\)。若 \(K_0\) 为空，可省去这一内部项。

在开集 \(N_0=\bigcup_{j=0}^N V_j\) 上取从属于这有限覆盖的光滑单位分解，再取 \(\chi\in C_c^\infty(N_0)\)，使它在 \(\overline\Omega\) 的一个邻域 \(N_1\) 中等于 \(1\)。单位分解局部有限，所以只有有限个分解函数的支集与 \(\operatorname{supp}\chi\) 相交。将这些函数乘以 \(\chi\)，并按所属的 \(V_j\) 合并，得到
\[
 \eta_j\in C_c^\infty(V_j),\quad
 \sum_{j=0}^N\eta_j=1
 \quad(x\in N_1).
\]
这一步只依赖区域和图表，不依赖待延拓的函数。

内部函数 \(\eta_0u\) 由\cref{b1:lem:compact-sobolev-zero-extension}补零，记为 \(e_0\)。对每个边界项应用\cref{b1:lem:localized-graph-extension}，记其延拓为 \(E_{\eta_j}u\)，并规定
\[
 Eu=e_0+\sum_{j=1}^N E_{\eta_j}u.
\]
限制到 \(\Omega\) 后，各项之和为 \(\sum_j\eta_j u=u\)。所有构造均线性，所以 \(E\) 线性。用三角不等式及各局部估计，得到
\[
 \begin{split}
 \|Eu\|_{W^{1,p}(\mathbb R^d)}
 &\le\|e_0\|_{W^{1,p}(\mathbb R^d)}
       +\sum_{j=1}^N\|E_{\eta_j}u\|_{W^{1,p}(\mathbb R^d)}\\
 &\le\left(C_{0,p}+\sum_{j=1}^N C_{\eta_j,F_j,p}\right)
       \|u\|_{W^{1,p}(\Omega)}.
 \end{split}
\]
常数包含有限个图表的 Lipschitz 界与截断函数的一阶导数界，这就证明\eqref{b1:eq:lipschitz-domain-extension-bound}。各项支集都包含在预先固定的紧集中，有限并仍紧。整个公式不随 \(p\) 改变，因此在不同指数空间的交上也给出同一个延拓。
\end{proof}

半空间的精确因子 \(2^{1/p}\) 来自标准坐标中的测度保持反射；一般区域还经过坐标变换和截断，不能期待相同常数。延拓算子也不是唯一的，改变外部截断或边界图表通常就会改变开集外的函数，但不会改变它在 \(\Omega\) 上的限制。

一阶构造不能直接重复用于高阶空间。取 \(\zeta\in C_c^\infty(\mathbb R)\)，使它在零点附近等于 \(1\)，并令 \(u(t)=t\zeta(t)\) 于 \(t>0\)。这个函数属于半轴上的每个整数阶 Sobolev 空间，但其偶反射在零点附近等于 \(|t|\)，二阶分布导数含有 \(2\delta_0\)，不是局部可积函数。另一方面，Lipschitz 图函数本身也没有足够的高阶导数供坐标链式法则使用。因此这里只证明一阶延拓；Lipschitz 区域上的高阶延拓需要另外的构造，不能把普通反射的失败解释为高阶延拓不可能。

对于有限 \(p\)，将 \(Eu\) 在整个空间磨光再限制回 \(\Omega\)，便可得到跨越边界的光滑近似；\(p=\infty\) 时，延拓的存在并不使卷积自动在一阶 Sobolev 范数下收敛。两个问题分别是能否补出受控函数、能否在指定范数中逼近，不能因前者成立而省略后者的端点条件。

% !TeX root = ../../Book1.tex
\section{局部化}
\label{b1:sec:2204}

前三节反复使用“先乘截断，再磨光或延拓”。现在把其中的估计单独整理出来。局部化不仅用来拼接函数，也能把一个只允许紧支集测试函数的积分恒等式，转成较小区域上的定量控制。需要留意的是截断函数的导数：它们记录了较小区域与外部边界之间的距离，也是局部估计中常数的来源。

\subsection{截断函数}
\label{b1:subsec:220401}

设 \(U,V\subseteq\Omega\) 为开集，且 \(\overline U\Subset V\)。由光滑截断引理，可选 \(0\le\eta\le1\)，使 \(\eta\in C_c^\infty(V)\)，并在 \(\overline U\) 的邻域内等于 \(1\)。对 \(u\in W^{k,p}_{\mathrm{loc}}(\Omega)\)，函数 \(\eta\cdot u\) 在 \(U\) 上与 \(u\) 相同，却具有内部紧支集。这个替换保留了所关注区域中的函数和全部弱导数。

\begin{proposition}\label{b1:prop:sobolev-cutoff-estimate}
设 \(k\in\mathbb N_0\)、\(1\le p\le\infty\)，且 \(\eta\in C_c^\infty(V)\)。对 \(u\in W^{k,p}(V)\)，把 \(\eta\cdot u\) 在 \(V\) 外补零，得到
\[
 \|\widetilde{\eta\cdot u}\|_{W^{k,p}(\mathbb R^d)}
 \le C_{d,k}
 \left(\max_{|\beta|\le k}\|\partial^\beta\eta\|_\infty\right)
 \|u\|_{W^{k,p}(V)}.
\]
更精细地，对每个 \(|\alpha|\le k\)，
\begin{equation}\label{b1:eq:sobolev-cutoff-components}
 \|\partial^\alpha(\eta\cdot u)\|_{L^p(V)}
 \le\sum_{\beta\le\alpha}\binom\alpha\beta
     \|\partial^\beta\eta\|_\infty
     \cdot\|\partial^{\alpha-\beta}u\|_{L^p(V)}.
\end{equation}
\end{proposition}
\begin{proof}
弱 Leibniz 公式给出
\[
 \partial^\alpha(\eta\cdot u)
 =\sum_{\beta\le\alpha}\binom\alpha\beta
       (\partial^\beta\eta)\cdot(\partial^{\alpha-\beta}u).
\]
对每一项使用有界函数乘法的 \(L^p\) 估计，再用三角不等式，即得分量估计。这对 \(p=\infty\) 也成立。多重指标的数目和二项式系数仅依赖于 \(d,k\)，有限维范数比较于是给出总范数估计；所用比较常数可统一取在 \(1\le p\le\infty\) 上有界。\cref{b1:lem:compact-sobolev-zero-extension}保证补零后的导数不增加边界项，也不改变范数。
\end{proof}

在两个同心球之间，截断常数可以写得更具体。若 \(0<r<R\)，存在 \(\eta\in C_c^\infty(B(x_0,R))\)，使 \(0\le\eta\le1\)、\(\eta=1\) 于 \(\overline B(x_0,r)\)，并且
\begin{equation}\label{b1:eq:ball-cutoff-derivatives}
 \|\partial^\beta\eta\|_\infty
 \le C_{d,\beta}(R-r)^{-|\beta|}
 \quad(|\beta|\ge1).
\end{equation}
例如取 \(\delta=(R-r)/4\)，将 \(B(x_0,r+2\delta)\) 的示性函数与 \(\rho_\delta\) 卷积。卷积在 \(\overline B(x_0,r)\) 附近等于 \(1\)，支集包含于 \(\overline B(x_0,r+3\delta)\subset B(x_0,R)\)。将导数放在核上，并用
\[
 \|\partial^\beta\rho_\delta\|_{L^1}
 =\delta^{-|\beta|}\|\partial^\beta\rho\|_{L^1},
\]
就得到上述界。

\ProseParagraphBreak

因此，局部化不是无代价的限制运算。例如一阶时，
\[
 \|\nabla(\eta\cdot u)\|_{L^p}
 \le \|\eta\cdot\nabla u\|_{L^p}
      +\|u\cdot\nabla\eta\|_{L^p}.
\]
第二项含有函数本身，而不只含梯度；当 \(R-r\) 缩小时，它通常会增大。把一个全域估计用于 \(\eta\cdot u\) 后，必须保留这项，不能因 \(\eta=1\) 于内球就将它删去。

\subsection{单位分解}
\label{b1:subsec:220402}

截断处理一个区域；单位分解则使相容的局部函数成为一个整体。这里先区分“局部属于 Sobolev 空间”和“具有有限全域范数”。

\begin{proposition}\label{b1:prop:sobolev-local-gluing}
设 \(k\in\mathbb N_0\)、\(1\le p\le\infty\)，且 \((U_j)_{j\in J}\) 为开集 \(\Omega\) 的开覆盖。若 \(u\in L^1_{\mathrm{loc}}(\Omega)\)，且每个限制 \(u|_{U_j}\) 都属于 \(W^{k,p}_{\mathrm{loc}}(U_j)\)，则 \(u\in W^{k,p}_{\mathrm{loc}}(\Omega)\)。

\ProseParagraphBreak

若进一步 \(u\in L^p(\Omega)\)，且各个局部弱导数拼成的函数 \(g_\alpha\) 满足 \(g_\alpha\in L^p(\Omega)\)（\(1\le|\alpha|\le k\)），则 \(u\in W^{k,p}(\Omega)\)，并有 \(\partial^\alpha u=g_\alpha\)。
\end{proposition}
\begin{proof}
在两个开集的交上，同阶弱导数由唯一性相等。因此各局部导数可以拼接。为避免不可数个零测集的问题，先利用欧氏空间的可数基，从原覆盖中取出可数子覆盖；再在这个可数族内去掉各交集上导数不一致的零测集。由此得到几乎处处定义且可测的 \(g_\alpha\)。

任取 \(K\Subset\Omega\)。可以用有限个相对紧于某些 \(U_j\) 的小开集覆盖 \(K\)，所以 \(g_\alpha\) 在 \(K\) 附近属于 \(L^p\)。现取 \(\varphi\in C_c^\infty(\Omega)\)，并将其按从属于覆盖的光滑单位分解写为有限和
\[
 \varphi=\sum_\ell\psi_\ell\varphi.
\]
每个非零分量都紧支于某个 \(U_j\)。在那里应用弱导数恒等式，求和得到
\[
 \begin{aligned}
 \int_\Omega u\,\partial^\alpha\varphi
 &=\sum_\ell\int_\Omega u\,\partial^\alpha(\psi_\ell\varphi)\\
 &=(-1)^{|\alpha|}\sum_\ell\int_\Omega g_\alpha\psi_\ell\varphi
 =(-1)^{|\alpha|}\int_\Omega g_\alpha\varphi.
 \end{aligned}
\]
第一处等号使用的是测试函数分解后的导数之和；单位分解导数产生的项已经在这个和中相消。于是 \(g_\alpha\) 是全域弱导数。局部或全域的可积性分别给出两种结论。
\end{proof}

同样，若存在 \(\Omega\) 中的局部有限闭集族 \((K_j)\)，使每个 \(v_j\in W^{k,p}_{\mathrm{loc}}(\Omega)\) 都在 \(\Omega\setminus K_j\) 上几乎处处为零，那么局部有限和
\[
 v=\sum_jv_j
\]
属于 \(W^{k,p}_{\mathrm{loc}}(\Omega)\)，其导数可以逐项相加，因为每个点附近实际只有有限项。这里没有保证 \(\|v\|_{W^{k,p}(\Omega)}<\infty\)。例如在全空间取光滑单位分解，所有分解函数都在各自的紧集外为零，但它们的和为 \(1\)，在有限 \(p\) 时不属于 \(L^p(\mathbb R^d)\)。

\ProseParagraphBreak

获得全域界需要额外估计。一种简单而充分的条件是
\[
 \sum_j\|v_j\|_{W^{k,p}(\Omega)}<\infty.
\]
此时完备性使级数在全域范数中收敛，其极限与局部有限和相同，而且总范数不超过右侧之和。Meyers--Serrin 定理正是对磨光误差建立了这样的可和界；不必要求单位分解后的原函数各项也满足它。

\subsection{局部 Sobolev 空间}
\label{b1:subsec:220403}

\secref{b1:sec:2101}的局部空间定义，要求在每个内部相对紧开集上控制函数及导数。用刚才的截断估计，可以把这个条件改写成全空间中的一族条件。

\begin{proposition}\label{b1:prop:local-sobolev-cutoff-characterization}
设 \(k\in\mathbb N_0\)、\(1\le p\le\infty\)，且 \(u\in L^1_{\mathrm{loc}}(\Omega)\)。下列条件等价：
\begin{equivalences}
\item\label{b1:item:local-sobolev-membership}
\(u\in W^{k,p}_{\mathrm{loc}}(\Omega)\)。
\item\label{b1:item:all-sobolev-cutoffs}
对每个 \(\chi\in C_c^\infty(\Omega)\)，补零函数 \(\widetilde{\chi\cdot u}\) 属于 \(W^{k,p}(\mathbb R^d)\)。
\end{equivalences}
若 \(u_n,u\in W^{k,p}_{\mathrm{loc}}(\Omega)\)，则在每个 \(\overline U\Subset\Omega\) 的开集 \(U\) 上有 \(u_n\to u\) 于 \(W^{k,p}(U)\)，当且仅当对每个上述 \(\chi\)，
\[
 \widetilde{\chi\cdot u_n}\longrightarrow\widetilde{\chi\cdot u}
 \quad\text{于 }W^{k,p}(\mathbb R^d).
\]
\end{proposition}
\begin{proof}
由局部可积性，先取一个内部开集包含 \(\operatorname{supp}\chi\)，再应用截断估计和补零引理，得到从\itemref{b1:item:local-sobolev-membership}到\itemref{b1:item:all-sobolev-cutoffs}的蕴含。反过来，对任意 \(\overline U\Subset\Omega\)，选取在 \(\overline U\) 附近为一的 \(\chi\)。由于 \(\widetilde{\chi\cdot u}|_U=u|_U\)，限制全空间弱导数便得到 \(u|_U\in W^{k,p}(U)\)。

收敛部分使用相同的两个操作。由局部收敛及固定 \(\chi\) 的截断估计，得到全空间乘积的收敛；反过来，取 \(\chi=1\) 于 \(\overline U\) 附近，限制不会增大范数，故得到 \(U\) 上的收敛。
\end{proof}

这里的“每个截断函数”不意味着实际工作中必须检验不可数多项。取递增的内部相对紧开集 \(U_j\) 穷竭 \(\Omega\)，再固定 \(\chi_j=1\) 于 \(\overline{U_j}\) 附近，即可只用这一列截断检验局部条件。任意内部紧集都包含在某个 \(U_j\) 中，因而不会漏掉局部信息。

\ProseParagraphBreak

局部有界性还足以在自反指数范围内抽取局部弱极限。具体说，若 \(1<p<\infty\)，且 \((u_n)\) 在每个 \(W^{k,p}(U_j)\) 中有界，则在 \(U_1\) 上用自反性抽取弱收敛子列，再在 \(U_2\) 上从中继续抽取，依次进行。对角子列在每个 \(U_j\) 上弱收敛。限制算子是有界线性的，故相邻区域上的极限限制相同；这些极限拼成 \(u\in W^{k,p}_{\mathrm{loc}}(\Omega)\)。这只给出局部弱收敛，不保证全域范数有界，更不保证强收敛。以互相远离的平移函数为例，函数可以在每个固定内部区域上消失，而全域范数一直保持不变。

\subsection{局部能量估计与 Caccioppoli 不等式}
\label{b1:subsec:220404}

局部化的一个重要用途，是允许测试函数依赖于待研究的 Sobolev 函数。若积分恒等式最初只对 \(C_c^\infty\) 函数成立，不能直接把一个不光滑的乘积代入；内部紧支集的范数逼近提供了所需过渡。

\ProseParagraphBreak

以下在实数范围内讨论。设 \(u\in H^1_{\mathrm{loc}}(\Omega)\)，满足
\begin{equation}\label{b1:eq:weak-harmonic-condition}
 \int_\Omega\nabla u\cdot\nabla\varphi\,\mathrm dx=0
 \quad\bigl(\varphi\in C_c^\infty(\Omega)\bigr).
\end{equation}
这个恒等式是方程 \(-\Delta u=0\) 的弱形式，并不预先要求 \(u\) 有经典二阶导数。下面的\term[caccioppoli-budengshi]{Caccioppoli 不等式}[Caccioppoli inequality]用较大球中的函数大小控制较小球中的梯度能量。

\begin{theorem}[Caccioppoli]\label{b1:thm:caccioppoli}
若 \(u\in H^1_{\mathrm{loc}}(\Omega)\) 满足\eqref{b1:eq:weak-harmonic-condition}，则对实常数 \(c\) 及实值 \(\eta\in C_c^\infty(\Omega)\)，
\begin{equation}\label{b1:eq:caccioppoli-cutoff}
 \int_\Omega\eta^2|\nabla u|^2\,\mathrm dx
 \le4\int_\Omega|u-c|^2\cdot|\nabla\eta|^2\,\mathrm dx.
\end{equation}
特别地，若 \(0<r<R\)、\(\overline B(x_0,R)\Subset\Omega\)，则
\begin{equation}\label{b1:eq:caccioppoli-balls}
 \int_{B(x_0,r)}|\nabla u|^2\,\mathrm dx
 \le\frac{C_d}{(R-r)^2}
       \int_{B(x_0,R)}|u-c|^2\,\mathrm dx.
\end{equation}
\end{theorem}
\begin{proof}
令 \(v=\eta^2(u-c)\)。弱乘积公式保证 \(v\in H^1(\Omega)\)，且它具有内部紧支集。选一个开集 \(V\)，使 \(\operatorname{supp}\eta\Subset V\)、\(\overline V\Subset\Omega\)。由\cref{b1:prop:compact-sobolev-smooth-approximation}，存在 \(v_j\in C_c^\infty(V)\)，使 \(v_j\to v\) 于 \(H^1(\Omega)\)。由于 \(\nabla u\in L^2(V)\)，Cauchy--Schwarz 不等式给出
\[
 \left|\int_V\nabla u\cdot\nabla(v_j-v)\,\mathrm dx\right|
 \le\|\nabla u\|_{L^2(V)}
      \cdot\|\nabla(v_j-v)\|_{L^2(V)}
 \longrightarrow0.
\]
因此弱恒等式也允许使用测试函数 \(v\)。展开
\[
 \nabla v=\eta^2\nabla u+2\eta(u-c)\nabla\eta
\]
并代入，得到
\[
 \int_\Omega\eta^2|\nabla u|^2
 =-2\int_\Omega\eta(u-c)\nabla u\cdot\nabla\eta.
\]
用向量 Cauchy--Schwarz 不等式及 \(2ab\le a^2/2+2b^2\)，得到
\[
 \int_\Omega\eta^2|\nabla u|^2
 \le\frac12\int_\Omega\eta^2|\nabla u|^2
       +2\int_\Omega|u-c|^2\cdot|\nabla\eta|^2.
\]
吸收左侧的一半即得\eqref{b1:eq:caccioppoli-cutoff}。最后选取\eqref{b1:eq:ball-cutoff-derivatives}中的球间截断，使用 \(\eta=1\) 于内球及 \(|\nabla\eta|\le C_d/(R-r)\)，便得\eqref{b1:eq:caccioppoli-balls}。
\end{proof}

常数 \(c\) 可以自由选取，因为方程只涉及梯度。取 \(c=0\) 得到直接的 \(L^2\) 控制；取 \(c\) 为较大球上的平均值，则右侧只测量函数相对于常数的振幅。当 \(u\) 本来就是常数时，左右两侧都为零，这也说明估计保留了方程对加常数的不变性。

\ProseParagraphBreak

本证明没有从 \(\|u\|_{L^2}\) 推出任意函数的梯度估计，而是先用方程将梯度能量变成截断交叉项。没有弱恒等式时，这一步并不存在。另一方面，截断只在内部起作用，所以定理没有使用边界值或迹；靠近边界的估计则需要额外信息。下一章将给出 Sobolev 函数的边界值应如何定义，以及这种边界值与 \(W_0^{1,p}\) 的关系。

单位分解的归一化与支撑安排见\cref{b1:fig:visual-f16}。
% F16：本书原创矢量示意；依据所在节的定义、证明或明确模型。
\begin{figure}[htbp]
\centering
\begin{tikzpicture}[x=1cm,y=1cm,>=stealth,every node/.style={font=\small},line width=.65pt]
\node at(6,3.7){开覆盖与截断函数};
\draw[AnalysisBlue,thick](1.0,2.7)--(5.0,2.7);\node[AnalysisBlue] at(3.0,3.1){\(U_1\)};
\draw[AnalysisBlue,thick] plot coordinates {(1.0000,0.0000) (1.0400,0.0000) (1.0800,0.0000) (1.1200,0.0174) (1.1600,0.0521) (1.2000,0.0868) (1.2400,0.1216) (1.2800,0.1563) (1.3200,0.1911) (1.3600,0.2258) (1.4000,0.2605) (1.4400,0.2953) (1.4800,0.3300) (1.5200,0.3647) (1.5600,0.3995) (1.6000,0.4342) (1.6400,0.4689) (1.6800,0.5037) (1.7200,0.5384) (1.7600,0.5732) (1.8000,0.6079) (1.8400,0.6426) (1.8800,0.6774) (1.9200,0.7121) (1.9600,0.7468) (2.0000,0.7816) (2.0400,0.8163) (2.0800,0.8511) (2.1200,0.8858) (2.1600,0.9205) (2.2000,0.9553) (2.2400,0.9900) (2.2800,1.0247) (2.3200,1.0595) (2.3600,1.0942) (2.4000,1.1289) (2.4400,1.1637) (2.4800,1.1984) (2.5200,1.2332) (2.5600,1.2679) (2.6000,1.3026) (2.6400,1.3374) (2.6800,1.3721) (2.7200,1.4068) (2.7600,1.4416) (2.8000,1.4763) (2.8400,1.5111) (2.8800,1.5458) (2.9200,1.5805) (2.9600,1.6153) (3.0000,1.6500) (3.0400,1.6153) (3.0800,1.5805) (3.1200,1.5458) (3.1600,1.5111) (3.2000,1.4763) (3.2400,1.4416) (3.2800,1.4068) (3.3200,1.3721) (3.3600,1.3374) (3.4000,1.3026) (3.4400,1.2679) (3.4800,1.2332) (3.5200,1.1984) (3.5600,1.1637) (3.6000,1.1289) (3.6400,1.0942) (3.6800,1.0595) (3.7200,1.0247) (3.7600,0.9900) (3.8000,0.9553) (3.8400,0.9205) (3.8800,0.8858) (3.9200,0.8511) (3.9600,0.8163) (4.0000,0.7816) (4.0400,0.7468) (4.0800,0.7121) (4.1200,0.6774) (4.1600,0.6426) (4.2000,0.6079) (4.2400,0.5732) (4.2800,0.5384) (4.3200,0.5037) (4.3600,0.4689) (4.4000,0.4342) (4.4400,0.3995) (4.4800,0.3647) (4.5200,0.3300) (4.5600,0.2953) (4.6000,0.2605) (4.6400,0.2258) (4.6800,0.1911) (4.7200,0.1563) (4.7600,0.1216) (4.8000,0.0868) (4.8400,0.0521) (4.8800,0.0174) (4.9200,0.0000) (4.9600,0.0000) (5.0000,0.0000)};
\draw[orange!80!black,thick](4.5,2.7)--(8.5,2.7);\node[orange!80!black] at(6.5,3.1){\(U_2\)};
\draw[orange!80!black,thick] plot coordinates {(4.5000,0.0000) (4.5400,0.0000) (4.5800,0.0000) (4.6200,0.0174) (4.6600,0.0521) (4.7000,0.0868) (4.7400,0.1216) (4.7800,0.1563) (4.8200,0.1911) (4.8600,0.2258) (4.9000,0.2605) (4.9400,0.2953) (4.9800,0.3300) (5.0200,0.3647) (5.0600,0.3995) (5.1000,0.4342) (5.1400,0.4689) (5.1800,0.5037) (5.2200,0.5384) (5.2600,0.5732) (5.3000,0.6079) (5.3400,0.6426) (5.3800,0.6774) (5.4200,0.7121) (5.4600,0.7468) (5.5000,0.7816) (5.5400,0.8163) (5.5800,0.8511) (5.6200,0.8858) (5.6600,0.9205) (5.7000,0.9553) (5.7400,0.9900) (5.7800,1.0247) (5.8200,1.0595) (5.8600,1.0942) (5.9000,1.1289) (5.9400,1.1637) (5.9800,1.1984) (6.0200,1.2332) (6.0600,1.2679) (6.1000,1.3026) (6.1400,1.3374) (6.1800,1.3721) (6.2200,1.4068) (6.2600,1.4416) (6.3000,1.4763) (6.3400,1.5111) (6.3800,1.5458) (6.4200,1.5805) (6.4600,1.6153) (6.5000,1.6500) (6.5400,1.6153) (6.5800,1.5805) (6.6200,1.5458) (6.6600,1.5111) (6.7000,1.4763) (6.7400,1.4416) (6.7800,1.4068) (6.8200,1.3721) (6.8600,1.3374) (6.9000,1.3026) (6.9400,1.2679) (6.9800,1.2332) (7.0200,1.1984) (7.0600,1.1637) (7.1000,1.1289) (7.1400,1.0942) (7.1800,1.0595) (7.2200,1.0247) (7.2600,0.9900) (7.3000,0.9553) (7.3400,0.9205) (7.3800,0.8858) (7.4200,0.8511) (7.4600,0.8163) (7.5000,0.7816) (7.5400,0.7468) (7.5800,0.7121) (7.6200,0.6774) (7.6600,0.6426) (7.7000,0.6079) (7.7400,0.5732) (7.7800,0.5384) (7.8200,0.5037) (7.8600,0.4689) (7.9000,0.4342) (7.9400,0.3995) (7.9800,0.3647) (8.0200,0.3300) (8.0600,0.2953) (8.1000,0.2605) (8.1400,0.2258) (8.1800,0.1911) (8.2200,0.1563) (8.2600,0.1216) (8.3000,0.0868) (8.3400,0.0521) (8.3800,0.0174) (8.4200,0.0000) (8.4600,0.0000) (8.5000,0.0000)};
\draw[black,thick](8.0,2.7)--(12.0,2.7);\node[black] at(10.0,3.1){\(U_3\)};
\draw[black,thick] plot coordinates {(8.0000,0.0000) (8.0400,0.0000) (8.0800,0.0000) (8.1200,0.0174) (8.1600,0.0521) (8.2000,0.0868) (8.2400,0.1216) (8.2800,0.1563) (8.3200,0.1911) (8.3600,0.2258) (8.4000,0.2605) (8.4400,0.2953) (8.4800,0.3300) (8.5200,0.3647) (8.5600,0.3995) (8.6000,0.4342) (8.6400,0.4689) (8.6800,0.5037) (8.7200,0.5384) (8.7600,0.5732) (8.8000,0.6079) (8.8400,0.6426) (8.8800,0.6774) (8.9200,0.7121) (8.9600,0.7468) (9.0000,0.7816) (9.0400,0.8163) (9.0800,0.8511) (9.1200,0.8858) (9.1600,0.9205) (9.2000,0.9553) (9.2400,0.9900) (9.2800,1.0247) (9.3200,1.0595) (9.3600,1.0942) (9.4000,1.1289) (9.4400,1.1637) (9.4800,1.1984) (9.5200,1.2332) (9.5600,1.2679) (9.6000,1.3026) (9.6400,1.3374) (9.6800,1.3721) (9.7200,1.4068) (9.7600,1.4416) (9.8000,1.4763) (9.8400,1.5111) (9.8800,1.5458) (9.9200,1.5805) (9.9600,1.6153) (10.0000,1.6500) (10.0400,1.6153) (10.0800,1.5805) (10.1200,1.5458) (10.1600,1.5111) (10.2000,1.4763) (10.2400,1.4416) (10.2800,1.4068) (10.3200,1.3721) (10.3600,1.3374) (10.4000,1.3026) (10.4400,1.2679) (10.4800,1.2332) (10.5200,1.1984) (10.5600,1.1637) (10.6000,1.1289) (10.6400,1.0942) (10.6800,1.0595) (10.7200,1.0247) (10.7600,0.9900) (10.8000,0.9553) (10.8400,0.9205) (10.8800,0.8858) (10.9200,0.8511) (10.9600,0.8163) (11.0000,0.7816) (11.0400,0.7468) (11.0800,0.7121) (11.1200,0.6774) (11.1600,0.6426) (11.2000,0.6079) (11.2400,0.5732) (11.2800,0.5384) (11.3200,0.5037) (11.3600,0.4689) (11.4000,0.4342) (11.4400,0.3995) (11.4800,0.3647) (11.5200,0.3300) (11.5600,0.2953) (11.6000,0.2605) (11.6400,0.2258) (11.6800,0.1911) (11.7200,0.1563) (11.7600,0.1216) (11.8000,0.0868) (11.8400,0.0521) (11.8800,0.0174) (11.9200,0.0000) (11.9600,0.0000) (12.0000,0.0000)};
\node at(6,-.65){\(\psi_j=\eta_j/\sum_i\eta_i\)，在工作区域上 \(\sum_j\psi_j=1\)};
\draw[AnalysisBlue,thick] plot coordinates {(1.2000,-1.3000) (1.2463,-1.3000) (1.2926,-1.3000) (1.3389,-1.3000) (1.3852,-1.3000) (1.4314,-1.3000) (1.4777,-1.3000) (1.5240,-1.3000) (1.5703,-1.3000) (1.6166,-1.3000) (1.6629,-1.3000) (1.7092,-1.3000) (1.7555,-1.3000) (1.8017,-1.3000) (1.8480,-1.3000) (1.8943,-1.3000) (1.9406,-1.3000) (1.9869,-1.3000) (2.0332,-1.3000) (2.0795,-1.3000) (2.1258,-1.3000) (2.1721,-1.3000) (2.2183,-1.3000) (2.2646,-1.3000) (2.3109,-1.3000) (2.3572,-1.3000) (2.4035,-1.3000) (2.4498,-1.3000) (2.4961,-1.3000) (2.5424,-1.3000) (2.5886,-1.3000) (2.6349,-1.3000) (2.6812,-1.3000) (2.7275,-1.3000) (2.7738,-1.3000) (2.8201,-1.3000) (2.8664,-1.3000) (2.9127,-1.3000) (2.9590,-1.3000) (3.0052,-1.3000) (3.0515,-1.3000) (3.0978,-1.3000) (3.1441,-1.3000) (3.1904,-1.3000) (3.2367,-1.3000) (3.2830,-1.3000) (3.3293,-1.3000) (3.3755,-1.3000) (3.4218,-1.3000) (3.4681,-1.3000) (3.5144,-1.3000) (3.5607,-1.3000) (3.6070,-1.3000) (3.6533,-1.3000) (3.6996,-1.3000) (3.7459,-1.3000) (3.7921,-1.3000) (3.8384,-1.3000) (3.8847,-1.3000) (3.9310,-1.3000) (3.9773,-1.3000) (4.0236,-1.3000) (4.0699,-1.3000) (4.1162,-1.3000) (4.1624,-1.3000) (4.2087,-1.3000) (4.2550,-1.3000) (4.3013,-1.3000) (4.3476,-1.3000) (4.3939,-1.3000) (4.4402,-1.3000) (4.4865,-1.3000) (4.5328,-1.3000) (4.5790,-1.3000) (4.6253,-1.4266) (4.6716,-1.6581) (4.7179,-1.8895) (4.7642,-2.1210) (4.8105,-2.3524) (4.8568,-2.5838) (4.9031,-2.8000) (4.9493,-2.8000) (4.9956,-2.8000) (5.0419,-2.8000) (5.0882,-2.8000) (5.1345,-2.8000) (5.1808,-2.8000) (5.2271,-2.8000) (5.2734,-2.8000) (5.3197,-2.8000) (5.3659,-2.8000) (5.4122,-2.8000) (5.4585,-2.8000) (5.5048,-2.8000) (5.5511,-2.8000) (5.5974,-2.8000) (5.6437,-2.8000) (5.6900,-2.8000) (5.7362,-2.8000) (5.7825,-2.8000) (5.8288,-2.8000) (5.8751,-2.8000) (5.9214,-2.8000) (5.9677,-2.8000) (6.0140,-2.8000) (6.0603,-2.8000) (6.1066,-2.8000) (6.1528,-2.8000) (6.1991,-2.8000) (6.2454,-2.8000) (6.2917,-2.8000) (6.3380,-2.8000) (6.3843,-2.8000) (6.4306,-2.8000) (6.4769,-2.8000) (6.5231,-2.8000) (6.5694,-2.8000) (6.6157,-2.8000) (6.6620,-2.8000) (6.7083,-2.8000) (6.7546,-2.8000) (6.8009,-2.8000) (6.8472,-2.8000) (6.8934,-2.8000) (6.9397,-2.8000) (6.9860,-2.8000) (7.0323,-2.8000) (7.0786,-2.8000) (7.1249,-2.8000) (7.1712,-2.8000) (7.2175,-2.8000) (7.2638,-2.8000) (7.3100,-2.8000) (7.3563,-2.8000) (7.4026,-2.8000) (7.4489,-2.8000) (7.4952,-2.8000) (7.5415,-2.8000) (7.5878,-2.8000) (7.6341,-2.8000) (7.6803,-2.8000) (7.7266,-2.8000) (7.7729,-2.8000) (7.8192,-2.8000) (7.8655,-2.8000) (7.9118,-2.8000) (7.9581,-2.8000) (8.0044,-2.8000) (8.0507,-2.8000) (8.0969,-2.8000) (8.1432,-2.8000) (8.1895,-2.8000) (8.2358,-2.8000) (8.2821,-2.8000) (8.3284,-2.8000) (8.3747,-2.8000) (8.4210,-2.8000) (8.4672,-2.8000) (8.5135,-2.8000) (8.5598,-2.8000) (8.6061,-2.8000) (8.6524,-2.8000) (8.6987,-2.8000) (8.7450,-2.8000) (8.7913,-2.8000) (8.8376,-2.8000) (8.8838,-2.8000) (8.9301,-2.8000) (8.9764,-2.8000) (9.0227,-2.8000) (9.0690,-2.8000) (9.1153,-2.8000) (9.1616,-2.8000) (9.2079,-2.8000) (9.2541,-2.8000) (9.3004,-2.8000) (9.3467,-2.8000) (9.3930,-2.8000) (9.4393,-2.8000) (9.4856,-2.8000) (9.5319,-2.8000) (9.5782,-2.8000) (9.6245,-2.8000) (9.6707,-2.8000) (9.7170,-2.8000) (9.7633,-2.8000) (9.8096,-2.8000) (9.8559,-2.8000) (9.9022,-2.8000) (9.9485,-2.8000) (9.9948,-2.8000) (10.0410,-2.8000) (10.0873,-2.8000) (10.1336,-2.8000) (10.1799,-2.8000) (10.2262,-2.8000) (10.2725,-2.8000) (10.3188,-2.8000) (10.3651,-2.8000) (10.4114,-2.8000) (10.4576,-2.8000) (10.5039,-2.8000) (10.5502,-2.8000) (10.5965,-2.8000) (10.6428,-2.8000) (10.6891,-2.8000) (10.7354,-2.8000) (10.7817,-2.8000) (10.8279,-2.8000) (10.8742,-2.8000) (10.9205,-2.8000) (10.9668,-2.8000) (11.0131,-2.8000) (11.0594,-2.8000) (11.1057,-2.8000) (11.1520,-2.8000) (11.1983,-2.8000) (11.2445,-2.8000) (11.2908,-2.8000) (11.3371,-2.8000) (11.3834,-2.8000) (11.4297,-2.8000) (11.4760,-2.8000) (11.5223,-2.8000) (11.5686,-2.8000) (11.6148,-2.8000) (11.6611,-2.8000) (11.7074,-2.8000) (11.7537,-2.8000) (11.8000,-2.8000)};
\draw[orange!80!black,thick] plot coordinates {(1.2000,-2.8000) (1.2463,-2.8000) (1.2926,-2.8000) (1.3389,-2.8000) (1.3852,-2.8000) (1.4314,-2.8000) (1.4777,-2.8000) (1.5240,-2.8000) (1.5703,-2.8000) (1.6166,-2.8000) (1.6629,-2.8000) (1.7092,-2.8000) (1.7555,-2.8000) (1.8017,-2.8000) (1.8480,-2.8000) (1.8943,-2.8000) (1.9406,-2.8000) (1.9869,-2.8000) (2.0332,-2.8000) (2.0795,-2.8000) (2.1258,-2.8000) (2.1721,-2.8000) (2.2183,-2.8000) (2.2646,-2.8000) (2.3109,-2.8000) (2.3572,-2.8000) (2.4035,-2.8000) (2.4498,-2.8000) (2.4961,-2.8000) (2.5424,-2.8000) (2.5886,-2.8000) (2.6349,-2.8000) (2.6812,-2.8000) (2.7275,-2.8000) (2.7738,-2.8000) (2.8201,-2.8000) (2.8664,-2.8000) (2.9127,-2.8000) (2.9590,-2.8000) (3.0052,-2.8000) (3.0515,-2.8000) (3.0978,-2.8000) (3.1441,-2.8000) (3.1904,-2.8000) (3.2367,-2.8000) (3.2830,-2.8000) (3.3293,-2.8000) (3.3755,-2.8000) (3.4218,-2.8000) (3.4681,-2.8000) (3.5144,-2.8000) (3.5607,-2.8000) (3.6070,-2.8000) (3.6533,-2.8000) (3.6996,-2.8000) (3.7459,-2.8000) (3.7921,-2.8000) (3.8384,-2.8000) (3.8847,-2.8000) (3.9310,-2.8000) (3.9773,-2.8000) (4.0236,-2.8000) (4.0699,-2.8000) (4.1162,-2.8000) (4.1624,-2.8000) (4.2087,-2.8000) (4.2550,-2.8000) (4.3013,-2.8000) (4.3476,-2.8000) (4.3939,-2.8000) (4.4402,-2.8000) (4.4865,-2.8000) (4.5328,-2.8000) (4.5790,-2.8000) (4.6253,-2.6734) (4.6716,-2.4419) (4.7179,-2.2105) (4.7642,-1.9790) (4.8105,-1.7476) (4.8568,-1.5162) (4.9031,-1.3000) (4.9493,-1.3000) (4.9956,-1.3000) (5.0419,-1.3000) (5.0882,-1.3000) (5.1345,-1.3000) (5.1808,-1.3000) (5.2271,-1.3000) (5.2734,-1.3000) (5.3197,-1.3000) (5.3659,-1.3000) (5.4122,-1.3000) (5.4585,-1.3000) (5.5048,-1.3000) (5.5511,-1.3000) (5.5974,-1.3000) (5.6437,-1.3000) (5.6900,-1.3000) (5.7362,-1.3000) (5.7825,-1.3000) (5.8288,-1.3000) (5.8751,-1.3000) (5.9214,-1.3000) (5.9677,-1.3000) (6.0140,-1.3000) (6.0603,-1.3000) (6.1066,-1.3000) (6.1528,-1.3000) (6.1991,-1.3000) (6.2454,-1.3000) (6.2917,-1.3000) (6.3380,-1.3000) (6.3843,-1.3000) (6.4306,-1.3000) (6.4769,-1.3000) (6.5231,-1.3000) (6.5694,-1.3000) (6.6157,-1.3000) (6.6620,-1.3000) (6.7083,-1.3000) (6.7546,-1.3000) (6.8009,-1.3000) (6.8472,-1.3000) (6.8934,-1.3000) (6.9397,-1.3000) (6.9860,-1.3000) (7.0323,-1.3000) (7.0786,-1.3000) (7.1249,-1.3000) (7.1712,-1.3000) (7.2175,-1.3000) (7.2638,-1.3000) (7.3100,-1.3000) (7.3563,-1.3000) (7.4026,-1.3000) (7.4489,-1.3000) (7.4952,-1.3000) (7.5415,-1.3000) (7.5878,-1.3000) (7.6341,-1.3000) (7.6803,-1.3000) (7.7266,-1.3000) (7.7729,-1.3000) (7.8192,-1.3000) (7.8655,-1.3000) (7.9118,-1.3000) (7.9581,-1.3000) (8.0044,-1.3000) (8.0507,-1.3000) (8.0969,-1.3000) (8.1432,-1.5162) (8.1895,-1.7476) (8.2358,-1.9790) (8.2821,-2.2105) (8.3284,-2.4419) (8.3747,-2.6734) (8.4210,-2.8000) (8.4672,-2.8000) (8.5135,-2.8000) (8.5598,-2.8000) (8.6061,-2.8000) (8.6524,-2.8000) (8.6987,-2.8000) (8.7450,-2.8000) (8.7913,-2.8000) (8.8376,-2.8000) (8.8838,-2.8000) (8.9301,-2.8000) (8.9764,-2.8000) (9.0227,-2.8000) (9.0690,-2.8000) (9.1153,-2.8000) (9.1616,-2.8000) (9.2079,-2.8000) (9.2541,-2.8000) (9.3004,-2.8000) (9.3467,-2.8000) (9.3930,-2.8000) (9.4393,-2.8000) (9.4856,-2.8000) (9.5319,-2.8000) (9.5782,-2.8000) (9.6245,-2.8000) (9.6707,-2.8000) (9.7170,-2.8000) (9.7633,-2.8000) (9.8096,-2.8000) (9.8559,-2.8000) (9.9022,-2.8000) (9.9485,-2.8000) (9.9948,-2.8000) (10.0410,-2.8000) (10.0873,-2.8000) (10.1336,-2.8000) (10.1799,-2.8000) (10.2262,-2.8000) (10.2725,-2.8000) (10.3188,-2.8000) (10.3651,-2.8000) (10.4114,-2.8000) (10.4576,-2.8000) (10.5039,-2.8000) (10.5502,-2.8000) (10.5965,-2.8000) (10.6428,-2.8000) (10.6891,-2.8000) (10.7354,-2.8000) (10.7817,-2.8000) (10.8279,-2.8000) (10.8742,-2.8000) (10.9205,-2.8000) (10.9668,-2.8000) (11.0131,-2.8000) (11.0594,-2.8000) (11.1057,-2.8000) (11.1520,-2.8000) (11.1983,-2.8000) (11.2445,-2.8000) (11.2908,-2.8000) (11.3371,-2.8000) (11.3834,-2.8000) (11.4297,-2.8000) (11.4760,-2.8000) (11.5223,-2.8000) (11.5686,-2.8000) (11.6148,-2.8000) (11.6611,-2.8000) (11.7074,-2.8000) (11.7537,-2.8000) (11.8000,-2.8000)};
\draw[black,thick] plot coordinates {(1.2000,-2.8000) (1.2463,-2.8000) (1.2926,-2.8000) (1.3389,-2.8000) (1.3852,-2.8000) (1.4314,-2.8000) (1.4777,-2.8000) (1.5240,-2.8000) (1.5703,-2.8000) (1.6166,-2.8000) (1.6629,-2.8000) (1.7092,-2.8000) (1.7555,-2.8000) (1.8017,-2.8000) (1.8480,-2.8000) (1.8943,-2.8000) (1.9406,-2.8000) (1.9869,-2.8000) (2.0332,-2.8000) (2.0795,-2.8000) (2.1258,-2.8000) (2.1721,-2.8000) (2.2183,-2.8000) (2.2646,-2.8000) (2.3109,-2.8000) (2.3572,-2.8000) (2.4035,-2.8000) (2.4498,-2.8000) (2.4961,-2.8000) (2.5424,-2.8000) (2.5886,-2.8000) (2.6349,-2.8000) (2.6812,-2.8000) (2.7275,-2.8000) (2.7738,-2.8000) (2.8201,-2.8000) (2.8664,-2.8000) (2.9127,-2.8000) (2.9590,-2.8000) (3.0052,-2.8000) (3.0515,-2.8000) (3.0978,-2.8000) (3.1441,-2.8000) (3.1904,-2.8000) (3.2367,-2.8000) (3.2830,-2.8000) (3.3293,-2.8000) (3.3755,-2.8000) (3.4218,-2.8000) (3.4681,-2.8000) (3.5144,-2.8000) (3.5607,-2.8000) (3.6070,-2.8000) (3.6533,-2.8000) (3.6996,-2.8000) (3.7459,-2.8000) (3.7921,-2.8000) (3.8384,-2.8000) (3.8847,-2.8000) (3.9310,-2.8000) (3.9773,-2.8000) (4.0236,-2.8000) (4.0699,-2.8000) (4.1162,-2.8000) (4.1624,-2.8000) (4.2087,-2.8000) (4.2550,-2.8000) (4.3013,-2.8000) (4.3476,-2.8000) (4.3939,-2.8000) (4.4402,-2.8000) (4.4865,-2.8000) (4.5328,-2.8000) (4.5790,-2.8000) (4.6253,-2.8000) (4.6716,-2.8000) (4.7179,-2.8000) (4.7642,-2.8000) (4.8105,-2.8000) (4.8568,-2.8000) (4.9031,-2.8000) (4.9493,-2.8000) (4.9956,-2.8000) (5.0419,-2.8000) (5.0882,-2.8000) (5.1345,-2.8000) (5.1808,-2.8000) (5.2271,-2.8000) (5.2734,-2.8000) (5.3197,-2.8000) (5.3659,-2.8000) (5.4122,-2.8000) (5.4585,-2.8000) (5.5048,-2.8000) (5.5511,-2.8000) (5.5974,-2.8000) (5.6437,-2.8000) (5.6900,-2.8000) (5.7362,-2.8000) (5.7825,-2.8000) (5.8288,-2.8000) (5.8751,-2.8000) (5.9214,-2.8000) (5.9677,-2.8000) (6.0140,-2.8000) (6.0603,-2.8000) (6.1066,-2.8000) (6.1528,-2.8000) (6.1991,-2.8000) (6.2454,-2.8000) (6.2917,-2.8000) (6.3380,-2.8000) (6.3843,-2.8000) (6.4306,-2.8000) (6.4769,-2.8000) (6.5231,-2.8000) (6.5694,-2.8000) (6.6157,-2.8000) (6.6620,-2.8000) (6.7083,-2.8000) (6.7546,-2.8000) (6.8009,-2.8000) (6.8472,-2.8000) (6.8934,-2.8000) (6.9397,-2.8000) (6.9860,-2.8000) (7.0323,-2.8000) (7.0786,-2.8000) (7.1249,-2.8000) (7.1712,-2.8000) (7.2175,-2.8000) (7.2638,-2.8000) (7.3100,-2.8000) (7.3563,-2.8000) (7.4026,-2.8000) (7.4489,-2.8000) (7.4952,-2.8000) (7.5415,-2.8000) (7.5878,-2.8000) (7.6341,-2.8000) (7.6803,-2.8000) (7.7266,-2.8000) (7.7729,-2.8000) (7.8192,-2.8000) (7.8655,-2.8000) (7.9118,-2.8000) (7.9581,-2.8000) (8.0044,-2.8000) (8.0507,-2.8000) (8.0969,-2.8000) (8.1432,-2.5838) (8.1895,-2.3524) (8.2358,-2.1210) (8.2821,-1.8895) (8.3284,-1.6581) (8.3747,-1.4266) (8.4210,-1.3000) (8.4672,-1.3000) (8.5135,-1.3000) (8.5598,-1.3000) (8.6061,-1.3000) (8.6524,-1.3000) (8.6987,-1.3000) (8.7450,-1.3000) (8.7913,-1.3000) (8.8376,-1.3000) (8.8838,-1.3000) (8.9301,-1.3000) (8.9764,-1.3000) (9.0227,-1.3000) (9.0690,-1.3000) (9.1153,-1.3000) (9.1616,-1.3000) (9.2079,-1.3000) (9.2541,-1.3000) (9.3004,-1.3000) (9.3467,-1.3000) (9.3930,-1.3000) (9.4393,-1.3000) (9.4856,-1.3000) (9.5319,-1.3000) (9.5782,-1.3000) (9.6245,-1.3000) (9.6707,-1.3000) (9.7170,-1.3000) (9.7633,-1.3000) (9.8096,-1.3000) (9.8559,-1.3000) (9.9022,-1.3000) (9.9485,-1.3000) (9.9948,-1.3000) (10.0410,-1.3000) (10.0873,-1.3000) (10.1336,-1.3000) (10.1799,-1.3000) (10.2262,-1.3000) (10.2725,-1.3000) (10.3188,-1.3000) (10.3651,-1.3000) (10.4114,-1.3000) (10.4576,-1.3000) (10.5039,-1.3000) (10.5502,-1.3000) (10.5965,-1.3000) (10.6428,-1.3000) (10.6891,-1.3000) (10.7354,-1.3000) (10.7817,-1.3000) (10.8279,-1.3000) (10.8742,-1.3000) (10.9205,-1.3000) (10.9668,-1.3000) (11.0131,-1.3000) (11.0594,-1.3000) (11.1057,-1.3000) (11.1520,-1.3000) (11.1983,-1.3000) (11.2445,-1.3000) (11.2908,-1.3000) (11.3371,-1.3000) (11.3834,-1.3000) (11.4297,-1.3000) (11.4760,-1.3000) (11.5223,-1.3000) (11.5686,-1.3000) (11.6148,-1.3000) (11.6611,-1.3000) (11.7074,-1.3000) (11.7537,-1.3000) (11.8000,-1.3000)};
\draw[gray,dashed](1.2,-1.3)--(11.8,-1.3);\node[anchor=east] at(1,-1.3){\(1\)};\node[anchor=east] at(1,-2.8){\(0\)};\node at(6,-3.3){归一化后的三个权重};\end{tikzpicture}
\caption[局部覆盖与单位分解拼接]{局部覆盖与单位分解拼接。折线仅示意非负截断的支撑与重叠，实际选取光滑截断。分母在工作区域上为正，归一化后权重和为一；可数情形还需局部有限性。}
\label{b1:fig:visual-f16}
\end{figure}



\begin{chaptersummary}
内部紧支集使补零与弱微分相容；有限指数的卷积逼近随后给出光滑替代。任意开集上的 Meyers--Serrin 定理还需要逐层控制支集及误差，不能把单位分解的局部有限性误当作全域范数收敛。

而 \(W_0^{1,p}\) 用范数闭包定义，不依赖预先存在的边界值。它的元素可以补零，但任意开集上不能由补零反推出闭包条件。有限区间的零端点刻画说明了边界层估计如何起作用，也显示无穷指数与有限指数不能混用。

Lipschitz 区域的一阶延拓来自相容的反射、坐标变换和有限拼接。局部化则把这些构造中反复出现的截断整理为估计工具，使弱恒等式可以接受内部紧支集的 Sobolev 测试函数。下一章将正式定义迹，并说明在适当边界上，范数闭包意义的零边界条件怎样成为零迹条件。
\end{chaptersummary}

\BookExercises

% \chapref{b1:ch:22}习题临时稿；由章聚合文件已有的 BookExercises 入口统一接入。
% 十题均为自拟的基础训练或工具变式，来源说明记录实际前置，不冒认教材原题编号。
% 不使用第23章的迹或容量，也不使用后续嵌入及一般椭圆方程理论。

\begin{exercise}\label{b1:ex:22-basic-sobolev-mollification}
设 \(1\le p<\infty\)、\(u\in W^{1,p}(\mathbb R^d)\)，并令
\(u_\varepsilon=\rho_\varepsilon*u\)。逐个导数分量证明
\[
 \partial_i u_\varepsilon=\rho_\varepsilon*\partial_i u,qquad
 u_\varepsilon\longrightarrow u\quad\text{于 }W^{1,p}(\mathbb R^d).
\]
说明为什么同一论证不能把收敛结论直接扩展到 \(p=\infty\)。
\end{exercise}
% 来源：自拟；把正文全阶磨光命题降到一阶逐分量复算，训练最基本的Sobolev范数逼近。
\begin{solution}
由分布卷积的求导公式，
\[
 \partial_i u_\varepsilon
 =\partial_i(\rho_\varepsilon*u)
 =\rho_\varepsilon*\partial_i u.
\]
有限指数的逼近恒等分别给出
\[
 \|\rho_\varepsilon*u-u\|_p\longrightarrow0,qquad
 \|\rho_\varepsilon*\partial_i u-\partial_i u\|_p\longrightarrow0
 \quad(1\le i\le d).
\]
按本书的 \(W^{1,p}\) 范数把这有限多个分量的 \(p\) 次方相加，便得
\(\|u_\varepsilon-u\|_{W^{1,p}}\to0\)。无穷指数下仍有卷积范数不增，但一般 \(L^\infty\) 函数不具平移连续性；即使函数本身可以一致逼近，一阶导数的跳跃也可能使 \(W^{1,\infty}\) 误差不趋零。因此上述最后一步所用的有限指数逼近恒等不能照搬。
\end{solution}

\begin{exercise}\label{b1:ex:22-basic-even-reflection}
在半直线 \((0,\infty)\) 上取 \(u(t)=e^{-t}\)，并在全实轴定义偶反射
\(Eu(t)=e^{-|t|}\)。直接按测试函数定义证明
\[
 (Eu)'(t)=-\operatorname{sgn}(t)e^{-|t|}
 \quad\text{几乎处处},
\]
特别说明原点处没有 Dirac 项。再验证有限 \(p\) 时
\(\|Eu\|_{W^{1,p}(\mathbb R)}=2^{1/p}\|u\|_{W^{1,p}(0,\infty)}\)，而 \(p=\infty\) 时两侧相等。
\end{exercise}
% 来源：自拟；用显式偶反射训练跨边界弱求导及范数计算，不重复一般维半空间证明。
\begin{solution}
任取 \(\varphi\in C_c^\infty(\mathbb R)\)。分别在两侧积分并换元，得到
\[
 \begin{aligned}
 \int_{\mathbb R}e^{-|t|}\varphi'(t)\,\mathrm dt
 &=\int_0^\infty e^{-t}\varphi'(t)\,\mathrm dt
   +\int_0^\infty e^{-t}\varphi'(-t)\,\mathrm dt\\
 &=-\int_0^\infty(-e^{-t})\varphi(t)\,\mathrm dt
   -\int_{-\infty}^0 e^t\varphi(t)\,\mathrm dt.
 \end{aligned}
\]
两次分部积分在原点产生的边界项互相抵消；这正是偶反射两侧函数值相接所排除的 Dirac 项。右端等于
\(-\int_{\mathbb R}[-\operatorname{sgn}(t)e^{-|t|}]\varphi(t)\,\mathrm dt\)，故得到所列弱导数。

函数及弱导数的绝对值都等于 \(e^{-|t|}\)。有限 \(p\) 时，它们各自在全轴上的 \(p\) 次积分都是半轴上的两倍，所以 Sobolev 范数乘以 \(2^{1/p}\)；无穷端点的本性上确界在反射后不变。
\end{solution}

\begin{exercise}\label{b1:ex:22-strip-poincare}
设非空开集 \(\Omega\subset(a,b)\times\mathbb R^{d-1}\)，其中 \(L=b-a<\infty\)。证明对每个 \(1\le p\le\infty\) 及 \(u\in W_0^{1,p}(\Omega)\)，
\[
 \|u\|_{L^p(\Omega)}\le L\|\partial_1u\|_{L^p(\Omega)}.
\]
据此说明，仅由全部一阶弱导数组成的半范数在 \(W_0^{1,p}(\Omega)\) 上是范数，并与原 Sobolev 范数等价。不要求 \(\Omega\) 在其余坐标方向有界。
\end{exercise}
% 来源：自拟；紧支集光滑函数的沿线积分、Hölder及闭包定义给出粗Poincare估计，不预借后续同名定理。
\begin{solution}
先取 \(\varphi\in C_c^\infty(\Omega)\)，并把它补零为全空间上的紧支集光滑函数。写 \(x=(x_1,x')\)。沿第一坐标应用微积分基本定理，
\[
 \varphi(x_1,x')=\int_a^{x_1}\partial_1\varphi(t,x')\,\mathrm dt
 \quad(a<x_1<b).
\]
若 \(1\le p<\infty\)，Hölder 不等式给出
\[
 |\varphi(x_1,x')|^p
 \le L^{p-1}\int_a^b|\partial_1\varphi(t,x')|^p\,\mathrm dt.
\]
这里 \(p=1\) 时直接使用积分的三角不等式。先对 \(x_1\in(a,b)\) 积分，再对 \(x'\) 积分；补零函数及其导数在 \(\Omega\) 外都为零，所以
\[
 \|\varphi\|_{L^p(\Omega)}^p
 \le L^p\|\partial_1\varphi\|_{L^p(\Omega)}^p.
\]
若 \(p=\infty\)，直接由沿线积分取得
\(\|\varphi\|_\infty\le L\|\partial_1\varphi\|_\infty\)。

对一般 \(u\in W_0^{1,p}(\Omega)\)，按\cref{b1:def:sobolev-zero-boundary}取 \(\varphi_n\in C_c^\infty(\Omega)\)，使 \(\varphi_n\to u\) 于 \(W^{1,p}\)。函数及第一偏导数的范数均随之收敛，在刚证的不等式中取极限即可。这个步骤也适用于 \(p=\infty\)，因为使用的是闭包定义给定的逼近，而非一般函数的无穷范数磨光。

记 \(M_p(u)=\bigl(\sum_{j=1}^d\|\partial_ju\|_p^p\bigr)^{1/p}\)，其中 \(p<\infty\)；无穷端点改记 \(M_\infty(u)=\max_j\|\partial_ju\|_\infty\)。如果 \(M_p(u)=0\)，所求不等式使 \(\|u\|_p=0\)，故它确为范数。有限 \(p\) 时，
\[
 M_p(u)\le\|u\|_{W^{1,p}}
 \le(1+L^p)^{1/p}M_p(u);
\]
无穷端点则有
\[
 M_\infty(u)\le\|u\|_{W^{1,\infty}}
 \le\max(1,L)M_\infty(u).
\]
这些常数只取决于条带宽度和指数，并不要求条带中每条截线与 \(\Omega\) 的交集连通。
\end{solution}

\begin{exercise}[\(\star\) 选读]\label{b1:ex:22-w0-infty-interval}
设 \(I=(0,1)\)，\(u\in W^{1,\infty}(I)\)。证明下列两条等价：
\begin{equivalences}
\item\label{b1:item:22-w0-infty-closure}
\(u\in W_0^{1,\infty}(I)\)。
\item\label{b1:item:22-w0-infty-c1-endpoints}
等价类 \(u\) 有一个代表 \(U\in C^1([0,1])\)，且
\[
 U(0)=U(1)=U'(0)=U'(1)=0.
\]
\end{equivalences}
这里端点导数按单侧导数解释。
\end{exercise}
% 来源：自拟；用C1一致极限刻画无穷端点的闭包，进一步说明正文单个端点反例背后的完整条件。
\begin{solution}
先设 \itemref{b1:item:22-w0-infty-closure}成立，选取紧支集光滑函数 \(\varphi_n\to u\) 于 \(W^{1,\infty}(I)\)。每个 \(\varphi_n\) 在端点附近为零，可以看作 \([0,1]\) 上的光滑函数。连续函数在这个闭区间上的上确界等于其在 \(I\) 上的本性上确界。因此 \((\varphi_n)\) 与 \((\varphi_n')\) 都在一致范数下为 Cauchy 列，分别趋于连续函数 \(U,G\)，并满足
\[
 U(0)=U(1)=G(0)=G(1)=0.
\]
在等式 \(\varphi_n(x)=\int_0^x\varphi_n'(t)\,\mathrm dt\) 中取一致极限，得到 \(U(x)=\int_0^xG(t)\,\mathrm dt\)。所以 \(U\in C^1([0,1])\)，且 \(U'=G\)。由于 \(\varphi_n\) 同时在 \(L^\infty(I)\) 中趋于 \(u\)，这个 \(U\) 正是 \(u\) 的一个代表，第二条成立。

反过来，设第二条成立。取 \(\theta\in C^\infty(\mathbb R)\)，使 \(0\le\theta\le1\)，在 \((-\infty,1]\) 上为零，在 \([2,\infty)\) 上为一，并记 \(C=\|\theta'\|_\infty\)。对 \(0<\varepsilon<1/4\)，令
\[
 \chi_\varepsilon(x)=\theta(x/\varepsilon)
                 \theta\bigl((1-x)/\varepsilon\bigr).
\]
它在两端的 \(\varepsilon\) 邻域为零，在 \([2\varepsilon,1-2\varepsilon]\) 上为一，且
\(\|\chi_\varepsilon'\|_\infty\le2C/\varepsilon\)。
记
\[
 \omega_\varepsilon
 =\sup_{x\in[0,2\varepsilon]\cup[1-2\varepsilon,1]}|U'(x)|.
\]
由 \(U'\) 连续且端点值为零，\(\omega_\varepsilon\to0\)。利用 \(U\) 的两个端点值也为零，在这些边界区间上有
\[
 |U(x)|\le2\varepsilon\omega_\varepsilon.
\]
于是
\[
 \|(1-\chi_\varepsilon)U\|_\infty
 \le2\varepsilon\omega_\varepsilon,
 \quad
 \|\bigl((1-\chi_\varepsilon)U\bigr)'\|_\infty
 \le(1+4C)\omega_\varepsilon.
\]
因此 \(\chi_\varepsilon U\to U\) 于 \(W^{1,\infty}(I)\)。

每个 \(\chi_\varepsilon U\) 是内部紧支集的 \(C^1\) 函数，补零后仍属于 \(C_c^1(\mathbb R)\)。固定 \(\varepsilon\)，再选择充分小的磨光尺度，使卷积支集仍紧含于 \(I\)。命题\ref{b1:prop:local-smoothing-approximation}的 \(C^1\) 部分保证函数及其一阶导数一致逼近。依次选择这两个参数，便得到 \(C_c^\infty(I)\) 中依 \(W^{1,\infty}\) 逼近 \(u\) 的序列。

这一步只对已有 \(C^1\) 代表的函数使用逐阶一致磨光，并未声称任意 \(W^{1,\infty}\) 函数都有这样的逼近。
\end{solution}

\begin{exercise}\label{b1:ex:22-w0-order-ideal}
设 \(1\le p<\infty\)，实值 \(u\in W_0^{1,p}(\Omega)\)、\(v\in W^{1,p}(\Omega)\)，且 \(u\ge0\)、\(|v|\le u\) 几乎处处。
证明 \(v\in W_0^{1,p}(\Omega)\)。再在 \(I=(0,1)\) 上构造满足相同大小关系的 \(u\in C_c^\infty(I)\) 和内部紧支集 \(v\in W^{1,\infty}(I)\)，使 \(v\notin W_0^{1,\infty}(I)\)。
\end{exercise}
% 来源：自拟；用第21章正部强连续性制造保序逼近，无穷端点反例是内部尖点，不重复正文边界障碍。
\begin{solution}
取实值 \(\varphi_n\in C_c^\infty(\Omega)\)，使 \(\varphi_n\to u\) 于 \(W^{1,p}\)，并令 \(s_n=\varphi_n^+\)。由\cref{b1:prop:sobolev-lattice-strong-continuity}，\(s_n\to u^+=u\) 于同一空间。再令
\[
 v_n=\max\bigl(-s_n,\min(v,s_n)\bigr).
\]
正部运算、加法和减法都在有限指数的 \(W^{1,p}\) 中连续，而
\[
 \min(a,b)=a-(a-b)^+,\quad
 \max(a,b)=b+(a-b)^+.
\]
因此 \(v_n\in W^{1,p}(\Omega)\)，并在这个空间中有
\[
 v_n\longrightarrow\max\bigl(-u,\min(v,u)\bigr)=v.
\]
最后一个等式使用了 \(|v|\le u\)。

由于 \(-s_n\le v_n\le s_n\)，函数 \(v_n\) 在紧集 \(\operatorname{supp}\varphi_n\Subset\Omega\) 外几乎处处为零。它虽不一定光滑，却具有正文所定义的内部紧支集。由\cref{b1:prop:compact-sobolev-smooth-approximation}，每个 \(v_n\) 都能由内部紧支集光滑函数在 \(W^{1,p}\) 中逼近，故 \(v_n\in W_0^{1,p}(\Omega)\)。闭包空间本身是闭的，再由 \(v_n\to v\) 得到所求结论。

对无穷端点，取非负 \(\chi\in C_c^\infty(I)\)，使它在 \(1/2\) 的一个邻域等于一，令
\[
 u=\chi,\quad v(x)=\chi(x)|x-1/2|.
\]
则 \(0\le v\le u/2\le u\)。函数 \(v\) 为内部紧支集 Lipschitz 函数，因而属于 \(W^{1,\infty}(I)\)；但它在 \(1/2\) 处有尖点，不属于 \(C^1(I)\)。若它另有一个 \(C^1\) 代表，该代表与连续的 \(v\) 几乎处处相等，便由连续性处处相等，仍无法消除尖点。由\cref{b1:ex:22-w0-infty-interval}，\(v\notin W_0^{1,\infty}(I)\)。

所以内部紧支集本身不能保证无穷指数的光滑范数逼近；有限指数证明所用的非线性强连续性与内部磨光都确有指数范围。
\end{solution}

\begin{exercise}\label{b1:ex:22-halfspace-zero-extension-converse}
令 \(H=\{\,x\in\mathbb R^d\mid x_d>0\,\}\)，\(1\le p<\infty\)，并设 \(u\in W^{1,p}(H)\)。记 \(\widetilde u\) 为它在 \(H\) 外的零延拓。证明下列两条等价：
\begin{equivalences}
\item\label{b1:item:22-halfspace-w0}
\(u\in W_0^{1,p}(H)\)。
\item\label{b1:item:22-halfspace-global-zero-extension}
\(\widetilde u\in W^{1,p}(\mathbb R^d)\)。
\end{equivalences}
反向证明须实际构造支集留在半空间内部的光滑逼近，而不使用迹定理。
\end{exercise}
% 来源：自拟；由向内平移、远处截断和内部磨光证明半空间零延拓逆命题，正文只给正向与任意域的失败边界。
\begin{solution}
第一条蕴涵第二条是命题\ref{b1:prop:w0-zero-extension}。现设第二条成立，记 \(F=\widetilde u\)。对 \(h>0\) 定义
\[
 F_h(x)=F(x-he_d).
\]
这个函数在 \(\{x_d<h\}\) 上几乎处处为零。通过弱导数测试中的平移换元可得
\[
 \partial_jF_h(x)=(\partial_jF)(x-he_d).
\]
对函数和全部一阶导数分别应用\cref{b1:lem:lp-translation-continuity}，得到
\[
 \|F_h-F\|_{W^{1,p}(\mathbb R^d)}\longrightarrow0
 \quad(h\downarrow0).
\]
因而可以先把函数的支集与边界隔开，再处理它在无穷远处的部分。

固定 \(h\)。取 \(0\le\chi\le1\)、\(\chi\in C_c^\infty(\mathbb R^d)\)，使 \(\chi=1\) 于单位球，置 \(\chi_R(x)=\chi(x/R)\)、\(G_{h,R}=\chi_RF_h\)。当 \(R\to\infty\) 时，函数及导数中的
\[
 (\chi_R-1)F_h,\quad(\chi_R-1)\partial_jF_h
\]
都由控制收敛在 \(L^p\) 中趋零，而额外项满足
\[
 \|(\partial_j\chi_R)F_h\|_p
 \le R^{-1}\|\partial_j\chi\|_\infty\cdot\|F_h\|_p
 \longrightarrow0.
\]
所以 \(G_{h,R}\to F_h\) 于 \(W^{1,p}(\mathbb R^d)\)。同时 \(G_{h,R}\) 有紧支集，并在 \(\{x_d<h\}\) 上几乎处处为零。

最后固定 \(h,R\)，令
\[
 \varphi_{h,R,\varepsilon}=G_{h,R}*\rho_\varepsilon,
 \quad 0<\varepsilon<h/2.
\]
其支集有界且包含在 \(\{x_d\ge h-\varepsilon\}\) 中，因此 \(\varphi_{h,R,\varepsilon}\in C_c^\infty(H)\)。卷积微分公式以及有限指数的逼近恒等分别作用于 \(G_{h,R}\) 和各一阶导数，给出它在全空间 \(W^{1,p}\) 中的收敛：
\[
 \varphi_{h,R,\varepsilon}\longrightarrow G_{h,R}.
\]
给定任意误差要求，依次选择充分小的 \(h\)、充分大的 \(R\)、再选小于 \(h/2\) 的充分小 \(\varepsilon\)，使上述三段误差之和满足要求。限制到 \(H\) 不会增大 Sobolev 范数，所得函数便在 \(W^{1,p}(H)\) 中逼近 \(u\)。这证明 \(u\in W_0^{1,p}(H)\)。

半空间的作用在第一步：统一向内平移能为整个支集留下正距离。这个几何步骤在一般开集上未必可行，因此不能据本题撤去正文任意域逆命题的限制。
\end{solution}

\begin{exercise}[\(\star\) 选读]\label{b1:ex:22-second-order-reflection}
设 \(1\le p\le\infty\)，记 \(I=(0,\infty)\)。对 \(u\in W^{2,p}(I)\)，定义
\[
 Eu(x)=
 \begin{cases}
 u(x),&x>0,\\
 3u(-x)-2u(-2x),&x<0.
 \end{cases}
\]
原点处的值不影响等价类。证明 \(E\) 是从 \(W^{2,p}(I)\) 到 \(W^{2,p}(\mathbb R)\) 的有界线性延拓算子，并说明系数 \(3,-2\) 如何由原点处的两项匹配条件确定。
\end{exercise}
% 来源：自拟；修正正文偶反射的一阶构造，在半直线上显式匹配函数及一阶导数，不预借高阶迹理论。
\begin{solution}
先说明匹配条件有意义。固定 \(R>0\)，由有限区间上的 Hölder 不等式，\(u,u',u''\) 都在 \((0,R)\) 上可积。对 \(u'\) 应用\cref{b1:thm:one-dimensional-weak-ac}，其连续代表可写为某个常数加上 \(u''\) 的积分。由于 \(u''\in L^1(0,R)\)，这个代表绝对连续地延伸到原点，记为 \(G\)。同样，\(u\) 有一个代表 \(U\)，满足
\[
 U(x)=U(x_0)+\int_{x_0}^xG(t)\,\mathrm dt,
\]
并延伸为 \([0,R]\) 上的 \(C^1\) 函数，\(U'=G\)。在不同 \(R\) 上，连续代表由唯一性一致，故 \(U(0)\)、\(G(0)\) 都有明确含义。

假设先在负半轴使用 \(AU(-x)+BU(-2x)\)。其原点左极限为 \((A+B)U(0)\)，一阶导数的左极限为 \((-A-2B)G(0)\)。要对任意 \(u\) 分别与右侧的函数和导数匹配，应有
\[
 A+B=1,\quad -A-2B=1.
\]
解得 \(A=3\)、\(B=-2\)。因此题中延拓函数在原点连续，其分段一阶导数
\[
 Q(x)=
 \begin{cases}
 G(x),&x>0,\\
 -3G(-x)+4G(-2x),&x<0
 \end{cases}
\]
也在原点连续。

两侧绝对连续函数只要端点值相同，就能绝对连续地拼合：把两侧的积分表示以共同端点值为常数连接即可。由此，\(EU\) 在每个紧区间上绝对连续，弱导数是 \(Q\)；而 \(Q\) 同样绝对连续，其弱导数在负半轴上为 \(3u''(-x)-8u''(-2x)\)，在正半轴上为 \(u''(x)\)。两次拼合均无跳跃，因而没有遗漏集中于原点的分布项。

统一地，对 \(j=0,1,2\)，负半轴上的第 \(j\) 阶导数为
\[
 \partial^j(Eu)(x)
 =3(-1)^ju^{(j)}(-x)-2(-2)^ju^{(j)}(-2x).
\]
记 \(I_-=(-\infty,0)\)。换元给出
\[
 \|u^{(j)}(-2\,\cdot)\|_{L^p(I_-)}
 =2^{-1/p}\|u^{(j)}\|_{L^p(I)},
\]
其中 \(p=\infty\) 时约定 \(1/p=0\)。由三角不等式，负半轴上的导数范数至多是
\(\bigl(3+2^{j+1-1/p}\bigr)\|u^{(j)}\|_{L^p(I)}\)。
再加上正半轴的贡献，可使用统一粗界
\[
 \|\partial^j(Eu)\|_{L^p(\mathbb R)}
 \le12\|u^{(j)}\|_{L^p(I)}.
\]
对三个导数分量取相应的有限 \(p\) 和范数或无穷端点最大值，得到
\[
 \|Eu\|_{W^{2,p}(\mathbb R)}
 \le12\|u\|_{W^{2,p}(I)}.
\]
定义显然对等价类良定义并线性，且在正半轴保留 \(u\)，因此确为所需延拓算子。这里建立的是半直线上的二阶构造，不是任意 Lipschitz 开集的高阶延拓结论。
\end{solution}

\begin{exercise}\label{b1:ex:22-quadratic-localization}
设 \(\Omega\subset\mathbb R^d\) 为开集，实光滑函数 \(\eta_1,\ldots,\eta_N\) 满足
\[
 \sum_{j=1}^N\eta_j^2=1,\quad
 A=\left\|\sum_{j=1}^N|\nabla\eta_j|^2\right\|_{L^\infty(\Omega)}<\infty.
\]
证明对实值或复值 \(u\in H^1(\Omega)\)，每个 \(\eta_ju\) 都属于 \(H^1(\Omega)\)，并有
\[
 \sum_{j=1}^N\|\eta_ju\|_{H^1(\Omega)}^2
 =\|u\|_{H^1(\Omega)}^2
   +\int_\Omega\left(\sum_{j=1}^N|\nabla\eta_j|^2\right)|u|^2\,\mathrm dx.
\]
据此给出局部化范数和原范数的双向估计，并用一个具体例子说明右侧的额外项不能一般删除。
\end{exercise}
% 来源：自拟；平方和分解使交叉项精确抵消，新增可计算的局部化能量误差，不另命名未经核定的专门公式。
\begin{solution}
由平方和为一，\(|\eta_j|\le1\)；由 \(A<\infty\)，每个光滑 \(\nabla\eta_j\) 也有界。光滑乘法的弱 Leibniz 公式给出
\[
 \nabla(\eta_ju)=\eta_j\nabla u+u\nabla\eta_j.
\]
右侧两项均属于 \(L^2(\Omega)\)，而 \(\eta_ju\in L^2(\Omega)\)，所以乘积属于 \(H^1\)。这里不要求 \(u\) 本性有界：起作用的是光滑乘子的函数和一阶导数有界。

几乎处处展开欧氏长度的平方，
\[
 |\nabla(\eta_ju)|^2
 =\eta_j^2|\nabla u|^2+|u|^2\cdot|\nabla\eta_j|^2
   +2\eta_j\operatorname{Re}(\overline u\nabla u)\cdot\nabla\eta_j.
\]
对于实值函数，实部与共轭不改变表达式。对恒等式 \(\sum_j\eta_j^2=1\) 求导，得到
\[
 \sum_{j=1}^N\eta_j\nabla\eta_j=0.
\]
因此将前一等式求和时，全部交叉项相消，得到
\[
 \sum_{j=1}^N|\nabla(\eta_ju)|^2
 =|\nabla u|^2+\left(\sum_{j=1}^N|\nabla\eta_j|^2\right)|u|^2.
\]
零阶部分另有 \(\sum_j|\eta_ju|^2=|u|^2\)。将两条等式积分，便得所求能量恒等式。

额外项非负且至多为 \(A\|u\|_2^2\)，所以
\[
 \|u\|_{H^1}
 \le\left(\sum_{j=1}^N\|\eta_ju\|_{H^1}^2\right)^{1/2}
 \le(1+A)^{1/2}\|u\|_{H^1}.
\]
取 \(\Omega=(0,1)\)、\(\eta_1(x)=\cos x\)、\(\eta_2(x)=\sin x\)，以及 \(u=1\)。此时平方和为一，梯度平方和也为一，额外项恰等于 \(1\)，严格为正。局部化保持零阶平方和，却会为分割函数的变化付出一阶能量。
\end{solution}

\begin{exercise}\label{b1:ex:22-caccioppoli-reaction}
设实值 \(u\in H^1_{\mathrm{loc}}(\Omega)\)、\(f\in L^2_{\mathrm{loc}}(\Omega)\)，常数 \(\lambda>0\)，且对每个实值 \(\varphi\in C_c^\infty(\Omega)\) 都有
\[
 \int_\Omega\nabla u\cdot\nabla\varphi\,\mathrm dx
 +\lambda\int_\Omega u\varphi\,\mathrm dx
 =\int_\Omega f\varphi\,\mathrm dx.
\]
证明对每个非负 \(\eta\in C_c^\infty(\Omega)\)，
\[
 \int_\Omega\eta^2|\nabla u|^2\,\mathrm dx
 +\lambda\int_\Omega\eta^2u^2\,\mathrm dx
 \le4\int_\Omega u^2|\nabla\eta|^2\,\mathrm dx
     +\lambda^{-1}\int_\Omega\eta^2f^2\,\mathrm dx.
\]
须先说明为什么可以用通常不光滑的 \(\eta^2u\) 代入测试恒等式。
\end{exercise}
% 来源：自拟；把正文弱调和Caccioppoli方法扩展到源项及正零阶项，独立说明测试合法化及吸收常数。
\begin{solution}
令 \(\psi=\eta^2u\)。光滑乘法的弱求导公式说明 \(\psi\in H^1(\Omega)\)，且它具有内部紧支集。由\cref{b1:prop:compact-sobolev-smooth-approximation}，存在 \(\psi_n\in C_c^\infty(\Omega)\) 在 \(H^1\) 中逼近它。还可以要求这些测试函数的支集包含在同一个开集 \(V\) 中，并满足 \(\overline V\Subset\Omega\)：先选 \(\zeta\in C_c^\infty(V)\)，使 \(\zeta=1\) 于 \(\operatorname{supp}\eta\) 的邻域，再把任意所述逼近列换成 \(\zeta\psi_n\)。乘子及其梯度有界，所以替换后仍趋于 \(\zeta\psi=\psi\)。

函数 \(u,\nabla u,f\) 在 \(V\) 上均平方可积。测试恒等式三项之差的绝对值分别受
\[
 \begin{gathered}
 \|\nabla u\|_{L^2(V)}\cdot\|\nabla(\psi_n-\psi)\|_{L^2(V)},\\
 \lambda\|u\|_{L^2(V)}\cdot\|\psi_n-\psi\|_{L^2(V)},\\
 \|f\|_{L^2(V)}\cdot\|\psi_n-\psi\|_{L^2(V)}
 \end{gathered}
\]
控制，都趋于零。因而恒等式对 \(\psi=\eta^2u\) 成立。

利用 \(\nabla(\eta^2u)=\eta^2\nabla u+2\eta u\nabla\eta\)，代入得
\[
 \int_\Omega\eta^2|\nabla u|^2
 +\lambda\int_\Omega\eta^2u^2
 =-2\int_\Omega\eta u\nabla u\cdot\nabla\eta
   +\int_\Omega\eta^2fu.
\]
分别使用两个带参数的平方不等式，
\[
 2|\eta|\cdot|u|\cdot|\nabla u|\cdot|\nabla\eta|
 \le\frac12\eta^2|\nabla u|^2+2u^2|\nabla\eta|^2,
\]
\[
 \eta^2|f|\cdot|u|
 \le\frac{\lambda}{2}\eta^2u^2+\frac1{2\lambda}\eta^2f^2.
\]
将它们积分并代回，梯度项和零阶项各有一半可以吸收到左端；再将两边乘以 \(2\)，恰得题中的估计。常数 \(\lambda>0\) 在第二个吸收步骤中使用，不能在最终公式里直接令 \(\lambda=0\)。
\end{solution}

\begin{exercise}[\(\star\) 选读]\label{b1:ex:22-dirichlet-minimum}
设 \(\Omega\subset\mathbb R^d\) 为非空有界开集，实值 \(g\in H^1(\Omega)\)，记
\[
 \mathcal A=g+H_0^1(\Omega),\quad
 \mathcal E(v)=\int_\Omega|\nabla v|^2\,\mathrm dx.
\]
证明对 \(u\in\mathcal A\)，下列两条等价：
\begin{equivalences}
\item\label{b1:item:22-energy-minimum}
\(\mathcal E(u)=\inf_{v\in\mathcal A}\mathcal E(v)\)。
\item\label{b1:item:22-energy-orthogonality}
对每个实值 \(\varphi\in H_0^1(\Omega)\)，都有
\[
 \int_\Omega\nabla u\cdot\nabla\varphi\,\mathrm dx=0.
\]
\end{equivalences}
再证明这样的 \(u\) 存在且唯一，并说明这里的边界条件为何不需要先定义边界迹。
\end{exercise}
% 来源：自拟；题1恢复纯梯度内积的正定性，再用第8章Hilbert Riesz求仿射闭包类的最小能量元素。
\begin{solution}
先证明等价性。若 \(u\) 为极小点，对任意 \(\varphi\in H_0^1(\Omega)\) 和实数 \(t\)，仍有 \(u+t\varphi\in\mathcal A\)，而
\[
 \mathcal E(u+t\varphi)
 =\mathcal E(u)+2t\int_\Omega\nabla u\cdot\nabla\varphi
                 +t^2\|\nabla\varphi\|_2^2.
\]
这个关于 \(t\) 的二次多项式在 \(t=0\) 取极小值，其一次项系数必须为零，所以第二条成立。

反过来，若第二条成立，任意 \(v\in\mathcal A\) 都可写成 \(v=u+\psi\)，其中 \(\psi\in H_0^1(\Omega)\)。交叉项为零，故
\[
 \mathcal E(v)=\mathcal E(u)+\|\nabla(v-u)\|_2^2
 \ge\mathcal E(u).
\]
这就证明了第一条。

为证明存在性，在 \(V=H_0^1(\Omega)\) 上定义实内积
\[
 (v,w)_V=\int_\Omega\nabla v\cdot\nabla w\,\mathrm dx.
\]
因为 \(\Omega\) 有界，可以把它放入某个有限宽度的坐标条带。习题\ref{b1:ex:22-strip-poincare}说明 \(\|\nabla v\|_2\) 在 \(V\) 上确为范数，并与 \(H^1\) 范数等价。空间 \(V\) 按定义为 \(H^1(\Omega)\) 的闭子空间，故关于原范数完备；范数等价使其关于梯度内积也完备。因此这是一个 Hilbert 空间。

定义连续线性泛函
\[
 \Lambda(\varphi)=-\int_\Omega\nabla g\cdot\nabla\varphi\,\mathrm dx.
\]
由 Cauchy--Schwarz 不等式，
\[
 |\Lambda(\varphi)|\le\|\nabla g\|_2\cdot\|\varphi\|_V.
\]
Hilbert 空间的 Riesz 表示定理\ref{b1:thm:hilbert-riesz}给出 \(h\in V\)，使
\[
 \int_\Omega\nabla h\cdot\nabla\varphi\,\mathrm dx
 =-\int_\Omega\nabla g\cdot\nabla\varphi\,\mathrm dx
 \quad(\varphi\in V).
\]
令 \(u=g+h\)，则 \(u\in\mathcal A\) 且满足正交条件，因而是极小点。

若 \(u_1,u_2\) 都为极小点，上面的能量分解给出
\(\|\nabla(u_1-u_2)\|_2=0\)。其差属于 \(H_0^1(\Omega)\)，再次使用题1的不等式便得 \(u_1=u_2\) 几乎处处。这证明在 Sobolev 等价类意义下唯一；连通性并不是这里所需的额外假设。

本题的边界条件就是 \(u-g\in H_0^1(\Omega)\)，即这个差能由内部紧支集光滑函数依 \(H^1\) 范数逼近。整个构造没有给每个边界点赋值，也没有将任意开集上的闭包条件认作某个尚未建立的零迹条件。若 \(g\) 换成同一仿射类的另一个代表，集合 \(\mathcal A\) 不变，所得唯一极小点也不变。
\end{solution}

% !TeX root = ../../Book1.tex
\chapter{迹理论与分数阶 Sobolev 空间}
\label{b1:ch:23}
\BookChapterNumber{b1:ch:23}
\begin{chapterabstract}
本章从一个边界值问题出发：域内几乎处处等价的 Sobolev 函数，怎样决定边界上的数据？先通过沿线端点与范数估计构造迹算子，证明有限指数下零迹与紧支集光滑闭包等价；再以点对差分定义分数阶 Sobolev 空间，刻画一阶迹的真实目标，并构造有界线性右逆。最后讨论平坦边界上的高阶数据相容性、法向导数在折角处的限制，以及 Besov 空间与本章差分范数的关系。
\end{chapterabstract}

\begin{learninggoals}
\begin{enumerate}
\item\label{b1:goal:23-trace}区分可测代表的边界赋值与由内部正则性确定的迹，完成半空间和有界 Lipschitz 区域上的 \(L^p\) 迹构造。
\item\label{b1:goal:23-zero}通过局部压平、零端点拼合及内部逼近，证明有限指数下的零迹刻画，说明其中的几何和指数条件。
\item\label{b1:goal:23-fractional}掌握分数阶双积分、差分表示、完备性、乘子、坐标变换及光滑稠密性，追踪非局部截断的额外项。
\item\label{b1:goal:23-boundary-space}证明 \(1<p<\infty\) 时的一阶分数迹定理与满射性，理解右逆的尺度卷积构造及边界范数的图表独立性。
\item\label{b1:goal:23-higher}在选读部分辨认高阶边界数据之间的切向相容关系，区分平坦法向和变化的法向场，并理解当前分数阶空间的 Besov 表述。
\end{enumerate}
\end{learninggoals}

本章使用\chapref{b1:ch:21}的沿线绝对连续代表，以及\chapref{b1:ch:22}的逼近、补零和一阶延拓。边界面积通过 Lipschitz 图像公式计算。因此先前的几何测度工具并未退出分析问题：它们保证不同边界参数中所说的“几乎处处”相容，也使边界积分得到正确的面积因子。

初读可先读完\secref{b1:sec:2301}，建立零迹与 \(W_0^{1,p}\) 的联系；再依次读分数阶空间的基本工具和完整的一阶迹定理。\secref{b1:sec:2304}为选读，后续内部嵌入理论不依赖其中未展开的高阶右逆或一般 Besov 理论。特别应留意两个范围：本章的 \(L^p\) 迹与零迹刻画包含 \(p=1\)，精确分数阶目标及所构造的线性右逆则只在 \(1<p<\infty\) 证明。

中文伴读可选王明新的《索伯列夫空间》\cite{Wang2013SobolevChinese}，先对照第2.4节、第2.10节的边界迹，再读第3.5节中的实指数空间与迹。该书书名采用人名的中文音译，本书正文仍沿用 Sobolev 原拼。Di Nezza、Palatucci 与 Valdinoci 的《Hitchhiker's guide to the fractional Sobolev spaces》\cite{DiNezza2012Hitchhikers}适合伴读本章的差分部分；涉及原文仅陈述或采用不同前置工具的地方，各节另行说明，不把一条参考文献当作未写出的证明。

% !TeX root = ../../Book1.tex
\section{Sobolev 迹}
\label{b1:sec:2301}

Sobolev 函数首先是域内的几乎处处等价类。给这个等价类随意补上边界值，不会改变任何体积积分，却会任意改变所谓“边界限制”。所以迹不是把某个可测代表直接放到边界上，而是一个由内部正则性决定、并对 Sobolev 范数连续的运算。

本节先把目标空间取为边界上的 \(L^p\)，并处理全部有限指数 \(1\le p<\infty\)。更精细的分数阶目标留到\secref{b1:sec:2303}。暂时不要求读者预先理解分数阶空间，也不把有限指数的零迹刻画推广到 \(p=\infty\)。

\subsection{光滑函数的边界限制}
\label{b1:subsec:230101}

设 \(\Omega\subset\mathbb R^d\) 为有界 Lipschitz 区域，记
\[
 \sigma(A)=\mathcal H^{d-1}(A\cap\partial\Omega).
\]
这里表示将 \((d-1)\) 维 Hausdorff 测度限制到边界；以后直接写 \(\int_{\partial\Omega}f\,\mathrm d\sigma\)。在\secref{b1:sec:2203}的图表中，边界参数为
\[
 \Gamma(z)=(z,\varphi(z)).
\]
由 Lipschitz 图像的面积公式，
\begin{equation}\label{b1:eq:lipschitz-boundary-measure}
 \int_{\Gamma(A)}f\,\mathrm d\sigma
 =\int_A f(\Gamma(z))\sqrt{1+|\nabla\varphi(z)|^2}\,\mathrm dz.
\end{equation}
梯度在几乎处处存在，面积因子处于 \(1\) 与 \(\sqrt{1+L^2}\) 之间。因此图表中的零测集、可积性和 \(L^p\) 范数都能与横向 Lebesgue 测度比较。有限图表覆盖又保证 \(\sigma(\partial\Omega)<\infty\)。维数为一时，\(\mathcal H^0\) 是计数测度，有界 Lipschitz 区域的边界是有限点集。

若 \(v\) 在 \(\overline\Omega\) 的一个邻域内光滑，则 \(v|_{\partial\Omega}\) 没有歧义。把它延伸到一般 Sobolev 元素之前，需要先知道这样的函数确实足够多。

\begin{proposition}\label{b1:prop:sobolev-smooth-up-to-boundary}
设 \(1\le p<\infty\)，且 \(\Omega\) 是有界 Lipschitz 区域。则
\[
 \{\,v|_\Omega\mid v\in C_c^\infty(\mathbb R^d)\,\}
\]
在 \(W^{1,p}(\Omega)\) 中稠密。
\end{proposition}
\begin{proof}
取\cref{b1:thm:lipschitz-domain-extension}的有界延拓 \(Eu\in W^{1,p}(\mathbb R^d)\)。由全空间光滑稠密性，选 \(v_j\in C_c^\infty(\mathbb R^d)\)，使 \(v_j\to Eu\) 于 \(W^{1,p}(\mathbb R^d)\)。限制到 \(\Omega\) 不增加范数，且 \((Eu)|_\Omega=u\)，所以 \(v_j|_\Omega\to u\)。
\end{proof}

这里的紧支集属于整个空间，不必留在 \(\Omega\) 内；这些近似函数可以在边界处非零。这与 \(W_0^{1,p}\) 的定义有本质区别。稠密性本身还不足以保证它们的边界限制收敛，必须另外建立边界范数估计。

\subsection{迹算子}
\label{b1:subsec:230102}

先在 \(H_+=\mathbb R^{d-1}\times(0,\infty)\) 上说明边界估计的来源。横向变量写作 \(x'\)，纵向变量写作 \(t\)。沿几乎每条竖线，Sobolev 函数具有绝对连续代表，因而有机会通过 \(t\downarrow0\) 确定边界值。

\begin{theorem}[半空间的 \(L^p\) 迹]\label{b1:thm:half-space-lp-trace}
设 \(1\le p<\infty\)。存在唯一有界线性算子
\[
 \operatorname{Tr}_+:W^{1,p}(H_+)\longrightarrow L^p(\mathbb R^{d-1}),
\]
使它对 \(C_c^\infty(\mathbb R^d)\) 的限制等于 \(u(x',0)\)。它由几乎每条竖线的绝对连续端点值给出，并满足
\begin{equation}\label{b1:eq:half-space-lp-trace}
 \|\operatorname{Tr}_+u\|_{L^p(\mathbb R^{d-1})}
 \le\|u\|_{L^p(\mathbb R^{d-1}\times(0,1))}
      +\|\partial_du\|_{L^p(\mathbb R^{d-1}\times(0,1))}.
\end{equation}
选择这些沿线代表后，对每个 \(0<t<1\)，
\begin{equation}\label{b1:eq:trace-slice-convergence}
 \|u(\cdot,t)-\operatorname{Tr}_+u\|_{L^p}
 \le t^{1-1/p}
       \|\partial_du\|_{L^p(\mathbb R^{d-1}\times(0,t))}.
\end{equation}
特别地，切片在 \(L^p\) 中趋于迹；\(p=1\) 时右侧也趋零。
\end{theorem}
\begin{proof}
取\cref{b1:thm:sobolev-acl}的共同代表。Fubini 定理保证，对几乎每个 \(x'\)，函数 \(u(x',\cdot)\) 与 \(\partial_du(x',\cdot)\) 都在 \((0,1)\) 上可积。因此其绝对连续代表延伸到端点，并可写成
\[
 u(x',t)=a(x')+\int_0^t\partial_du(x',s)\,\mathrm ds.
\]
在例外的 \(x'\) 上赋零；端点值可由 \(u(x',1/j)\) 的极限取得，故可取为可测函数。由上式，对 \(0<t<1\)，
\[
 |a(x')|\le |u(x',t)|
                  +\int_0^1|\partial_du(x',s)|\,\mathrm ds.
\]
对 \(t\in(0,1)\) 平均，再对 \(x'\) 取 \(L^p\) 范数。积分 Minkowski 不等式和区间长度为一给出\eqref{b1:eq:half-space-lp-trace}。因此 \(a\in L^p\)，令 \(\operatorname{Tr}_+u=a\)。两代表几乎处处相同，则对几乎每个 \(x'\) 沿线几乎处处相同，绝对连续代表及其端点值因而一致。这使定义在等价类层面良定义并线性。

在端点积分表示中用 Hölder 不等式，
\[
 |u(x',t)-a(x')|^p
 \le t^{p-1}\int_0^t|\partial_du(x',s)|^p\,\mathrm ds.
\]
对 \(x'\) 积分即得\eqref{b1:eq:trace-slice-convergence}。当 \(p=1\) 时，用积分三角不等式取得同一结论，右端由可积函数在缩小边界层上的积分趋零。沿线代表的选择还使这些切片对每个 \(t\) 都可测并属于 \(L^p\)。

对光滑函数，端点值当然就是通常限制。最后，将 \(u\) 偶反射到全空间，再用全空间 \(C_c^\infty\) 稠密性，说明光滑限制在 \(W^{1,p}(H_+)\) 中稠密。任意两个有界线性算子若在这一稠密类上相同，就在全空间相同，得到唯一性。
\end{proof}

这种由连续性延伸边界限制的算子，称为\term[ji-suanzi]{迹算子}[trace operator]，其输出称为函数的\term[ji]{迹}[trace]，常记为 \(\operatorname{Tr}u\) 或 \(\gamma_0u\)。本书统一以前一记号表示零阶边界数据。“迹”不是函数在任意可测代表上的边界赋值；上述定理通过内部沿线积分把它唯一地确定下来。
\symidx{\(\operatorname{Tr}u\)：Sobolev 函数的边界迹}

\subsection{\texorpdfstring{\(W^{1,p}\)}{W1,p} 的迹}
\label{b1:subsec:230103}

对于 Lipschitz 边界，竖线改成图表中横向位置固定的线段。图函数只需 Lipschitz，因为本章首先处理一阶空间，双 Lipschitz 换元已有相应弱导数规则。

\begin{theorem}[Lipschitz 区域的 \(L^p\) 迹]\label{b1:thm:lipschitz-lp-trace}
设 \(\Omega\subset\mathbb R^d\) 为有界 Lipschitz 区域，\(1\le p<\infty\)。存在唯一有界线性算子
\[
 \operatorname{Tr}:W^{1,p}(\Omega)\longrightarrow L^p(\partial\Omega,\sigma)
\]
延伸光滑函数的通常边界限制，并满足
\[
 \|\operatorname{Tr}u\|_{L^p(\partial\Omega)}
 \le C_{\Omega,p}\|u\|_{W^{1,p}(\Omega)}.
\]
若 \(u\in W^{1,p}(\Omega)\cap C(\overline\Omega)\)，则迹与该连续代表的边界限制相同。在各个局部图表内，迹也等于相应竖线的绝对连续端点值。
\end{theorem}
\begin{proof}
先设 \(d\ge2\)。取\secref{b1:sec:2203}的有限内外图表、压平映射 \(F_j\) 及光滑分割函数 \(\eta_j\)。内部项 \(\eta_0\) 在边界附近为零，其他项满足 \(\sum_{j=1}^N\eta_j=1\) 于边界。记 \(\Gamma_j(z)=F_j(z,0)\)。

对 \(u\in W^{1,p}(\Omega)\)，先将 \(\eta_ju\) 拉回正半盒。它在侧面及顶面附近为零，可以像局部延拓引理那样补零为 \(H_+\) 上的 \(W^{1,p}\) 函数 \(v_j\)。这里不跨过 \(t=0\) 补零，故不需要预先假设零边界值。定义
\[
 T_ju=(\operatorname{Tr}_+v_j)\circ\Gamma_j^{-1}
\]
于该边界片，并在其余边界上取零。由半空间迹估计、坐标换元、光滑乘法及\eqref{b1:eq:lipschitz-boundary-measure}，
\[
 \|T_ju\|_{L^p(\partial\Omega)}
 \le C_j\|u\|_{W^{1,p}(\Omega)}.
\]
有限和 \(Tu=\sum_{j=1}^N T_ju\) 因而是有界线性映射。

若 \(u\) 光滑到边界，则沿图表竖线有通常连续端点值，即
\[
 T_ju=(\eta_j|_{\partial\Omega})(u|_{\partial\Omega}).
\]
分割函数之和在边界为一，所以 \(Tu=u|_{\partial\Omega}\)。由\cref{b1:prop:sobolev-smooth-up-to-boundary}的稠密性，\(T\) 是唯一延伸光滑限制的有界算子，记为 \(\operatorname{Tr}\)。不同图表或分割函数得到的构造必然一致。

还需说明局部相容性并不限于本次选定的分割函数。给定一个图表及支集留在其内的光滑截断 \(\chi\)，取全空间光滑限制 \(u_m\to u\) 于 \(W^{1,p}(\Omega)\)。弱乘积与双 Lipschitz 换元保证 \((\chi u_m)\circ F\to(\chi u)\circ F\) 于相应正半盒。对各 \(u_m\)，两种端点解释一致；两边的 \(L^p\) 迹映射都连续，因此取极限仍一致。在任意较小边界片选取 \(\chi=1\)，即可得到所述局部端点刻画。

若 \(u\) 还有闭域连续代表，则沿这些图表竖线的端点值就是它的连续边界值，故两者一致。最后，\(d=1\) 时区域是有限个有界区间之并；在每个区间使用\secref{b1:sec:2202}的端点连续映射，再对有限端点求和，给出同一结论。
\end{proof}

\begin{proposition}\label{b1:prop:trace-multiplication}
若 \(\chi\) 在 \(\overline\Omega\) 的一个邻域内光滑，则
\[
 \operatorname{Tr}(\chi u)
 =(\chi|_{\partial\Omega})\operatorname{Tr}u
 \quad\text{于 }L^p(\partial\Omega).
\]
\end{proposition}
\begin{proof}
对光滑限制该等式只是通常乘法。对一般 \(u\)，用光滑限制依 \(W^{1,p}\) 逼近。由于 \(\overline\Omega\) 紧，\(\chi\) 及其一阶导数有界，乘法在 \(W^{1,p}(\Omega)\) 中连续；边界乘法及迹也连续，故等式可以取极限。
\end{proof}

迹定理并不说每个 \(L^p\) 边界函数都是某个 \(W^{1,p}\) 函数的迹。在 \(1<p<\infty\) 时，梯度可积性还强迫边界数据具有分数阶正则性；后两节将刻画这一附加条件。本节只先建立一个足以比较边界值、且涵盖 \(p=1\) 的连续目标空间。

\subsection{零迹与 \texorpdfstring{\(W_0^{1,p}\)}{W01,p}}
\label{b1:subsec:230104}

现在可以把上一章的闭包条件与真正的边界数据联系起来。零迹条件比“在边界上给代表赋零”强得多：它要求由内部弱导数确定的端点值为零。

\begin{lemma}\label{b1:lem:half-space-zero-trace-extension}
对 \(u\in W^{1,p}(H_+)\)、\(1\le p<\infty\)，若 \(\operatorname{Tr}_+u=0\)，则零延拓 \(\widetilde u\) 属于 \(W^{1,p}(\mathbb R^d)\)，并有
\[
 \partial_i\widetilde u=\widetilde{\partial_i u}.
\]
此外，\(u\in W_0^{1,p}(H_+)\)。
\end{lemma}
\begin{proof}
沿几乎每条竖线，\(u\) 的绝对连续代表在零点取值零。把它在负半线补零，两侧绝对连续函数的端点值相同，因而在跨过零点的紧区间上仍绝对连续；导数就是原导数的补零。沿其他坐标方向，正半空间保持原来的沿线性质，负半空间恒为零，中间整层直线在横向参数中是零测集。因此零延拓具有一个共同 ACL 代表，函数和候选导数又都属于 \(L^p\)。由 ACL 判据得到所求弱导数等式。

为得到内部光滑逼近，令 \(F=\widetilde u\)，取向内平移 \(F_h(x)=F(x-he_d)\)。函数及其弱导数的 \(L^p\) 平移连续性给出 \(F_h\to F\) 于 \(W^{1,p}(\mathbb R^d)\)，而 \(F_h=0\) 于 \(\{x_d<h\}\)。固定 \(h\)，用远处截断 \(\chi_R\) 令 \(\chi_RF_h\to F_h\)；这与全空间稠密性证明中的尾部估计相同。再以小于 \(h/2\) 的尺度磨光，所得函数的紧支集留在 \(H_+\) 内，且在全空间 Sobolev 范数中逼近 \(\chi_RF_h\)。依次令平移、尾部及磨光误差趋零，再限制回半空间，即得 \(W_0^{1,p}\) 的闭包条件。
\end{proof}

\begin{theorem}[零迹刻画]\label{b1:thm:zero-trace-characterization}
设 \(\Omega\) 为有界 Lipschitz 区域，\(1\le p<\infty\)。则
\[
 W_0^{1,p}(\Omega)
 =\{\,u\in W^{1,p}(\Omega)\mid\operatorname{Tr}u=0\,\}
 =\ker\operatorname{Tr}.
\]
\end{theorem}
\begin{proof}
若 \(u\in W_0^{1,p}(\Omega)\)，取 \(\varphi_m\in C_c^\infty(\Omega)\) 在 \(W^{1,p}\) 中逼近它。各 \(\varphi_m\) 在边界附近为零，故其迹均为零。迹的连续性给出 \(\operatorname{Tr}u=0\)。

反过来，设 \(\operatorname{Tr}u=0\)。维数为一的情形已由区间零端点刻画证明；以下设 \(d\ge2\)。仍用有限分解 \(u=\eta_0u+\sum_{j=1}^N\eta_ju\)。内部项具有内部紧支集，故由有限指数的内部磨光属于 \(W_0^{1,p}(\Omega)\)。

对一个边界项，将 \(\eta_ju\) 拉回正半盒并在半空间内补零，记为 \(v_j\)。迹的乘法及图表相容性说明 \(\operatorname{Tr}_+v_j=0\)，所以前一引理给出 \(v_j\in W_0^{1,p}(H_+)\)。取 \(a_m\in C_c^\infty(H_+)\)，使 \(a_m\to v_j\) 于 \(W^{1,p}(H_+)\)。

为在原图表内推回，选 \(\zeta\in C_c^\infty(D_j)\)，使它在 \(v_j\) 的闭支集附近等于一；这里支集可以接触 \(t=0\)，但仍紧含于外盒 \(D_j\)。光滑乘法连续，故 \(\zeta a_m\to v_j\) 于半空间 Sobolev 范数。每个 \(\zeta a_m\) 的支集同时避开 \(t=0\) 与外盒边界。

将 \(\zeta a_m\) 经 \(F_j^{-1}\) 推回，并在图表外补零，得到 \(\Omega\) 内部紧支集的 \(W^{1,p}\) 函数 \(b_m\)。虽然 Lipschitz 坐标变换不保证它光滑，但内部紧支集逼近保证 \(b_m\in W_0^{1,p}(\Omega)\)。由双 Lipschitz 换元的范数估计，\(b_m\to\eta_ju\) 于 \(W^{1,p}(\Omega)\)。闭性于是给出 \(\eta_ju\in W_0^{1,p}(\Omega)\)。有限个边界项与内部项相加，完成证明。
\end{proof}

证明中不能省略“推回后再次作内部光滑逼近”：图函数通常并不光滑，拉回坐标中的光滑函数推回去未必光滑。真正保留下来的是一阶 Sobolev 范数和内部紧支集，这已足以使用上一章的逼近定理。

有限 \(p\) 的范围也不可省去。区间函数 \(u(x)=x(1-x)\) 在两个端点取零，但\secref{b1:sec:2202}已经证明它不属于本书按范数闭包定义的 \(W_0^{1,\infty}(0,1)\)。另一方面，任意开集未必有单侧 Lipschitz 图表，上一章去掉中点的区间反例也不能被当前定理覆盖。零迹与闭包条件等价，既使用了指数范围，也使用了边界几何。

% !TeX root = ../../Book1.tex
\section{分数阶 Sobolev 空间}
\label{b1:sec:2302}

整数阶 Sobolev 范数通过弱导数衡量函数的变化。分数阶空间则直接比较两点的函数值：差值除以距离的一定幂，再对所有点对积分。这样既不必先定义一个“半阶导数”，也能记录介于函数值与一阶导数之间的正则性。稍后的迹空间将使用这一尺度，因为限制到边界后，正则性通常不再是整数阶。

Di Nezza、Palatucci 与 Valdinoci 的《Hitchhiker's guide to the fractional Sobolev spaces》\cite[Section 2 and Lemmas 5.1--5.3]{DiNezza2012Hitchhikers}给出双积分定义、指数比较及局部化工具。本节固定 \(0<s<1\)、\(1\leq p<\infty\)，函数可以取实值或复值；\(\Omega\subseteq\mathbb R^d\) 为非空开集，\(d\geq1\)。所有积分均使用 Lebesgue 测度，不在双积分前加入依赖于 \(s\) 的归一化系数。

% 实读 DiNezza2012Hitchhikers 的作者 arXiv:1104.4345v3 副本：
% 内页5--10，式(2.1)--(2.4)、Proposition2.1/2.2及Theorem2.4；
% 内页33--36，Lemma5.1--5.3及其证明。页码为作者副本内部页码，非期刊页码。
% 2.4仅陈述稠密性并转引Adams，未将其视为已读取的完整证明。
% 本节稠密性用差分映射的Lp平移连续性与远处截断另行完整展开；
% 双Lipschitz换元用本书已证测度变换界，不借后文迹理论。

\subsection{Slobodeckij 半范数}
\label{b1:subsec:230201}

\begin{definition}\label{b1:def:slobodeckij-seminorm}
设 \(u:\Omega\rightarrow\mathbb K\) 可测。定义其\term[slobodeckij-banfanshu]{Slobodeckij 半范数}[Slobodeckij seminorm]为
\begin{equation}\label{b1:eq:slobodeckij-seminorm}
 [u]_{s,p;\Omega}
 \coloneqq\left(\int_\Omega\int_\Omega
       \frac{|u(x)-u(y)|^p}{|x-y|^{d+sp}}\,\mathrm dy\,\mathrm dx\right)^{1/p},
\end{equation}
允许取值 \(+\infty\)。对角线 \(x=y\) 上的被积函数约定为零。全空间的半范数简记为 \([u]_{s,p}\)。
\end{definition}
\symidx{\([u]_{s,p;\Omega}\)：Slobodeckij 半范数}

Di Nezza 等在所引文章的 Section 2 将同一双积分称为 Gagliardo semi-norm，亦即以 Gagliardo 命名的半范数。\index[foreign]{Gagliardo semi-norm}这里的两个名称指向同一个差分条件；本书保持上述主称及不附加归一化系数的约定。

改变 \(u\) 在一个零测集 \(N\) 上的值，只会改变 \((N\times\Omega)\cup(\Omega\times N)\) 上的被积函数；这个集合由 Tonelli 定理为零测集。因此半范数只依赖几乎处处等价类。它的齐次性来自绝对值，三角不等式则是对函数
\[
 (x,y)\longmapsto\frac{u(x)-u(y)}{|x-y|^{d/p+s}}
\]
应用 \(L^p(\Omega\times\Omega)\) 中的 Minkowski 不等式。若半范数为零，Fubini 定理给出某个 \(y\in\Omega\)，使 \(u(x)=u(y)\) 对几乎每个 \(x\) 成立；反过来，常数函数的半范数为零。

这个结论不需要 \(\Omega\) 连通，因为积分也比较不同连通分支上的点。例如取 \(\Omega=(-2,-1)\cup(1,2)\)，令 \(u\) 在左区间等于 \(1\)，在右区间等于 \(0\)。它的弱导数在 \(\Omega\) 内为零，但
\[
 [u]_{s,p;\Omega}^p
 =2\int_{-2}^{-1}\int_1^2\frac{1}{|x-y|^{1+sp}}\,\mathrm dy\,\mathrm dx>0.
\]
两个区间的距离为 \(2\)，所以该积分有限。分数阶半范数不仅记录每个小邻域中的变化，还记录相隔位置之间的差异。

\begin{proposition}\label{b1:prop:fractional-difference-formula}
对 \(u\in L^p(\mathbb R^d)\)，沿用平移约定
\[
 \tau_hu(x)\coloneqq u(x-h),\qquad
 \Delta_hu\coloneqq\tau_hu-u.
\]
也就是说，\(\Delta_hu(x)=u(x-h)-u(x)\)。则
\begin{equation}\label{b1:eq:fractional-difference-formula}
 [u]_{s,p}^p
 =\int_{\mathbb R^d}\frac{\|\Delta_hu\|_p^p}{|h|^{d+sp}}\,\mathrm dh,
\end{equation}
两端允许同时为无穷。
\end{proposition}
\symidx{\(\tau_hu\)：平移，\(\tau_hu(x)=u(x-h)\)}
\symidx{\(\Delta_hu\)：有限差分，\(\Delta_hu=\tau_hu-u\)}
\begin{proof}
双积分中的函数非负。对每个固定 \(x\) 作平移换元 \(y=x-h\)，再用 Tonelli 定理交换 \(x,h\) 的积分次序，内层积分便是 \(\|\Delta_hu\|_p^p\)。无需预先假定半范数有限。
\end{proof}

因此，分数阶控制可以理解为：以不同距离平移函数，把各个平移误差按距离加权后合在一起。单个 \(h\) 上误差趋零，只是 \(L^p\) 平移连续性；这里还要求这些误差满足一个关于全部小尺度的可积条件。

\subsection{\texorpdfstring{\(W^{s,p}\)}{Ws,p}，\texorpdfstring{\(0<s<1\)}{0<s<1}}
\label{b1:subsec:230202}

\begin{definition}\label{b1:def:fractional-sobolev}
设 \(u\in L^p(\Omega)\)。若 \([u]_{s,p;\Omega}<\infty\)，则称 \(u\) 属于\term[fenshujie-sobolev-kongjian]{分数阶 Sobolev 空间}[fractional Sobolev space] \(W^{s,p}(\Omega)\)。该空间按几乎处处相等取商，赋予范数
\[
 \|u\|_{W^{s,p}(\Omega)}
 \coloneqq\left(\|u\|_{L^p(\Omega)}^p+[u]_{s,p;\Omega}^p\right)^{1/p}.
\]
\end{definition}
\symidx{\(W^{s,p}(\Omega)\)：分数阶 Sobolev 空间}

零阶项保证正定性。在无穷测度域上，它也控制远处的函数值，不能单凭双积分代替。

\begin{theorem}\label{b1:thm:fractional-sobolev-complete}
上述 \(W^{s,p}(\Omega)\) 是 Banach 空间。
\end{theorem}
\begin{proof}
记 \(G_su(x,y)=\bigl(u(x)-u(y)\bigr)/|x-y|^{d/p+s}\)，对角线上补零。映射 \(u\mapsto(u,G_su)\) 将本空间等距嵌入
\(L^p(\Omega)\times L^p(\Omega\times\Omega)\)，乘积使用两个范数的 \(p\) 次和。

设 \((u_j)\) 在 Sobolev 范数中为 Cauchy 列。由 \(L^p\) 完备性，存在 \(u\in L^p(\Omega)\) 和 \(F\in L^p(\Omega\times\Omega)\)，使 \(u_j\to u\)、\(G_su_j\to F\) 于相应 \(L^p\) 空间。可选共同子列 \((u_{j_\ell})\)，使
\[
 \sum_\ell\left(\|u_{j_\ell}-u\|_p^p+
                       \|G_su_{j_\ell}-F\|_p^p\right)<\infty.
\]
Tonelli 定理说明两组函数值误差分别几乎处处趋于零。删去 \(u_{j_\ell}\) 不趋于 \(u\) 的零测集所形成的两条乘积带后，\(G_su_{j_\ell}(x,y)\to G_su(x,y)\) 对几乎每个点对成立。因此 \(F=G_su\) 几乎处处，\(u\in W^{s,p}(\Omega)\)，而最初的整列在该范数中趋于 \(u\)。
\end{proof}

下面的三个工具用于把这个空间放到边界图表上。它们均直接控制点对积分，不要求对 \(u\) 作任何求导。

\begin{proposition}\label{b1:prop:fractional-lipschitz-multiplier}
设 \(a:\Omega\rightarrow\mathbb K\) 满足 \(|a|\leq A\) 且 Lipschitz 常数不超过 \(L\)。则乘法 \(u\mapsto au\) 在 \(W^{s,p}(\Omega)\) 上有界，具体有
\[
 [au]_{s,p;\Omega}^p
 \leq2^{p-1}A^p[u]_{s,p;\Omega}^p
       +C(d,s,p)(L^p+A^p)\|u\|_p^p,
 \quad \|au\|_p\leq A\|u\|_p.
\]
\end{proposition}
\begin{proof}
分解
\[
 a(x)u(x)-a(y)u(y)
 =a(x)\bigl(u(x)-u(y)\bigr)+u(y)\bigl(a(x)-a(y)\bigr).
\]
使用 \(|b+c|^p\leq2^{p-1}(|b|^p+|c|^p)\)，第一部分给出所列半范数项。对第二部分，利用
\(\bigl|a(x)-a(y)\bigr|\leq\min\{L|x-y|,2A\}\)。对每个 \(y\)，内层核积分不超过
\[
 L^p\int_{|h|<1}|h|^{p-d-sp}\,\mathrm dh
 +(2A)^p\int_{|h|\geq1}|h|^{-d-sp}\,\mathrm dh
 =d\omega_d\left(\frac{L^p}{p(1-s)}+\frac{(2A)^p}{sp}\right).
\]
这里用了\eqref{b1:eq:radial-one-dimensional-integral}，两个积分恰因 \(0<s<1\) 而有限。再乘以 \(|u(y)|^p\) 并积分，得到结论。
\end{proof}

\begin{proposition}\label{b1:prop:fractional-bilipschitz-change}
设 \(U,V\subseteq\mathbb R^d\) 为开集，\(\Phi:U\rightarrow V\) 为双 Lipschitz 同胚，\(\Phi\) 及其逆的 Lipschitz 常数分别不超过 \(L,M\)。则复合 \(u\mapsto u\circ\Phi\) 给出 \(W^{s,p}(V)\) 与 \(W^{s,p}(U)\) 之间的有界线性同构。正向估计为
\[
 \|u\circ\Phi\|_{L^p(U)}^p\leq M^d\|u\|_{L^p(V)}^p,
 \quad
 [u\circ\Phi]_{s,p;U}^p\leq L^{d+sp}M^{2d}[u]_{s,p;V}^p.
\]
\end{proposition}
\begin{proof}
逆映射保持 Lebesgue 零测集，故复合不依赖代表。由\eqref{b1:eq:sobolev-change-measure-bound}，对任意非负可测 \(b\)，有 \(\int_U b\bigl(\Phi(x)\bigr)\,\mathrm dx\leq M^d\int_Vb\)，从而得到零阶估计。

距离不等式 \(|\Phi(x)-\Phi(y)|\leq L|x-y|\) 给出
\[
 \frac{\bigl|u\bigl(\Phi(x)\bigr)-u\bigl(\Phi(y)\bigr)\bigr|^p}{|x-y|^{d+sp}}
 \leq L^{d+sp}
       \frac{\bigl|u\bigl(\Phi(x)\bigr)-u\bigl(\Phi(y)\bigr)\bigr|^p}{|\Phi(x)-\Phi(y)|^{d+sp}}.
\]
对右侧先在 \(x\) 变量使用测度变换界，再在 \(y\) 变量使用一次；Tonelli 定理允许这些非负积分运算，得到因子 \(M^{2d}\) 及所列半范数估计。对 \(\Phi^{-1}\) 重复论证，便得到逆方向的有界性。
\end{proof}

\begin{lemma}\label{b1:lem:compact-fractional-zero-extension}
设 \(u\in W^{s,p}(\Omega)\)，在 \(\Omega\setminus K\) 上几乎处处为零，其中 \(K\Subset\Omega\) 非空。记 \(\delta=\operatorname{dist}(K,\Omega^c)>0\)，在域外补零得到 \(\widetilde u\)。则 \(\widetilde u\in W^{s,p}(\mathbb R^d)\)，且
\begin{equation}\label{b1:eq:compact-fractional-zero-extension}
 [\widetilde u]_{s,p}^p
 \leq[u]_{s,p;\Omega}^p
       +\frac{2d\omega_d}{sp}\delta^{-sp}\|u\|_p^p,
 \quad\|\widetilde u\|_p=\|u\|_p.
\end{equation}
若 \(\Omega=\mathbb R^d\)，约定 \(\delta=\infty\)，附加项为零。
\end{lemma}
\begin{proof}
把全空间点对分成域内、域外及两个交叉部分，得到
\[
 [\widetilde u]_{s,p}^p
 =[u]_{s,p;\Omega}^p
   +2\int_K |u(x)|^p
      \left(\int_{\Omega^c}\frac{\mathrm dy}{|x-y|^{d+sp}}\right)\mathrm dx.
\]
若 \(\Omega^c\ne\varnothing\)，对 \(x\in K\) 和 \(y\in\Omega^c\) 有 \(|x-y|\geq\delta\)，故括号内不超过
\(\int_{|h|\geq\delta}|h|^{-d-sp}\,\mathrm dh
=d\omega_d\delta^{-sp}/(sp)\)。由此即得结论。空补集情形没有交叉积分，零阶范数等式则直接来自补零的定义。
\end{proof}

与整数阶内部补零不同，这里一般会增加半范数。原因是新加入的域外零值仍要与域内非零值比较；支集到边界的距离为这种相互作用提供了统一控制。边界图表中的截断必须在图表侧边留出余量，正是为了能够使用这一估计。

\subsection{基本嵌入}
\label{b1:subsec:230203}

\begin{proposition}\label{b1:prop:integer-to-fractional-sobolev}
全空间上有连续包含 \(W^{1,p}(\mathbb R^d)\subseteq W^{s,p}(\mathbb R^d)\)。具体地，
\begin{equation}\label{b1:eq:integer-fractional-estimate}
 [u]_{s,p}^p
 \leq d\omega_d\left(
    \frac{\|\nabla u\|_p^p}{p(1-s)}
    +\frac{2^p\|u\|_p^p}{sp}\right),
\end{equation}
其中 \(\|\nabla u\|_p\) 使用梯度的欧氏长度，与各偏导数 \(L^p\) 范数的有限乘积范数等价。
\end{proposition}
\begin{proof}
先设 \(u\in C_c^\infty(\mathbb R^d)\)。线段上的积分公式及 Hölder 不等式给出
\[
 \begin{aligned}
 u(x-h)-u(x)&=-\int_0^1\nabla u(x-th)\cdot h\,\mathrm dt,\\
 |u(x-h)-u(x)|^p
 &\leq |h|^p\int_0^1|\nabla u(x-th)|^p\,\mathrm dt.
 \end{aligned}
\]
对 \(x\) 积分，得到 \(\|\Delta_hu\|_p\leq|h|\cdot\|\nabla u\|_p\)。利用\cref{b1:thm:whole-space-sobolev-density}取 \(W^{1,p}\) 中的光滑逼近，函数差及梯度都在 \(L^p\) 中收敛，因而对每个固定 \(h\) 可把这一不等式传给一般 \(u\)。另一方面，平移不改变 \(L^p\) 范数，所以
\[
 \|\Delta_hu\|_p
 \leq\min\{\,|h|\cdot\|\nabla u\|_p,\ 2\|u\|_p\,\}.
\]
把它代入\eqref{b1:eq:fractional-difference-formula}。在 \(|h|<1\) 上用梯度界，在 \(|h|\geq1\) 上用函数值界；相应径向积分分别为 \(1/[p(1-s)]\) 和 \(1/(sp)\)，得到所列估计。
\end{proof}

这一论证先在全空间中进行，线段和平移都不会离开定义域。对已有有界 \(W^{1,p}\) 延拓算子的域，可以先延拓、再使用该命题并限制回原域；不能在任意开集上未经说明就沿两点间的整条线段积分。

估计中的两个分母说明了指数范围的作用：\(s\uparrow1\) 时，小尺度的径向积分接近 \(\int_0^1r^{-1}\,\mathrm dr\)；\(s\downarrow0\) 时，远处的径向积分接近 \(\int_1^\infty r^{-1}\,\mathrm dr\)。所以不能把这两个端点直接代入当前定义并认作原来的整数阶范数。适当归一化后的极限是另一项需要证明的结论。

\begin{proposition}\label{b1:prop:fractional-order-embedding}
若 \(0<t<s<1\)，则对任意开集 \(\Omega\)，有连续包含 \(W^{s,p}(\Omega)\subseteq W^{t,p}(\Omega)\)，并满足
\[
 [u]_{t,p;\Omega}^p
 \leq[u]_{s,p;\Omega}^p+\frac{2^pd\omega_d}{tp}\|u\|_p^p.
\]
若 \(\Omega\) 有界，另有更直接的估计
\( [u]_{t,p;\Omega}\leq(\operatorname{diam}\Omega)^{s-t}[u]_{s,p;\Omega}\)。
\end{proposition}
\begin{proof}
在 \(|x-y|<1\) 上，\(|x-y|^{-d-tp}\leq|x-y|^{-d-sp}\)，所以这一部分不超过 \(s\) 阶半范数的 \(p\) 次方。在余下部分，用
\( |u(x)-u(y)|^p\leq2^{p-1}(|u(x)|^p+|u(y)|^p)\)，并把每个核积分扩大到 \(|h|\geq1\)，便得附加项 \(2^pd\omega_d\|u\|_p^p/(tp)\)。如果直径有限，则直接在全部点对上使用 \(|x-y|^{(s-t)p}\leq(\operatorname{diam}\Omega)^{(s-t)p}\)，得到第二个估计。
\end{proof}

\subsection{光滑函数的稠密性}
\label{b1:subsec:230204}

对每个 \(h\)，普通 \(L^p\) 平移连续性已经成立，但若要在带奇异权的所有 \(h\) 上积分，还需要统一处理这些差分。把它们放进一个乘积空间即可完成这一点。

\begin{proposition}\label{b1:prop:fractional-translation-mollification}
设 \(u\in W^{s,p}(\mathbb R^d)\)，对上述平移有
\[
 \|\tau_zu-u\|_{W^{s,p}(\mathbb R^d)}\longrightarrow0\quad(z\to0).
\]
对\secref{b1:sec:2003}的磨光函数，还成立
\[
 u_\varepsilon=\rho_\varepsilon*u\in C^\infty(\mathbb R^d)\cap W^{s,p}(\mathbb R^d),
 \quad\|u_\varepsilon\|_{W^{s,p}}\leq\|u\|_{W^{s,p}},
 \quad u_\varepsilon\to u\text{ 于 }W^{s,p}.
\]
\end{proposition}
\begin{proof}
在 \(h\ne0\) 时定义
\[
 D_su(x,h)=\frac{\Delta_hu(x)}{|h|^{d/p+s}}
 =\frac{u(x-h)-u(x)}{|h|^{d/p+s}},
\]
在 \(h=0\) 时补零。差分积分公式说明 \(D_su\in L^p(\mathbb R^{2d})\)，其范数为 \([u]_{s,p}\)。而
\[
 D_s(\tau_zu)(x,h)=D_su(x-z,h).
\]
因此 \([\tau_zu-u]_{s,p}\to0\) 就是 \(L^p(\mathbb R^{2d})\) 中沿向量 \((z,0)\) 的平移连续性。对零阶项同样应用\cref{b1:lem:lp-translation-continuity}，得到第一条结论。

卷积的光滑性来自\secref{b1:sec:2003}。差分与积分相交换，给出
\[
 D_su_\varepsilon(x,h)
 =\int\rho_\varepsilon(z)D_su(x-z,h)\,\mathrm dz
\]
几乎处处成立；其绝对可积性可先用概率测度上的 Hölder 不等式，再用 Tonelli 定理对 \((x,h)\) 积分确认。分别对零阶项及上述差分项使用同一估计，得到
\[
 \|u_\varepsilon-u\|_{W^{s,p}}^p
 \leq\int\rho_\varepsilon(z)\|\tau_zu-u\|_{W^{s,p}}^p\,\mathrm dz.
\]
右侧只使用 \(|z|\leq\varepsilon\) 的平移误差，故由第一条结论趋于零。不减去 \(u\) 时，同一估计与平移不改变范数的事实给出 \(\|u_\varepsilon\|_{W^{s,p}}\leq\|u\|_{W^{s,p}}\)。
\end{proof}

\begin{theorem}\label{b1:thm:fractional-sobolev-smooth-density}
空间 \(C_c^\infty(\mathbb R^d)\) 在 \(W^{s,p}(\mathbb R^d)\) 中稠密。
\end{theorem}
\begin{proof}
先证明可以截去无穷远处的尾部。取 \(0\leq\chi\leq1\)、\(\chi\in C_c^\infty\bigl(B(0,2)\bigr)\)，使 \(\chi=1\) 于 \(\overline B(0,1)\)，并置 \(\chi_R(x)=\chi(x/R)\)、\(R\geq1\)。令 \(w_R=(1-\chi_R)u\)。零阶范数 \(\|w_R\|_p\) 由控制收敛定理趋于零。对差分，写成
\[
 \Delta_hw_R(x)
 =\bigl(1-\chi_R(x)\bigr)\Delta_hu(x)
   +u(x-h)\bigl(\chi_R(x)-\chi_R(x-h)\bigr).
\]
第一项除以 \(|h|^{d/p+s}\) 后，在 \((x,h)\) 上几乎处处趋于零，其 \(p\) 次方受 \(|D_su|^p\) 控制，因此相应积分趋零。

设 \(L\) 为 \(\chi\) 的 Lipschitz 常数。第二项的加权 \(p\) 次积分不超过
\[
 \begin{aligned}
 &\int_{\mathbb R^d}\int_{\mathbb R^d}
  |u(x-h)|^p\frac{\min\{L|h|/R,2\}^p}{|h|^{d+sp}}\,\mathrm dx\,\mathrm dh\\
 &=R^{-sp}\|u\|_p^p
      \int_{\mathbb R^d}\frac{\min\{L|z|,2\}^p}{|z|^{d+sp}}\,\mathrm dz.
 \end{aligned}
\]
等式先在 \(x\) 中使用平移不变性，再令 \(h=Rz\)。最后的核积分在原点附近由 \(|z|^{p-d-sp}\) 控制，在远处由 \(2^p|z|^{-d-sp}\) 控制，故有限。因 \(sp>0\)，这一项也趋于零。合并两项并使用 \(|a+b|^p\leq2^{p-1}(|a|^p+|b|^p)\)，得到 \([w_R]_{s,p}\to0\)。所以 \(\chi_Ru\to u\) 于 \(W^{s,p}\)。

给定误差 \(\eta>0\)，先取 \(R\) 使截断误差小于 \(\eta/2\)，再用前一命题磨光 \(\chi_Ru\)，使磨光误差小于 \(\eta/2\)。卷积的支集包含于 \(\operatorname{supp}\chi_R+\overline B(0,\varepsilon)\)，仍为紧集。所得函数属于 \(C_c^\infty\)，三角不等式完成证明。
\end{proof}

截断证明的第二项不能省去：即使 \(u\) 在远处很小，截断函数的两个取值不同仍会参与点对积分。估计中的 \(R^{-sp}\) 把这种非局部误差控制住，而不是把它当成普通的函数值尾部。

\begin{corollary}\label{b1:cor:compact-fractional-smooth-approximation}
设 \(u\in W^{s,p}(\Omega)\) 在内部紧集 \(K\) 外几乎处处为零。对任意包含 \(K\) 的开集 \(V\subseteq\Omega\)，存在 \(u_j\in C_c^\infty(V)\)，使 \(u_j\to u\) 于 \(W^{s,p}(\Omega)\)。
\end{corollary}
\begin{proof}
由\cref{b1:lem:compact-fractional-zero-extension}先补零得到 \(\widetilde u\in W^{s,p}(\mathbb R^d)\)。前面的磨光命题保证 \(\rho_\varepsilon*\widetilde u\to\widetilde u\) 于全空间范数。取 \(\varepsilon<\operatorname{dist}(K,V^c)\)，其支集便留在 \(V\) 内；限制回 \(\Omega\) 不增加范数误差。若 \(K\) 为空，直接取零函数列。
\end{proof}

\begin{corollary}\label{b1:cor:fractional-domain-smooth-density}
对任意开集 \(\Omega\)，空间 \(C^\infty(\Omega)\cap W^{s,p}(\Omega)\) 在 \(W^{s,p}(\Omega)\) 中稠密。
\end{corollary}
\begin{proof}
取\cref{b1:cor:precompact-smooth-partition}给出的光滑单位分解 \((\psi_i)\) 与局部有限开覆盖 \((V_i)\)，其中 \(\psi_i\in C_c^\infty(V_i)\)、\(V_i\Subset\Omega\)。每个 \(\psi_i\) 是有界 Lipschitz 函数，故由乘子命题，\(\psi_i u\in W^{s,p}(\Omega)\)，且在紧集 \(\operatorname{supp}\psi_i\) 外几乎处处为零。给定 \(\eta>0\)，前一推论允许取 \(v_i\in C_c^\infty(V_i)\)，使
\[
 \|v_i-\psi_i u\|_{W^{s,p}(\Omega)}<\eta2^{-i-1}.
\]
局部有限和 \(v=\sum_i v_i\) 在 \(\Omega\) 内光滑。另一方面，由完备性，误差级数 \(\sum_i(v_i-\psi_i u)\) 在 \(W^{s,p}(\Omega)\) 中收敛到某个 \(e\)，且 \(\|e\|_{W^{s,p}}\leq\eta/2\)。在每个内部相对紧开集上，级数只有有限个非零项，故其和等于 \(v-u\)；范数收敛蕴含局部 \(L^p\) 收敛，因而 \(e=v-u\) 几乎处处。于是 \(v\in W^{s,p}(\Omega)\)，并达到所需误差。
\end{proof}

这里的局部有限光滑和未必具有紧支集。因此最后一个推论不能直接改写成 \(C_c^\infty(\Omega)\) 的稠密性；域内光滑逼近与要求逼近函数在边界附近为零，仍是不同的问题。下一节将把本节的乘子、换元和内部补零估计用于边界图表，使分数阶范数成为迹空间的具体描述。

% !TeX root = ../../Book1.tex
\section{Lipschitz 边界上的迹空间}
\label{b1:sec:2303}

迹属于边界上的 \(L^p\)，只说明它的大小可以控制，还没有说明边界不同位置的函数值如何相互约束。一阶内部导数会给边界留下一个分数阶的差分条件。本节证明：当 \(1<p<\infty\) 时，有界 Lipschitz 区域的一阶迹空间恰为 \(W^{1-1/p,p}(\partial\Omega)\)，而且每个这样的边界数据都有线性、有界地选取的内部延拓。

\ProseParagraphBreak

证明按四步推进：先确认边界分数阶范数不依赖图表与分割函数的选择，再把半空间迹估计增强到 \(W^{1-1/p,p}\)；随后在半空间中令磨光尺度随高度趋零，构造迹的右逆；最后用有限图表与截断把这一构造拼回 Lipschitz 区域。也就是
\[
 \text{图表独立性}
 \longrightarrow \text{半空间分数阶迹}
 \longrightarrow \text{半空间右逆}
 \longrightarrow \text{区域上的有限拼接}.
\]

Di Nezza、Palatucci 与 Valdinoci 的《Hitchhiker's guide to the fractional Sobolev spaces》\cite[Section 2, Proposition 3.8]{DiNezza2012Hitchhikers}介绍了这一联系，并用 Fourier 方法证明二次可积情形。本节改用差分、一个一维积分不等式及尺度卷积，直接处理 \(1<p<\infty\)。
% 实读：DiNezza2012Hitchhikers，作者公开副本 arXiv:1104.4345，
% 第2节末迹空间说明（副本第11页），Proposition 3.8及完整证明（副本第19--21页）。
% 该命题的证明限于p=2；第2节末的一般p陈述没有给出这里所需的完整证明。
% 图表组织另据已实读的第5节；本节Hardy--差分与尺度卷积路线在正文独立展开，
% 不把上述p=2证明写成一般p证明的来源，也不调用后续嵌入或高阶边界理论。

\subsection{局部图表示}
\label{b1:subsec:230301}

先设 \(d\ge2\)，记 \(n=d-1\)，并令
\(\sigma=\mathcal H^n\llcorner\partial\Omega\)。取\secref{b1:sec:2203}的有限内外图表：在刚性坐标中，
\[
 F_j(z,t)=(z,\varphi_j(z)+t),\quad
 D_j=B_j\times(-2h_j,2h_j),\quad
 B_j=(-2\rho_j,2\rho_j)^n,
\]
其中 \(F_j(D_j)\cap\Omega=F_j(D_j\cap H_+)\)。边界参数化为
\(\Gamma_j(z)=F_j(z,0)\)。固定非负函数 \(\chi_j\in C_c^\infty(V_j)\)，其中 \(V_j\) 是对应内图表，使
\[
 \sum_{j=1}^N\chi_j(x)=1\quad(x\in\partial\Omega).
\]
这些函数可由前节构造所用的边界分割项取得；内部项在边界上为零。

\ProseParagraphBreak

\cref{b1:thm:graph-area-formula}给出
\[
 \int_{\Gamma_j(B_j)}q\,\mathrm d\sigma
 =\int_{B_j}q\bigl(\Gamma_j(z)\bigr)
       \sqrt{1+|\nabla\varphi_j(z)|^2}\,\mathrm dz
\]
对非负可测 \(q\) 成立。若 \(\varphi_j\) 的 Lipschitz 常数为 \(L_j\)，面积因子几乎处处介于 \(1\) 与 \(\sqrt{1+L_j^2}\) 之间。因此图上与参数域上的零测集相互对应，\(L^p\) 范数也只相差受控常数。

\begin{definition}\label{b1:def:boundary-fractional-sobolev}
给定 \(0<s<1\)、\(1\le p<\infty\) 及 \(f\in L^p(\partial\Omega,\sigma)\)，将
\[
 f_j(z)=\chi_j\bigl(\Gamma_j(z)\bigr)f\bigl(\Gamma_j(z)\bigr)
 \quad(z\in B_j)
\]
在 \(B_j\) 外补零。若每个 \(f_j\) 都属于 \(W^{s,p}(\mathbb R^n)\)，则称 \(f\) 属于\term[bianjie-fenshujie-sobolev-kongjian]{边界分数阶 Sobolev 空间}[fractional Sobolev space on the boundary]，记为 \(W^{s,p}(\partial\Omega)\)，并规定
\[
 \|f\|_{W^{s,p}(\partial\Omega)}
 \coloneqq\left(\sum_{j=1}^N\|f_j\|_{W^{s,p}(\mathbb R^n)}^p\right)^{1/p}.
\]
\end{definition}
\symidx{\(W^{s,p}(\partial\Omega)\)：用局部图表定义的边界分数阶空间}

分割系数不能从定义中直接删去：将图表边缘处不消失的函数补零，可能引入原函数没有的跳跃。这里每个 \(\chi_j\) 的支集都与外图表的边缘保持正距离。

\begin{proposition}\label{b1:prop:trace-chart-independence}
上述空间不依赖有限图表与分割函数的选择，不同选择给出的范数等价。它在所定范数下完备，而且连续嵌入 \(L^p(\partial\Omega,\sigma)\)。
\end{proposition}
\begin{proof}
由 \(f=\sum_j\chi_j f\)、有限和的不等式及面积因子的上下界，
\[
 \|f\|_{L^p(\sigma)}^p
 \le C\sum_j\|f_j\|_{L^p(\mathbb R^n)}^p
 \le C\|f\|_{W^{s,p}(\partial\Omega)}^p.
\]
这也说明定义中的表达式是范数。

考虑另一套图表 \(\widetilde\Gamma_k\) 与分割函数 \(\widetilde\chi_k\)。在两个边界片相交处，过渡坐标
\[
 \Psi_{jk}=\Gamma_j^{-1}\circ\widetilde\Gamma_k
\]
是两个开子集之间的双 Lipschitz 同胚：图参数化为 Lipschitz 映射，其逆就是相应刚性坐标中的横向投影。待比较的第 \(k\) 个局部函数是有限和
\[
 (\widetilde\chi_k f)\circ\widetilde\Gamma_k
 =\sum_j
   (\widetilde\chi_k\circ\widetilde\Gamma_k)(f_j\circ\Psi_{jk}),
\]
其中每项只需在相应的重叠参数域中解释。

两处分割函数的支集相交部分是重叠图表内的紧集。因而可以在重叠参数域选一个光滑紧支集函数，在该紧集上等于 \(1\)。在右端对应项前乘上这个辅助函数并不改变该项，却使它的支集与重叠域边缘保持正距离。依次使用\cref{b1:prop:fractional-bilipschitz-change}、\cref{b1:prop:fractional-lipschitz-multiplier}及\cref{b1:lem:compact-fractional-zero-extension}，得到该项补零后的 \(W^{s,p}(\mathbb R^n)\) 范数不超过 \(C_{jk}\|f_j\|_{W^{s,p}(\mathbb R^n)}\)。有限求和得到一方向的范数界，交换两套图表得到反向界。

最后，若 \(f^{(m)}\) 在边界范数下为 Cauchy 列，第一步保证它在 \(L^p(\sigma)\) 中收敛到某个 \(f\)。每个局部函数 \(f_j^{(m)}\) 在 \(W^{s,p}(\mathbb R^n)\) 中为 Cauchy 列，由\cref{b1:thm:fractional-sobolev-complete}存在极限。其 \(L^p\) 极限必为 \((\chi_j f)\circ\Gamma_j\) 的零延拓，所以 \(f\) 属于边界空间，且 \(f^{(m)}\) 在定义的范数中收敛到 \(f\)。
\end{proof}

\subsection{半空间迹定理}
\label{b1:subsec:230302}

本节在半空间公式中将 \(\operatorname{Tr}_+\) 简记为 \(\operatorname{Tr}\)。以下令 \(H_+=\mathbb R^n\times(0,\infty)\)，并固定 \(1<p<\infty\) 及 \(s=1-1/p\)。估计边界差分时，把两个边界点抬到与其距离相同的高度。两段竖向差需要下面的\term[hardy-budengshi]{Hardy 不等式}[Hardy inequality]。

\begin{lemma}[Hardy]\label{b1:lem:trace-hardy-inequality}
若 \(a\ge0\) 且 \(a\in L^p(0,\infty)\)，其中 \(1<p<\infty\)，则
\[
 \int_0^\infty r^{-p}\left(\int_0^r a(t)\,\mathrm dt\right)^p\mathrm dr
 \le\left(\frac p{p-1}\right)^p\int_0^\infty a(r)^p\,\mathrm dr.
\]
\end{lemma}
\begin{proof}
先设 \(a\) 有界且支集紧含于 \((0,\infty)\)，令 \(A(r)=\int_0^r a(t)\,\mathrm dt\)。在零点附近 \(A=0\)，在充分大的 \(r\) 处 \(A\) 为常数；由于 \(p>1\)，分部积分的两端边界项都为零。记左端为 \(I\)，则
\[
 I=\frac p{p-1}\int_0^\infty
       \left(\frac{A(r)}r\right)^{p-1}a(r)\,\mathrm dr
 \le\frac p{p-1}I^{(p-1)/p}\|a\|_p.
\]
若 \(I>0\)，约去相应幂次即得结论；\(I=0\) 时也成立。对一般 \(a\)，取
\[
 a_m=\min\{a,m\}\mathbf1_{(1/m,m)},
\]
再由单调收敛定理取极限。
\end{proof}

\begin{theorem}[半空间分数阶迹]\label{b1:thm:half-space-fractional-trace}
设 \(n\ge1\)、\(1<p<\infty\)。半空间的迹算子满足
\[
 \operatorname{Tr}:W^{1,p}(H_+)
 \longrightarrow W^{1-1/p,p}(\mathbb R^n),\quad
 \|\operatorname{Tr}u\|_{W^{1-1/p,p}}
 \le C_{n,p}\|u\|_{W^{1,p}(H_+)}.
\]
\end{theorem}
\begin{proof}
先取整个空间的光滑紧支集函数在闭半空间上的限制，记 \(f(x)=u(x,0)\)。给定 \(h=r\omega\)，其中 \(r>0\)、\(|\omega|=1\)，分解
\[
 \begin{split}
 \Delta_{r\omega}f(x)
 ={}&f(x-r\omega)-u(x-r\omega,r)\\
    &+u(x-r\omega,r)-u(x,r)+u(x,r)-f(x).
 \end{split}
\]
令 \(a(t)=\|\partial_tu(\cdot,t)\|_p\)，并令
\(b(t)=\sum_{i=1}^n\|\partial_i u(\cdot,t)\|_p\)。由微积分基本定理、积分的 Minkowski 不等式及平移不变性，两段竖向差的 \(L^p\) 范数各不超过 \(\int_0^r a(t)\,\mathrm dt\)，而水平差满足
\[
 \|u(\cdot-r\omega,r)-u(\cdot,r)\|_p
 \le r\,b(r).
\]
所以
\[
 \|\Delta_{r\omega}f\|_p^p
 \le C_p\left(\int_0^r a(t)\,\mathrm dt\right)^p
       +C_p r^p b(r)^p,
\]
这里沿用\cref{b1:prop:fractional-difference-formula}的记号
\(\Delta_hf=\tau_hf-f\)，不再把 \(\Delta_h\) 临时改作平移算子。

由于 \(n+sp=n+p-1\)，\cref{b1:prop:fractional-difference-formula}与极坐标变换给出
\[
 [f]_{s,p}^p
 =\int_{\mathbb S^{n-1}}\int_0^\infty
       r^{-p}\|\Delta_{r\omega}f\|_p^p\,\mathrm dr\,
       \mathrm d\mathcal H^{n-1}(\omega).
\]
代入上面的三项估计并用\cref{b1:lem:trace-hardy-inequality}，可得
\[
 [f]_{s,p}^p
 \le C_{n,p}\int_0^\infty\bigl(a(t)^p+b(t)^p\bigr)\,\mathrm dt
 \le C_{n,p}\sum_{i=1}^{n+1}\|\partial_i u\|_{L^p(H_+)}^p.
\]
再加上\cref{b1:thm:half-space-lp-trace}的零阶迹估计，就得到所需全范数界。

对一般 \(u\in W^{1,p}(H_+)\)，先偶反射，再在整个空间中作光滑紧支集逼近，最后限制回半空间。这是\secref{b1:sec:2301}构造 \(L^p\) 迹时使用的逼近。上述估计作用于两项之差，使光滑迹成为 \(W^{s,p}(\mathbb R^n)\) 中的 Cauchy 列；其极限在 \(L^p\) 中又等于既有的 \(\operatorname{Tr}u\)。完备性于是给出结论，并保留同一个范数界。
\end{proof}

指数 \(1-1/p\) 已在积分中出现：边界的径向体积因子将差分核变成 \(r^{-p}\)，恰好与竖向积分的 Hardy 估计匹配。若把 \(p=1\) 直接代入，所用 Hardy 常数发散，而 \(s=0\) 也离开本章分数阶双积分定义的范围；这不是把目标空间写成 \(W^{0,1}\) 就能补上的端点论证。

\subsection{迹的满射性与右逆}
\label{b1:subsec:230303}

仅有迹估计还不能说明每个边界函数都能实现。下面先在半空间中令“磨光尺度等于高度”，同时恢复原数据并控制竖向导数；再把所得右逆逐图表搬回、截断并作有限求和，得到一般有界 Lipschitz 区域上的右逆。

\begin{theorem}\label{b1:thm:half-space-trace-right-inverse}
设 \(n\ge1\)、\(1<p<\infty\)，并令 \(s=1-1/p\)。存在有界线性映射
\[
 R_+:W^{s,p}(\mathbb R^n)\longrightarrow W^{1,p}(H_+),
 \quad \operatorname{Tr}(R_+f)=f.
\]
因此半空间分数阶迹算子为满射。
\end{theorem}
\begin{proof}
固定非负 \(\rho\in C_c^\infty(B(0,1))\)，使 \(\int\rho=1\)，并取 \(\eta\in C_c^\infty(\mathbb R)\)，使 \(\eta=1\) 于 \([-1,1]\)，支集包含于 \((-2,2)\)。对 \(t>0\) 规定
\[
 \rho_t(h)=t^{-n}\rho(h/t),\quad
 v(x,t)=(\rho_t*f)(x),\quad
 (R_+f)(x,t)=\eta(t)v(x,t).
\]
由卷积不等式，\(\|v(\cdot,t)\|_p\le\|f\|_p\)，所以零阶项和竖向求导产生的 \(\eta'(t)v(x,t)\) 都有由 \(\|f\|_p\) 控制的全域 \(L^p\) 范数。

对切向导数，核为 \(t^{-n-1}(\partial_i\rho)(h/t)\)；对高度导数，核为
\[
 \partial_t\rho_t(h)=t^{-n-1}\psi(h/t),\quad
 \psi(z)=-n\rho(z)-z\cdot\nabla\rho(z).
\]
这些核都支撑在 \(|h|<t\) 内，绝对值不超过 \(Ct^{-n-1}\)，且积分为零。最后一点对 \(\psi\) 由紧支集函数的分部积分得到。若 \(K_t\) 是其中任一导数核，因而可以写
\[
 (K_t*f)(x)=\int K_t(h)\Delta_hf(x)\,\mathrm dh.
\]
Hölder 不等式给出
\[
 \begin{split}
 |(K_t*f)(x)|^p
 &\le\left(\int|K_t(h)|\,\mathrm dh\right)^{p-1}
       \int|K_t(h)|\cdot|f(x-h)-f(x)|^p\,\mathrm dh\\
 &\le Ct^{-n-p}\int_{|h|<t}|f(x-h)-f(x)|^p\,\mathrm dh.
 \end{split}
\]
先对 \(x\) 积分，再对 \(0<t<2\) 积分，由 Tonelli 定理，
\[
 \begin{split}
 \int_0^2\|K_t*f\|_p^p\,\mathrm dt
 &\le C\int_{|h|<2}\|\Delta_hf\|_p^p
                  \int_{|h|}^2t^{-n-p}\,\mathrm dt\,\mathrm dh\\
 &\le C_{n,p}\int_{\mathbb R^n}
        \frac{\|\Delta_hf\|_p^p}{|h|^{n+p-1}}\,\mathrm dh
 =C_{n,p}[f]_{s,p}^p.
 \end{split}
\]
函数 \(v\) 在 \(t>0\) 处光滑，因而这些可积的经典导数也是弱导数。结合零阶项及 \(\eta'\) 项的估计，得到
\[
 \|R_+f\|_{W^{1,p}(H_+)}
 \le C_{n,p}\|f\|_{W^{s,p}(\mathbb R^n)}.
\]
所有核与截断都已固定，所以该映射线性。最后，近似恒等核保证 \(v(\cdot,t)\to f\) 于 \(L^p\)，且 \(\eta(t)=1\) 于充分小的正高度。由半空间 \(L^p\) 迹的切片刻画，\(\operatorname{Tr}(R_+f)=f\)。
\end{proof}

这里的 \(R_+\) 是迹算子的右逆：先延入半空间，再取迹，恢复边界数据；反向复合不必恢复任意内部函数，因为零迹函数还可以自由加入。

\begin{theorem}[Lipschitz 边界的分数阶迹]\label{b1:thm:lipschitz-fractional-trace}
设 \(d\ge2\)，\(\Omega\subset\mathbb R^d\) 为有界 Lipschitz 区域，且 \(1<p<\infty\)。既有的迹算子增强为有界满射
\[
 \operatorname{Tr}:W^{1,p}(\Omega)
 \longrightarrow W^{1-1/p,p}(\partial\Omega).
\]
它有一个有界线性右逆 \(R_\Omega\)。此外，边界范数等价于最小延拓范数：
\[
 \|f\|_{W^{1-1/p,p}(\partial\Omega)}
 \asymp
 \inf\bigl\{\,\|u\|_{W^{1,p}(\Omega)}
        \mid \operatorname{Tr}u=f\,\bigr\}.
\]
\end{theorem}
\begin{proof}
先证明有界性。对每个固定边界系数 \(\chi_j\)，将 \(\chi_j u\) 拉回 \(D_j\cap H_+\)，再在侧面和顶面之外补零，得到整个半空间上的函数 \(v_j\)。由边界截断的构造，存在固定紧集 \(L_j\Subset D_j\)，使 \(v_j=0\) 几乎处处成立于 \(H_+\setminus L_j\)；\(L_j\) 可以接触 \(t=0\)，但与图表盒的侧面和顶面有正距离。\secref{b1:sec:2203}局部图表延拓证明中的测试分解因而适用于此处。因此
\[
 \|v_j\|_{W^{1,p}(H_+)}
 \le C_j\|u\|_{W^{1,p}(\Omega)}.
\]
由\cref{b1:thm:lipschitz-lp-trace}的局部端点刻画，它的半空间迹是
\((\chi_j\operatorname{Tr}u)\circ\Gamma_j\) 的零延拓。该相容性已经通过跨边界光滑逼近建立，不要求图函数光滑。半空间分数阶迹估计于是逐项给出
\[
 \|\operatorname{Tr}u\|_{W^{1-1/p,p}(\partial\Omega)}^p
 \le C\|u\|_{W^{1,p}(\Omega)}^p.
\]

现给定 \(f\in W^{s,p}(\partial\Omega)\)，其中 \(s=1-1/p\)。令 \(f_j\) 为定义中的局部零延拓，取 \(w_j=R_+f_j\)。记
\[
 K_j=\operatorname{supp}(\chi_j\circ\Gamma_j)\Subset B_j,
\]
并固定 \(\beta_j\in C_c^\infty(D_j)\)，使
\(\beta_j(z,0)=1\) 于 \(K_j\)。函数
\[
 q_j(z,t)=\beta_j(z,t)w_j(z,t),\quad t>0,
\]
在侧面和顶面附近为零。由 \(w_j(\cdot,t)\to f_j\) 于 \(L^p\)，以及 \(\beta_j(\cdot,t)\) 一致趋于 \(\beta_j(\cdot,0)\)，可得
\(\operatorname{Tr}q_j=\beta_j(\cdot,0)f_j=f_j\)。
把 \(q_j\) 经 \(F_j^{-1}\) 推回 \(\Omega\cap F_j(D_j)\)，再在 \(\Omega\) 的其余部分补零，记结果为 \(Q_jf\)。由于 \(\beta_j\in C_c^\infty(D_j)\)，\(q_j\) 在固定紧集 \(\operatorname{supp}\beta_j\) 外几乎处处为零；这一紧集与侧面和顶面有正距离，故局部测试恒等式可以拼接。这里始终只在 \(\Omega\) 内补零，并没有跨越 \(\partial\Omega\) 作零延拓。故
\[
 \|Q_jf\|_{W^{1,p}(\Omega)}
 \le C_j\|f_j\|_{W^{s,p}(\mathbb R^n)}.
\]
迹与局部坐标的相容性给出
\(\operatorname{Tr}(Q_jf)=\chi_j f\)，包括图表外的零部分。

规定
\[
 R_\Omega f=\sum_{j=1}^N Q_jf.
\]
所有图表、边界系数、卷积核与 \(\beta_j\) 都固定不变，所以它是线性的。有限求和与上述估计给出有界性，而
\[
 \operatorname{Tr}(R_\Omega f)
 =\sum_{j=1}^N\chi_j f=f
\]
证明右逆性质和满射性。最后，对任意迹为 \(f\) 的 \(u\)，迹估计给出边界范数不超过 \(C\|u\|_{W^{1,p}}\)；反向取 \(u=R_\Omega f\) 即得最小延拓范数的等价关系。
\end{proof}

这个迹空间允许的边界不连续性取决于 \(p\)。例如在二维半空间上取边界函数 \(f=\mathbf1_{(0,1)}\)。\cref{b1:ex:23-interval-characteristic}的差分计算表明，它属于 \(W^{s,p}(\mathbb R)\) 当且仅当 \(sp<1\)。代入 \(s=1-1/p\)，就知道它能作为 \(W^{1,p}(H_+)\) 函数的迹，当且仅当 \(1<p<2\)。因此不能从内部具有一阶弱导数，直接推断边界数据连续。

\ProseParagraphBreak

本节没有把两个端点纳入分数阶结论。对 \(p=1\)，\secref{b1:sec:2301}仍给出取值于 \(L^1(\partial\Omega)\) 的连续迹与零迹刻画；本章不据此声称 \(L^1(\partial\Omega)\) 是迹的精确像空间，也不声称该迹满射或具有同样的线性右逆。对 \(p=\infty\)，既不能代入有限 \(p\) 的双积分范数，也不能用该范数的稠密性取极限；应回到局部 Lipschitz 代表与边界几何讨论。

\ProseParagraphBreak

若 \(d=1\)，有界 Lipschitz 开集的每个边界点都在边界中孤立，边界又紧，故只有有限多个端点。迹空间此时就是这些端点上的有限维数据空间：在每个区间分量上作连接两端数据的仿射函数，便得到线性右逆，其范数由有限多个区间长度控制。这个论证也处理相应端点指数，但不把 \(n=0\) 代入上面的欧氏差分核。

% !TeX root = ../../Book1.tex
\OptionalSection{高阶迹理论}{b1:sec:2304}

一阶迹决定函数在边界上的零阶数据。若域内还有更多弱导数，便可以继续对这些导数取迹；但所得数据不是一组互不相关的边界函数。切向导数必须与原函数迹的切向变化一致，法向导数则承担另外的边界信息。本节先在平坦边界上证明这种结构，再说明弯曲边界的几何限制。完整的高阶边界数据延拓不在本节证明范围内。

\subsection{\texorpdfstring{\(W^{k,p}\)}{Wk,p} 的边界数据}
\label{b1:subsec:230401}

设 \(d\ge2\)、\(n=d-1\)、\(H_+=\mathbb R^n\times(0,\infty)\)，并取整数 \(k\ge1\)、\(1<p<\infty\)。若 \(u\in W^{k,p}(H_+)\)，那么对每个 \(|\alpha|\le k-1\)，都有 \(\partial^\alpha u\in W^{1,p}(H_+)\)，所以
\[
 g_\alpha=\operatorname{Tr}_+(\partial^\alpha u)
\]
已经有明确定义。各个 \(g_\alpha\) 至少属于 \(W^{s,p}(\mathbb R^n)\)，其中 \(s=1-1/p\)；这是把一阶分数迹定理用于 \(\partial^\alpha u\) 的直接结果。下面还要利用不同阶数据之间的相容关系。

先约定非整数高阶空间的记号。对整数 \(m\ge0\)、\(0<s<1\)，令
\[
 W^{m+s,p}(\mathbb R^n)
 =\{\,f\in W^{m,p}(\mathbb R^n)
       \mid \partial^\beta f\in W^{s,p}(\mathbb R^n)
               \text{ 对所有 }|\beta|=m\,\},
\]
其范数取为
\[
 \|f\|_{W^{m+s,p}}^p
 =\|f\|_{W^{m,p}}^p+
     \sum_{|\beta|=m}[\partial^\beta f]_{s,p}^p.
\]
当 \(m=0\) 时，这正是\secref{b1:sec:2302}的定义。这里对最高整数阶弱导数增加分数阶差分条件，而不是把原先的整数阶定义中的“阶数”形式地换成实数。

\begin{proposition}[平坦边界的高阶迹]\label{b1:prop:flat-higher-traces}
设 \(u\in W^{k,p}(H_+)\)、\(1<p<\infty\)。对 \(|\alpha|\le k-2\) 及切向方向 \(1\le i\le n\)，有
\begin{equation}\label{b1:eq:trace-tangential-compatibility}
 \partial_i g_\alpha=g_{\alpha+e_i}
 \quad\text{作为 }\mathbb R^n\text{ 上的弱导数}.
\end{equation}
若记
\[
 \gamma_j u=\operatorname{Tr}_+(\partial_d^j u),
 \quad 0\le j\le k-1,
\]
则
\begin{equation}\label{b1:eq:higher-flat-trace-bound}
 \gamma_j u\in W^{k-j-1/p,p}(\mathbb R^n),\qquad
 \|\gamma_j u\|_{W^{k-j-1/p,p}}
 \le C_{d,k,p}\|u\|_{W^{k,p}(H_+)}.
\end{equation}
\end{proposition}
\begin{proof}
先证切向相容性。固定 \(|\alpha|\le k-2\)，令 \(v=\partial^\alpha u\)、\(w=\partial_i v\)。两者都属于 \(W^{1,p}(H_+)\)。由 Fubini 定理及切向弱导数恒等式，对几乎每个 \(t>0\)，切片满足
\[
 \int_{\mathbb R^n}v(x',t)\partial_i\phi(x')\,\mathrm dx'
 =-\int_{\mathbb R^n}w(x',t)\phi(x')\,\mathrm dx'
 \quad\bigl(\phi\in C_c^\infty(\mathbb R^n)\bigr).
\]
这里可以先对可数稠密的测试函数族取得共同的好切片，再由局部积分连续性扩展到全部测试函数。选取这样的 \(t_\ell\downarrow0\)。由\eqref{b1:eq:trace-slice-convergence}，\(v(\cdot,t_\ell)\to g_\alpha\)、\(w(\cdot,t_\ell)\to g_{\alpha+e_i}\) 于 \(L^p(\mathbb R^n)\)。在积分恒等式中取极限，便得\eqref{b1:eq:trace-tangential-compatibility}。

令 \(m=k-1-j\)。反复使用切向相容性，得到
\[
 \partial_{x'}^\beta\gamma_j u
 =\operatorname{Tr}_+
       (\partial_{x'}^\beta\partial_d^j u)
 \quad(|\beta|\le m).
\]
右侧由半空间的 \(L^p\) 迹估计控制，因为内部导数总阶数至多 \(k-1\)，还能再取一阶弱导数。这证明 \(\gamma_ju\in W^{m,p}(\mathbb R^n)\) 并控制其整数阶范数。

对 \(|\beta|=m\)，内部函数
\(\partial_{x'}^\beta\partial_d^j u\in W^{1,p}(H_+)\)。再用\cref{b1:thm:half-space-fractional-trace}，得到它的迹属于 \(W^{s,p}\)，且相应半范数受 \(\|u\|_{W^{k,p}}\) 控制。有限多个 \(\beta\) 求和，即得 \(\gamma_ju\in W^{m+s,p}\) 及所需估计，因为 \(m+s=k-j-1/p\)。
\end{proof}

例如 \(u\in W^{2,p}(H_+)\) 时，零阶迹属于 \(W^{2-1/p,p}\)，法向方向的一阶导数迹属于 \(W^{1-1/p,p}\)，而每个切向导数迹必须等于零阶迹的对应弱导数。不能先任意指定函数值及全部一阶偏导数的边界值，再期待存在一个内部函数实现它们。

上述证明只需对已有整数阶导数重复使用一阶迹，不依赖高阶函数在闭半空间上的光滑稠密性。它给出的是连续性和数据相容性；同时为所有 \(\gamma_j\) 构造有界右逆，还需要能匹配多项边界数据的延拓公式，不能把零阶的尺度卷积原样使用 \(k\) 次就视为完成。

\subsection{法向导数}
\label{b1:subsec:230402}

在上半空间的边界上，外单位法向为 \(\nu=-e_d\)。因此对 \(u\in W^{2,p}(H_+)\)，\term[faxiang-daoshu]{法向导数}[normal derivative]的边界迹为
\[
 \partial_\nu u=-\operatorname{Tr}_+(\partial_du)=-\gamma_1u.
\]
更高阶沿固定法向的导数相应带有 \((-1)^j\)。本节的 \(\gamma_j\) 始终按正 \(x_d\) 方向定义；写外法向数据时另加这个符号，避免把方向约定隐藏在记号中。

设 \(\Omega\) 为有界 Lipschitz 区域。边界图表几乎处处给出外单位法向 \(\nu\)，其分量本性有界。对 \(u\in W^{2,p}(\Omega)\)，各一阶导数属于 \(W^{1,p}(\Omega)\)，所以可以在 \(L^p(\partial\Omega)\) 中定义
\[
 \partial_\nu u
 =\sum_{i=1}^d\nu_i\operatorname{Tr}(\partial_i u).
\]
由 \(|\nu|=1\) 及 \(L^p\) 迹估计，这是受 \(\|u\|_{W^{2,p}(\Omega)}\) 控制的边界函数。若 \(u\) 与边界足够光滑，它与通常的方向导数相同。

但是，本性有界乘子未必保持分数阶 Sobolev 空间。\secref{b1:sec:2302}的乘子结论还要求 Lipschitz 控制；在折角处，\(\nu\) 可以跳跃。因此不能仅凭每个 \(\operatorname{Tr}(\partial_i u)\) 属于 \(W^{1-1/p,p}(\partial\Omega)\)，就断言它们与 \(\nu_i\) 相乘后仍属于该空间。

一个具体例子是正方形 \(\Omega=(0,1)^2\) 与光滑函数 \(u(x,y)=x\)。在下边上 \(\partial_\nu u=0\)，在左边上 \(\partial_\nu u=-1\)。在原点附近用折线弧长作参数，这个边界函数有一个大小为 \(1\) 的跳跃。若它属于 \(W^{s,p}\)，其半范数必须包含有限的跨角点积分；然而该积分从下方控制
\[
 \int_0^\varepsilon\int_0^\varepsilon
       \frac{\mathrm da\,\mathrm db}{(a+b)^{1+sp}}.
\]
在 \(0<a<\varepsilon/2\)、\(a<b<2a\) 上积分，便得到常数倍的
\(\int_0^{\varepsilon/2}a^{-sp}\,\mathrm da\)，当 \(sp\ge1\) 时发散。因此取 \(s=1-1/p\)、\(p\ge2\)，这个法向数据不属于预期的分数阶空间，尽管 \(u\) 在整个平面上光滑。

若边界具有更多光滑性，法向及坐标过渡也具有更多可控导数，平坦边界的高阶估计才有望在图表之间保持。对高阶法向导数，还须说明是在边界点固定方向作高阶导数，还是把法向场延到邻域后反复作用；后者可能引入法向场的导数。高阶迹理论同时涉及函数正则性、边界正则性和数据约定，三者不能只靠一个“边界光滑”的笼统表述代替。

\subsection{Besov 空间作为自然迹空间的预告}
\label{b1:subsec:230403}

一阶迹的指数 \(1-1/p\)，来自对法向厚度的积分与切向差分尺度之间的配合。允许更一般的光滑阶数和尺度求和方式后，便进入\term[besov-kongjian]{Besov 空间}[Besov space]。这里只介绍本章已经能够理解的一种差分定义，不预用 Fourier 变换。

先设 \(0<s<1\)、\(1\le p<\infty\) 及 \(1\le q<\infty\)。对 \(f\in L^p(\mathbb R^n)\)，令
\[
 [f]_{B^s_{p,q}}
 =\left(
   \int_{\mathbb R^n}
    \frac{\|\Delta_hf\|_{L^p}^q}{|h|^{sq}}
    \,\frac{\mathrm dh}{|h|^n}
   \right)^{1/q}.
\]
要求这个量有限，并以 \(\|f\|_p+[f]_{B^s_{p,q}}\) 为范数，得到本节所用的非齐次 Besov 空间 \(B^s_{p,q}(\mathbb R^n)\)。另外，当 \(q=\infty\) 时规定
\[
 [f]_{B^s_{p,\infty}}
 =\operatorname*{ess\,sup}_{h\ne0}
       \frac{\|\Delta_hf\|_{L^p}}{|h|^s}.
\]
由于大尺度差分受 \(2\|f\|_p\) 控制，也可以只在 \(|h|<1\) 上积分或取本性上确界，得到等价范数。

两个指数承担不同的平均：\(p\) 先度量固定平移 \(h\) 下、空间变量中的差分大小；\(q\) 再度量这些大小随平移尺度的分布。当 \(q=p\) 时，Tonelli 定理直接给出
\[
 [f]_{B^s_{p,p}}^p
 =\int_{\mathbb R^n}\int_{\mathbb R^n}
       \frac{|\Delta_hf(x)|^p}{|h|^{n+sp}}\,\mathrm dx\,\mathrm dh
 =[f]_{s,p}^p.
\]
因此在当前 \(0<s<1\) 的范围内，
\[
 B^s_{p,p}(\mathbb R^n)=W^{s,p}(\mathbb R^n)
\]
作为集合相同，所列非齐次范数等价。一阶迹空间也就可以写成 \(B^{1-1/p}_{p,p}\)，并不是又多出了一种不同的边界条件。

对更高的非整数阶数，需要使用高阶差分或先取整数阶导数；这些定义之间的等价、一般 Besov 空间的插值以及完整高阶迹右逆，属于进一步的函数空间理论。可从 Triebel 的《Theory of Function Spaces》\cite{Triebel1983Theory}继续阅读，也可在 Adams 与 Fournier 的《Sobolev Spaces》\cite{Adams2003Sobolev}中衔接经典 Sobolev 理论。这里不将尚未展开的一般 Besov 迹定理用作后续正文的前置。

至此，边界值的定义、其一阶分数正则性及零边界闭包已连成一条证明链。下一章回到区域内部：在这些逼近与延拓工具准备好以后，将研究梯度怎样控制函数的振幅和可积性。这些嵌入估计与本章的迹估计相互补充，但处理的是不同的测度和不同位置上的函数信息。


\begin{chaptersummary}
迹由内部沿线绝对连续性和 Sobolev 范数估计确定，而不是任意边界赋值。对有界 Lipschitz 区域，有限指数下迹的核正是 \(W_0^{1,p}\)；反向证明必须在图表中控制支集，再将不一定光滑的坐标推回函数作内部逼近。

分数阶范数比较所有点对，因而能检测跨越边界或连通分支的差异。乘子、补零、截断和坐标变换都要控制这种非局部比较。对于 \(1<p<\infty\)，法向积分和切向尺度配合给出 \(1-1/p\) 阶边界正则性；高度随磨光尺度变化的构造，又把任意这样的数据延回内部。

高阶数据进一步具有切向相容性。平坦边界的结论不能不加条件地转移到有折角的法向场上。本章因此把已证明的一阶边界理论、平坦高阶估计和进一步的 Besov 理论分开。接下来的嵌入与紧性，将主要利用前两章和本章已经完成的逼近、延拓及零边界工具。
\end{chaptersummary}

\BookExercises

% \chapref{b1:ch:23}习题临时稿；不另设 BookExercises 入口。
% 八题均为自拟的工具变式；来源为已建立的正文接口与此处独立展开的计算，不冒认教材原题编号。
% 仅用本章及前章工具，不使用后续嵌入、Fourier或尚未建立的Besov表示定理。

\begin{exercise}\label{b1:ex:23-interval-characteristic}
设 \(0<s<1\)、\(1\le p<\infty\)，令 \(u=\mathbf1_{(0,1)}\)，把它看作整个 \(\mathbb R\) 上的函数。证明
\[
 u\in W^{s,p}(\mathbb R)\quad\Longleftrightarrow\quad sp<1,
\]
并在有限情形计算半范数 \([u]_{s,p}^p\) 的精确值。说明临界条件 \(sp=1\) 不能归入有限的一侧。
\end{exercise}
% 来源：自拟；直接计算定义中的跨跳跃双积分，区别于原区间内常函数的零半范数。
\begin{solution}
函数的 \(L^p\) 范数为 \(1\)，问题完全在于半范数。只有一个变量在 \((0,1)\)、另一个变量在其外时，函数差才非零。由对称性及 Tonelli 定理，
\[
 \begin{aligned}
 [u]_{s,p}^p
 &=2\int_0^1\left(
     \int_{-\infty}^0\frac{\mathrm dy}{(x-y)^{1+sp}}
    +\int_1^\infty\frac{\mathrm dy}{(y-x)^{1+sp}}
             \right)\mathrm dx\\
 &=\frac2{sp}\int_0^1
       \bigl(x^{-sp}+(1-x)^{-sp}\bigr)\,\mathrm dx.
 \end{aligned}
\]
当 \(sp<1\) 时，两个端点积分各等于 \(1/(1-sp)\)，故
\[
 [u]_{s,p}^p=\frac4{sp(1-sp)}.
\]
若 \(sp\ge1\)，至少一个端点积分发散；在 \(sp=1\) 时就是 \(\int_0^1x^{-1}\,\mathrm dx\) 的对数发散。因此条件为严格不等式。

若仅在 \((0,1)\) 内计算，\(u\) 是常函数，双积分为零。全空间半范数还比较边界两侧的函数值，正是这些跨越端点的配对检测到了跳跃。
\end{solution}

\begin{exercise}\label{b1:ex:23-smooth-integer-limit}
设 \(d\ge2\)、\(1\le p<\infty\)，实值 \(u\in C_c^\infty(\mathbb R^d)\)。记 \(\mathrm d\sigma\) 为单位球面上的 \(\mathcal H^{d-1}\) 测度。证明
\[
 \lim_{s\uparrow1}(1-s)[u]_{s,p}^p
 =\frac1p\int_{\mathbb S^{d-1}}\int_{\mathbb R^d}
            |\nabla u(x)\cdot\omega|^p\,\mathrm dx\,\mathrm d\sigma(\omega).
\]
特别地，\(p=2\) 时右端等于
\[
 \frac{\sigma(\mathbb S^{d-1})}{2d}
       \int_{\mathbb R^d}|\nabla u|^2\,\mathrm dx.
\]
本题只要求证明光滑紧支集情形，不把所得极限未经说明地推广到一般函数。
\end{exercise}
% 来源：自拟；差分公式、Taylor余项与径向积分的组合计算，不预借一般整数阶极限定理。
\begin{solution}
由\cref{b1:prop:fractional-difference-formula}，半范数可写成平移差的积分。先估计 \(|h|\ge1\) 的部分。平移等距性与三角不等式给出
\[
 \|u(\,\cdot+h)-u\|_p^p\le2^p\|u\|_p^p,
\]
所以
\[
 (1-s)\int_{|h|\ge1}
 \frac{\|u(\,\cdot+h)-u\|_p^p}{|h|^{d+sp}}\,\mathrm dh
 \le\frac{(1-s)2^p\sigma(\mathbb S^{d-1})}{sp}\|u\|_p^p
 \longrightarrow0.
\]

对 \(0<r<1\)、\(\omega\in\mathbb S^{d-1}\)，记
\[
 A(r,\omega)=
 \left\|\frac{u(\,\cdot+r\omega)-u}{r}\right\|_p^p,
 \quad A(0,\omega)=\|\nabla u\cdot\omega\|_p^p.
\]
微积分基本定理给出
\[
 \frac{u(x+r\omega)-u(x)}r
 =\int_0^1\nabla u(x+tr\omega)\cdot\omega\,\mathrm dt.
\]
因为 \(\nabla u\) 一致连续且有紧支集，右端与
\(\nabla u(x)\cdot\omega\) 之差在 \(L^p(\mathbb R^d)\) 中趋零，而且对全部单位向量 \(\omega\) 一致：函数差的支集包含在一个固定有界集内，其上确界受 \(\nabla u\) 在距离不超过 \(r\) 的两点间的振幅控制。各差商的 \(L^p\) 范数又一致有界，故
\[
 \sup_\omega|A(r,\omega)-A(0,\omega)|\longrightarrow0.
\]

令 \(B(r)=\int_{\mathbb S^{d-1}}A(r,\omega)\,\mathrm d\sigma(\omega)\)，则 \(B\) 有界且 \(B(r)\to B(0)\)。用径向积分处理 \(|h|<1\) 的部分，得到
\[
 (1-s)\int_0^1r^{p(1-s)-1}B(r)\,\mathrm dr.
\]
这段权重的总积分为 \(1/p\)。给定 \(0<\delta<1\)，在 \((0,\delta)\) 上用 \(B(r)\) 接近 \(B(0)\) 控制误差；在 \((\delta,1)\) 上，权重质量为
\[
 \frac{1-\delta^{p(1-s)}}p\longrightarrow0.
\]
先选小 \(\delta\)，再令 \(s\uparrow1\)，便知所列积分趋于 \(B(0)/p\)。结合远处部分趋零，得到第一条公式。

当 \(p=2\) 时，球面对称性给出
\(\int\omega_i\omega_j\,\mathrm d\sigma=0\)（\(i\ne j\)），各
\(\int\omega_i^2\,\mathrm d\sigma\) 相等，而它们之和为
\(\sigma(\mathbb S^{d-1})\)。展开 \(|\nabla u\cdot\omega|^2\) 后逐项积分，便得第二条公式。有限差分在小尺度下读出方向导数，系数 \(1/p\) 来自径向权重，而不是可以任意省略的归一化。
\end{solution}

\begin{exercise}\label{b1:ex:23-fractional-modulus-truncation}
设 \(0<s<1\)、\(1\le p<\infty\)，实值 \(u\in W^{s,p}(\Omega)\)。
\begin{enumerate}
\item\label{b1:item:23-fractional-modulus}
证明 \(|u|\in W^{s,p}(\Omega)\)，且 \([\,|u|\,]_{s,p;\Omega}\le[u]_{s,p;\Omega}\)。当 \(p=2\) 时，求两个半范数平方之差，并证明半范数相等当且仅当 \(u\) 几乎处处不变号。
\item\label{b1:item:23-fractional-truncation}
令 \(T_M(t)=\max\bigl(-M,\min(t,M)\bigr)\)。证明 \(T_M(u)\to u\) 于 \(W^{s,p}(\Omega)\)。
\end{enumerate}
\end{exercise}
% 来源：自拟；非线性收缩、变号能量损失与可积双积分上的控制收敛，不重复正文线性乘子及光滑稠密性。
\begin{solution}
由反三角不等式，
\[
 \bigl||u(x)|-|u(y)|\bigr|\le|u(x)-u(y)|.
\]
将其 \(p\) 次幂除以核的分母并积分，即得半范数估计；零阶 \(L^p\) 范数则完全相等。

若 \(p=2\)，对任意实数 \(a,b\)，
\[
 |a-b|^2-\bigl||a|-|b|\bigr|^2
 =4|a|\cdot|b|\,\mathbf1_{\{ab<0\}}.
\]
两边积分得到
\[
 [u]_{s,2;\Omega}^2-[\,|u|\,]_{s,2;\Omega}^2
 =4\iint_{\{u(x)u(y)<0\}}
       \frac{|u(x)|\cdot|u(y)|}{|x-y|^{d+2s}}\,\mathrm dx\,\mathrm dy.
\]
两个半范数都有限，所以这里没有无穷量相减。若 \(u\) 几乎处处不变号，右侧为零。反之，若正值集和负值集都具有正测度，则其乘积集具有正的乘积测度，核上的被积函数在那里严格为正，故积分严格为正。这证明等号的刻画。

对于第二问，记 \(R_M(t)=t-T_M(t)\)。这个实函数连续、分段线性，各段斜率为 \(0\) 或 \(1\)，因此
\[
 |R_M(a)-R_M(b)|\le|a-b|,\quad |R_M(a)|\le|a|.
\]
又因为有限实数 \(a\) 满足 \(R_M(a)\to0\)，控制收敛先给出
\(\|R_M(u)\|_p\to0\)。对于双积分，几乎每一对 \((x,y)\) 的分子也趋于零，并且
\[
 \frac{\bigl|R_M\bigl(u(x)\bigr)-R_M\bigl(u(y)\bigr)\bigr|^p}{|x-y|^{d+sp}}
 \le\frac{|u(x)-u(y)|^p}{|x-y|^{d+sp}}.
\]
右端按假设可积，再次应用控制收敛，得到
\([R_M(u)]_{s,p;\Omega}\to0\)。这就是所求范数收敛。两处收敛都使用有限 \(p\) 的积分控制，并非对所有函数值统一取上确界。
\end{solution}

\begin{exercise}\label{b1:ex:23-boundary-concentration}
设 \(d\ge2\)、\(1\le p<\infty\)，并记
\(\Omega=(0,1)^{d-1}\times(0,1)\)、\(B'=(0,1)^{d-1}\)。
取非零 \(f\in C_c^\infty(B')\) 及满足 \(\chi(0)=1\)、\(\operatorname{supp}\chi\subset(-1,1)\) 的 \(\chi\in C_c^\infty(\mathbb R)\)。对 \(0<\varepsilon<1\)，令
\[
 u_\varepsilon(x',t)=f(x')\chi(t/\varepsilon).
\]
计算 \(u_\varepsilon\) 的体积 \(L^p\) 范数、边界迹的 \(L^p\) 范数以及
\(\|\partial_tu_\varepsilon\|_{L^p(\Omega)}\) 随 \(\varepsilon\) 的变化。由此证明：不存在与 \(u\) 无关的常数 \(C\)，使所有在 \(\overline\Omega\) 附近光滑的函数都满足
\[
 \|\operatorname{Tr}u\|_{L^p(\partial\Omega)}
 \le C\|u\|_{L^p(\Omega)}.
\]
说明这个例子为什么不反驳本章的 \(W^{1,p}\) 迹估计。
\end{exercise}
% 来源：自拟；边界层缩放例，用来分开体积收敛、迹收敛与一阶导数控制。
\begin{solution}
记两个严格为正的常数
\[
 A_p=\left(\int_0^1|\chi(r)|^p\,\mathrm dr\right)^{1/p},
 \quad
 B_p=\left(\int_0^1|\chi'(r)|^p\,\mathrm dr\right)^{1/p}.
\]
\(A_p>0\) 来自 \(\chi(0)=1\) 及连续性；\(B_p>0\) 是因为 \(\chi\) 在 \(1\) 的邻域为零，不可能在 \([0,1]\) 上恒为常数。代换 \(t=\varepsilon r\) 得
\[
 \|u_\varepsilon\|_{L^p(\Omega)}
   =\varepsilon^{1/p}A_p\|f\|_{L^p(B')},
 \quad
 \|\partial_tu_\varepsilon\|_{L^p(\Omega)}
   =\varepsilon^{1/p-1}B_p\|f\|_{L^p(B')}.
\]
底面上的限制为 \(f\)，顶面上的限制为零；由于 \(f\) 的支集包含在 \(B'\) 的内部，各侧面上的限制也为零。棱边的 \(\mathcal H^{d-1}\) 测度为零，故光滑函数的迹满足
\[
 \|\operatorname{Tr}u_\varepsilon\|_{L^p(\partial\Omega)}
 =\|f\|_{L^p(B')}>0.
\]
若题设估计成立，约去这个正数就会得到
\(1\le C A_p\varepsilon^{1/p}\)，令 \(\varepsilon\downarrow0\) 即矛盾。

另一方面，法向导数的范数在 \(p>1\) 时趋于无穷，在 \(p=1\) 时保持为严格正的常数；切向导数则满足
\[
 \|\partial_{x_i}u_\varepsilon\|_{L^p(\Omega)}
 =\varepsilon^{1/p}A_p\|\partial_{x_i}f\|_{L^p(B')}.
\]
因此 \(u_\varepsilon\) 虽然在体积 \(L^p\) 中趋于零，却不在 \(W^{1,p}\) 中趋于零。迹估计使用的一阶导数项正好检测到了这层越来越薄的过渡区域。
\end{solution}

\begin{exercise}\label{b1:ex:23-trace-right-inverse-freedom}
设 \(\Omega\subset\mathbb R^d\) 为非空有界 Lipschitz 开集，\(d\ge2\)、\(1<p<\infty\)。记
\[
 X=W^{1,p}(\Omega),\quad
 Y=W^{1-1/p,p}(\partial\Omega),\quad
 T=\operatorname{Tr}:X\longrightarrow Y,
\]
并取\cref{b1:thm:lipschitz-fractional-trace}给出的一个有界线性右逆 \(R:Y\to X\)。
\begin{enumerate}
 \item\label{b1:item:23-trace-right-inverse-difference}
 证明有界线性算子 \(S:Y\to X\) 也是 \(T\) 的右逆，当且仅当
 \(S=R+K\)，其中 \(K:Y\to W_0^{1,p}(\Omega)\) 是有界线性算子。实际构造一个非零的 \(K\)，说明右逆一般不唯一。
 \item\label{b1:item:23-trace-affine-class}
 对任意 \(g\in Y\)，证明
 \[
  \{u\in X:Tu=g\}=Rg+W_0^{1,p}(\Omega).
 \]
 说明为什么用另一个右逆替换 \(R\)，不会改变给定边界值的这一函数类。
\end{enumerate}
\end{exercise}
% 来源：自拟；把迹满射和零迹刻画组合为仿射 Dirichlet 类，不预用后续偏微分方程存在定理。
\begin{solution}
由\cref{b1:thm:zero-trace-characterization}，
\(\ker T=W_0^{1,p}(\Omega)\)。若 \(TS=\operatorname{id}_Y=TR\)，则
\[
 T(S-R)=0,
\]
所以 \(K=S-R\) 的像包含于 \(W_0^{1,p}(\Omega)\)。这个子空间采用 \(X\) 的限制范数，因此 \(K\) 作为取值于该子空间的算子仍然有界。反过来，若 \(K\) 的像包含于 \(\ker T\)，则 \(T(R+K)=\operatorname{id}_Y\)。

下面给出非零的选择。记 \(\sigma=\mathcal H^{d-1}|_{\partial\Omega}\)，取非零 \(w\in C_c^\infty(\Omega)\)，并令
\[
 \ell(g)=\int_{\partial\Omega}g\,\mathrm d\sigma,
 \quad K(g)=\ell(g)w.
\]
有界 Lipschitz 开集的边界具有有限且严格为正的 \(\sigma\) 测度，而本章选定的边界范数控制 \(L^p(\partial\Omega)\) 范数。因此 Hölder 不等式给出
\[
 |\ell(g)|\le
 \sigma(\partial\Omega)^{1-1/p}\|g\|_{L^p(\partial\Omega)}
 \le C\|g\|_Y.
\]
常函数 \(1\) 是 \(X\) 中常函数的迹，故属于 \(Y\)，并且
\(\ell(1)=\sigma(\partial\Omega)>0\)。于是 \(K\ne0\)，且其像包含于
\(\operatorname{span}\{w\}\subset W_0^{1,p}(\Omega)\)。这便构造出两个不同的右逆 \(R\) 与 \(R+K\)。

最后，\(Tu=g\) 等价于 \(T(u-Rg)=0\)，也就等价于
\(u-Rg\in W_0^{1,p}(\Omega)\)，这证明第二问的集合恒等式。如果换用 \(S\)，则
\(Sg-Rg\in W_0^{1,p}(\Omega)\)；用子空间中的一个元素平移这个子空间不会改变所得陪集。因此边界条件本身是确定的，选择右逆只是选择了其中一个起点。
\end{solution}

\begin{exercise}\label{b1:ex:23-circle-energy}
把单位圆盘 \(D\subset\mathbb R^2\) 记为复平面中的 \(\{z:|z|<1\}\)。对 \(n\ge1\)，令
\[
 U_n(z)=z^n,\quad g_n(e^{i\theta})=e^{in\theta}.
\]
本题单独采用圆周上的弦长双积分
\[
 Q(g)=\int_{\mathbb S^1}\int_{\mathbb S^1}
        \frac{|g(x)-g(y)|^2}{|x-y|^2}\,
                 \mathrm d\mathcal H^1(x)\,\mathrm d\mathcal H^1(y).
\]
证明
\[
 \|g_n\|_{L^2(\mathbb S^1)}^2=2\pi,\quad
 Q(g_n)=4\pi^2 n,\quad
 \int_D|\nabla U_n|^2\,\mathrm dx=2\pi n,
\]
其中复值函数的梯度平方是
\(|\partial_xU_n|^2+|\partial_yU_n|^2\)。
只用有限几何级数与积分计算完成证明，不使用 Fourier 展开定理。
\end{exercise}
% 来源：自拟；圆周高频数据和显式光滑延拓的精确计算。Q 为题内指定双积分，不冒认为正文固定图表范数的原值。
\begin{solution}
因为 \(|g_n|=1\)，第一个等式直接由圆周长度得到。在几乎每一对参数 \((\theta,\varphi)\) 处，有限几何级数给出
\[
 \frac{|e^{in\theta}-e^{in\varphi}|^2}
      {|e^{i\theta}-e^{i\varphi}|^2}
 =\left|\sum_{j=0}^{n-1}e^{ij(\theta-\varphi)}\right|^2.
\]
对于整数 \(j,\ell\)，初等积分满足
\[
 \int_0^{2\pi}e^{i(j-\ell)h}\,\mathrm dh
 =
 \begin{cases}
  2\pi,&j=\ell,\\
  0,&j\ne\ell.
 \end{cases}
\]
展开有限和的模平方，再逐项积分，得到
\[
 \int_0^{2\pi}\left|\sum_{j=0}^{n-1}e^{ijh}\right|^2
                        \,\mathrm dh=2\pi n.
\]
对固定的 \(\varphi\)，以 \(h=\theta-\varphi\) 代换并用周期性，内层积分仍等于 \(2\pi n\)。外层再积分一次便得 \(Q(g_n)=4\pi^2n\)。

在圆盘内直接求导，
\[
 \partial_x U_n=n z^{n-1},\quad
 \partial_y U_n=in z^{n-1}.
\]
所以两项的模平方之和为 \(2n^2|z|^{2n-2}\)，极坐标积分给出
\[
 \int_D|\nabla U_n|^2\,\mathrm dx
 =2n^2\int_0^{2\pi}\int_0^1r^{2n-2}r\,\mathrm dr\,\mathrm d\theta
 =2\pi n.
\]
在这个显式函数族上，\(Q(g_n)\) 恰为内部梯度能量的 \(2\pi\) 倍，而边界 \(L^2\) 范数不随 \(n\) 增长。这反映了分数阶边界量对振荡的敏感性。本题计算的是指定的弦长双积分；正文通过固定图表和分割选取的边界范数，数值还取决于那些选择。
\end{solution}

\begin{exercise}\label{b1:ex:23-higher-zero-endpoints}
设 \(I=(0,1)\)、\(k\ge1\) 为整数、\(1\le p<\infty\)。以
\[
 W_0^{k,p}(I)=\overline{C_c^\infty(I)}^{\,W^{k,p}(I)}
\]
定义高阶零边界空间。给定 \(u\in W^{k,p}(I)\)，对 \(0\le j<k\)，用
\(u^{(j)}\) 的绝对连续代表定义它在 \(0,1\) 处的端点值。证明下列条件等价。
\begin{equivalences}
 \item\label{b1:item:23-higher-zero-closure}
 \(u\in W_0^{k,p}(I)\)。
 \item\label{b1:item:23-higher-zero-jets}
 对所有 \(0\le j<k\)，都有 \(u^{(j)}(0)=u^{(j)}(1)=0\)。
\end{equivalences}
证明充分性时，应给出边界截断的高阶导数估计；仅仅说“在端点附近截去一小段”不足以推出 \(W^{k,p}\) 收敛。
\end{exercise}
% 来源：自拟；一维绝对连续积分表示与内部紧支集逼近的高阶组合，不借助一般高阶迹空间。
\begin{solution}
每个 \(u^{(j)}\)、\(j<k\)，都属于 \(W^{1,p}(I)\)，因而具有题设的绝对连续代表。若 \(u_n\in C_c^\infty(I)\) 且 \(u_n\to u\) 于 \(W^{k,p}(I)\)，则
\[
 u_n^{(j)}\longrightarrow u^{(j)}
       \quad\text{于 }W^{1,p}(I),\quad 0\le j<k.
\]
一维代表估计，即命题\ref{b1:prop:sobolev-interval-representative}，使这个收敛也成为闭区间上的一致收敛。由于 \(u_n^{(j)}\) 的两端值都是零，极限的两端值也为零，必要性得证。

反过来，设所有这些端点值都为零。反复使用微积分基本定理，对于 \(0\le r<k\)，有
\[
 u^{(r)}(x)=\frac1{(k-r-1)!}
             \int_0^x(x-t)^{k-r-1}u^{(k)}(t)\,\mathrm dt.
\]
将右端视为截断积分核与 \(u^{(k)}\) 的卷积，并使用 Minkowski 积分不等式，可得对 \(0<a<1\)，
\[
 \|u^{(r)}\|_{L^p(0,a)}
 \le\frac{a^{k-r}}{(k-r)!}\,
             \|u^{(k)}\|_{L^p(0,a)}.
\]
具体地，只需把 \(u^{(k)}|_{(0,a)}\) 向外补零；核
\(t^{k-r-1}\mathbf1_{(0,a)}(t)/(k-r-1)!\) 的 \(L^1\) 范数为
\(a^{k-r}/(k-r)!\)，平移保持 \(L^p\) 范数，于是逐个平移积分便得到上式。这一步同样适用于 \(p=1\)。从右端点向左积分，还得到
\[
 \|u^{(r)}\|_{L^p(1-a,1)}
 \le\frac{a^{k-r}}{(k-r)!}\,
             \|u^{(k)}\|_{L^p(1-a,1)}.
\]

取 \(0<\varepsilon<1/4\)，选取 \(0\le\eta_\varepsilon\le1\) 的光滑截断，使它在
\([2\varepsilon,1-2\varepsilon]\) 上为 \(1\)，在两个端点的 \(\varepsilon\) 邻域内为零，并满足
\[
 \|\eta_\varepsilon^{(\ell)}\|_\infty
          \le C_\ell\varepsilon^{-\ell},\quad 1\le\ell\le k.
\]
记 \(A_\varepsilon=(0,2\varepsilon)\cup(1-2\varepsilon,1)\)。
对 \(0\le j\le k\)，Leibniz 公式给出
\[
 \bigl((1-\eta_\varepsilon)u\bigr)^{(j)}
 =(1-\eta_\varepsilon)u^{(j)}
  -\sum_{\ell=1}^j\binom j\ell
           \eta_\varepsilon^{(\ell)}u^{(j-\ell)}.
\]
第一项的 \(L^p\) 范数不超过
\(\|u^{(j)}\|_{L^p(A_\varepsilon)}\)，因 \(p<\infty\) 而趋于零。和式中
\(j-\ell<k\)，可以应用刚才的端点估计，得到
\[
 \begin{aligned}
 \|\eta_\varepsilon^{(\ell)}u^{(j-\ell)}\|_{L^p(I)}
 &\le C\varepsilon^{-\ell}
       \|u^{(j-\ell)}\|_{L^p(A_\varepsilon)}\\
 &\le C\varepsilon^{k-j}
       \|u^{(k)}\|_{L^p(A_\varepsilon)}
 \longrightarrow0.
 \end{aligned}
\]
即使 \(j=k\)，最后一项也因 \(u^{(k)}\in L^p(I)\) 的积分绝对连续性而趋于零。因此
\(\eta_\varepsilon u\to u\) 于 \(W^{k,p}(I)\)。

每个 \(\eta_\varepsilon u\) 都有包含在 \(I\) 内部的紧支集。由命题%
\ref{b1:prop:compact-sobolev-smooth-approximation}，它可以在 \(W^{k,p}(I)\) 中由
\(C_c^\infty(I)\) 函数逼近。先选 \(\varepsilon\) 控制截断误差，再作光滑逼近，便得到
\(u\in W_0^{k,p}(I)\)。这里有限 \(p\) 保证了边界层上的最高阶导数范数趋零；证明并未涵盖 \(p=\infty\)。
\end{solution}

\begin{exercise}\label{b1:ex:23-smooth-normal-data}
设 \(d\ge2\)，\(H=\mathbb R^{d-1}\times(0,\infty)\)，其边界外法向为
\(\nu=-e_d\)。本题只考虑光滑边界数据。
\begin{enumerate}
 \item\label{b1:item:23-independent-smooth-normal-data}
 给定任意 \(a,b\in C_c^\infty(\mathbb R^{d-1})\)，构造一个在
 \(\overline H\) 附近光滑、支集在 \(\overline H\) 中紧的函数 \(U\)，满足
 \[
  U(x',0)=a(x'),\quad \partial_\nu U(x',0)=b(x').
 \]
 说明该函数属于每个整数阶的 \(W^{k,p}(H)\)，其中 \(1\le p<\infty\)。
 \item\label{b1:item:23-smooth-tangential-compatibility}
 若还要求对 \(1\le i<d\) 有
 \(\partial_{x_i}U(x',0)=c_i(x')\)，其中 \(c_i\in C_c^\infty(\mathbb R^{d-1})\)，证明这些要求可同时满足，当且仅当
 \(c_i=\partial_{x_i}a\) 对每个 \(i\) 成立。
\end{enumerate}
\end{exercise}
% 来源：自拟；显式光滑边界数据延拓，分开自由法向数据与受约束的切向数据，不声称给出了粗糙数据的高阶右逆。
\begin{solution}
取 \(\chi\in C_c^\infty(\mathbb R)\)，使 \(\chi\) 在 \(0\) 的邻域恒等于 \(1\)。令
\[
 U(x',t)=\chi(t)\bigl(a(x')-t b(x')\bigr),\quad t\ge0.
\]
这个公式也在 \(t<0\) 定义光滑函数，因而确实给出了边界附近的光滑延拓。它在
\(\overline H\) 中的支集包含于
\[
 \bigl(\operatorname{supp}a\cup\operatorname{supp}b\bigr)
       \times\bigl(\operatorname{supp}\chi\cap[0,\infty)\bigr),
\]
后者为紧集。边界处 \(\chi(0)=1\)、\(\chi'(0)=0\)，所以
\[
 U(x',0)=a(x'),\quad
 \partial_tU(x',0)=-b(x'),\quad
 \partial_\nu U(x',0)=-\partial_tU(x',0)=b(x').
\]
所有阶导数都是紧支集的有界光滑函数，因此属于每个有限指数的 \(L^p(H)\)，也就得到题设的全部整数阶 Sobolev 正则性。

对任何在边界附近光滑的 \(U\)，把恒等式 \(U(x',0)=a(x')\) 沿边界变量求导，必有
\[
 \partial_{x_i}U(x',0)=\partial_{x_i}a(x'),\quad 1\le i<d.
\]
这证明条件的必要性。若 \(c_i=\partial_{x_i}a\)，第一问构造的 \(U\) 已经同时满足它们，故条件也充分。由此可见，光滑函数的边界值与法向导数可以独立指定，切向导数却已由边界值决定。整个构造只处理光滑数据，不据此推断一般分数阶数据所需的高阶延拓估计。
\end{solution}

% !TeX root = ../../Book1.tex
\chapter{Poincaré 不等式与 Sobolev 嵌入}
\label{b1:ch:24}
\BookChapterNumber{b1:ch:24}
\begin{chapterabstract}
本章研究弱梯度能够控制哪些函数性质。减去均值或施加零边界条件后，Poincaré 不等式控制函数振幅；当导数可积指数低于维数时，Sobolev 不等式提高函数的可积指数；超过维数时，Morrey 估计给出具有确定 Hölder 指数的连续代表。高阶结论通过迭代与局部化建立，临界情形则转向有限指数常数的增长和指数可积性。不同结论对区域几何、边界条件与指数端点各有要求。
\end{chapterabstract}

\begin{learninggoals}
\begin{enumerate}
\item\label{b1:goal:24-poincare}证明球上的均值估计、任意有界域的零边界估计及连通 Lipschitz 区域的均值估计，明确常数核和几何常数的作用。
\item\label{b1:goal:24-sobolev}从逐坐标积分与任意维乘积估计出发证明 Gagliardo--Nirenberg--Sobolev 不等式，再用幂次、稠密性和延拓处理一般指数与区域。
\item\label{b1:goal:24-morrey}通过逐尺度均值比较在每个点构造连续代表，推导局部、全空间及闭区域上的 Morrey 估计，并辨认指数的适用边界。
\item\label{b1:goal:24-higher}把一阶结论组织成高阶和分数阶嵌入，区分经典连续导数、弱导数与高阶边界延拓的额外要求。
\item\label{b1:goal:24-critical}在选读部分追踪临界 \(L^q\) 常数随 \(q\) 的增长，证明非最优常数的 Trudinger--Moser 不等式，并用具体函数辨认临界有界性的失败。
\end{enumerate}
\end{learninggoals}

本章以\chapref{b1:ch:21}的弱导数运算和代表理论为基础，反复使用\chapref{b1:ch:22}的内部逼近、零边界闭包及一阶延拓。\chapref{b1:ch:23}的分数阶空间用于高阶一节末尾，迹则帮助解释零边界条件。这里先直接证明均值 Poincaré 不等式，不把下一章的紧嵌入提前作为工具。

读到每个估计时，宜先检查右端是否包含函数的零阶范数，以及常数依赖于什么。全空间的紧支集路线能给出只含梯度的估计；一般有界区域若不减均值、不施加零边界条件，就必须保留常数函数的贡献。次临界可积性和超临界连续性虽都来自一阶导数，却不能通过简单替换指数相互推导。

中文伴读可参阅王明新的《索伯列夫空间》\cite[第2.6--2.11节]{Wang2013SobolevChinese}。Nirenberg 的《On elliptic partial differential equations》\cite[Lecture II]{Nirenberg1959Elliptic}提供 Sobolev 不等式的原始证明路线；本章展开其中任意维 Hölder 归纳和幂次代换。Morrey 部分参照 Hajłasz 的《Sobolev spaces on metric-measure spaces》\cite{Hajlasz2003SobolevMetric}中球均值的组织方式，但这里的证明始终在欧氏空间中完成。

% !TeX root = ../../Book1.tex
\section{Poincaré 型不等式}
\label{b1:sec:2401}

梯度不能区分两个相差常数的函数。因此，在一般函数空间中用梯度控制函数本身，必须先消除常数这一自由度。减去平均值是一种办法，要求零边界条件是另一种办法；它们需要的几何假设并不相同。本节先在球上建立尺度明确的估计，再说明连通性和边界条件分别在哪里起作用。

\subsection{平均值版本}
\label{b1:subsec:240101}

对有限正测度的集合 \(A\)，记
\[
 u_A=\frac1{|A|}\int_Au(x)\,\mathrm dx.
\]
本章 \(|A|\) 表示 Lebesgue 测度，\(|\nabla u|\) 表示弱梯度的欧氏长度。后者的 \(L^p\) 范数与本书逐偏导数求和的范数只相差维数常数。

\ProseParagraphBreak

\term[poincare-budengshi]{Poincaré 不等式}[Poincaré inequality]的球上版本如下。

\begin{theorem}[Poincaré]\label{b1:thm:ball-poincare-lp}
设 \(B=B(x_0,r)\subset\mathbb R^d\)、\(1\le p\le\infty\)。对 \(u\in W^{1,p}(B)\)，有
\begin{equation}\label{b1:eq:ball-poincare-lp}
 \|u-u_B\|_{L^p(B)}
 \le C_d r\|\nabla u\|_{L^p(B)}.
\end{equation}
常数可以只依赖于维数，不依赖于球的位置、半径或指数 \(p\)。
\end{theorem}
\begin{proof}
先取 \(u\in C^\infty(B)\cap W^{1,p}(B)\)。球的凸性保证两点间的线段留在球内，故
\[
 |u(x)-u(y)|
 \le\int_0^1|\nabla u(x+t(y-x))|\cdot|x-y|\,\mathrm dt.
\]
对 \(y\) 平均并用积分三角不等式，得到
\[
 |u(x)-u_B|
 \le\frac1{|B|}\int_0^1\int_B
   |\nabla u(x+t(y-x))|\cdot|x-y|\,\mathrm dy\,\mathrm dt.
\]
对固定 \(x,t\) 作代换 \(z=x+t(y-x)\)。Jacobi 因子为 \(t^d\)，而 \(|x-y|=|x-z|/t\)。若 \(z\in x+t(B-x)\)，则 \(|x-z|<2rt\)，因此
\[
 \begin{aligned}
 |u(x)-u_B|
 &\le\frac1{|B|}\int_B|\nabla u(z)|\cdot|x-z|
       \int_{|x-z|/(2r)}^1t^{-d-1}\,\mathrm dt\,\mathrm dz\\
 &\le C_d\int_B\frac{|\nabla u(z)|}{|x-z|^{d-1}}\,\mathrm dz.
 \end{aligned}
\]
所有被积函数非负，故换元后的次序交换由 Tonelli 定理合法；对角线上的值不影响积分。这里
\[
 \int_{a}^1t^{-d-1}\,\mathrm dt\le\frac{a^{-d}}d,
 \quad |B|=\omega_dr^d,
\]
所以半径在系数中相消。

在域外将 \(|\nabla u|\) 补零，右端受卷积核
\[
 K_r(h)=|h|^{1-d}\mathbf1_{\{|h|<2r\}}
\]
控制。这个核可积，并满足 \(\|K_r\|_1=C_dr\)。卷积的 \(L^1*L^p\) 估计于是给出\eqref{b1:eq:ball-poincare-lp}，其中常数不依赖 \(p\)。

若 \(p<\infty\)，由 Meyers--Serrin 定理在 \(W^{1,p}(B)\) 中逼近一般 \(u\)。均值是 \(L^p(B)\) 上的连续线性泛函，所以函数、均值和梯度都可取极限，得到一般情形。若 \(p=\infty\)，则 \(u\in W^{1,q}(B)\) 对每个有限 \(q\) 成立。应用刚证且常数与 \(q\) 无关的估计，再令 \(q\to\infty\)，利用有限测度集上 \(L^q\) 范数趋于本性上确界，得到无穷端点。
\end{proof}

这个证明同时给出一个有用的凸集版本。若 \(G\) 是非空有界凸开集，直径为 \(D\)，则同样的线段换元把系数换成 \(C_dD^d/|G|\)。截断核的 \(L^1\) 范数再贡献 \(C_dD\)，因此
\begin{equation}\label{b1:eq:convex-poincare-shape}
 \|u-u_G\|_{L^p(G)}
 \le C_d\frac{D^{d+1}}{|G|}\|\nabla u\|_{L^p(G)}.
\end{equation}
这里给出的是充分使用的粗常数，不主张它对所有凸集都最优。对长宽比固定的盒子，\(|G|\) 与 \(D^d\) 可比，便仍得到“直径乘梯度”的尺度。

\ProseParagraphBreak

平均值不是随意插入的修饰。若改成固定常数零，非零常函数会立即反驳估计；若把半径因子删去，伸缩 \(u(x)=v((x-x_0)/r)\) 又会使函数和梯度两边的尺度失配。估计既消除了梯度的常数核，也保留了空间尺度。

\subsection{零边界版本}
\label{b1:subsec:240102}

零边界条件已经排除非零常数，因此可以直接估计函数本身。这里不要求开集连通，也不要求它的边界具有图表。

\begin{theorem}[零边界 Poincaré 不等式]\label{b1:thm:zero-boundary-poincare}
设非空开集 \(\Omega\subset(a,b)\times\mathbb R^{d-1}\)，其中 \(b-a<\infty\)。则对 \(1\le p\le\infty\)、\(u\in W_0^{1,p}(\Omega)\)，有
\begin{equation}\label{b1:eq:zero-boundary-poincare}
 \|u\|_{L^p(\Omega)}
 \le(b-a)\|\partial_1u\|_{L^p(\Omega)}
 \le(b-a)\|\nabla u\|_{L^p(\Omega)}.
\end{equation}
特别地，每个有界开集都满足相应的零边界估计。
\end{theorem}
\begin{proof}
先取 \(\varphi\in C_c^\infty(\Omega)\)，并把它补零到全空间。对固定的其余坐标 \(x'\)，
\[
 \varphi(x_1,x')=\int_a^{x_1}\partial_1\varphi(t,x')\,\mathrm dt.
\]
有限 \(p\) 时，Hölder 不等式给出
\[
 |\varphi(x_1,x')|^p
 \le(b-a)^{p-1}\int_a^b|\partial_1\varphi(t,x')|^p\,\mathrm dt.
\]
先对 \(x_1\in(a,b)\) 积分，再对 \(x'\) 积分，所得即为\eqref{b1:eq:zero-boundary-poincare}的 \(p\) 次方。对于 \(p=1\)，沿线的积分三角不等式就是同一计算；对于 \(p=\infty\)，直接对沿线表示取上确界即可。

一般 \(u\in W_0^{1,p}(\Omega)\) 按定义有内部紧支集光滑逼近。将刚证估计传给这列逼近的极限，即得结论，包括按闭包定义给定的无穷端点。最后，\(|\partial_1u|\le|\nabla u|\) 给出第二个不等式。
\end{proof}

这一计算在\cref{b1:ex:22-strip-poincare}中已作为工具出现，现在将它列为后续反复使用的定理。它证明梯度半范数在 \(W_0^{1,p}\) 上是范数，并与原 Sobolev 范数等价。在有界 Lipschitz 区域和有限 \(p\) 时，可用上一章的零迹刻画把条件改写成 \(\operatorname{Tr}u=0\)；对任意开集或无穷端点，仍应保留闭包定义。

\subsection{Poincaré–Wirtinger 不等式}
\label{b1:subsec:240103}

在一般连通区域上，均值版本通常称为\term[poincare-wirtinger-budengshi]{Poincaré--Wirtinger 不等式}[Poincaré--Wirtinger inequality]。连通性阻止函数在不同分支上分别取互不相同的常数；边界图表则提供有限个可比较的局部区域。

\begin{theorem}[Poincaré--Wirtinger]\label{b1:thm:connected-domain-poincare}
设 \(\Omega\subset\mathbb R^d\) 为非空、有界、连通的 Lipschitz 区域。对 \(1\le p\le\infty\)，存在常数 \(C_{\Omega,p}\)，使
\[
 \|u-u_\Omega\|_{L^p(\Omega)}
 \le C_{\Omega,p}\|\nabla u\|_{L^p(\Omega)}
 \quad\bigl(u\in W^{1,p}(\Omega)\bigr).
\]
\end{theorem}
\begin{proof}
先构造有限局部覆盖。边界附近取\secref{b1:sec:2203}的有限内图表；它们与 \(\Omega\) 的交都是正半盒在双 Lipschitz 映射下的像。图表未覆盖的剩余部分是内部紧集，再用有限个内部球覆盖。由此得到有限个非空开集 \(G_1,\ldots,G_N\)，覆盖 \(\Omega\)。

每个 \(G_j\) 上都有
\[
 \|u-u_{G_j}\|_{L^p(G_j)}
 \le C_j\|\nabla u\|_{L^p(G_j)}.
\]
对于球，这是前面的定理。对于图表像，将函数拉回凸正半盒，使用\eqref{b1:eq:convex-poincare-shape}，再通过双 Lipschitz 换元控制函数和梯度。拉回的均值是一个常数 \(c_j\)，虽未必等于物理区域均值，但
\[
 \|u-u_{G_j}\|_{L^p(G_j)}
 \le2\|u-c_j\|_{L^p(G_j)}.
\]
这个估计来自
\(|u_{G_j}-c_j|\le |G_j|^{-1/p}\|u-c_j\|_p\)，无穷端点按 \(1/p=0\) 解释。故局部均值估计确实成立。

以 \(G_j\) 为顶点，若 \(G_i\cap G_j\ne\varnothing\) 就连一条边。这个有限图连通；否则两组顶点的并会把连通的 \(\Omega\) 分成两个不交非空开集。每个非空交集也是开集，测度严格为正。记 \(m_j=u_{G_j}\)，对相邻顶点及 \(A_{ij}=G_i\cap G_j\)，有
\[
 |m_i-m_j|
 \le |A_{ij}|^{-1/p}
       \bigl(\|u-m_i\|_{L^p(G_i)}
              +\|u-m_j\|_{L^p(G_j)}\bigr).
\]
这是把常数 \(m_i-m_j=(m_i-u)+(u-m_j)\) 限制到交集后取范数。

在有限连通图中固定从顶点1到各顶点的路径，沿路径相加得到
\[
 |m_j-m_1|\le C_{\Omega,p}\|\nabla u\|_{L^p(\Omega)}.
\]
常数包含局部图表界、交集测度及有限条路径。再在每个 \(G_j\) 上写
\[
 \|u-m_1\|_{L^p(G_j)}
 \le\|u-m_j\|_{L^p(G_j)}+|G_j|^{1/p}|m_j-m_1|.
\]
有限覆盖允许将这些局部估计合为
\(\|u-m_1\|_{L^p(\Omega)}\le C_{\Omega,p}\|\nabla u\|_p\)。
最后用
\(\|u-u_\Omega\|_p\le2\|u-m_1\|_p\)，完成证明。上述有限和及交集估计也适用于 \(p=\infty\)。
\end{proof}

这个证明没有用到下一章的紧嵌入定理。相反，它具体显示了常数如何受几何影响：图表越来越窄或相邻块的交集越来越小，所给的统一常数都可能恶化。仅知道区域有界，并不足以对一族变化的区域取得相同常数。

\ProseParagraphBreak

若 \(\Omega\) 不连通，取函数在两个分量上分别为不同常数，则弱梯度为零，减去全域均值后却不为零。因此通常只能逐连通分量减去各自均值，或另加能连接这些自由常数的约束。零边界版本没有这一障碍，因为每个分量上的非零常数都已被闭包条件排除。

\subsection{局部形式}
\label{b1:subsec:240104}

将球上不等式除以 \(|B|^{1/p}\)，得到尺度归一化的形式
\begin{equation}\label{b1:eq:poincare-average-form}
 \left(\frac1{|B|}\int_B|u-u_B|^p\,\mathrm dx\right)^{1/p}
 \le C_dr
       \left(\frac1{|B|}\int_B|\nabla u|^p\,\mathrm dx\right)^{1/p}.
\end{equation}
在有限 \(p\) 时，这个式子同时比较平均振幅与平均梯度大小。它适用于内部球，不要求全域边界规则，也不包含尚未指定的边界条件。

\ProseParagraphBreak

两个不同半径的均值也可以比较。若 \(B_r=B(x,r)\)、\(B_{2r}=B(x,2r)\subset\Omega\)，则
\[
 \begin{aligned}
 |u_{B_r}-u_{B_{2r}}|
 &\le \frac1{|B_r|}\int_{B_r}|u-u_{B_{2r}}|\\
 &\le C_d r^{1-d/p}\|\nabla u\|_{L^p(B_{2r})}.
 \end{aligned}
\]
第二步先扩大积分区域，再使用 \(L^1\) Poincaré 与 Hölder 不等式。这个半径幂是\secref{b1:sec:2403} Morrey 理论的关键：当 \(p>d\) 时，沿半径逐次减半得到一个可和的几何级数；当 \(p=d\) 时，仅凭这个界不再获得衰减。

\ProseParagraphBreak

另一种局部形式来自零边界估计。对 \(\eta\in C_c^\infty(B_R)\)、\(u\in W^{1,p}(B_R)\)，有限 \(p\) 时 \(\eta u\in W_0^{1,p}(B_R)\)，因此
\[
 \|\eta u\|_{L^p(B_R)}
 \le C_dR\bigl(\|\eta\nabla u\|_{L^p(B_R)}
                   +\|u\nabla\eta\|_{L^p(B_R)}\bigr).
\]
要使右侧的零阶项更小，常把 \(u\) 换成 \(u-u_{B_R}\)，再用均值版本控制它。这种减常数再截断的组织，与上一章 Caccioppoli 不等式中任意常数 \(c\) 的选择一致；局部化引入的截断导数仍必须保留。

% !TeX root = ../../Book1.tex
\section{Sobolev 不等式}
\label{b1:sec:2402}

一阶导数可积，能否使函数本身具有更高的可积指数？Poincaré 型估计控制函数减去常数后的大小，而本节要提高左端的积分指数。证明从紧支集函数开始：函数沿每条坐标线都在两端消失，因而可以用该方向的导数恢复函数值。把各方向的估计合在一起，维数便进入了可积指数。

\ProseParagraphBreak

Nirenberg 的《On elliptic partial differential equations》\cite[Lecture II, pp. 128--129]{Nirenberg1959Elliptic}给出逐坐标积分的证明及幂次代换。本节把其中的乘积积分步骤写成任意维归纳，再用已经建立的稠密性和延拓定理进入 Sobolev 空间。除最后的一维说明外，始终设 \(d\geq2\)；标量域为 \(\mathbb K=\mathbb R\) 或 \(\mathbb C\)，梯度长度采用
\[
 |\nabla u|=\left(\sum_{i=1}^d|\partial_i u|^2\right)^{1/2}.
\]
其 \(L^p\) 范数与一阶导数分量的有限乘积范数等价，所需比较常数只依赖于 \(d,p\)。

% 实读正式来源 Nirenberg1959Elliptic，Numdam 开放扫描：
% https://numdam.org/item/ASNSP_1959_3_13_2_115_0.pdf
% 印刷128--129页（PDF第15--16页），式(2.4)、(2.4)'及三维逐坐标证明。
% 原文一般维数写为同样使用Holder，一般p仅给幂次；本节展开归纳、幂次计算与复值C1细节。
% 一般W1p推广、W0补零和Lipschitz域推广调用本书第22章，不将原文的光滑域讨论扩写成任意域结论。

\subsection{\texorpdfstring{\(1\leq p<d\)}{1≤p<d}}
\label{b1:subsec:240201}

先处理 \(p=1\)。对 \(x=(x_1,\ldots,x_d)\)，以 \(\widehat{x}_i\) 表示删去第 \(i\) 个坐标所得的向量。沿第 \(i\) 条坐标线积分导数后，得到的函数不再依赖 \(x_i\)。下面的不等式正是为这种结构准备的。

\begin{lemma}\label{b1:lem:coordinate-product-integral}
设 \(f_i\in L^1(\mathbb R^{d-1})\) 非负，\(1\leq i\leq d\)。则
\begin{equation}\label{b1:eq:coordinate-product-integral}
 \int_{\mathbb R^d}\prod_{i=1}^d
          f_i(\widehat{x}_i)^{1/(d-1)}\,\mathrm dx
 \leq\prod_{i=1}^d
          \left(\int_{\mathbb R^{d-1}}f_i\right)^{1/(d-1)}.
\end{equation}
\end{lemma}
\begin{proof}
对维数归纳。当 \(d=2\) 时，左端为
\(\int_{\mathbb R^2}f_1(x_2)\cdot f_2(x_1)\,\mathrm dx\)，由 Tonelli 定理恰好等于两个积分的乘积。

设 \(d\geq3\)，且结论在 \(d-1\) 维成立。写 \(x'=(x_1,\ldots,x_{d-1})\)，对 \(i<d\) 置
\[
 F_i(\widehat{x'}_i)=\int_{\mathbb R}f_i(\widehat{x}_i)\,\mathrm dx_d.
\]
它是 \(d-2\) 个变量的非负可积函数，而且
\(\int F_i=\int f_i\)。对固定 \(x'\)，在 \(x_d\) 中使用 \(d-1\) 个指数都为 \(d-1\) 的 Hölder 不等式，得
\[
 \int_{\mathbb R}\prod_{i=1}^{d-1}
       f_i(\widehat{x}_i)^{1/(d-1)}\,\mathrm dx_d
 \leq\prod_{i=1}^{d-1}
       F_i(\widehat{x'}_i)^{1/(d-1)}.
\]
这里的多因子 Hölder 可由两因子形式归纳得到：分出第一个因子，用指数 \(N,N/(N-1)\)，再对后 \(N-1\) 个因子的 \(N/(N-1)\) 次幂应用归纳假设，即得 \(N\) 个指数都为 \(N\) 的形式。

记\eqref{b1:eq:coordinate-product-integral}左端为 \(I\)。由于 \(f_d\) 与 \(x_d\) 无关，上述估计给出
\[
 \begin{aligned}
 I
 &\leq\int_{\mathbb R^{d-1}}f_d(x')^{1/(d-1)}
             \prod_{i=1}^{d-1}F_i(\widehat{x'}_i)^{1/(d-1)}\,\mathrm dx'\\
 &\leq\left(\int_{\mathbb R^{d-1}}f_d\right)^{1/(d-1)}
       \left(\int_{\mathbb R^{d-1}}
             \prod_{i=1}^{d-1}F_i(\widehat{x'}_i)^{1/(d-2)}
             \,\mathrm dx'\right)^{(d-2)/(d-1)}.
 \end{aligned}
\]
第二步的两个 Hölder 指数为 \(d-1\) 和 \((d-1)/(d-2)\)。对最后的积分应用 \(d-1\) 维归纳假设，再取外面的幂，得到
\[
 I\leq\left(\int f_d\right)^{1/(d-1)}
             \prod_{i=1}^{d-1}\left(\int F_i\right)^{1/(d-1)}.
\]
这就是所需估计。全部被积函数非负，交换积分由 Tonelli 定理保证，无须先假定乘积可积。
\end{proof}

把它用于各坐标方向的导数，就得到\term[gagliardo-nirenberg-sobolev-budengshi]{Gagliardo–Nirenberg–Sobolev 不等式}[Gagliardo--Nirenberg--Sobolev inequality]的端点版本。

\begin{theorem}[Gagliardo–Nirenberg–Sobolev，\(p=1\)]\label{b1:thm:gns-endpoint}
若 \(u\in C_c^1(\mathbb R^d;\mathbb K)\)，则
\begin{equation}\label{b1:eq:gns-endpoint}
 \|u\|_{d/(d-1)}
 \leq\frac12\prod_{i=1}^d\|\partial_i u\|_1^{1/d}
 \leq\frac12\|\nabla u\|_1.
\end{equation}
\end{theorem}
\begin{proof}
固定其余坐标。由于 \(u\) 在该坐标线两端为零，微积分基本定理分别给出
\[
 u(x)=\int_{-\infty}^{x_i}\partial_i u(x_1,\ldots,t,\ldots,x_d)\,\mathrm dt
     =-\int_{x_i}^{\infty}\partial_i u(x_1,\ldots,t,\ldots,x_d)\,\mathrm dt.
\]
这对复值函数也成立，只须按实部和虚部积分。对两式分别取绝对值并相加，得
\[
 2|u(x)|\leq
 A_i(\widehat{x}_i),\quad
 A_i(\widehat{x}_i)
 =\int_{\mathbb R}|\partial_i u(x_1,\ldots,t,\ldots,x_d)|\,\mathrm dt.
\]
对 \(i=1,\ldots,d\) 相乘，再取 \(1/(d-1)\) 次幂并积分。由前一引理，
\[
 \int |u|^{d/(d-1)}
 \leq2^{-d/(d-1)}
       \prod_{i=1}^d\left(\int A_i\right)^{1/(d-1)}
 =2^{-d/(d-1)}
       \prod_{i=1}^d\|\partial_i u\|_1^{1/(d-1)}.
\]
取 \((d-1)/d\) 次幂即得第一个不等式。又因每个
\(\|\partial_i u\|_1\leq\|\nabla u\|_1\)，得到第二个不等式。
\end{proof}

这里并未寻求最佳常数。重要的是右端只含一阶导数，而左端指数已经大于一；各坐标方向都参与估计，不能只在一个方向积分后便把其余变量略去。

\subsection{临界指数}
\label{b1:subsec:240202}

\begin{definition}\label{b1:def:sobolev-critical-exponent}
设 \(1\leq p<d\)。定义相应的\term[sobolev-linjie-zhishu]{Sobolev 临界指数}[Sobolev critical exponent]为
\begin{equation}\label{b1:eq:sobolev-critical-exponent}
 p^*\coloneqq\frac{dp}{d-p},\quad
 \frac1{p^*}=\frac1p-\frac1d.
\end{equation}
\end{definition}
\symidx{\(p^*\)：Sobolev 临界指数}

这里的“临界”指给定 \(p<d\) 后，函数值可以达到的边界指数 \(p^*\)，不是把导数指数设成 \(p=d\)。后者将在\secref{b1:sec:2405}另行讨论。比如 \(d=3,p=2\) 时，\(p^*=6\)：平方可积的一阶导数将给出函数的六次可积性。

\ProseParagraphBreak

端点估计应用于函数的一定幂，就能提高导数端的指数。设 \(1<p<d\)，希望把\eqref{b1:eq:gns-endpoint}用于 \(v=|u|^\gamma\)。左端积分中的幂为 \(\gamma d/(d-1)\)，而右端使用 \(p,p'=p/(p-1)\) 的 Hölder 不等式后，会出现 \((\gamma-1)p'\)。要让这两者相同，须取
\begin{equation}\label{b1:eq:sobolev-power-choice}
 \gamma=\frac{p(d-1)}{d-p}>1,\quad
 \frac{\gamma d}{d-1}=(\gamma-1)p'=p^*.
\end{equation}
因此这个幂不是独立的技巧性参数，而是由两端积分必须相容决定的。

\begin{lemma}\label{b1:lem:gns-power-estimate}
设 \(1<p<d\)，\(u\in C_c^1(\mathbb R^d;\mathbb K)\)，并令 \(\gamma\) 如\eqref{b1:eq:sobolev-power-choice}。则
\[
 \|u\|_{p^*}\leq\frac{\gamma}{2}\|\nabla u\|_p.
\]
\end{lemma}
\begin{proof}
先说明幂次代换的正则性。把 \(\mathbb C\) 视为 \(\mathbb R^2\)，映射 \(F(z)=|z|^\gamma\) 在 \(z\ne0\) 处的导数长度为 \(\gamma|z|^{\gamma-1}\)。在原点，
\[
 \frac{|F(z)-F(0)|}{|z|}=|z|^{\gamma-1}\longrightarrow0,
\]
所以导数为零；上述导数长度也趋于零，故 \(F\) 在原点为 \(C^1\)。实值情形是同一论证的一维版本。因此 \(v=|u|^\gamma\) 属于 \(C_c^1\)，不必把它误当成无限次光滑函数。

经典链式法则及欧氏范数不等式给出
\[
 |\nabla v|\leq\gamma|u|^{\gamma-1}\cdot|\nabla u|.
\]
复值情况下，每个偏导数为
\(\partial_i v=\gamma|u|^{\gamma-2}\operatorname{Re}
(\overline u\,\partial_i u)\)（在 \(u\ne0\) 处），在 \(u=0\) 处为零；用
\(\bigl|\operatorname{Re}(\overline u\,\partial_i u)\bigr|
\leq|u|\cdot|\partial_i u|\) 后求平方和，即得同一梯度界。

将 \(v\) 代入端点估计，再用 Hölder 不等式和\eqref{b1:eq:sobolev-power-choice}，得到
\[
 \begin{aligned}
 \|u\|_{p^*}^{\gamma}
 &=\|v\|_{d/(d-1)}
 \leq\frac{\gamma}{2}
          \int |u|^{\gamma-1}\cdot|\nabla u|\\
 &\leq\frac{\gamma}{2}
       \left(\int|u|^{(\gamma-1)p'}\right)^{1/p'}
       \|\nabla u\|_p
 =\frac{\gamma}{2}\|u\|_{p^*}^{\gamma-1}
                  \cdot\|\nabla u\|_p.
 \end{aligned}
\]
紧支集和连续性保证出现的范数有限。若 \(u\ne0\)，约去
\(\|u\|_{p^*}^{\gamma-1}\)；零函数则直接满足结论。
\end{proof}

\subsection{Gagliardo–Nirenberg–Sobolev 不等式}
\label{b1:subsec:240203}

接下来把光滑函数上的估计传给弱导数。需要辨明左端极限：逼近原先只在 \(W^{1,p}\) 中给出，而提高指数后的极限由刚才的估计产生，不能预先假定二者相同。

\begin{theorem}[Gagliardo–Nirenberg–Sobolev]\label{b1:thm:sobolev-subcritical-embedding}
设 \(d\geq2\)、\(1\leq p<d\)。则每个 \(u\in W^{1,p}(\mathbb R^d;\mathbb K)\) 都属于 \(L^{p^*}(\mathbb R^d)\)，并有
\begin{equation}\label{b1:eq:sobolev-global-embedding}
 \|u\|_{p^*}\leq C_{d,p}\|\nabla u\|_p.
\end{equation}
可以取 \(C_{d,1}=1/2\)，当 \(p>1\) 时取
\(C_{d,p}=p(d-1)/[2(d-p)]\)。本节不主张这些常数最佳。
\end{theorem}
\begin{proof}
由\cref{b1:thm:whole-space-sobolev-density}取
\(u_j\in C_c^\infty(\mathbb R^d;\mathbb K)\)，使 \(u_j\to u\) 于 \(W^{1,p}\)。将已经证明的估计用于 \(u_j-u_\ell\)，得
\[
 \|u_j-u_\ell\|_{p^*}
 \leq C_{d,p}\|\nabla u_j-\nabla u_\ell\|_p\longrightarrow0.
\]
由 \(L^{p^*}\) 完备性，存在 \(v\in L^{p^*}\)，使 \(u_j\to v\) 于该空间。取一条子列，使它同时几乎处处收敛到原来的 \(L^p\) 极限 \(u\) 和新的 \(L^{p^*}\) 极限 \(v\)：例如令两种范数误差的相应次幂之和沿该子列可和，再用 Tonelli 定理。因此 \(u=v\) 几乎处处。最后在
\(\|u_j\|_{p^*}\leq C_{d,p}\|\nabla u_j\|_p\) 中取极限，得到结论。这里 \(p^*<\infty\)，所以两种有限指数的逼近都适用。
\end{proof}

\begin{corollary}\label{b1:cor:domain-sobolev-embedding}
设 \(1\leq p<d\)，则有以下两种域上的结论。
\begin{enumerate}
\item\label{b1:item:domain-sobolev-zero-boundary}
对任意非空开集 \(\Omega\subseteq\mathbb R^d\) 及
\(u\in W_0^{1,p}(\Omega)\)，有
\[
 \|u\|_{L^{p^*}(\Omega)}
 \leq C_{d,p}\|\nabla u\|_{L^p(\Omega)}.
\]
常数与 \(\Omega\) 无关。
\item\label{b1:item:domain-sobolev-extension}
若 \(\Omega\) 为有界 Lipschitz 区域，则每个
\(u\in W^{1,p}(\Omega)\) 满足
\[
 \|u\|_{L^{p^*}(\Omega)}
 \leq C_{\Omega,p}\|u\|_{W^{1,p}(\Omega)}.
\]
\end{enumerate}
若相应的 \(\Omega\) 具有有限测度，这两条结论还分别给出所有
\(1\leq q\leq p^*\) 的 \(L^q\) 估计，只须在左端变换指数时加入因子
\(m_d(\Omega)^{1/q-1/p^*}\)。
\end{corollary}
\begin{proof}
第一条使用\cref{b1:prop:w0-zero-extension}：补零函数
\(E_0u\in W^{1,p}(\mathbb R^d)\)，其梯度也由域内梯度补零得到。对它应用\eqref{b1:eq:sobolev-global-embedding}，两端范数分别等于域内范数。

第二条使用\cref{b1:thm:lipschitz-domain-extension}的有界延拓 \(E\)，依次得到
\[
 \|u\|_{L^{p^*}(\Omega)}
 \leq\|Eu\|_{L^{p^*}(\mathbb R^d)}
 \leq C_{d,p}\|\nabla(Eu)\|_{L^p(\mathbb R^d)}
 \leq C_{\Omega,p}\|u\|_{W^{1,p}(\Omega)}.
\]
最后，有限测度下直接对 \(|u|^q\cdot1\) 使用 Hölder 不等式，有
\(\|u\|_q\leq m_d(\Omega)^{1/q-1/p^*}\|u\|_{p^*}\)，即得较低指数的估计。
\end{proof}

一般开集的第一条结论依赖零边界闭包，而非边界光滑性。第二条没有零边界条件，右端因此保留零阶项：非零常数函数在有界区域内属于 \(W^{1,p}\)，梯度却为零，已经排除了对所有函数只用梯度控制的可能。

\begin{proposition}\label{b1:prop:sobolev-subcritical-interpolation}
设 \(u\in W^{1,p}(\mathbb R^d)\)，\(p\leq q\leq p^*\)，并置
\(\theta=d(1/p-1/q)\in[0,1]\)。则
\begin{equation}\label{b1:eq:sobolev-subcritical-interpolation}
 \|u\|_q
 \leq \|u\|_p^{1-\theta}\cdot\|u\|_{p^*}^{\theta}
 \leq C_{d,p}^{\theta}
       \|u\|_p^{1-\theta}\cdot\|\nabla u\|_p^{\theta}.
\end{equation}
端点 \(\theta=0,1\) 按相应单个范数理解。
\end{proposition}
\begin{proof}
只需处理 \(0<\theta<1\)。关系
\[
 \frac1q=\frac{1-\theta}{p}+\frac{\theta}{p^*}
\]
使 \(p/[(1-\theta)q]\) 与 \(p^*/(\theta q)\) 成为共轭指数。将
\(|u|^q=|u|^{(1-\theta)q}\cdot|u|^{\theta q}\) 代入 Hölder 不等式，并取 \(q\) 次根，得到第一个估计；再代入全空间的 Sobolev 不等式，得到第二个。两个端点分别是恒等式和\eqref{b1:eq:sobolev-global-embedding}。
\end{proof}

全空间的测度无限，所以不能像有界域那样，仅凭 \(L^{p^*}\) 就向所有较低指数下降。这里 \(L^p\) 项与梯度项共同控制中间指数；\(q<p\) 不在这个插值结论的范围内。

\subsection{缩放}
\label{b1:subsec:240204}

临界指数还可以从尺度直接读出。固定非零
\(u\in C_c^\infty(\mathbb R^d)\)，令 \(u_\lambda(x)=u(\lambda x)\)，其中 \(\lambda>0\)。换元给出
\[
 \|u_\lambda\|_q=\lambda^{-d/q}\|u\|_q,\quad
 \|\nabla u_\lambda\|_p=\lambda^{1-d/p}\|\nabla u\|_p,
\]
\(q=\infty\) 时按 \(1/q=0\) 理解。如果存在不随函数改变的常数，使所有紧支集光滑函数都满足
\(\|u\|_q\leq C\|\nabla u\|_p\)，那么必有
\[
 \lambda^{d/p-1-d/q}\|u\|_q\leq C\|\nabla u\|_p
 \quad(\lambda>0).
\]
指数为正时令 \(\lambda\to\infty\)，为负时令 \(\lambda\to0\)，都会矛盾。因此必须
\(1/q=1/p-1/d\)，即 \(q=p^*\)。尺度只给出必要条件；前面的坐标积分证明才保证这个指数确实能够达到。

\ProseParagraphBreak

即使定义域固定为一个小球，非齐次估计也不能越过 \(p^*\)。取
\(x_0\in\Omega\)、非零 \(\varphi\in C_c^\infty\bigl(B(0,1)\bigr)\)，并在
\(\varepsilon<\operatorname{dist}(x_0,\Omega^c)\) 时置
\[
 v_\varepsilon(x)
 =\varepsilon^{1-d/p}
       \varphi\left(\frac{x-x_0}{\varepsilon}\right).
\]
这些函数属于 \(C_c^\infty(\Omega)\)，而
\[
 \|v_\varepsilon\|_p=\varepsilon\|\varphi\|_p,\quad
 \|\nabla v_\varepsilon\|_p=\|\nabla\varphi\|_p,\quad
 \|v_\varepsilon\|_q
 =\varepsilon^{1-d/p+d/q}\|\varphi\|_q.
\]
当 \(q>p^*\) 时，最后的指数为负，左端趋于无穷，前两个量却有界。这排除了任何统一的
\(W^{1,p}(\Omega)\rightarrow L^q(\Omega)\) 范数估计；它只检验可积性估计的边界，不需要引入后文的紧性理论。

\ProseParagraphBreak

一维情形须单独读。条件 \(1\leq p<d\) 在 \(d=1\) 时为空，但沿整条实轴的基本定理仍给出另一个端点：
\[
 \|u\|_\infty\leq\frac12\|u'\|_1
 \quad\bigl(u\in W^{1,1}(\mathbb R)\bigr).
\]
对 \(C_c^1\) 函数，它就是从左右两端积分后相加的估计。对一般
\(W^{1,1}\) 函数，取紧支集光滑逼近；将估计用于两项之差，得到一致 Cauchy 列，其连续极限与 \(L^1\) 极限几乎处处相同，故同一界传到极限。这里的结论是本性上确界控制；\secref{b1:sec:2403}讨论 \(p>d\) 时更精确的连续性估计。

% !TeX root = ../../Book1.tex
\section{Morrey 理论}
\label{b1:sec:2403}

当梯度的可积指数超过空间维数时，局部平均振荡不仅趋零，而且在逐次缩小半径后可以求和。这使 Sobolev 函数获得连续代表，并给出两点之间的定量控制。以下先讨论 \(d\ge1\)、\(d<p<\infty\)，令
\[
 \alpha=1-\frac dp\in(0,1).
\]

为把本节放回整章的位置，\cref{b1:tab:first-order-sobolev-regimes}列出 \(d\ge2\) 时一阶理论的三个指数区间。它只给出典型目标；区域、边界条件与端点的精确假设仍以各正式定理为准。
\begin{table}[htbp]
  \centering
  \caption[一阶 Sobolev 的三个区间]{一阶 Sobolev 的三个区间。梯度尺度从积分增益，经临界端点，过渡到 Hölder 连续性。}
  \label{b1:tab:first-order-sobolev-regimes}
  \begin{tabularx}{\textwidth}{@{}p{2.7cm}p{3.6cm}X@{}}
    \toprule
    指数范围 & 临界尺度 & 典型结论 \\
    \midrule
    \(1\le p<d\) & \(p^*=dp/(d-p)\) & 嵌入到 \(L^{p^*}\)，紧性只在严格低于临界指数时出现 \\
    \(p=d\) & \(p^*=\infty\) 的形式端点 & 一般无 \(L^\infty\) 控制；有界域上有全部有限 \(L^q\) 控制，零边界时还有指数可积性 \\
    \(p>d\) & \(\alpha=1-d/p>0\) & 存在 \(C^{0,\alpha}\) 代表，并有 Morrey 型双点估计 \\
    \bottomrule
  \end{tabularx}
\end{table}

以下以实值函数为主，复值或有限维向量值函数可按分量处理。记 \(\|\nabla u\|_{L^p}\) 为弱梯度的欧氏长度的 \(L^p\) 范数；它与本书按偏导分量计算的梯度范数只相差依赖维数及 \(p\) 的等价常数。

Hajłasz 的《Sobolev spaces on metric-measure spaces》\cite[Section 9]{Hajlasz2003SobolevMetric}用球均值的望远镜求和把平均振荡变成双点估计。下面保留这一组织，但利用欧氏体积与 \(p>d\)，直接让球均值在每个点收敛，不先把讨论限制在 Lebesgue 点上。
% 实读 Hajlasz2003SobolevMetric：作者正式公开PDF，§9 Theorem 9.4及完整证明，
% 印刷204--205页；另读Theorem 9.7(3)的Hölder结论陈述，印刷206页。
% 这里改编逐半径球均值比较；一般度量空间的先修与9.7完整证明不作为本节依赖。
% 欧氏球Poincare由本章24.1提供；每一点的极限、任意半径一致性及代表辨析另行展开。

\subsection{\texorpdfstring{\(p>d\)}{p>d}}
\label{b1:subsec:240301}

对球 \(B=B(x,r)\)，记 \(u_B=m_d(B)^{-1}\int_Bu\)。球上的 Poincaré 不等式与 Hölder 不等式给出
\begin{equation}\label{b1:eq:morrey-average-oscillation}
 \begin{split}
 \frac1{m_d(B)}\int_B|u-u_B|
 &\le m_d(B)^{-1/p}\|u-u_B\|_{L^p(B)}\\
 &\le C_{d,p}r^{1-d/p}\|\nabla u\|_{L^p(B)}.
 \end{split}
\end{equation}
这里使用的是\cref{b1:thm:ball-poincare-lp}，无需先知道 \(u\) 连续。幂次 \(1-d/p\) 是导数贡献的一个长度因子，与平均化引入的 \(r^{-d/p}\) 相乘所得。当它为正时，各个尺度的误差形成收敛的几何级数。

\begin{lemma}\label{b1:lem:morrey-ball-average-limit}
若 \(u\in W^{1,p}_{\mathrm{loc}}(\Omega)\)，其中 \(\Omega\subset\mathbb R^d\) 为开集、\(d<p<\infty\)，则对每个 \(x\in\Omega\)，有限极限
\[
 u^*(x)=\lim_{r\downarrow0}u_{B(x,r)}
\]
都存在，且 \(u^*=u\) 几乎处处。对每个满足 \(\overline{B(x,R)}\Subset\Omega\) 的球，还有
\begin{equation}\label{b1:eq:morrey-average-error}
 |u^*(x)-u_{B(x,R)}|
 \le C_{d,p}R^\alpha\|\nabla u\|_{L^p(B(x,R))}.
\end{equation}
\end{lemma}
\begin{proof}
固定 \(x,R\)，置 \(r_j=2^{-j}R\) 和 \(B_j=B(x,r_j)\)。因为两相邻球的体积比为 \(2^d\)，
\[
 \begin{split}
 |u_{B_{j+1}}-u_{B_j}|
 &\le\frac1{m_d(B_{j+1})}\int_{B_j}|u-u_{B_j}|\\
 &\le C_{d,p}r_j^\alpha\|\nabla u\|_{L^p(B_j)}
 \le C_{d,p}2^{-j\alpha}R^\alpha
                         \|\nabla u\|_{L^p(B_0)}.
 \end{split}
\]
右端可求和，所以均值序列为 Cauchy 列，存在有限极限；从 \(j=0\) 起求和就得到\eqref{b1:eq:morrey-average-error}。

还要排除极限对选定半径列的依赖。若 \(r_{j+1}<r\le r_j\)，则 \(B(x,r)\subset B_j\)，且体积比仍不超过 \(2^d\)。同一估计给出
\[
 |u_{B(x,r)}-u_{B_j}|
 \le C_{d,p}r_j^\alpha\|\nabla u\|_{L^p(B_0)}
 \longrightarrow0.
\]
因此所有趋零半径的均值都趋于同一极限，它与最初的 \(R\) 无关。以上论证对每个固定点成立，没有删除例外点。最后由\cref{b1:thm:lebesgue-point}，在几乎每个点，该极限等于原函数值，故确实得到 \(u\) 的一个代表。
\end{proof}

当 \(p=d\) 时，上面用整个固定球上的梯度范数估计的几何级数失去衰减。这个观察只说明当前证明不能给出正的 Hölder 指数。\secref{b1:sec:2405}将讨论临界情形中的其他控制。

\subsection{全空间 Morrey 不等式}
\label{b1:subsec:240302}

\begin{definition}\label{b1:def:holder-space}
给定非空集合 \(A\subset\mathbb R^d\) 及 \(0<\gamma\le1\)，若函数 \(v:A\rightarrow\mathbb R\) 满足
\[
 [v]_{0,\gamma;A}
 \coloneqq\sup_{\substack{x,y\in A\\x\ne y}}
       \frac{|v(x)-v(y)|}{|x-y|^\gamma}<\infty,
\]
则称它在 \(A\) 上为 \(\gamma\) 阶\term[holder-lianxu]{Hölder 连续}[Hölder continuous]。若还满足 \(\sup_A|v|<\infty\)，记 \(v\in C^{0,\gamma}(A)\)，并规定
\[
 \|v\|_{C^{0,\gamma}(A)}
 \coloneqq\sup_A|v|+[v]_{0,\gamma;A}.
\]
集合只有一个点时，半范数中的上确界约定为零。
\end{definition}
\symidx{\(C^{0,\gamma}(A)\)：有界Hölder连续函数空间}
\symidx{\([v]_{0,\gamma;A}\)：Hölder半范数}

指数为 \(1\) 时就是 Lipschitz 条件。对紧集，有限半范数和任意一个有限点值已保证有界；在全空间上则不能省去零阶条件，例如 \(v(x)=x_1\) 虽然 Lipschitz，却不属于这里的 \(C^{0,1}(\mathbb R^d)\)。

\begin{theorem}[Morrey]\label{b1:thm:morrey-whole-space}
设 \(d<p<\infty\) 且 \(u\in W^{1,p}(\mathbb R^d)\)。球平均极限 \(u^*\) 是唯一的连续代表，属于 \(C^{0,\alpha}(\mathbb R^d)\)，其中 \(\alpha=1-d/p\)，并满足
\begin{equation}\label{b1:eq:morrey-whole-space}
 [u^*]_{0,\alpha;\mathbb R^d}
 \le C_{d,p}\|\nabla u\|_{L^p(\mathbb R^d)},
 \quad
 \|u^*\|_{C^{0,\alpha}(\mathbb R^d)}
 \le C_{d,p}\|u\|_{W^{1,p}(\mathbb R^d)}.
\end{equation}
\end{theorem}
\begin{proof}
给定 \(x\ne y\)，令 \(r=|x-y|\)，并取
\[
 B_x=B(x,r),\quad B_y=B(y,r),\quad B_*=B(x,3r).
\]
两个小球都包含于 \(B_*\)，体积比只依赖于 \(d\)。由\eqref{b1:eq:morrey-average-oscillation}，
\[
 |u_{B_x}-u_{B_*}|+|u_{B_y}-u_{B_*}|
 \le C_{d,p}r^\alpha\|\nabla u\|_{L^p(B_*)}.
\]
再在 \(x,y\) 两点使用\eqref{b1:eq:morrey-average-error}，得到
\begin{equation}\label{b1:eq:morrey-local-pair}
 |u^*(x)-u^*(y)|
 \le C_{d,p}|x-y|^\alpha\cdot
                 \|\nabla u\|_{L^p(B(x,3|x-y|))}.
\end{equation}
放大积分域到整个空间，即得 Hölder 半范数界，也证明 \(u^*\) 处处连续。

半范数不能单独控制常数，因此还须估计一个平均值。对任意 \(x\)，用半径为 \(1\) 的球，
\[
 \begin{split}
 |u^*(x)|
 &\le|u^*(x)-u_{B(x,1)}|+|u_{B(x,1)}|\\
 &\le C_{d,p}\|\nabla u\|_{L^p(\mathbb R^d)}
       +\omega_d^{-1/p}\|u\|_{L^p(\mathbb R^d)}.
 \end{split}
\]
于是得到全范数估计。若另一个连续函数与 \(u\) 几乎处处相等，它与 \(u^*\) 的差连续且几乎处处为零；一旦在某点非零，就会在一个正测度小球内非零，矛盾。因此连续代表唯一。
\end{proof}

上述\term[morrey-budengshi]{Morrey 不等式}[Morrey inequality]把弱梯度的积分控制转成点值的变化控制。证明中的每一个平均值只依赖等价类，连续代表不是通过任意指定某个零测集上的值获得的。

\subsection{局部与区域版本}
\label{b1:subsec:240303}

全空间估计也给出内部估计。这里必须留意参与均值比较的球是否仍在原开集内，而不能直接对靠近边界的点使用越出开集的球。

\begin{corollary}\label{b1:cor:morrey-domain}
设 \(d<p<\infty\)、\(\alpha=1-d/p\)。
若 \(u\in W^{1,p}_{\mathrm{loc}}(\Omega)\)，则其球平均代表局部 Hölder 连续。具体地，若 \(\overline{B(a,8R)}\Subset\Omega\)，则
\begin{equation}\label{b1:eq:morrey-interior-ball}
 \begin{split}
 [u^*]_{0,\alpha;\overline{B(a,R)}}
 &\le C_{d,p}\|\nabla u\|_{L^p(B(a,8R))},\\
 \sup_{\overline{B(a,R)}}|u^*|
 &\le C_{d,p}\left(
        R^{-d/p}\|u\|_{L^p(B(a,8R))}
        +R^\alpha\|\nabla u\|_{L^p(B(a,8R))}\right).
 \end{split}
\end{equation}
若 \(\overline U\Subset V\subset\Omega\)，其中 \(U,V\) 为开集，且 \(u\in W^{1,p}(V)\)，则
\[
 \|u^*\|_{C^{0,\alpha}(\overline U)}
 \le C_{U,V,d,p}\|u\|_{W^{1,p}(V)}.
\]
若 \(\Omega\) 为有界 Lipschitz 区域，则 \(u\in W^{1,p}(\Omega)\) 有唯一的连续延伸到 \(\overline\Omega\) 的代表，且
\begin{equation}\label{b1:eq:morrey-lipschitz-domain}
 \|u^*\|_{C^{0,\alpha}(\overline\Omega)}
 \le C_{\Omega,p}\|u\|_{W^{1,p}(\Omega)}.
\end{equation}
\end{corollary}
\begin{proof}
对 \(x,y\in\overline{B(a,R)}\)，有 \(|x-y|\le2R\)，故
\(B(x,3|x-y|)\subset B(a,8R)\)。主定理的双点比较只在这些内部球中进行，所以\eqref{b1:eq:morrey-local-pair}仍成立，给出第一个估计。对 \(x\) 使用 \(B(x,R)\) 的平均值，并以 Hölder 不等式控制该平均值，就得到第二个估计。这同时证明代表的局部 Hölder 连续性。

给定 \(\overline U\Subset V\)，取固定 \(\chi\in C_c^\infty(V)\)，使它在 \(\overline U\) 的一个邻域内等于 \(1\)。将 \(\chi u\) 补零到整个空间。由弱乘积公式及\cref{b1:lem:compact-sobolev-zero-extension}，所得 \(w\) 满足
\[
 \|w\|_{W^{1,p}(\mathbb R^d)}
 \le C_{\chi,p}\|u\|_{W^{1,p}(V)}.
\]
在 \(\overline U\) 附近的充分小球上，\(u\) 与 \(w\) 的平均值相同，故两者的精确代表也相同。应用全空间估计即得第二项结论。

最后，若 \(\Omega\) 有界且为 Lipschitz 区域，用\cref{b1:thm:lipschitz-domain-extension}取得全空间延拓 \(Eu\)，再把其 Hölder 连续代表限制到 \(\overline\Omega\)。在内部，它与原来的球平均代表一致；边界值由内部连续趋近唯一确定。因此结果不依赖所选延拓算子，且延拓范数界给出\eqref{b1:eq:morrey-lipschitz-domain}。
\end{proof}

区域不必连通，但区域版本使用的是完整 \(W^{1,p}\) 范数，不能一律只保留 \(\|\nabla u\|_p\)。例如在两个分离的球上分别取常数 \(0\) 与 \(1\)，弱梯度为零，跨两个分量的 Hölder 半范数却非零。全空间延拓同时处理各分量之间的变化，相关常数包含区域的几何信息。

当 \(p=\infty\) 时，结论中的指数改为 \(1\)，但证明不通过把 \(p\) 代入有限指数公式来完成。\cref{b1:thm:w1infty-lipschitz}已经给出局部 Lipschitz 代表；在全空间中，把连接两点的线段分为有限个落在局部球内的小段并求和，就得到由全域弱梯度本性界控制的 Lipschitz 常数。其上确界由 \(\|u\|_\infty\) 控制。再使用截断或有界 Lipschitz 区域的延拓算子，分别得到局部和闭区域上的 \(C^{0,1}\) 估计。

指数 \(\alpha=1-d/p\) 在统一范数估计中不能提高。固定非恒定函数 \(\phi\in C_c^\infty(B(0,1))\)，并令
\[
 u_\varepsilon(x)=\varepsilon^\alpha\phi(x/\varepsilon),
 \quad 0<\varepsilon<1.
\]
换元直接给出
\[
 \|u_\varepsilon\|_p=\varepsilon\|\phi\|_p,\quad
 \|\nabla u_\varepsilon\|_p=\|\nabla\phi\|_p.
\]
这些函数在 \(W^{1,p}\) 中一致有界。然而，选取 \(\phi(a)\ne\phi(b)\)，对每个 \(\alpha<\beta\le1\)，
\[
 [u_\varepsilon]_{0,\beta;\mathbb R^d}
 \ge\varepsilon^{\alpha-\beta}
       \frac{|\phi(a)-\phi(b)|}{|a-b|^\beta}
 \longrightarrow\infty.
\]
所以不存在把所有 \(W^{1,p}\) 函数的 \(C^{0,\beta}\) 范数控制住的统一常数；同一例子也适用于包含这些小支集的固定球。每个 \(u_\varepsilon\) 本身仍然光滑，这里排除的是更强的连续嵌入估计，而不是说单个函数不能具有更好的正则性。

\subsection{精确代表}
\label{b1:subsec:240304}

\secref{b1:sec:1404}曾由球平均定义精确代表，并在极限不存在处另作赋值。现在，超临界可积性使这一例外情形在内部完全消失。

\begin{proposition}\label{b1:prop:morrey-precise-representative}
若 \(u\in W^{1,p}_{\mathrm{loc}}(\Omega)\)、\(d<p<\infty\)，则其精确代表 \(u^*\) 是唯一的连续代表，并且对每个 \(x\in\Omega\)，
\[
 \lim_{r\downarrow0}\frac1{m_d(B(x,r))}
       \int_{B(x,r)}|u(y)-u^*(x)|\,\mathrm dy=0.
\]
特别地，\(u^*\) 的每个内部点都是 Lebesgue 点。若 \(u_j\to u\) 于 \(W^{1,p}_{\mathrm{loc}}(\Omega)\)，则对每个相对紧开集 \(U\)，精确代表在 \(C^{0,\alpha}(\overline U)\) 中收敛；这里取 \(\overline U\Subset\Omega\)，且 \(\alpha=1-d/p\)。
\end{proposition}
\begin{proof}
有限极限已由\cref{b1:lem:morrey-ball-average-limit}在每个点构造，局部 Hölder 连续性由前一推论给出。连续代表的唯一性仍由非空开球具有正测度得到。对充分小的 \(r\)，
\[
 \begin{split}
 \frac1{m_d(B(x,r))}\int_{B(x,r)}|u-u^*(x)|
 &\le\frac1{m_d(B(x,r))}\int_{B(x,r)}|u-u_{B(x,r)}|
       +|u_{B(x,r)}-u^*(x)|\\
 &\le C_{d,p}r^\alpha
                 \|\nabla u\|_{L^p(B(x,r))}
 \longrightarrow0.
 \end{split}
\]
将积分中的 \(u\) 换成与它几乎处处相等的 \(u^*\)，即得每点的 Lebesgue 点性质。

最后取 \(\overline U\Subset V\) 且 \(\overline V\Subset\Omega\)。球均值的线性及每点极限说明
\((u_j-u)^*=u_j^*-u^*\)。对差函数使用前一推论，便有
\[
 \|u_j^*-u^*\|_{C^{0,\alpha}(\overline U)}
 \le C_{U,V,d,p}\|u_j-u\|_{W^{1,p}(V)}
 \longrightarrow0.
\]
\end{proof}

连续性属于这个被确定的代表，而不是等价类中的每个可测函数。即使只在一点改变 \(u^*\) 的值，也不改变任何 \(L^p\) 范数或弱导数，却可能破坏该点的连续性。另一方面，在有界 Lipschitz 区域上，\cref{b1:thm:lipschitz-lp-trace}说明这个闭域连续代表的边界限制恰是前章的迹；积分方式定义的边界值与通常点值在这里相遇。

本节得到的是函数值的 Hölder 控制，并不使弱梯度自动连续。把同一方法应用到哪些低阶导数、怎样组织高阶连续代表，是\secref{b1:sec:2404}高阶嵌入要处理的问题。

% !TeX root = ../../Book1.tex
\section{高阶 Sobolev 嵌入}
\label{b1:sec:2404}

一阶嵌入可以作用于函数，也可以作用于它的低阶弱导数。逐次使用这一事实，就能把高阶可积性转化为更高的积分指数，或转化为经典导数的 Hölder 连续性。整个过程有两个需要区分的步骤：先提高每个弱导数的可积性，再证明所得连续代表确实互为经典导数。

\ProseParagraphBreak

本节首先取整数 \(k\ge1\) 和有限指数 \(1\le p<\infty\)，在全空间中讨论。一般开集上的内部结论可通过截断取得；涉及整个边界时，所需延拓必须控制到 \(k\) 阶。最后单独从双积分定义证明 \(0<s<1\) 的分数阶版本。
% 整数阶迭代由本章Gagliardo--Nirenberg--Sobolev与Morrey定理自足推出。
% 分数阶均值比较沿用已实读Hajlasz2003SobolevMetric §9 Theorem 9.4
% 的望远镜组织（印刷204--205页）；原文9.7(3)仅作为已读的范围说明，
% 不把该文未展开于本书的度量空间理论当作高阶或分数阶证明前置。

\subsection{高阶导数}
\label{b1:subsec:240401}

若 \(u\in W^{k,p}(\mathbb R^d)\)，则每个 \(|\beta|\le k-1\) 的弱导数 \(\partial^\beta u\) 都属于 \(W^{1,p}\)。在 \(p<d\) 时，对这些函数分别使用一阶嵌入，便得到
\[
 W^{k,p}(\mathbb R^d)\subset W^{k-1,p_1}(\mathbb R^d),
 \quad \frac1{p_1}=\frac1p-\frac1d.
\]
右端包含所有不超过 \(k-1\) 阶的弱导数，不只是最高阶的某几个分量。因此下一次迭代仍有完整的 Sobolev 空间可用。

\begin{theorem}\label{b1:thm:higher-order-subcritical-embedding}
设 \(1\le p<\infty\)、\(k\ge1\) 且 \(kp<d\)，并令
\[
 q=\frac{dp}{d-kp}.
\]
则有连续嵌入
\[
 W^{k,p}(\mathbb R^d)\subset L^q(\mathbb R^d),
 \quad
 \|u\|_q\le C_{d,k,p}\|u\|_{W^{k,p}(\mathbb R^d)}.
\]
\end{theorem}
\begin{proof}
规定 \(p_0=p\) 及
\[
 \frac1{p_j}=\frac1p-\frac jd,\quad 0\le j\le k.
\]
条件 \(kp<d\) 保证这些指数有限，而且每一步开始时 \(p_j<d\)，其中 \(0\le j<k\)。若已知 \(u\in W^{k-j,p_j}\)，则对每个 \(|\beta|\le k-j-1\)，由\cref{b1:thm:sobolev-subcritical-embedding}，
\[
 \|\partial^\beta u\|_{p_{j+1}}
 \le C_{d,p_j}\sum_{i=1}^d
                  \|\partial^{\beta+e_i}u\|_{p_j}.
\]
这些弱导数仍是原来的分布导数；改变可积指数不改变其定义。对有限多个 \(\beta\) 求和，得到
\[
 \|u\|_{W^{k-j-1,p_{j+1}}}
 \le C_{d,k,p_j}\|u\|_{W^{k-j,p_j}}.
\]
从 \(j=0\) 迭代到 \(j=k-1\)，便到达 \(W^{0,p_k}=L^q\)，并得到范数界。
\end{proof}

每提高一次积分指数，消耗一阶导数，同时使 \(1/p\) 减少 \(1/d\)。例如在五维空间中，\(W^{2,2}\) 先嵌入 \(W^{1,10/3}\)，再嵌入 \(L^{10}\)。若中途指数达到 \(d\)，就不能继续把分母等于零的表达式当作下一步的一阶公式；若指数超过 \(d\)，则应转入 Morrey 理论。

\subsection{Hölder 空间}
\label{b1:subsec:240402}

\begin{definition}\label{b1:def:higher-holder-space}
给定开集 \(U\subset\mathbb R^d\)、整数 \(m\ge0\) 及 \(0<\alpha\le1\)，若 \(v\in C^m(U)\)，所有不超过 \(m\) 阶的经典偏导数均有界，且每个 \(m\) 阶偏导数都为 \(\alpha\) 阶 Hölder 连续，则记 \(v\in C^{m,\alpha}(U)\)，并规定
\[
 \|v\|_{C^{m,\alpha}(U)}
 \coloneqq\sum_{|\beta|\le m}\sup_U|\partial^\beta v|
   +\sum_{|\beta|=m}[\partial^\beta v]_{0,\alpha;U}.
\]
记号 \(C^{m,\alpha}(\overline U)\) 还要求这些导数连续延伸到 \(\overline U\)，上式中的上确界与 Hölder 半范数也在闭包上计算。当 \(m=0\) 时与前节的定义一致。
\end{definition}
\symidx{\(C^{m,\alpha}(U)\)：最高\(m\)阶经典导数为Hölder连续的空间}

\begin{theorem}\label{b1:thm:higher-order-holder-embedding}
设 \(k\ge1\)、\(1\le p<\infty\)，且 \(k-d/p>0\) 不是整数。写成
\[
 k-\frac dp=m+\alpha,\quad m\in\mathbb N_0,\quad 0<\alpha<1.
\]
则每个 \(u\in W^{k,p}(\mathbb R^d)\) 都有唯一的 \(C^m\) 代表 \(v\)，满足
\[
 v\in C^{m,\alpha}(\mathbb R^d),\quad
 \|v\|_{C^{m,\alpha}(\mathbb R^d)}
 \le C_{d,k,p}\|u\|_{W^{k,p}(\mathbb R^d)}.
\]
\end{theorem}
\begin{proof}
令 \(\ell=k-m-1\ge0\)，并规定
\[
 \frac1q=\frac1p-\frac\ell d=\frac{1-\alpha}{d}.
\]
故 \(q>d\)。若 \(\ell=0\)，已知 \(u\in W^{m+1,q}\)。若 \(\ell>0\)，前一证明中的一阶迭代可以进行 \(\ell\) 次：最后一次开始前的指数仍小于 \(d\)，因为对应的 \(d/p_{\ell-1}=2-\alpha>1\)。因此两种情形均得到
\[
 \|u\|_{W^{m+1,q}(\mathbb R^d)}
 \le C_{d,k,p}\|u\|_{W^{k,p}(\mathbb R^d)}.
\]
对每个 \(|\beta|\le m\)，函数 \(\partial^\beta u\) 属于 \(W^{1,q}\)。由\cref{b1:thm:morrey-whole-space}，它有 \(C^{0,\alpha}\) 代表 \(v_\beta\)，且相应范数受上式控制。

现在核实这些代表确实是同一经典函数的导数。取局部矩形及其稍大的内部邻域，令 \(u_\varepsilon=\rho_\varepsilon*u\)。由\cref{b1:prop:sobolev-friedrichs-smoothing}，
\[
 \partial^\beta u_\varepsilon
 =\rho_\varepsilon*\partial^\beta u
 =\rho_\varepsilon*v_\beta,\quad |\beta|\le m.
\]
因为 \(v_\beta\) 连续，右端在每个紧集上一致趋于 \(v_\beta\)。对 \(|\beta|<m\)，在矩形内的坐标线段上有
\[
 \partial^\beta u_\varepsilon(x+te_i)-\partial^\beta u_\varepsilon(x)
 =\int_0^t\partial^{\beta+e_i}u_\varepsilon(x+re_i)\,\mathrm dr.
\]
让 \(\varepsilon\) 趋零，得到相同恒等式，其两侧分别换成 \(v_\beta\) 与 \(v_{\beta+e_i}\)。由于被积函数连续，\(v_\beta\) 的经典偏导数就是 \(v_{\beta+e_i}\)。连续偏导保证经典可微；逐阶使用便得到
\(v_0\in C^m\) 且 \(\partial^\beta v_0=v_\beta\)。当 \(m=0\) 时，只需已经得到的连续代表。

对所有低阶上确界及最高阶 Hölder 半范数求和，得到所需估计。连续代表唯一，因而 \(C^m\) 代表也唯一。
\end{proof}

正整数差不能不加论证地解释为端点 Lipschitz 正则性。下面给出足以覆盖任意较小 Hölder 指数的结论。

\begin{proposition}\label{b1:prop:integer-gap-holder-embedding}
设 \(1<p<\infty\)，且 \(j=k-d/p\) 为正整数。对每个 \(0<\alpha<1\)，均有连续嵌入
\[
 W^{k,p}(\mathbb R^d)\subset C^{j-1,\alpha}(\mathbb R^d),
\]
其中右端取经典代表。
\end{proposition}
\begin{proof}
选取充分接近 \(p\) 的 \(p_0\in(1,p)\)，使
\[
 k-\frac d{p_0}=j-1+\alpha_0,\quad \alpha<\alpha_0<1.
\]
固定 \(\chi\in C_c^\infty\bigl(B(0,2)\bigr)\)，使它在 \(B(0,1)\) 上等于 \(1\)，并令
\(\chi_a(x)=\chi(x-a)\)。由于支集体积固定，Hölder 不等式及高阶弱 Leibniz 公式给出与中心 \(a\) 无关的界
\[
 \|\chi_a u\|_{W^{k,p_0}(\mathbb R^d)}
 \le C_{d,k,p,p_0,\chi}\|u\|_{W^{k,p}(\mathbb R^d)}.
\]
应用前一定理，\(\chi_a u\) 在 \(B(a,1)\) 上的经典代表具有一致的 \(C^{j-1,\alpha_0}\) 界。不同球上的代表在交集几乎处处相等，因连续而处处相等，其经典导数也相同，所以这些局部代表拼成同一个 \(C^{j-1}\) 函数。

所有不超过 \(j-1\) 阶的导数由局部界得到全域上确界。对最高阶导数的两点差，若 \(|x-y|<1\)，取中心 \(a=x\)，局部估计与
\(|x-y|^{\alpha_0}\le|x-y|^\alpha\) 给出所需界；若 \(|x-y|\ge1\)，则用两倍上确界控制差值。两种距离范围合起来得到全域的 \(C^{j-1,\alpha}\) 范数界。
\end{proof}

这个证明中的降指数只在固定大小的支集上进行，没有使用全空间 \(L^p\subset L^{p_0}\) 这一错误包含关系。它也没有证明 \(\alpha=1\) 的结论。若 \(p=1\)，就不能再降到允许范围内的 \(p_0<p\)；该临界端点需要专门论证，本节不从上述方法推断其最强结论。差 \(k-d/p=0\) 同样不在这个正整数差命题中。无穷指数若另作讨论，应逐阶使用 Lipschitz 代表的判据，而不是代入有限指数迭代。

\ProseParagraphBreak

一般开集上的内部版本不需要边界正则性：若
\(\overline U\Subset V\)、\(u\in W^{k,p}(V)\)，选取在 \(\overline U\) 附近等于 \(1\) 的 \(\chi\in C_c^\infty(V)\)，以高阶弱乘积公式估计 \(\chi\cdot u\)，再用\cref{b1:lem:compact-sobolev-zero-extension}补零即可应用上述全空间结论。由此取得内部的 \(L^q\) 或 \(C^{m,\alpha}(\overline U)\) 估计。

\ProseParagraphBreak

若要把这些估计延至整个 \(\overline\Omega\)，可以假设存在有界线性延拓
\[
 E_k:W^{k,p}(\Omega)\longrightarrow W^{k,p}(\mathbb R^d),
 \quad (E_ku)|_\Omega=u,
\]
再应用全空间定理。\secref{b1:sec:2203}只为 Lipschitz 区域证明了一阶延拓；在这里讨论 \(k>1\) 时，不能把那个一阶算子当作已经具备高阶有界性。

\subsection{整数阶与非整数阶}
\label{b1:subsec:240403}

整数阶迭代使用了弱导数链。对于 \(0<s<1\) 的空间，正则性由差分双积分度量，必须重新从该定义取得平均振荡估计。两种推导最后会汇合到相同的球均值方法。

\begin{theorem}[分数阶 Morrey]\label{b1:thm:fractional-morrey}
设 \(0<s<1\)、\(1\le p<\infty\) 且 \(sp>d\)，令
\(\alpha=s-d/p\in(0,1)\)。每个 \(u\in W^{s,p}(\mathbb R^d)\) 都有唯一的连续代表 \(u^*\)，其值在每个点由球平均极限给出，且
\[
 [u^*]_{0,\alpha;\mathbb R^d}\le C_{d,s,p}[u]_{s,p},
 \quad
 \|u^*\|_{C^{0,\alpha}(\mathbb R^d)}
 \le C_{d,s,p}\|u\|_{W^{s,p}(\mathbb R^d)}.
\]
\end{theorem}
\begin{proof}
先取 \(B=B(x,r)\)。由平均值的 Jensen 不等式，
\[
 \begin{split}
 \frac1{m_d(B)}\int_B|u-u_B|^p
 &\le\frac1{m_d(B)^2}\int_B\int_B|u(z)-u(y)|^p\,\mathrm dy\,\mathrm dz\\
 &\le\frac{(2r)^{d+sp}}{m_d(B)^2}
       \int_B\int_B
          \frac{|u(z)-u(y)|^p}{|z-y|^{d+sp}}\,\mathrm dy\,\mathrm dz\\
 &\le C_{d,s,p}r^{sp-d}[u]_{s,p;B}^p.
 \end{split}
\]
再用 Hölder 不等式，得到
\[
 \frac1{m_d(B)}\int_B|u-u_B|
 \le C_{d,s,p}r^\alpha[u]_{s,p;B}.
\]
固定任意 \(x,R\)，取 \(B_i=B(x,2^{-i}R)\)。与相邻球的体积比较给出
\[
 |u_{B_{i+1}}-u_{B_i}|
 \le C_{d,s,p}2^{-i\alpha}R^\alpha[u]_{s,p;B(x,R)}.
\]
因为 \(\alpha>0\)，级数收敛。对夹在两个相邻二分半径之间的任意半径，仍可用体积比及刚才的平均振荡界比较均值。因此球平均在每个点存在有限极限，而且
\[
 |u^*(x)-u_{B(x,R)}|
 \le C_{d,s,p}R^\alpha[u]_{s,p;B(x,R)}.
\]
这正是\cref{b1:lem:morrey-ball-average-limit}证明中的求和与任意半径夹逼步骤；这里只将弱梯度范数换成了由双积分直接得到的控制量。Lebesgue 点定理保证该极限几乎处处等于 \(u\)。

给定 \(x\ne y\)，令 \(r=|x-y|\)，比较 \(B(x,r)\)、\(B(y,r)\) 与 \(B(x,3r)\) 的均值。两个小球与大球的体积比有界，所以它们的均值差由大球的平均振荡控制。再加上两个点各自的均值误差，得到
\[
 |u^*(x)-u^*(y)|
 \le C_{d,s,p}|x-y|^\alpha\cdot[u]_{s,p;B(x,3|x-y|)}
 \le C_{d,s,p}|x-y|^\alpha\cdot[u]_{s,p}.
\]
于是 \(u^*\) 连续，并取得所需 Hölder 半范数界。取 \(R=1\)，以 \(\|u\|_p\) 控制单位球均值，便得到全范数界。连续代表的唯一性与前节相同。
\end{proof}

上述论证仍在每个点构造极限，而不只是在几乎处处的点获得一个差分估计。它也说明分数阶嵌入的尺度来自 \(sp-d>0\)：双积分能量控制平均振荡，再由几何级数恢复点值。这里没有通过把整数阶公式中的 \(k\) 直接替换为 \(s\) 来定义分数阶导数。

\ProseParagraphBreak

Hajłasz 的《Sobolev spaces on metric-measure spaces》\cite[Theorem 9.4]{Hajlasz2003SobolevMetric}提供了这种均值比较的系统组织。整数阶和分数阶条件如何产生平均振荡界，各有自己的证明；将这些条件放在同一尺度观点下理解，也有助于辨清\secref{b1:sec:2405}临界情形中哪些步骤不再成立。

高阶嵌入应同时核对指数余量与区域的延拓性质，见\cref{b1:tab:visual-t13}。
% Source fingerprint: b498fd314c29f9c0aadfeeb4aca942a02d5a32b6174e8cf3848275fbd1ed6b88
% T13: raw TeX cells preserved; no numeric highlighting.
\begin{table}[htbp]
\centering\small\renewcommand{\arraystretch}{1.18}\setlength{\tabcolsep}{4pt}
\caption[高阶 Sobolev 嵌入速查]{高阶 Sobolev 嵌入速查。本表补充前面的一阶指数表。函数空间的范数与代表约定均沿用本节，不把内部正则性写成跨边界结论。}
\label{b1:tab:visual-t13}
\begin{tabularx}{\linewidth}{@{}>{\raggedright\arraybackslash}p{3.1cm}>{\raggedright\arraybackslash}p{4.7cm}>{\raggedright\arraybackslash}X@{}}
\toprule
参数区间 & 本节得到的目标 & 适用说明 \\
\midrule
\(kp<d\)，\(1\le p<\infty\) & \(L^q\)，\(q=dp/(d-kp)\) & 逐次一阶嵌入；分母必须严格为正 \\
\(kp=d\) & 不能把上行公式代入为无穷指数 & 需按相应临界理论处理，不由形式代入宣称有界连续 \\
\(k-d/p=m+\alpha\)，\(0<\alpha<1\) & 具有 \(C^{m,\alpha}\) 代表 & 有限 p；余量非整数，经典导数的相容性由正文证明 \\
\(j=k-d/p\in\mathbb N_{>0}\) & \(C^{j-1,\alpha}\)，任意 \(0<\alpha<1\) & 本节该命题要求 \(1<p<\infty\)；不擅取 \(\alpha=1\) \\
一般区域 & 内部结论通过局部化取得；全域结论需相应高阶延拓 & 一阶延拓不能直接替代 \(k\) 阶延拓假设 \\
\bottomrule
\end{tabularx}
\end{table}


% !TeX root = ../../Book1.tex
\OptionalSection{一阶临界嵌入：\texorpdfstring{\(W^{1,d}\)}{W1,d}}{b1:sec:2405}

本节只处理一阶临界情形 \(W^{1,d}\)，不讨论一般 \(W^{k,p}\) 在 \(kp=d\) 时的端点。此时 Sobolev 次临界指数的分母为零，而 Morrey 指数也恰好降为零。这并不意味着所有估计都消失了；它意味着不能继续用前面的幂次公式描述端点。以下设 \(d\ge2\)，先证明零边界函数在每个有限 \(L^q\) 中都有控制，再追踪常数随 \(q\) 的增长，把这些估计合成为指数可积性。

\ProseParagraphBreak

Moser 的《A Sharp Form of an Inequality by N. Trudinger》\cite[1. Statements, p. 1077]{Moser1971Trudinger}说明了临界指数可积性的背景，并进一步研究最佳常数。本节采用较直接的卷积与指数级数路线，只证明足以展示临界现象的非最优常数版本；不把下述常数认作原文的尖锐常数。
% 实读范围：Moser1971Trudinger官方首页PDF，印刷1077页第1节。
% 完整Theorem1跨至1078页，本次未读全定理或尖锐常数证明。
% 下文势表示、有限指数卷积估计与级数求和均在本书展开；
% 只引用实读首页作历史定位与进一步阅读，不声称复现Moser最佳常数。

\subsection{\texorpdfstring{\(p=d\)}{p=d}}
\label{b1:subsec:240501}

为了在卷积估计中让输出指数增大，使用\cref{b1:lem:young-convolution-finite}的有限指数 Young 卷积不等式。该基础工具已在卷积理论建立处证明，不把它的唯一正式出处留在本选读节中。

\begin{proposition}\label{b1:prop:critical-finite-q-estimate}
设 \(\Omega\subset\mathbb R^d\) 为非空有界开集，\(D=\operatorname{diam}\Omega\)，且 \(d\ge2\)。对 \(d\le q<\infty\)、\(u\in W_0^{1,d}(\Omega)\)，有
\begin{equation}\label{b1:eq:critical-q-growth}
 \|u\|_{L^q(\Omega)}
 \le C_dD^{d/q}q^{1-1/d}\|\nabla u\|_{L^d(\Omega)}.
\end{equation}
其中 \(C_d\) 不依赖 \(q,u,\Omega\)。
\end{proposition}
\begin{proof}
先取 \(u\in C_c^\infty(\Omega)\)，在全空间补零。固定 \(x\) 及单位向量 \(\omega\)，沿射线积分得
\[
 u(x)=-\int_0^\infty\nabla u(x+r\omega)\cdot\omega\,\mathrm dr.
\]
对单位球面平均，再使用极坐标换元，得到
\[
 |u(x)|\le\frac1{d\omega_d}
        \int_{\mathbb R^d}\frac{|\nabla u(y)|}{|x-y|^{d-1}}\,\mathrm dy.
\]
若 \(x\in\Omega\)，只有 \(y\in\Omega\) 对积分有贡献，故 \(|x-y|\le D\)。记
\[
 K_D(h)=|h|^{1-d}\mathbf1_{\{|h|\le D\}},\qquad
 b_q=\frac{dq}{(d-1)q+d}.
\]
则 \(1\le b_q<d/(d-1)\)，并且
\[
 1+\frac1q=\frac1d+\frac1{b_q}.
\]
Young 卷积不等式给出
\[
 \|u\|_{L^q(\Omega)}
 \le C_d\|K_D\|_{L^{b_q}}\cdot\|\nabla u\|_{L^d(\Omega)}.
\]
核的范数可以直接计算：
\[
 \begin{aligned}
 \|K_D\|_{L^{b_q}}
 &=\left(\frac{d\omega_d}
                    {d-(d-1)b_q}\right)^{1/b_q}
        D^{\,d/b_q-(d-1)}\\
 &\le C_d q^{1/b_q}D^{d/q}
 \le C_d q^{1-1/d}D^{d/q}.
 \end{aligned}
\]
这里 \(d/b_q-(d-1)=d/q\)，
\(d-(d-1)b_q=db_q/q\)，且
\(q^{1/b_q}=q^{1-1/d}q^{1/q}\)，最后一因子在 \(q\ge1\) 上有界。这解释了常数中 \(q^{1-1/d}\) 的来源。

对一般 \(u\in W_0^{1,d}(\Omega)\)，取内部紧支集光滑逼近。将已证估计用于逼近列两项之差，可知它在每个固定有限 \(L^q\) 中为 Cauchy 列；其极限与原来的 \(L^d\) 极限一致，所以就是 \(u\)。取极限完成证明。
\end{proof}

在有界 Lipschitz 区域上，一般 \(u\in W^{1,d}(\Omega)\) 也属于每个有限 \(L^q(\Omega)\)。事实上，\cref{b1:thm:lipschitz-domain-extension}给出支集留在固定紧集中的全空间延拓；把这个紧集放进一个更大的球，延拓函数属于该球的 \(W_0^{1,d}\)，便可应用上面的估计。所得界包含 \(\|u\|_{W^{1,d}(\Omega)}\)，不能只用梯度，因为非零常数也在此空间中。

\ProseParagraphBreak

不能在\eqref{b1:eq:critical-q-growth}中令 \(q=\infty\)：常数随 \(q\) 增大而发散。接下来的求和恰好利用这一增长速度，而不是把它忽略。

\subsection{Trudinger–Moser 不等式}
\label{b1:subsec:240502}

记 \(d'=d/(d-1)\)。\term[trudinger-moser-budengshi]{Trudinger--Moser 不等式}[Trudinger--Moser inequality]的一个非最优常数版本如下。

\begin{theorem}[Trudinger--Moser，非最优常数版]\label{b1:thm:trudinger-moser}
设 \(d\ge2\)，\(\Omega\subset\mathbb R^d\) 为非空有界开集，\(D=\operatorname{diam}\Omega\)。存在只依赖于 \(d\) 的正数 \(c_d,C_d\)，使每个 \(u\in W_0^{1,d}(\Omega)\) 满足
\begin{equation}\label{b1:eq:trudinger-moser}
 \int_\Omega
  \exp\left(c_d
       \left|\frac{u(x)}{\|\nabla u\|_{L^d(\Omega)}}\right|^{d'}\right)
       \mathrm dx
 \le C_dD^d
\end{equation}
当 \(\|\nabla u\|_d>0\)。若梯度范数为零，则 \(u=0\)，将被积函数解释为 \(1\)。这里不声称 \(c_d\) 是最大的允许常数。
\end{theorem}
\begin{proof}
由零边界 Poincaré 不等式，梯度为零时 \(u=0\)，所列约定成立。其余情形可通过数乘归一化为 \(\|\nabla u\|_d=1\)。展开非负指数级数，单调收敛定理给出
\[
 \int_\Omega e^{c|u|^{d'}}
 =|\Omega|+
   \sum_{j=1}^\infty\frac{c^j}{j!}
          \|u\|_{L^{d'j}(\Omega)}^{d'j}.
\]
对 \(d'j\ge d\)，应用\eqref{b1:eq:critical-q-growth}，
\[
 \|u\|_{L^{d'j}}^{d'j}
 \le D^d C_d^{d'j}(d'j)^j.
\]
由
\[
 \log(j!)=\sum_{\ell=1}^j\log\ell
 \ge\int_1^j\log t\,\mathrm dt
 =j\log j-j+1
\]
得到 \(j!\ge(j/e)^j\)。因此上述级数的相应项不超过
\[
 D^d\bigl(c\,C_d^{d'}d'e\bigr)^j.
\]
选取只依赖于 \(d\) 的充分小 \(c=c_d>0\)，使括号不超过 \(1/2\)，这些高次项之和便受常数倍的 \(D^d\) 控制。

剩下满足 \(d'j<d\) 的项只有有限个。由 Hölder 不等式及零边界 Poincaré 估计，
\[
 \|u\|_{L^{d'j}}^{d'j}
 \le|\Omega|^{\,1-d'j/d}\|u\|_{L^d}^{d'j}
 \le C_dD^d.
\]
这里把 \(\Omega\) 放入直径不超过常数倍 \(D\) 的球，使用
\(|\Omega|\le\omega_dD^d\) 与 \(\|u\|_d\le C_dD\)。
有限低次项、常数项和高次几何级数合起来，给出所需估计。最后恢复原函数的梯度范数，完成证明。
\end{proof}

这个定理是对梯度范数有界函数族的统一估计，不只是说每个固定函数在某个很小参数下可积。指数幂 \(d'\) 使高次 \(L^q\) 常数产生的 \(j^j\) 与阶乘的增长相抵；若把指数幂再增大，后者便未必能够控制前者。习题中的集中函数列将直接说明这种更强统一估计为何失败。

\ProseParagraphBreak

一般 \(W^{1,d}\) 函数若没有零边界条件，不能在分母中只放梯度；常数函数已足以显示问题。可在有界连通 Lipschitz 区域上先减去均值，再使用 Poincaré 和固定支集延拓取得相应的指数估计，但其常数会随区域几何改变。本文所证明的\eqref{b1:eq:trudinger-moser}保留了一个便于检查的版本：任意有界开集、零边界闭包、非最优常数和明确的直径尺度。

\subsection{Morrey–Campanato 观点}
\label{b1:subsec:240503}

Morrey 证明中的决定性步骤不是逐点微分，而是球平均的振幅随半径衰减。把导数暂时拿掉，只保留这种平均条件，就得到 Campanato 观点的基本形式。

\begin{proposition}[Campanato 型平均条件]\label{b1:prop:campanato-holder}
设 \(u\in L^1_{\mathrm{loc}}(\mathbb R^d)\)、\(0<\alpha\le1\)。若存在 \(A<\infty\)，使每个球 \(B(x,r)\) 都满足
\begin{equation}\label{b1:eq:campanato-mean-condition}
 \frac1{|B(x,r)|}\int_{B(x,r)}
       |u-u_{B(x,r)}|\,\mathrm dy\le Ar^\alpha,
\end{equation}
则 \(u\) 有唯一连续代表 \(u^*\)，并有
\[
 |u^*(x)-u^*(y)|\le C_{d,\alpha}A|x-y|^\alpha.
\]
反过来，任意满足全域 \(\alpha\)-Hölder 估计的函数都满足相应的平均条件。
\end{proposition}
\begin{proof}
固定 \(x,R\)。两相邻同心球的体积比为 \(2^d\)，故
\[
 |u_{B(x,2^{-j-1}R)}-u_{B(x,2^{-j}R)}|
 \le2^dA(2^{-j}R)^\alpha.
\]
右侧为可和几何级数，所以球均值在每个点有有限极限，并有
\[
 |u^*(x)-u_{B(x,R)}|
 \le C_{d,\alpha}AR^\alpha.
\]
任意半径可夹在两个相邻二分半径之间，使用相同的体积比估计，说明极限不依赖这列半径。Lebesgue 微分定理保证 \(u^*=u\) 几乎处处。

对 \(x\ne y\)，令 \(R=|x-y|\)。两个球 \(B(x,R)\)、\(B(y,R)\) 都包含在 \(B(x,3R)\) 内，且体积比受维数控制。因此各自均值与大球均值之差均不超过 \(C_dAR^\alpha\)。再加上两个点到各自均值的误差，得到所求双点估计。它同时证明连续性；两个连续代表若几乎处处相等，就处处相等，因此代表唯一。

反过来，若 \(|u(z)-u(w)|\le L|z-w|^\alpha\)，则
\[
 \begin{aligned}
 \frac1{|B|}\int_B|u-u_B|
 &\le\frac1{|B|^2}\int_B\int_B|u(z)-u(w)|\,\mathrm dz\,\mathrm dw\\
 &\le L(2r)^\alpha,
 \end{aligned}
\]
这就是平均条件。
\end{proof}

更常见的表述把左侧换成 \(q\) 次平均振幅，\(1\le q<\infty\)。它由 Hölder 不等式控制这里的一次平均，而 Hölder 连续函数也控制每个这样的 \(q\) 次平均，因此在正衰减指数范围内仍刻画相同的连续性。一般 Campanato 空间允许更多参数和多项式近似；这里仅需平均常数的基本情形。

\ProseParagraphBreak

在临界指数 \(p=d\) 时，Poincaré 不等式只给出
\[
 \frac1{|B|}\int_B|u-u_B|
 \le C_d\left(\int_B|\nabla u|^d\right)^{1/d}.
\]
右侧没有固定的正半径幂。对固定的可积梯度，它会随球的体积缩小而趋零，但这些误差未必沿二分尺度可求和。因而“平均振幅趋零”与“可以把球均值求和为连续代表”仍有距离；仅把前一命题的 \(\alpha\) 设为零，会破坏其关键级数。

\subsection{临界现象}
\label{b1:subsec:240504}

下面构造一个 \(W_0^{1,d}\) 中本性无界的函数，说明临界指数的确不能一般给出连续或有界代表。取
\[
 0<\beta<1-\frac1d,
\]
并选光滑径向截断 \(\chi\)，使它在原点附近等于 \(1\)、支集包含于 \(B(0,1/2)\)。对 \(0<|x|<1/2\)，令
\[
 u(x)=\chi(x)\bigl(\log(e/|x|)\bigr)^\beta,
\]
在支集外取零，在原点任意赋一个有限值。显然它在每个原点邻域中本性无界。

\ProseParagraphBreak

函数本身属于 \(L^d(B(0,1))\)，因为径向体积因子 \(r^{d-1}\) 控制任意固定的对数幂。在 \(\chi=1\) 的小球内，梯度长度为
\[
 |\nabla u(x)|
 =\frac{\beta}{|x|}
       \bigl(\log(e/|x|)\bigr)^{\beta-1},
\]
故其 \(d\) 次积分等价于
\[
 \int_0^\varepsilon
      \frac{\mathrm dr}
           {r\bigl(\log(e/r)\bigr)^{d(1-\beta)}}.
\]
代换 \(t=\log(e/r)\) 后，这是尾部的
\(\int^\infty t^{-d(1-\beta)}\,\mathrm dt\)，由于 \(d(1-\beta)>1\) 而收敛。截断的变化只发生在避开原点的紧环带上，不增加奇性。

\ProseParagraphBreak

还须确认这些经典导数确实是弱导数。对测试函数在挖去半径 \(\varepsilon\) 小球的区域上分部积分，内球面边界项的绝对值受
\[
 C\varepsilon^{d-1}
       \bigl(\log(e/\varepsilon)\bigr)^\beta
       \|\varphi\|_\infty
\]
控制，因 \(d\ge2\) 而趋零。函数和候选导数局部可积，故令 \(\varepsilon\downarrow0\) 得到全域弱导数恒等式。于是 \(u\in W^{1,d}\)，又因为它具有内部紧支集，有限指数的内部逼近给出 \(u\in W_0^{1,d}(B(0,1))\)。

\ProseParagraphBreak

它没有有界代表，但仍服从本节证明的指数估计。这并不矛盾：指数可积性限制的是大值集合的体积，而非禁止函数取得任意大的值。另一方面，当 \(d=1\) 时，\(W^{1,1}\) 的绝对连续代表已经由积分基本定理得到有界控制，本节的临界反例和指数幂 \(d/(d-1)\) 都不适用。

\ProseParagraphBreak

本章至此建立了梯度对均值振幅、可积性与连续性的不同控制。下一章将询问一个新的问题：若一列函数都服从这些控制，能否从中抽出强收敛子列？单个函数的嵌入估计不会自动给出序列紧性；还须排除质量逃向无穷远、平移失控和临界尺度的集中。


\begin{chaptersummary}
梯度首先只能控制函数相对于常数的变化。均值与零边界条件以不同方式消除了这一自由度：前者在一般区域上需要连通性及合适的几何控制，后者可以在任意有界开集上通过紧支集逼近建立。

指数低于维数时，逐坐标积分与幂次代换提高了函数值的可积性；指数高于维数时，可和的球平均误差产生了连续代表。两条路线都保留了尺度，并各有不能直接纳入的端点。高阶与分数阶嵌入需要把这些工具实际作用于相应导数或差分，不能仅对定理中的阶数作形式替换。

临界指数没有一般的 \(L^\infty\) 控制，却仍能给出指数可积性。常数的增长、平均振幅的衰减和集中反例共同说明了这种边界。下一章将把注意力从单个函数转到函数列，研究这些有界控制在什么条件下能够产生强收敛子列。
\end{chaptersummary}

\BookExercises

% \chapref{b1:ch:24}习题临时稿；不另设 BookExercises 入口。
% 八题均为自拟的正文工具变式；各解答展开新增计算或构造，不冒认教材原题编号。
% 依赖止于本章，不调用\chapref{b1:ch:25}的紧嵌入或平移紧性判据。

\begin{exercise}\label{b1:ex:24-componentwise-poincare}
设 \(1\le p<\infty\)，\(\Omega_1,\ldots,\Omega_N\subset\mathbb R^d\) 为非空、有界、连通的 Lipschitz 开集，且它们的闭包两两不交。令
\(\Omega=\bigcup_{j=1}^N\Omega_j\)，并对 \(u\in W^{1,p}(\Omega)\) 定义
\[
 Pu=\sum_{j=1}^Nu_{\Omega_j}\mathbf1_{\Omega_j},
 \quad u_{\Omega_j}=\frac1{|\Omega_j|}\int_{\Omega_j}u.
\]
\begin{enumerate}
 \item\label{b1:item:24-componentwise-estimate}
 证明 \(P\) 是 \(L^p(\Omega)\) 上的压缩投影，且
 \[
  \|u-Pu\|_{L^p(\Omega)}
       \le C\|\nabla u\|_{L^p(\Omega)}.
 \]
 若各分量均值 Poincaré 不等式的最优常数为 \(C_j\)，证明此处最优常数为
 \(\max_j C_j\)。
 \item\label{b1:item:24-componentwise-kernel}
 刻画 \(W^{1,p}(\Omega)\) 中满足 \(\nabla u=0\) 的全部函数，并证明当 \(N\ge2\) 时，不能把上式中的 \(Pu\) 换成全域平均值 \(u_\Omega\)。
\end{enumerate}
\end{exercise}
% 来源：自拟；有限分量的直和估计及梯度核，补充连通域均值 Poincaré 的假设边界。
\begin{solution}
由 Jensen 不等式，
\[
 |\Omega_j|\cdot|u_{\Omega_j}|^p
       \le\int_{\Omega_j}|u|^p.
\]
求和得到 \(\|Pu\|_p\le\|u\|_p\)。取过一次分量平均后，每个分量上已是常函数，所以 \(P^2=P\)；线性也由积分的线性得到。

分量上的常函数在 \(\Omega\) 内没有弱导数。事实上，测试函数的支集与
\(\partial\Omega\) 保持正距离，在每个分量内积分其偏导都为零；这里没有跨越域外边界的跳跃项。因此 \(Pu\in W^{1,p}(\Omega)\)，且
\(\nabla(u-Pu)=\nabla u\)。在每个分量应用定理%
\ref{b1:thm:connected-domain-poincare}，再把 \(p\) 次方相加，得到
\[
 \begin{aligned}
 \|u-Pu\|_p^p
 &=\sum_{j=1}^N\|u-u_{\Omega_j}\|_{L^p(\Omega_j)}^p\\
 &\le\left(\max_j C_j\right)^p
                       \sum_{j=1}^N\|\nabla u\|_{L^p(\Omega_j)}^p.
 \end{aligned}
\]
最优常数的下界可以逐个分量测试：在 \(\Omega_j\) 上任选一个函数，在其他分量上置零，整体比值恰好成为该分量的比值。取上确界并遍历 \(j\)，就得到所求的精确常数，不要求最优值由某个函数取得。

若 \(\nabla u=0\)，刚才的估计给出 \(u=Pu\) 几乎处处，故 \(u\) 在每个分量上分别为常数；反向已在前面说明。若 \(N\ge2\)，取
\[
 u=|\Omega_2|\mathbf1_{\Omega_1}
           -|\Omega_1|\mathbf1_{\Omega_2},
\]
并在其余分量上置零。此函数的全域平均为零、梯度为零，却不是零函数。因此全域平均不能消除梯度核中的全部自由度；需要消除的是每个分量上的一个常数。
\end{solution}

\begin{exercise}\label{b1:ex:24-thin-neck-constant}
设 \(d\ge2\)、\(1\le p<\infty\)、\(0<\varepsilon<1/2\)。考虑两个盒子由细颈连接而成的开集
\[
 \begin{aligned}
 \Omega_\varepsilon={}&
   \bigl((-2,-1/2)\times(-1,1)^{d-1}\bigr)\\
 &{}\cup\bigl((-1,1)\times(-\varepsilon,\varepsilon)^{d-1}\bigr)\\
 &{}\cup\bigl((1/2,2)\times(-1,1)^{d-1}\bigr).
 \end{aligned}
\]
令 \(C_\varepsilon\) 为该连通 Lipschitz 开集上均值 Poincaré 不等式的最优常数。证明
\[
 C_\varepsilon\ge
       \frac{3^{1/p}}2\,\varepsilon^{-(d-1)/p}.
\]
据此说明：对全部有界连通 Lipschitz 开集，常数不能只由维数、指数以及直径和体积的统一上下界控制。
\end{exercise}
% 来源：自拟；固定体积盒子与退化细颈的显式测试，不借用特征值或谱理论。
\begin{solution}
定义分段线性函数
\[
 h(s)=
 \begin{cases}
  -1,&s\le-1/2,\\
  2s,&-1/2<s<1/2,\\
  1,&s\ge1/2,
 \end{cases}
 \quad u_\varepsilon(x)=h(x_1).
\]
这个函数在全空间 Lipschitz，因而其在 \(\Omega_\varepsilon\) 上的限制属于
\(W^{1,p}\)。区域关于 \(x_1\mapsto-x_1\) 对称，而函数关于此变换变号，所以
\((u_\varepsilon)_{\Omega_\varepsilon}=0\)。

两个大盒子上的函数模都等于 \(1\)。它们各自的体积为
\((3/2)2^{d-1}\)，故
\[
 \|u_\varepsilon\|_{L^p(\Omega_\varepsilon)}^p
              \ge3\cdot2^{d-1}.
\]
唯一非零的偏导是中央细颈上的 \(\partial_1u_\varepsilon=2\)。这一部分的长度为 \(1\)，横截面体积为 \((2\varepsilon)^{d-1}\)，因而
\[
 \|\nabla u_\varepsilon\|_{L^p(\Omega_\varepsilon)}^p
              =2^p(2\varepsilon)^{d-1}.
\]
将这个函数代入最优常数的定义，即得
\[
 C_\varepsilon^p\ge
     \frac{3\cdot2^{d-1}}{2^p(2\varepsilon)^{d-1}}
     =\frac3{2^p}\varepsilon^{-(d-1)}.
\]

这些区域均包含固定的两个大盒子，又都包含在
\((-2,2)\times(-1,1)^{d-1}\) 中，故体积和直径均有与 \(\varepsilon\) 无关的正下界及有限上界。对每个固定 \(\varepsilon\)，盒子的平面边界与有限个转角给出 Lipschitz 边界；开集连通，因为细颈与两个盒子分别重叠。随 \(\varepsilon\) 缩小，退化的是连接几何，而不是这些粗略尺寸。上面的下界趋于无穷，证明了最后的断言。
\end{solution}

\begin{exercise}\label{b1:ex:24-linear-normalization}
设 \(\Omega\subset\mathbb R^d\) 为非空、有界、连通的 Lipschitz 开集，
\(1\le p<\infty\)。本题取实值函数，并令
\(\ell:W^{1,p}(\Omega)\to\mathbb R\) 为连续线性泛函。证明下列条件等价。
\begin{equivalences}
 \item\label{b1:item:24-normalization-detects-constants}
 \(\ell(1)\ne0\)。
 \item\label{b1:item:24-normalization-poincare}
 存在有限常数 \(C_\ell\)，使所有满足 \(\ell(u)=0\) 的 \(u\) 都有
 \(\|u\|_{L^p(\Omega)}\le C_\ell\|\nabla u\|_{L^p(\Omega)}\)。
\end{equivalences}
进一步，若 \(A\subset\Omega\) 可测且 \(|A|>0\)，证明条件 \(\int_Au=0\) 蕴含
\[
 \|u\|_{L^p(\Omega)}
 \le C_P\Bigl(1+(|\Omega|/|A|)^{1/p}\Bigr)
          \|\nabla u\|_{L^p(\Omega)},
\]
其中 \(C_P\) 是通常均值 Poincaré 不等式的任一个可用常数。
\end{exercise}
% 来源：自拟；用一般连续线性约束代替全域均值，再给出可测子集上的明确常数。
\begin{solution}
若 \(\ell(1)=0\)，常函数 \(u=1\) 就属于 \(\ker\ell\)，但其梯度为零、函数范数非零，所以第二个条件不可能成立。

现在设 \(\ell(1)\ne0\)，并取 \(\ell(u)=0\)。写成
\[
 u=v+c,\quad c=u_\Omega,\quad v=u-u_\Omega.
\]
由\cref{b1:thm:connected-domain-poincare}及弱导数范数的有限维等价性，有
\[
 \|v\|_{L^p(\Omega)}\le C_P\|\nabla u\|_p,
 \quad \|v\|_{W^{1,p}(\Omega)}\le C_0\|\nabla u\|_p,
\]
其中 \(C_0\) 只依赖于 \(\Omega,d,p\)。因为
\(0=\ell(v)+c\ell(1)\)，连续性给出
\[
 |c|\le\frac{\|\ell\|}{|\ell(1)|}
                   \|v\|_{W^{1,p}(\Omega)}
       \le\frac{C_0\|\ell\|}{|\ell(1)|}\|\nabla u\|_p.
\]
于是
\[
 \|u\|_p\le\|v\|_p+|\Omega|^{1/p}\cdot|c|
 \le\left(C_P+
       \frac{|\Omega|^{1/p}C_0\|\ell\|}{|\ell(1)|}\right)
                 \|\nabla u\|_p,
\]
这就证明等价性。

对于最后的条件，仍用同一分解；从 \(\int_Au=0\) 得到
\[
 |c|=\left|\frac1{|A|}\int_Av\right|
       \le |A|^{-1/p}\|v\|_{L^p(A)}.
\]
这里 \(p=1\) 时直接使用积分三角不等式，\(p>1\) 时使用 Hölder 不等式。因此
\[
 \|u\|_p\le
     \Bigl(1+(|\Omega|/|A|)^{1/p}\Bigr)
                   \|v\|_p,
\]
再代入通常的 Poincaré 估计，便得到所列常数。这个约束不要求 \(A\) 是开集，也不要求它包含一个球。
\end{solution}

\begin{exercise}\label{b1:ex:24-supercritical-algebra}
设 \(\Omega\subset\mathbb R^d\) 为非空有界 Lipschitz 开集，\(d<p<\infty\)，本题取实值函数。
\begin{enumerate}
 \item\label{b1:item:24-sobolev-algebra}
 证明乘法满足
 \[
  \|uv\|_{W^{1,p}(\Omega)}
       \le C\|u\|_{W^{1,p}(\Omega)}
                 \cdot\|v\|_{W^{1,p}(\Omega)}.
 \]
 由此证明，适当放大原范数后，新范数仍完备且满足次乘性
 \(\|uv\|_{\mathrm A}\le\|u\|_{\mathrm A}\cdot\|v\|_{\mathrm A}\)，常函数 \(1\) 为乘法单位元。
 \item\label{b1:item:24-sobolev-inversion}
 若 \(u\in W^{1,p}(\Omega)\) 且 \(|u|\ge\delta>0\) 几乎处处，证明
 \(1/u\in W^{1,p}(\Omega)\)，并给出它的函数项与梯度项估计。进一步证明：在这类远离零的函数中，若 \(u_n\to u\) 于 \(W^{1,p}\)，且所有 \(|u_n|\ge\delta\)，则
 \(1/u_n\to1/u\) 于 \(W^{1,p}\)。
\end{enumerate}
\end{exercise}
% 来源：自拟；Morrey 嵌入与已建立的 Leibniz、链式法则组合，强调代数性所需的上确界控制。
\begin{solution}
由\cref{b1:cor:morrey-domain}，
\[
 \|w\|_\infty\le C_M\|w\|_{W^{1,p}(\Omega)}.
\]
所以 \(u,v\) 都本性有界，可以使用\cref{b1:prop:sobolev-product}的乘积公式。
函数项和梯度项分别满足
\[
 \|uv\|_p\le\|u\|_\infty\cdot\|v\|_p,
 \quad
 \|\nabla(uv)\|_p
 \le\|u\|_\infty\cdot\|\nabla v\|_p
       +\|v\|_\infty\cdot\|\nabla u\|_p.
\]
代入上确界估计并合并各项，就得到第一问。取一个不小于 \(1\) 的可用常数 \(C\)，将范数改为
\(\|w\|_{\mathrm A}=C\|w\|_{W^{1,p}}\)，则
\(\|uv\|_{\mathrm A}\le\|u\|_{\mathrm A}\cdot\|v\|_{\mathrm A}\)。
等价范数保持完备性，而有界域上的常函数 \(1\) 属于这个空间，并且是乘法单位元。

对第二问，取光滑函数 \(\zeta\)，使它在 \([-\delta/2,\delta/2]\) 上为零，在
\((-\infty,-\delta]\cup[\delta,\infty)\) 上为一，令
\(F(t)=\zeta(t)/t\) 并规定 \(F(0)=0\)。这样 \(F\) 在全实轴光滑，导数有界，且
\(F(u)=1/u\)。由\cref{b1:thm:sobolev-chain}，
\[
 \nabla(1/u)=-u^{-2}\nabla u
       \quad\text{几乎处处},
\]
且
\[
 \|1/u\|_p\le|\Omega|^{1/p}\delta^{-1},
 \quad
 \|\nabla(1/u)\|_p\le\delta^{-2}\|\nabla u\|_p.
\]
这里有界域使函数项可积，不仅仅是导数项可积。

最后，Morrey 估计使 \(u_n\to u\) 也成为连续代表的一致收敛。倒数之差满足
\[
 \|1/u_n-1/u\|_p\le\delta^{-2}\|u_n-u\|_p.
\]
在集合 \(\{\,t:|t|\ge\delta\,\}\) 上，映射 \(t\mapsto t^{-2}\) 满足
\[
 |a^{-2}-b^{-2}|\le2\delta^{-3}|a-b|.
\]
当 \(a,b\) 同号时这来自导数界；异号时可先比较 \(|a|,|b|\)，并使用
\(\bigl||a|-|b|\bigr|\le|a-b|\)。因此
\[
 \begin{aligned}
 \|\nabla(1/u_n)-\nabla(1/u)\|_p
 &\le\delta^{-2}\|\nabla u_n-\nabla u\|_p\\
 &\quad+2\delta^{-3}\|u_n-u\|_\infty
                         \cdot\|\nabla u\|_p
 \longrightarrow0.
 \end{aligned}
\]
两项合起来便给出所求 Sobolev 范数收敛。
\end{solution}

\begin{exercise}\label{b1:ex:24-morrey-logarithmic-cusp}
设 \(d\ge2\)、\(d<p<\infty\)，并记 \(\alpha=1-d/p\)。
对于实数 \(\gamma\)，在单位球 \(B\) 上定义
\[
 u(0)=0,\quad
 u(x)=|x|^\alpha\bigl(\log(e/|x|)\bigr)^{-\gamma}
                  \quad(0<|x|<1).
\]
证明
\[
 u\in W^{1,p}(B)\quad\Longleftrightarrow\quad\gamma p>1.
\]
在 \(\gamma p>1\) 时，证明 \(u\) 的连续代表属于 \(C^{0,\alpha}(\overline B)\)，但不属于任何
\(C^{0,\beta}(\overline B)\)，其中 \(\alpha<\beta\le1\)。
这里要检验一个固定函数的正则性，不能只用函数族的缩放排除统一估计。
\end{exercise}
% 来源：自拟；临界幂次乘对数因子的可积性计算，与正文的统一 Hölder 估计尖锐性区分。
\begin{solution}
令 \(L(r)=\log(e/r)\)。由于任意正幂的衰减快于对数的增长，
\(r^\alpha L(r)^{-\gamma}\to0\)，所以给定函数在 \(\overline B\) 连续，特别属于
\(L^p(B)\)。在原点以外，径向求导得到
\[
 \frac{\mathrm d}{\mathrm dr}
       \bigl(r^\alpha L(r)^{-\gamma}\bigr)
 =r^{\alpha-1}L(r)^{-\gamma}
                     \left(\alpha+\frac\gamma{L(r)}\right).
\]
当 \(r\) 足够小时，最后括号中的量处于两个严格为正的常数之间。因此经典梯度的
\(p\) 次积分在原点附近有限，当且仅当
\[
 \int_0^{1/2}r^{d-1+p(\alpha-1)}L(r)^{-\gamma p}\,\mathrm dr
 =\int_0^{1/2}\frac{\mathrm dr}{rL(r)^{\gamma p}}<\infty.
\]
代换 \(t=L(r)\) 后，这等价于 \(\int_{\log(2e)}^\infty t^{-\gamma p}\,\mathrm dt<\infty\)，也就是 \(\gamma p>1\)。
如果 \(u\in W^{1,p}(B)\)，它的弱梯度在不含原点的区域必与这个经典梯度一致，故该条件必要。

为证明充分性，还须核实原点没有遗漏的分布项。设 \(\gamma p>1\)，取光滑
\(\eta:[0,\infty)\to[0,1]\)，使它在 \([0,1]\) 上为零，在 \([2,\infty)\) 上为一，令
\[
 u_\varepsilon(x)=\eta(|x|/\varepsilon)u(x).
\]
每个 \(u_\varepsilon\) 在 \(B\) 内光滑。记上面算出的经典梯度为 \(g\)，并在原点随意赋值。已经知道 \(g\in L^p(B)\)，且
\[
 \nabla u_\varepsilon
  =\eta(|x|/\varepsilon)g
       +\varepsilon^{-1}\eta'(|x|/\varepsilon)
                         u(x)\frac{x}{|x|}
\]
在原点以外成立。第一项在 \(L^p\) 中趋于 \(g\)。第二项仅支于
\(\{\varepsilon<|x|<2\varepsilon\}\)，其范数的 \(p\) 次方不超过
\[
 C\varepsilon^{-p}\varepsilon^{d+\alpha p}
                  L(\varepsilon)^{-\gamma p}
   =C L(\varepsilon)^{-\gamma p}\longrightarrow0.
\]
这里用到了环带上 \(L(|x|)\) 与 \(L(\varepsilon)\) 的可比性，以及
\(d+\alpha p-p=0\)。同时 \(u_\varepsilon\to u\) 于 \(L^p\)。在弱导数恒等式中通过极限，得到 \(u\in W^{1,p}(B)\) 且 \(\nabla u=g\)。

由\cref{b1:cor:morrey-domain}，这个空间的元素有 \(\alpha\)-Hölder 连续代表。它与题中已有的连续函数几乎处处相同，因而处处相同。可是对任意 \(\beta>\alpha\)，
\[
 \frac{|u(re_1)-u(0)|}{|re_1|^\beta}
       =r^{\alpha-\beta}L(r)^{-\gamma}
                  \longrightarrow\infty.
\]
所以任何更大的 Hölder 指数都不能用于这个固定函数。对数衰减使临界梯度可积，却不能补偿正幂之间的差距。
\end{solution}

\begin{exercise}\label{b1:ex:24-moser-concentration}
设 \(d\ge2\)，记 \(\omega_d=|B(0,1)|\)、\(\sigma_{d-1}=d\omega_d\)，固定 \(R=1/2\)。对 \(L\ge1\)，在 \(\mathbb R^d\) 上定义
\[
 u_L(x)=\sigma_{d-1}^{-1/d}
 \begin{cases}
  L^{(d-1)/d},& |x|\le R e^{-L},\\
  L^{-1/d}\log(R/|x|),& R e^{-L}<|x|<R,\\
  0,&|x|\ge R.
 \end{cases}
\]
\begin{enumerate}
 \item\label{b1:item:24-moser-norms}
 证明 \(u_L\in W_0^{1,d}(B(0,1))\)，并且
 \(\|\nabla u_L\|_d=1\)、\(\|u_L\|_d\to0\)。
 \item\label{b1:item:24-moser-supercritical-exponent}
 证明若 \(q>d/(d-1)\)，则对于每个 \(a>0\)，都有
 \[
  \int_{B(0,1)}\exp(a|u_L|^q)\,\mathrm dx
                       \longrightarrow\infty.
 \]
 在 \(q=d/(d-1)\) 时，进一步证明同样的发散对于
 \(a>d\,\sigma_{d-1}^{1/(d-1)}\) 成立。
\end{enumerate}
由此说明超临界指数不可能具有梯度单位球上的统一指数积分界；本题不要求证明临界系数上界能够取得。
\end{exercise}
% 来源：自拟整理经典径向对数集中机制；这里逐项证明所需结论，不冒认未实读的 Moser 原题或最优常数证明。
\begin{solution}
对于每个固定 \(L\)，径向剖面连续且分段光滑，导数有界。因此 \(u_L\) 是全空间的
Lipschitz 函数，支集包含于 \(\overline{B(0,R)}\Subset B(0,1)\)。它属于
\(W^{1,d}\)，由\cref{b1:prop:compact-sobolev-smooth-approximation}的内部紧支集逼近，得到 \(u_L\in W_0^{1,d}(B(0,1))\)。

梯度只在中间环带非零，其模为
\(\sigma_{d-1}^{-1/d}L^{-1/d}|x|^{-1}\)。极坐标积分给出
\[
 \int_{\mathbb R^d}|\nabla u_L|^d\,\mathrm dx
       =L^{-1}\int_{Re^{-L}}^R\frac{\mathrm dr}{r}=1.
\]
对函数本身，在环带使用 \(t=\log(R/r)\) 代换，得到
\[
 \|u_L\|_d^d
 =\frac{\omega_dR^d}{\sigma_{d-1}}L^{d-1}e^{-dL}
       +\frac{R^d}{L}\int_0^L t^d e^{-dt}\,\mathrm dt.
\]
第一个项趋于零；第二个项也趋于零，因为
\(\int_0^\infty t^d e^{-dt}\,\mathrm dt<\infty\)。这证明第一问的两个范数结论。

内球 \(B(0,Re^{-L})\) 上的函数值是常数，故
\[
 \int_{B(0,1)}\exp(a|u_L|^q)\,\mathrm dx
 \ge\omega_dR^d
       \exp\left(-dL+
          a\sigma_{d-1}^{-q/d}L^{q(d-1)/d}\right).
\]
若 \(q>d/(d-1)\)，最后的正项关于 \(L\) 的幂次严格大于 \(1\)，因而它超过线性负项，右端趋于无穷。若 \(q=d/(d-1)\)，指数变为
\[
 \left(a\sigma_{d-1}^{-1/(d-1)}-d\right)L.
\]
题设对 \(a\) 的条件恰使括号为正，所以仍然发散。这给出了临界指数下系数的一个必要上界；证明没有给出该界处的统一可积估计。
\end{solution}

\begin{exercise}\label{b1:ex:24-higher-order-scaling}
设 \(B=B(0,1)\subset\mathbb R^d\)、\(k\ge2\) 为整数、\(1\le p<\infty\)。本题只测试
\(C_c^\infty(B)\) 中的函数，不使用高阶延拓定理。
\begin{enumerate}
 \item\label{b1:item:24-higher-order-lq-obstruction}
 当 \(kp<d\) 时，设 \(1\le q\le\infty\)。证明若
 \[
  \|u\|_{L^q(B)}\le C\|u\|_{W^{k,p}(B)}
       \quad\text{对全部 }u\in C_c^\infty(B)
 \]
 成立，则必须有 \(q\le dp/(d-kp)\)。这里允许 \(q=\infty\)。
 \item\label{b1:item:24-higher-order-holder-obstruction}
 设 \(0\le j<k\)、\(0<\beta\le1\)。若存在有限常数 \(C\)，使所有
 \(u\in C_c^\infty(B)\) 都满足
 \[
  \max_{|\kappa|=j}[D^\kappa u]_{0,\beta;\overline B}
        \le C\|u\|_{W^{k,p}(B)},
 \]
 证明必要关系为 \(j+\beta\le k-d/p\)。
\end{enumerate}
说明为什么第一问的同一种集中计算，在 \(kp=d\) 时不能单独排除目标空间
\(L^\infty\)。
\end{exercise}
% 来源：自拟；高阶非齐次范数中各阶项的尺度分离，以及高阶 Hölder 目标的必要限制。
\begin{solution}
取非零 \(\varphi\in C_c^\infty(B)\)，并对 \(0<\varepsilon<1\) 令
\[
 u_\varepsilon(x)=\varepsilon^{k-d/p}\varphi(x/\varepsilon).
\]
它的支集包含于 \(\varepsilon\operatorname{supp}\varphi\Subset B\)。对每个
\(|\kappa|\le k\)，变量代换给出
\[
 \|D^\kappa u_\varepsilon\|_{L^p(B)}
       =\varepsilon^{k-|\kappa|}
                    \|D^\kappa\varphi\|_{L^p(\mathbb R^d)}.
\]
因此这列函数的 \(W^{k,p}(B)\) 范数一致有界；最高阶项保持不变，低阶项趋于零。
另一方面，对 \(q<\infty\)，
\[
 \|u_\varepsilon\|_{L^q(B)}
       =\varepsilon^{k-d/p+d/q}\|\varphi\|_q,
\]
而 \(q=\infty\) 时同式按 \(d/q=0\) 理解。若 \(kp<d\) 且
\(q>dp/(d-kp)\)，这个幂次严格为负，左端趋于无穷，与统一估计矛盾。

为检验第二问，选取 \(\varphi\)、一个 \(|\kappa|=j\) 的多重指标，以及
\(a,b\in B\)，使 \(a\ne b\) 且
\(D^\kappa\varphi(a)\ne D^\kappa\varphi(b)\)。例如取一个相应导数不恒定的光滑紧支集函数即可。对同一集中族，在两个点
\(\varepsilon a,\varepsilon b\) 处计算得到
\[
 \frac{|D^\kappa u_\varepsilon(\varepsilon a)
             -D^\kappa u_\varepsilon(\varepsilon b)|}
      {|\varepsilon a-\varepsilon b|^\beta}
 =\varepsilon^{k-d/p-j-\beta}
       \frac{|D^\kappa\varphi(a)-D^\kappa\varphi(b)|}
            {|a-b|^\beta}.
\]
最后一个分式是严格为正的固定数。因此若
\(j+\beta>k-d/p\)，左侧趋于无穷，而右侧所允许的 Sobolev 范数仍然有界，矛盾。

在临界关系 \(kp=d\) 下，第一问的函数族变成
\(u_\varepsilon(x)=\varphi(x/\varepsilon)\)，其上确界范数不随 \(\varepsilon\) 改变。这个计算既没有制造上确界发散，也没有证明统一 \(L^\infty\) 嵌入成立。临界情形需要对数剖面等不同于单纯重缩放的构造。
\end{solution}

\begin{exercise}\label{b1:ex:24-morrey-integrability-balance}
设 \(d<p<\infty\)、\(1\le r<\infty\)，令
\[
 \alpha=1-\frac dp,\quad
 a=\frac d{d+\alpha r}.
\]
证明对于每个
\(u\in W^{1,p}(\mathbb R^d)\cap L^r(\mathbb R^d)\)，其连续代表满足
\[
 \|u\|_\infty
    \le C_{d,p,r}\|\nabla u\|_p^a
                    \cdot\|u\|_r^{\,1-a}.
\]
提示：比较一个点的函数值与半径 \(R\) 的球平均，再依两个范数选择 \(R\)。
还应处理梯度范数为零的情形，并说明为何任何具有相同两因子形式、对全空间全部函数成立的估计，都必须采用这里的指数 \(a\)。
\end{exercise}
% 来源：自拟；将全空间 Morrey 振荡控制与额外 Lr 可积性平衡，不使用紧性或插值空间理论。
\begin{solution}
记 \(G=\|\nabla u\|_p\)、\(A=\|u\|_r\)。由定理%
\ref{b1:thm:morrey-whole-space}，存在唯一连续代表，使
\[
 |u(x)-u(y)|\le C_{d,p}G|x-y|^\alpha
                  \quad(x,y\in\mathbb R^d).
\]
若 \(G=0\)，这个代表为全空间常函数；有限 \(r\) 的可积性迫使它等于零，所求估计成立。
若 \(A=0\)，函数同样为零。因此以下只考虑 \(G,A>0\)。

固定 \(x\in\mathbb R^d\) 及任意 \(R>0\)，对 \(y\in B(x,R)\) 平均上述振荡界，再使用 Hölder 不等式，得到
\[
 \begin{aligned}
 |u(x)|
 &\le\frac1{|B(x,R)|}\int_{B(x,R)}|u(y)|\,\mathrm dy
                                  +C_{d,p}G R^\alpha\\
 &\le\omega_d^{-1/r}R^{-d/r}A+C_{d,p}G R^\alpha.
 \end{aligned}
\]
在 \(r=1\) 时，第二行只是把局部积分放大为全空间积分。选择
\[
 R=(A/G)^{1/(\alpha+d/r)}
\]
后，两个项都由一个只依赖 \(d,p,r\) 的常数乘
\[
 G^{\,d/(d+\alpha r)}A^{\,\alpha r/(d+\alpha r)}
\]
控制。此界与 \(x\) 无关，取上确界即得结论。

最后，假设同样的估计能以某个固定指数 \(b\) 代替 \(a\)。取非零
\(\varphi\in C_c^\infty(\mathbb R^d)\)，对全部 \(\lambda>0\) 测试
\(\varphi_\lambda(x)=\varphi(\lambda x)\)。其上确界不变，而
\[
 \|\nabla\varphi_\lambda\|_p
        =\lambda^\alpha\|\nabla\varphi\|_p,\quad
 \|\varphi_\lambda\|_r
        =\lambda^{-d/r}\|\varphi\|_r.
\]
估计右端中的尺度因子为 \(\lambda^{b\alpha-(1-b)d/r}\)。若幂次非零，令
\(\lambda\) 趋于零或无穷即可使右端趋于零，违反左端严格为正这一事实。所以
\(b\alpha=(1-b)d/r\)，解得 \(b=a\)。这一必要性来自尺度平衡，并未使用任何极值函数或紧性结论。
\end{solution}

% !TeX root = ../../Book1.tex
\chapter{Sobolev 空间中的紧性}
\label{b1:ch:25}
\BookChapterNumber{b1:ch:25}
\begin{chapterabstract}
本章把 Sobolev 范数的统一估计转化为函数列的收敛结论。先以有限网格平均证明有限指数 \(L^p\) 的平移紧性判据，再结合延拓、梯度平移估计及插值得到 Rellich–Kondrachov 定理。随后区分弱极限的识别、较弱范数中的强收敛和原 Sobolev 范数的强收敛，并在具体变分问题中说明各自的用途。临界集中和粗糙边界的例子用来检验结论的适用范围。
\end{chapterabstract}

\begin{learninggoals}
\begin{enumerate}
\item\label{b1:goal:25-fk}用尾部和一致平移连续性构造有限维逼近，证明 Fréchet–Kolmogorov 判据，辨别单个函数的连续性与函数族的一致条件。
\item\label{b1:goal:25-rellich}从弱梯度控制小平移，证明积分范数和连续函数范数中的紧嵌入，并说明严格低于临界指数的必要性。
\item\label{b1:goal:25-boundary}区分一般 Lipschitz 区域、任意有界开集上的零边界闭包及内部紧性，高阶结论保留相应的延拓假设。
\item\label{b1:goal:25-weak}识别 Sobolev 弱极限，分别通过紧映射和标准范数的收敛恢复强收敛，避免对非线性表达式直接代入弱极限。
\item\label{b1:goal:25-variation}完整实施一个凸能量的极小值证明，并在习题中用紧性保留非凸的单位范数约束。
\end{enumerate}
\end{learninggoals}

本章以前一章的连续嵌入为起点，但“连续”与“紧”所回答的问题不同。连续嵌入对每个函数给出同一常数的估计；紧嵌入则使任意有界列都能抽出强收敛子列。有限维空间中有界列自然具有这一性质，无限维函数空间却存在振荡、逃逸和集中，必须逐一控制。

阅读\secref{b1:sec:2501}时可以暂时忘记弱导数，只关心怎样用有限个平均值逼近整族函数。\secref{b1:sec:2502}再将这个判据用于 Sobolev 空间，所需延拓已在第22章建立。最后一节接回第9章的自反性与弱下半连续性：先问近似问题给了哪种有界估计，再问通过极限究竟需要哪种收敛，而不是一开始就把所有收敛方式混为一谈。

平移判据可对照 Hanche-Olsen 与 Holden 的《The Kolmogorov--Riesz compactness theorem》\cite[Section 3]{HancheOlsenHolden2010Compactness}阅读；本章把其有限维投影路线中的网格平均误差和远处截断分别展开。Sobolev 紧嵌入及变分方法的系统伴读可选 Adams 与 Fournier 的《Sobolev Spaces》\cite{Adams2003Sobolev}和 Brezis 的《Functional Analysis, Sobolev Spaces and Partial Differential Equations》\cite{Brezis2011Functional}。本章的证明始终从前面已经建立的工具出发，不要求读者预先掌握第二卷的偏微分方程理论。

% !TeX root = ../../Book1.tex
\section{\texorpdfstring{\(L^p\)}{Lᵖ} 中的平移判据}
\label{b1:sec:2501}

范数有界只能把函数限制在一个大球内，不能阻止它们向无穷远移动，也不能阻止固定区域内越来越快的振荡。强紧性需要同时限制这些变化。本节用远处的积分尾部和小平移误差刻画它，并把证明落实为有限个网格平均值的逼近；暂不要求函数具有弱导数。

Hanche-Olsen 与 Holden 的《The Kolmogorov--Riesz compactness theorem》\cite[Section 3, Theorem 5]{HancheOlsenHolden2010Compactness}采用有限维投影处理这一问题，并讨论不同文献中的定理命名。本节固定 \(d\geq1\)、\(1\leq p<\infty\) 和 \(\mathbb K=\mathbb R\) 或 \(\mathbb C\)。所有紧性均指 \(L^p\) 范数拓扑下的紧性，不是弱紧性。

% 实读 HancheOlsenHolden2010Compactness 的作者公开副本：
% https://www.math.ntnu.no/conservation/2009/037.pdf ，第3节Theorem5及完整证明，内页3--4。
% 正式发表：Expositiones Mathematicae 28(4) (2010),385--394，
% DOI10.1016/j.exmath.2010.03.001；作者副本首页2009是稿本日期，不作发表年。
% 本节将有限投影拆成全空间网格平均与有限cell截断，保留差集体积因子2^d。
% 必要性用有限网；任意域的零延拓结论只在Lp层面，未借后节Sobolev紧嵌入。

\subsection{平移连续性}
\label{b1:subsec:250101}

沿用 \(\tau_hu(x)=u(x-h)\)。平移是 \(L^p(\mathbb R^d)\) 上的线性等距映射，并满足
\[
 \tau_h\tau_k=\tau_{h+k},\quad
 \|\tau_hu-\tau_hv\|_p=\|u-v\|_p.
\]
由\cref{b1:lem:lp-translation-continuity}，每个固定的 \(u\in L^p\) 都有
\(\|\tau_hu-u\|_p\to0\)。但这里允许控制 \(h\) 的尺度随 \(u\) 改变；若面对一整个函数族，就需要一个适用于所有成员的尺度。

为把量词写清楚，对非空族 \(\mathcal F\subseteq L^p(\mathbb R^d)\) 记
\[
 \omega_{\mathcal F}(\delta)
 =\sup_{u\in\mathcal F}\sup_{|h|\leq\delta}
       \|\tau_hu-u\|_p.
\]
于是所需条件为 \(\omega_{\mathcal F}(\delta)\to0\)：给定误差后，先选同一个 \(\delta\)，再允许任取 \(u\in\mathcal F\) 和 \(|h|\leq\delta\)。这比逐个函数的平移连续性更强。

这种一致性可以控制函数在每个小网格内被常数代替时的误差。固定 \(\delta>0\)，用半开立方体
\[
 Q_{\delta,k}=\delta\bigl(k+[0,1)^d\bigr),
 \quad k\in\mathbb Z^d
\]
划分全空间；这些集合两两不交，并覆盖每个点。定义
\[
 (P_\delta u)(x)=\frac1{\delta^d}\int_{Q_{\delta,k}}u(y)\,\mathrm dy
 \quad(x\in Q_{\delta,k}).
\]
有限测度上的 Hölder 不等式保证每个平均值有定义。即使 \(u\) 是复值函数，平均值仍按普通复积分定义，不取共轭。

\begin{lemma}\label{b1:lem:lp-cell-average}
算子 \(P_\delta\) 是 \(L^p(\mathbb R^d)\) 上的线性投影，且
\(\|P_\delta u\|_p\leq\|u\|_p\)。此外，
\begin{equation}\label{b1:eq:lp-cell-average-error}
 \|u-P_\delta u\|_p^p
 \leq\frac1{\delta^d}
       \int_{[-\delta,\delta]^d}\|\tau_hu-u\|_p^p\,\mathrm dh.
\end{equation}
因此
\[
 \sup_{u\in\mathcal F}\|u-P_\delta u\|_p
 \leq2^{d/p}\omega_{\mathcal F}(\sqrt d\,\delta).
\]
\end{lemma}
\begin{proof}
对每个网格 \(Q=Q_{\delta,k}\)，平均值的 Hölder 估计给出
\[
 \delta^d\left|\frac1{\delta^d}\int_Q u\right|^p
 \leq\int_Q|u|^p.
\]
对 \(k\) 求和，得到收缩性。平均一个已在每个网格上为常数的函数，不会改变它，故 \(P_\delta^2=P_\delta\)。

对 \(x\in Q\)，同一估计应用于 \(u(x)-u(y)\)，得
\[
 |u(x)-P_\delta u(x)|^p
 \leq\frac1{\delta^d}\int_Q|u(x)-u(y)|^p\,\mathrm dy.
\]
先在 \(x\) 上积分并对网格求和，再作换元 \(y=x+h\)。同一个网格中的两个点满足 \(h\in[-\delta,\delta]^d\)，所以
\[
 \begin{aligned}
 \|u-P_\delta u\|_p^p
 &\leq\frac1{\delta^d}
      \sum_k\int_{Q_{\delta,k}}\int_{Q_{\delta,k}}
             |u(x)-u(y)|^p\,\mathrm dy\,\mathrm dx\\
 &\leq\frac1{\delta^d}
      \sum_k\int_{Q_{\delta,k}}\int_{[-\delta,\delta]^d}
             |u(x)-u(x+h)|^p\,\mathrm dh\,\mathrm dx\\
 &=\frac1{\delta^d}
      \int_{[-\delta,\delta]^d}\|\tau_{-h}u-u\|_p^p\,\mathrm dh.
 \end{aligned}
\]
这些是非负积分，Tonelli 定理允许求和及换序。积分区域关于原点对称，令 \(h\) 换成 \(-h\) 即得\eqref{b1:eq:lp-cell-average-error}。最后，该区域体积为 \((2\delta)^d\)，其中每个向量满足 \(|h|\leq\sqrt d\,\delta\)，从而得到因子 \(2^{d/p}\)。
\end{proof}

差向量所在的立方体边长为 \(2\delta\)，不是 \(\delta\)。因此上式右端不是一个已经归一化为一的平均；这个体积差正是因子 \(2^d\) 的来源。

\subsection{Fréchet–Kolmogorov 定理}
\label{b1:subsec:250102}

先回顾证明中需要的度量事实。若对每个 \(\varepsilon>0\)，一个集合 \(A\) 都能被有限个半径为 \(\varepsilon\) 的球覆盖，则称 \(A\) 为\term[quanyoujie]{全有界}[totally bounded]集；这些球的中心构成一个有限 \(\varepsilon\)-网。

\begin{lemma}\label{b1:lem:complete-total-bounded}
在完备度量空间中，一个子集相对紧，当且仅当它全有界。
\end{lemma}
\begin{proof}
紧闭包的任意固定半径球覆盖都有有限子覆盖，故必要性成立。反过来，若 \(A\) 全有界，则闭包 \(K=\overline A\) 也全有界：先用半径为 \(\varepsilon/2\) 的有限网覆盖 \(A\)，把半径扩大到 \(\varepsilon\) 即覆盖闭包。

对 \(K\) 中任意序列，逐次使用半径 \(2^{-j}\) 的有限覆盖，保留落在同一球中的无限子列，再对角抽选，所得子列为 Cauchy 列。由于 \(K\) 闭且环境完备，它在 \(K\) 中收敛。

这也给出开覆盖意义下的紧性。对 \(K\) 的任意开覆盖，必有 \(\eta>0\)，使每个 \(x\in K\) 的相对球 \(B_K(x,\eta)\) 包含于某个覆盖成员。否则可选 \(x_j\)，使 \(B_K(x_j,1/j)\) 不包含于任何单个成员；取收敛子列及包含其极限的开集，就会在充分大时得到相反结论。有限网的中心也可选在 \(K\) 中：从每个与 \(K\) 相交的网球中取一点作新中心，至多把半径扩大一倍。因此可用中心在 \(K\) 内的有限 \(\eta/2\)-网覆盖 \(K\)，为每个中心选择一个容纳其 \(\eta\)-球的覆盖成员，便得到有限子覆盖。
\end{proof}

现在引入\term[frechet-kolmogorov-dingli]{Fréchet–Kolmogorov 定理}[Fréchet--Kolmogorov compactness theorem]。条件分别写成范数、尾部与平移的形式，便于在具体问题中使用。

\begin{theorem}[Fréchet–Kolmogorov]\label{b1:thm:frechet-kolmogorov-lp}
设 \(1\leq p<\infty\)，\(\mathcal F\subseteq L^p(\mathbb R^d;\mathbb K)\) 非空。则 \(\mathcal F\) 相对紧，当且仅当下列条件都成立。
\begin{enumerate}
\item\label{b1:item:fk-bounded}
范数一致有界：
\[
 \sup_{u\in\mathcal F}\|u\|_p<\infty.
\]
\item\label{b1:item:fk-tail}
远处尾部一致趋零：
\[
 \lim_{R\to\infty}\sup_{u\in\mathcal F}
       \int_{|x|>R}|u(x)|^p\,\mathrm dx=0.
\]
\item\label{b1:item:fk-translation}
小平移误差一致趋零：
\[
 \lim_{\delta\downarrow0}\omega_{\mathcal F}(\delta)=0.
\]
\end{enumerate}
等价地，\(\mathcal F\) 中每个序列都有一个在 \(L^p\) 范数中收敛的子列，极限允许在 \(\overline{\mathcal F}\) 中。
\end{theorem}
\begin{proof}
先证明必要性。相对紧性给出任意尺度的有限网。一个有限 \(1\)-网已经说明范数有界，因为每个 \(u\) 的范数不超过其所靠近的网点范数加一。

固定 \(\varepsilon>0\)，取有限 \(\varepsilon\)-网
\(g_1,\ldots,g_N\)。每个 \(g_j\in L^p\) 都满足
\(\|\mathbf1_{\{|x|>R\}}g_j\|_p\to0\)。网点有限，可以选同一个 \(R\)，使这些尾部范数都小于 \(\varepsilon\)。给定 \(u\in\mathcal F\)，选取
\(\|u-g_j\|_p<\varepsilon\) 的网点，则
\[
 \|\mathbf1_{\{|x|>R\}}u\|_p
 \leq\|u-g_j\|_p+\|\mathbf1_{\{|x|>R\}}g_j\|_p<2\varepsilon.
\]
由此得到尾部条件。

同样，由每个网点的平移连续性，可选同一个 \(\delta>0\)，使
\(\|\tau_hg_j-g_j\|_p<\varepsilon\) 对全部 \(j\) 及
\(|h|\leq\delta\) 成立。平移等距性给出
\[
 \|\tau_hu-u\|_p
 \leq2\|u-g_j\|_p+\|\tau_hg_j-g_j\|_p<3\varepsilon.
\]
这就把有限多个函数的连续性变成了整族的一致连续性。

再证充分性。给定 \(\varepsilon>0\)，由平移条件选 \(\delta>0\)，使
\[
 \sup_{u\in\mathcal F}\|u-P_\delta u\|_p<\varepsilon/3.
\]
这是网格平均引理的直接应用，选取时保留其中的
\(2^{d/p}\) 及 \(\sqrt d\)。再由尾部条件选 \(R\)，使
\[
 \sup_{u\in\mathcal F}\|\mathbf1_{\{|x|>R\}}u\|_p<\varepsilon/3.
\]
保留所有与 \(B(0,R)\) 相交的网格，指标集合记为 \(J_{\delta,R}\)，它是有限集；令
\[
 S_{\delta,R}=\bigcup_{k\in J_{\delta,R}}Q_{\delta,k},
 \quad A_{\delta,R}u=\mathbf1_{S_{\delta,R}}P_\delta u.
\]
未保留的网格全部位于球外，至多涉及一个零测球面。逐网格使用平均值收缩估计，得到
\begin{equation}\label{b1:eq:cell-projection-tail}
 \begin{aligned}
 \|P_\delta u-A_{\delta,R}u\|_p^p
 &=\sum_{k\notin J_{\delta,R}}
       \delta^d\left|\frac1{\delta^d}\int_{Q_{\delta,k}}u\right|^p\\
 &\leq\sum_{k\notin J_{\delta,R}}\int_{Q_{\delta,k}}|u|^p
 \leq\int_{|x|>R}|u|^p.
 \end{aligned}
\end{equation}
因此
\(\|u-A_{\delta,R}u\|_p<2\varepsilon/3\)，对整族一致成立。

所有 \(A_{\delta,R}u\) 落在有限维空间
\[
 V_{\delta,R}
 =\operatorname{span}\{\,\mathbf1_{Q_{\delta,k}}:k\in J_{\delta,R}\,\},
\]
且范数一致有界。具体地，若 \(v=\sum_{k\in J_{\delta,R}}a_k\mathbf1_{Q_{\delta,k}}\)，则
\[
 \|v\|_p^p=\delta^d\sum_{k\in J_{\delta,R}}|a_k|^p.
\]
于是每个系数都落在固定的有界区间或有界复平面方框中。将这些有限个坐标分别用足够细的有限网格逼近，即得
\(A_{\delta,R}\mathcal F\) 的有限 \(\varepsilon/3\)-网；每个坐标误差不超过 \(\eta\) 时，总误差不超过
\(\delta^{d/p}|J_{\delta,R}|^{1/p}\eta\)。

三角不等式说明同一组网点构成 \(\mathcal F\) 的有限
\(\varepsilon\)-网。因此 \(\mathcal F\) 全有界。由\cref{b1:thm:lp-complete}及前一引理，闭包紧。引理证明中的抽选过程同时给出所述子列结论；反过来，若不存在某个尺度的有限网，就可逐个选出彼此距离至少为该尺度的序列，它不可能有收敛子列。
\end{proof}

证明始终先对原函数作全空间网格平均，然后截去平均函数的远处部分。\eqref{b1:eq:cell-projection-tail}控制了第二步；没有把“先截断函数、再平移”产生的边界差当成零。范数一致有界用于保证保留下来的有限个平均值有界，尾部与平移条件则决定需要保留多少网格、网格应当多细。

\subsection{尾部一致控制与支集}
\label{b1:subsec:250103}

若所有函数都在同一个有界集外几乎处处为零，尾部条件自动满足。反过来，尾部条件不要求严格的共同紧支集。例如，一个已经在 \(L^p\) 中收敛的序列连同极限是紧集，定理保证其尾部一致小，而各成员的支集完全可以无界。

固定支集不能排除振荡。令 \(Q=(0,1)^d\)，定义
\[
 u_n(x)=\mathbf1_Q(x)(-1)^{\lfloor2^n x_1\rfloor}.
\]
这些函数的 \(L^p\) 范数均为一，且在同一个紧集外为零。当 \(m>n\) 时，在每个较粗的二进小区间内，两种符号相同与相反的部分各占一半，故
\[
 \|u_m-u_n\|_p^p=2^{p-1}.
\]
所以任何子列都不是 Cauchy 列。还可直接看出平移条件失败：取
\(h_n=2^{-n}e_1\)，在
\((2^{-n},1)\times(0,1)^{d-1}\) 上，\(u_n(x-h_n)=-u_n(x)\)，从而
\[
 \|\tau_{h_n}u_n-u_n\|_p^p
 \geq2^p(1-2^{-n}).
\]
每个固定的 \(u_n\) 仍有自己的平移连续性，但这些尺度不能统一。

另一种失败是向远处移动。取 \(g=\mathbf1_Q\)，令
\(v_n=\tau_{3ne_1}g\)。由于平移可交换，
\[
 \|\tau_hv_n-v_n\|_p=\|\tau_hg-g\|_p,
\]
小平移条件对整列一致成立。然而对每个固定 \(R\)，充分远的
\(v_n\) 把全部 \(L^p\) 质量放在 \(B(0,R)\) 外，尾部条件失败。其支集两两不交，故不同项之间的距离为 \(2^{1/p}\)。这与前面的振荡例子是不同的障碍。

现在把判据用于一般定义域。这里只需要 Lebesgue 可测性，不要求边界光滑。

\begin{corollary}\label{b1:cor:lp-domain-zero-extension-compactness}
设 \(\Omega\subseteq\mathbb R^d\) 为 Lebesgue 可测集，且
\(\mathcal F\subseteq L^p(\Omega;\mathbb K)\)。把每个 \(u\) 在域外补零，记为 \(E_0u\)。则 \(\mathcal F\) 在 \(L^p(\Omega)\) 中相对紧，当且仅当它范数一致有界，且
\[
 \lim_{R\to\infty}\sup_{u\in\mathcal F}
      \int_{\Omega\cap\{|x|>R\}}|u|^p=0,\quad
 \lim_{\delta\downarrow0}\sup_{u\in\mathcal F}
      \sup_{|h|\leq\delta}\|\tau_hE_0u-E_0u\|_{L^p(\mathbb R^d)}=0.
\]
若 \(\Omega\) 有界，第一个极限条件自动成立。空族按平凡情形处理。
\end{corollary}
\begin{proof}
补零是 \(L^p(\Omega)\) 到 \(L^p(\mathbb R^d)\) 的线性等距嵌入，其像为在 \(\Omega^c\) 上几乎处处为零的函数所成的闭子空间。闭性可直接检查：若 \(E_0u_j\to f\) 于全空间 \(L^p\)，则
\[
 \|\mathbf1_{\Omega^c}f\|_p
 \leq\|f-E_0u_j\|_p\longrightarrow0.
\]
因此 \(\mathcal F\) 相对紧与 \(E_0\mathcal F\) 相对紧等价。把前一定理用于补零后的族，三个条件正好成为所列条件。
\end{proof}

其中的平移误差必须在整个空间计算。若只在两份定义域的交集上比较函数，就会漏掉跨越边界的部分。事实上，记
\(\Omega+h=\{x+h:x\in\Omega\}\)，分割积分区域可得
\[
 \begin{aligned}
 \|\tau_hE_0u-E_0u\|_p^p
 &=\int_{\Omega\cap(\Omega+h)}
         |u(x-h)-u(x)|^p\,\mathrm dx\\
 &\quad+\int_{\Omega\setminus(\Omega+h)}|u(x)|^p\,\mathrm dx
       +\int_{\Omega\setminus(\Omega-h)}|u(x)|^p\,\mathrm dx.
 \end{aligned}
\]
最后一项由 \(x\in(\Omega+h)\setminus\Omega\) 的积分作平移换元得到。这个恒等式既记录了域内变化，也记录了靠近边界的质量；此处从未假设弱导数可以补零。后续处理 Sobolev 函数时，还须另用延拓定理或零边界条件建立所需的平移估计。

有限 \(p\) 的限制不能直接删去。\secref{b1:sec:1304}已经指出，
\(u=\mathbf1_{(0,1)}\) 满足
\(\|\tau_hu-u\|_\infty=1\)（\(0<|h|<1\)）。因此连单元素紧集
\(\{u\}\) 都不满足上述平移条件的无穷范数版本。网格平均的积分误差证明也没有提供对应的上确界逼近，不能把本定理机械改成 \(p=\infty\)。

% !TeX root = ../../Book1.tex
\section{Rellich–Kondrachov 定理}
\label{b1:sec:2502}

第24章的嵌入估计说明，一族函数的弱导数若有统一积分界，函数本身就有更强的可积性或连续性。但连续嵌入只给出范数估计，尚未保证有收敛子列。本节再用梯度控制小平移，从而把 Sobolev 有界性接到上一节的紧性判据。区域有界提供尾部控制，边界延拓则使平移能够在全空间进行。

两项限制会贯穿证明。第一，一般区域上的函数不能直接补零后求弱导数，必须使用第22章建立的延拓。第二，从连续嵌入得到紧嵌入时，目标指数通常要留出严格余量；临界指数处的缩放会保留一部分范数，而使支集越来越小。

\subsection{有界区域}
\label{b1:subsec:250201}

先明确梯度所提供的统一模。以下有限 \(p\) 的估计已用于第23章的整数阶到分数阶嵌入；把它单列出来，便于识别紧性证明的输入。

\begin{lemma}\label{b1:lem:sobolev-translation-gradient}
若 \(1\le p<\infty\)、\(u\in W^{1,p}(\mathbb R^d)\)，则
\begin{equation}\label{b1:eq:sobolev-translation-gradient}
 \|\tau_hu-u\|_{L^p(\mathbb R^d)}
 \le |h|\cdot\bigl\|\,|\nabla u|\,\bigr\|_{L^p(\mathbb R^d)}
 \quad(h\in\mathbb R^d).
\end{equation}
\end{lemma}
\begin{proof}
先设 \(u\in C_c^\infty(\mathbb R^d)\)。沿线段积分得
\[
 u(x-h)-u(x)=-\int_0^1\nabla u(x-th)\cdot h\,\mathrm dt.
\]
对 \(t\) 的概率积分使用 Jensen 不等式，再对 \(x\) 积分，得到
\[
 \int_{\mathbb R^d}|u(x-h)-u(x)|^p\,\mathrm dx
 \le |h|^p\int_0^1\int_{\mathbb R^d}|\nabla u(x-th)|^p
                \,\mathrm dx\,\mathrm dt
 =|h|^p\int_{\mathbb R^d}|\nabla u|^p.
\]
由\cref{b1:thm:whole-space-sobolev-density}取光滑逼近；平移等距，梯度的欧氏 \(L^p\) 范数又受标准 Sobolev 范数控制，故两端都能通过极限。
\end{proof}

\begin{proposition}\label{b1:prop:rellich-base-lp}
设 \(\Omega\subset\mathbb R^d\) 为有界 Lipschitz 区域，\(1\le p<\infty\)。则 \(W^{1,p}(\Omega)\) 的任意有界子集在 \(L^p(\Omega)\) 中相对紧。
\end{proposition}
\begin{proof}
由\cref{b1:thm:lipschitz-domain-extension}，存在固定的有界线性延拓算子 \(E\)，其像中所有函数都在一个与函数无关的紧集 \(K\) 外为零。若 \(\mathcal F\) 在 \(W^{1,p}(\Omega)\) 中有界，则 \(E\mathcal F\) 在全空间 \(W^{1,p}\) 中有界，因而在 \(L^p\) 中有界；共同支集保证远处尾部为零；上一引理又给出
\[
 \sup_{u\in\mathcal F}\|\tau_hEu-Eu\|_p\le C|h|.
\]
Fréchet–Kolmogorov 定理的三个条件都满足。因此任取 \(u_j\in\mathcal F\)，可抽取子列使 \(Eu_j\) 在全空间 \(L^p\) 中收敛。限制到 \(\Omega\) 不增加 \(L^p\) 距离，就得到所需收敛子列。
\end{proof}

证明不要求 \(\Omega\) 连通，因为所用的一阶延拓不要求连通。它也不要求事先知道序列有弱极限，因而包括 \(p=1\)。不过，把“有界 Lipschitz 区域”仅改成“有界开集”，对一般 \(W^{1,p}\) 就不成立；第\ref{b1:ex:25-infinitely-many-components}题用越来越小的连通分支给出反例。

\subsection{次临界嵌入}
\label{b1:subsec:250202}

若赋范空间之间的包含映射连续，并把有界集映为相对紧集，就称这个嵌入为\term[jinqianru]{紧嵌入}[compact embedding]。下面的结果通常称为\term[rellich-kondrachov-dingli]{Rellich–Kondrachov 定理}[Rellich--Kondrachov theorem]。这里分别列出积分范数与连续函数范数的结论。

\begin{theorem}[Rellich–Kondrachov，有限目标指数]\label{b1:thm:rellich-subcritical}
设 \(\Omega\subset\mathbb R^d\) 为有界 Lipschitz 区域，\(1\le p<\infty\)。下列嵌入是紧的：
\begin{enumerate}
\item 若 \(p<d\)，则 \(W^{1,p}(\Omega)\subset L^q(\Omega)\)，其中 \(1\le q<dp/(d-p)\)。
\item 若 \(p=d\)，则 \(W^{1,p}(\Omega)\subset L^q(\Omega)\)，其中 \(1\le q<\infty\)。
\item 若 \(p>d\)，则 \(W^{1,p}(\Omega)\subset L^q(\Omega)\)，其中 \(1\le q<\infty\)；实际上连续代表还满足下面的更强结论。
\end{enumerate}
\end{theorem}
\begin{proof}
取 Sobolev 有界列。前一命题给出在 \(L^p\) 中收敛的子列。当 \(q=p\) 时，这正是前一命题的结论。以下只需处理 \(q\ne p\)。若 \(q<p\)，有限测度上的 Hölder 不等式
\[
 \|v\|_q\le m_d(\Omega)^{1/q-1/p}\|v\|_p
\]
已经给出 \(L^q\) 收敛。若存在 \(r>q>p\)，使原列的 \(L^r\) 范数一致有界，则对任意两项之差使用第10章的插值估计，得到
\[
 \|u_j-u_k\|_q
 \le \|u_j-u_k\|_p^\theta
       \cdot\|u_j-u_k\|_r^{1-\theta},
 \qquad
 \frac1q=\frac\theta p+\frac{1-\theta}r.
\]
于是子列在 \(L^q\) 中也是 Cauchy 列。由完备性它有 \(L^q\) 极限；有限测度上两种范数收敛都蕴含依测度收敛，故极限相同。

当 \(p<d\) 时，由第24章的域上 Sobolev 嵌入可取 \(r=p^*=dp/(d-p)\)。当 \(p=d\ge2\) 时，\secref{b1:sec:2405}已经证明任意有限指数的连续嵌入，任选有限的 \(r>q\) 即可。若不读该选读节，也可先在有限测度的 \(\Omega\) 上降到 \(p_0<d\)，选 \(p_0\) 足够接近 \(d\) 使 \(dp_0/(d-p_0)>q\)，再用次临界嵌入获得同一 \(L^r\) 界。

剩下 \(d=p=1\)。利用一阶延拓取得共同有界支集的 \(W^{1,1}(\mathbb R)\) 函数。取一个包含这些支集的有限区间；由主线\secref{b1:sec:1604}中的绝对连续代表与微积分基本定理（\cref{b1:thm:ac-fundamental-calculus}），其代表满足
\[
 |Eu(x)|\le\int_{\mathbb R}|(Eu)'(t)|\,\mathrm dt.
\]
这是从支集左侧的零值沿区间积分得到的。因此有统一 \(L^\infty\) 界，在有限测度的 \(\Omega\) 上也有任意有限 \(L^r\) 界，同样完成插值。

若 \(p>d\)，Morrey 不等式给出连续代表的统一上确界界，再取有限 \(r>q\) 即可。上述各步也同时给出包含映射的连续性。
\end{proof}

这个证明的紧性来源是 \(L^p\) 中的平移控制；更高指数来自第24章已经建立的连续估计。插值中的 \(\theta>0\) 很重要：正是这一幂次把已知趋零的 \(L^p\) 差保留下来。到 \(q=p^*\) 时，不能再由此获得趋零因子。

\begin{theorem}[Rellich–Kondrachov，连续代表]\label{b1:thm:rellich-holder}
设 \(\Omega\) 为有界 Lipschitz 区域。
若 \(d<p<\infty\)，记 \(\alpha=1-d/p\)；若 \(p=\infty\)，记 \(\alpha=1\)。
则连续代表给出的映射
\[
 W^{1,p}(\Omega)\longrightarrow C(\overline\Omega)
 \quad\text{及}\quad
 W^{1,p}(\Omega)\longrightarrow C^{0,\beta}(\overline\Omega)
 \quad(0<\beta<\alpha)
\]
都是紧嵌入。其中 \(C(\overline\Omega)\) 取上确界范数。
\end{theorem}
\begin{proof}
由延拓及 Morrey 估计，有限 \(p>d\) 时，有界列的连续代表满足
\[
 \sup_j\|u_j\|_{C^{0,\alpha}(\overline\Omega)}\le M.
\]
若 \(p=\infty\)，第21章的 Lipschitz 代表与第22章的全空间延拓给出同样的结论，指数取 \(1\)。这时并未把有限 \(p\) 的平移判据推广到无穷范数。

下面直接证明所需的一致收敛子列。紧集 \(\overline\Omega\) 在每个尺度 \(1/n\) 都有有限网；把这些网点排成一个可数集。每个网点处的数列有界，对网点依次抽取子列并取对角列，便使所有网点处的值都收敛。给定 \(\varepsilon>0\)，先取一个足够细的有限网，使 \(2M|x-y|^\alpha<\varepsilon/2\) 对相邻点与网点成立；再把这有限个网点处的两项差控制在 \(\varepsilon/2\) 内。三角不等式说明对角列在整个 \(\overline\Omega\) 上一致 Cauchy。逐点取极限并使用这一统一控制，所得极限连续，子列一致收敛。

还须增强到较低阶 Hölder 范数。对任意 \(v\in C^{0,\alpha}(\overline\Omega)\)，令 \(A=\|v\|_\infty\)、\(B=[v]_{0,\alpha}\)。当 \(A,B>0\) 时，对不同的两点及 \(r=|x-y|\)，
\[
 |v(x)-v(y)|\le\min\{2A,Br^\alpha\}.
\]
在 \(r=(2A/B)^{1/\alpha}\) 两侧分别使用两个估计，得到
\begin{equation}\label{b1:eq:holder-interpolation-compact}
 [v]_{0,\beta}\le
 (2A)^{1-\beta/\alpha}B^{\beta/\alpha}.
\end{equation}
若 \(A=0\) 或 \(B=0\)，结论直接成立。子列的一致极限仍有 Hölder 半范数不超过 \(M\)，因为每个固定两点的估计都能通过极限。对 \(v=u_j-u\) 使用\eqref{b1:eq:holder-interpolation-compact}，右端的 \(A\) 趋零而 \(B\le2M\)，于是得到 \(C^{0,\beta}\) 收敛。
\end{proof}

这里需要 \(\beta<\alpha\)，与积分情形的严格余量相对应。第\ref{b1:ex:25-critical-mass-concentration}题和第\ref{b1:ex:25-holder-endpoint}题分别计算临界积分指数与临界 Hölder 指数的集中例子：有界性仍在，但紧性丧失。另一方面，\(W^{1,1}\) 在区间上虽由绝对连续代表连续嵌入 \(L^\infty\)，其有界列并没有统一的连续模，因此本定理没有把 \(p=d=1\) 列为连续函数空间中的紧嵌入。

\subsection{边界条件}
\label{b1:subsec:250203}

区域边界的规则性与函数自身的边界条件，可以分别提供进入全空间所需的通道。对零边界闭包，补零已由第22章证明能够保留弱导数，因此不必再次要求 Lipschitz 边界。

\begin{corollary}\label{b1:cor:rellich-zero-boundary}
设 \(\Omega\subset\mathbb R^d\) 为任意有界开集。若 \(1\le p<\infty\)，则下列嵌入是紧的：
\begin{enumerate}
\item 当 \(p<d\) 时，\(W_0^{1,p}(\Omega)\subset L^q(\Omega)\)，其中 \(1\le q<dp/(d-p)\)；
\item 当 \(p=d\) 时，\(W_0^{1,p}(\Omega)\subset L^q(\Omega)\)，其中 \(1\le q<\infty\)；
\item 当 \(p>d\) 时，\(W_0^{1,p}(\Omega)\subset L^q(\Omega)\)，其中 \(1\le q<\infty\)。
\end{enumerate}
若 \(p=\infty\)，则 \(W_0^{1,\infty}(\Omega)\subset L^q(\Omega)\) 对每个 \(1\le q<\infty\) 也是紧嵌入。进一步，当 \(d<p<\infty\) 时令 \(\alpha=1-d/p\)，当 \(p=\infty\) 时令 \(\alpha=1\)；补零后的连续代表限制到 \(\overline\Omega\)，给出
\[
 W_0^{1,p}(\Omega)\longrightarrow C(\overline\Omega),
 \qquad
 W_0^{1,p}(\Omega)\longrightarrow C^{0,\beta}(\overline\Omega)
 \quad(0<\beta<\alpha)
\]
这两类紧嵌入。
\end{corollary}
\begin{proof}
有限 \(p\) 时，\cref{b1:prop:w0-zero-extension}把有界列等距送入 \(W^{1,p}(\mathbb R^d)\)，共同支集包含于紧集 \(\overline\Omega\)。梯度平移估计和 Fréchet–Kolmogorov 定理直接给出全空间 \(L^p\) 收敛子列。更高积分指数的界，取第24章的零边界版本；临界 \(p=d\ge2\) 也可先在包含 \(\overline\Omega\) 的球上降低指数，再用该球的连续嵌入。其余插值步骤完全相同。

当 \(p>d\) 时，对全空间补零函数使用 Morrey 估计；\(p=\infty\) 时使用全空间 Lipschitz 代表。于是有限网抽列及 Hölder 插值的证明都直接用于紧集 \(\overline\Omega\)。所得一致收敛又在有限测度的 \(\Omega\) 上蕴含每个有限 \(L^q\) 范数中的收敛。
\end{proof}

上述连续代表确实在每个边界点为零，而不只是在域外几乎处处为零。由 \(W_0^{1,p}\) 的闭包定义，可用 \(C_c^\infty(\Omega)\) 函数在 Sobolev 范数下逼近。其补零函数在全空间有同一范数收敛；对 \(p>d\) 的 Morrey 估计或 \(p=\infty\) 的上确界控制，使这些补零函数一致趋于连续代表。每个逼近函数在 \(\partial\Omega\) 都为零，极限也为零。这个理由包括边界中可能存在的裂缝，不能只凭“域外几乎处处为零”略过。

在有界 Lipschitz 区域上，有限 \(p\) 的零迹刻画把这个推论与第23章接起来。固定边界迹也不会破坏紧性：若 \(g\in W^{1,p}(\Omega)\) 固定，而 \(u_j-g\in W_0^{1,p}(\Omega)\) 且 \(u_j\) 有界，就对 \(u_j-g\) 使用推论，再把 \(g\) 加回去。在任意有界开集上，同样可以用这个闭仿射条件表述边界约束；不必先给粗糙边界定义迹。

内部紧性则完全不需要全域边界条件。若 \(1\le p<\infty\)，且 \((u_j)\) 在 \(W^{1,p}_{\mathrm{loc}}(\Omega)\) 中局部有界，则可抽取共同子列，使它在每个 \(U\Subset\Omega\) 上强收敛于 \(L^p(U)\)。利用同一子列及连续嵌入与插值，还可得到：当 \(p<d\) 时，对每个 \(1\le q<dp/(d-p)\) 强收敛于 \(L^q(U)\)；当 \(p=d\) 时，对每个有限 \(q\) 成立；当 \(p>d\) 时，同样对每个有限 \(q\) 成立。若 \(d<p<\infty\)，记 \(\alpha=1-d/p\)；另若 \((u_j)\) 在 \(W^{1,\infty}_{\mathrm{loc}}(\Omega)\) 中局部有界，则记 \(\alpha=1\)。在这两种情形中可取连续代表，使相应的共同子列在每个 \(U\Subset\Omega\) 上局部一致收敛，并在每个 \(C^{0,\beta}(\overline U)\)、\(0<\beta<\alpha\) 中收敛。

证明这些局部结论时，固定 \(\overline U\Subset\Omega\)，取等于一于 \(\overline U\) 邻域的光滑紧支集截断。乘积补零后形成一个具有共同紧支集的全空间 Sobolev 有界族，故可使用前述相应结论；再对可数内部穷尽作对角抽选。所得结论不保证原区域的全局范数收敛，因为质量仍可能靠近边界或移向无穷远。

\subsection{高阶版本}
\label{b1:subsec:250204}

高阶紧性首先表现为损失一阶导数后的强收敛。为了不预借尚未证明的高阶边界延拓，先把所需条件写在命题中。

\begin{proposition}\label{b1:prop:rellich-higher-order}
设 \(k\ge2\)、\(1\le p<\infty\)，\(\Omega\) 为有界开集。若存在有界线性延拓
\[
 E_k:W^{k,p}(\Omega)\longrightarrow W^{k,p}(\mathbb R^d),
 \qquad (E_ku)|_\Omega=u,
\]
则 \(W^{k,p}(\Omega)\subset W^{k-1,p}(\Omega)\) 是紧嵌入。
无论是否存在这个延拓，令
\[
 W_0^{k,p}(\Omega)
 =\overline{C_c^\infty(\Omega)}^{\,W^{k,p}(\Omega)},
\]
则 \(W_0^{k,p}(\Omega)\subset W^{k-1,p}(\Omega)\) 总是紧嵌入。
\end{proposition}
\begin{proof}
先处理延拓情形。取等于一于 \(\overline\Omega\) 邻域的固定光滑紧支集函数 \(\chi\)。乘子估计说明 \(u\mapsto\chi E_ku\) 有界进入全空间 \(W^{k,p}\)，而且所有像具有共同紧支集。对 \(|\gamma|\le k-1\)，函数 \(\partial^\gamma(\chi E_ku_j)\) 在全空间 \(W^{1,p}\) 中有界，故由梯度平移界和紧性判据可抽出 \(L^p\) 收敛子列。这样的 \(\gamma\) 只有有限多个，逐次抽取可使它们同时收敛。

限制到 \(\Omega\) 后，所得子列在 \(W^{k-1,p}(\Omega)\) 范数中 Cauchy，因为这个范数由刚才有限多个分量的 \(p\) 次和构成。完备性给出该空间中的极限；或者逐项对测试函数通过积分极限，也能确认低阶导数的相容关系。

再看零边界闭包。对 \(u\in W_0^{k,p}(\Omega)\)，取 \(C_c^\infty(\Omega)\) 逼近列。每一项及其各阶导数补零后，在全空间的 \(W^{k,p}\) 距离与原域距离相同，因而为 Cauchy 列。由全空间完备性得到极限；其零阶分量是 \(E_0u\)，所有导数均为原弱导数的补零。因此补零是所需的等距延拓，并具有共同有界支集，前面的有限分量抽列再次适用。
\end{proof}

若关心更高的积分指数，可在上述证明中对每个低阶导数使用第24章的高阶连续嵌入，再与已获得的强 \(L^p\) 收敛插值。具体地，若 \(kp<d\) 且命题中的有界延拓 \(E_k\) 存在，则
\[
 W^{k,p}(\Omega)\subset L^q(\Omega),
 \qquad 1\le q<\frac{dp}{d-kp}
\]
为紧嵌入；对任意有界开集，不要求存在 \(E_k\)，也有
\[
 W_0^{k,p}(\Omega)\subset L^q(\Omega),
 \qquad 1\le q<\frac{dp}{d-kp}
\]
为紧嵌入。右端的严格不等式仍不能省去。

同样，设 \(k-d/p=m+\alpha>0\)，其中 \(m\) 为非负整数、\(0<\alpha<1\)。若有上述延拓 \(E_k\)，则
\[
 W^{k,p}(\Omega)\longrightarrow C^{m,\beta}(\overline\Omega)
 \quad(0<\beta<\alpha)
\]
为紧嵌入；对任意有界开集，相应的
\[
 W_0^{k,p}(\Omega)\longrightarrow C^{m,\beta}(\overline\Omega)
 \quad(0<\beta<\alpha)
\]
也为紧嵌入。第24章给出全部不超过 \(m\) 阶导数的统一连续界及 \(m\) 阶导数的统一 Hölder 界；对这有限多个分量使用有限网抽列和\eqref{b1:eq:holder-interpolation-compact}即可。极限分量的导数关系由内部线段积分通过一致极限确认，与第24章相同。若 \(k-d/p\) 为正整数，则须在该章已证明的低于 \(1\) 的 Hölder 指数中另留严格余量，并分别保留上述延拓或零边界条件。

这些表述中的高阶延拓假设没有因为使用了第24章而消失。对于任意开集，可以始终取得内部版本；对于零边界闭包，可以始终使用补零；对于一般高阶 Sobolev 函数的全边界版本，还必须给出对应的延拓。把不同通道分开，才能看清紧性究竟由哪些已建立的条件保证。

% !TeX root = ../../Book1.tex
\section{弱收敛方法}
\label{b1:sec:2503}

紧嵌入使有界序列在较弱的范数下具有收敛子列，Sobolev 空间的自反性则在原空间中提供弱收敛子列。这两种结论经常一起使用，但承担不同的任务：弱收敛保留导数的积分恒等式，紧嵌入加强函数本身的收敛。至于某个非线性能量能否通过极限，还要看它的凸性、增长条件以及所需的收敛方式。

除最后的实值变分问题外，本节允许实值或复值函数。主要指数范围为 \(1<p<\infty\)，并记 \(p'=p/(p-1)\)。
% 本节证明直接接本书第9章弱闭性、凸泛函弱下半连续性，
% 第21章Sobolev分量结构及本章Rellich结论。
% Brezis2011Functional在来源台账中只有出版社信息与目录证据；
% 文末仅作系统伴读推荐，不将未取得的全文列为某项证明的实读来源。

\subsection{有界序列}
\label{b1:subsec:250301}

设 \(\Omega\subset\mathbb R^d\) 为开集，并已从一个近似问题中取得估计
\[
 \sup_j\left(
   \|u_j\|_{L^p(\Omega)}^p+
   \sum_{i=1}^d\|\partial_i u_j\|_{L^p(\Omega)}^p
 \right)<\infty.
\]
由\cref{b1:prop:sobolev-separable-reflexive}与\ref{b1:cor:sobolev-weak-subsequence}，可抽取子列，使
\[
 u_{j_\ell}\rightharpoonup u
 \quad\text{于 }W^{1,p}(\Omega).
\]
这一抽列不要求 \(\Omega\) 有界，也不要求边界光滑。它来自空间的自反性，而不是来自区域上的紧嵌入。

计算中也常先把函数和导数分别处理。由于分量只有有限个，对
\(u_j,\partial_1u_j,\ldots,\partial_du_j\) 逐次抽取子列，就能得到一个共同子列和函数 \(v_0,v_1,\ldots,v_d\in L^p(\Omega)\)，使各分量分别弱收敛到这些函数。此时还不能把 \(v_i\) 直接记成 \(\partial_i v_0\)：它们最初只是不同 \(L^p\) 空间中的极限，导数关系需要由测试函数确认。

在两个端点，上述自反性理由不再适用。当 \(p=1\) 时，有界导数的积分可能集中；\(p=\infty\) 时，利用 \(L^1\) 测试得到的弱-*收敛也不等于由整个连续对偶测试的弱收敛。前者的集中问题可回看\secref{b1:sec:1005}，后者的区别见\cref{b1:prop:sobolev-weak-components}后的说明。

\subsection{弱极限}
\label{b1:subsec:250302}

\begin{proposition}\label{b1:prop:sobolev-weak-limit-identification}
设 \(1<p<\infty\)、\(u_j\in W^{1,p}(\Omega)\)，并有
\[
 u_j\rightharpoonup v_0,\quad
 \partial_i u_j\rightharpoonup v_i
 \quad\text{于 }L^p(\Omega),\quad 1\le i\le d.
\]
则 \(v_0\in W^{1,p}(\Omega)\)、\(\partial_i v_0=v_i\)，且
\(u_j\rightharpoonup v_0\) 于 \(W^{1,p}(\Omega)\)。
若另有固定的 \(g\in W^{1,p}(\Omega)\) 及范数闭线性子空间
\(M\subset W^{1,p}(\Omega)\)，并且所有 \(u_j\in g+M\)，则 \(v_0\in g+M\)。
\end{proposition}
\begin{proof}
对任意 \(\varphi\in C_c^\infty(\Omega)\)，函数 \(\varphi,\partial_i\varphi\) 都属于 \(L^{p'}(\Omega)\)。将它们代入两个弱极限，得到
\[
 \int_\Omega v_0\partial_i\varphi\,\mathrm dx
 =\lim_{j\to\infty}\int_\Omega u_j\partial_i\varphi\,\mathrm dx
 =-\lim_{j\to\infty}\int_\Omega\partial_i u_j\varphi\,\mathrm dx
 =-\int_\Omega v_i\varphi\,\mathrm dx.
\]
这正是 \(v_i\) 为 \(v_0\) 的弱偏导数的定义。由\cref{b1:prop:sobolev-weak-components}，全部分量的弱收敛给出 Sobolev 空间中的弱收敛。

仿射集 \(g+M\) 范数闭且凸，故由\cref{b1:thm:convex-weak-closure}也是弱闭的。因此弱极限仍属于这个集合。
\end{proof}

取 \(M=W_0^{1,p}(\Omega)\)，便知边界约束
\[
 u_j-g\in W_0^{1,p}(\Omega)
\]
能够随弱极限保留。这里使用的是\cref{b1:def:sobolev-zero-boundary}中的范数闭包，不需要给任意粗糙边界上的函数指定逐点值。在有界 Lipschitz 区域上，若再调用迹定理，才可把这一约束解释为与 \(g\) 具有相同的迹。

导数的识别只用了线性积分关系。若近似方程中还出现
\(|\nabla u_j|^{p-2}\nabla u_j\) 等非线性表达式，不能仅凭
\(\nabla u_j\rightharpoonup\nabla u\) 就把它的弱极限写成对应的极限表达式。后面的变分论证将先比较能量，再从极小性质推导方程。

\subsection{强收敛的恢复}
\label{b1:subsec:250303}

已有弱极限时，紧性不只是再给出一条强收敛子列。它能确定像序列的全部极限，从而使整列收敛。第\ref{b1:ex:09-compact-improves}题中的抽列方法在此成为使用紧嵌入的基本步骤。

\begin{proposition}\label{b1:prop:compact-weak-to-strong}
设 \(X,Y\) 为赋范空间，\(T:X\rightarrow Y\) 为紧线性算子。若
\(x_j\rightharpoonup x\) 于 \(X\)，则 \(Tx_j\to Tx\) 于 \(Y\) 的范数。
\end{proposition}
\begin{proof}
弱收敛列有界。又因每个 \(y^*\in Y^*\) 与 \(T\) 的复合都属于 \(X^*\)，有
\(Tx_j\rightharpoonup Tx\)。若范数收敛不成立，可选 \(\varepsilon>0\) 及子列，使
\[
 \|Tx_{j_\ell}-Tx\|_Y\ge\varepsilon
 \quad\text{对所有 }\ell.
\]
紧性使这条子列还有范数收敛子列，设其极限为 \(y\)。范数收敛蕴含弱收敛，弱极限唯一，故 \(y=Tx\)，与上式矛盾。
\end{proof}

例如，设 \(\Omega\) 为有界 Lipschitz 区域，且
\(u_j\rightharpoonup u\) 于 \(W^{1,p}(\Omega)\)。由\cref{b1:thm:rellich-subcritical}，当 \(p<d\) 时，对每个 \(1\le q<p^*=dp/(d-p)\)，整列 \(u_j\to u\) 于 \(L^q(\Omega)\)；当 \(p=d\) 时，这对每个有限 \(q\) 成立。若 \(p>d\)，则\cref{b1:thm:rellich-holder}使连续代表在
\(C(\overline\Omega)\) 以及每个 \(0<\beta<1-d/p\) 对应的
\(C^{0,\beta}(\overline\Omega)\) 中整列收敛。这些结论没有使梯度自动强收敛。

要恢复原 Sobolev 范数的收敛，可以补充范数大小的信息。先证明所需的标量事实。

\begin{lemma}\label{b1:lem:lp-weak-norm-strong}
设 \(1<p<\infty\)。若 \(f_j\rightharpoonup f\) 于实或复 \(L^p(\mu)\)，且
\(\|f_j\|_p\to\|f\|_p\)，则 \(\|f_j-f\|_p\to0\)。
\end{lemma}
\begin{proof}
标量函数 \(a\mapsto|a|^p\) 严格凸：三角不等式的等号首先要求两个非零标量同向，而实函数 \(t\mapsto t^p\) 在非负半轴上的严格凸性又排除不同的长度。因此
\[
 D(a,b)=\frac{|a|^p+|b|^p}{2}
              -\left|\frac{a+b}{2}\right|^p
\]
非负，且仅在 \(a=b\) 时为零。固定 \(0<\varepsilon<1\)。在有限维紧集
\[
 \left\{\,(a,b):|a|^p+|b|^p=1,\ |a-b|\ge\varepsilon\,\right\}
\]
上，\(D\) 有正的最小值 \(\delta_\varepsilon\)。由齐次性，若
\(S=|a|^p+|b|^p\) 且 \(|a-b|\ge\varepsilon S^{1/p}\)，便有
\(D(a,b)\ge\delta_\varepsilon S\)。

另一方面，\((f_j+f)/2\rightharpoonup f\)，故范数弱下半连续性与三角不等式给出
\[
 \|f\|_p
 \le\liminf_j\left\|\frac{f_j+f}{2}\right\|_p
 \le\limsup_j\left\|\frac{f_j+f}{2}\right\|_p
 \le\|f\|_p.
\]
因此 \(\int D(f_j,f)\,\mathrm d\mu\to0\)。在
\(|f_j-f|<\varepsilon\bigl(|f_j|^p+|f|^p\bigr)^{1/p}\) 的部分，用该不等式控制差；在其余部分，使用刚才的亏损估计以及
\(|a-b|^p\le2^{p-1}S\)。积分后得到
\[
 \|f_j-f\|_p^p
 \le\varepsilon^p\bigl(\|f_j\|_p^p+\|f\|_p^p\bigr)
   +\frac{2^{p-1}}{\delta_\varepsilon}
       \int D(f_j,f)\,\mathrm d\mu.
\]
先让 \(j\to\infty\)，再让 \(\varepsilon\downarrow0\)，即得结论。
\end{proof}

\begin{proposition}\label{b1:prop:sobolev-weak-norm-strong}
设 \(1<p<\infty\)，且 \(u_j\rightharpoonup u\) 于 \(W^{1,p}(\Omega)\)。若本书规定的标准范数满足
\[
 \|u_j\|_{W^{1,p}(\Omega)}
 \longrightarrow\|u\|_{W^{1,p}(\Omega)},
\]
则 \(u_j\to u\) 于 \(W^{1,p}(\Omega)\)。
\end{proposition}
\begin{proof}
在 \(Z=\Omega\times\{0,\ldots,d\}\) 上取 Lebesgue 测度与有限计数测度的乘积，令
\[
 F_j(x,0)=u_j(x),\quad F_j(x,i)=\partial_i u_j(x)
 \quad(1\le i\le d),
\]
并由 \(u\) 定义 \(F\)。每个 \(L^{p'}(Z)\) 测试函数只有有限多个分量，故 Sobolev 弱收敛蕴含 \(F_j\rightharpoonup F\) 于 \(L^p(Z)\)。而标准范数恰好满足
\[
 \|F_j\|_{L^p(Z)}=\|u_j\|_{W^{1,p}(\Omega)},\quad
 \|F_j-F\|_{L^p(Z)}=\|u_j-u\|_{W^{1,p}(\Omega)}.
\]
应用前一引理即可。
\end{proof}

这里使用了标准范数的有限分量 \(p\) 次和形式。等价范数保持强、弱拓扑，却不保证“某一范数的数值收敛”对另一个范数仍然成立；不能把上一命题的附加条件换成任意等价范数而不另作证明。

在 \(\Omega=(0,2\pi)\) 上，取
\[
 u_j(x)=\frac{\sin(jx)}j,\quad u_j'(x)=\cos(jx).
\]
函数一致趋于零，由零端点及\cref{b1:thm:interval-zero-boundary}也知它们属于 \(W_0^{1,p}(\Omega)\)，但
\[
 \|u_j'\|_{L^p(0,2\pi)}^p
 =\int_0^{2\pi}|\cos x|^p\,\mathrm dx>0.
\]
同时梯度弱趋于零：对 \(\psi\in C_c^\infty(0,2\pi)\)，分部积分给出
\[
 \left|\int_0^{2\pi}\psi(x)\cos(jx)\,\mathrm dx\right|
 =\frac1j\left|\int_0^{2\pi}\psi'(x)\sin(jx)\,\mathrm dx\right|
 \le\frac{\|\psi'\|_1}{j}\longrightarrow0.
\]
再用这类函数在 \(L^{p'}\) 中的稠密性和余弦函数的统一 \(L^p\) 界，把结论推广到全部测试函数。于是 \(u_j\rightharpoonup0\) 于 \(W^{1,p}\)，却不强收敛。这个例子中的障碍是梯度振荡；函数值的强收敛没有消除导数的范数。

\subsection{变分问题中的应用}
\label{b1:subsec:250304}

下面在一个具体能量上实施\secref{b1:sec:0905}的变分直接法。零阶正项同时控制函数大小，因此不需要借助 Poincaré 不等式，也不必为区域增加边界正则性。

\begin{theorem}\label{b1:thm:sobolev-variational-minimum}
设 \(\Omega\subset\mathbb R^d\) 为有界开集，\(1<p<\infty\)，
\(g\in W^{1,p}(\Omega)\)、\(f\in L^{p'}(\Omega)\) 均为实值。则泛函
\[
 J(u)=\frac1p\int_\Omega\bigl(|\nabla u|^p+|u|^p\bigr)\,\mathrm dx
                -\int_\Omega fu\,\mathrm dx
\]
在仿射集 \(\mathcal A=g+W_0^{1,p}(\Omega)\) 上有唯一极小点 \(u\)。这里
\(|\nabla u|\) 取欧氏长度。极小点也恰是 \(\mathcal A\) 中满足下列积分恒等式的函数：
\begin{equation}\label{b1:eq:sobolev-euler-lagrange}
 \int_\Omega\left(
   |\nabla u|^{p-2}\nabla u\cdot\nabla\varphi
       +|u|^{p-2}u\varphi
 \right)\,\mathrm dx
 =\int_\Omega f\varphi\,\mathrm dx
 \quad\bigl(\varphi\in W_0^{1,p}(\Omega)\bigr).
\end{equation}
其中 \(|z|^{p-2}z\) 在 \(z=0\) 处规定为零。
\end{theorem}
\begin{proof}
记
\[
 A(u)=\int_\Omega\bigl(|\nabla u|^p+|u|^p\bigr)\,\mathrm dx.
\]
有限维范数的等价性给出 \(c_{d,p}\|u\|_{W^{1,p}}^p\le A(u)\)。当 \(p\ne2\) 时，不把这里的欧氏梯度能量与逐分量 \(p\) 次和认作数值相等。由 Hölder 不等式及带参数的 Young 不等式，
\[
 \left|\int_\Omega fu\,\mathrm dx\right|
 \le\|f\|_{p'}\cdot\|u\|_p
 \le\frac1{2p}\|u\|_p^p+C_p\|f\|_{p'}^{p'}.
\]
所以
\[
 J(u)\ge\frac1{2p}A(u)-C_p\|f\|_{p'}^{p'}
 \ge c\|u\|_{W^{1,p}}^p-C_p\|f\|_{p'}^{p'}.
\]
这既给出下界，也保证能量有统一上界的序列在 \(W^{1,p}\) 中有界。又 \(\mathcal A\) 非空且 \(J(g)\) 有限，故其下确界是实数，可以选取极小化序列。

再看弱下半连续性。映射 \(z\mapsto|z|^p\) 在欧氏空间中凸，因此 \(A\) 凸。它也在 Sobolev 范数下连续：梯度的欧氏 \(L^p\) 范数之差不超过梯度差的欧氏 \(L^p\) 范数，后者由标准 Sobolev 范数控制；零阶项相同。取 \(p\) 次幂仍保持连续。由\cref{b1:thm:convex-weak-lsc}，\(A\) 弱下半连续，而 \(u\mapsto\int fu\) 弱连续。因此 \(J\) 弱下半连续。

从极小化序列中抽出 Sobolev 弱收敛子列。由\cref{b1:prop:sobolev-weak-limit-identification}，极限仍在 \(\mathcal A\) 中；由弱下半连续性，它的能量不超过原下确界，因而取得极小值。若两个实函数代表不同的等价类，它们在正测度集合上不同，严格凸的零阶函数 \(s\mapsto|s|^p\) 使中点能量严格小于两个能量的平均。梯度项凸、线性项不影响这个严格不等式，故极小点唯一。

最后推导方程，并核实积分求导。对固定 \(\varphi\in W_0^{1,p}(\Omega)\)，所有实数 \(t\) 都满足 \(u+t\varphi\in\mathcal A\)。欧氏函数 \(z\mapsto|z|^p/p\) 在包括零点在内的每一点可微，导数为 \(|z|^{p-2}z\)。对 \(0<|t|\le1\)，沿线段积分其导数，得到
\[
 \left|\frac{|z+tw|^p-|z|^p}{pt}\right|
 \le C_p\bigl(|z|^{p-1}+|w|^{p-1}\bigr)\cdot|w|.
\]
分别取 \((z,w)=(\nabla u,\nabla\varphi)\) 与 \((u,\varphi)\)。右端由 Hölder 不等式可积，故控制收敛定理允许在两个能量积分内通过差商极限。由
\[
 \left.\frac{\mathrm d}{\mathrm dt}J(u+t\varphi)\right|_{t=0}=0
\]
即得\eqref{b1:eq:sobolev-euler-lagrange}。

反过来，若 \(u\in\mathcal A\) 满足该恒等式，对任意 \(v\in\mathcal A\) 取 \(\varphi=v-u\)。凸函数的一阶支撑不等式给出
\[
 \begin{split}
 J(v)-J(u)\ge
 \int_\Omega\bigl(
  |\nabla u|^{p-2}\nabla u\cdot\nabla(v-u)
       +|u|^{p-2}u(v-u)-f(v-u)
 \bigr)\,\mathrm dx=0.
 \end{split}
\]
因此 \(u\) 是极小点，且由已经证明的严格凸性唯一。
\end{proof}

式\eqref{b1:eq:sobolev-euler-lagrange}称为这一能量的%
\term*[euler-lagrange-fangcheng]{Euler–Lagrange 方程}[Euler–Lagrange equation]的弱形式。它在推导时只要求一阶弱导数，没有预先假定极小点具有二阶经典导数。

这个存在性证明没有使用 Rellich 紧性：自反性给出弱子列，闭仿射约束保留边界条件，凸性给出能量的弱下半连续性。若再加入关于 \(u\) 的非凸低阶项，弱下半连续性可能失效，此时紧嵌入提供的 \(L^q\) 强收敛才可能用来控制该项，还须核对它的连续性和增长条件。不能仅因已经抽取了弱子列，就宣称全部非线性项都能通过极限。

希望把这些步骤放在 Sobolev 空间与偏微分方程的系统课程中学习，可伴读 Brezis 的《Functional Analysis, Sobolev Spaces and Partial Differential Equations》\cite{Brezis2011Functional}。本节的极小值论证也说明，使用紧性之前先辨认方程中究竟需要哪一种收敛，往往能够减少不必要的几何假设。


\begin{chaptersummary}
有限指数 \(L^p\) 的强紧性可以落实为一个具体逼近过程：先用小平移控制网格内的变化，再用尾部控制把无穷多个网格缩减为有限多个。有界弱梯度提供了整族的一致平移模；有界区域及适当延拓提供尾部控制，由此得到最基本的 Sobolev 紧嵌入。更高积分指数和较低 Hölder 指数的紧性，通过连续估计与插值进一步获得。

边界规则性不是所有版本都需要的条件。一般函数可借 Lipschitz 延拓，零边界闭包可以补零，局部问题可以先截断；高阶问题则必须有真正控制相应导数的延拓。临界指数处不再保留插值余量，集中例子说明这会实质性地破坏紧性。

弱收敛保留线性导数关系与闭凸约束，紧性能够加强像空间中的收敛，而原 Sobolev 范数的强收敛还需要额外信息。凸能量的极小值问题可以仅用弱下半连续性；需要保留单位范数等非弱闭约束，或处理合适的非凸低阶项时，强紧性才承担不可替代的任务。下一章进一步研究等价类在多小的例外集外能够确定点值，这与本章的积分范数紧性是另一层问题。
\end{chaptersummary}

\BookExercises

% \chapref{b1:ch:25}习题临时稿；不另设 BookExercises 入口。
% 八题均按本章及前章已建立的工具自拟，来源与实质增量逐题说明，不冒认教材原题号。
% 不使用容量、谱完备性或尚未建立的偏微分方程正则性。

\begin{exercise}\label{b1:ex:25-fk-redundant-boundedness}
设 \(1\le p<\infty\)，\(\mathcal F\subset L^p(\mathbb R^d)\) 满足
\[
 \lim_{R\to\infty}\sup_{f\in\mathcal F}
       \|f\|_{L^p(\{|x|>R\})}=0,
 \quad
 \lim_{\delta\downarrow0}\sup_{f\in\mathcal F}
       \sup_{|h|<\delta}\|\tau_hf-f\|_p=0,
\]
其中 \(\tau_hf(x)=f(x-h)\)。证明 \(\mathcal F\) 的 \(L^p\) 范数自动一致有界。
因此，Fréchet–Kolmogorov 判据虽然可以列成三个条件，却不能据此断言三个条件在逻辑上相互独立。
\end{exercise}
% 来源：自拟推导；用长位移的有限分步把球内质量与球外尾部相比较，不借用未实读论文的证明。
\begin{solution}
取 \(R\ge1\)，使每个 \(f\in\mathcal F\) 的球外范数不超过 \(1\)，再取
\(\delta>0\)，使所有 \(|h|<\delta\) 的平移误差都不超过 \(1\)。固定
\(a=3Re_1\)，选一个正整数 \(N\)，使 \(|a|/N<\delta\)，并记 \(h=a/N\)。
平移等距性与望远镜求和给出
\[
 \tau_af-f=\sum_{j=0}^{N-1}\tau_{jh}(\tau_hf-f),
 \quad
 \|\tau_af-f\|_p\le N.
\]
对 \(x\in B(0,R)\)，有 \(|x-a|>R\)，因此
\[
 \|\tau_af\|_{L^p(B(0,R))}
   =\|f\|_{L^p(B(0,R)-a)}\le1.
\]
于是
\[
 \begin{aligned}
 \|f\|_p
 &\le\|f\|_{L^p(B(0,R))}
              +\|f\|_{L^p(\{|x|>R\})}\\
 &\le\|f-\tau_af\|_p
        +\|\tau_af\|_{L^p(B(0,R))}+1
 \le N+2.
 \end{aligned}
\]
球面为零测集，开闭球的差别不影响这些范数。数 \(N\) 只由两个统一条件选出的
\(R,\delta\) 决定，不依赖 \(f\)，故得到一致界。再结合定理%
\ref{b1:thm:frechet-kolmogorov-lp}，这两个条件本身已经足以推出相对紧性。
\end{solution}

\begin{exercise}\label{b1:ex:25-weighted-confinement}
设 \(1\le p<\infty\)、\(a>0\)、\(0\le M<\infty\)。证明集合
\[
 \mathcal K=\left\{u\in W^{1,p}(\mathbb R^d):
   \|u\|_p^p+\|\nabla u\|_p^p+
        \int_{\mathbb R^d}|x|^a\cdot|u(x)|^p\,\mathrm dx
                               \le M^p\right\}
\]
在 \(L^p(\mathbb R^d)\) 中相对紧。解释最后的加权项如何替代有界区域或共同紧支集的条件。
\end{exercise}
% 来源：自拟；全空间的矩控制给出尾部，弱导数给出一致平移控制，组合为非有界区域上的紧性。
\begin{solution}
集合中的函数首先满足 \(\|u\|_p\le M\)。对 \(R>0\)，由
\(|x|^a\ge R^a\) 在球外成立，得到
\[
 \int_{|x|>R}|u(x)|^p\,\mathrm dx
 \le R^{-a}\int_{\mathbb R^d}|x|^a\cdot|u(x)|^p\,\mathrm dx
 \le M^pR^{-a}.
\]
因此尾部一致趋于零。

还需控制平移。若 \(u\in C_c^\infty(\mathbb R^d)\)，线段积分给出
\[
 u(x-h)-u(x)
       =-\int_0^1\nabla u(x-th)\cdot h\,\mathrm dt,
\]
从而 Minkowski 积分不等式及平移等距性推出
\[
 \|\tau_hu-u\|_p\le|h|\cdot\|\nabla u\|_p.
\]
由\cref{b1:thm:whole-space-sobolev-density}，对一般
\(u\in W^{1,p}(\mathbb R^d)\) 作光滑紧支集逼近；等式两侧所需的
\(L^p\) 范数都可通过极限，上式仍成立。因此对 \(\mathcal K\) 有
\[
 \sup_{u\in\mathcal K}\|\tau_hu-u\|_p\le M|h|\longrightarrow0.
\]
定理\ref{b1:thm:frechet-kolmogorov-lp}的三个条件均已满足，所求相对紧性随之成立。

加权项并不要求任何一个函数具有紧支集；它要求所有函数在远处放置的质量都受到同一个界限制。导数界约束局部振荡，而这个额外条件阻止质量整体移向无穷远，两者承担不同的任务。
\end{solution}

\begin{exercise}\label{b1:ex:25-critical-mass-concentration}
设 \(d\ge2\)、\(1\le p<d\)、\(p^*=dp/(d-p)\)，取非零
\(\varphi\in C_c^\infty(B(0,1))\)。对 \(0<\varepsilon<1\)，令
\[
 u_\varepsilon(x)=\varepsilon^{1-d/p}\varphi(x/\varepsilon),
 \quad x\in\mathbb R^d.
\]
\begin{enumerate}
 \item\label{b1:item:25-critical-no-strong-subsequence}
 证明 \(u_\varepsilon\) 在 \(W_0^{1,p}(B(0,1))\) 中有界，在每个
 \(1\le q<p^*\) 的 \(L^q\) 中趋于零；但对于任意 \(\varepsilon_n\downarrow0\)，
 \((u_{\varepsilon_n})\) 都没有 \(L^{p^*}\) 强收敛子列。
 \item\label{b1:item:25-critical-defect-measures}
 把
 \[
  \mu_\varepsilon=|u_\varepsilon|^{p^*}\,\mathrm dx,\quad
  \nu_\varepsilon=|\nabla u_\varepsilon|^p\,\mathrm dx
 \]
 看成全空间上的有限正 Borel 测度。证明
 \[
  \mu_\varepsilon\Rightarrow
      \|\varphi\|_{p^*}^{p^*}\delta_0,
  \quad
  \nu_\varepsilon\Rightarrow
      \|\nabla\varphi\|_p^p\delta_0.
 \]
\end{enumerate}
\end{exercise}
% 来源：自拟；临界重缩放后进一步计算有限测度的弱极限，不止于前章排除过强范数估计。
\begin{solution}
每个函数光滑且具有包含在单位球内部的紧支集。变量代换给出
\[
 \|\nabla u_\varepsilon\|_p=\|\nabla\varphi\|_p,
 \quad
 \|u_\varepsilon\|_q
  =\varepsilon^{1-d/p+d/q}\|\varphi\|_q
                      \quad(1\le q<\infty).
\]
特别地，\(\|u_\varepsilon\|_p=\varepsilon\|\varphi\|_p\)，所以
\(W^{1,p}\) 范数一致有界。当 \(q<p^*\) 时，指数
\(1-d/p+d/q\) 为正，得到强 \(L^q\) 收敛；在 \(q=p^*\) 时，这个指数等于零，因而
\[
 \|u_\varepsilon\|_{p^*}=\|\varphi\|_{p^*}>0.
\]
对每个 \(x\ne0\)，当 \(\varepsilon\) 足够小时，\(x\) 已在
\(\varepsilon\operatorname{supp}\varphi\) 之外，所以 \(u_\varepsilon(x)=0\)。
假如某个子列在 \(L^{p^*}\) 中强收敛，则它还能抽出几乎处处收敛的子列；上述逐点结论使其极限只能为零。但强收敛应使范数趋于零，与固定的正范数矛盾。这排除了全部强收敛子列。

对第二问，任取 \(\psi\in C_b(\mathbb R^d)\)。临界指数恰好消去变量代换产生的尺度因子，故
\[
 \int\psi\,\mathrm d\mu_\varepsilon
    =\int_{\mathbb R^d}\psi(\varepsilon y)|\varphi(y)|^{p^*}\,\mathrm dy
    \longrightarrow\psi(0)\|\varphi\|_{p^*}^{p^*}.
\]
控制收敛适用，因为 \(\psi\) 有界，而 \(|\varphi|^{p^*}\) 可积。同样，
\[
 \int\psi\,\mathrm d\nu_\varepsilon
    =\int_{\mathbb R^d}\psi(\varepsilon y)|\nabla\varphi(y)|^p\,\mathrm dy
    \longrightarrow\psi(0)\|\nabla\varphi\|_p^p.
\]
这就是有限测度弱收敛的定义。函数在次临界范数中消失，并不意味着这些临界总量消失；它们集中到了一个点上。
\end{solution}

\begin{exercise}\label{b1:ex:25-holder-endpoint}
设 \(d<p<\infty\)、\(\alpha=1-d/p\)，取
\(\varphi\in C_c^\infty(B(0,1/4))\) 且 \(\varphi(0)=1\)。对
\(0<\varepsilon<1\)，在单位球 \(B\) 上令
\[
 u_\varepsilon(x)=\varepsilon^\alpha\varphi(x/\varepsilon).
\]
证明这是一族有界的 \(W_0^{1,p}(B)\) 函数，并且对每个
\(0<\beta<\alpha\)，有 \(u_\varepsilon\to0\) 于
\(C^{0,\beta}(\overline B)\)。另一方面，证明对任何
\(\varepsilon_n\downarrow0\)，它都没有 \(C^{0,\alpha}(\overline B)\) 强收敛子列。
由此区分 Morrey 的连续嵌入指数与 Rellich 的紧嵌入指数。
\end{exercise}
% 来源：自拟；在可用的临界 Hölder 估计内证明失去相对紧性，不再讨论更高指数是否有统一界。
\begin{solution}
与上一题相同的变量代换给出
\[
 \|u_\varepsilon\|_p=\varepsilon\|\varphi\|_p,
 \quad
 \|\nabla u_\varepsilon\|_p=\|\nabla\varphi\|_p,
 \quad
 \|u_\varepsilon\|_\infty=\varepsilon^\alpha\|\varphi\|_\infty.
\]
函数均光滑且内部紧支集，所以属于 \(W_0^{1,p}(B)\)，并在其中有界；最后一式还给出一致趋零。

把 \(\varphi\) 视为全空间光滑紧支集函数。它的全空间
\(\beta\)-Hölder 半范数有限：距离不超过 \(1\) 时用梯度上界，距离超过 \(1\) 时用两倍上确界。对任意 \(x,y\in B\) 作缩放比较，得到
\[
 [u_\varepsilon]_{0,\beta;\overline B}
     \le\varepsilon^{\alpha-\beta}
                   [\varphi]_{0,\beta;\mathbb R^d}.
\]
若 \(\beta<\alpha\)，这个半范数也趋于零，因而得到所求的强
\(C^{0,\beta}\) 收敛。

但取 \(a=e_1/2\)，有 \(\varphi(a)=0\)。在 \(0,\varepsilon a\) 两点比较，
\[
 [u_\varepsilon]_{0,\alpha;\overline B}
 \ge\frac{|u_\varepsilon(0)-u_\varepsilon(\varepsilon a)|}
             {|\varepsilon a|^\alpha}
 =2^\alpha.
\]
若某个子列在 \(C^{0,\alpha}\) 中强收敛，它也应一致收敛；唯一可能的极限是零。
然而强 \(C^{0,\alpha}\) 收敛到零又要求上述半范数趋于零，矛盾。
所以临界指数处有连续嵌入，并不足以推出同一目标空间中的相对紧性。
\end{solution}

\begin{exercise}\label{b1:ex:25-infinitely-many-components}
设 \(d\ge1\)、\(1\le p<\infty\)，令
\[
 x_n=2^{-n}e_1,\quad r_n=2^{-n-3},\quad
 B_n=B(x_n,r_n),\quad
 \Omega=\bigcup_{n=1}^\infty B_n.
\]
证明 \(\Omega\) 是有界开集，但嵌入
\(W^{1,p}(\Omega)\to L^p(\Omega)\) 不紧。具体地，考虑
\[
 u_n=|B_n|^{-1/p}\mathbf1_{B_n}
                  \quad\text{作为 }\Omega\text{ 上的函数},
\]
计算它们的弱导数、Sobolev 范数和两两 \(L^p\) 距离。
解释这为什么不反驳任意有界开集上的 \(W_0^{1,p}\) 紧嵌入。
\end{exercise}
% 来源：自拟；无限分量上的归一化常数，区别无边界条件空间与零边界闭包。
\begin{solution}
相邻两个中心的距离为 \(2^{-n-1}\)，而半径之和为
\(3\cdot2^{-n-4}\)，前者严格较大；不相邻的球距离更大。因此各球闭包两两不交。
所有球都包含在单位球内，故 \(\Omega\) 是有界开集。中心列趋于原点，而原点不属于 \(\Omega\)。

函数 \(u_n\) 在每个分量内都是常数，因此其弱梯度在 \(\Omega\) 内为零。也可以直接检验：若 \(\phi\in C_c^\infty(\Omega)\)，则
\(\phi|_{B_n}\) 的支集与 \(\partial B_n\) 保持正距离，故对每个坐标 \(i\)，
\[
 \int_\Omega u_n\partial_i\phi
      =|B_n|^{-1/p}\int_{B_n}\partial_i\phi=0.
\]
这里是在 \(\Omega\) 内求弱导数，没有把 \(u_n\) 补零后跨越分量边界求导。
由归一化可得
\[
 \|u_n\|_{L^p(\Omega)}=1,\quad
 \nabla u_n=0,\quad
 \|u_n\|_{W^{1,p}(\Omega)}=1.
\]
当 \(n\ne m\) 时，两个函数支集不交，因此
\[
 \|u_n-u_m\|_{L^p(\Omega)}^p
          =\|u_n\|_p^p+\|u_m\|_p^p=2.
\]
这是一列有界 Sobolev 函数，却没有 \(L^p\) Cauchy 子列，故嵌入不紧。

最后，任何一个 \(u_n\) 都不属于 \(W_0^{1,p}(\Omega)\)。否则定理%
\ref{b1:thm:zero-boundary-poincare}会使
\(\|u_n\|_p\le C\|\nabla u_n\|_p=0\)，与其范数等于 \(1\) 矛盾。
因此这个例子检验的是不加边界条件且不加边界规则性时的失败，不涉及零边界版本的反例。
\end{solution}

\begin{exercise}\label{b1:ex:25-nonlinear-reaction-limit}
设 \(\Omega\subset\mathbb R^d\) 为有界 Lipschitz 开集，
\(1\le p<d\)、\(1<q<p^*=dp/(d-p)\)，并令 \(q'=q/(q-1)\)。
设实值函数 \(u_n\rightharpoonup u\) 于 \(W^{1,p}(\Omega)\)。
\begin{enumerate}
 \item\label{b1:item:25-reaction-strong-limit}
 证明
 \[
  |u_n|^{q-2}u_n\longrightarrow |u|^{q-2}u
                       \quad\text{于 }L^{q'}(\Omega),
 \]
 其中在函数值为零处规定相应表达式为零。分别给出 \(q\ge2\) 和
 \(1<q<2\) 所需的点态估计。
 \item\label{b1:item:25-reaction-strong-weak-pairing}
 若另有 \(v_n\rightharpoonup v\) 于 \(L^q(\Omega)\)，证明
 \[
  \int_\Omega |u_n|^{q-2}u_n v_n
          \longrightarrow\int_\Omega |u|^{q-2}u v.
 \]
\end{enumerate}
\end{exercise}
% 来源：自拟；次临界紧嵌入与幂次映射的定量连续性，再组合强弱配对，不把非线性操作直接套在弱极限上。
\begin{solution}
由\cref{b1:thm:rellich-subcritical}，嵌入
\(W^{1,p}(\Omega)\to L^q(\Omega)\) 是紧线性算子。命题%
\ref{b1:prop:compact-weak-to-strong}因而给出整列
\(u_n\to u\) 于 \(L^q(\Omega)\)，不需要再为不同测试函数分别抽子列。
记 \(F(t)=|t|^{q-2}t\)。首先
\[
 \|F(w)\|_{q'}=\|w\|_q^{q-1},
\]
所以题中的非线性项都属于所需对偶指数空间。

当 \(q>2\) 时，由 \(F'(t)=(q-1)|t|^{q-2}\) 及一维积分公式，
\[
 |F(a)-F(b)|
       \le(q-1)(|a|+|b|)^{q-2}\cdot|a-b|.
\]
以指数 \(q/(q-2)\) 和 \(q\) 对乘积使用 Hölder 不等式，得到
\[
 \|F(u_n)-F(u)\|_{q'}
 \le(q-1)(\|u_n\|_q+\|u\|_q)^{q-2}
                              \cdot\|u_n-u\|_q
 \longrightarrow0.
\]
在 \(q=2\) 时，\(F(t)=t\)，结论直接就是 \(L^2\) 强收敛。

当 \(1<q<2\) 时，记 \(\theta=q-1\in(0,1)\)。同号的 \(a,b\) 可用
\(t^\theta\) 的次可加性得到
\(|F(a)-F(b)|\le|a-b|^\theta\)；异号时，由凹性有
\[
 |F(a)-F(b)|
   =|a|^\theta+|b|^\theta
   \le2^{1-\theta}(|a|+|b|)^\theta
   =2^{1-\theta}|a-b|^\theta.
\]
因而对所有实数都有
\[
 |F(a)-F(b)|\le2^{2-q}|a-b|^{q-1},
 \quad
 \|F(u_n)-F(u)\|_{q'}
                  \le2^{2-q}\|u_n-u\|_q^{q-1}\longrightarrow0.
\]
这完成第一问。

对于第二问，弱收敛列 \((v_n)\) 的 \(L^q\) 范数一致有界。分解为
\[
 \int_\Omega F(u_n)v_n-\int_\Omega F(u)v
   =\int_\Omega\bigl(F(u_n)-F(u)\bigr)v_n
                       +\int_\Omega F(u)(v_n-v).
\]
第一项的绝对值不超过
\(\|F(u_n)-F(u)\|_{q'}\cdot\|v_n\|_q\)，所以趋于零；第二项则因
\(F(u)\in L^{q'}\) 是一个固定的弱收敛测试函数而趋于零。
这里先由紧性恢复了 \(u_n\) 的强收敛，才有资格把非线性项传到极限。
\end{solution}

\begin{exercise}\label{b1:ex:25-first-dirichlet-eigenvalue}
设 \(\Omega\subset\mathbb R^d\) 为非空有界开集，本题取实值函数，并定义
\[
 \lambda_1=\inf\left\{
       \int_\Omega|\nabla v|^2:
       v\in H_0^1(\Omega),\ \|v\|_{L^2(\Omega)}=1
                    \right\}.
\]
\begin{enumerate}
 \item\label{b1:item:25-first-eigenvalue-attainment}
 证明 \(0<\lambda_1<\infty\)，且存在 \(u\in H_0^1(\Omega)\) 取得这个下确界。
 \item\label{b1:item:25-first-eigenvalue-equation}
 不借助抽象乘子定理，沿保持约束的曲线证明
 \[
  \int_\Omega\nabla u\cdot\nabla\phi
       =\lambda_1\int_\Omega u\phi
              \quad\bigl(\phi\in H_0^1(\Omega)\bigr).
 \]
 若非零 \(w\in H_0^1(\Omega)\) 也满足这个恒等式，但参数换成实数
 \(\lambda\)，证明 \(\lambda\ge\lambda_1\)。
\end{enumerate}
本题的方程按弱导数解释，不要求证明特征函数的经典光滑性或全部特征函数的完备性。
\end{exercise}
% 来源：自拟；L2 单位球这一非凸约束下的直接法与归一化变分，不重复正文带正零阶项的无约束凸能量问题。
\begin{solution}
由\cref{b1:thm:zero-boundary-poincare}，存在 \(C_P<\infty\)，使
\(\|v\|_2\le C_P\|\nabla v\|_2\) 对所有 \(v\in H_0^1(\Omega)\) 成立。
于是约束集中的每个函数都有梯度能量至少 \(C_P^{-2}\)，故
\(\lambda_1>0\)。又因为非空开集包含一个球，可在其中取非零光滑紧支集函数并归一化，得到约束集非空且 \(\lambda_1<\infty\)。

取极小化序列 \(u_n\)，满足 \(\|u_n\|_2=1\) 及
\(\|\nabla u_n\|_2^2\to\lambda_1\)。这列函数在 Hilbert 空间
\(H_0^1(\Omega)\) 中有界。由\cref{b1:thm:reflexive-weak-subsequence}，可以抽出一个子列在该空间中弱收敛到某个 \(u\)。
定理\ref{b1:thm:rellich-subcritical}与\ref{b1:thm:rellich-holder}的零边界版本给出
\(H_0^1(\Omega)\to L^2(\Omega)\) 的紧性：在 \(d>2\) 时
\(2<2d/(d-2)\)，在 \(d=2\) 时使用临界的有限指数结论；在 \(d=1\) 时，紧 Hölder 嵌入先给出一致收敛，再由区域有界性推出 \(L^2\) 收敛。
结合命题\ref{b1:prop:compact-weak-to-strong}，该子列在 \(L^2\) 中强收敛到
\(u\)，所以 \(\|u\|_2=1\)。

梯度范数的弱下半连续性给出
\[
 \|\nabla u\|_2^2
      \le\liminf_n\|\nabla u_n\|_2^2=\lambda_1.
\]
由于 \(u\) 已满足约束，下确界的定义又给出反向不等式，因此
\(\|\nabla u\|_2^2=\lambda_1\)。强 \(L^2\) 收敛在这里保住了单位范数约束；仅有弱收敛并不能保证这一点。

对任意 \(\phi\in H_0^1(\Omega)\)，令
\[
 v_t=\frac{u+t\phi}{\|u+t\phi\|_2}
\]
并取充分小的实数 \(t\)，使分母非零。这个曲线始终属于约束集。
记 \(A=\int_\Omega\nabla u\cdot\nabla\phi\)、\(B=\int_\Omega u\phi\)，则
\[
 \int_\Omega|\nabla v_t|^2
 =\frac{\lambda_1+2tA+t^2\|\nabla\phi\|_2^2}
        {1+2tB+t^2\|\phi\|_2^2}.
\]
左侧在 \(t=0\) 取得极小值，故其导数为零，即
\(2A-2\lambda_1B=0\)。这正是所求恒等式。特别地，对光滑紧支集测试函数成立，所以它表示分布意义的
\(-\Delta u=\lambda_1u\)，零边界条件则由 \(u\in H_0^1(\Omega)\) 表达。

最后，若非零 \(w\) 满足参数为 \(\lambda\) 的同类弱方程，取测试函数
\(\phi=w\)，便有
\[
 \lambda=\frac{\|\nabla w\|_2^2}{\|w\|_2^2}\ge\lambda_1.
\]
因此取得的参数确实是所有这种弱 Dirichlet 特征值中的最小值；这个结论不需要预先建立谱分解。
\end{solution}

\begin{exercise}\label{b1:ex:25-recover-weak-limit}
设 \(X\) 为自反 Banach 空间，\(Y\) 为 Banach 空间，
\(T:X\to Y\) 为有界线性单射。设 \((x_n)\) 在 \(X\) 中有界，且
\(Tx_n\to y\) 于 \(Y\) 的范数。
\begin{enumerate}
 \item\label{b1:item:25-reflexive-recovery}
 证明存在唯一 \(x\in X\)，使 \(Tx=y\)，并且整列
 \(x_n\rightharpoonup x\) 于 \(X\)。注意本方向没有假设 \(T\) 为紧算子。
 \item\label{b1:item:25-nonreflexive-recovery-failure}
 证明去掉 \(X\) 的自反性后，结论即使对紧线性单射也会失败。具体检验
 \[
  T:\ell^1\longrightarrow\ell^2,\quad
  T(a_1,a_2,\ldots)=(a_1,a_2/2,a_3/3,\ldots)
 \]
 及 \(\ell^1\) 的标准单位向量列。
\end{enumerate}
\end{exercise}
% 来源：自拟；由较弱空间中的强极限识别原空间弱极限，配合紧对角算子展示自反性不可由紧性替代。
\begin{solution}
由\cref{b1:thm:reflexive-weak-subsequence}，有界列可以抽出弱收敛子列
\(x_{n_j}\rightharpoonup x\)。有界线性算子保持弱收敛，因为对每个
\(f\in Y^*\)，\(f\circ T\in X^*\)，所以 \(Tx_{n_j}\rightharpoonup Tx\)。
另一方面，它已在范数中趋于 \(y\)，故也弱收敛到 \(y\)。弱极限唯一，得到
\(Tx=y\)。单射性使这样的 \(x\) 至多有一个。

还须从一个子列推进到整列。假设整列不弱收敛到 \(x\)，则存在
\(g\in X^*\)、\(\varepsilon>0\) 及一个子列，满足
\[
 |g(x_{n_j})-g(x)|\ge\varepsilon\quad\text{对每个 }j.
\]
这个子列仍然有界，再用自反性抽出弱收敛子列，记其弱极限为 \(z\)。
按前面的论证，必有 \(Tz=y\)，从而单射性给出 \(z=x\)。这与所列标量距离始终不小于 \(\varepsilon\) 矛盾，因此整列确实弱收敛。

对第二问，所给对角算子线性且单射，并满足
\(\|Ta\|_2\le\|a\|_1\)。记 \(T_Na\) 为只保留前 \(N\) 个坐标的部分，则
\[
 \|(T-T_N)a\|_2
       \le\frac1{N+1}\|a\|_2
       \le\frac1{N+1}\|a\|_1.
\]
对 \(\ell^1\) 单位球，前 \(N\) 个坐标的像是有限维有界集，具有有限的任意精度网；尾项又由上式一致控制。因此 \(T\) 的单位球像全有界，在完备空间
\(\ell^2\) 中相对紧，说明 \(T\) 是紧算子。

可是对标准单位向量 \(e_n\)，有
\[
 \|e_n\|_1=1,\quad \|Te_n\|_2=1/n\longrightarrow0.
\]
若原向量列弱收敛，其极限由 \(T\) 的连续性和单射性只能为零。但连续线性泛函
\[
 g(a)=\sum_{j=1}^\infty a_j,\quad |g(a)|\le\|a\|_1,
\]
始终满足 \(g(e_n)=1\)，故这个列并不弱收敛到零。紧性能够控制像空间中的收敛，不能替代原空间中抽取弱子列所需的条件。
\end{solution}

% !TeX root = ../../Book1.tex
\OptionalChapter{Sobolev 容量与精细性质}{b1:ch:26}
\BookChapterNumber{b1:ch:26}
\begin{chapterabstract}
本章研究 Sobolev 等价类能够在多小的例外集之外确定点值。以邻域上的最小能量定义容量后，证明它的次可加性，构造拟连续代表，并建立拟处处唯一性及范数收敛的逐点子列结论。容量与 Hausdorff 测度的比较进一步给出异常集的维数限制。最后介绍细拓扑及 Wiener 判据。基础容量性质在正文中证明；涉及非线性势论的深层等价放在选读部分，只给出证明概要。
\end{chapterabstract}

\begin{learninggoals}
\begin{enumerate}
\item\label{b1:goal:26-capacity}从开邻域上的几乎处处约束定义非齐次容量，证明单调性、可列次可加性，并辨别零测与零容量。
\item\label{b1:goal:26-representative}用快速光滑逼近构造拟连续代表，证明超水平容量估计及拟处处唯一性，不用尚未建立的点值概念反过来定义容量。
\item\label{b1:goal:26-dimension}通过小球覆盖、弱紧性与平均振荡估计，把容量零集与 Hausdorff 维数联系起来，单独处理临界指数和对数尺度。
\item\label{b1:goal:26-fine}区分全空间非齐次容量与相对纯梯度容量，理解薄性和细拓扑的定义，以及 Wiener 边界判据还需要哪些势论工具。
\item\label{b1:goal:26-boundaries}辨认拟连续、通常连续与 Lebesgue 精确代表之间已证的关系，避免把几乎处处结论未经证明提升为拟处处结论。
\end{enumerate}
\end{learninggoals}

前几章在积分范数中讨论弱导数、逼近和紧性，函数始终先作为几乎处处等价类出现。第24章在 \(p>d\) 时已经得到唯一连续代表，但当 \(p\le d\) 时，普通连续性不再普遍成立。另一方面，完全放弃点值又会掩盖零边界条件、薄集合和异常集的几何。容量提供了介于这两种处理之间的方式：不要求每个点都被确定，却要求例外集比任意 Lebesgue 零测集更小。

\ProseParagraphBreak

本章属于选读。主线固定 \(1<p<\infty\)，因为次可加性和若干极限步骤使用自反性；端点 \(p=1\) 需要另一套细节，本章不作讨论。前三节主要依赖第9章的 Mazur 定理与凸泛函弱下半连续性、第13章的 \(5r\) 覆盖、第15章的 Hausdorff 测度、第16章和第21章的一维绝对连续理论与 ACL 代表，以及第24章的 Poincaré 和 Morrey 不等式。弱导数、光滑稠密性与 Sobolev 弱紧性仍是贯穿其中的基础。

\ProseParagraphBreak

最后一节把这些结果放入非线性势论的图景。Wiener 判据、Cartan 性质及超调和函数刻画均作为外部输入，并注明出处；它们不由前三节推出。

\ProseParagraphBreak

Kinnunen 与 Martio 的《The Sobolev capacity on metric spaces》\cite{KinnunenMartio1996Capacity}适合对照容量性质和快速逼近的组织；其中的度量空间范数与这里的经典能量仅等价，正文会指出具体差别。系统的欧氏精细性质可继续参阅 Ziemer 的《Weakly Differentiable Functions》\cite{Ziemer1989Weakly}；容量与非线性椭圆方程之间的联系，可伴读 Heinonen、Kilpeläinen 与 Martio 的《Nonlinear Potential Theory of Degenerate Elliptic Equations》\cite{Heinonen2006Nonlinear}。这些专著的范围超过本章，不要求读者在开始阅读前掌握全部势论。

% !TeX root = ../../Book1.tex
\section{\texorpdfstring{\(p\)}{p}-容量}
\label{b1:sec:2601}

一个集合的 Lebesgue 测度为零，意味着在它上面更改函数值不影响任何 \(L^p\) 等价类。但如果要求函数在这个集合的整个邻域中保持某种大小，更改就不再是免费的：过渡到邻域外部需要弱梯度，而梯度也计入 Sobolev 范数。容量用完成这种过渡所需的最小能量度量集合。

本章固定 \(d\ge1\)、\(1<p<\infty\)，容许函数先取实值。这个范围使我们可以使用已经证明的自反性；它不包括 \(p=1\) 的另行处理，也不把 \(p=\infty\) 纳入同一能量定义。记
\begin{equation}\label{b1:eq:sobolev-capacity-energy}
 E_p(u)=\int_{\mathbb R^d}
       \bigl(|u(x)|^p+|\nabla u(x)|^p\bigr)\,\mathrm dx.
\end{equation}
梯度取欧氏长度，故 \(E_p(u)^{1/p}\) 是一个旋转不变的范数，与第21章逐分量 \(p\) 次和定义的标准范数等价；当 \(p\ne2\) 时，二者通常不相等。保留零阶项，使这里的容量成为非齐次全空间容量。\secref{b1:sec:2604}的相对纯梯度容量将另行定义。

Kinnunen 与 Martio 的《The Sobolev capacity on metric spaces》\cite[Sections 1 and 3, p. 368]{KinnunenMartio1996Capacity}提供开邻域容量及其基本性质的组织方式。其主体使用度量空间中的另一种 Sobolev 范数，并明确区分它与经典弱梯度能量；本节直接使用\eqref{b1:eq:sobolev-capacity-energy}，以本书已经建立的截断、弱紧性和弱下半连续性完成证明。
% 正式论文368页(1.3)--(1.5)及371--375页3.1--3.7、4.1已由协作撰写者实读。
% 本节经典欧氏格能量路线独立展开，不把Hajlasz型梯度当作普通弱梯度。

\subsection{定义}
\label{b1:subsec:260101}

\begin{definition}\label{b1:def:sobolev-capacity}
对开集 \(G\subseteq\mathbb R^d\)，定义
\[
 C_p(G)=\inf\left\{E_p(u):
 u\in W^{1,p}(\mathbb R^d),\
 u\ge1\text{ 几乎处处于 }G\right\}.
\]
若容许类为空，规定下确界为 \(+\infty\)。对任意集合 \(A\subseteq\mathbb R^d\)，定义其 \term[sobolev-rongliang]{Sobolev 容量}[Sobolev capacity]，亦称 \(p\)-容量，为
\begin{equation}\label{b1:eq:sobolev-outer-capacity}
 C_p(A)=\inf\{C_p(G): A\subseteq G,\ G\text{ 开}\}.
\end{equation}
\end{definition}
\symidx{\(C_p(A)\)：非齐次 Sobolev \(p\)-容量}

定义分为两步。第一步只要求开集上的几乎处处不等式，因此容许性不依赖 Sobolev 代表的选择。第二步以真正包含 \(A\) 的开集作外逼近，所以即使 \(A\) 本身不可测，也仍有明确的容量。开集的单调性直接来自容许类的包含关系；因此当 \(A\) 本身开时，第二步所得数值与第一步相同。

不能把上述条件改成“选某个代表使其在 \(A\) 上至少为一”。若 \(A\) 是零测集，零函数类可以任意在 \(A\) 上改成一；这种定义会令全部零测集都没有代价，失去我们要研究的区别。也不能一开始就要求“拟处处至少为一”，因为拟处处正要依赖容量才能定义。

截断使容许函数具有统一的数值范围。以下格运算的能量公式也将保证我们能处理多个集合。

\begin{lemma}\label{b1:lem:sobolev-lattice-energy}
若 \(u,v\in W^{1,p}(\mathbb R^d)\) 为实值，则 \(u\vee v=\max(u,v)\)、\(u\wedge v=\min(u,v)\) 都属于该空间，且
\begin{equation}\label{b1:eq:sobolev-lattice-energy}
 E_p(u\vee v)+E_p(u\wedge v)=E_p(u)+E_p(v).
\end{equation}
此外，\(T(u)=\min(1,\max(0,u))\) 满足 \(E_p(T(u))\le E_p(u)\)。若 \(u\) 在开集 \(G\) 上容许，则 \(T(u)\) 仍容许，且 \(0\le T(u)\le1\)，在 \(G\) 上几乎处处等于一。
\end{lemma}
\begin{proof}
由第21章的截断公式，\(u\vee v=v+(u-v)_+\) 属于 Sobolev 空间。在 \(\{u>v\}\) 上，它的弱梯度为 \(\nabla u\)；在 \(\{u<v\}\) 上为 \(\nabla v\)。在相等集合上，\cref{b1:prop:sobolev-gradient-on-level-sets}给出 \(\nabla(u-v)=0\) 几乎处处，两个选择也一致。最小值的公式相反。

于是对几乎所有点，最大值与最小值只是重新排列了两个函数值和两个梯度向量。分别取绝对值或欧氏长度的 \(p\) 次幂后相加，再积分，就得\eqref{b1:eq:sobolev-lattice-energy}。截断的值不增大绝对值，梯度只在 \(0<u<1\) 上保留原梯度，故能量不增。最后的容许性由截断在 \([1,\infty)\) 上恒等于一立即得到。
\end{proof}

\subsection{单调性}
\label{b1:subsec:260102}

\begin{proposition}\label{b1:prop:sobolev-capacity-basic}
容量具有以下性质：
\begin{enumerate}
\item \(C_p(\varnothing)=0\)，且 \(A\subseteq B\) 蕴含 \(C_p(A)\le C_p(B)\)。
\item 容量满足\eqref{b1:eq:sobolev-outer-capacity}的开外正则性，并且
\[
 m_d^*(A)\le C_p(A)
 \quad\text{对任意 }A\subseteq\mathbb R^d.
\]
\item 每个有界集具有有限容量。平移或正交变换不改变容量。
\end{enumerate}
\end{proposition}
\begin{proof}
零函数在空集上容许，所以空集容量为零。开集之间的单调性已由容许类包含得到；再对开邻域取下确界，便得任意集合的单调性。

若 \(A\subseteq G\) 且 \(u\) 在开集 \(G\) 上容许，则
\[
 m_d^*(A)\le m_d(G)\le\int_G|u|^p
 \le E_p(u).
\]
先对容许函数、再对开邻域取下确界，即得外测度下界。开外正则性本身就是第二步定义。

有界集包含在一个球内，取在该球附近等于一的光滑紧支集函数，其有限能量给出容量上界。平移、正交变换保持 Lebesgue 测度和欧氏梯度长度，把一个集合的容许函数与开邻域双向变换，就得到容量不变。
\end{proof}

零阶项使一个简单结论变得明显：无限 Lebesgue 测度的集合具有无限容量。这与只看梯度的全空间齐次问题不同。另一方面，面积为零的薄片仍可能具有正容量，稍后会给出直接证明。

若把 \(E_p\) 换成标准 Sobolev 范数的 \(p\) 次幂，由两种能量的正常数比较，所得容量也双向受常数控制。因此零容量集相同；\secref{b1:sec:2602}按任意小容量开集定义的拟连续性也相同。但对于一个具体集合，容量的数值可能改变，不能在精确常数的计算中随意更换能量。

\subsection{次可加性}
\label{b1:subsec:260103}

本节所说的可列次可加性 (countable subadditivity) 是并集容量不超过各项容量之和；下文简称次可加性。它不同于互不相交并上的可列可加等式。把两个容许函数相加虽然能覆盖两个集合，却会在能量中产生不必要的幂次常数。取最大值并使用格能量公式，可以保留系数一；处理一列集合时，再用弱紧性通过极限。

\begin{theorem}\label{b1:thm:sobolev-capacity-subadditive}
对任意集合列 \((A_j)_{j\ge1}\)，
\begin{equation}\label{b1:eq:capacity-countable-subadditivity}
 C_p\left(\bigcup_{j=1}^\infty A_j\right)
 \le\sum_{j=1}^\infty C_p(A_j).
\end{equation}
\end{theorem}
\begin{proof}
若右端无限，结论显然。以下设其有限，固定 \(\varepsilon>0\)。逐项选择开邻域 \(G_j\supseteq A_j\) 及截断后的容许函数 \(u_j\)，使
\[
 0\le u_j\le1,\qquad
 E_p(u_j)\le C_p(A_j)+\varepsilon2^{-j}.
\]
这里可在选择开集与选择容许函数时分别分配半份误差。令
\[
 v_N=\max(u_1,\ldots,u_N),\qquad v=\sup_j u_j.
\]
反复应用格能量公式可得
\[
 E_p(v_N)\le\sum_{j=1}^N E_p(u_j)
 \le\sum_jC_p(A_j)+\varepsilon.
\]
同时 \(0\le v_N\uparrow v\)，且 \(v^p\le\sum_j|u_j|^p\)，右端可积。控制收敛因此给出 \(v_N\to v\) 于 \(L^p\)。

有界能量等价于 Sobolev 范数有界。由于 \(1<p<\infty\)，可抽取 Sobolev 弱收敛子列；其零阶弱极限由已经取得的 \(L^p\) 强极限确定，必为 \(v\)。所以 \(v\in W^{1,p}\)。能量是凸且范数连续的泛函，故由第9章的凸泛函弱下半连续性，
\[
 E_p(v)\le\liminf_N E_p(v_N)
 \le\sum_jC_p(A_j)+\varepsilon.
\]
每个 \(u_j\) 在 \(G_j\) 上几乎处处等于一。合并可数个例外零集后，\(v\) 在开集 \(\bigcup_jG_j\) 上几乎处处等于一，因此是这个开并的容许函数。该开并包含 \(\bigcup_jA_j\)，故得到\eqref{b1:eq:capacity-countable-subadditivity}，最后让 \(\varepsilon\downarrow0\)。
\end{proof}

至此，\(C_p\) 具有外测度的三个公理。不过它一般不是通常意义的可加测度；不能因其定义在全部集合上，就认定所有 Borel 集自动满足 Carathéodory 可测性。第\ref{b1:ex:26-one-dimensional-exact-capacity}题将具体算出可加性的失败。我们在后续只使用已经证明的单调性、次可加性和开外正则性。

下一条估计把能量与点值联系起来，但暂时只对连续函数陈述。

\begin{lemma}\label{b1:lem:continuous-capacity-level-set}
设 \(v\in W^{1,p}(\mathbb R^d)\cap C(\mathbb R^d)\)，\(t>0\)。则
\[
 C_p(\{|v|>t\})\le t^{-p}E_p(v).
\]
\end{lemma}
\begin{proof}
连续性使超水平集开；函数 \(\min(1,|v|/t)\) 在其上几乎处处等于一，且能量不超过 \(t^{-p}E_p(v)\)，因而是所需容许函数。
\end{proof}

这不是对任意代表都成立的 Chebyshev 不等式。若只知道 \(v\) 是一个可测代表，超水平集未必开，而且在零测集上任意改变它会改变该集合的容量。下一节先构造拟连续代表，再为那一类代表证明对应估计。

\subsection{容量零集}
\label{b1:subsec:260104}

若 \(C_p(N)=0\)，则称 \(N\) 为容量零集。它的任意子集仍为容量零集，可数个容量零集的并也为容量零集；这些集合又都是 Lebesgue 零测集。为了了解这个例外集尺度到底有多细，先考察小球与单点。

\begin{proposition}\label{b1:prop:sobolev-capacity-balls}
存在只依赖 \(d,p\) 的常数，使所有 \(x\in\mathbb R^d\)、\(r>0\) 满足
\[
 C_p(B(x,r))\le C_{d,p}(r^d+r^{d-p}).
\]
\end{proposition}
\begin{proof}
固定 \(0\le\eta\le1\) 的光滑函数，等于一于 \(B(0,1)\) 的一个邻域，且支集包含于 \(B(0,2)\)。函数 \(\eta((\,\cdot-x)/r)\) 在 \(B(x,r)\) 上容许，变量代换给出其零阶能量为 \(r^d\|\eta\|_p^p\)，梯度能量为 \(r^{d-p}\|\nabla\eta\|_p^p\)。两者相加即得。
\end{proof}

\begin{proposition}\label{b1:prop:sobolev-point-capacity}
若 \(1<p\le d\)，每个单点以及每个可数集的 \(p\)-容量均为零。
若 \(p>d\)，存在 \(c_{d,p}>0\)，使每个非空集合 \(A\) 都满足 \(C_p(A)\ge c_{d,p}\)。
\end{proposition}
\begin{proof}
当 \(p<d\) 时，小球上界在 \(r\downarrow0\) 时趋零，单点结论由单调性得到，再用次可加性处理可数并。

当 \(p=d>1\) 时，普通缩放的梯度能量不趋零，需要较慢的对数过渡。固定 \(R>0\)，取 \(0<\varepsilon<R\)，定义径向函数
\[
 v_\varepsilon(x)=
 \begin{cases}
 1,&|x|\le\varepsilon,\\
 \displaystyle\frac{\log(R/|x|)}{\log(R/\varepsilon)},
       &\varepsilon<|x|<R,\\
 0,&|x|\ge R.
 \end{cases}
\]
对每个固定的 \(\varepsilon\)，这是紧支集 Lipschitz 函数，故属于 \(W^{1,d}\)。记 \(\sigma_{d-1}=d\omega_d\) 为单位球面的面积，径向积分给出
\[
 \int|\nabla v_\varepsilon|^d
 =\frac{\sigma_{d-1}}{\log(R/\varepsilon)^d}
            \int_\varepsilon^R\frac{\mathrm dr}{r}
 =\sigma_{d-1}\log(R/\varepsilon)^{1-d}\longrightarrow0.
\]
同时 \(0\le v_\varepsilon\le\mathbf1_{B(0,R)}\)，并在每个非零点趋于零，控制收敛给出其零阶能量趋零。它在原点的开邻域上容许，因此 \(C_d(\{0\})=0\)。平移及次可加性完成本情形。

最后设 \(p>d\)。由 Morrey 估计与范数等价性，每个 Sobolev 函数都有连续代表，并满足
\[
 |u^*(x)|\le C_{d,p}E_p(u)^{1/p}
 \quad(x\in\mathbb R^d).
\]
若 \(u\) 在含 \(x\) 的开集上几乎处处至少为一，连续性说明 \(u^*(x)\ge1\)：否则某个小球上都小于一，与几乎处处条件矛盾。所以每个这样的容许函数都有统一的正能量下界。两次取下确界得到单点正下界，再用单调性得到全部非空集合的下界。
\end{proof}

在 \(d\ge2\) 时，一个平面片给出更直观的零测而正容量集合。令
\[
 S=[0,1]^{d-1}\times\{0\}.
\]
若开集 \(G\) 包含 \(S\)，紧性保证存在 \(\delta>0\)，使 \([0,1]^{d-1}\times(-\delta,\delta)\subset G\)。对在 \(G\) 上容许的 \(u\)，ACL 性质与 Fubini 定理表明，对几乎所有 \(x'\in(0,1)^{d-1}\)，竖线上的一维绝对连续代表在高度零处至少为一。一维点值估计给出
\[
 1\le C_p\int_{\mathbb R}
          \bigl(|u(x',t)|^p+|\partial_du(x',t)|^p\bigr)\,\mathrm dt.
\]
为看出该估计的常数不依赖 \(\delta\)，可任取 \(s\in(0,1)\)，由
\[
 |u(x',0)|\le |u(x',s)|+\int_0^1|\partial_du(x',t)|\,\mathrm dt
\]
对 \(s\) 平均，再用 Hölder 不等式。随后对 \(x'\) 积分，得到 \(1\le C_p E_p(u)\)。因此 \(C_p(S)>0\)，尽管 \(m_d(S)=0\)。邻域厚度可以任意小，保持片上值所需的能量却不能消失。

这些例子揭示了容量与代表理论之间的联系。超过维数的 Sobolev 指数已经使点值连续受控，于是一个点也不能被视为可忽略的容量例外。指数不超过维数时，单点可以忽略，但不是所有 Lebesgue 零测集都可以忽略；下一节将在这个更严格的例外集尺度下讨论代表的唯一性。

% !TeX root = ../../Book1.tex
\section{拟处处性质}
\label{b1:sec:2602}

弱导数只由几乎处处等价类决定，而点值需要指定代表。容量提供了比 Lebesgue 零测集更细的例外集尺度：我们希望在这个尺度下选出合适的代表，并说明哪些逐点结论不再依赖选择。本节固定 \(1<p<\infty\)，先讨论实值 \(W^{1,p}(\mathbb R^d)\)，容量及能量 \(E_p\) 均沿用\cref{b1:def:sobolev-capacity}。特别地，\(E_p^{1/p}\) 与通常 Sobolev 范数等价，梯度项用欧氏模。

\ProseParagraphBreak

Kinnunen 与 Martio 的《The Sobolev capacity on metric spaces》\cite[Section 3, 3.6. Theorem and 3.7. Corollary, p. 368]{KinnunenMartio1996Capacity}用可和的坏集容量，从连续函数逼近构造拟连续代表。这条路线也适用于这里已经建立的经典容量。该文研究的度量空间范数与本书能量不同，并明确说明，在欧氏空间中两种范数及相应容量等价，而非数值相同。

\ProseParagraphBreak

“拟处处”与“拟连续”的中文名称可对照王巍翰、赵唯的《Finsler 度量测度空间上的容量》\cite[定义4.4]{WangZhao2023Capacity}。那里允许删去一般的小容量集；借助本章的开外正则性，可以把它包含在容量仍任意小的开集中。本书在定义时直接要求坏集开放，以便随后使用闭补集上的相对连续性；名称对照不意味着把其 Finsler 空间结论未经证明移到这里。

% 实读正式期刊公开全文 https://www.acadsci.fi/mathematica/Vol21/kinnunen.pdf：
% 第368页(1.3)--(1.5)区分Hajlasz型范数与经典欧氏能量；
% 第371--374页3.1--3.7，重点3.6的快速逼近完整证明及3.7；
% 第375页4.1的测度下界。第374页另核看实际页面。
% 本节不照搬该文容量数值；拟连续超水平估计及唯一性按经典开集容量独立证明。
% 前置：26.1格能量、次可加性、连续函数超水平估计、m*<=C；
% 22.1全空间光滑稠密。未调用一般精细拓扑或容量版Lebesgue点定理。

\subsection{拟处处}
\label{b1:subsec:260201}

\begin{definition}\label{b1:def:capacity-quasi-everywhere}
设 \(E\subseteq\mathbb R^d\)。若某性质在 \(E\) 上除去一个 \(C_p\)-容量为零的集合后成立，则称它在 \(E\) 上\term[nichuchu]{拟处处}[quasi-everywhere]成立，记作 \(p\)-q.e.；固定 \(p\) 后也简写为 q.e.
\end{definition}
\symidx{\(p\text{-q.e.}\)：除去 \(p\)-容量零集后成立}

由\cref{b1:prop:sobolev-capacity-basic}中的 \(m_d^*(N)\leq C_p(N)\)，容量零集一定是 Lebesgue 零测集，所以拟处处成立蕴含几乎处处成立。这是容量估计的后果，不是两种例外集的定义相同。反向蕴含不能由此推出。

\ProseParagraphBreak

可数个拟处处结论可以放在同一个例外集外使用。事实上，若第 \(j\) 条结论的例外集为 \(N_j\)，则次可加性给出
\[
 C_p\left(\bigcup_{j=1}^{\infty}N_j\right)
 \leq\sum_{j=1}^{\infty}C_p(N_j)=0.
\]
同样，如果一列集合 \(A_j\) 满足 \(\sum_j C_p(A_j)<\infty\)，则
\begin{equation}\label{b1:eq:capacity-summable-exceptions}
 C_p\left(\bigcap_{N=1}^{\infty}\bigcup_{j\geq N}A_j\right)
 \leq\inf_N\sum_{j\geq N}C_p(A_j)=0.
\end{equation}
这只用了单调性与次可加性，没有把容量当成可数可加测度。其含义是：拟处处的点只会落入有限多个 \(A_j\)。

\ProseParagraphBreak

还应分清集合上的两类条件。容量定义中的“\(u\geq1\) 几乎处处于某个开邻域”只涉及 Sobolev 等价类；“某个代表在 \(E\) 上逐点大于等于 \(1\)”涉及额外选择。下面先构造代表，再证明可对这些代表作点值容量估计，不倒过来用点值解释原定义。

\subsection{拟连续代表}
\label{b1:subsec:260202}

\begin{definition}\label{b1:def:sobolev-quasicontinuous}
设 \(v\colon\mathbb R^d\rightarrow\mathbb R\) 为有限值函数。若对每个 \(\varepsilon>0\)，都存在开集 \(G\)，使
\[
 C_p(G)<\varepsilon,\quad
 v|_{\mathbb R^d\setminus G}\ \text{连续},
\]
则称 \(v\) 为 \(p\)-\term[nilianxu]{拟连续}[quasicontinuous]函数。这里的连续性是补集上的相对连续性。若还有 \(v=u\) 几乎处处，其中 \(u\in W^{1,p}(\mathbb R^d)\)，则称 \(v\) 为 \(u\) 的拟连续代表。
\end{definition}

本节在容量零集上可以补以有限值，例如补零，不要求那些点具有由局部平均恢复的值。连续函数当然拟连续，只需取 \(G=\varnothing\)。另一个直接事实是：有限个拟连续函数的线性组合仍拟连续。为每个函数选取足够小的坏开集，取它们的并；在并集外，各项均连续，容量则由各项容量之和控制。

\begin{theorem}\label{b1:thm:sobolev-quasicontinuous-representative}
每个 \(u\in W^{1,p}(\mathbb R^d)\) 都有有限值 Borel 拟连续代表 \(\widetilde u\)。而且可以选取
\(\varphi_j\in C_c^\infty(\mathbb R^d)\)，使 \(\varphi_j\to u\) 于 \(W^{1,p}\)，并使这列函数在拟处处收敛到 \(\widetilde u\)。对每个 \(\varepsilon>0\)，还能删去一个容量小于 \(\varepsilon\) 的开集，使同一列函数在其补集上一致收敛。
\end{theorem}
\begin{proof}
由\cref{b1:thm:whole-space-sobolev-density}及范数等价性，选择
\[
 E_p(\varphi_j-u)^{1/p}\leq2^{-3j-2}.
\]
能量的 \(p\) 次方根满足三角不等式，因此
\[
 E_p(\varphi_{j+1}-\varphi_j)^{1/p}\leq2^{-3j}.
\]
定义坏开集
\[
 A_j=\{\,x:|\varphi_{j+1}(x)-\varphi_j(x)|>2^{-j}\,\},
 \quad G_N=\bigcup_{j\geq N}A_j.
\]
函数之差连续，故这些集合确实是开集。由\cref{b1:lem:continuous-capacity-level-set}和容量的次可加性，
\begin{equation}\label{b1:eq:quasicontinuous-bad-open-sets}
 C_p(G_N)\leq\sum_{j\geq N}C_p(A_j)
 \leq\sum_{j\geq N}2^{jp}E_p(\varphi_{j+1}-\varphi_j)
 \leq\sum_{j\geq N}2^{-2jp}.
\end{equation}
右端趋于零。令 \(Z=\bigcap_NG_N\)，由单调性得到 \(C_p(Z)=0\)。

固定 \(N\)。在闭集 \(F_N=\mathbb R^d\setminus G_N\) 上，对 \(k>j\geq N\)，
\[
 |\varphi_k(x)-\varphi_j(x)|
 \leq\sum_{i=j}^{k-1}|\varphi_{i+1}(x)-\varphi_i(x)|
 \leq\sum_{i=j}^{k-1}2^{-i}\leq2^{1-j}.
\]
所以这列连续函数在 \(F_N\) 上一致为 Cauchy 列，其极限有限且在 \(F_N\) 上连续。不同 \(F_N\) 上的极限在重叠处相同，因为它们来自同一列逐点极限。又有
\(\bigcup_NF_N=\mathbb R^d\setminus Z\)，故可定义
\[
 \widetilde u(x)=
 \begin{cases}
 \displaystyle\lim_{j\to\infty}\varphi_j(x),&x\notin Z,\\
 0,&x\in Z.
 \end{cases}
\]
集合 \(Z\) 为 Borel 集，在其补集上取 Borel 函数列的有限极限，再在 \(Z\) 上补零，所得函数仍为 Borel 函数。\eqref{b1:eq:quasicontinuous-bad-open-sets}及 \(F_N\) 上的一致收敛证明了拟连续性与定理最后一项。

还须核实这是原来 Sobolev 类的代表，而非仅仅某个点态极限。上述逼近速度给出
\[
 \sum_j\|\varphi_j-u\|_p^p<\infty.
\]
对非负项积分求和，得到
\(\sum_j|\varphi_j(x)-u(x)|^p<\infty\) 几乎处处，因而
\(\varphi_j(x)\to u(x)\) 几乎处处。另一方面，\(C_p(Z)=0\) 已由测度下界推出 \(m_d(Z)=0\)。在两个满测度集合的交上，点态极限只能是 \(u\)，所以 \(\widetilde u=u\) 几乎处处。
\end{proof}

这里的坏开集按容量控制，集合 \(F_N\) 不必有内点，也不必包含每个点的邻域。因此，补集上的一致极限只能直接给出相对连续性，不能将结论改写为通常连续性。光滑逼近的速度也不可省略：普通范数收敛只说明相邻差趋零，而证明需要加权后的能量级数可和。

\subsection{拟连续代表的唯一性}
\label{b1:subsec:260203}

先把连续函数的容量估计推广到拟连续代表。此时超水平集未必开，需要借助定义中的坏开集，把点态条件放入一个真正的开邻域。

\begin{proposition}\label{b1:prop:quasicontinuous-capacity-level-set}
设 \(v\) 是 \(u\in W^{1,p}(\mathbb R^d)\) 的拟连续代表。对每个 \(t>0\)，有
\begin{equation}\label{b1:eq:quasicontinuous-capacity-level-set}
 C_p\bigl(\{\,x:|v(x)|>t\,\}\bigr)
 \leq t^{-p}E_p(u).
\end{equation}
\end{proposition}
\begin{proof}
固定 \(\varepsilon>0\) 及 \(0<a<t\)。选取开集 \(G\)，使
\(C_p(G)<\varepsilon\)，且 \(v\) 在闭集 \(F=G^c\) 上连续。相对开集
\(\{\,x\in F:|v(x)|>a\,\}\) 可以写成 \(O\cap F\)，其中 \(O\) 在全空间开。于是
\[
 \{\,|v|>t\,\}\subseteq O\cup G,\quad
 |v|>a\ \text{于 }O\setminus G.
\]
由开集容量定义，取实值 \(w_G\in W^{1,p}\)，使
\[
 0\leq w_G\leq1,\quad
 w_G=1\ \text{几乎处处于 }G,\quad
 E_p(w_G)<C_p(G)+\varepsilon<2\varepsilon.
\]
截断到 \([0,1]\) 不增能量，故前两项可以同时要求。

在 Sobolev 等价类中定义
\[
 z=\max\{\,|u|/a,w_G\,\}.
\]
由\cref{b1:lem:sobolev-lattice-energy}，
\[
 z\in W^{1,p},\quad
 E_p(z)\leq a^{-p}E_p(u)+E_p(w_G)
 <a^{-p}E_p(u)+2\varepsilon.
\]
在 \(G\) 上，函数 \(w_G\) 几乎处处等于 \(1\)；在 \(O\setminus G\) 上，由 \(u=v\) 几乎处处可知 \(|u|/a>1\) 几乎处处。因此 \(z\) 是开集 \(O\cup G\) 的容许函数，进而
\[
 C_p\bigl(\{\,|v|>t\,\}\bigr)
 \leq C_p(O\cup G)
 \leq a^{-p}E_p(u)+2\varepsilon.
\]
先令 \(\varepsilon\downarrow0\)，再令 \(a\uparrow t\)，得到所需估计。
\end{proof}

估计中的能量由等价类 \(u\) 决定，左端的超水平集却由指定代表 \(v\) 决定。把两者连接起来的是拟连续性；如果给任意几乎处处代表套用同一估计，就会错误地认定所有 Lebesgue 零测集都具有零容量。

\begin{theorem}\label{b1:thm:quasicontinuous-uniqueness}
同一 Sobolev 类的两个有限值拟连续代表拟处处相等。反过来，在容量零集上任意改变一个拟连续代表的有限点值，所得函数仍是同一类的拟连续代表。
\end{theorem}
\begin{proof}
设 \(v,w\) 为同一类的拟连续代表。它们的差拟连续，并且几乎处处为零，故由前一命题，对每个正整数 \(j\)，
\[
 C_p\bigl(\{\,|v-w|>1/j\,\}\bigr)
 \leq j^p E_p(0)=0.
\]
取这些集合的可数并，得到 \(C_p(\{v\ne w\})=0\)。

再设 \(N\) 容量为零，函数 \(v'\) 在 \(N^c\) 上与 \(v\) 相同，在 \(N\) 上仍取有限值。给定 \(\varepsilon>0\)，先删去开集 \(G\)，使
\(C_p(G)<\varepsilon/2\)，且 \(v|_{G^c}\) 连续；再由外正则性选开集 \(H\supseteq N\)，使 \(C_p(H)<\varepsilon/2\)。在
\((G\cup H)^c\) 上，两个函数相同，故 \(v'\) 也连续，而
\(C_p(G\cup H)<\varepsilon\)。这就证明拟连续性。又因容量零集为 Lebesgue 零测集，且 Lebesgue 测度完备，任意这样的有限值改动仍可测，且不改变 Sobolev 类。
\end{proof}

因此可以用 \(\widetilde u\) 表示 Sobolev 类的拟连续代表，只要涉及的点态等式均在拟处处意义下理解。\symidx{\(\widetilde u\)：Sobolev 函数的拟连续代表，拟处处唯一}%
这个记号没有把某个光滑逼近序列或零容量集上的补值固定为额外结构。

\subsection{范数收敛与拟处处收敛}
\label{b1:subsec:260204}

拟连续代表的选择还与 Sobolev 范数中的极限相容。这比单纯的 \(L^p\) 几乎处处抽列结论更细，所用超水平估计则必须是刚刚证明的拟连续版本。

\begin{theorem}\label{b1:thm:sobolev-quasi-everywhere-subsequence}
若 \(u_n\to u\) 于 \(W^{1,p}(\mathbb R^d)\)，为各项及极限任取有限值拟连续代表 \(\widetilde u_n,\widetilde u\)，则存在子列
\(\widetilde u_{n_j}\) 拟处处收敛到 \(\widetilde u\)。对每个 \(\varepsilon>0\)，还能删去一个容量小于 \(\varepsilon\) 的开集，使该子列在补集上一致收敛。
\end{theorem}
\begin{proof}
抽取子列，使
\[
 E_p(u_{n_j}-u)\leq2^{-2jp}.
\]
差函数的拟连续代表可取
\(\widetilde u_{n_j}-\widetilde u\)。由超水平估计，
\[
 C_p\bigl(\{\,|\widetilde u_{n_j}-\widetilde u|>2^{-j}\,\}\bigr)
 \leq2^{-jp}.
\]
外正则性允许把每个超水平集包含在开集 \(O_j\) 中，并要求
\(C_p(O_j)<2^{1-jp}\)。令
\[
 V_N=\bigcup_{j\geq N}O_j,\quad Z=\bigcap_NV_N.
\]
则 \(C_p(V_N)\leq2\sum_{j\geq N}2^{-jp}\to0\)，因而 \(C_p(Z)=0\)。在 \(V_N^c\) 上，对每个 \(j\geq N\) 都有
\[
 |\widetilde u_{n_j}(x)-\widetilde u(x)|\leq2^{-j}.
\]
这同时证明补集上的一致收敛和 \(Z^c\) 上的逐点收敛。给定 \(\varepsilon\) 后取足够大的 \(N\)，即得最后一项。
\end{proof}

对整个收敛序列，虽未必有逐点收敛，上述命题仍给出每个固定 \(t>0\) 的定量结论：
\[
 C_p\bigl(\{\,|\widetilde u_n-\widetilde u|>t\,\}\bigr)
 \leq t^{-p}E_p(u_n-u)\longrightarrow0.
\]
这里不能用任意的可测代表替换拟连续代表，也不能只假设零阶 \(L^p\) 收敛而删去梯度控制。

\ProseParagraphBreak

再与\secref{b1:sec:1404}的精确代表 \(u^*\) 比较。那里的定义由缩小球的平均值决定，与可测代表的任意改动无关；本节的 \(\widetilde u\) 则由容量和连续逼近确定，只在拟处处意义下唯一。由于两者均与 \(u\) 几乎处处相等，已有
\[
 \widetilde u=u^*\quad\text{几乎处处}.
\]
若 \(\widetilde u\) 在某点 \(x\) 通常连续，则对充分小的球，
\[
 \frac1{m_d\bigl(B(x,r)\bigr)}
 \int_{B(x,r)}|u(y)-\widetilde u(x)|\,\mathrm dy
 \leq\sup_{y\in B(x,r)}
          |\widetilde u(y)-\widetilde u(x)|\longrightarrow0.
\]
所以这个点具有由平均绝对偏差确定的值，且 \(u^*(x)=\widetilde u(x)\)。对于一般拟连续点，删去小容量集后的相对连续性本身并未估计球内落在该坏集中的积分。本节因而不把上述几乎处处等式升级为“每个拟连续代表拟处处都是 Lebesgue 点”；这需要进一步的容量微分估计，不能仅从普通 Lebesgue 点定理推出。

\ProseParagraphBreak

单点容量把两种连续性的边界具体化。若 \(p>d\)，由\cref{b1:prop:sobolev-point-capacity}，每个非空集合的容量都有同一个正下界。拟连续定义中把 \(\varepsilon\) 取得小于这个下界时，坏开集只能为空，故拟连续函数实际上通常连续。若 \(1<p\leq d\)，单点容量为零，次可加性说明 \(C_p(\mathbb Q^d)=0\)。于是 \(\mathbf1_{\mathbb Q^d}\) 是零 Sobolev 类的一个拟连续代表，却处处不连续；它与连续代表零函数恰好在一个容量零集上不同。这也说明拟处处唯一不能替换为逐点唯一。

\ProseParagraphBreak

复值情形可按实部、虚部分别构造代表，再把两组坏开集合并。两个分量的唯一性给出复代表的拟处处唯一；范数收敛抽列时先后为两个分量选子列，仍得到共同坏开集外的一致收敛。容量的容许函数依旧是实值非负函数，不因讨论复值代表而改变。

% !TeX root = ../../Book1.tex
\section{容量与 Hausdorff 测度}
\label{b1:sec:2603}

容量由函数能量定义，Hausdorff 测度由小集合覆盖定义。两者相遇的地方是球截断：半径为 \(r\) 的球可以用一个梯度能量约为 \(r^{d-p}\) 的函数包住。这个尺度解释了 \(d-p\) 为何成为临界维数，但它本身还不能处理临界测度有限、却不为零的集合。

Kinnunen 与 Martio 的《The Sobolev capacity on metric spaces》\cite[pp. 380--381, 4.13. Theorem and 4.15. Theorem]{KinnunenMartio1996Capacity}组织了容量与 Hausdorff 覆盖之间的两个方向。本节使用本书的经典弱梯度，分别补出临界测度有限时的凸组合论证，以及反方向中的 \(5r\) 覆盖步骤。除非另行限制，以下集合不必为 Borel 集；\(\mathcal H^s\) 与 \(C_p\) 均按外测度值理解。

% 实读 KinnunenMartio1996Capacity 官方全文第380--381页4.13--4.15及完整证明，
% 并核看第381页。该文的g是Hajlasz梯度，本节不将其当作弱梯度。
% 球覆盖方向参考4.13的组织；反向以本书球Poincare重写4.15，
% 并用13章5r引理展开原文援引Federer的最后一步。
% H^{d-p}有限/σ有限推出零容量是本节独立展开的经典Sobolev凸组合论证。
% 前置：26.1格能量与球截断；9章Mazur；13章5r；15章Hausdorff外测度；
% 18章径向积分；21章Sobolev弱子列；24章球Poincare。不用Frostman或一般可容性。

\subsection{维数估计}
\label{b1:subsec:260301}

先设 \(1<p<d\)，记 \(t=d-p>0\)。由\cref{b1:prop:sobolev-capacity-balls}，当 \(0<r\leq1\) 时，
\[
 C_p\bigl(B(x,r)\bigr)\leq C_{d,p}(r^d+r^{d-p})
 \leq2C_{d,p}r^t.
\]
它立即给出一个覆盖上界。

\begin{proposition}\label{b1:prop:capacity-hausdorff-upper}
设 \(1<p<d\)。存在仅依赖 \(d,p\) 的常数 \(C\)，使每个
\(E\subseteq\mathbb R^d\) 都满足
\[
 C_p(E)\leq C\mathcal H^{d-p}(E).
\]
\end{proposition}
\begin{proof}
令 \(t=d-p\)，并假设右端有限。由 Hausdorff 覆盖定义，对每个
\(0<\delta<1\) 及 \(\eta>0\)，可以用开球 \(B(x_i,r_i)\) 覆盖 \(E\)，使
\begin{equation}\label{b1:eq:critical-hausdorff-ball-cover}
 0<r_i\leq\delta,\quad
 \sum_i r_i^t\leq C_t\mathcal H^t(E)+\eta.
\end{equation}
具体地，先取直径不超过 \(\delta/4\) 的可数覆盖。直径 \(a_i>0\) 的非空成员，可放入以其中一点为中心、半径 \(2a_i\) 的开球；直径为零的成员至多是单点，为这些单点另选正半径，使半径的 \(t\) 次幂之和小于任意给定误差。再利用归一化常数 \(c_t>0\)，便得到所列球覆盖。这里使用了 \(t>0\)。

单调性、可数次可加性及球上界给出
\[
 C_p(E)\leq\sum_iC_p\bigl(B(x_i,r_i)\bigr)
 \leq C_{d,p}\sum_i r_i^t
 \leq C\mathcal H^t(E)+C_{d,p}\eta.
\]
令 \(\eta\downarrow0\) 即得结论。右端无穷时不需另证。
\end{proof}

这个命题已说明 \(\mathcal H^{d-p}(E)=0\) 足以推出零容量。若临界测度只是有限，上界却只给出有限容量；要进一步得到零，还必须利用覆盖半径能够趋于零。

\begin{theorem}\label{b1:thm:finite-critical-hausdorff-zero-capacity}
设 \(1<p<d\)。若 \(E\subseteq\mathbb R^d\) 可以写成可数个具有有限
\(\mathcal H^{d-p}\) 外测度的集合之并，则 \(C_p(E)=0\)。特别地，
\(\mathcal H^{d-p}(E)<\infty\) 时容量为零。
\end{theorem}
\begin{proof}
先设 \(\mathcal H^t(E)<\infty\)，其中 \(t=d-p\)。取 \(0<\delta_n<1\)、\(\delta_n\downarrow0\)，并用\eqref{b1:eq:critical-hausdorff-ball-cover}选取开球覆盖
\[
 E\subseteq G_n=\bigcup_i B(x_{n,i},r_{n,i}),\quad
 r_{n,i}\leq\delta_n,\quad
 \sum_i r_{n,i}^t\leq M.
\]
常数 \(M<\infty\) 与 \(n\) 无关，可把每次覆盖的额外误差取为 \(1\)。

固定 \(0\leq\chi\leq1\)、\(\chi\in C_c^\infty\bigl(B(0,2)\bigr)\)，使 \(\chi=1\) 于 \(B(0,1)\)，并令
\[
 \chi_{n,i}(x)=\chi\left(\frac{x-x_{n,i}}{r_{n,i}}\right).
\]
直接伸缩计算给出
\[
 E_p(\chi_{n,i})
 \leq C_{d,p}(r_{n,i}^d+r_{n,i}^{d-p}).
\]
我们需要这些截断的可数最大值，而不是它们的和；求和可能在高度上重复付出代价。对固定 \(n\)，置
\[
 v_{n,k}=\max_{1\leq i\leq k}\chi_{n,i},\quad
 u_n=\sup_{i\geq1}\chi_{n,i},\quad
 S_n=\bigcup_i B(x_{n,i},2r_{n,i}).
\]
由\cref{b1:lem:sobolev-lattice-energy}，有
\[
 E_p(v_{n,k})\leq\sum_iE_p(\chi_{n,i})\leq C_{d,p}M.
\]
又有 \(0\leq v_{n,k}\uparrow u_n\leq\mathbf1_{S_n}\)，并且
\begin{equation}\label{b1:eq:critical-cover-small-volume}
 m_d(S_n)\leq2^d\omega_d\sum_i r_{n,i}^d
 \leq2^d\omega_d\delta_n^p\sum_i r_{n,i}^{d-p}
 \leq C_dM\delta_n^p.
\end{equation}
所以 \(v_{n,k}\to u_n\) 于 \(L^p\)。由\cref{b1:cor:sobolev-weak-subsequence}取一个 Sobolev 弱收敛子列，其零阶极限只能为 \(u_n\)。能量是一个等价范数的 \(p\) 次幂，故弱下半连续，得到
\[
 u_n\in W^{1,p}(\mathbb R^d),\quad
 E_p(u_n)\leq C_{d,p}M,\quad
 u_n=1\ \text{几乎处处于 }G_n.
\]
最后一项也可从逐点最大值直接读出；作为 Sobolev 类使用时，只要求这个几乎处处条件。

由\eqref{b1:eq:critical-cover-small-volume}，
\(\|u_n\|_p^p\leq C_dM\delta_n^p\to0\)。再从能量有界列
\((u_n)\) 中抽取 Sobolev 弱收敛子列，零阶强极限把其弱极限确定为零。仍把该子列记作 \((u_n)\)，并保留相应的 \(G_n\)。

由\cref{b1:thm:mazur}，每个尾部的有限凸组合可以在 Sobolev 范数中任意接近零。于是存在
\[
 z_j=\sum_{n=j}^{N_j}a_{j,n}u_n,\quad
 a_{j,n}\geq0,\quad \sum_{n=j}^{N_j}a_{j,n}=1,\quad
 E_p(z_j)\longrightarrow0.
\]
有限交 \(U_j=\bigcap_{n=j}^{N_j}G_n\) 仍是包含 \(E\) 的开集。在这个交集上，所有参与组合的 \(u_n\) 几乎处处等于 \(1\)，所以
\(z_j=1\) 几乎处处于 \(U_j\)。因此
\[
 C_p(E)\leq C_p(U_j)\leq E_p(z_j)\longrightarrow0.
\]
此处必须用有限凸组合；无限多个开邻域的交未必还是开邻域。

最后，若 \(E=\bigcup_\ell E_\ell\)，每个 \(E_\ell\) 的临界外测度有限，就由刚证的结论及容量的可数次可加性得到 \(C_p(E)=0\)。
\end{proof}

这个论证不要求 \(E\) 紧，也不要求从 \(E\) 上取得一个 Frostman 测度。它控制的是开集 \(S_n\) 的体积，不是 \(u_n\) 的拓扑支集；当覆盖的中心稠密时，\(S_n\) 的闭包可能很大。

\subsection{零容量集}
\label{b1:subsec:260302}

反方向的目标不同：不是从覆盖造容许函数，而是用一组能量很小的容许函数造出一个函数，使它在给定集合每一点附近的平均值都发散。这迫使梯度能量在那些点周围聚集。

\begin{theorem}\label{b1:thm:zero-capacity-hausdorff-dimension}
设 \(1<p\leq d\)，且 \(E\subseteq\mathbb R^d\) 满足 \(C_p(E)=0\)。则对每个 \(s>d-p\)，都有
\[
 \mathcal H^s(E)=0.
\]
特别地，零容量集的 Hausdorff 维数至多为 \(d-p\)。
\end{theorem}
\begin{proof}
由容量为零，选取开集 \(G_j\supseteq E\) 和非负函数
\(u_j\in W^{1,p}(\mathbb R^d)\)，使
\[
 u_j\geq1\ \text{几乎处处于 }G_j,\quad
 E_p(u_j)^{1/p}\leq2^{-j}.
\]
能量范数的可和性和 Sobolev 完备性说明
\[
 w=\sum_{j=1}^{\infty}u_j\in W^{1,p}(\mathbb R^d).
\]
这里可取非负逐点和作为几乎处处代表：部分和在 \(L^p\) 中收敛，又单调递增，故逐点和有限几乎处处，并与范数极限相同。

任取 \(x\in E\) 及正整数 \(N\)。有限交
\(\bigcap_{j=1}^NG_j\) 是 \(x\) 的开邻域，其中 \(w\geq N\) 几乎处处。因此
\begin{equation}\label{b1:eq:zero-capacity-divergent-means}
 \lim_{r\downarrow0}w_{B(x,r)}=+\infty
 \quad(x\in E).
\end{equation}
每个固定球上的平均仍有限，因为 \(w\in L^p\subseteq L^1_{\mathrm{loc}}\)。这个发散结论由开邻域上的几乎处处不等式推出，不需要预先给 \(w\) 选拟连续点值。

令 \(\mu=|\nabla w|^p\,m_d\)，它是有限正 Borel 测度。固定 \(s>d-p\)。假设某个 \(x\in E\) 满足
\[
 \limsup_{r\downarrow0}\frac{\mu\bigl(B(x,r)\bigr)}{r^s}<\infty.
\]
则存在 \(R>0\) 和 \(A<\infty\)，使
\(\mu\bigl(B(x,r)\bigr)\leq Ar^s\) 对 \(0<r\leq R\) 均成立。记
\(r_i=2^{-i}R\)、\(B_i=B(x,r_i)\)。由 Hölder 不等式及\cref{b1:thm:ball-poincare-lp}，
\[
 \begin{aligned}
 |w_{B_{i+1}}-w_{B_i}|
 &\leq\frac1{m_d(B_{i+1})}\int_{B_{i+1}}|w-w_{B_i}|\\
 &\leq C_d\,m_d(B_i)^{-1/p}\|w-w_{B_i}\|_{L^p(B_i)}\\
 &\leq C_d r_i^{1-d/p}\|\nabla w\|_{L^p(B_i)}
 \leq C_dA^{1/p}r_i^{(p-d+s)/p}.
 \end{aligned}
\]
指数 \((p-d+s)/p\) 严格为正，所以右端对 \(i\) 可和。相邻平均值之差可和，说明 \((w_{B_i})\) 收敛到有限值，与\eqref{b1:eq:zero-capacity-divergent-means}矛盾。因此
\begin{equation}\label{b1:eq:capacity-zero-gradient-density}
 E\subseteq
 \left\{\,x:\limsup_{r\downarrow0}
               \frac{\mu\bigl(B(x,r)\bigr)}{r^s}=+\infty\,\right\}.
\end{equation}

最后把这个密度条件变成覆盖估计。固定 \(M>0\)、\(\delta>0\) 和 \(R_0>0\)。对每个 \(x\in E\cap B(0,R_0)\)，由\eqref{b1:eq:capacity-zero-gradient-density}，存在任意小的半径
\(0<r<\delta/10\)，使
\(\mu\bigl(B(x,r)\bigr)>Mr^s\)。把全部这样的开球组成一个族。由\cref{b1:lem:five-r-covering}，可选至多可数个两两不交的球 \(B_i=B(x_i,r_i)\)，使
\[
 E\cap B(0,R_0)\subseteq\bigcup_i5B_i,\quad
 \operatorname{diam}(5B_i)=10r_i<\delta.
\]
于是
\[
 \begin{aligned}
 \mathcal H^s_\delta\bigl(E\cap B(0,R_0)\bigr)
 &\leq c_s10^s\sum_i r_i^s\\
 &\leq\frac{c_s10^s}{M}\sum_i\mu(B_i)
 \leq\frac{c_s10^s}{M}\mu(\mathbb R^d).
 \end{aligned}
\]
此处测度只作用于所选开球，不要求 \(E\) 可测。先令
\(\delta\downarrow0\)，再令 \(M\to\infty\)，得到
\(\mathcal H^s(E\cap B(0,R_0))=0\)。最后对整数 \(R_0\) 取可数并，便有 \(\mathcal H^s(E)=0\)。
\end{proof}

证明中严格不等式 \(s>d-p\) 用在几何级数的收敛上。若直接取 \(s=d-p\)，相邻平均值的估计只剩一个不随尺度衰减的常数，不能据此推出平均值收敛。

\subsection{奇异集}
\label{b1:subsec:260303}

上一节产生的例外集，现在可以附加几何大小的说明。这里所说的异常，是某条拟处处结论的例外，不是所有通常不连续点的统称。

\begin{corollary}\label{b1:cor:capacity-exception-dimension}
设 \(1<p\leq d\)。同一 Sobolev 类的两个拟连续代表不相等的集合，对每个 \(s>d-p\) 都具有零 \(\mathcal H^s\) 测度。若一列 Sobolev 函数在范数中收敛，则可以选取子列，使其拟连续代表不逐点收敛到极限代表的集合也满足同一结论。
\end{corollary}
\begin{proof}
由\cref{b1:thm:quasicontinuous-uniqueness}，第一类集合容量为零。由\cref{b1:thm:sobolev-quasi-everywhere-subsequence}，第二类集合可以包含在一个容量零集中。分别应用前一定理即可。
\end{proof}

例如，在 \(\mathbb R^3\) 中取 \(p=2\)，一条有限线段的
\(\mathcal H^1\) 测度有限，故其 \(2\)-容量为零。这个结论正好落在临界维数 \(d-p=1\) 上，不能仅凭“维数较低”而跳过有限临界测度的论证。

更一般地，在 \(1<p<d\) 时，若
\(\dim_{\mathrm H}E<d-p\)，则 \(\mathcal H^{d-p}(E)=0\)，所以容量为零；若
\(\dim_{\mathrm H}E>d-p\)，容量便不可能为零。两种严格不等式给出清楚的判别，而维数恰为 \(d-p\) 时仍须查看更细的测度或覆盖信息。特别是，已证的充分条件包括临界测度 \(\sigma\)-有限，并不只包括临界测度为零。

这些结果没有把普通 Lebesgue 点定理的例外集改成容量零集。\secref{b1:sec:2602}只在已经选定的拟连续代表与容量例外集之间建立了关系；由此得到的维数结论也只适用于这些例外集。例如当 \(1<p\leq d\) 时，零类的代表 \(\mathbf1_{\mathbb Q^d}\) 拟连续，却通常处处不连续，其通常不连续点集是整个空间。它不与本推论冲突，因为与零代表不相等的集合只是 \(\mathbb Q^d\)。

\subsection{临界维数}
\label{b1:subsec:260304}

当 \(p=d>1\) 时，幂次 \(r^{d-p}\) 变为常数，普通伸缩截断不能反映小球容量趋零的速度。这时可以把高度下降分散到许多同心环带，得到一个对数尺度。

\begin{proposition}\label{b1:prop:critical-logarithmic-capacity}
设 \(d>1\)。存在仅依赖于 \(d\) 的常数 \(\kappa_d\)，使
\[
 C_d\bigl(B(x,r)\bigr)
 \leq\kappa_d\bigl(\log(1/r)\bigr)^{1-d},
 \quad 0<r\leq\tfrac12.
\]
这里左端 \(C_d(\,\cdot\,)\) 是指数为 \(d\) 的容量。
\end{proposition}
\begin{proof}
平移不改变能量，故只考虑球心为零。令 \(L=\log(1/r)\)，取径向函数
\[
 v_r(x)=
 \begin{cases}
 1,&|x|\leq r,\\
 L^{-1}\log(1/|x|),&r<|x|<1,\\
 0,&|x|\geq1.
 \end{cases}
\]
对固定 \(r>0\)，该函数为全空间 Lipschitz 函数，紧支集且在
\(B(0,r)\) 上等于 \(1\)，因而是容许函数。在中间环带上，
\(|\nabla v_r(x)|=1/(L|x|)\) 几乎处处；其余部分梯度为零。
由\cref{b1:prop:radial-shell-integration}，
\[
 \int_{\mathbb R^d}|\nabla v_r|^d
 =d\omega_d L^{-d}\int_r^1\frac{\mathrm d\rho}{\rho}
 =d\omega_d L^{1-d}.
\]
非齐次容量还包括零阶项，不能只估计这一梯度积分。作代换
\(t=\log(1/\rho)\)，得到
\[
 \begin{aligned}
 \int_{\mathbb R^d}|v_r|^d
 &=\omega_dr^d+d\omega_dL^{-d}
       \int_r^1\bigl(\log(1/\rho)\bigr)^d\rho^{d-1}\,\mathrm d\rho\\
 &=\omega_de^{-dL}
   +d\omega_dL^{-d}\int_0^L t^de^{-dt}\,\mathrm dt
 \leq A_dL^{-d}.
 \end{aligned}
\]
最后一步使用
\(\sup_{L\geq\log2}L^de^{-dL}<\infty\) 及
\(\int_0^\infty t^de^{-dt}\,\mathrm dt<\infty\)；指数衰减使这两个量有限。因此，当 \(L\geq\log2\) 时，
\[
 E_d(v_r)\leq d\omega_dL^{1-d}+A_dL^{-d}
 \leq\kappa_dL^{1-d}.
\]
以这个能量作为球容量的上界，即得结论。
\end{proof}

为记录相应的覆盖判据，记
\[
 \ell_d(r)=\bigl(\log(1/r)\bigr)^{1-d},
 \quad 0<r\leq\tfrac12.
\]
对 \(0<\delta\leq1/2\)，直接用可数开球覆盖定义对数覆盖量
\[
 \mathcal H^{\ell_d}_\delta(E)
 =\inf\left\{\,
    \sum_i\ell_d(r_i):
    E\subseteq\bigcup_i B(x_i,r_i),\quad 0<r_i\leq\delta
   \,\right\}.
\]
再令
\[
 \mathcal H^{\ell_d}(E)
 =\lim_{\delta\downarrow0}\mathcal H^{\ell_d}_\delta(E).
\]
这个极限由尺度单调性存在。本式以半径为代价的自变量，没有另加
\(c_s\) 归一化；它是对数代价下的 Hausdorff 覆盖构造，不是
\(\mathcal H^0\) 的另一种记法。

\begin{corollary}\label{b1:cor:logarithmic-cover-zero-capacity}
若 \(d>1\)，则对任意 \(E\subseteq\mathbb R^d\)，有
\[
 C_d(E)\leq\kappa_d\mathcal H^{\ell_d}(E).
\]
特别地，对数覆盖量为零的集合具有零 \(d\)-容量。
\end{corollary}
\begin{proof}
对任意半径不超过 \(\delta\leq1/2\) 的开球覆盖，用容量的次可加性和前一命题，得到
\[
 C_d(E)\leq\sum_iC_d\bigl(B(x_i,r_i)\bigr)
 \leq\kappa_d\sum_i\ell_d(r_i).
\]
先对覆盖取下确界，再令 \(\delta\downarrow0\) 即可。
\end{proof}

由于 \(\ell_d(r)\to0\)，单点的对数覆盖量为零；可数集也可逐点分配任意小的可和代价，所以仍为零。另一方面，\(\mathcal H^0\) 是计数测度，单点已有 \(\mathcal H^0\) 值 \(1\)。这说明在临界指数处，以零维测度替代对数尺度会丢掉信息。

零 \(d\)-容量仍由前面定理推出 \(\mathcal H^s(E)=0\) 对全部 \(s>0\) 成立，但“维数为零”本身没有估计这里的对数覆盖量。也不能把 \(p<d\) 时有限临界测度的凸组合证明原样搬来断言有限对数覆盖量必然为零容量：上面的对数截断延伸到固定大小的球，未使容许函数在一个体积趋零的区域外为零。本节在此只使用已经证明的零覆盖量判据。

% !TeX root = ../../Book1.tex
\section{细拓扑}
\label{b1:sec:2604}

拟连续性允许删去容量很小的开集，却没有把余集变成通常的开集。若希望直接谈论这些代表的连续性，就需要一种能够辨认容量大小的邻域。本节先从相对容量构造这种拓扑，再说明它与边界值问题的联系。相对容量的基本性质和拓扑公理将在这里证明；Cartan 性质、拟连续与细连续的联系以及非线性 Wiener 判别属于进一步的势论，均明确给出证明概要，不把它们当成前面 Sobolev 理论的直接推论。

“细拓扑”与“薄性”采用 Doob 的《经典位势论与概率位势论（上册）》中文目录中的名称。\footnote{科学出版社的\href{https://www.ecsponline.com/goods.php?id=6199}{官方书目}第XI章及第XI.2、第XIII.17节采用这些名称；这里仅据目录核对称谓，不以古典位势论的目录代替非线性 \(p\)-容量的定义或证明。}本节的“细”指一套特定的势论拓扑，并不是泛指任意比原拓扑更细的拓扑。

% 中文“细拓扑”“薄性”据科学出版社《经典位势论与概率位势论（上册）》
% 官方目录第XI章、XI.2及XIII.17；目录只作名称证据，不作非线性证明来源。
% 实读 Heinonen--Kilpeläinen--Martio 1989, AIF 39(2), 293--318：
% 293--295页引言，298--299页第3节定义及Theorem 3.2；
% 300--303页只读相关证明路线，不声称完整重建全部势论。
% 实读 Björn--Björn--Latvala 2018, J. Anal. Math. 135(1), 59--83：
% 62页Theorem 1.4，64页拟开/拟连续定义，71页Definitions 6.1--6.2，
% 75页第8节Theorem 1.4的证明；未读完其所调用的所有先前论文。
% 实读 Kilpeläinen--Malý 1994, Acta Math. 172(1), 137--161：
% 137--139页Theorems 1.1、1.3、1.6的陈述；
% 155--158页第5节Cartan性质、Perron比较与Wiener判别证明路线。
% 未完整读第4节非线性势估计的证明；本节将该估计明确列为未展开输入。
% 相对开集a.e.定义、格能量证明和离散尺度整理按本书前章工具重写。

\subsection{薄性}
\label{b1:subsec:260401}

本节仍取 \(1<p<\infty\)。全空间容量 \(C_p\) 同时计入函数和梯度的能量；在一个不断缩小的球内考察边界障碍时，更合适的是固定外边界、只计梯度的相对容量。两个对象控制相同的局部零容量现象，但尺度公式不同，不能混用。

\begin{definition}\label{b1:def:relative-variational-capacity}
设 \(D\subset\mathbb R^d\) 为有界开集。给定开集 \(G\subset D\)，定义它相对于 \(D\) 的 \term[xiangdui-bianfen-rongliang]{相对变分容量}[relative variational capacity]为
\[
 \operatorname{cap}_p(G,D)
 =\inf\left\{\int_D|\nabla v|^p\,\mathrm dx:
 v\in W_0^{1,p}(D),\
 v\ge1\ \text{几乎处处于 }G\right\}.
\]
空容许类的下确界为 \(+\infty\)。对任意 \(E\subset D\)，再令
\[
 \operatorname{cap}_p(E,D)
 =\inf\{\,\operatorname{cap}_p(G,D):
 E\subset G\subset D,\ G\text{开}\,\}.
\]
这里的 \(W_0^{1,p}(D)\) 仍是内部紧支集光滑函数的闭包，不预先用边界点值解释。
\end{definition}
\symidx{\(\operatorname{cap}_p(E,D)\)：相对于开集的纯梯度容量}

\begin{lemma}\label{b1:lem:relative-capacity-basic}
相对容量关于 \(E\) 单调且可数次可加，关于外域 \(D\) 反向单调。容许函数可以限制为 \(0\le v\le1\)。对紧集 \(K\Subset D\)，还有
\begin{equation}\label{b1:eq:relative-compact-capacity}
 \operatorname{cap}_p(K,D)
 =\inf\left\{\int_D|\nabla\varphi|^p\,\mathrm dx:
 \varphi\in C_c^\infty(D),\
 \varphi\ge1\text{ 于 }K\text{ 的某个开邻域}\right\}.
\end{equation}
\end{lemma}
\begin{proof}
先说明格运算不会破坏零边界闭包。以 \(C_c^\infty(D)\) 逼近 \(v\in W_0^{1,p}(D)\)，对逼近函数取正部或截断。所得函数有内部紧支集，仍可在 \(W^{1,p}\) 中用内部光滑函数逼近。截断的梯度界给有界性，零阶项则收敛到相应的截断；弱极限识别及 \(W_0^{1,p}(D)\) 的弱闭性说明极限仍在这个空间。由正部公式又得到最大值、最小值的闭包性质。于是前面的格能量论证同样给出
\[
 \int_D|\nabla(v\vee w)|^p
 +\int_D|\nabla(v\wedge w)|^p
 =\int_D|\nabla v|^p+\int_D|\nabla w|^p.
\]
取值截断到 \([0,1]\) 不增加能量，也不改变容许性。

集合单调性来自容许类的包含；外域增大时，利用\cref{b1:prop:w0-zero-extension}把容许函数补零即可。证明可数次可加性时，只须考虑容量之和有限的情形。为各集选取开覆盖及近乎最优的容许函数 \(v_j\in[0,1]\)，使梯度能量之和有限，并令 \(w_N=\max_{j\le N}v_j\)。格公式控制其梯度，\cref{b1:thm:zero-boundary-poincare}控制其零阶项；而
\[
 0\le w_N\uparrow w,\quad
 w^p\le\sum_j v_j^p\in L^1(D).
\]
故 \(w_N\to w\) 于 \(L^p(D)\)。抽取 Sobolev 弱子列，极限仍在 \(W_0^{1,p}(D)\)，且梯度能量不超过各项能量之和。它在开覆盖之并上几乎处处等于 \(1\)。最后让选取容许函数时的总误差趋零。

式\eqref{b1:eq:relative-compact-capacity}右边至少等于左边。反过来，选取开集 \(G\supset K\) 和在 \(G\) 上几乎处处等于 \(1\) 的容许函数 \(v\in[0,1]\)。取 \(\psi\in C_c^\infty(G)\)，使 \(0\le\psi\le1\) 且在 \(K\) 的邻域等于 \(1\)，再取 \(v_j\in C_c^\infty(D)\) 趋于 \(v\)。函数
\[
 \varphi_j=\psi+(1-\psi)v_j
\]
在该邻域等于 \(1\)，并且
\[
 \varphi_j-v=(1-\psi)(v_j-v),
\]
因为 \(\psi(1-v)=0\) 几乎处处。光滑乘法的有界性给出 \(\varphi_j\to v\) 于 \(W^{1,p}(D)\)，所以光滑容许函数的能量可逼近原下确界。
\end{proof}

这也说明本定义与经典势论先在紧集上定义容量的口径一致。具体地，取开集 \(G\) 的递增紧穷竭 \(K_j\)，使 \(\bigcup_j\operatorname{int}K_j=G\)。若这些紧集的容量有统一上界，对其光滑容许函数取弱极限，每个固定 \(K_j\) 上的几乎处处约束都保留下来，因而极限在 \(G\) 上几乎处处至少为 \(1\)。这就证明
\(\operatorname{cap}_p(G,D)=\sup_{K\Subset G}\operatorname{cap}_p(K,D)\)；最后再对任意集取开外正则化。此处不需要先定义某个 Sobolev 代表在任意集合上的逐点约束。

\begin{proposition}\label{b1:prop:relative-capacity-scale}
平移、伸缩给出
\[
 \operatorname{cap}_p(x+rE,x+rD)
 =r^{d-p}\operatorname{cap}_p(E,D),\quad r>0.
\]
存在仅依赖 \(d,p\) 的正常数 \(c,C\)，使
\begin{equation}\label{b1:eq:relative-ball-capacity}
 cr^{d-p}\le
 \operatorname{cap}_p\bigl(B(x,r),B(x,2r)\bigr)
 \le Cr^{d-p}.
\end{equation}
此外，对任意 \(F\subset B(x,r)\)，有
\begin{equation}\label{b1:eq:relative-outer-ball-comparison}
 \operatorname{cap}_p\bigl(F,B(x,4r)\bigr)
 \le\operatorname{cap}_p\bigl(F,B(x,2r)\bigr)
 \le C\operatorname{cap}_p\bigl(F,B(x,4r)\bigr).
\end{equation}
\end{proposition}
\begin{proof}
作变量代换 \(y=x+rz\)，梯度产生因子 \(r^{-1}\)，体积产生因子 \(r^d\)；该变换及其逆保持零边界闭包和开邻域约束，所以先对开集、再对任意集得到精确的尺度公式。

固定一个在单位球上等于 \(1\)、支集含于二倍球的光滑截断，伸缩后得到球对的上界。若 \(v\) 是内球的容许函数，则零边界 Poincaré 不等式给出
\[
 \omega_d r^d
 \le\int_{B(x,2r)}|v|^p
 \le(4r)^p\int_{B(x,2r)}|\nabla v|^p,
\]
从而得到下界。

最后一个式子的左侧来自外域单调性。证明右侧时，取在 \(B(x,r)\) 上等于 \(1\)、支集含于 \(B(x,2r)\) 的截断 \(\eta\)，其中 \(|\nabla\eta|\le C/r\)。将 \(B(x,4r)\) 中的容许函数 \(v\) 乘以 \(\eta\)，便得到较小外球中的容许函数。乘积法则和零边界 Poincaré 不等式给出
\[
 \int|\nabla(\eta v)|^p
 \le C\int|\nabla v|^p+Cr^{-p}\int|v|^p
 \le C'\int|\nabla v|^p.
\]
若原约束在开集 \(G\supset F\) 上成立，乘积的约束就在 \(G\cap B(x,r)\) 上成立，仍是 \(F\) 的开邻域。取下确界即得结论。
\end{proof}

特别地，\(p=d\) 时球对容量有与半径无关的正上下界。这不与全空间容量 \(C_d(B(x,r))\to0\) 矛盾：相对容量要求函数在距离中心仅为 \(2r\) 的地方已经满足零边界条件。

\begin{definition}\label{b1:def:p-thin-set}
给定 \(E\subset\mathbb R^d\) 和 \(x\in\mathbb R^d\)，令 \(r_j=2^{-j}\)。若
\begin{equation}\label{b1:eq:p-wiener-series}
 \sum_{j=0}^{\infty}
 \left(
 \frac{\operatorname{cap}_p\bigl(E\cap B(x,r_j),B(x,2r_j)\bigr)}
 {r_j^{d-p}}
 \right)^{1/(p-1)}<\infty,
\end{equation}
则称 \(E\) 在 \(x\) 处具有 \(p\)-\term[boxing]{薄性}[thinness]，或称 \(E\) 在 \(x\) 处是 \(p\)-薄的；级数发散时则称它在该点不是 \(p\)-薄的。
\end{definition}

分母也可以换成球对容量，因\eqref{b1:eq:relative-ball-capacity}只会改变正常数。有限个大尺度项不影响收敛性。外球的倍数也不是定义的本质：\eqref{b1:eq:relative-outer-ball-comparison}及其相同的截断证明允许把倍数 \(2\) 换成任意固定的 \(\lambda>1\)。

文献通常在连续的对数尺度上写 Wiener 积分；本节采用二进级数，使任意集合 \(E\) 都可直接进入定义。两种尺度判据的比较只需前述估计：当 \(r/2\le t\le r\) 时，利用内集单调性、外球比较以及 \(t^{d-p}\) 与 \(r^{d-p}\) 的可比性，中间尺度的归一化容量被相邻两个二进尺度的量从上下控制。因此逐个对数尺度求和与通常的 Wiener 积分给出同一收敛判据，也不依赖起始半径的选择。相对容量及这一薄性条件可与 Heinonen、Kilpeläinen 和 Martio 的《Fine topology and quasilinear elliptic equations》\cite[Section 3]{Heinonen1989Fine}对读。

\subsection{细连续性}
\label{b1:subsec:260402}

\begin{definition}\label{b1:def:p-fine-topology}
设 \(U\subset\mathbb R^d\)。若 \(\mathbb R^d\setminus U\) 在每个 \(x\in U\) 处都是 \(p\)-薄的，则称 \(U\) 为 \(p\)-细开集。所有这样的集合组成的拓扑称为 \(p\)-\term[xituopu]{细拓扑}[fine topology]。
\end{definition}

这里先给出了开集的判据，仍须验证它确实定义一个拓扑。还要留意量词：补集只在一个点 \(x\) 处是薄的，并不根据这一定义就成为 \(x\) 的细邻域；定义要求对细开集中的每个点检查薄性。二者之间的联系将在最后一小节作为势论定理陈述。

\begin{proposition}\label{b1:prop:p-fine-topology-basic}
细开集组成一个比欧氏拓扑更细的拓扑。若 \(C_p(N)=0\)，则 \(\mathbb R^d\setminus N\) 为细开集。若 \(p>d\)，细拓扑恰好等于欧氏拓扑。
\end{proposition}
\begin{proof}
薄性对取子集保持。相对容量的次可加性与
\[
 (a+b)^q\le\max(1,2^{q-1})(a^q+b^q),\quad
 a,b\ge0,\quad q=\frac1{p-1},
\]
又说明有限个在 \(x\) 处薄的集合之并仍在 \(x\) 处薄。因此有限个细开集的交是细开的。若 \(x\in\bigcup_\lambda U_\lambda\)，选取一个包含 \(x\) 的 \(U_\lambda\)；并集的补集包含于这个 \(U_\lambda\) 的补集，所以也在 \(x\) 处薄。这证明任意并的封闭性，空集与全空间则直接满足条件。普通开集的补集在每个内点附近为空，故一定细开。

若 \(C_p(N)=0\)，固定 \(B(x,r)\)。将全空间中覆盖 \(N\) 的容许函数乘以在该球上等于 \(1\)、支集含于 \(B(x,2r)\) 的截断，得到
\[
 \operatorname{cap}_p\bigl(N\cap B(x,r),B(x,2r)\bigr)
 \le C_{r,p}\,C_p(N)=0.
\]
于是 \(N\) 在每一点都薄，结论成立。

最后设 \(p>d\)，并令 \(y\in B(x,r)\)。对 \(B(x,2r)\) 中覆盖单点 \(y\) 的任一容许函数，先补零到全空间。\cref{b1:thm:morrey-whole-space}给出其连续代表；它在容许开邻域上几乎处处至少为 \(1\)，故在 \(y\) 处也至少为 \(1\)。取 \(z=x+3re_1\)，补零代表在 \(z\) 处为零，因而
\[
 1\le |v(y)-v(z)|
 \le C(4r)^{1-d/p}\|\nabla v\|_{L^p}.
\]
所以相对单点容量至少为 \(cr^{d-p}\)，常数对 \(y\in B(x,r)\) 一致。若 \(x\in\overline E\)，每个小球都含有 \(E\) 的一点，\eqref{b1:eq:p-wiener-series}的每项便有统一正下界，不可能收敛。反之，若 \(x\notin\overline E\)，小尺度项最终为零。于是这个范围内，“在 \(x\) 处薄”等价于“在 \(x\) 的某个普通邻域中为空”，细开集恰好就是普通开集。
\end{proof}

当 \(1<p\le d\) 时则不同。由\cref{b1:sec:2601}的单点容量计算，\(\mathbb Q^d\) 的容量为零，因此 \(\mathbb R^d\setminus\mathbb Q^d\) 是细开集；它却不是欧氏开集，因为每个普通球都含有有理点。细拓扑所忽略的是某种容量很小的障碍，而不是距离上很远的点。

\begin{definition}\label{b1:def:p-fine-continuity}
设 \(v\colon\mathbb R^d\to\mathbb R\)。若对每个 \(\varepsilon>0\)，都存在含 \(x\) 的细开集 \(V\)，使
\[
 |v(y)-v(x)|<\varepsilon\quad(y\in V),
\]
则称 \(v\) 在 \(x\) 处 \(p\)-\term[xilianxu]{细连续}[finely continuous]。这就是定义域取细拓扑、值域取通常拓扑时的点态连续性。
\end{definition}

普通连续性蕴含细连续性。反过来没有这样的普遍蕴含；而拟连续性与细连续性的关系，也不是仅由“细拓扑更细”这一事实决定的。

\begin{theorem}\label{b1:thm:quasicontinuous-fine-continuity}
设 \(1<p<\infty\)。每个有限实值的 \(p\)-拟连续函数都在 \(p\)-拟处处细连续。特别地，\cref{b1:thm:sobolev-quasicontinuous-representative}所构造的 Sobolev 代表具有这一性质。
\end{theorem}
\begin{proofsketch}
需要的深层输入是以下集合刻画：若一个集合在加入容量任意小的适当开集后总能成为开集，则它可写成一个细开集与一个容量零集的并。这是非线性势论中的结果，其证明依赖 Cartan 型分离、Choquet 性质和相应的容量例外集控制；这些步骤不在本节展开。Björn、Björn 和 Latvala 的《The Cartan, Choquet and Kellogg properties for the fine topology on metric spaces》\cite[Theorem 1.4, Section 8]{Bjorn2018Cartan}给出该刻画与本定理的联系，这里仅使用其欧氏空间情形。

有了该输入，最后一步是直接的。拟连续函数 \(v\) 的每个超水平集 \(\{v>a\}\) 和次水平集 \(\{v<a\}\)，在加入使 \(v\) 于其补集连续的坏开集后，都成为开集。对所有有理数 \(a\) 应用集合刻画，合并所得的可数个容量零集为 \(N\)。在 \(N\) 外，这些水平集都由相应的细开集给出。又因 \(\mathbb R^d\setminus N\) 细开，围绕 \(v(x)\) 选取两个有理阈值，就得到细连续性所需的邻域。

因此，未展开的关键在于上述集合刻画，而不在有理阈值这一步。前一节在坏开集外的一致逼近本身，没有证明那些补集在其各点附近都是细邻域。
\end{proofsketch}

\subsection{Wiener 型思想}
\label{b1:subsec:260403}

薄性之所以采用指数 \(1/(p-1)\)，最终与 \(p\) 次梯度能量的非线性边界行为有关。为说明这点，先写出相应的方程。

\begin{definition}\label{b1:def:p-harmonic}
设 \(\Omega\subset\mathbb R^d\) 开。给定连续函数 \(u\in W_{\mathrm{loc}}^{1,p}(\Omega)\)，若
\[
 \int_\Omega |\nabla u|^{p-2}\nabla u\cdot\nabla\varphi\,\mathrm dx=0
 \quad(\varphi\in C_c^\infty(\Omega)),
\]
则称 \(u\) 为 \term[p-tiaohe-hanshu]{\(p\)-调和函数}[p-harmonic function]。当梯度为零时，向量 \( |\nabla u|^{p-2}\nabla u\) 约定为零。
\end{definition}

形式上，该方程写成 \(\Delta_pu=0\)，其中
\[
 \Delta_pu=\operatorname{div}\bigl(|\nabla u|^{p-2}\nabla u\bigr).
\]
\symidx{\(\Delta_pu\)：由 \(p\) 次梯度能量产生的非线性算子}
当 \(p=2\) 时，它是通常的 Laplace 方程；一般 \(p\) 时不能利用解的叠加。

设 \(\Omega\) 有界，且 \(g\in W^{1,p}(\Omega)\cap C(\overline\Omega)\)。在仿射类 \(g+W_0^{1,p}(\Omega)\) 上最小化 \(\int_\Omega|\nabla u|^p\)，确实有唯一极小元：对 \(u-g\) 用零边界 Poincaré 不等式，便由梯度界得到 Sobolev 有界性；弱子列、仿射类弱闭性和凸能量的弱下半连续性给出存在。严格凸性先给极小元的梯度相等，再由零边界 Poincaré 不等式得到函数相等。沿内部测试函数作变分，则得到上面的弱方程。这部分与\cref{b1:thm:sobolev-variational-minimum}使用同一直接法。

但极小元具有内部连续代表，尤其在 \(p\le d\) 时，不是 Morrey 不等式的推论，而是 \(p\)-调和方程的内正则性定理。以下边界判别把这项正则性作为经典非线性势论的输入，并用 \(u_g\) 表示该连续代表。

\begin{definition}\label{b1:def:p-dirichlet-regular-point}
设 \(\Omega\) 为有界开集，\(x_0\in\partial\Omega\)。若对每个 \(g\in W^{1,p}(\Omega)\cap C(\overline\Omega)\)，上述变分解都满足
\[
 \lim_{\substack{x\to x_0\\x\in\Omega}}u_g(x)=g(x_0),
\]
则称 \(x_0\) 为这个 \(p\)-Dirichlet 问题的\term[zhengze-bianjiedian]{正则边界点}[regular boundary point]。
\end{definition}

边界值在这里不是任意指定的可测代表值。闭包条件给出变分意义的边界约束，而正则边界点要求这一约束能够由内部连续解的极限恢复。

\term[wiener-panbie-zhunze]{Wiener 判别准则}[Wiener criterion]把这个分析问题转化为补集在每个小尺度上的容量问题。
\begin{theorem}[Wiener]\label{b1:thm:p-wiener-criterion}
设 \(d\ge2\)、\(1<p\le d\)，\(\Omega\subset\mathbb R^d\) 为有界开集，且 \(x_0\in\partial\Omega\)。则 \(x_0\) 是正则边界点，当且仅当 \(\mathbb R^d\setminus\Omega\) 在 \(x_0\) 处不是 \(p\)-薄的，即\eqref{b1:eq:p-wiener-series}中取 \(E=\mathbb R^d\setminus\Omega\) 所得级数发散。
\end{theorem}
\begin{proofsketch}
充分性需要把补集的相对容量转化为变分解在相邻尺度间的边界振荡衰减，再迭代这些估计。发散条件恰好使未被边界约束消除的振荡趋于零。这里需要方程的比较原理和容量控制的边界衰减估计，不能用普通球均值 Poincaré 不等式直接代替。

必要性更深：当补集薄时，借助相对容量的极小函数、非线性势的点态估计及 Cartan 性质，构造在该点不能恢复指定边值的解。Kilpeläinen 与 Malý 的《The Wiener test and potential estimates for quasilinear elliptic equations》\cite[THEOREM 1.1, Section 5]{Kilpelainen1994Wiener}给出上述完整范围，并说明 Cartan 性质、比较方法和 Wiener 判别之间的联系。本节不展开支撑必要性的非线性势估计，也不另行建立 Perron 方法及其与变分解的对应。定理在这里用于指出容量理论通向边界正则性的方向，而不是宣称这些势论工具已在前文证明。
\end{proofsketch}

这个判据可以立即解释两种相反的几何情形。若存在固定 \(a>0\)，使每个充分小的 \(B(x_0,r)\setminus\Omega\) 都包含一个半径为 \(ar\) 的球，那么任何容许函数在该小球上至少为 \(1\)。对外球中的容许函数用零边界 Poincaré 不等式，得到
\[
 \operatorname{cap}_p\bigl(B(x_0,r)\setminus\Omega,B(x_0,2r)\bigr)
 \ge c_a r^{d-p}.
\]
故 Wiener 级数发散。Lipschitz 图表的外侧含有这种按比例缩小的球，因此它的边界点符合该判据。相反，在穿孔球 \(\Omega=B(0,1)\setminus\{0\}\) 中，\(1<p\le d\) 时内边界点 \(0\) 附近的补集只有容量为零的单点，级数收敛，于是该点不正则。这说明一个拓扑边界点未必足以对变分解施加可见的点态边值。

\subsection{势论联系}
\label{b1:subsec:260404}

线性势论以调和函数及超调和函数组织比较原理；非线性情形保留比较的思想，但不能把超调和函数都当作局部 Sobolev 函数。

\begin{definition}\label{b1:def:p-superharmonic}
设 \(\Omega\subset\mathbb R^d\) 开。给定下半连续、在一个稠密子集上有限的函数 \(v\colon\Omega\to(-\infty,+\infty]\)。若对每个满足 \(\overline D\Subset\Omega\) 的开集 \(D\)，以及每个在 \(D\) 上 \(p\)-调和且属于 \(C(\overline D)\) 的 \(h\)，都有
\[
 h\le v\text{ 于 }\partial D
 \quad\Longrightarrow\quad h\le v\text{ 于 }D,
\]
则称 \(v\) 为 \term[p-chaotiaohe-hanshu]{\(p\)-超调和函数}[p-superharmonic function]。
\end{definition}

定义允许取 \(+\infty\)，也不先要求 \(v\in W_{\mathrm{loc}}^{1,p}\)。因此不能未经论证就把每个这样的 \(v\) 代入弱方程并写出其梯度。超调和函数的截断正则性、下半连续代表与弱不等式的对应，都属于这套理论需要另外建立的内容。

\term[cartan-xingzhi]{Cartan 性质}[Cartan property]说明，薄性可以由一个势函数在该点与沿障碍集合的值之间的严格间隙检测。
\begin{theorem}[Cartan]\label{b1:thm:p-cartan-characterization}
设 \(d\ge2\)、\(1<p\le d\)，且 \(x_0\in\overline E\setminus E\)。集合 \(E\) 在 \(x_0\) 处是 \(p\)-薄的，当且仅当在 \(x_0\) 的某个普通邻域中存在非负 \(p\)-超调和函数 \(v\)，使
\[
 v(x_0)<
 \liminf_{\substack{x\to x_0\\x\in E}}v(x).
\]
相应地，对任意 \(U\ni x_0\)，\(U\) 包含 \(x_0\) 的一个细开邻域，当且仅当 \(\mathbb R^d\setminus U\) 在 \(x_0\) 处是 \(p\)-薄的。细拓扑也正是使所有局部 \(p\)-超调和函数连续的最弱拓扑，允许函数值为 \(+\infty\) 时在值域采用通常的扩充实数拓扑。
\end{theorem}
\begin{proofsketch}
由势函数的比较性质和容量估计，可以从严格间隙推出薄性。反向需要将小尺度的容量障碍组合成超调和势，同时控制它在 \(x_0\) 的值。在线性情形可以使用势的叠加；一般 \(p\) 时必须用非线性势估计与比较方法替代，不能把线性证明中的级数逐项照搬。

Kilpeläinen 与 Malý 的《The Wiener test and potential estimates for quasilinear elliptic equations》\cite[THEOREM 1.3, Section 5]{Kilpelainen1994Wiener}建立这种点态分离；由超调和函数定义的拓扑与薄性邻域的等价刻画，可参阅 Heinonen、Kilpeläinen 和 Martio 的《Fine topology and quasilinear elliptic equations》\cite[3.2. THEOREM]{Heinonen1989Fine}。这些结论包含了前文尚未证明的邻域提取步骤。若一个集合已经含有细开邻域，其补集薄只需单调性；从仅在一个点成立的薄性反过来提取细开邻域，才是此处需要势论的方向。
\end{proofsketch}

可以用一个公式看清这些深层估计增加了什么。给定正 Radon 测度 \(\sigma\)，相关的局部非线性势积分具有形式
\[
 \int_0^R
 \left(\frac{\sigma(B(x,r))}{r^{d-p}}\right)^{1/(p-1)}
 \frac{\mathrm dr}{r}.
\]
它把测度在所有小尺度上的集中程度累积起来。完整理论先为容量障碍构造极小函数和相应测度，再用这类势积分估计超调和函数的点值；由此才能控制 Cartan 分离和边界极限。本节没有证明该测度构造及点态估计，公式本身也不能代替它们。

到这里，三种“忽略异常”的方式便有了明确分工。几乎处处只识别积分等价类；拟处处用容量限制可忽略的点集；细拓扑则让容量所控制的局部障碍进入邻域结构。前两者的基本 Sobolev 结论已在本章建立，后者与非线性方程之间的深层联系在本节给出了准确入口。继续阅读时，可先沿 Cartan 性质与 Wiener 判别的两篇原始论文补足势估计和比较方法，再回看拟连续代表为何能在拟处处成为细连续函数。


\begin{chaptersummary}
容量所度量的是在集合邻域保持函数大小所需的能量，而不是单独在集合上任意填写点值的代价。开集上的几乎处处约束与开外正则化，使这一概念能先于代表的选择建立。次可加性随后把快速逼近中的坏集汇合起来，得到拟连续代表及其拟处处唯一性。

\ProseParagraphBreak

零容量集一定零测，反向一般不成立。小球缩放、临界对数过渡与球平均的累积，把这种区别同 Hausdorff 维数联系起来；等于临界维数时，还要考察比单一维数更细的覆盖信息。拟连续性给出了小容量开集外的相对连续性，并不自动说明所有点都通常连续，也不能单靠普通微分定理确定其容量异常集。

\ProseParagraphBreak

细拓扑让邻域概念也适应容量尺度。它与边界可达值、超解及非线性势估计之间有深层联系，本章在相应证明概要中列明了尚未展开的输入。下一章回到可以直接由分布导数测度描述的对象：当弱导数不再有可积密度，而只是一份有限测度时，便进入多维有界变差函数理论。
\end{chaptersummary}



容量的单点行为随指数跨过维数而改变，见\cref{b1:tab:visual-t14}。
% Source fingerprint: 88b089aaadd01ba58e73e9a153026309df6f8ac9e9e35b19f12fcd69aca8859f
% T14: raw TeX cells preserved; no numeric highlighting.
\begin{table}[htbp]
\centering\small\renewcommand{\arraystretch}{1.18}\setlength{\tabcolsep}{4pt}
\caption[容量与维数的三个指数区间]{容量与维数的三个指数区间。前三行讨论本章全空间非齐次容量。最后一行是第26.4节建立的相对容量；两种能量的尺度不同。}
\label{b1:tab:visual-t14}
\begin{tabularx}{\linewidth}{@{}>{\raggedright\arraybackslash}p{3.1cm}>{\raggedright\arraybackslash}p{4.7cm}>{\raggedright\arraybackslash}X@{}}
\toprule
指数 & 小球与单点 & 零测集与零容量集 \\
\midrule
\(1<p<d\) & \(C_p(B_r)\le C(r^d+r^{d-p})\)；单点容量零 & 单点给出非空零容量集；平面片可为零测而正容量 \\
\(p=d>1\) & 需用对数型截断证明单点容量零；上一行幂估计本身不够 & 可数集容量零，但不能把所有 Lebesgue 零测集都忽略 \\
\(p>d\) & 单点容量有统一正下界 & 每个非空集合容量为正；只有空集可为零容量集 \\
相对容量的球对 & \(\operatorname{cap}_p(B_r,B_{2r})\) 与 \(r^{d-p}\) 可比 & 纯梯度能量，不能与全空间非齐次容量混为同一数值 \\
\bottomrule
\end{tabularx}
\end{table}


\BookExercises


% \chapref{b1:ch:26}习题临时稿；不另设 BookExercises 入口。
% 全部题目按已建立的欧氏 Sobolev 工具自拟，不冒认未读教材的原题或题号。
% 固定非齐次能量 E_p(u)=int(|u|^p+|grad u|^p)，梯度取欧氏模，1<p<infty。
% 层次：1 为基础辨析；2--4 为容量比较与约束逼近，其中3为提高题；
% 5--6 为集合极限与反例；7--9 为局部化及细拓扑；10--11 为迹与边界空间。
% 不使用一般 Choquet 定理、容量版 Lebesgue 点定理或任意域的零拟处处刻画。

\begin{exercise}\label{b1:ex:26-basic-capacity-properties}
直接利用容量的定义和本节已经证明的结论，完成下列基本辨析。
\begin{enumerate}
 \item 若 \(A\subseteq B\subseteq\mathbb R^d\)，证明 \(C_p(A)\le C_p(B)\)。
 \item 对集合列 \((A_j)\)，写出容量的可列次可加不等式，并说明它不是可列可加等式。
 \item 选取一个固定的光滑截断函数，重新得到
 \[
  C_p(B(x,r))\le C_{d,p}(r^d+r^{d-p}).
 \]
 \item 当 \(1<p\le d\) 时，分别举出非空的零容量 Lebesgue 零测集与正容量 Lebesgue 零测集。
 当 \(p>d\) 时，证明不存在非空零容量集合。说明这些结论与 \(m_d^*(E)\le C_p(E)\) 的关系。
\end{enumerate}
\end{exercise}
% 来源：自拟基础辨析；只调用本节定义、次可加性、球估计和已证反例。
\begin{solution}
第一项由开邻域类的包含关系直接得到：每个包含 \(B\) 的开集也包含 \(A\)，先对容许函数、再对开邻域取下确界即可。

第二项是
\[
 C_p\left(\bigcup_{j=1}^{\infty}A_j\right)
 \le\sum_{j=1}^{\infty}C_p(A_j).
\]
左端即使对两两不交的集合也未必等于右端；本章后面的两点集合给出严格不等式。因此这里是可列次可加性，而不是测度的可列可加性。

取 \(\eta\in C_c^\infty(B(0,2))\)，使 \(0\le\eta\le1\) 且 \(\eta=1\) 于 \(B(0,1)\) 的一个邻域。函数
\[
 u_{x,r}(y)=\eta\left(\frac{y-x}{r}\right)
\]
在 \(B(x,r)\) 上容许，而变量代换给出
\[
 E_p(u_{x,r})
 =r^d\|\eta\|_p^p+r^{d-p}\|\nabla\eta\|_p^p.
\]
取只依赖 \(d,p,\eta\) 的常数便得第三项。

当 \(1<p\le d\) 时，由\cref{b1:prop:sobolev-point-capacity}，单点是非空、Lebesgue 零测且容量为零的集合。此时必有 \(d\ge2\)，平面片
\[
 [0,1]^{d-1}\times\{0\}
\]
的 \(d\) 维 Lebesgue 测度为零，但由\secref{b1:sec:2601}中竖线上的一维点值估计具有正容量。

当 \(p>d\) 时，对任意非空集合 \(E\) 取 \(x\in E\)。由单调性及\cref{b1:prop:sobolev-point-capacity}，
\[
 C_p(E)\ge C_p(\{x\})>0.
\]
因此不存在非空零容量集合；空集满足 \(C_p(\varnothing)=0\)。单点仍是零测而正容量的例子，特别包括 \(d=1\) 时本题允许的全部指数。不等式 \(m_d^*(E)\le C_p(E)\) 只说明零容量蕴含零测；它并未断言零测蕴含零容量，所以没有矛盾。
\end{solution}

\begin{exercise}\label{b1:ex:26-affine-capacity-comparison}
设 \(A\colon\mathbb R^d\rightarrow\mathbb R^d\) 为可逆线性映射，
\(b\in\mathbb R^d\)。证明对任意集合 \(E\subseteq\mathbb R^d\)，都有
\[
 |\det A|\min\{1,\|A\|^{-p}\}C_p(E)
 \le C_p(AE+b)
 \le|\det A|\max\{1,\|A^{-1}\|^p\}C_p(E).
\]
这里算子范数由欧氏范数诱导。由此证明容量在平移与正交变换下不变，并写出
\(E\mapsto tE\)、\(t>0\) 时的双边估计。解释为什么这些估计不同于齐次容量的单一幂次伸缩公式。
\end{exercise}
% 来源：自拟；把线性换元与非齐次能量结合，训练零阶项和梯度项的不同尺度。
\begin{solution}
先对光滑函数令 \(v(y)=u\bigl(A^{-1}(y-b)\bigr)\)。链式法则及换元给出
\[
 E_p(v)=|\det A|\int_{\mathbb R^d}
       \bigl(|u(x)|^p+|A^{-\mathsf T}\nabla u(x)|^p\bigr)\,\mathrm dx.
\]
对任意向量 \(\xi\)，有
\[
 \|A\|^{-1}\cdot|\xi|\le |A^{-\mathsf T}\xi|
                         \le\|A^{-1}\|\cdot|\xi|.
\]
因此所显示的两个常数分别给出 \(E_p(v)\) 相对于 \(E_p(u)\) 的下界和上界。
用全空间光滑稠密性逼近 \(u\)，这些估计同时保证复合函数及其候选弱梯度收敛，
所以同一公式和估计适用于所有 \(u\in W^{1,p}(\mathbb R^d)\)。

仿射映射把开集 \(G\) 双射到开集 \(AG+b\)，并保持 Lebesgue 零测集。
因而 \(u\ge1\) 几乎处处于 \(G\)，当且仅当相应的 \(v\ge1\) 几乎处处于 \(AG+b\)。
容许函数与开邻域在这两个映射下均一一对应。先对函数、再对开邻域取下确界，便得题设双边界；
不要求 \(E\) 可测。

当 \(A\) 正交时，两个常数都等于 \(1\)；取 \(A=I\) 同时得到平移不变性。
当 \(A=tI\) 时，估计化为
\[
 \min\{t^d,t^{d-p}\}C_p(E)
 \le C_p(tE)
 \le\max\{t^d,t^{d-p}\}C_p(E).
\]
零阶积分乘以 \(t^d\)，梯度积分乘以 \(t^{d-p}\)；一般不能把这两个因子合成一个相同的幂。
特别地，可逆仿射变换保持零容量集，但不保持一般集合的容量数值。
\end{solution}

\begin{optionalexercise}\label{b1:ex:26-one-dimensional-exact-capacity}
在 \(\mathbb R\) 上取 \(p=2\)，并设 \(L>0\)。计算
\[
 C_2(\{0\}),\quad C_2([0,L]),\quad C_2(\{0,L\}).
\]
证明两点集合的容量严格小于两个单点容量之和，由此说明 \(C_2\) 不是可加测度。
\end{optionalexercise}
% 来源：自拟；以一维能量配平方求精确值，区别非齐次容量与点质量的可加性。
\begin{solution}
先证半直线上的能量界。若 \(f\in H^1(0,\infty)\)，取其在有限区间上的绝对连续代表。
因为 \(f\in L^2(0,\infty)\)，可以选 \(R_j\to\infty\)，使 \(f(R_j)\to0\)。
在 \([0,R_j]\) 上展开平方并积分，有
\[
 \int_0^{R_j}(|f'|^2+|f|^2)\,\mathrm dx
 =\int_0^{R_j}|f'+f|^2\,\mathrm dx+|f(0)|^2-|f(R_j)|^2.
\]
令 \(j\to\infty\)，得到能量至少为 \(|f(0)|^2\)。反射后，左半直线也有相同结论。
在一维，\(H^1\) 函数的局部绝对连续代表连续；
若它在某点的开邻域几乎处处至少为 \(1\)，则该点值至少为 \(1\)。
因此，任意单点容许函数的总能量至少为 \(2\)。
函数 \(e^{-|x|}\) 的能量恰为 \(2\)，而
\((1+\varepsilon)e^{-|x|}\) 在原点的某个邻域大于 \(1\)。
让 \(\varepsilon\downarrow0\)，即得
\[
 C_2(\{0\})=2.
\]

对区间 \([0,L]\)，内部的零阶能量至少为 \(L\)，两侧的能量至少各为 \(1\)。
函数
\[
 v(x)=
 \begin{cases}
 e^x,&x<0,\\
 1,&0\le x\le L,\\
 e^{L-x},&x>L
 \end{cases}
\]
连续、分段光滑且属于 \(H^1(\mathbb R)\)，能量为 \(L+2\)。
用 \((1+\varepsilon)v\) 作开邻域容许函数，得到
\[
 C_2([0,L])=L+2.
\]

两点之间不再要求函数保持大于等于 \(1\)。令
\[
 h(x)=\frac{\cosh(x-L/2)}{\cosh(L/2)},\quad 0\le x\le L.
\]
它满足 \(h''=h\)、\(h(0)=h(L)=1\)，以及
\(h'(L)=-h'(0)=\tanh(L/2)\)。
任取两点集合的容许函数 \(u\)，在 \([0,L]\) 上写 \(u=h+w\)。
此时 \(w(0),w(L)\ge0\)，分部积分给出
\[
 \begin{aligned}
 \int_0^L(|u'|^2+|u|^2)\,\mathrm dx
 &=\int_0^L(|h'|^2+|h|^2)\,\mathrm dx
       +2[h'w]_0^L+\int_0^L(|w'|^2+|w|^2)\,\mathrm dx\\
 &\ge[h'h]_0^L=2\tanh(L/2).
 \end{aligned}
\]
加上左右两侧的能量下界，得到总能量至少为 \(2+2\tanh(L/2)\)。
将上面的 \(v\) 在中间区间改成 \(h\)，所得连续分段光滑函数达到这个能量；
再乘 \(1+\varepsilon\)，便在两个端点的开邻域满足容许条件。因此
\[
 C_2(\{0,L\})=2+2\tanh(L/2)<4=C_2(\{0\})+C_2(\{L\}).
\]
这里试探函数本身在端点取值恰为 \(1\)，乘以 \(1+\varepsilon\) 的步骤保证了容量定义要求的开邻域约束。
\end{solution}

\begin{exercise}\label{b1:ex:26-compact-smooth-admissibility}
设 \(K\subset\mathbb R^d\) 为紧集。证明
\[
 C_p(K)=\inf\left\{\,E_p(\varphi):
   \varphi\in C_c^\infty(\mathbb R^d),\
   0\le\varphi\le1,\
   \varphi=1\ \text{于 }K\text{ 的某个开邻域}\,\right\}.
\]
再在 \(1<p\le d\) 时，分别用一个无界可数集和一个有界稠密可数集，说明这个公式不能原样推广到任意集合。
\end{exercise}
% 来源：自拟；全空间稠密性的约束保持版本，新增远处截断和磨光尺度的选择。
\begin{solution}
记右端为 \(I(K)\)。其中每个光滑函数都是容量定义中的容许函数，所以 \(C_p(K)\le I(K)\)。
反向不等式不能仅由范数稠密性推出，因为任意光滑逼近未必保留邻域上的约束。

固定 \(\varepsilon>0\)。选开集 \(G\supseteq K\) 及容许函数 \(u\)，使
\[
 E_p(u)<C_p(K)+\varepsilon.
\]
这里紧集容量有限，可先用一个光滑截断函数得到有限上界。
将 \(u\) 截断到 \([0,1]\)，能量不增加，且 \(u=1\) 几乎处处于 \(G\)。
取 \(0\le\chi_R\le1\)，使 \(\chi_R=1\) 于 \(B(0,R)\)，
\(\operatorname{supp}\chi_R\subset B(0,2R)\)，并使
\(|\nabla\chi_R|\le C/R\)。乘积公式及 \(L^p\) 尾部收敛给出
\[
 \chi_R\cdot u\longrightarrow u\quad\text{于 }W^{1,p}(\mathbb R^d).
\]
确实，梯度差由 \((\chi_R-1)\cdot\nabla u\) 和 \(u\cdot\nabla\chi_R\) 两项组成；
前者由尾部积分趋零，后者的 \(L^p\) 范数至多为 \(C\|u\|_p/R\)。
选足够大的 \(R\)，使 \(K\subset B(0,R)\) 且
\(E_p(\chi_R\cdot u)<E_p(u)+\varepsilon\)。

紧集 \(K\) 到 \(\bigl(G\cap B(0,R)\bigr)^c\) 有正距离。
因此可选足够小的磨光尺度 \(\delta\)，使
\(\varphi=\rho_\delta*(\chi_R\cdot u)\) 在 \(K\) 的一个开邻域恒等于 \(1\)。
函数 \(\varphi\) 光滑紧支集，且 \(0\le\varphi\le1\)。
磨光对函数和向量梯度的 \(L^p\) 范数均为压缩，故
\[
 E_p(\varphi)\le E_p(\chi_R\cdot u)<C_p(K)+2\varepsilon.
\]
令 \(\varepsilon\downarrow0\)，证明了公式。

当 \(1<p\le d\) 时，可数集容量为零。对无界集合 \(\mathbb Z^d\)，
不存在紧支集函数能在其全体点的邻域恒等于 \(1\)，故对应的光滑下确界为无穷。
即使只考虑有界集合，也不能取消紧性条件。
取 \(E=\mathbb Q^d\cap B(0,1)\)。任何连续函数若在 \(E\) 的邻域恒等于 \(1\)，
则由稠密性在 \(\overline{B(0,1)}\) 上恒等于 \(1\)，能量至少为 \(m_d\bigl(B(0,1)\bigr)\)。
但 \(C_p(E)=0\)。后一个例子说明，连续容许函数会额外感受到集合的闭包。
\end{solution}

\begin{exercise}\label{b1:ex:26-capacity-hulls-and-decreasing-sets}
证明下列结论，并辨明最后一个极限为何不与第二项矛盾。
\begin{enumerate}
 \item\label{b1:item:26-capacity-gdelta-hull}
 每个集合 \(E\subseteq\mathbb R^d\) 都包含于一个 \(G_\delta\) 集 \(N\)，使
 \(C_p(N)=C_p(E)\)。特别地，零容量集有 Borel 零容量包络。
 \item\label{b1:item:26-capacity-decreasing-compact}
 若紧集列满足 \(K_1\supseteq K_2\supseteq\cdots\)，则
 \[
  C_p\left(\bigcap_{j=1}^\infty K_j\right)=\lim_{j\to\infty}C_p(K_j).
 \]
 \item\label{b1:item:26-capacity-decreasing-open}
 在 \(\mathbb R\) 上取 \(p=2\)。证明 \(G_j=(0,1/j)\) 递减到空集，而
 \(C_2(G_j)\to2\)。
\end{enumerate}
\end{exercise}
% 来源：自拟；从开外正则性推出限定的集合极限性质，不调用一般 Choquet 定理。
\begin{solution}
先设 \(C_p(E)<\infty\)。由外正则性选开集 \(O_j\supseteq E\)，使
\(C_p(O_j)<C_p(E)+1/j\)，并令 \(N=\bigcap_jO_j\)。
单调性给出
\[
 C_p(E)\le C_p(N)\le C_p(O_j)<C_p(E)+1/j.
\]
令 \(j\to\infty\)，得到第一项。若 \(C_p(E)=\infty\)，取 \(N=\mathbb R^d\) 即可。
这里仅断言存在同容量的 Borel 包络，没有断言容量在所有 Borel 集上可加。

令 \(K=\bigcap_jK_j\)。单调性给出 \(C_p(K)\le\lim_jC_p(K_j)\)。
任取包含 \(K\) 的开集 \(O\)。必有某个 \(j_0\)，使 \(K_j\subseteq O\) 对所有
\(j\ge j_0\) 成立：否则可以选 \(x_j\in K_j\setminus O\)，
由 \(K_1\) 紧性取收敛子列。固定任意 \(i\) 后，该子列的尾部落在闭集
\(K_i\setminus O\) 中，故其极限属于所有 \(K_i\) 却不属于 \(O\)，矛盾。
于是
\[
 \lim_jC_p(K_j)\le C_p(O).
\]
对开集 \(O\supseteq K\) 取下确界，即得第二项。

最后，\((0,1/j)\) 包含单点 \(1/(2j)\)，并包含于闭区间 \([0,1/j]\)。
由\cref{b1:ex:26-one-dimensional-exact-capacity}及平移不变性，
\[
 2\le C_2\bigl((0,1/j)\bigr)\le2+1/j.
\]
这些开区间的交为空集，容量却趋于 \(2\)。
它们不是紧集；其闭包都有共同点 \(0\)，而这个点未进入任何 \(G_j\)。
第二项证明中将极限留在每个闭集内的步骤，在这里恰好失效。
\end{solution}

\begin{exercise}\label{b1:ex:26-strong-convergence-without-full-pointwise-limit}
设 \(1<p<d\)。构造 \(u_n\in C_c^\infty(\mathbb R^d)\)，使
\[
 0\le u_n\le1,\quad u_n\longrightarrow0\ \text{于 }W^{1,p}(\mathbb R^d),
\]
但对每个 \(x\in[0,1]^d\)，都有
\[
 \liminf_{n\to\infty}u_n(x)=0,\quad
 \limsup_{n\to\infty}u_n(x)=1.
\]
由此说明范数收敛抽取拟处处收敛子列的结论，不能把“子列”删去。
\end{exercise}
% 来源：自拟移动峰；在每一层用有限个小立方体覆盖全部点，并插入零项固定下极限。
\begin{solution}
选取 \(0\le\eta\le1\) 的光滑函数，使它在 \([-1/2,1/2]^d\) 上恒等于 \(1\)，
支集包含于 \((-1,1)^d\)。对第 \(k\) 层的每个小立方体中心
\[
 c_{k,j}=2^{-k}\bigl(j+(1/2,\ldots,1/2)\bigr),
 \quad j\in\{0,\ldots,2^k-1\}^d,
\]
定义
\[
 v_{k,j}(x)=\eta\bigl(2^k(x-c_{k,j})\bigr).
\]
伸缩计算给出与 \(j\) 无关的界
\[
 E_p(v_{k,j})=
 2^{-kd}\|\eta\|_p^p+
 2^{-k(d-p)}\|\nabla\eta\|_p^p\longrightarrow0.
\]
按 \(k=1,2,\ldots\) 的次序列出每一层的全部 \(v_{k,j}\)，
并在相邻两项之间插入一个恒零函数，所得序列记为 \(u_n\)。
每层只有有限项，所以当 \(n\to\infty\) 时，非零项所在的层数趋于无穷。
上式因此证明了整列的 Sobolev 强收敛。

固定 \(x\in[0,1]^d\)。每一层都有一个闭小立方体包含 \(x\)，
它对应的函数在 \(x\) 处等于 \(1\)；立方体公共边界上的点也满足这一点。
故 \(u_n(x)=1\) 无穷多次，而插入的零项使 \(u_n(x)=0\) 也无穷多次。
结合 \(0\le u_n\le1\)，得到所要求的上下极限。
最后，由测度下界 \(C_p([0,1]^d)\ge1\)，这个不收敛集合不是零容量集。
所有 \(u_n\) 都是通常连续的代表，失败并非来自在零测集上任意改值。
\end{solution}

\begin{exercise}\label{b1:ex:26-local-quasicontinuous-gluing}
设 \(\Omega\subseteq\mathbb R^d\) 为开集。
对有限值函数 \(v\colon\Omega\rightarrow\mathbb R\)，把拟连续条件理解为：
每个 \(\varepsilon>0\) 都可找到全空间开集 \(G\)，满足 \(C_p(G)<\varepsilon\)，
且 \(v|_{\Omega\setminus G}\) 相对连续。
\begin{enumerate}
 \item\label{b1:item:26-local-quasicontinuous-cover}
 若 \(\Omega=\bigcup_jU_j\)，各 \(U_j\) 开，且 \(v|_{U_j}\) 均满足上述条件，
 则 \(v\) 在 \(\Omega\) 上拟连续。
 \item\label{b1:item:26-local-quasicontinuous-sobolev}
 每个 \(u\in W_{\mathrm{loc}}^{1,p}(\Omega)\) 都有有限值 Borel 拟连续代表，
 且两个这样的代表在 \(\Omega\) 上拟处处相等。
\end{enumerate}
\end{exercise}
% 来源：自拟局部化；使用光滑划分与全空间代表，不假定局部函数有有限的全空间能量。
\begin{solution}
对第一项，给定 \(\varepsilon>0\)，为第 \(j\) 个限制函数选坏开集 \(G_j\)，
使 \(C_p(G_j)<\varepsilon2^{-j-1}\)。集合 \(G=\bigcup_jG_j\) 的容量小于 \(\varepsilon\)。
固定 \(x\in\Omega\setminus G\)，选 \(j\) 使 \(x\in U_j\)。
在这个相对开邻域内，\(v\) 是连续函数
\(v|_{U_j\setminus G_j}\) 的限制，故 \(v\) 在 \(\Omega\setminus G\) 上连续。
连续性是局部性质，而统一的坏集由次可加性控制。

对第二项，取\cref{b1:thm:smooth-partition-unity}给出的非负光滑划分
\((\chi_j)\)，使其支集均紧含于 \(\Omega\)，在 \(\Omega\) 内局部有限，且
\(\sum_j\chi_j=1\)。
对每个 \(j\)，另取 \(\eta_j\in C_c^\infty(\Omega)\)，使
\(\eta_j=1\) 于 \(\operatorname{supp}\chi_j\) 的邻域。
函数 \(\eta_j\cdot u\) 的支集远离边界，由乘积公式和内部零延拓，
其全空间零延拓 \(u_j\) 属于 \(W^{1,p}(\mathbb R^d)\)。
为它选一个有限值 Borel 拟连续代表 \(\widetilde u_j\)，并在 \(\Omega\) 上令
\[
 v(x)=\sum_j\chi_j(x)\cdot\widetilde u_j(x).
\]
局部有限性保证这个和在每点附近都是有限和，因此有限且 Borel 可测。
在共同的满测度集合上，各 \(\widetilde u_j=u_j\)，于是
\[
 v=\sum_j\chi_j\cdot\eta_j\cdot u=\sum_j\chi_j\cdot u=u
 \quad\text{几乎处处于 }\Omega.
\]
再为各 \(\widetilde u_j\) 选容量小于 \(\varepsilon\cdot2^{-j-1}\) 的坏开集。
它们的并以外，每项 \(\chi_j\cdot\widetilde u_j\) 连续，且和仍局部有限，
故 \(v\) 拟连续。

设 \(v,w\) 是 \(u\) 的两个有限值拟连续代表。
对每个 \(j\)，将 \(\chi_j\cdot(v-w)\) 在 \(\Omega\) 外补零。
这个函数全空间拟连续：在所选坏开集外，乘积在 \(\Omega\) 内连续；
而 \(\operatorname{supp}\chi_j\) 紧含于 \(\Omega\)，所以在每个边界点附近乘积恒为零。
它几乎处处等于零，是零 Sobolev 类的拟连续代表。
由\cref{b1:thm:quasicontinuous-uniqueness}，存在容量零集 \(N_j\)，使
\(\chi_j\cdot(v-w)=0\) 于 \(N_j^c\)。
删去 \(N=\bigcup_jN_j\) 后，对每点至少有一个 \(\chi_j>0\)，从而 \(v=w\)。
证明没有把仅具局部能量的 \(u\) 直接代入全空间超水平估计。
\end{solution}

\begin{optionalexercise}\label{b1:ex:26-zero-capacity-thinness}
设 \(N\subseteq\mathbb R^d\) 且 \(C_p(N)=0\)。
\begin{enumerate}
 \item 证明 \(N\) 在每一点都是 \(p\)-薄的，并由此证明 \(\mathbb R^d\setminus N\) 是细开集。
 \item 设 \(d\ge2\)、\(1<p\le d\)，并令
 \(\Omega=B(0,1)\setminus\{0\}\)。直接判断补集在 \(0\) 处的 Wiener 级数是收敛还是发散；再利用选读的 Wiener 判据判断 \(0\) 是否为正则边界点。
\end{enumerate}
\end{optionalexercise}
% 来源：自拟；训练全空间零容量到相对薄性的截断接口及穿孔球判据。
\begin{solution}
固定 \(x\) 与 \(r>0\)。取一个在 \(B(x,r)\) 上等于一、支集含于 \(B(x,2r)\) 的光滑截断 \(\eta\)。若 \(u\) 是覆盖 \(N\) 的全空间容许函数，则 \(\eta\cdot u\) 在 \(B(x,2r)\) 中零边界，并覆盖 \(N\cap B(x,r)\)。乘积估计给出
\[
 \operatorname{cap}_p\bigl(N\cap B(x,r),B(x,2r)\bigr)
 \le C_{x,r,p}\cdot E_p(u).
\]
依次对容许函数和开邻域取下确界，右端为零。故薄性定义中的每一项都为零，\(N\) 在每一点都薄。对 \(N\) 的补集应用细开集定义，即得第一项。

对穿孔球，在所有充分小的半径上，
\[
 (\mathbb R^d\setminus\Omega)\cap B(0,r)=\{0\}.
\]
单点全空间容量为零，第一项的论证说明相应相对容量也为零，故 Wiener 级数收敛，补集在 \(0\) 处是薄的。选读的 Wiener 判据因此说明 \(0\) 不是正则边界点。
\end{solution}

\begin{optionalexercise}\label{b1:ex:26-exterior-corkscrew-wiener}
设 \(d\ge2\)、\(1<p\le d\)，\(\Omega\subset\mathbb R^d\) 为有界开集，且 \(x_0\in\partial\Omega\)。假设存在 \(a\in(0,1/4)\) 与 \(r_0>0\)，使每个 \(0<r<r_0\) 都有
\[
 B(y_r,ar)\subseteq B(x_0,r)\setminus\Omega.
\]
用零边界 Poincaré 不等式证明
\[
 \operatorname{cap}_p\bigl(B(x_0,r)\setminus\Omega,B(x_0,2r)\bigr)
 \ge c_{d,p,a}r^{d-p}.
\]
由此验证 Wiener 级数发散；再利用选读的 Wiener 判据说明 \(x_0\) 正则。
\end{optionalexercise}
% 来源：自拟；把外部软木塞条件直接转成 Wiener 级数的尺度下界。
\begin{solution}
设 \(v\in W_0^{1,p}(B(x_0,2r))\) 在包含 \(B(y_r,ar)\) 的开集上几乎处处至少为一。截断后可设 \(0\le v\le1\)，于是
\[
 \omega_d a^dr^d
 \le\int_{B(x_0,2r)}|v|^p
 \le (4r)^p\int_{B(x_0,2r)}|\nabla v|^p,
\]
其中第二个不等式是零边界 Poincaré 不等式。对所有容许函数取下确界，并用内集单调性，得到题设下界。

当 \(r=2^{-j}<r_0\) 时，Wiener 级数的第 \(j\) 项至少为正常数
\(c_{d,p,a}^{1/(p-1)}\)。因此级数发散，补集在 \(x_0\) 处不是薄的；选读的 Wiener 判据给出 \(x_0\) 正则。
\end{solution}

\begin{exercise}\label{b1:ex:26-quasicontinuous-plane-trace}
设 \(d\ge2\)，\(P=\mathbb R^{d-1}\times\{0\}\)，以 \(\sigma\) 表示 \(P\) 上的
\((d-1)\) 维 Lebesgue 测度，\(\sigma^*\) 表示相应外测度。
\begin{enumerate}
 \item\label{b1:item:26-plane-capacity-control}
 证明存在仅依赖 \(p\) 的常数 \(A_p\)，使对任意集合 \(E\subseteq\mathbb R^d\)，
 \[
  \sigma^*(E\cap P)\le A_p\cdot C_p(E).
 \]
 \item\label{b1:item:26-plane-trace-representative}
 设 \(u\in W^{1,p}(\mathbb R^d)\)，\(\widetilde u\) 是任一有限值拟连续代表。
 证明 \(\widetilde u|_P\) 与上半空间限制的 \(L^p\) 迹
 \(\operatorname{Tr}_P u\) 在 \(\sigma\)-几乎处处意义下相等。
 因而迹虽然不能由任意 Lebesgue 代表的平面点值给出，却可由拟连续代表给出。
\end{enumerate}
\end{exercise}
% 来源：自拟接口题；ACL 的一维点值界控制平面测度，再以拟处处逼近辨认迹。
\begin{solution}
对绝对连续函数 \(f\colon[0,1]\rightarrow\mathbb R\)，先由基本定理对终点平均，得到
\[
 |f(0)|\le\int_0^1|f(t)|\,\mathrm dt+\int_0^1|f'(t)|\,\mathrm dt.
\]
Hölder 不等式及标量幂次估计给出
\[
 |f(0)|^p\le2^{p-1}\int_0^1\bigl(|f(t)|^p+|f'(t)|^p\bigr)\,\mathrm dt.
\]
沿几乎每条竖线使用 ACL 代表，再对 \(x'\in\mathbb R^{d-1}\) 积分，
便得到平面迹的界
\[
 \|\operatorname{Tr}_P u\|_{L^p(P)}^p\le A_p\cdot E_p(u),
 \quad A_p=2^{p-1}.
\]
这个迹就是\cref{b1:thm:half-space-lp-trace}的端点迹；
该不等式也说明它对全空间 Sobolev 范数连续。

设 \(G\supseteq E\) 开，\(u\) 是 \(G\) 的容量容许函数。
Fubini 定理允许先排除一个横向零测集，使 \(u\ge1\) 在每个相关竖截面内几乎处处成立，
并使这些竖线上的 ACL 代表都存在。
若 \((x',0)\in G\)，开性保证该竖截面含有原点的一个区间。
沿线连续性于是推出 \(\operatorname{Tr}_P u(x')\ge1\)。
所以
\[
 \sigma^*(E\cap P)\le\sigma(G\cap P)
 \le\|\operatorname{Tr}_P u\|_{L^p(P)}^p
 \le A_p\cdot E_p(u).
\]
依次对容许函数及开邻域取下确界，得到第一项。

由全空间光滑稠密性及\cref{b1:thm:sobolev-quasi-everywhere-subsequence}，
可选 \(\varphi_j\in C_c^\infty(\mathbb R^d)\)，使
\(\varphi_j\to u\) 于 \(W^{1,p}\)，并拟处处收敛到 \(\widetilde u\)。
第一项说明这个收敛的例外集在 \(P\) 上具有 \(\sigma\) 外测度零。
另一方面，迹的连续性给出
\[
 \varphi_j|_P=\operatorname{Tr}_P\varphi_j
       \longrightarrow\operatorname{Tr}_P u
       \quad\text{于 }L^p(P).
\]
再抽一个子列，使平面上的收敛几乎处处成立。
同一子列的两个点态极限必须相同，所以
\(\widetilde u|_P=\operatorname{Tr}_P u\) 在 \(\sigma\)-几乎处处成立。
这也保证了限制函数在平面完备化测度下可测。
若只指定任意 Lebesgue 代表，则可以在整个零体积的平面上随意改值，
因而不能得到同样结论。
\end{solution}

\begin{exercise}\label{b1:ex:26-slit-zero-boundary-space}
设 \(d\ge2\)，
\[
 \Omega=(-2,2)^d,\quad
 K=[-1,1]^{d-1}\times\{0\},\quad D=\Omega\setminus K.
\]
证明下列结论。
\begin{enumerate}
 \item\label{b1:item:26-slit-positive-capacity}
 集合 \(K\) 的 \(d\) 维 Lebesgue 测度为零，但 \(C_p(K)>0\)。
 \item\label{b1:item:26-slit-strict-w0}
 经零延拓并限制到 \(\Omega\)，空间 \(W_0^{1,p}(D)\) 是
 \(W_0^{1,p}(\Omega)\) 的闭真子空间。
 \item\label{b1:item:26-slit-zero-extension-insufficient}
 存在 \(v\in W^{1,p}(D)\)，使其全空间零延拓属于
 \(W^{1,p}(\mathbb R^d)\)，但 \(v\notin W_0^{1,p}(D)\)。
\end{enumerate}
解答中只用 \(W_0^{1,p}\) 的闭包定义及已经建立的迹，不调用一般域上的零拟处处刻画。
\end{exercise}
% 来源：自拟裂缝域；零体积障碍仍可留下非零迹约束，连接容量与闭包定义。
\begin{solution}
集合 \(K\) 位于一个超平面内，故 \(m_d(K)=0\)，而
\(\sigma(K)=2^{d-1}\)。
由\cref{b1:ex:26-quasicontinuous-plane-trace}，
\[
 C_p(K)\ge A_p^{-1}2^{d-1}>0.
\]

若 \(f\in W_0^{1,p}(D)\)，取 \(f_j\in C_c^\infty(D)\) 在该空间范数中逼近 \(f\)。
每个 \(f_j\) 在裂缝及 \(\partial\Omega\) 附近均为零，因此在 \(\Omega\) 上补零后属于
\(C_c^\infty(\Omega)\)。
由零延拓\cref{b1:prop:w0-zero-extension}，这些延拓在全空间 Sobolev 范数中收敛，
其极限限制到 \(\Omega\) 属于 \(W_0^{1,p}(\Omega)\)。
由于 \(K\) 的体积为零，这个映射保持 Sobolev 范数。
源空间完备，故它的像是闭子空间。

为了证明像不是整个空间，选 \(\varphi\in C_c^\infty(\Omega)\)，使
\(\varphi=1\) 于 \(K\) 的某个邻域，并令 \(v=\varphi|_D\)。
假设 \(v\in W_0^{1,p}(D)\)，则存在 \(f_j\in C_c^\infty(D)\) 满足
\(f_j\to v\) 于 \(W^{1,p}(D)\)。
将 \(f_j\) 延拓为全空间光滑紧支集函数；由于 \(m_d(K)=0\)，有
\[
 \|f_j-\varphi\|_{W^{1,p}(\mathbb R^d)}
       =\|f_j-v\|_{W^{1,p}(D)}\longrightarrow0.
\]
平面迹连续，故 \(\operatorname{Tr}_P f_j\to\operatorname{Tr}_P\varphi\) 于 \(L^p(P)\)。
但各 \(f_j\) 在 \(K\) 上恒为零，而 \(\varphi|_K=1\)，于是
\[
 \|\operatorname{Tr}_P f_j-\operatorname{Tr}_P\varphi\|_{L^p(P)}^p
       \ge \sigma(K)=2^{d-1},
\]
矛盾。故 \(\varphi\) 不在上述像中，第二项得证。

同一个 \(v\) 也证明第三项：它在 \(D\) 外补零所得函数与全空间的
\(\varphi\) 仅在 \(K\) 上可能不同，因而是同一个 \(W^{1,p}(\mathbb R^d)\) 类；
但是刚才已经证明 \(v\notin W_0^{1,p}(D)\)。
零延拓属于 Sobolev 空间，只约束体积意义下的弱导数；
裂缝上的额外零迹约束则由内部光滑紧支集逼近保留下来。
\end{solution}

% !TeX root = ../../Book1.tex
\chapter[有界变差函数与有限周长集合]{有界变差函数\\与有限周长集合}
\label{b1:ch:27}
\BookChapterNumber{b1:ch:27}
\begin{chapterabstract}
本章把可积弱梯度推广为有限向量测度，建立有界变差函数的逼近、紧性及余面积公式。
随后从集合示性函数的分布导数定义周长，在约化边界处证明半空间缩放与可整流性，
得到 De Giorgi 结构定理及 Gauss–Green 公式。最后把有界变差函数的 Sobolev 端点估计与集合的等周估计联系起来。
\end{chapterabstract}

\begin{learninggoals}
\begin{enumerate}
\item 从向量测度的欧氏全变差和测试场对偶式识别有界变差函数，区分它与 \(W^{1,1}\) 函数的导数类型。
\item 掌握严格光滑逼近中的单位分解误差，区分局部紧性、全域紧性、弱-*收敛和范数收敛。
\item 从层集截断推导有界变差函数的余面积公式，说明这一推导为何不依赖光滑边界或结构定理。
\item 区分测度论边界、约化边界和任意代表的拓扑边界，从半空间缩放恢复法向及周长面积。
\item 理解函数与集合不等式之间的双向传递，区分本章的有效估计与带最佳常数的尖锐等周定理。
\end{enumerate}
\end{learninggoals}

在一维，阶跃函数虽然没有可积的弱导数，却有一个十分明确的导数：跳跃点上的点质量。
在高维，两个常值区域之间的一道界面也应当贡献变化，而不是因为经典梯度几乎处处为零就被忽略。
有界变差函数空间保留了这类变化，把导数放进有限测度而非仅仅可积函数的范围。
因此它既与前面的 Sobolev 理论连续衔接，又能直接处理不连续函数和集合边界。

本章的证明次序依赖这一选择。先建立变差的分析工具，再用示性函数定义周长；
余面积公式只凭这些工具就能完成。随后再讨论周长究竟集中在哪一种边界上。
“周长等于面积”是这一阶段的结论，不参与此前的定义与逼近。
结构定理中最重要的桥梁是两侧体积下界：方向集中虽然指向一个半空间，
但若没有这项下界，缩放极限仍可能退化为全空间或空集。

Simon 的《Lectures on Geometric Measure Theory》%
\cite[\S\S6 and 14]{Simon1983Lectures}是本章局部变差及边界结构的对读资料。
本书展开其中需要承接的单位分解、平移紧性和图分片步骤；
更系统的有界变差函数理论可参阅 Ambrosio、Fusco 与 Pallara 的%
《Functions of Bounded Variation and Free Discontinuity Problems》%
\cite{Ambrosio2000Functions}。
后者还讨论下一章的跳跃、Cantor 部分与自由间断问题，初读本章不必先掌握这些内容。

% !TeX root = ../../Book1.tex
\section{有界变差函数}
\label{b1:sec:2701}

一维阶跃函数的经典导数几乎处处为零，但它的分布导数在跳跃点留下一个点质量。高维函数也可能把变化集中在超曲面上。若只允许弱导数由 \(L^1\) 函数表示，就会排除这些对象；若允许导数是一份有限 Radon 测度，就能同时保留通常的梯度和集中在薄集上的变化。

\ProseParagraphBreak

本章的函数除非另行说明均为实值，\(\Omega\subseteq\mathbb R^d\) 为开集。导数 \(Du\) 取值于 \(\mathbb R^d\)，其长度由欧氏范数衡量。这里 \(D\) 表示分布导数测度，\(\nabla u\) 则用于确有可积弱梯度的情形；二者不能不加说明地互换。

\ProseParagraphBreak

Simon 的《Lectures on Geometric Measure Theory》\cite[\S6, 6.1]{Simon1983Lectures}从散度测试泛函引入局部有界变差函数。本节先补齐有限维向量测度的变差与极分解，再把这种测试表述同分布导数逐项连接起来。它只调用第6章的 Radon–Nikodym 表示和第11章的 Radon 测度表示，不预借本章后面的几何结构定理。
% 实读本地正式版Simon第6节印刷35--36页，式6.1与局部变差测试表述；
% 该书把向量泛函表示接第4节，本书改经已有标量Riesz--Markov逐分量构造。
% 向量变差的有限分割与连续单位球测试证明在本节完整展开。

\subsection{分布导数为 Radon 测度}
\label{b1:subsec:270101}

先取一个有限向量 Radon 测度 \(\mu=(\mu_1,\ldots,\mu_d)\)，即各分量均为有限实 Radon 符号测度。对 Borel 集 \(A\subseteq\Omega\)，定义其欧氏全变差为
\[
 |\mu|(A)=
 \sup\left\{\sum_{j=1}^N|\mu(A_j)|:
       A=\bigsqcup_{j=1}^N A_j,\ A_j\text{ Borel}\right\}.
\]
这里 \(|\mu(A_j)|\) 是向量的欧氏长度，而不是分量绝对值之和。

\begin{proposition}\label{b1:prop:vector-measure-polar}
令 \(\lambda=\sum_{i=1}^d|\mu_i|\)，并由 Radon–Nikodym 定理写 \(\mu_i=g_i\lambda\)，\(g=(g_1,\ldots,g_d)\)。则
\[
 |\mu|=|g|\lambda.
\]
因此 \(|\mu|\) 是有限正 Radon 测度，且有
\[
 \mu=\sigma_\mu|\mu|,\qquad |\sigma_\mu|=1
 \quad |\mu|\text{-几乎处处}.
\]
这个方向函数在 \(|\mu|\)-几乎处处意义下唯一。上述结论在局部有限情形可逐个内部区域使用，并在重叠处相容。
\end{proposition}
\begin{proof}
对任意有限 Borel 分割，由向量积分的三角不等式，
\[
 \sum_j|\mu(A_j)|=\sum_j\left|\int_{A_j}g\,\mathrm d\lambda\right|
 \le\int_A|g|\,\mathrm d\lambda.
\]
反向取一个有限值向量简单函数 \(s=\sum_j a_j\cdot\mathbf1_{A_j}\)，使
\(\int_A|g-s|\,\mathrm d\lambda<\varepsilon\)。这可以逐分量简单逼近，再取共同的有限分割。于是
\[
 \begin{aligned}
 \sum_j|\mu(A_j)|
 &\ge\sum_j |a_j|\cdot \lambda(A_j)
           -\sum_j\left|\int_{A_j}(g-a_j)\,\mathrm d\lambda\right|\\
 &\ge\int_A|s|\,\mathrm d\lambda-\varepsilon
 \ge\int_A|g|\,\mathrm d\lambda-2\varepsilon.
 \end{aligned}
\]
让误差趋零，得到变差公式，因而也证明了可数可加性。有限 Radon 测度乘以非负可积密度仍为 Radon 测度，这是第11章的密度稳定性。

在 \(g\ne0\) 处取 \(\sigma_\mu=g/|g|\)，其余处取零。测度 \(|\mu|\) 在 \(\{g=0\}\) 上为零，所以所得方向在其几乎处处长度为一，并满足表示式。两个方向若给出同一个向量测度，逐分量用 Radon–Nikodym 唯一性即可证明它们相等。对局部有限分量，把论证限制到内部相对紧开集；唯一性保证局部得到的变差和方向在相交处一致。
\end{proof}

这就是向量测度的极分解，与第6章复测度的方向分解相似。用于控制的 \(\lambda\) 很方便，但真正的欧氏变差是 \(|g|\lambda\)。例如，\(\mu=(m_d,m_d)\) 限制到一个有限体积区域时，其变差是 \(\sqrt2m_d\)，不是 \(2m_d\)。

\begin{definition}\label{b1:def:bv-function}
给定 \(u\in L^1(\Omega)\)。若对每个 \(i=1,\ldots,d\)，分布偏导数都由有限 Radon 符号测度 \(D_i u\) 表示，即
\[
 \int_\Omega u\,\partial_i\varphi\,\mathrm dx
 =-\int_\Omega\varphi\,\mathrm dD_i u
 \quad(\varphi\in C_c^\infty(\Omega)),
\]
则称 \(u\) 为 \(\Omega\) 上的\term[youjie-biancha-hanshu]{有界变差函数}[function of bounded variation]，记 \(u\in BV(\Omega)\)。\foreignidx{bounded variation}\foreignidx{BV}记
\[
 Du\coloneqq(D_1u,\ldots,D_du),
\]
并以 \(|Du|\) 表示其欧氏全变差测度。
\end{definition}
\symidx{\(BV(\Omega)\)：分布导数为有限向量测度的可积函数空间}
\symidx{\(Du\)：有界变差函数的分布导数测度}
\symidx{\(\lvert Du\rvert\)：\(Du\) 的欧氏全变差测度}

若只要求 \(u\in L^1_{\mathrm{loc}}(\Omega)\)，且上述导数测度在每个紧集上具有有限变差，则记 \(u\in BV_{\mathrm{loc}}(\Omega)\)。在内部区域上施加分布测试即可定义这些测度；相邻区域上的表示由测试函数的唯一性相互一致。一个局部有界变差函数不一定在整个 \(\Omega\) 上可积，也不一定有有限的总变差。

\subsection{全变差}
\label{b1:subsec:270102}

变差不仅能由测度分割定义，也能直接从函数通过测试读出。对于任意 \(u\in L^1_{\mathrm{loc}}(\Omega)\) 和开集 \(U\subseteq\Omega\)，规定
\begin{equation}\label{b1:eq:bv-variational-definition}
 V(u,U)\coloneqq\sup\left\{\int_Uu\,\operatorname{div}\varphi\,\mathrm dx:
 \varphi\in C_c^\infty(U;\mathbb R^d),\
 \sup_U|\varphi|\le1\right\}\in[0,\infty].
\end{equation}
每个积分有限，因为测试函数的支集紧含于 \(U\)。把 \(\varphi\) 换成 \(-\varphi\)，说明上确界中也可以对积分取绝对值。

\begin{proposition}\label{b1:prop:bv-variation-duality}
设 \(u\in L^1_{\mathrm{loc}}(\Omega)\)。则 \(V(u,\Omega)<\infty\)，当且仅当分布导数 \(Du\) 是有限向量 Radon 测度。此时对每个开集 \(U\subseteq\Omega\)，
\begin{equation}\label{b1:eq:bv-variation-open-set}
 V(u,U)=|Du|(U).
\end{equation}
若只要求每个 \(\overline U\Subset\Omega\) 的 \(V(u,U)\) 有限，则相应地得到 \(BV_{\mathrm{loc}}\) 的判别和同一局部公式。
\end{proposition}
\begin{proof}
先设 \(V(u,\Omega)<\infty\)。对每个坐标 \(i\)，实线性泛函
\[
 T_i(\psi)=-\int_\Omega u\,\partial_i\psi\,\mathrm dx
 \quad(\psi\in C_c^\infty(\Omega))
\]
满足 \(|T_i(\psi)|\le V(u,\Omega)\|\psi\|_\infty\)，只须在\eqref{b1:eq:bv-variational-definition}中取 \(\varphi=\psi e_i/\|\psi\|_\infty\)。光滑紧支集函数在 \(C_0(\Omega)\) 的上确界范数中稠密：先以紧支集截断，再在离边界有正距离处磨光即可。故 \(T_i\) 唯一延拓为 \(C_0(\Omega)\) 上的有界实泛函。由\cref{b1:thm:riesz-markov-c0}，它由有限 Radon 符号测度 \(D_i u\) 表示，正好满足分布导数关系。

反过来，若 \(Du\) 为有限向量测度，则对任意容许测试场，
\[
 \left|\int_Uu\,\operatorname{div}\varphi\,\mathrm dx\right|
 =\left|\int_U\varphi\cdot\mathrm dDu\right|
 \le\int_U|\varphi|\,\mathrm d|Du|\le|Du|(U).
\]
这先给出 \(V(u,U)\le|Du|(U)\)。

要取得反向不等式，不能直接把可测方向 \(-\sigma_u\) 当作光滑测试场。由有限 Radon 测度的 \(C_c\) 稠密性，逐分量在 \(L^1(|Du|\!\llcorner U)\) 中用连续紧支集向量场逼近 \(-\sigma_u\)。再投影到闭单位球：映射 \(z\mapsto z/\max(1,|z|)\) 连续，且对任意 \(|w|\le1\) 满足
\[
 \left|\frac z{\max(1,|z|)}-w\right|\le2|z-w|.
\]
因此投影后仍能以任意小的积分误差逼近方向，支集仍紧含于 \(U\)，欧氏长度不超过一。将其补零后用非负核小尺度磨光，支集仍在 \(U\) 内，长度由凸性仍不超过一；连续紧支集函数的一致逼近又使积分误差趋零。于是可取光滑场 \(\varphi\)，使
\[
 -\int_U\varphi\cdot\sigma_u\,\mathrm d|Du|
 \ge |Du|(U)-\varepsilon.
\]
由分布恒等式，左端就是 \(\int_Uu\operatorname{div}\varphi\)，从而得到反向不等式。

局部情形在各内部区域实施相同证明。表示测度在重叠处由分布唯一性相等，故可以拼成局部有限的向量 Radon 测度。若开集 \(U\) 上总变差无穷，就先对其中的内部穷尽应用有限情形；下界可以任意大，仍有 \(V(u,U)=|Du|(U)=\infty\)。这也证明了允许无穷值的公式。
\end{proof}

对于 \(u\in W^{1,1}(\Omega)\)，各测度 \(D_i u\) 都有密度 \(\partial_i u\)，前面的向量分解给出
\begin{equation}\label{b1:eq:sobolev-bv-variation}
 Du=\nabla u\,m_d,\qquad
 |Du|=|\nabla u|\,m_d.
\end{equation}
因此 \(W^{1,1}(\Omega)\subseteq BV(\Omega)\)，且光滑函数的变差就是梯度长度的积分。一般有界变差函数的导数测度可能含有相对于体积的奇异部分，不能把\eqref{b1:eq:sobolev-bv-variation}不加条件地用于所有有界变差函数。

\ProseParagraphBreak

这一测试表述也是下半连续性的来源：固定测试场时，积分只依赖 \(u\) 在一个内部紧集上的 \(L^1\) 值；取所有这类连续表达式的上确界，便得到可通过极限保留的下界。下一节将把这个事实连同光滑逼近一起用于紧性。

\subsection{有界变差函数空间的范数}
\label{b1:subsec:270103}

在 \(BV(\Omega)\) 上规定
\[
 \|u\|_{BV(\Omega)}\coloneqq\|u\|_{L^1(\Omega)}+|Du|(\Omega).
\]
向量测度的三角不等式保证它满足范数三角不等式，零阶项则保证范数为零只对应零等价类。

\begin{proposition}\label{b1:prop:bv-complete}
上述范数使 \(BV(\Omega)\) 成为 Banach 空间。
\end{proposition}
\begin{proof}
设 \((u_n)\) 是 Cauchy 列。由 \(L^1\) 完备性，存在 \(u\in L^1(\Omega)\)，使 \(u_n\to u\) 于 \(L^1\)。同时
\[
 |D_i u_n-D_i u_m|(\Omega)
 \le|D(u_n-u_m)|(\Omega)\longrightarrow0.
\]
第11章的有限 Radon 测度空间在全变差范数下完备，因此每个分量都趋于某个有限 Radon 符号测度 \(\mu_i\)。对任意 \(\varphi\in C_c^\infty(\Omega)\)，
\[
 \int_\Omega u\,\partial_i\varphi
 =\lim_n\int_\Omega u_n\partial_i\varphi
 =-\lim_n\int_\Omega\varphi\,\mathrm dD_i u_n
 =-\int_\Omega\varphi\,\mathrm d\mu_i.
\]
故 \(D_i u=\mu_i\)，而 \(u\in BV(\Omega)\)。再由
\[
 |D(u_n-u)|(\Omega)
 \le\sum_{i=1}^d|D_i u_n-\mu_i|(\Omega)\longrightarrow0,
\]
得到原范数中的收敛。
\end{proof}

局部化仍然可用，但乘法公式中需要同时保留测度项与体积项。

\begin{proposition}\label{b1:prop:bv-product-rule}
设 \(u\in BV(\Omega)\)。若 \(v\) 在 \(\Omega\) 中局部 Lipschitz，并满足 \(v,\nabla v\in L^\infty(\Omega)\)，则
\[
 v\cdot u\in BV(\Omega),\qquad
 D(v\cdot u)=v\cdot Du+u\cdot\nabla v\,m_d,
\]
且
\[
 |D(vu)|(\Omega)
 \le\|v\|_\infty|Du|(\Omega)
       +\|\nabla v\|_\infty\|u\|_1.
\]
第一项使用 \(v\) 的连续点值，第二项使用其几乎处处弱梯度。相应局部结论也成立。
\end{proposition}
\begin{proof}
先取光滑 \(v\)。对任意光滑紧支集测试函数，将乘积 \(v\cdot\varphi\) 代入 \(Du\) 的定义，再用普通乘积求导：
\[
 \int_\Omega u\cdot v\,\partial_i\varphi
 =\int_\Omega u\,\partial_i(v\cdot\varphi)
       -\int_\Omega u\cdot\varphi\,\partial_i v
 =-\int_\Omega v\cdot\varphi\,\mathrm dD_i u
       -\int_\Omega u\cdot\varphi\,\partial_i v.
\]
这就是所需测度公式。若 \(v\) 只局部 Lipschitz，固定测试函数的支集，在其邻域中磨光 \(v\)。磨光函数局部一致趋于 \(v\)，梯度局部一致有界，且沿一个子列几乎处处趋于 \(\nabla v\)。因此测度积分由局部一致收敛通过极限，体积积分由 \(|u|\) 控制收敛通过极限。每个测试函数都得到相同公式，故分布表示成立。

最后，对 Borel 集上的向量积分取变差，得到
\[
 |D(vu)|\le |v|\,|Du|+|u|\cdot|\nabla v|\,m_d.
\]
积分即得估计，且零阶积 \(v\cdot u\) 可积。局部版本限制到内部区域即可。
\end{proof}

特别地，乘以内部光滑紧支集截断后，可以在域外补零，得到全空间有界变差函数。证明只需在截断支集附近代入测试函数，因而不会产生边界项。这个内部补零与在任意边界直接补零不同；后者可能增加一份边界测度，甚至使总变差失去控制。

\subsection{基本例子}
\label{b1:subsec:270104}

在区间 \((-1,1)\) 上取 \(u=\mathbf1_{(0,1)}\)。对任意测试函数，
\[
 \int_{-1}^1u\varphi'\,\mathrm dx
 =\int_0^1\varphi'\,\mathrm dx=-\varphi(0).
\]
因此 \(Du=\delta_0\)，\(|Du|((-1,1))=1\)。这个函数属于 \(BV\)，却不属于 \(W^{1,1}\)，因为点质量没有 \(L^1\) 密度。导数测度记录了一次跳跃的大小，而不是把它藏进“经典导数几乎处处为零”的表述中。

\ProseParagraphBreak

高维的长方体可用相同计算处理。令 \(Q=\prod_{i=1}^d(a_i,b_i)\)。在全空间中，
\[
 \int_Q\partial_i\varphi\,\mathrm dx
 =\int_{\prod_{j\ne i}(a_j,b_j)}
       \bigl(\varphi(x',b_i)-\varphi(x',a_i)\bigr)\,\mathrm dx'.
\]
所以 \(D_i\mathbf1_Q\) 在下侧面取正的面测度，在上侧面取负的面测度。除边棱这个面测度零集以外，各侧面互不重叠；导数方向在每个面上均为单位法向的负号。因此
\[
 |D\mathbf1_Q|(\mathbb R^d)
 =2\sum_{i=1}^d\prod_{j\ne i}(b_j-a_j).
\]
这是边界各面的面积之和。当 \(d=1\) 时空乘积取一，恢复一个有界区间两端的总变差二。到\secref{b1:sec:2704}，我们会把这种现象扩展为一般集合的周长。

\ProseParagraphBreak

变化也可能既不由体积密度给出，又不集中在离散跳点。令 \(\mu_C\) 为标准 Cantor 概率测度，\(c(x)=\mu_C([0,x])\)，\(0<x<1\)。对 \(\varphi\in C_c^\infty(0,1)\)，Fubini 定理给出
\[
 \begin{aligned}
 \int_0^1c(x)\varphi'(x)\,\mathrm dx
 &=\int_{[0,1]}\left(\int_t^1\varphi'(x)\,\mathrm dx\right)
                                      \mathrm d\mu_C(t)\\
 &=-\int_{[0,1]}\varphi(t)\,\mathrm d\mu_C(t).
 \end{aligned}
\]
端点不带原子，故在开区间内有 \(Dc=\mu_C\)，总变差为一。函数 \(c\) 连续，在 Cantor 集的补区间上局部常数，所以经典导数几乎处处为零；但导数测度仍保留了非零的奇异连续变化。这是后面 Cantor 部分的最基本模型。

\ProseParagraphBreak

这些例子也说明，有界变差函数空间由积分与分布刻画，不是对任意点态代表计算振幅总和。把零函数在有理点上改为一，不改变其 \(L^1\) 类和分布导数，所以在多维定义下仍代表 \(BV\) 空间的零元。若回到一维的点态分割变差，则需先选择相应的规范代表；不能把任意零测改动后的点态分割和与这里的 \(|Du|\) 混为一谈。

% !TeX root = ../../Book1.tex
\section{有界变差函数的逼近与紧性}
\label{b1:sec:2702}

对 Sobolev 函数，磨光可以同时逼近函数及其弱梯度；对一般有界变差函数，导数可能含有集中在零测集上的部分，光滑梯度不能按总变差范数逼近这种测度。本节保留函数的 \(L^1\) 收敛和导数的总量收敛，并用导数测度控制平移，得到局部紧性。所有函数均取实值，\(|Du|\) 使用前一节的欧氏全变差。

Simon 的《Lectures on Geometric Measure Theory》\cite[\S6, 6.2 LEMMA and 6.3 THEOREM]{Simon1983Lectures}以磨光及变差对偶式组织有界变差函数的逼近与紧性。下面的全域严格逼近另外展开单位分解的误差控制；局部紧性的抽选步骤则调用已经证明的 \(L^1\) 平移判据。

% 实读 Simon1983Lectures 正式公开本§6，印刷35--38页：
% 6.1局部BV的测度表示、6.2局部磨光及变差测度收敛、6.3局部紧性；
% 核看第37页。原书6.3将平滑函数的紧性细节交给GT，本节改用已证FK。
% 任意开集的全域严格逼近按22.1局部有限分区路线独立展开，
% 特别控制rho_i*(u grad psi_i)-u grad psi_i，不声称直接来自Simon完整证明。
% 不使用一般BV边界延拓；补零只对内部光滑截断或另列已知BV补零假设。

\subsection{光滑逼近}
\label{b1:subsec:270201}

\begin{definition}\label{b1:def:bv-strict-convergence}
设 \(u_j,u\in BV(\Omega)\)。若
\[
 \|u_j-u\|_{L^1(\Omega)}\longrightarrow0,\quad
 |Du_j|(\Omega)\longrightarrow|Du|(\Omega),
\]
则称 \(u_j\) 严格收敛到 \(u\)；这种收敛称为\term[youjie-biancha-kongjian-yange-shoulian]{有界变差空间中的严格收敛}[strict convergence in BV]。
\end{definition}

这一定义比较的是两个导数测度的总变差数值，并没有要求
\(|D(u_j-u)|(\Omega)\to0\)。先记录内部截断及全空间磨光的公式。

\begin{lemma}\label{b1:lem:bv-cutoff-and-mollification}
设 \(u\in BV_{\mathrm{loc}}(\Omega)\)、\(\eta\in C_c^\infty(\Omega)\)，把 \(\eta u\) 在 \(\Omega\) 外补零，记为 \(v\)。则 \(v\in BV(\mathbb R^d)\)，且
\begin{equation}\label{b1:eq:bv-cutoff-derivative}
 Dv=\eta Du+u\nabla\eta\,m_d,
 \quad
 |Dv|(\mathbb R^d)
 \leq\int_\Omega|\eta|\,\mathrm d|Du|
       +\int_\Omega|u|\cdot|\nabla\eta|\,\mathrm dx,
\end{equation}
右侧两项测度均在域外补零。

若 \(v\in BV(\mathbb R^d)\)，用非负、偶、单位积分的磨光函数定义
\(v_\varepsilon=\rho_\varepsilon*v\)，则
\[
 v_\varepsilon\in C^\infty(\mathbb R^d)\cap W^{1,1}(\mathbb R^d),
 \quad v_\varepsilon\longrightarrow v\ \text{于 }L^1,
\]
并有
\begin{equation}\label{b1:eq:bv-convolution-gradient}
 \nabla v_\varepsilon=\rho_\varepsilon*Dv,\quad
 \int_{\mathbb R^d}|\nabla v_\varepsilon|\,\mathrm dx
 \leq|Dv|(\mathbb R^d).
\end{equation}
\end{lemma}
\begin{proof}
对任意 \(\Phi\in C_c^\infty(\mathbb R^d;\mathbb R^d)\)，函数
\(\eta\Phi\) 的支集紧含于 \(\Omega\)。乘积公式给出
\[
 \begin{aligned}
 \int_{\mathbb R^d}v\,\operatorname{div}\Phi
 &=\int_\Omega u\,\operatorname{div}(\eta\Phi)
       -\int_\Omega u\nabla\eta\cdot\Phi\\
 &=-\int_\Omega\eta\Phi\cdot\mathrm dDu
       -\int_\Omega u\nabla\eta\cdot\Phi.
 \end{aligned}
\]
两项导数测度均有限，得到补零后的导数公式；总变差的三角不等式给出估计。这里没有跨越 \(\partial\Omega\) 的边界项，因为 \(\eta\) 的支集留在域内。

对任意有限向量测度 \(\mu\)，定义
\[
 (\rho_\varepsilon*\mu)(x)
 =\int_{\mathbb R^d}\rho_\varepsilon(x-y)\,\mathrm d\mu(y).
\]
由向量测度的极分解和积分三角不等式，
\[
 \bigl|(\rho_\varepsilon*\mu)(x)\bigr|
 \leq\int\rho_\varepsilon(x-y)\,\mathrm d|\mu|(y),
 \quad
 \int|\rho_\varepsilon*\mu|\,\mathrm dx\leq|\mu|(\mathbb R^d).
\]
后一式由 Tonelli 定理及核的单位积分得到。

函数 \(v_\varepsilon\) 的光滑性由卷积求导给出，而 \(L^1\) 收敛来自\cref{b1:prop:approximate-identity-lp}。利用核的偶性及 Fubini 定理，对测试向量场 \(\Phi\) 有
\[
 \begin{aligned}
 \int v_\varepsilon\operatorname{div}\Phi
 &=\int v\,\operatorname{div}(\rho_\varepsilon*\Phi)\\
 &=-\int(\rho_\varepsilon*\Phi)\cdot\mathrm dDv
 =-\int\Phi(x)\cdot(\rho_\varepsilon*Dv)(x)\,\mathrm dx.
 \end{aligned}
\]
因此经典梯度为 \(\rho_\varepsilon*Dv\)。刚证的向量测度卷积估计说明它可积，并给出\eqref{b1:eq:bv-convolution-gradient}。
\end{proof}

\begin{theorem}\label{b1:thm:bv-strict-approximation}
设 \(\Omega\subseteq\mathbb R^d\) 为任意非空开集，\(u\in BV(\Omega)\)。存在
\[
 u_j\in C^\infty(\Omega)\cap W^{1,1}(\Omega)
\]
严格收敛到 \(u\)，也就是
\[
 u_j\to u\ \text{于 }L^1(\Omega),\quad
 \int_\Omega|\nabla u_j|\,\mathrm dx\longrightarrow|Du|(\Omega).
\]
这里不要求 \(u_j\) 紧支集，也不要求边界正则。
\end{theorem}
\begin{proof}
给定 \(\varepsilon>0\)。沿用\cref{b1:thm:meyers-serrin}证明中的局部有限开层及光滑单位分解，取
\[
 \psi_i\geq0,\quad \sum_i\psi_i=1,\quad
 K_i=\operatorname{supp}\psi_i\Subset V_{j(i)},\quad
 \overline{V_{j(i)}}\Subset\Omega.
\]
每个紧集只遇到有限个开层 \(V_j\)，每一层只分配有限多个 \(\psi_i\)。

令 \(f_i=\psi_i u\)，并补零到全空间。由前一引理，
\[
 Df_i=\mu_i+g_i\,m_d,\quad
 \mu_i=\psi_iDu,\quad g_i=u\nabla\psi_i.
\]
每个 \(f_i,g_i\) 都属于全空间 \(L^1\)，且支集包含于 \(K_i\)。选择足够小的半径 \(\delta_i>0\)，同时满足
\[
 K_i+\overline B(0,\delta_i)\subset V_{j(i)},\quad
 \|\rho_{\delta_i}*f_i-f_i\|_1<\varepsilon2^{-i-1},\quad
 \|\rho_{\delta_i}*g_i-g_i\|_1<\varepsilon2^{-i-1}.
\]
支集条件由紧集到补集的正距离保证；两个误差条件由 \(L^1\) 磨光逼近保证，因而可以通过进一步缩小同一个半径同时实现。

定义
\[
 v=\sum_i\rho_{\delta_i}*f_i.
\]
每个紧集只遇到有限个磨光后的支集，所以此和在局部为有限个光滑函数之和，故 \(v\in C^\infty(\Omega)\)。又因
\[
 \sum_i\|f_i\|_1=\int_\Omega|u|\sum_i\psi_i=\|u\|_1,
\]
卷积后的级数也在 \(L^1\) 中绝对收敛，并且
\[
 \|v-u\|_1
 \leq\sum_i\|\rho_{\delta_i}*f_i-f_i\|_1<\varepsilon/2.
\]

梯度需要另作估计。在任意内部紧集上，单位分解满足
\(\sum_i\nabla\psi_i=0\)。因此局部逐项求导给出
\begin{equation}\label{b1:eq:bv-partition-commutator}
 \nabla v
 =\sum_i\rho_{\delta_i}*\mu_i
       +\sum_i(\rho_{\delta_i}*g_i-g_i).
\end{equation}
第一组项的全域 \(L^1\) 范数可求和，因为非负权重给出
\[
 \sum_i\int_{\mathbb R^d}|\rho_{\delta_i}*\mu_i|
 \leq\sum_i|\mu_i|(\mathbb R^d)
 =\sum_i\int_\Omega\psi_i\,\mathrm d|Du|
 =|Du|(\Omega).
\]
第二组误差的 \(L^1\) 范数之和小于 \(\varepsilon/2\)。所以
\[
 v\in W^{1,1}(\Omega),\quad
 \int_\Omega|\nabla v|\leq|Du|(\Omega)+\varepsilon/2.
\]
这里没有要求 \(\sum_i\|u\nabla\psi_i\|_1<\infty\)；相消先在局部进行，全域求和只作用于已经控制的误差。

令 \(\varepsilon\downarrow0\)，得到所需的 \(L^1\) 逼近列及梯度积分的上极限估计。为取得反向不等式，对每个
\(\Phi\in C_c^\infty(\Omega;\mathbb R^d)\)、\(|\Phi|\leq1\)，有
\[
 \int_\Omega u\,\operatorname{div}\Phi
 =\lim_j\int_\Omega u_j\operatorname{div}\Phi
 \leq\liminf_j\int_\Omega|\nabla u_j|.
\]
对 \(\Phi\) 取上确界，由变差对偶式得到
\(|Du|(\Omega)\leq\liminf_j\int_\Omega|\nabla u_j|\)，完成证明。
\end{proof}

有界变差空间中的严格收敛确实弱于 \(BV\) 范数收敛。例如在 \((-1,1)\) 上取
\(u=\mathbf1_{(0,1)}\)，则 \(Du=\delta_0\)。令
\[
 u_\varepsilon(x)=\int_{-\infty}^x\rho_\varepsilon(t)\,\mathrm dt,
 \quad 0<\varepsilon<1.
\]
这些函数光滑，与 \(u\) 只在 \((-\varepsilon,\varepsilon)\) 内不同，且
\(\int_{-1}^1|u_\varepsilon'|=1=|Du|\bigl((-1,1)\bigr)\)。因此它们严格收敛，但
\(\rho_\varepsilon\,m_1\) 与 \(\delta_0\) 相互奇异，所以
\[
 |D(u_\varepsilon-u)|\bigl((-1,1)\bigr)=2.
\]
不能把光滑严格逼近改称一般有界变差函数的光滑范数稠密性。

\subsection{下半连续性}
\label{b1:subsec:270202}

前面最后一步只使用了紧支集测试，因此它实际上是一个局部结论。

\begin{theorem}[变差的下半连续性]\label{b1:thm:bv-lower-semicontinuity}
设 \(u_j,u\in L^1_{\mathrm{loc}}(\Omega)\)，且
\(u_j\to u\) 于 \(L^1_{\mathrm{loc}}(\Omega)\)。对任意开集 \(U\subseteq\Omega\)，有
\[
 V(u,U)\leq\liminf_{j\to\infty}V(u_j,U),
\]
其中 \(V\) 是\cref{b1:prop:bv-variation-duality}中的扩展变差，允许取无穷。

若右端在每个内部相对紧开集上有限，则 \(u\in BV_{\mathrm{loc}}(\Omega)\)。在导数测度存在时，上式就是
\[
 |Du|(U)\leq\liminf_j|Du_j|(U).
\]
若另有 \(u\in L^1(\Omega)\) 及 \(\liminf_jV(u_j,\Omega)<\infty\)，则 \(u\in BV(\Omega)\)。
\end{theorem}
\begin{proof}
任取 \(\Phi\in C_c^\infty(U;\mathbb R^d)\)、\(|\Phi|\leq1\)。因为
\(\operatorname{div}\Phi\) 有界且紧支集，局部 \(L^1\) 收敛给出
\[
 \int_Uu\,\operatorname{div}\Phi
 =\lim_j\int_Uu_j\operatorname{div}\Phi
 \leq\liminf_jV(u_j,U).
\]
右端不依赖 \(\Phi\)，故对测试场取上确界，得到第一项。其余断言由局部与全局的有限变差判别及 \(V=|Du|\) 得到。
\end{proof}

对层集示性函数也可以直接用这个版本：即使 \(\{u>t\}\) 的体积无穷，
\(\mathbf1_{\{u>t\}}\) 仍局部可积，扩展变差依然有定义。全域
\(L^1\) 是进入 \(BV(\Omega)\) 的额外要求，不是定义局部周长的必要条件。

\subsection{有界变差函数的紧性定理}
\label{b1:subsec:270203}

导数是测度时，平移仍能由它的总变差控制。

\begin{lemma}\label{b1:lem:bv-translation-estimate}
若 \(v\in BV(\mathbb R^d)\)，则对每个 \(h\in\mathbb R^d\)，有
\[
 \|\tau_hv-v\|_{L^1(\mathbb R^d)}
 \leq |h|\cdot|Dv|(\mathbb R^d),
 \quad \tau_hv(x)=v(x-h).
\]
\end{lemma}
\begin{proof}
先对全空间磨光 \(v_\varepsilon\) 使用沿线段的微积分基本定理：
\[
 |v_\varepsilon(x-h)-v_\varepsilon(x)|
 \leq|h|\int_0^1|\nabla v_\varepsilon(x-th)|\,\mathrm dt.
\]
积分换序及平移不变性给出
\[
 \|\tau_hv_\varepsilon-v_\varepsilon\|_1
 \leq|h|\int|\nabla v_\varepsilon|
 \leq|h|\cdot|Dv|(\mathbb R^d).
\]
由 \(v_\varepsilon\to v\) 于 \(L^1\)，以及平移的等距性，左端在
\(\varepsilon\downarrow0\) 时趋于 \(\|\tau_hv-v\|_1\)，从而完成证明。
\end{proof}

\begin{theorem}[有界变差函数的局部紧性]\label{b1:thm:bv-local-compactness}
设 \(u_j\in BV_{\mathrm{loc}}(\Omega)\)，且对每个
\(\overline U\Subset\Omega\) 的开集 \(U\)，都有
\[
 \sup_j\bigl(\|u_j\|_{L^1(U)}+|Du_j|(U)\bigr)<\infty.
\]
则存在子列及 \(u\in BV_{\mathrm{loc}}(\Omega)\)，使
\[
 u_{j_k}\longrightarrow u\quad\text{于 }L^1_{\mathrm{loc}}(\Omega).
\]
对任意开集 \(U\subseteq\Omega\)，所得极限满足
\(|Du|(U)\leq\liminf_k|Du_{j_k}|(U)\)。
\end{theorem}
\begin{proof}
取递增开集穷竭 \(U_m\)，使
\(\overline{U_m}\Subset U_{m+1}\)，并取
\(\eta_m\in C_c^\infty(U_{m+1})\)、\(0\leq\eta_m\leq1\)，在 \(U_m\) 上等于 \(1\)。把 \(\eta_mu_j\) 补零到全空间，记为 \(v_{m,j}\)。由截断引理及假设，
\[
 \sup_j\bigl(\|v_{m,j}\|_1+|Dv_{m,j}|(\mathbb R^d)\bigr)<\infty
\]
对每个固定 \(m\) 成立，且这组函数都在
\(\operatorname{supp}\eta_m\) 外为零。

共同紧支集保证尾部条件，前一引理保证平移误差一致趋零。由\cref{b1:thm:frechet-kolmogorov-lp}取 \(p=1\)，对固定 \(m\)，族
\((v_{m,j})_j\) 在全空间 \(L^1\) 中相对紧，因而 \((u_j)_j\) 在
\(L^1(U_m)\) 中有收敛子列。

逐次为 \(m=1,2,\ldots\) 选取子列，再取对角子列。其在每个 \(U_m\) 上都有 \(L^1\) 极限，不同区域的极限在交集上几乎处处相同，因为
\(L^1\) 极限唯一。因此这些极限定义了一个 \(u\in L^1_{\mathrm{loc}}(\Omega)\)。局部变差界及下半连续性保证 \(u\in BV_{\mathrm{loc}}\)，并给出全部开集上的所述不等式。
\end{proof}

若原序列的 \(BV(\Omega)\) 范数一致有界，局部紧性得到的极限也属于全域
\(BV(\Omega)\)：可以再抽出几乎处处收敛子列，用 Fatou 引理控制其
\(L^1\) 范数，再用变差下半连续性控制 \(|Du|(\Omega)\)。但这个结论尚未给出全域 \(L^1\) 收敛，质量仍可能移到穷竭所遗漏的区域。

\begin{corollary}\label{b1:cor:bv-global-compactness-with-tails}
设 \((u_j)\) 在 \(BV(\Omega)\) 中一致有界。若存在递增紧集
\(K_m\Subset\Omega\)，使
\[
 \lim_{m\to\infty}\sup_j
       \int_{\Omega\setminus K_m}|u_j|\,\mathrm dx=0,
\]
则它有全域 \(L^1(\Omega)\) 收敛子列。

特别地，若 \(\Omega\) 有界，只需另有边界附近的统一质量控制
\[
 \lim_{\delta\downarrow0}\sup_j
 \int_{\{\,x\in\Omega:\operatorname{dist}(x,\Omega^c)<\delta\,\}}
       |u_j(x)|\,\mathrm dx=0.
\]
另一项充分条件是：\(\Omega\) 有界，而且补零函数 \(E_0u_j\) 属于
\(BV(\mathbb R^d)\)，并在该空间中一致有界。
\end{corollary}
\begin{proof}
先取局部 \(L^1\) 收敛子列，再抽出几乎处处收敛子列，极限记为
\(u\in BV(\Omega)\)。Fatou 引理使极限继承同一个尾部界。给定
\(\varepsilon>0\)，取 \(m\)，使序列及极限在 \(\Omega\setminus K_m\) 上的 \(L^1\) 范数都小于 \(\varepsilon\)。于是
\[
 \|u_j-u\|_{L^1(\Omega)}
 \leq\|u_j-u\|_{L^1(K_m)}+2\varepsilon.
\]
第一项由局部收敛趋于零，再令 \(\varepsilon\downarrow0\)。

当 \(\Omega\) 有界时，可取
\(K_m=\{\,x\in\Omega:\operatorname{dist}(x,\Omega^c)\geq1/m\,\}\)；
这些集合紧含于 \(\Omega\)，并穷竭它，所以边界条件推出前述尾部条件。

最后一种情形直接在全空间应用平移引理和 Fréchet--Kolmogorov 定理：补零函数具有共同有界支承区域，且全空间导数总变差一致有界。所得全空间 \(L^1\) 收敛限制到 \(\Omega\) 即可。
\end{proof}

仅有“\(\Omega\) 有界”不能代替这些条件。第\ref{b1:ex:25-infinitely-many-components}题中，取 \(p=1\)，在越来越小的不同连通分量上放置归一化常数，得到
\[
 \|u_j\|_{L^1(\Omega)}=1,\quad Du_j=0,\quad
 \|u_j-u_k\|_{L^1(\Omega)}=2\quad(j\ne k).
\]
这同时是有界变差函数空间中的有界列，却没有全域 \(L^1\) 收敛子列。对有界
Lipschitz 区域，若要通过延拓证明无附加条件的有界变差函数的全局紧性，还须先建立有界变差函数的延拓定理；第22章的 \(W^{1,1}\) 延拓算子并未定义在一般有界变差函数上，不能直接套用。

\subsection{弱-*收敛}
\label{b1:subsec:270204}

局部 \(L^1\) 收敛首先决定分布导数的极限，变差界再把光滑测试扩展为连续测试。

\begin{proposition}\label{b1:prop:bv-distributional-weak-star}
设 \(u_j,u\in BV_{\mathrm{loc}}(\Omega)\)，\(u_j\to u\) 于
\(L^1_{\mathrm{loc}}(\Omega)\)，且导数变差在每个内部紧集的邻域上一致有界。则对每个
\(\Phi\in C_c(\Omega;\mathbb R^d)\)，有
\[
 \int_\Omega\Phi\cdot\mathrm dDu_j
 \longrightarrow\int_\Omega\Phi\cdot\mathrm dDu.
\]
若进一步有 \(\sup_j|Du_j|(\Omega)<\infty\)，则 \(|Du|(\Omega)<\infty\)，并且测试函数可以取为任意 \(\Phi\in C_0(\Omega;\mathbb R^d)\)。
\end{proposition}
\begin{proof}
若 \(\Phi\) 光滑且紧支集，由分布导数定义，
\[
 \int\Phi\cdot\mathrm dDu_j
 =-\int u_j\operatorname{div}\Phi
 \longrightarrow-\int u\operatorname{div}\Phi
 =\int\Phi\cdot\mathrm dDu.
\]
给定连续紧支集 \(\Phi\)，可把它补零后磨光，得到统一收敛于它的光滑向量场 \(\Phi_\ell\)，并使这些场的支集留在某个固定
\(K\Subset\Omega\) 内。若 \(M_K=\sup_j|Du_j|(K)<\infty\)，则
\[
 \left|\int(\Phi-\Phi_\ell)\cdot\mathrm dDu_j\right|
 \leq M_K\|\Phi-\Phi_\ell\|_\infty.
\]
极限测度有同样的估计。先固定 \(\ell\) 令 \(j\to\infty\)，再令
\(\ell\to\infty\)，得到连续紧支集测试的结论。

若总变差一致有界，下半连续性给出 \(|Du|(\Omega)<\infty\)。由\cref{b1:prop:cc-dense-c0}，连续紧支集函数在 \(C_0\) 中一致稠密；用全域总变差界重复刚才的误差估计，就得到全部 \(C_0\) 测试。
\end{proof}

\begin{definition}\label{b1:def:bv-weak-star-convergence}
设 \(u_j,u\in BV(\Omega)\)。若 \(u_j\to u\) 于 \(L^1(\Omega)\)，且有限向量测度 \(Du_j\) 对全部 \(C_0(\Omega;\mathbb R^d)\) 测试弱-*收敛到 \(Du\)，则称这种收敛为\term[youjie-biancha-kongjian-ruoxing-shoulian]{有界变差空间中的弱-*收敛}[weak-* convergence in BV]。
\end{definition}

因此，对于有界变差函数空间中的有界列，一旦已经取得全域 \(L^1\) 收敛，导数测度的弱-*收敛就随之成立。只取得局部 \(L^1\) 收敛时，仍有前一命题的测度测试结论，却不能省去本定义中的全域 \(L^1\) 条件。

弱-*收敛不保证总变差数值收敛。边界逃逸给出一个例子：在
\(\Omega=(0,1)\) 上，令 \(u_j=\mathbf1_{(0,1/j)}\)。则
\[
 u_j\to0\ \text{于 }L^1,\quad
 Du_j=-\delta_{1/j}\overset{*}{\rightharpoonup}0,\quad
 |Du_j|(\Omega)=1.
\]
这里 \(C_0\bigl((0,1)\bigr)\) 测试在边界附近趋零；常数函数 \(1\) 不属于这个测试空间。若改用全部 \(C_b\) 测试，就换成了\chapref{b1:ch:12}的另一种测度收敛。

振荡也会造成变差损失，即使不把质量集中到边界。取
\(\Omega=(0,2\pi)\)、\(u_j(x)=j^{-1}\sin(jx)\)。函数在 \(L^1\) 中趋零，导数为 \(\cos(jx)\,m_1\)，故由前一命题弱-*趋零，但
\[
 |Du_j|(\Omega)=\int_0^{2\pi}|\cos(jx)|\,\mathrm dx=4.
\]
这些变差测度在边界宽度为 \(\delta\) 的区域内质量至多 \(2\delta\)，所以即使控制了边界质量，也不能从导数的弱-*收敛推出总变差收敛。严格逼近定理额外取得的，正是后一个数值极限。

% !TeX root = ../../Book1.tex
\section{BV 的余面积公式}
\label{b1:sec:2703}

一个实值函数可以由全部超水平集重建。对于 Sobolev 函数，第%
\ref{b1:sec:1802}节用梯度长度补偿水平集上的面积；对于 BV 函数，
梯度不再一定是函数，但超水平集的边界仍可由示性函数的分布导数描述。
本节证明：函数的总变差恰好是各层周长的积分。证明先用窄条带截断处理
\(W^{1,1}\) 函数，再借严格逼近通过极限，不以有限周长集合的深层边界结构为前提。

% 实际来源范围：本地 Simon1983Lectures 正式扫描本，第6节印刷34--40页，
% 6.1的局部变差对偶、6.2的磨光及6.4--6.6的Poincare型估计。
% 该范围没有提供本节BV余面积的完整证明；下面条带截断、可数测试族和
% 有符号分层的证明由本书已证工具独立组织，不冒称来自原书某一未读定理。
% 依赖：27.1变差对偶；27.2扩展变差L1loc下半连续与任意开集严格逼近；
% 21章Sobolev截断。第18章仅作几何对照，不作为本证明的必要黑箱。

\subsection{超水平集}
\label{b1:subsec:270301}

以下固定开集 \(\Omega\subseteq\mathbb R^d\)。为实值可测函数 \(u\) 记
\[
 E_t=\{\,x\in\Omega:u(x)>t\,\},\quad t\in\mathbb R.
\]
改动 \(u\) 的一个 Lebesgue 零测集，只会在同一个零测集内改动每个 \(E_t\)。
示性函数的分布导数不受这种改动影响，所以后面的周长由原来的几乎处处类决定。
需要联合可测性时，可以先选一个 Borel 代表。

对有限实数 \(a,b\)，有两个基本恒等式：
\begin{equation}\label{b1:eq:bv-signed-layer-cake}
 a=\int_{\mathbb R}
       \bigl(\mathbf1_{\{a>t\}}-\mathbf1_{\{0>t\}}\bigr)\,\mathrm dt,
 \quad
 |a-b|=\int_{\mathbb R}
       \bigl|\mathbf1_{\{a>t\}}-\mathbf1_{\{b>t\}}\bigr|\,\mathrm dt.
\end{equation}
第一式在 \(a\ge0\) 时积分区间为 \((0,a)\)，在 \(a<0\) 时则在
\((a,0)\) 上积分 \(-1\)；第二式的被积函数只在两个端点之间等于 \(1\)。
负值情形必须保留第一式中的基准项，不能将有符号函数写成所有超水平集示性函数的直接积分。

由第二式及 Tonelli 定理，对 \(u,v\in L^1(U)\)，有
\begin{equation}\label{b1:eq:bv-level-set-l1-distance}
 \int_{\mathbb R}
 m_d\bigl(\{u>t\}\mathbin{\triangle}\{v>t\}\bigr)\,\mathrm dt
   =\|u-v\|_{L^1(U)},
\end{equation}
其中所有层集均取在 \(U\) 内。若 \(u_j\to u\) 于 \(L^1(U)\)，
抽取满足 \(\sum_j\|u_j-u\|_1<\infty\) 的子列，就得到对几乎每个 \(t\)，
\[
 \mathbf1_{\{u_j>t\}}\longrightarrow\mathbf1_{\{u>t\}}
       \quad\text{于 }L^1(U).
\]
这里的子列对几乎所有层值共同有效；它不是为每个 \(t\) 分别抽取的子列。

\subsection{周长}
\label{b1:subsec:270302}

在任意开集 \(U\subseteq\Omega\) 上，沿用前节的扩展变差记号
\[
 V(f,U)=\sup\left\{\,
     \int_U f\,\operatorname{div}\varphi\,\mathrm dx:
     \varphi\in C_c^1(U;\mathbb R^d),\
     \|\varphi\|_\infty\le1\,\right\}.
\]
\(C_c^1\) 测试场可在共同紧集内磨光，并在必要时按上确界范数归一化；
因此这个上确界与前节使用 \(C_c^\infty\) 的定义相同。
这个数允许为无穷；\(f\) 只需局部可积。测试场的大小始终用欧氏范数。

\begin{definition}\label{b1:def:perimeter}
设 \(E\subseteq\Omega\) 为 Lebesgue 可测集。定义它在开集 \(U\subseteq\Omega\) 内的%
\term[zhouchang]{周长}[perimeter]为
\[
 P(E,U)=V(\mathbf1_E,U).
\]
若 \(P(E,\Omega)<\infty\)，则称 \(E\) 为 \(\Omega\) 内的%
\term[youxian-zhouchang-ji]{有限周长集}[set of finite perimeter]。
若对每个满足 \(\overline U\Subset\Omega\) 的开集 \(U\) 都有 \(P(E,U)<\infty\)，
则称 \(E\) 在 \(\Omega\) 内具有局部有限周长。
后一条件下，\(D\mathbf1_E\) 是局部有限的向量 Radon 测度；对 Borel 集 \(A\subseteq\Omega\)，记
\[
 P(E,A)=|D\mathbf1_E|(A).
\]
全空间周长简记为 \(P(E)=P(E,\mathbb R^d)\)。
\end{definition}
\symidx{\(P(E,U)\)：集合在开集内的相对周长}
\symidx{\(P(E,\cdot)\)：有限周长集合的周长测度}

由变差对偶，测度记号在开集上与前面的定义一致。
示性函数总是局部可积，但在无界区域上未必属于 \(L^1(\Omega)\)。
因此，有限周长不自动意味着 \(\mathbf1_E\in BV(\Omega)\)；还需要 \(m_d(E)<\infty\)。
例如 \(E=\Omega\) 的相对周长为零，无论 \(\Omega\) 的体积是否有限。
相对周长也不计入位于 \(\partial\Omega\) 上的部分，不能与全空间周长混用。
补集与原集合具有相同的相对周长，因为
\(D\mathbf1_{\Omega\setminus E}=-D\mathbf1_E\) 于 \(\Omega\)。

\begin{lemma}[层集周长的可测性]\label{b1:lem:level-perimeter-measurable}
若 \(u\) 在 \(\Omega\) 上可测，则对每个开集 \(U\subseteq\Omega\)，
函数 \(t\mapsto P(E_t,U)\) 可测，允许取值 \(+\infty\)。
\end{lemma}
\begin{proof}
用紧集 \(K_j\) 穷竭 \(U\)，并要求每个紧支集测试场的支集最终包含在某个 \(K_j\) 内。
固定 \(j\)，考虑支集包含于 \(K_j\)、且 \(\|\varphi\|_\infty\le1\) 的
\(C_c^1\) 测试场。以函数及一阶导数的上确界之和作距离，这个集合可分：
将场及其全部一阶导数限制到 \(K_j\)，它嵌入有限个 \(C(K_j)\) 的乘积，
而紧度量空间上的连续函数空间可分。于是可以在这组测试场内部选一个可数稠密族。
将各 \(j\) 的稠密族合并，记为 \((\varphi_k)\)。

对支集固定的测试场，\(C^1\) 收敛使
\(\int_{E_t}\operatorname{div}\varphi\) 收敛，因为示性函数有界且支集体积有限。
因而对每个 \(t\)，都有
\[
 P(E_t,U)=\sup_k\int_U\mathbf1_{\{u(x)>t\}}
                         \operatorname{div}\varphi_k(x)\,\mathrm dx.
\]
每一项都是可测的含参变量积分：被积函数在共同的紧支集上绝对可积。
可数上确界仍可测，结论成立。
\end{proof}

\subsection{余面积公式}
\label{b1:subsec:270303}

\begin{lemma}[层集积分对变差的控制]\label{b1:lem:layer-perimeter-controls-variation}
设 \(u\in L^1_{\mathrm{loc}}(\Omega)\)。对任意开集 \(U\subseteq\Omega\)，有
\[
 V(u,U)\le\int_{\mathbb R}P(E_t,U)\,\mathrm dt.
\]
\end{lemma}
\begin{proof}
任取 \(\varphi\in C_c^1(U;\mathbb R^d)\)，\(\|\varphi\|_\infty\le1\)。
由\eqref{b1:eq:bv-signed-layer-cake}，且
\[
 \int_{\mathbb R}\int_U
 \bigl|\mathbf1_{\{u>t\}}-\mathbf1_{\{0>t\}}\bigr|
                 \cdot|\operatorname{div}\varphi|\,\mathrm dx\,\mathrm dt
 =\int_U|u|\cdot|\operatorname{div}\varphi|\,\mathrm dx<\infty,
\]
可以使用 Fubini 定理。紧支集又保证
\(\int_U\operatorname{div}\varphi\,\mathrm dx=0\)，故
\[
 \int_Uu\,\operatorname{div}\varphi\,\mathrm dx
 =\int_{\mathbb R}\left(\int_{E_t}\operatorname{div}\varphi\,\mathrm dx\right)\,\mathrm dt
 \le\int_{\mathbb R}P(E_t,U)\,\mathrm dt.
\]
先保留基准项验证绝对可积，再利用散度积分为零将它消去，是有符号分层的关键。
对测试场取上确界即得结论。
\end{proof}

\begin{lemma}[\(W^{1,1}\) 的层集周长公式]\label{b1:lem:w11-perimeter-coarea}
设 \(U\) 为有界开集，\(v\in W^{1,1}(U)\)。则
\[
 \int_{\mathbb R}P(\{v>t\},U)\,\mathrm dt
       =\int_U|\nabla v|\,\mathrm dx.
\]
\end{lemma}
\begin{proof}
取 \(\varepsilon>0\)，定义分段线性截断
\[
 H_\varepsilon(s)=\min\bigl\{1,\max\{0,s/\varepsilon\}\bigr\}.
\]
Sobolev 截断公式及水平集上的梯度为零给出
\[
 H_\varepsilon(v-t)\in W^{1,1}(U),\quad
 \bigl|\nabla H_\varepsilon(v-t)\bigr|
   =\varepsilon^{-1}\mathbf1_{\{t<v<t+\varepsilon\}}|\nabla v|
       \quad\text{几乎处处}.
\]
在固定的层值 \(t\) 上，\(H_\varepsilon(v-t)\) 逐点收敛到
\(\mathbf1_{\{v>t\}}\)，且值在 \([0,1]\) 内。
由于 \(U\) 有限体积，控制收敛给出 \(L^1(U)\) 收敛。
变差下半连续性于是给出
\[
 P(\{v>t\},U)
 \le\liminf_{\varepsilon\downarrow0}
      \varepsilon^{-1}\int_{\{t<v<t+\varepsilon\}}|\nabla v|\,\mathrm dx.
\]
沿任意趋零的正数列应用 Fatou 引理，再用 Tonelli 定理，得到
\[
 \begin{aligned}
 \int_{\mathbb R}P(\{v>t\},U)\,\mathrm dt
 &\le\liminf_{\varepsilon\downarrow0}
   \int_U|\nabla v(x)|\left(
    \varepsilon^{-1}\int_{\mathbb R}
       \mathbf1_{\{t<v(x)<t+\varepsilon\}}\,\mathrm dt\right)\,\mathrm dx\\
 &=\int_U|\nabla v|\,\mathrm dx.
 \end{aligned}
\]
括号中的积分等于 \(1\)，因为固定 \(x\) 后，所允许的 \(t\) 构成长度为
\(\varepsilon\) 的区间。反向不等式由前一引理和
\(V(v,U)=\int_U|\nabla v|\) 得到。
\end{proof}

\begin{theorem}[BV 余面积公式]\label{b1:thm:bv-coarea}
设 \(u\in BV(\Omega)\) 为实值函数，\(E_t=\{u>t\}\)。
则几乎每个 \(E_t\) 在 \(\Omega\) 内具有有限周长，且对每个开集 \(U\subseteq\Omega\)，
\begin{equation}\label{b1:eq:bv-coarea-open}
 |Du|(U)=\int_{\mathbb R}P(E_t,U)\,\mathrm dt.
\end{equation}
在周长无限的零测层值集合上，将周长测度统一记为零，则对任意 Borel 集
\(A\subseteq\Omega\)，都有
\begin{equation}\label{b1:eq:bv-coarea-measure}
 |Du|(A)=\int_{\mathbb R}|D\mathbf1_{E_t}|(A)\,\mathrm dt.
\end{equation}
\end{theorem}
\begin{proof}
先设 \(U\subseteq\Omega\) 有界。将 \(u\) 限制到 \(U\)，用%
\cref{b1:thm:bv-strict-approximation}选
\(u_j\in C^\infty(U)\cap W^{1,1}(U)\)，使
\[
 u_j\to u\quad\text{于 }L^1(U),\quad
 \int_U|\nabla u_j|\,\mathrm dx\longrightarrow|Du|(U).
\]
按\eqref{b1:eq:bv-level-set-l1-distance}抽一个共同子列，使几乎每个层值都有
\(\mathbf1_{\{u_j>t\}}\to\mathbf1_{E_t}\) 于 \(L^1(U)\)。
变差下半连续性和 Fatou 引理给出
\[
 \begin{aligned}
 \int_{\mathbb R}P(E_t,U)\,\mathrm dt
 &\le\liminf_j\int_{\mathbb R}P(\{u_j>t\},U)\,\mathrm dt\\
 &=\lim_j\int_U|\nabla u_j|\,\mathrm dx=|Du|(U).
 \end{aligned}
\]
最后的等号使用前一引理。由层集积分控制变差的引理，又有反向不等式。

对一般开集 \(U\)，用递增有界开集 \(U_j\) 穷竭它。
每个紧支集测试场最终支承于某个 \(U_j\)，故对所有 \(t\)，
\[
 P(E_t,U)=\lim_jP(E_t,U_j).
\]
对周长一侧用单调收敛定理，对 \(|Du|\) 一侧用测度从下连续性，
便得到\eqref{b1:eq:bv-coarea-open}。
取 \(U=\Omega\)，右端有限，所以周长无限的层值构成一个 Lebesgue 零测集 \(Z\)。

还须将开集等式推广到 Borel 集，而不为每个集合另选一组例外层值。
对 \(t\notin Z\)，令 \(\mu_t=|D\mathbf1_{E_t}|\)，对 \(t\in Z\) 令 \(\mu_t=0\)。
这些都是 \(\Omega\) 上的有限 Borel 测度，且对每个开集 \(U\)，
\(t\mapsto\mu_t(U)\) 可测。
使 \(t\mapsto\mu_t(A)\) 可测的 Borel 集 \(A\) 构成一个 Dynkin 系统：
补集利用 \(\mu_t(\Omega)<\infty\) 作差，可数不交并利用单调收敛。
它包含开集，而开集构成生成 Borel 集的 \(\pi\)-系统，故包含全部 Borel 集。
因此
\[
 \lambda(A)=\int_{\mathbb R}\mu_t(A)\,\mathrm dt
\]
定义一个有限 Borel 测度；可数可加性由 Tonelli 定理给出。
开集等式说明 \(\lambda\) 与 \(|Du|\) 在所有开集上一致，
有限测度唯一性遂给出它们在全部 Borel 集上一致。
这就是\eqref{b1:eq:bv-coarea-measure}。
\end{proof}

周长本身始终用分布导数定义。在光滑情形，它与通常边界面积的联系还需相应的
Gauss 公式；上述论证没有将这个联系预先放入定义。
对于一般 BV 函数，几乎所有层集具有有限周长并不意味着每个层集都有光滑边界，
也不意味着例外层值不存在。

\subsection{Cavalieri 型公式}
\label{b1:subsec:270304}

普通的 Cavalieri 分层公式分解体积，BV 余面积公式分解变差。
例如对 \(u\in BV(\Omega)\)，前者给出
\[
 \|u\|_{L^1(\Omega)}
 =\int_0^\infty m_d(\{u>t\})\,\mathrm dt
    +\int_0^\infty m_d(\{u<-t\})\,\mathrm dt,
\]
而变差等式对全部实数层值积分，不能在有符号情形只留下正层值。
例如取 \(u=c\mathbf1_E\)，其中 \(\mathbf1_E\in BV(\Omega)\)、\(c\ne0\)。
若 \(c>0\)，有周长贡献的层位于 \((0,c)\)，这些层集均为 \(E\)；
若 \(c<0\)，对应区间为 \((c,0)\)，层集则为 \(\Omega\setminus E\)。
补集周长不变，两个方向都给出 \(|Du|=|c|\cdot|D\mathbf1_E|\)。

由\eqref{b1:eq:bv-coarea-measure}，对每个非负 Borel 函数 \(g\)，还有加权形式
\begin{equation}\label{b1:eq:bv-weighted-coarea}
 \int_\Omega g\,\mathrm d|Du|
 =\int_{\mathbb R}\left(\int_\Omega g\,\mathrm d|D\mathbf1_{E_t}|\right)\,\mathrm dt.
\end{equation}
先对示性函数和非负简单函数应用测度等式，再作单调逼近，即可同时得到内层积分的可测性与该公式。
取 \(g=\mathbf1_A\)，它描述指定区域内的变差来自哪些层；
取一般权重，则可以分别强调不同空间位置。

分层还精确描述了截断所保留的变差。固定 \(a<b\)，令
\[
 F_{a,b}(s)=\min\bigl\{\max\{s-a,0\},b-a\bigr\}
                -\min\bigl\{\max\{-a,0\},b-a\bigr\}.
\]
减去常数使 \(F_{a,b}(0)=0\)，所以
\(|F_{a,b}(u)|\le|u|\)，在无限体积区域上也保留 \(L^1\) 可积性。

\begin{proposition}[截断的变差分层]\label{b1:prop:bv-truncation-coarea}
若 \(u\in BV(\Omega)\)，则 \(F_{a,b}(u)\in BV(\Omega)\)，且对每个 Borel 集
\(A\subseteq\Omega\)，
\[
 |D F_{a,b}(u)|(A)=\int_a^b|D\mathbf1_{E_t}|(A)\,\mathrm dt.
\]
\end{proposition}
\begin{proof}
截断有有符号表示
\[
 F_{a,b}(u)=\int_a^b
      \bigl(\mathbf1_{\{u>t\}}-\mathbf1_{\{0>t\}}\bigr)\,\mathrm dt.
\]
按层集控制变差引理的测试计算，
\[
 V(F_{a,b}(u),\Omega)\le\int_a^bP(E_t,\Omega)\,\mathrm dt
       \le|Du|(\Omega).
\]
结合 \(L^1\) 可积性，变差对偶的刻画给出 \(F_{a,b}(u)\in BV(\Omega)\)。
记 \(c=\min\bigl\{\max\{-a,0\},b-a\bigr\}\)。
当 \(-c<s<b-a-c\) 时，
\[
 \{F_{a,b}(u)>s\}=\{u>a+c+s\}.
\]
区间以外的超水平集为全区域或空集，只有两个端点层值不必遵循这段表述，
而它们不影响层积分。将这一关系代入 \(F_{a,b}(u)\) 的测度余面积公式，
再作平移换元 \(t=a+c+s\)，即得结论。
\end{proof}

因此，把函数值限制在一段区间内，恰好保留这段层值所贡献的变差。
这种按值域而非按空间切分的方法，将在有限周长集合的逼近以及后续精细结构中继续使用。

% !TeX root = ../../Book1.tex
\section{有限周长集合}
\label{b1:sec:2704}

前一节用示性函数的分布导数定义了周长。这一定义不需要先知道边界是否光滑，甚至不要求集合有有限体积。现在的问题是：导数测度集中在哪里，它的方向能否解释为边界法向？为此必须区分普通拓扑边界、按体积密度定义的边界，以及能从周长测度读出单位方向的约化边界。

\ProseParagraphBreak

本节固定开集 \(\Omega\subset\mathbb R^d\)，令 \(E\subset\mathbb R^d\) 为在 \(\Omega\) 内局部有限周长的可测集。所有边界与支撑均只考察 \(\Omega\) 内的点。需要时可先将 \(E\) 换成几乎处处相等的 Borel 集；下面的对象不会因此改变。

% 实读本地 Simon1983Lectures，正式第14节印刷71--76页；
% 本节特别对应72页(14.1)--(14.2)及极分解后的微分论证。
% 原书用内法向表示Dχ_E/|Dχ_E|；本书外法向取其负号。
% 本节不提前把约化边界认作Hausdorff面积边界，几何识别在下一节证明。

\subsection{周长测度}
\label{b1:subsec:270401}

沿用\cref{b1:def:perimeter}，记
\[
 \mu_E=|D\chi_E|,\quad P(E,A)=\mu_E(A)
\]
对 Borel 集 \(A\subset\Omega\) 成立。局部有限性是说每个 \(K\Subset\Omega\) 上的 \(\mu_E(K)\) 有限，不等于 \(P(E,\Omega)<\infty\)，更不等于 \(m_d(E)<\infty\)。例如半空间的体积与全空间周长都无穷，但在每个有界内部区域中的周长有限。

\ProseParagraphBreak

周长只依赖集合的几乎处处类。若 \(m_d(E\mathbin{\triangle}F)=0\)，则 \(\chi_E=\chi_F\) 几乎处处，所以它们的分布导数及全变差完全相同。对于补集，内部测试函数的积分给出
\begin{equation}\label{b1:eq:perimeter-complement}
 D\chi_{\mathbb R^d\setminus E}=-D\chi_E,\quad
 \mu_{\mathbb R^d\setminus E}=\mu_E.
\end{equation}
这说明周长本身不区分内外，区分内外的是导数测度的方向。

\begin{proposition}\label{b1:prop:perimeter-support-locality}
对 \(x\in\Omega\)，下列条件等价：\(x\in\operatorname{supp}\mu_E\)；每个满足 \(\overline B(x,r)\Subset\Omega\) 的球都有
\[
 0<m_d\bigl(E\cap B(x,r)\bigr)<m_d\bigl(B(x,r)\bigr).
\]
\end{proposition}
\begin{proof}
若某个球内 \(\chi_E\) 几乎处处为 \(0\) 或 \(1\)，则内部所有分布偏导为零，周长测度在该球内也为零，故 \(x\) 不在支撑中。

反过来，若 \(\mu_E\bigl(B(x,r)\bigr)=0\)，则 \(\chi_E\) 在此球中的所有分布偏导为零。局部磨光后，所得光滑函数在每个较小的同心球上梯度为零，因而为常数。让磨光尺度趋零，常数的 \(L^1\) 极限说明 \(\chi_E\) 在每个较小球上几乎处处为常数；重叠处这些常数相同。由于示性函数只取 \(0,1\)，它在原球内也几乎处处取其中一个值。
\end{proof}

这个判据只说明每个球都看得到两侧的正体积，还没有给出两侧所占的比例。某一侧的比例可能随半径缩小而趋于零；不能把周长的支撑直接定义成下面的测度论边界。

\subsection{测度论内部与外部}
\label{b1:subsec:270402}

\begin{definition}\label{b1:def:measure-interior-exterior}
给定可测集 \(E\)，定义它在 \(\Omega\) 内的\term[cedulun-neibu]{测度论内部}[measure-theoretic interior]与\term[cedulun-waibu]{测度论外部}[measure-theoretic exterior]分别为
\[
 E^1\coloneqq\left\{x\in\Omega:
 \lim_{r\downarrow0}\frac{m_d\bigl(E\cap B(x,r)\bigr)}{\omega_dr^d}=1\right\},
\]
\[
 E^0\coloneqq\left\{x\in\Omega:
 \lim_{r\downarrow0}\frac{m_d\bigl(E\cap B(x,r)\bigr)}{\omega_dr^d}=0\right\}.
\]
\end{definition}
\symidx{\(E^1\)：集合的测度论内部}
\symidx{\(E^0\)：集合的测度论外部}

这是\chapref{b1:ch:14}密度点概念的集合记号，而不是另一种拓扑内部。球体积积分的 Borel 性及有理半径逼近说明 \(E^1,E^0\) 都是 Borel 集。\cref{b1:thm:lebesgue-density}给出
\[
 m_d\bigl((E\cap\Omega)\mathbin{\triangle}E^1\bigr)=0,\quad
 m_d\bigl((\Omega\setminus E)\mathbin{\triangle}E^0\bigr)=0.
\]
因此 \(E^1\) 是集合的一个由密度选定的代表，但它未必开。普通内点属于 \(E^1\)，普通外点属于 \(E^0\)；逆命题没有普遍保证。

\ProseParagraphBreak

例如在 \(\Omega=\mathbb R^d\) 中取 \(E=B(0,1)\setminus\mathbb Q^d\)。这个集合没有普通内点，却有 \(E^1=B(0,1)\)；它与开球具有完全相同的示性函数等价类和周长。若使用任意代表的拓扑边界去计算周长，就会把整个闭球误当成边界。

\subsection{测度论边界}
\label{b1:subsec:270403}

\begin{definition}\label{b1:def:measure-theoretic-boundary}
给定可测集 \(E\)，称
\[
 \partial^eE\coloneqq\Omega\setminus(E^0\cup E^1)
\]
为它在 \(\Omega\) 内的\term[cedulun-bianjie]{测度论边界}[measure-theoretic boundary]，也在外文索引中以 essential boundary 登记。\foreignidx{essential boundary}
\end{definition}
\symidx{\(\partial^eE\)：测度论边界}

定义允许密度不存在。因此 \(\partial^eE\) 并不只是密度介于 \(0\) 与 \(1\) 之间且极限存在的点集。它是 Borel 集，并且对任意可测 \(E\) 都有 \(m_d(\partial^eE)=0\)。但在余维一的尺度上，这个集合可以很大。

\ProseParagraphBreak

半空间边界上每个球被平面分成等体积的两半，故其边界点的密度为 \(1/2\)。对于第一象限 \(E=(0,\infty)^2\)，正坐标轴上的非原点边界点仍有密度 \(1/2\)，原点的密度却为 \(1/4\)。它们都属于测度论边界；仅知道某点在 \(\partial^eE\)，还不能给它指定唯一切线或法向。

\ProseParagraphBreak

对于任意选定的集合代表，\(\partial^eE\) 包含于它的普通边界在 \(\Omega\) 中的部分：普通内外点都已经被密度 \(1\) 或 \(0\) 排除。前述删去有理点的例子表明，这个包含可以非常严格。与普通边界不同，\(\partial^eE\) 在零测修改下保持不变。

\subsection{约化边界}
\label{b1:subsec:270404}

由\cref{b1:prop:vector-measure-polar}，写
\[
 D\chi_E=\sigma_E\mu_E,\quad |\sigma_E|=1
 \quad\mu_E\text{-几乎处处}.
\]
一个点附近的方向可能彼此抵消。只有当缩小球中方向的平均趋于单位向量时，才看得到一个没有被抵消的极限方向。

\begin{definition}\label{b1:def:reduced-boundary}
设 \(E\) 在 \(\Omega\) 内局部有限周长。若 \(x\in\operatorname{supp}\mu_E\)，并且极限
\begin{equation}\label{b1:eq:reduced-boundary-direction}
 a_E(x)\coloneqq\lim_{r\downarrow0}
 \frac{D\chi_E\bigl(\overline B(x,r)\bigr)}
 {\mu_E\bigl(\overline B(x,r)\bigr)}
\end{equation}
存在且长度为 \(1\)，则称 \(x\) 为 \(E\) 的约化边界点。全体这些点组成\term[yuehua-bianjie]{约化边界}[reduced boundary]，记为 \(\partial^*E\)。
\end{definition}
\symidx{\(\partial^*E\)：由周长方向极限定义的约化边界}

这里 \(\partial^*E\) 专指约化边界，\(\partial^eE\) 专指测度论边界；不能因文献记号不同而互换二者。式\eqref{b1:eq:reduced-boundary-direction}采用与 Radon 微分定理相同的闭球。将半径从内增大或从外减小，并对向量测度的各分量使用连续性，可知用开球得到的极限完全相同。在支撑上，两种球的分母在充分小的正半径下都为正。因此，这一定义与几何测度论文献中通常采用开球的约化边界定义等价，并未引入另一种边界概念。

\begin{proposition}\label{b1:prop:reduced-boundary-carries-perimeter}
约化边界为 Borel 集，\(a_E\) 可取为其上的 Borel 函数，并且
\[
 \mu_E(\Omega\setminus\partial^*E)=0.
\]
对每个 \(x\in\partial^*E\)，还有
\begin{equation}\label{b1:eq:reduced-direction-oscillation}
 \lim_{r\downarrow0}
 \frac{1}{\mu_E\bigl(B(x,r)\bigr)}
 \int_{B(x,r)}|\sigma_E(y)-a_E(x)|^2\,\mathrm d\mu_E(y)=0.
\end{equation}
\end{proposition}
\begin{proof}
球质量与向量分量的球积分都是 Borel 函数。对固定中心，将闭球半径从外作有理数逼近，利用局部有限测度的连续性，可把极限存在及极限长度为 \(1\) 的条件写成可数个 Borel 条件；支撑本身闭，故约化边界 Borel，极限函数也 Borel。

由\cref{b1:thm:radon-function-differentiation}逐分量作用于 \(\sigma_E\)，球平均在 \(\mu_E\)-几乎处处趋于 \(\sigma_E(x)\)。后者长度为 \(1\)，而支撑的补集为 \(\mu_E\)-零集，故周长集中在约化边界上。最后，固定该边界的一点，利用两个方向的单位长度得到
\[
 \frac{\int_{B(x,r)}|\sigma_E-a_E(x)|^2\,\mathrm d\mu_E}
 {\mu_E\bigl(B(x,r)\bigr)}
 =2-2a_E(x)\cdot
 \frac{D\chi_E\bigl(B(x,r)\bigr)}{\mu_E\bigl(B(x,r)\bigr)},
\]
右端趋于零，证明了最后一式。
\end{proof}

这一命题仍是向量测度的微分结论，没有证明 \(\mu_E\) 等于某种 Hausdorff 面积。下一节的结构定理将说明：约化边界处的集合缩放为半空间，周长缩放为其边界平面的面积测度。

\subsection{测度论法向}
\label{b1:subsec:270405}

\begin{definition}\label{b1:def:measure-theoretic-outer-normal}
对 \(x\in\partial^*E\)，定义 \(E\) 的\term[cedulun-waifaxiang]{测度论外法向}[measure-theoretic outer unit normal]为
\[
 \nu_E(x)\coloneqq-a_E(x).
\]
\end{definition}
\symidx{\(\nu_E(x)\)：有限周长集合的测度论外法向}

负号来自分布导数的约定。例如 \(E=\{x_d<0\}\) 时，直接沿 \(x_d\) 积分得到
\[
 \int_E\partial_d\varphi\,\mathrm dx
 =\int_{\mathbb R^{d-1}}\varphi(x',0)\,\mathrm dx',
\]
所以 \(D\chi_E=-e_d\mathcal H^{d-1}\llcorner\{x_d=0\}\)，而几何上的外法向是 \(e_d\)。一般情形因此写成
\[
 D\chi_E=-\nu_E\mu_E.
\]
法向只需在 \(\partial^*E\) 上指定；在其补集任意补值不会改变测度等式。补集的约化边界与原集相同，而外法向相反，这也与\eqref{b1:eq:perimeter-complement}一致。

\ProseParagraphBreak

第一象限的角点说明单位长度条件有实际作用。由两次一维分部积分，
\[
 D\chi_{(0,\infty)^2}
 =e_1\mathcal H^1\llcorner\bigl(\{0\}\times(0,\infty)\bigr)
 +e_2\mathcal H^1\llcorner\bigl((0,\infty)\times\{0\}\bigr).
\]
原点半径 \(r\) 的球中，两条射线各贡献长度 \(r\)，所以向量与全变差之比为 \((e_1+e_2)/2\)，长度为 \(1/\sqrt2\)。原点属于 \(\partial^eE\)，却不属于 \(\partial^*E\)。这不是在角点随意漏掉一个法向，而是定义准确地记录了两个方向的抵消。

\ProseParagraphBreak

Simon 的《Lectures on Geometric Measure Theory》以 \(D\chi_E/\mu_E\) 表示内法向\cite[\S14, 14.1--14.2]{Simon1983Lectures}。本书选择外法向，是为使下一节的 Gauss–Green 公式保持“区域内散度积分等于向外通量”的符号。阅读其他文献时，应先检查这一等式，再对照法向名称。

% !TeX root = ../../Book1.tex
\section{Gauss–Green 与 De Giorgi 结构}
\label{b1:sec:2705}

上一节已经得到带方向的周长测度，但尚未证明它是一份边界面积。本节把两件事接起来。先对光滑图域直接计算向外通量，再在约化边界点缩放集合；紧性给出极限，方向的集中使极限只能是半空间，而两侧体积下界使其边界必须穿过缩放中心。最后从这些半空间逼近中构造可数张 Lipschitz 图，并辨认周长的密度。

% 实读 Simon1983Lectures，正式第14节印刷71--76页：
% 72--73页(14.1)--(14.3)及定理；73--76页约化点方向、周长上界、
% BV紧性、半空间极限与两侧体积下界的完整证明。
% 本书外法向=-Dχ/|Dχ|，统一把法向e_d对应的极限写成{x_d<0}。
% 原书引用11.6与3.5完成可整流/测度识别；本稿改为显式锥排除图分片，
% 再调用本书14章Radon微分及19章已完整证明的带权密度命题。
% 相对等周用22章一阶球延拓、24章GNS与27.2严格逼近自足推出；
% ∂^eE\∂*E的Hausdorff零性另外给5r证明，不冒称Simon14单独陈述了全部扩充结论。

\subsection{Gauss–Green 公式}
\label{b1:subsec:270501}

先把分布恒等式写成\term*[gauss-green-gongshi]{Gauss--Green 公式}[Gauss--Green formula]。它此时涉及周长测度，面积形式还需要后面的几何识别。

\begin{proposition}[分布形式的 Gauss--Green 公式]\label{b1:prop:gauss-green-distributional}
设 \(E\) 在 \(\Omega\) 中局部有限周长，\(\Phi\in C_c^1(\Omega;\mathbb R^d)\)。则
\begin{equation}\label{b1:eq:gauss-green-perimeter-measure}
 \int_E\operatorname{div}\Phi\,\mathrm dx
 =\int_{\partial^*E}\Phi\cdot\nu_E\,\mathrm d\mu_E.
\end{equation}
\end{proposition}
\begin{proof}
分布导数的定义给出左侧等于 \(-\int\Phi\cdot\mathrm dD\chi_E\)。代入 \(D\chi_E=-\nu_E\mu_E\)，再用\cref{b1:prop:reduced-boundary-carries-perimeter}，即得公式。最初的光滑测试公式可以通过固定紧支集上的 \(C^1\) 逼近传给 \(C_c^1\) 向量场。
\end{proof}

此时边界积分仍是对 \(\mu_E\) 积分，不是已经知道的 Hausdorff 面积积分。对光滑边界，可以不依赖结构定理而直接识别它。

\begin{proposition}[光滑区域的周长]\label{b1:prop:smooth-domain-perimeter}
设开集 \(E\) 在 \(\Omega\) 内具有单侧 \(C^1\) 图边界：每个边界点附近，适当刚性变换后 \(E\) 都是某个 \(C^1\) 函数图像的一侧。则 \(E\) 在 \(\Omega\) 中局部有限周长，而且
\begin{equation}\label{b1:eq:smooth-domain-perimeter}
 D\chi_E=-\nu_E\mathcal H^{d-1}\llcorner(\partial E\cap\Omega),
 \quad
 P(E,A)=\mathcal H^{d-1}(\partial E\cap A)
\end{equation}
对 Borel 集 \(A\subset\Omega\) 成立。这里 \(\nu_E\) 是通常的外单位法向，且 \(\partial^*E=\partial E\cap\Omega\)。
\end{proposition}
\begin{proof}
先设 \(d\ge2\)，并在一个图表内写 \(x=(z,t)\)、\(E=\{t<g(z)\}\)，其中 \(g\in C^1\)。向量场支集紧含图表，所以沿竖直方向积分时，图表上下端没有额外项。微积分基本定理给出
\[
 \int_E\partial_d\Phi_d\,\mathrm dx
 =\int\Phi_d\bigl(z,g(z)\bigr)\,\mathrm dz.
\]
对 \(i<d\)，先对
\(\int_{-\infty}^{g(z)}\Phi_i(z,t)\,\mathrm dt\)
作 \(z_i\) 微分，再对 \(z\) 积分，得到
\[
 \int_E\partial_i\Phi_i\,\mathrm dx
 =-\int\Phi_i\bigl(z,g(z)\bigr)\,\partial_i g(z)\,\mathrm dz.
\]
因而图表内的散度积分等于
\[
 \int \Phi\bigl(z,g(z)\bigr)\cdot(-\nabla g(z),1)\,\mathrm dz.
\]
\cref{b1:thm:graph-area-formula}说明图面积元为
\(\sqrt{1+|\nabla g|^2}\,\mathrm dz\)，而
\[
 \nu_E\bigl(z,g(z)\bigr)
 =\frac{(-\nabla g(z),1)}{\sqrt{1+|\nabla g(z)|^2}}.
\]
这就得到局部公式。

对于一般紧支集测试场，以有限个边界图表及不交边界的内部区域覆盖其支集，再取光滑单位分解。分别计算各个分片后相加，分解系数的梯度项因其和恒为 \(1\) 而抵消。在不交边界的区域，示性函数局部常数，散度积分为零。每个紧集附近只需有限张图，其面积有限，所以所得公式同时证明局部有限周长。

向量测度表示的唯一性给出 \(D\chi_E\) 的等式；法向长度为 \(1\)，故其全变差就是边界面积。法向在每张图上连续，缩小球中的平均趋于该点的法向，于是每个光滑边界点都约化；内外点则不在周长支撑中。一维情形直接对区间使用微积分基本定理，每个内部端点贡献一个单位点质量，结论相同。
\end{proof}

“单侧”不能删除。如果只要求拓扑边界恰为一张光滑曲面，却允许集合占据曲面的两侧，例如删去一张超平面后的全空间，那么示性函数几乎处处恒为 \(1\)，周长实际上为零。

特别地，对球 \(E=B(a,R)\)，有
\[
 \nu_E(x)=\frac{x-a}{R}\quad(x\in\partial B(a,R)),\quad
 P\bigl(B(a,R),A\bigr)=\mathcal H^{d-1}\bigl(\partial B(a,R)\cap A\bigr).
\]
这个计算已经可以独立使用，不必等待一般有限周长集合的结构。

\subsection{约化边界的可整流性}
\label{b1:subsec:270502}

下面的证明先在 \(d\ge2\) 中进行。一维的结构将在定理后单独说明。第一个工具是球内的\term[xiangdui-dengzhou-budengshi]{相对等周不等式}[relative isoperimetric inequality]：如果球内两侧都有一定体积，它们之间就必须有一定周长。

\begin{lemma}[球内相对等周不等式]\label{b1:lem:relative-isoperimetric-ball}
设 \(d\ge2\)，\(\overline B(x,r)\Subset\Omega\)，且 \(E\) 在 \(\Omega\) 内局部有限周长。则
\begin{equation}\label{b1:eq:relative-isoperimetric-ball}
 \min\Bigl(m_d\bigl(E\cap B(x,r)\bigr),m_d\bigl(B(x,r)\setminus E\bigr)\Bigr)^{(d-1)/d}
 \le C_d\mu_E\bigl(B(x,r)\bigr).
\end{equation}
\end{lemma}
\begin{proof}
先对单位球内的 \(v\in W^{1,1}\) 证明
\[
 \|v-v_B\|_{L^{d/(d-1)}(B)}
 \le C_d\|\nabla v\|_{L^1(B)}.
\]
由\cref{b1:thm:ball-poincare-lp}，\(v-v_B\) 的 \(L^1\) 范数受梯度控制。用\cref{b1:thm:lipschitz-domain-extension}把它延拓到全空间，其梯度 \(L^1\) 范数仍有同一控制；再用\cref{b1:thm:sobolev-subcritical-embedding}的 \(p=1\) 情形即得结论。伸缩到任意球时，左侧与梯度积分都按 \(r^{d-1}\) 变化，故常数不依赖球。

将\cref{b1:thm:bv-strict-approximation}用于球中的 \(\chi_E\)。逼近函数在 \(L^1\) 中收敛，其均值也收敛，而梯度积分趋于 \(\mu_E(B)\)。抽取几乎处处收敛子列，用 Fatou 引理得到
\[
 \|\chi_E-a\|_{L^{d/(d-1)}(B)}
 \le C_d\mu_E(B),\quad a=\frac{m_d(E\cap B)}{m_d(B)}.
\]
若 \(a\le1/2\)，左侧至少为
\(\tfrac12m_d(E\cap B)^{(d-1)/d}\)；若 \(a\ge1/2\)，则改在补集上估计。这就证明了相对等周不等式。
\end{proof}

\begin{theorem}[约化边界处的半空间缩放]\label{b1:thm:reduced-boundary-blowup}
设 \(d\ge2\)，\(x\in\partial^*E\)，并记
\[
 E_{x,r}=\{z:x+rz\in E\},\quad
 H_x=\{z:z\cdot\nu_E(x)<0\}.
\]
则当 \(r\downarrow0\) 时，
\begin{equation}\label{b1:eq:reduced-boundary-halfspace}
 \chi_{E_{x,r}}\longrightarrow\chi_{H_x}
 \quad\text{于 }L^1_{\mathrm{loc}}(\mathbb R^d).
\end{equation}
而且对每个 \(\varphi\in C_c(\mathbb R^d)\)，
\begin{equation}\label{b1:eq:perimeter-flat-blowup}
 r^{1-d}\int_\Omega
 \varphi\left(\frac{y-x}{r}\right)\,\mathrm d\mu_E(y)
 \longrightarrow
 \int_{\nu_E(x)^\perp}\varphi\,\mathrm d\mathcal H^{d-1}.
\end{equation}
特别地，每个约化边界点都有集合密度 \(1/2\)，并且
\begin{equation}\label{b1:eq:reduced-perimeter-unit-density}
 \lim_{r\downarrow0}
 \frac{\mu_E\bigl(B(x,r)\bigr)}{\omega_{d-1}r^{d-1}}=1.
\end{equation}
\end{theorem}
\begin{proof}
经平移、旋转，设 \(x=0\)、\(\nu_E(0)=e_d\)。因此极方向的极限为 \(-e_d\)。以下简记 \(\mu=\mu_E\)、\(\sigma=\sigma_E\)，以及
\[
 V(r)=m_d\bigl(E\cap B(0,r)\bigr),\quad
 W(r)=m_d\bigl(B(0,r)\setminus E\bigr).
\]
球壳的体积界说明 \(V,W\) 绝对连续，且它们的导数几乎处处非负、不超过 \(d\omega_dr^{d-1}\)。方向极限给出：对所有充分小的半径，
\[
 \int_{B(0,r)}(-\sigma_d)\,\mathrm d\mu
 \ge\tfrac12\mu\bigl(B(0,r)\bigr).
\]

取在 \(B(0,r)\) 上为 \(1\)、在球壳 \(r<|y|<r+h\) 中线性降到零的径向截断，并用光滑截断逼近它。分布公式给出
\[
 \int(-\sigma_d)\zeta\,\mathrm d\mu
 =\int_E\partial_d\zeta\,\mathrm dx
 \le h^{-1}\bigl(V(r+h)-V(r)\bigr).
\]
在导数存在且球面不带 \(\mu\) 质量的半径令 \(h\downarrow0\)。这些半径占满一维测度，故得到
\begin{equation}\label{b1:eq:reduced-perimeter-radial-bound}
 \mu\bigl(B(0,r)\bigr)\le2V'(r)\le2d\omega_dr^{d-1}
 \quad\text{对几乎每个充分小的 }r.
\end{equation}
补集的外法向相反，同样有 \(\mu\bigl(B(0,r)\bigr)\le2W'(r)\)。由球的单调性，从较大半径向下逼近，还得到
\(\mu\bigl(B(0,r)\bigr)\le C_dr^{d-1}\)
对每个充分小的 \(r\) 成立。

现在排除退化的缩放极限。令 \(M(r)=\min(V(r),W(r))\)。由\cref{b1:prop:perimeter-support-locality}，\(M(r)>0\)；它绝对连续且 \(M(r)\to0\)。结合相对等周估计与刚才的两个导数界，得到
\[
 M(r)^{(d-1)/d}
 \le C_d\mu\bigl(B(0,r)\bigr)
 \le C'_d\min(V'(r),W'(r))
 \le C'_d M'(r)
\]
几乎处处成立。因此 \((M^{1/d})'\ge c_d>0\) 几乎处处。从正半径积分后再令下端趋零，得到
\begin{equation}\label{b1:eq:reduced-two-sided-volume}
 V(r)\ge c_dr^d,\quad W(r)\ge c_dr^d
\end{equation}
对全部充分小的 \(r\) 成立。这一步同时保证两侧不会消失。

对任意趋零序列 \(r_j\)，作变量代换可得
\[
 |D\chi_{E_{0,r_j}}|(B(0,R))
 =r_j^{1-d}\mu\bigl(B(0,r_jR)\bigr)\le C_dR^{d-1}
\]
对每个固定 \(R\) 及充分大的 \(j\) 成立。零阶 \(L^1\) 范数也受球体积控制。逐球使用\cref{b1:thm:bv-local-compactness}并对角抽列，得到
\(\chi_{E_{0,r_j}}\to\chi_F\)
于 \(L^1_{\mathrm{loc}}\)，其中 \(F\) 局部有限周长。这里各缩放定义域最终包含任意固定球；极限仍为示性函数，可由进一步的几乎处处收敛看出。

式\eqref{b1:eq:reduced-direction-oscillation}和 Cauchy 不等式说明
\[
 r^{1-d}\int_{B(0,rR)}|\sigma+e_d|\,\mathrm d\mu
 \longrightarrow0
\]
对固定 \(R\) 成立。于是缩放后的分布导数在前 \(d-1\) 个方向趋于零，在第 \(d\) 个方向的极限是非正分布。故
\[
 \partial_i\chi_F=0\quad(i<d),\quad
 \partial_d\chi_F\le0.
\]
将 \(\chi_F\) 磨光：前一组等式使光滑函数只依赖最后一个变量，后一不等式使它对这个变量非增。让磨光尺度趋零，就得到 \(\chi_F(z,t)=h(t)\) 几乎处处，其中 \(h\) 有非增代表且几乎处处只取 \(0,1\)。所以 \(F\) 几乎处处等于 \(\{t<b\}\)，其中暂允许 \(b=\pm\infty\)。

对任意固定 \(s>0\)，将\eqref{b1:eq:reduced-two-sided-volume}伸缩并传给极限，球 \(B(0,s)\) 中 \(F\) 及其补集的体积都至少为 \(c_ds^d\)。若 \(b\ne0\)，取足够小的 \(s\)，其中一侧就会为空，矛盾。因此 \(b=0\)。任意初始半径序列都有子列趋于同一个半空间，这也证明了全部半径下的\eqref{b1:eq:reduced-boundary-halfspace}。

最后记缩放周长测度为
\(\mu_r(A)=r^{1-d}\mu(rA)\)。
对 \(\varphi\in C_c^1\)，方向误差趋零给出
\[
 \int\varphi\,\mathrm d\mu_r
 =-\int\varphi\,\mathrm d(D_d\chi_{E_{0,r}})+o(1)
 =\int_{E_{0,r}}\partial_d\varphi\,\mathrm dx+o(1)
 \longrightarrow\int_{\{t=0\}}\varphi\,\mathrm d\mathcal H^{d-1}.
\]
紧集上的一致质量界允许用光滑函数一致逼近一般 \(C_c\) 测试函数，得到\eqref{b1:eq:perimeter-flat-blowup}。平面测度不给单位球面质量，故用内外连续截断逼近球的示性函数，可得\eqref{b1:eq:reduced-perimeter-unit-density}。集合密度 \(1/2\) 则直接来自半空间的体积及 \(L^1\) 收敛。
\end{proof}

半空间极限已经给出法向的几何意义，但可数图分片仍须构造。下面把每点成立的极限整理为统一尺度上的估计。

\begin{proposition}\label{b1:prop:reduced-boundary-rectifiable}
当 \(d\ge2\) 时，\(\partial^*E\) 可由可数张完整 Lipschitz 图覆盖，因而是可数 \(d-1\) 维可整流 Borel 集。
\end{proposition}
\begin{proof}
令 \(c_d\) 为\eqref{b1:eq:reduced-two-sided-volume}中的统一常数，固定
\(0<\varepsilon<c_d/8^d\)。
按极限开始满足估计的尺度，把约化边界分成可数个 Borel 片，同时要求 \(B(x,1/m)\subset\Omega\)，使每片中各点 \(x\) 对 \(0<r<1/m\) 都满足两侧体积下界，以及
\[
 m_d\bigl((E\mathbin{\triangle}\{y:(y-x)\cdot\nu_E(x)<0\})
                   \cap B(x,r)\bigr)\le\varepsilon r^d.
\]
球体积函数对半径连续，故只检查有理半径便可保证这些片 Borel。再按法向所落入的球面小邻域分片，使一片上
\(|\nu_E(x)-\nu_0|<1/8\)，其中 \(\nu_0\) 固定；最后与直径小于 \(1/(4m)\) 的可数小立方体相交。

取同一片中不同的 \(x,y\)，记 \(\ell=|x-y|\)。若
\(|(y-x)\cdot\nu_E(x)|>\ell/2\)，则 \(B(y,\ell/4)\) 完全位于以 \(x\) 为中心的极限半空间的一侧，并且包含于 \(B(x,2\ell)\)。在 \(y\) 处的两侧体积下界说明，这个小球对半空间的体积误差至少为 \(c_d(\ell/4)^d\)；在 \(x\) 处的误差估计却使它至多为 \(\varepsilon(2\ell)^d\)，与 \(\varepsilon\) 的选择矛盾。因此
\[
 |(y-x)\cdot\nu_0|
 \le |(y-x)\cdot\nu_E(x)|+\tfrac18|y-x|
 \le\tfrac58|y-x|.
\]
投影到 \(\nu_0^\perp\) 后，距离至少为 \(\sqrt{39}\,|y-x|/8\)。所以投影在该片上单射，法向坐标是投影坐标的 Lipschitz 函数，常数至多为 \(5/\sqrt{39}\)。由\cref{b1:thm:mcshane-extension}将这个实值函数延拓到整个 \(\nu_0^\perp\)，便得到包含该片的一张完整 Lipschitz 图。所有分片均可数，证明完成。
\end{proof}

这个论证没有从“极限是平面”直接跳到可整流性。关键是用两侧体积下界排除同一片上沿法向分离过大的两个点，进而让固定的正交投影成为单射。它也没有调用\chapref{b1:ch:19}只给出证明概要的一般密度逆判据。

\subsection{De Giorgi 结构定理}
\label{b1:subsec:270503}

\term*[de-giorgi-jiegou-dingli]{De Giorgi 结构定理}[De Giorgi structure theorem]把上述几何逼近与周长测度的表示合在一起。

\begin{theorem}[De Giorgi]\label{b1:thm:de-giorgi-structure}
设 \(E\) 在开集 \(\Omega\subset\mathbb R^d\) 中局部有限周长。当 \(d\ge2\) 时，\(\partial^*E\) 可数 \(d-1\) 维可整流；当 \(d=1\) 时，它是局部有限的点集。两种情形均有
\begin{equation}\label{b1:eq:de-giorgi-perimeter-measure}
 |D\chi_E|=\mathcal H^{d-1}\llcorner\partial^*E.
\end{equation}
此外，
\begin{equation}\label{b1:eq:essential-reduced-boundary}
 \partial^*E\subset\partial^eE,\quad
 \mathcal H^{d-1}(\partial^eE\setminus\partial^*E)=0.
\end{equation}
每个 \(x\in\partial^*E\) 的集合密度为 \(1/2\)。当 \(d\ge2\) 时，在\chapref{b1:ch:19}的缩放测试意义下，
\[
 \operatorname{Tan}^{d-1}(\partial^*E,x)=\nu_E(x)^\perp.
\]
\end{theorem}
\begin{proof}
先设 \(d\ge2\)，记 \(S=\partial^*E\)、\(k=d-1\)、\(\mu=\mu_E\)。可整流性和单位周长密度已经证明，且 \(\mu\) 集中在 \(S\) 上。先证明 \(\mu\ll\mathcal H^k\llcorner S\)。把 \(S\) 分成可数个片，在每片上都有
\[
 \mu\bigl(B(x,r)\bigr)\le C_dr^k\quad(0<r<1/m),
\]
并限制在距 \(\Omega\) 外部有正距离的有界区域内。若片中的集合 \(N\) 为 \(\mathcal H^k\)-零集，用直径小于 \(1/(2m)\) 的集合覆盖它，使直径的 \(k\) 次和任意小。每个覆盖集若与 \(N\) 相交，就以交点为中心取半径为其直径两倍的球；这些球仍覆盖 \(N\)，上式控制其总 \(\mu\) 质量。直径为零的单点没有质量，也由该上界看出。于是 \(\mu(N)=0\)。对各片取并，即得绝对连续性。

可数图覆盖保证 \(\mathcal H^k\llcorner S\) 是 \(\sigma\) 有限测度，故 Radon–Nikodym 定理给出
\[
 \mu=\theta\mathcal H^k\llcorner S.
\]
由已完整证明的\cref{b1:prop:rectifiable-weight-density}，\(\mu\) 的 \(k\) 维密度在 \(\mathcal H^k\)-几乎处处等于 \(\theta\)。而\eqref{b1:eq:reduced-perimeter-unit-density}说明该密度在每个 \(S\) 中的点都为 \(1\)，所以 \(\theta=1\) 几乎处处，证明了测度表示。特别地，未加权的约化边界面积现在也已知局部有限。将测度表示代入\eqref{b1:eq:perimeter-flat-blowup}，正好得到\chapref{b1:ch:19}定义中的单位重数切平面，故其方向为 \(\nu_E(x)^\perp\)。

半空间极限已经证明 \(S\subset\partial^eE\)。为了处理余项，注意：若 \(x\in\partial^eE\)，则存在 \(\delta_x>0\) 及任意小的半径 \(r\)，使球内两侧的体积分数均不小于 \(\delta_x\)。否则，连续函数
\[
 r\longmapsto\frac{m_d\bigl(E\cap B(x,r)\bigr)}{\omega_dr^d}
\]
在足够小半径下总落在分别靠近 \(0\)、\(1\) 的两个不交区间内；连续性使它只能停留在其中一侧，最终便趋于 \(0\) 或 \(1\)，与 \(x\in\partial^eE\) 矛盾。相对等周不等式于是给出任意小半径上的
\[
 \mu\bigl(B(x,r)\bigr)\ge c_xr^k,\quad c_x>0.
\]
将 \(N=\partial^eE\setminus S\) 按
\(\limsup_{r\downarrow0}\mu\bigl(B(x,r)\bigr)/r^k>1/j\)
分片，并与内部紧穷竭相交。球质量的 Borel 性及有理半径逼近说明所得各片 Borel；每片上都有任意小半径满足 \(\mu\bigl(B(x,r)\bigr)\ge r^k/j\)。固定其中一片 \(N_j\)，它是 \(\mu\)-零集。给定 \(\varepsilon>0\)，在一个内部有限质量的邻域中选取包含它的开集 \(G\)，使 \(\mu(G)<\varepsilon\)，再对每个点选取任意小、包含于 \(G\) 且满足上述下界的球。由\cref{b1:lem:five-r-covering}选出不交球 \(B_i=B(x_i,r_i)\)，其五倍球覆盖该片。因此在任意小的 Hausdorff 覆盖尺度上，
\[
 \mathcal H^k_{\delta}(N_j)
 \le C_d\sum_i r_i^k
 \le C_dj\sum_i\mu(B_i)
 \le C_dj\,\varepsilon.
\]
先令覆盖尺度趋零，再令 \(\varepsilon\) 趋零，得到 \(\mathcal H^k(N_j)=0\)。这证明了\eqref{b1:eq:essential-reduced-boundary}。

最后处理 \(d=1\)。在任意内部有界区间 \(I=(a,b)\) 上，记 \(m=D\chi_E\)，并令 \(F(t)=m((a,t])\)。Fubini 定理直接给出 \(DF=m\)，因此 \(\chi_E-F\) 的分布导数为零，几乎处处等于常数。这样得到 \(\chi_E\) 的右连续有限变差代表 \(u\)。它处处只能取 \(0,1\)：从任一点的右侧选取避开原零测例外集的点列，再用右连续性即可。

每次从 \(0\) 变到 \(1\) 或反过来都消耗至少一个单位的变差，所以任一紧区间上只能有有限次这样的变化。于是 \(E\) 在零测集修改后局部为有限个区间的并，\(D\chi_E\) 在左端点取 \(+1\)、右端点取 \(-1\)，全变差就是这些端点的计数测度。它们恰为约化边界及测度论边界，密度为 \(1/2\)，外法向分别为 \(-1,+1\)。这给出所有结论的一维版本。
\end{proof}

证明的主干可与 Simon 的《Lectures on Geometric Measure Theory》定理14.3及其证明对读\cite[第14节，第72--76页]{Simon1983Lectures}。其中“约化点处的方向集中—缩放紧性—半空间极限”的顺序尤其重要。本书另外把图分片构造、密度识别及测度论边界余项的覆盖论证展开，以便这些结论只依赖已经建立的工具。

\subsection{周长测度的表示}
\label{b1:subsec:270504}

现在可以把前面的分布公式写成熟悉的面积通量形式。

\begin{theorem}[Gauss--Green]\label{b1:thm:gauss-green-finite-perimeter}
设 \(E\) 在 \(\Omega\) 中局部有限周长。对任意 \(\Phi\in C_c^1(\Omega;\mathbb R^d)\)，有
\begin{equation}\label{b1:eq:gauss-green-hausdorff}
 \int_E\operatorname{div}\Phi\,\mathrm dx
 =\int_{\partial^*E}\Phi\cdot\nu_E\,\mathrm d\mathcal H^{d-1}.
\end{equation}
对任意 Borel 集 \(A\subset\Omega\)，
\[
 P(E,A)=\mathcal H^{d-1}(A\cap\partial^*E)
       =\mathcal H^{d-1}(A\cap\partial^eE).
\]
\end{theorem}
\begin{proof}
将\eqref{b1:eq:de-giorgi-perimeter-measure}代入\eqref{b1:eq:gauss-green-perimeter-measure}，得到通量公式。周长的两种表示分别来自同一测度等式及\eqref{b1:eq:essential-reduced-boundary}。
\end{proof}

第二种周长表示不意味着 \(\partial^eE\) 的每一点都有法向。第一象限的角点已经给出反例。积分能够忽略这些点，是因为结构定理证明它们在 \(\mathcal H^{d-1}\) 下构成零集；不是在定义法向时可以随意省去任意坏点。

也不能把 \(\partial^*E\) 换成任意代表的普通边界。给一个光滑集合删去内部稠密零测集，不会改变 \(D\chi_E\)、周长、约化边界或测度论边界，却可以让普通边界包含整个原区域。因此有限周长理论的几何对象是由等价类恢复的边界，而不是某次点集表示附带的拓扑边界。

这一表示还说明，示性函数的导数是集中在余维一可整流集上的测度；一般 BV 函数的变化则由全部水平集共同组成。与\cref{b1:thm:bv-coarea}合起来，函数的总变差便可以理解为各层约化边界面积的累积。不过本节没有把所有 BV 导数都断言为单独一张超曲面上的面积：不同水平的边界可以形成完全不同的整体分布。

% !TeX root = ../../Book1.tex
\section{等周不等式}
\label{b1:sec:2706}

周长描述边界的代价，体积描述集合占据的空间。两者的联系不仅属于几何：
将一个函数拆成超水平集，集合的体积估计就能转化为函数的可积性估计。
反过来，把函数不等式用于示性函数，又会恢复集合不等式。
本节在全空间完整建立这两个方向；所证明的是只依赖维数的常数，并不预先宣称其最优。

Simon 的《Lectures on Geometric Measure Theory》\cite[\S6, pp. 34--40]{Simon1983Lectures}
把局部变差、磨光与 Poincaré 型估计放在同一条线上。
这里将前节的严格逼近接到\secref{b1:sec:2402}已经证明的端点 Sobolev 不等式，
再用余面积公式说明其集合意义。

% 实读范围：本地 Simon1983Lectures 正式扫描本，第6节印刷34--40页，
% 尤其6.1、6.2及6.4--6.6。该范围用于局部变差/磨光/Poincare语境，
% 不冒称提供最佳欧氏等周常数或球等号唯一性的证明。
% 主证明直接依赖本书24章GNS端点、27.2严格逼近、27.3余面积。
% 全空间Cc∞严格逼近在下面实际补做远处截断，不把任意域光滑逼近写成紧支集逼近。

\subsection{有界变差函数形式}
\label{b1:subsec:270601}

\begin{theorem}[有界变差函数的 Sobolev 端点估计]\label{b1:thm:bv-sobolev-endpoint}
设 \(d\ge2\)，\(q=d/(d-1)\)。存在仅依赖维数的常数 \(S_d<\infty\)，使
每个 \(u\in BV(\mathbb R^d)\) 都属于 \(L^q(\mathbb R^d)\)，且
\begin{equation}\label{b1:eq:bv-sobolev-endpoint}
 \|u\|_{L^q(\mathbb R^d)}\le S_d|Du|(\mathbb R^d).
\end{equation}
\end{theorem}
\begin{proof}
先把严格逼近改成全空间紧支集逼近。由定理%
\ref{b1:thm:bv-strict-approximation}，可取
\(w_j\in C^\infty(\mathbb R^d)\cap W^{1,1}(\mathbb R^d)\)，使
\[
 \|w_j-u\|_1\to0,\quad
 \int_{\mathbb R^d}|\nabla w_j|\,\mathrm dx
       \longrightarrow|Du|(\mathbb R^d).
\]
选光滑截断 \(0\le\chi_R\le1\)，在 \(B(0,R)\) 上等于 \(1\)，
支集包含于 \(B(0,2R)\)，且 \(|\nabla\chi_R|\le C/R\)。
对固定 \(j\)，当 \(R\to\infty\) 时，
\[
 \|\chi_Rw_j-w_j\|_1\to0,\quad
 \int_{\mathbb R^d}|w_j|\cdot|\nabla\chi_R|\,\mathrm dx
       \le\frac C R\|w_j\|_1\to0.
\]
因此可取 \(R_j\ge j\)，使这两个量都小于 \(1/j\)。
令 \(v_j=\chi_{R_j}w_j\)，则 \(v_j\in C_c^\infty(\mathbb R^d)\)，
\(v_j\to u\) 于 \(L^1\)，且由乘积公式，
\[
 \limsup_j\int_{\mathbb R^d}|\nabla v_j|\,\mathrm dx
 \le\lim_j\left(\int_{\mathbb R^d}|\nabla w_j|\,\mathrm dx+\frac1j\right)
 =|Du|(\mathbb R^d).
\]
下半连续性给出反向的下极限不等式，所以
\(\int|\nabla v_j|\to|Du|(\mathbb R^d)\)。

现在对 \(v_j\) 应用\cref{b1:thm:gns-endpoint}：
\[
 \|v_j\|_q\le S_d\int_{\mathbb R^d}|\nabla v_j|\,\mathrm dx.
\]
再从 \(L^1\) 收敛中抽取几乎处处收敛子列，利用 Fatou 引理，得到
\[
 \|u\|_q
 \le\liminf_j\|v_j\|_q
 \le S_d|Du|(\mathbb R^d).
\]
这一步同时证明目标 \(L^q\) 可积性，不要求预先知道 \(u\in L^q\)。
\end{proof}

右端没有零阶项，但左端函数类仍要求 \(u\in L^1(\mathbb R^d)\)。
这一条件使远处截断的误差消失，也排除了全空间上的非零常数。
严格逼近在这里提供的是总变差的数值收敛，不是导数测度的变差范数收敛；
证明也没有推出 \(v_j\to u\) 于临界空间 \(L^q\)。

维数 \(d=1\) 时，对 \(v\in C_c^\infty(\mathbb R)\) 分别从左右无穷远积分，可得
\[
 |v(x)|\le\int_{-\infty}^x|v'|,\quad
 |v(x)|\le\int_x^\infty|v'|,
 \quad
 2\|v\|_\infty\le\int_{\mathbb R}|v'|.
\]
沿上面同一组紧支集严格逼近，并在几乎处处极限上取界，得到
\[
 \|u\|_{L^\infty(\mathbb R)}\le\frac12|Du|(\mathbb R),
 \quad u\in BV(\mathbb R).
\]
这是单独的一维端点，不把 \(d/(d-1)\) 当作有限指数使用。

\subsection{有限周长集合形式}
\label{b1:subsec:270602}

\begin{theorem}[有限周长集合的等周估计]\label{b1:thm:finite-perimeter-isoperimetric}
设 \(d\ge2\)，集合 \(E\subseteq\mathbb R^d\) 可测、体积有限且周长有限。
则
\begin{equation}\label{b1:eq:finite-perimeter-isoperimetric}
 m_d(E)^{(d-1)/d}\le S_dP(E).
\end{equation}
\end{theorem}
\begin{proof}
有限体积给出 \(\mathbf1_E\in L^1(\mathbb R^d)\)，有限周长给出
\(\mathbf1_E\in BV(\mathbb R^d)\)。将它代入\eqref{b1:eq:bv-sobolev-endpoint}，
并注意
\[
 \|\mathbf1_E\|_{d/(d-1)}=m_d(E)^{(d-1)/d},
 \quad |D\mathbf1_E|(\mathbb R^d)=P(E),
\]
便得结论。
\end{proof}

这类将体积与边界大小联系起来的估计称为%
\term[dengzhou-budengshi]{等周不等式}[isoperimetric inequality]。
这里使用的是全空间周长。若将右端直接改为 \(P(E,\Omega)\)，
取 \(E=\Omega\) 就会出现正体积与零相对周长，因此必须另加适合区域的归一化。
有限体积也不可删除：全空间自身的周长为零，却有无穷体积。
在一维，前面的端点界说明 \(0<m_1(E)<\infty\) 的有限周长集合满足 \(P(E)\ge2\)；
一个有界非空区间达到这个值。

球可以用来检查常数的尺度。记 \(\omega_d=m_d(B(0,1))\)。
由光滑域周长公式\ref{b1:prop:smooth-domain-perimeter}及球面积计算，
\[
 m_d(B(a,r))=\omega_dr^d,\quad
 P(B(a,r))=d\omega_dr^{d-1}.
\]
所以任何适用于所有集合的常数 \(S_d\) 都必须满足
\[
 S_d\ge\frac1{d\omega_d^{1/d}}.
\]
半径恰好消去，说明指数 \((d-1)/d\) 与边界尺度相匹配。
但求出球的比值，并不证明这个下界对所有集合都可作为不等式常数，
更不证明哪些集合能取等号。本节不从非最优的 Gagliardo–Nirenberg–Sobolev 估计
推断最佳等周常数或等号分类。

\subsection{Sobolev 与等周不等式的联系}
\label{b1:subsec:270603}

从函数不等式到集合不等式只需代入示性函数。反方向则需要把函数的全部层重新组合，
余面积公式正好将累积的周长转回总变差。

\begin{proposition}[集合与函数估计的常数传递]\label{b1:prop:isoperimetric-sobolev-equivalence}
固定 \(d\ge2\)、\(q=d/(d-1)\) 及 \(C>0\)。下列两条等价。
\begin{equivalences}
 \item\label{b1:item:bv-isoperimetric-set-bound}
 每个有限体积的有限周长集合 \(E\subseteq\mathbb R^d\) 都满足
 \(m_d(E)^{1/q}\le CP(E)\)。
 \item\label{b1:item:bv-isoperimetric-function-bound}
 每个实值 \(u\in BV(\mathbb R^d)\) 都满足
 \(\|u\|_q\le C|Du|(\mathbb R^d)\)。
\end{equivalences}
\end{proposition}
\begin{proof}
由\itemref{b1:item:bv-isoperimetric-function-bound}推出%
\itemref{b1:item:bv-isoperimetric-set-bound}已在前一定理中证明。
现假设集合估计成立。

先说明可以取绝对值而不增加总变差。用严格逼近取
\(u_j\in W^{1,1}(\mathbb R^d)\)，使 \(u_j\to u\) 于 \(L^1\)，且
\(\int|\nabla u_j|\to|Du|(\mathbb R^d)\)。
Sobolev 绝对值公式给出 \(|u_j|\in W^{1,1}\)，
\(\bigl|\nabla|u_j|\bigr|\le|\nabla u_j|\) 几乎处处。
又有 \(\||u_j|-|u|\|_1\le\|u_j-u\|_1\)，故变差下半连续性推出
\[
 w=|u|\in BV(\mathbb R^d),\quad
 |Dw|(\mathbb R^d)\le|Du|(\mathbb R^d).
\]

令 \(E_t=\{w>t\}\)，\(t>0\)。由于 \(w\in L^1\)，有
\(m_d(E_t)\le\|w\|_1/t<\infty\)；余面积公式又保证几乎每个 \(E_t\) 周长有限。
对 \(R,N>0\)，用积分 Minkowski 不等式及集合估计，得到
\[
 \begin{aligned}
 \|\mathbf1_{B(0,R)}\min\{w,N\}\|_q
 &=\left\|\int_0^N\mathbf1_{E_t\cap B(0,R)}\,\mathrm dt\right\|_q\\
 &\le\int_0^N m_d\bigl(E_t\cap B(0,R)\bigr)^{1/q}\,\mathrm dt\\
 &\le\int_0^N m_d(E_t)^{1/q}\,\mathrm dt\\
 &\le C\int_0^\infty P(E_t)\,\mathrm dt
   =C|Dw|(\mathbb R^d).
 \end{aligned}
\]
空间与层值的截断保证最初的函数有界且支集有限，因而没有预设 \(w\in L^q\)。
令 \(R,N\to\infty\)，单调收敛定理给出
\[
 \|u\|_q=\|w\|_q\le C|Dw|(\mathbb R^d)\le C|Du|(\mathbb R^d).
\]
同一个常数 \(C\) 在两个方向均未改变，证明完成。
\end{proof}

这个等价性说明，最佳集合常数与最佳有界变差函数常数若按各自下确界定义，必然相同；
但它本身并不计算那个下确界。
在本书当前的证明链中，先由坐标积分建立 Sobolev 估计，再由严格逼近与余面积转到集合，
所有实际使用的常数已经得到控制。若转而从尖锐的几何等周定理出发，
同一层积分论证也能将它的常数原样传给有界变差函数的端点估计。


\begin{chaptersummary}
有界变差函数以有限导数测度控制变化，既容纳可积梯度，也容纳跳跃及更分散的奇异变化。
光滑逼近保持函数的 \(L^1\) 极限和总变差数值，却通常不能保证导数测度在变差范数中收敛。
紧性来自平移控制。本章在任意开集上证明的全域版本还要求质量不逃向边界或无穷远。
对有界 Lipschitz 区域，标准紧嵌入可由另需建立的有界变差函数延拓定理得到，不必另加这种尾部假设。

余面积公式把导数变差分解为水平集周长。对每个有限周长集合，
De Giorgi 定理进一步把周长识别成约化边界的余维一面积，
其证明先取得非退化半空间极限，再构造可数 Lipschitz 图并辨认密度。
Gauss–Green 公式中的外法向负号须与分布导数的约定配合；换用其他文献的记号时，也应同步调整。

全空间有界变差函数的端点不等式与有限体积集合的等周不等式具有相同的常数传递机制。
本章给出有效常数，不讨论最佳常数。
一般有界变差函数并非单个集合的示性函数；各层边界的累积还会形成怎样的导数，
将由下一章的精细结构解释。
\end{chaptersummary}

\BookExercises

% \chapref{b1:ch:27}习题临时稿；不另设 BookExercises 入口。
% 九题均由本章与前章已证明工具自拟，逐题注明实质增量，不冒认教材原题号。
% 不调用第28章的一般有界变差函数链式法则、跳跃集结构或特殊有界变差函数理论。

\begin{exercise}\label{b1:ex:27-grid-interface-cancellation}
取正整数 \(m,n\)，令 \(\Omega=(0,m)\times(0,n)\)。
在每个单位格
\[
 Q_{ij}=(i-1,i)\times(j-1,j),\quad
 1\le i\le m,\quad1\le j\le n,
\]
上令 \(u=a_{ij}\)，其中 \(a_{ij}\in\mathbb R\)。
网格线上的值任取。求 \(Du\) 和 \(|Du|(\Omega)\)，并求全空间零延拓 \(E_0u\) 的总变差。
解释为什么内部界面的贡献是相邻值之差，而不是两个值的绝对值之和。
\end{exercise}
% 来源：自拟；把单个矩形的导数推广到共享界面，重点是先相消再取向量变差。
\begin{solution}
对光滑紧支集测试函数逐格积分。固定一条竖界面
\(V_{ij}=\{i\}\times(j-1,j)\)，左格在右侧贡献
\(-a_{ij}e_1\mathcal H^1\)，右格在左侧贡献
\(a_{i+1,j}e_1\mathcal H^1\)。
水平界面 \(H_{ij}=(i-1,i)\times\{j\}\) 同理。因此在 \(\Omega\) 中，
\[
 \begin{aligned}
 Du={}&\sum_{i=1}^{m-1}\sum_{j=1}^n
       (a_{i+1,j}-a_{ij})e_1\,\mathcal H^1\!\llcorner V_{ij}\\
 &+\sum_{i=1}^m\sum_{j=1}^{n-1}
       (a_{i,j+1}-a_{ij})e_2\,\mathcal H^1\!\llcorner H_{ij}.
 \end{aligned}
\]
不同线段只可能在端点相交，这些点的 \(\mathcal H^1\) 测度为零。
每条线段长度为一，所以
\[
 |Du|(\Omega)=
 \sum_{i=1}^{m-1}\sum_{j=1}^n|a_{i+1,j}-a_{ij}|
 +\sum_{i=1}^m\sum_{j=1}^{n-1}|a_{i,j+1}-a_{ij}|.
\]
外边界在相对导数中没有贡献。补零后，下侧与左侧的系数为内部值，
上侧与右侧的系数为内部值的负号；它们与内部界面互不重叠到面测度零集。
于是
\[
 \begin{aligned}
 |D(E_0u)|(\mathbb R^2)
 ={}&|Du|(\Omega)
    +\sum_{j=1}^n(|a_{1j}|+|a_{mj}|)\\
   &+\sum_{i=1}^m(|a_{i1}|+|a_{in}|).
 \end{aligned}
\]
若相邻两格取同一个值，跨界面并没有变化，故该界面必须完全抵消。
先对每格分别取变差再相加，会丢失导数测度的这个抵消过程。
\end{solution}

\begin{exercise}\label{b1:ex:27-small-balls-perimeter-threshold}
设 \(d\ge2\)、\(\alpha>1/d\)，并令
\[
 x_j=3je_1,\quad r_j=\frac14j^{-\alpha},\quad
 E=\bigcup_{j=1}^\infty B(x_j,r_j).
\]
证明 \(m_d(E)<\infty\)，并确定哪些 \(\alpha\) 使 \(E\) 有有限周长。
在允许无穷的意义下计算 \(P(E)\)，说明有限体积不能控制周长。
\end{exercise}
% 来源：自拟；用局部有限且相隔的小球把体积和边界的两个尺度化成不同级数。
\begin{solution}
这些闭球两两不交，而且每个紧集只与有限个球相交。
由体积的可数可加性，
\[
 m_d(E)=\omega_d4^{-d}\sum_{j=1}^\infty j^{-\alpha d}<\infty.
\]
对任意紧支集测试场，只有有限个球参与积分。
光滑域周长公式因而给出局部向量测度
\[
 D\mathbf1_E
 =-\sum_{j=1}^\infty
      \nu_j\,\mathcal H^{d-1}\!\llcorner\partial B(x_j,r_j),
 \quad \nu_j(x)=\frac{x-x_j}{r_j}.
\]
这是局部有限的和；各球面互不相交，故变差没有相消，且
\[
 P(E)=d\omega_d\sum_{j=1}^\infty r_j^{d-1}
     =d\omega_d4^{1-d}\sum_{j=1}^\infty j^{-\alpha(d-1)}.
\]
若级数收敛，右侧给出全空间有限导数测度，所以 \(E\) 有有限周长。
若级数发散，包含前 \(N\) 个球面的有界开集上的变差至少为前 \(N\) 项之和，
让 \(N\to\infty\) 即知全空间周长无穷。
因此
\[
 P(E)<\infty\quad\Longleftrightarrow\quad
       \alpha>\frac1{d-1}.
\]
特别地，\(1/d<\alpha\le1/(d-1)\) 时体积有限而周长无限。
两个端点不同，是因为球体积按 \(r^d\) 缩放，而球周长按 \(r^{d-1}\) 缩放。
\end{solution}

\begin{exercise}\label{b1:ex:27-basic-boundaries-and-blowup}
分别考察下列集合：
\[
 H=\{x\in\mathbb R^d:x_d<0\},\qquad
 Q=(0,\infty)^2\subset\mathbb R^2,
\]
以及
\[
 F=B(0,1)\setminus\mathbb Q^d\subset\mathbb R^d.
\]
对每个集合求其密度零点集、密度一点集、测度论边界与约化边界，并在约化边界上求测度论外法向。比较 \(F\) 的拓扑边界与测度论边界。最后直接验证：半空间在边界点的缩放恒为相应半空间；球在每个边界点的集合缩放与周长测度缩放分别趋于切半空间与切平面测度。说明这些结论为何也适用于 \(F\)。
\end{exercise}
% 来源：自拟；用半空间、角点及稠密零测修改训练三种边界与缩放极限的基本辨析。
\begin{solution}
先看半空间。直接由球的对称性，
\[
 H^1=H,\qquad H^0=\{x_d>0\},\qquad
 \partial^eH=\partial^*H=\{x_d=0\}.
\]
在边界上，外法向为 \(\nu_H=e_d\)，并且
\[
 D\mathbf1_H=-e_d\,\mathcal H^{d-1}
                    \!\llcorner\{x_d=0\}.
\]
若 \(x\in\partial H\)，则 \(H_{x,r}=H\) 对每个 \(r>0\) 都成立；相应的缩放周长测度也始终就是 \(\mathcal H^{d-1}\llcorner\{x_d=0\}\)。

对第一象限，
\[
 Q^1=Q,\qquad Q^0=\mathbb R^2\setminus\overline Q,\qquad
 \partial^eQ=\bigl([0,\infty)\times\{0\}\bigr)
              \cup\bigl(\{0\}\times[0,\infty)\bigr).
\]
正半轴上的非原点处密度为 \(1/2\)，原点处密度为 \(1/4\)，所以它们都属于测度论边界。两条开正半轴上的外法向分别为
\[
 \nu_Q(x_1,0)=-e_2\quad(x_1>0),qquad
 \nu_Q(0,x_2)=-e_1\quad(x_2>0).
\]
在原点，两个边界射线对 \(D\mathbf1_Q\) 的贡献分别沿 \(e_2\) 与 \(e_1\)，故
\[
 \frac{D\mathbf1_Q(B(0,r))}{|D\mathbf1_Q|(B(0,r))}
 =\frac{e_1+e_2}{2},
\]
其长度为 \(1/\sqrt2\)，不是 \(1\)。因此
\[
 \partial^*Q=\bigl((0,\infty)\times\{0\}\bigr)
              \cup\bigl(\{0\}\times(0,\infty)\bigr),
\]
角点属于 \(\partial^eQ\setminus\partial^*Q\)。

集合 \(F\) 与单位球只相差可数零测集，所以它们的密度点、分布导数、周长测度和约化边界完全相同：
\[
 F^1=B(0,1),\qquad F^0=\mathbb R^d\setminus\overline{B(0,1)},\qquad
 \partial^eF=\partial^*F=\partial B(0,1),\qquad
 \nu_F(x)=x.
\]
另一方面，\(F\) 与其补集都在单位球内稠密，故 \(F\) 的内部为空、闭包为 \(\overline{B(0,1)}\)，从而它的拓扑边界是整个闭球。这说明拓扑边界不具有零测修改不变性。

最后，对 \(x\in\partial B(0,1)\)，有
\[
 z\in B(0,1)_{x,r}
 \quad\Longleftrightarrow\quad
 z\cdot x<-\frac r2|z|^2.
\]
在每个固定球内，右侧阈值一致趋于零，边界超平面又是零测集，故示性函数在局部 \(L^1\) 中趋于 \(\mathbf1_{\{z\cdot x<0\}}\)。旋转到 \(x=e_d\) 后，靠近原点的缩放球面写成图
\[
 z_d=\frac{-1+\sqrt{1-r^2|z'|^2}}{r},
\]
它及其图面积因子在紧集上一致趋于 \(0\) 与 \(1\)；球面的另一支在固定紧集之外。由图面积公式或直接参数换元，缩放周长测度趋于
\(\mathcal H^{d-1}\llcorner x^\perp\)。由于 \(F\) 与球有相同的示性函数等价类及周长测度，同样的两种缩放结论也适用于 \(F\)。
\end{solution}

\begin{optionalexercise}\label{b1:ex:27-translation-bv-norm-characterization}
设 \(u\in BV(\mathbb R^d)\)，\(\tau_hu(x)=u(x-h)\)。
\begin{enumerate}
 \item\label{b1:item:27-translation-strict}
 证明当 \(h\to0\) 时，\(\tau_hu\) 总是严格收敛到 \(u\)。
 \item\label{b1:item:27-translation-norm}
 证明
 \[
  \|\tau_hu-u\|_{BV(\mathbb R^d)}\longrightarrow0
    \quad\Longleftrightarrow\quad u\in W^{1,1}(\mathbb R^d).
 \]
 \item\label{b1:item:27-translation-interval}
 对 \(u=\mathbf1_{(0,1)}\)，计算 \(0<h<1\) 时的
 \(\|\tau_hu-u\|_{BV(\mathbb R)}\)。
\end{enumerate}
\end{optionalexercise}
% 来源：自拟；以测度在全变差范数中的平移连续性辨认绝对连续导数，
% 不重复正文对阶跃函数作光滑严格逼近的例子。
\begin{solution}
函数的 \(L^1\) 平移连续，而
\[
 D(\tau_hu)=\tau_hDu,\quad
 |D(\tau_hu)|(\mathbb R^d)=|Du|(\mathbb R^d),
\]
其中 \((\tau_h\mu)(A)=\mu(A-h)\)。
这证明第一项。若 \(u\in W^{1,1}\)，则函数及各弱偏导数都在 \(L^1\) 中平移连续，
从而得到第二项的反向蕴涵。

现假设平移在 \(BV\) 范数中连续，记 \(\mu=Du\)。
于是 \(|\tau_h\mu-\mu|(\mathbb R^d)\to0\)。
用非负单位积分核磨光测度，记
\(\mu_\varepsilon=(\rho_\varepsilon*\mu)m_d\)。
对每个连续紧支集向量场 \(\Phi\)，Fubini 定理给出
\[
 \int\Phi\cdot\mathrm d(\mu_\varepsilon-\mu)
 =\int\rho_\varepsilon(h)
       \left(\int\Phi\cdot\mathrm d(\tau_h\mu-\mu)\right)\,\mathrm dh.
\]
向量测度全变差的连续测试对偶式于是给出
\[
 |\mu_\varepsilon-\mu|(\mathbb R^d)
 \le\sup_{|h|\le\varepsilon}
          |\tau_h\mu-\mu|(\mathbb R^d)\longrightarrow0.
\]
每个 \(\mu_\varepsilon\) 都对 Lebesgue 测度绝对连续。
若 \(N\) 为 Lebesgue 零测 Borel 集，则
\[
 |\mu|(N)\le|\mu-\mu_\varepsilon|(\mathbb R^d)
                     +|\mu_\varepsilon|(N)\longrightarrow0.
\]
所以 \(\mu\ll m_d\)。各分量由 Radon–Nikodym 定理写为 \(L^1\) 密度，
这些密度正是 \(u\) 的弱偏导数，因而 \(u\in W^{1,1}\)。
此处连续测试对偶式由极分解和连续紧支集函数的积分逼近得到，与变差对偶证明中的步骤相同。

最后，当 \(u=\mathbf1_{(0,1)}\) 时，
\[
 \|\tau_hu-u\|_1=2h,\quad
 D(\tau_hu-u)=\delta_h-\delta_{1+h}-\delta_0+\delta_1.
\]
四个点互不相同，故导数差的全变差为 \(4\)，于是
\[
 \|\tau_hu-u\|_{BV(\mathbb R)}=2h+4.
\]
函数与自身的小平移严格接近，却没有在 \(BV\) 范数中接近。
\end{solution}

\begin{optionalexercise}\label{b1:ex:27-finite-valued-strict-approximation}
设 \(m_d(\Omega)<\infty\)，\(u\in BV(\Omega)\)，且 \(0\le u\le M<\infty\) 几乎处处。
证明存在只取有限多个值的 \(v_j\in BV(\Omega)\)，使
\[
 \|v_j-u\|_{L^\infty(\Omega)}\longrightarrow0,\quad
 |Dv_j|(\Omega)\longrightarrow|Du|(\Omega).
\]
这里 \(L^\infty\) 指本性上确界范数。要求将每个 \(v_j\) 写成有限周长超水平集的示性函数之和，
并说明层值怎样选择。
\end{optionalexercise}
% 来源：自拟；在一组平移的等距层值中平均周长，选出变差受控的量化函数。
\begin{solution}
固定 \(h>0\)，对 \(0<\theta<h\) 定义
\[
 v_{h,\theta}
   =h\sum_{k=0}^\infty\mathbf1_{\{u>\theta+kh\}}.
\]
因为 \(u\le M\)，非空层只有有限个，所得函数取有限多个值。
对每个 \(s\in[0,M]\)，等距网格的计数误差小于或等于一个网格长度，
所以
\[
 |v_{h,\theta}-u|\le h\quad\text{几乎处处}.
\]
余面积公式说明，对几乎每个 \(\theta\)，上式涉及的全部层集都有有限周长；
只需排除每个平移网格命中坏层值的零测集，再取可数并。
在这些 \(\theta\) 上，
\[
 |Dv_{h,\theta}|(\Omega)
 \le S_h(\theta)=
        h\sum_{k=0}^\infty P(\{u>\theta+kh\},\Omega).
\]
由 Tonelli 定理及 \(u\ge0\)，
\[
 \frac1h\int_0^hS_h(\theta)\,\mathrm d\theta
 =\int_0^\infty P(\{u>t\},\Omega)\,\mathrm dt
 =|Du|(\Omega).
\]
因此可选一个容许的 \(\theta_h\)，使
\(S_h(\theta_h)\le|Du|(\Omega)+h\)。
取 \(h_j\downarrow0\)，令 \(v_j=v_{h_j,\theta_{h_j}}\)。
有限体积保证
\[
 \|v_j-u\|_1\le h_jm_d(\Omega)\longrightarrow0,
 \quad
 \limsup_j|Dv_j|(\Omega)\le|Du|(\Omega).
\]
下半连续性给出相反的下极限不等式，因而总变差收敛。
有限值逼近与光滑逼近采用不同的对象：前者以有限周长层集承载界面，
后者以光滑梯度承载变化；二者都只要求总变差数值收敛。
\end{solution}

\begin{exercise}\label{b1:ex:27-irregular-zero-extension}
令
\[
 a_j=2^{-j},\quad b_j=2^{-j}+2^{-j-2},\quad
 \Omega=\bigcup_{j=1}^\infty(a_j,b_j)\subset(0,1).
\]
在 \(\Omega\) 上取 \(u=1\)，并令 \(u_N\) 在前 \(N\) 个区间上等于 \(1\)，
在其余区间上等于 \(0\)。
证明 \(u_N\to u\) 于 \(BV(\Omega)\)，每个全空间零延拓 \(E_0u_N\) 都属于
\(BV(\mathbb R)\)，但 \(E_0u\notin BV(\mathbb R)\)。
计算 \(|D(E_0u_N)|(\mathbb R)\)。
\end{exercise}
% 来源：自拟；域内没有导数的常数分片，经补零后产生无穷多个端点代价。
\begin{solution}
这些区间两两分离，且长度为 \(2^{-j-2}\)。
函数 \(u_N\) 和 \(u\) 在每个连通分量上为常数，所以它们在 \(\Omega\) 中的分布导数均为零。
因此
\[
 \|u_N-u\|_{BV(\Omega)}
 =\sum_{j>N}(b_j-a_j)\longrightarrow0.
\]
另一方面，在全空间中逐区间计算，得到
\[
 D(E_0u_N)=\sum_{j=1}^N(\delta_{a_j}-\delta_{b_j}),
 \quad
 |D(E_0u_N)|(\mathbb R)=2N.
\]
所有端点互不相同，因此没有相消。

若 \(E_0u\) 有有限变差，则在每个端点的小邻域内，它的导数都必须与上述单个区间的导数一致。
每次固定有限多个端点，可以选取两两不交的小邻域并避开其余端点。
故对每个 \(N\)，其总变差至少为 \(2N\)，矛盾。
这个例子不否认内部紧支集截断的零延拓性质：\(u\) 的非零部分一直延伸到
无穷多个分量端点，没有同一个内部紧集隔开这些边界。
\end{solution}

\begin{exercise}\label{b1:ex:27-eccentric-hole-flux}
在 \(\mathbb R^2\) 中设 \(R,r>0\)、\(|a|+r<R\)，并令
\[
 E=B(0,R)\setminus\overline{B(a,r)}.
\]
求 \(E\) 的外法向、\(D\mathbf1_E\) 及 \(P(E)\)。
再对向量场
\[
 F(x)=\frac{x-a}{|x-a|^2}
\]
计算两条边界曲线各自的外向通量，解释为什么总通量为零。
不要求使用绕数或复分析。
\end{exercise}
% 来源：自拟；用偏心孔检验内边界方向，以光滑域Gauss计算奇点被排除后的通量。
\begin{solution}
外圆上的外法向为 \(x/R\)，内圆上的外法向指向孔内，即
\[
 \nu_E(x)=
 \begin{cases}
 x/R,&x\in\partial B(0,R),\\
 -(x-a)/r,&x\in\partial B(a,r).
 \end{cases}
\]
由光滑域周长公式，
\[
 D\mathbf1_E
 =-\frac{x}{R}\,\mathcal H^1\!\llcorner\partial B(0,R)
   +\frac{x-a}{r}\,\mathcal H^1\!\llcorner\partial B(a,r),
 \quad P(E)=2\pi(R+r).
\]
两条边界的周长相加，虽然它们的外法向相对于各自圆心取了不同的方向。

场 \(F\) 在 \(\overline E\) 的邻域内光滑，直接求导得
\[
 \operatorname{div}F(x)
 =\frac2{|x-a|^2}-\frac{2|x-a|^2}{|x-a|^4}=0.
\]
如需紧支集测试场，可乘以在 \(\overline E\) 的邻域等于 \(1\)、支集避开 \(a\) 的光滑截断。
在内圆上，\(F\cdot\nu_E=-1/r\)，故内圆通量为 \(-2\pi\)。
Gauss 公式给出两条曲线的通量之和等于 \(\int_E\operatorname{div}F=0\)，
所以外圆通量为 \(2\pi\)。
奇点位于被删去的孔中，不在积分区域内；把内孔也按远离其圆心的方向取法向，
会错误地丢失这个抵消。
\end{solution}

\begin{exercise}\label{b1:ex:27-signed-measure-primitive}
设 \(-\infty<a<b<\infty\)，\(\mu\) 是 \((a,b)\) 上的有限 Radon 符号测度，
并令
\[
 F(x)=\mu\bigl((a,x]\bigr),\quad a<x<b,\quad M=\mu\bigl((a,b)\bigr).
\]
\begin{enumerate}
 \item\label{b1:item:27-measure-primitive}
 证明 \(F\in BV(a,b)\)、\(DF=\mu\)，并证明所有分布导数为 \(\mu\) 的
 有界变差函数都与某个 \(F+c\) 几乎处处相等。
 \item\label{b1:item:27-measure-primitive-endpoints}
 求 \(F\) 的左右极限及跳跃值。将 \(F+c\) 在 \((a,b)\) 外补零，记为 \(G_c\)，
 计算 \(DG_c\) 和 \(|DG_c|(\mathbb R)\)。
 \item\label{b1:item:27-measure-primitive-optimal-constant}
 在全部常数 \(c\in\mathbb R\) 中，求上述零延拓总变差的最小值。
\end{enumerate}
\end{exercise}
% 来源：自拟；一般符号测度原函数及端点代价，不复述正文的Cantor概率测度例子。
\begin{solution}
有 \(|F(x)|\le|\mu|\bigl((a,b)\bigr)\)，故 \(F\in L^1(a,b)\)。
对 \(\varphi\in C_c^\infty(a,b)\)，Fubini 定理给出
\[
 \int_a^bF(x)\varphi'(x)\,\mathrm dx
 =\int_{(a,b)}\left(\int_t^b\varphi'(x)\,\mathrm dx\right)\,\mathrm d\mu(t)
 =-\int_{(a,b)}\varphi\,\mathrm d\mu.
\]
换序由有限变差及 \((b-a)\|\varphi'\|_\infty\) 的控制保证。
因此 \(DF=\mu\)，且 \(|DF|\bigl((a,b)\bigr)=|\mu|\bigl((a,b)\bigr)\)。

若 \(Dv=\mu\)，令 \(w=v-F\)，则 \(Dw=0\)。
任何积分为零的 \(\psi\in C_c^\infty(a,b)\) 都可写成某个
\(\phi\in C_c^\infty(a,b)\) 的导数，故 \(\int w\psi=0\)。
固定 \(\eta\in C_c^\infty(a,b)\)、\(\int\eta=1\)，令 \(c=\int w\eta\)。
将任意测试函数减去其积分乘以 \(\eta\)，可得
\(\int w\psi=c\int\psi\)，于是正则分布的唯一性给出 \(w=c\) 几乎处处。

有限测度的集合连续性给出
\[
 F(a+)=0,\quad F(b-)=M,\quad
 F(x+)=F(x),\quad F(x-)=\mu\bigl((a,x)\bigr).
\]
所以内部跳跃值为 \(F(x)-F(x-)=\mu(\{x\})\)。
对全空间测试函数重复上面的换序，但保留端点项，得到
\[
 DG_c=\mu+c\delta_a-(c+M)\delta_b.
\]
其中 \(\mu\) 在区间外补零。内部测度与两个端点质量相互奇异，故
\[
 |DG_c|(\mathbb R)=|\mu|\bigl((a,b)\bigr)+|c|+|c+M|.
\]
由三角不等式，\(|c|+|c+M|\ge|M|\)，
并且当且仅当 \(c\) 位于连接 \(0\) 与 \(-M\) 的闭区间内时取等号。
故最小值为
\[
 |\mu|\bigl((a,b)\bigr)+|M|.
\]
内部总变差由导数固定，常数的选择只改变两个端点之间分配的补零代价。
\end{solution}

\begin{optionalexercise}\label{b1:ex:27-confined-bv-minimization}
设 \(d\ge2\)，考虑
\[
 \mathcal A=\left\{\,u\in BV(\mathbb R^d):
      u\ge0\ \text{几乎处处},\quad\int_{\mathbb R^d}u\,\mathrm dx=1\,\right\},
\]
以及允许取值 \(+\infty\) 的泛函
\[
 J(u)=|Du|(\mathbb R^d)+\int_{\mathbb R^d}|x|u(x)\,\mathrm dx.
\]
证明 \(J\) 在 \(\mathcal A\) 上取得有限的极小值。
再证明，若删除最后的一阶矩项，只最小化总变差，则下确界为零但不能取得。
指出矩项在极限过程中保持单位质量的具体作用。
\end{optionalexercise}
% 来源：自拟；全空间BV直接法由矩控制防止质量逃逸，不使用最佳等周常数或解的正则性。
\begin{solution}
取非负光滑紧支集函数并归一化积分，得到 \(J<\infty\) 的容许函数。
故下确界 \(I\) 有限且非负。选 \(u_j\in\mathcal A\)，使 \(J(u_j)\to I\)，
并可要求 \(J(u_j)\le I+1\)。
于是 \(\|u_j\|_1=1\)、总变差一致有界，且
\[
 \int_{|x|>R}u_j(x)\,\mathrm dx
 \le R^{-1}\int|x|u_j(x)\,\mathrm dx\le\frac{I+1}{R}.
\]
由带尾部控制的有界变差函数紧性%
\cref{b1:cor:bv-global-compactness-with-tails}，可取全空间 \(L^1\) 收敛子列，
记极限为 \(u\)。再抽几乎处处收敛子列，有 \(u\ge0\)，而全域 \(L^1\) 收敛保证
\(\int u=1\)，并非只得到 \(\int u\le1\)。
变差下半连续性和 Fatou 引理分别给出
\[
 |Du|(\mathbb R^d)\le\liminf_j|Du_j|(\mathbb R^d),\quad
 \int|x|u(x)\,\mathrm dx\le\liminf_j\int|x|u_j(x)\,\mathrm dx.
\]
所以 \(u\in\mathcal A\)，且 \(J(u)\le\liminf_jJ(u_j)=I\)，即取得极小值。

删除矩项后，取固定的非负 \(\eta\in C_c^\infty(\mathbb R^d)\)、\(\int\eta=1\)，令
\[
 v_R(x)=R^{-d}\eta(x/R).
\]
则 \(v_R\in\mathcal A\)，且
\[
 |Dv_R|(\mathbb R^d)=R^{-1}\|\nabla\eta\|_1\longrightarrow0.
\]
因此总变差的下确界为零。如果某个 \(v\in\mathcal A\) 达到它，
由\hyperref[b1:thm:bv-sobolev-endpoint]{有界变差函数的 Sobolev 端点估计\ref*{b1:thm:bv-sobolev-endpoint}\CJKecglue}便有 \(\|v\|_{d/(d-1)}=0\)，
与 \(\int v=1\) 矛盾。
膨胀序列在每个有界区域内的质量趋零，但总质量始终为一；
矩项所给的统一尾部估计，恰好排除了这种单位质量在局部极限中丢失的情形。
\end{solution}

% !TeX root = ../../Book1.tex
\OptionalChapter{BV 函数的精细结构}{b1:ch:28}
\begin{chapterabstract}
本章把 BV 导数分成体积上的绝对连续部分、界面上的跳跃部分与剩余的 Cantor 部分。
通过近似上下极限、单侧值和跳跃法向，把等价类的点值恢复与测度分解联系起来。
最后介绍 SBV 空间，并在明确的值域与增长条件下推导 Mumford–Shah 型能量的极小元存在性。
精细迹、弥散结构及链律的深层输入单独标明，不以概览代替尚未展开的理论。
\end{chapterabstract}

\begin{learninggoals}
\begin{enumerate}
\item 用近似上下极限识别有限近似极限点，区分指定代表的近似连续与等价类自身的点态信息。
\item 确定跳跃的两侧值及法向，证明其唯一性与可测性，并用好层值解释跳跃集的可整流性。
\item 区分 \(D^au\)、\(D^ju\)、\(D^cu\)，辨认近似梯度、跳跃高度与奇异连续变化各自的作用。
\item 理解 SBV 为何排除 Cantor 部分，以及严格 BV 收敛为何仍不保证 SBV 闭合。
\item 在准确列出的深层输入下完成紧性、表面面积下半连续性和受限变分问题的存在性推导。
\end{enumerate}
\end{learninggoals}

上一章的余面积公式把一般函数还原为一族有限周长集合，
但这些层集的边界并不总是叠在同一张曲面上。
光滑函数的层面累积成体积中的梯度，阶跃函数的层面重叠成一道跳跃界面，
Cantor 函数则给出第三种可能：变化集中在零体积集合上，却没有任何非零跳跃。
导数测度的三部分分解正是对这三个现象的统一描述。

本章为选读，也是一卷实分析通向变分问题的收束。
近似极限的基础已在第十四章建立，这里增加上下极限与等价类代表的组织，
不重新改变“近似连续”对指定点值的含义。
读者应熟悉第六章的测度分解、第十九章的可整流性，以及上一章的 BV 逼近与余面积。
若暂时只想掌握可使用的结论，可以先看每个结构定理的准确陈述及一维例子，
再沿明确写出的证明概要追踪尚需学习的精细迹与切片工具。

Alberti 与 Mantegazza 的《A Note on the Theory of SBV Functions》%
\cite{AlbertiMantegazza1997SBV}提供了本章结构陈述与紧性路线的集中对照。
该文对基础精细结构另引专著，不能因为篇幅短就以为所有前置都已在几页内证明。
本书把能由已有工具推出的可整流性、测度分解后果、固定面积界下的闭合桥、
跳跃面积下半连续性和极小值论证写出；余下的深层输入明确标记。
进一步的系统阅读可选 Evans 与 Gariepy 的%
《Measure Theory and Fine Properties of Functions》\cite{Evans2015Measure}，
以及 Ambrosio、Fusco 与 Pallara 的%
《Functions of Bounded Variation and Free Discontinuity Problems》%
\cite{Ambrosio2000Functions}。

% !TeX root = ../../Book1.tex
\section{近似极限与近似极限代表}
\label{b1:sec:2801}

上一章把总变差解释为各层周长的累积。本章进一步问：这些周长在同一点附近怎样叠加？
有些点的函数值沿各个方向趋于同一个数，有些点则沿一个界面的两侧趋于不同的数。
若只查看任意选定代表的普通不连续点，零测修改就会掩盖这种区别。
因此先从邻近球中的体积比例定义点值，暂不使用导数测度。
本节的基础定义适用于一般实值可测函数；涉及有界变差函数的余维一结构时，再明确加入有限变差条件。

Alberti 与 Mantegazza 的《A Note on the Theory of SBV Functions》%
\cite[Remark 1.2]{AlbertiMantegazza1997SBV}把近似极限与导数分解放在一起介绍。
其文中把没有近似极限的集合记为 \(S_u\)，并在这个集合上给代表补零；
本章保留 \(S_u\) 表示近似不连续集，另以 \(J_u\) 专指有两个不同单侧值的跳跃点。
两者的关系是后面的结构定理，不能在定义时直接把它们当成同一个集合。
系统学习可继续参阅 Evans 与 Gariepy 的《Measure Theory and Fine Properties of Functions》%
\cite{Evans2015Measure}及 Ambrosio、Fusco 与 Pallara 的%
《Functions of Bounded Variation and Free Discontinuity Problems》%
\cite{Ambrosio2000Functions}。

% 实读AlbertiMantegazza1997SBV作者公开正式副本，刊印376--377页注1.2；
% 作者副本内页2--3。本文前节基础密度刻画独立证明，不冒称这些页含全证。
% 原文S_u与本书J_u的区别在正文说明；下节引述精细结构时另列证明义务。

\subsection{近似上、下极限}
\label{b1:subsec:280101}

固定开集 \(\Omega\subseteq\mathbb R^d\)，令 \(u:\Omega\to\mathbb R\) 可测。
只考察中心 \(x\in\Omega\) 且足够小的球。对可测集 \(A\subseteq\Omega\)，简记
\[
 \theta_r(A,x)=\frac{m_d(A\cap B(x,r))}{\omega_dr^d}.
\]
说一个性质在 \(x\) 附近以密度一成立，是说不满足它的点集之 \(\theta_r\) 趋于零。
这里的“密度一”不是“某个邻域中处处成立”，也不是“除掉一个固定的零测集后成立”。

\begin{definition}\label{b1:def:approximate-upper-lower-limits}
定义 \(u\) 在 \(x\) 的\term[jinsi-xia-jixian]{近似下极限}[approximate lower limit]与%
\term[jinsi-shang-jixian]{近似上极限}[approximate upper limit]为
\[
 u^\wedge(x)=
 \sup\{\,t\in\mathbb R:\lim_{r\downarrow0}\theta_r(\{u<t\},x)=0\,\},
\]
\[
 u^\vee(x)=
 \inf\{\,t\in\mathbb R:\lim_{r\downarrow0}\theta_r(\{u>t\},x)=0\,\}.
\]
允许取扩展实数值，约定 \(\sup\varnothing=-\infty\)、\(\inf\varnothing=+\infty\)。
\end{definition}
\symidx{\(u^\wedge\)：近似下极限}
\symidx{\(u^\vee\)：近似上极限}

下极限的容许阈值构成向下封闭的区间，上极限的容许阈值构成向上封闭的区间；
端点是否属于该区间不影响定义。不能把这里的体积密度换成小球中的本性上、下确界。
前者允许例外部分的比例逐渐趋零，后者只忽略严格零测的部分。

\begin{proposition}\label{b1:prop:approximate-limits-basic}
函数 \(u^\wedge,u^\vee\) 是扩展实值 Borel 函数，且
\[
 u^\wedge(x)\le u^\vee(x)\quad(x\in\Omega).
\]
它们只依赖 \(u\) 的几乎处处等价类。若 \(u\le v\) 几乎处处，则在每一点都有
\[
 u^\wedge\le v^\wedge,\quad u^\vee\le v^\vee.
\]
此外，\((-u)^\wedge=-u^\vee\)、\((-u)^\vee=-u^\wedge\)。
\end{proposition}
\begin{proof}
由\cref{b1:thm:lebesgue-borel-representative}，先把 \(u\) 换成几乎处处相等的 Borel 代表。
任意这样的替换都只在零测集内改变每个阈值集，
故不改变任何球中的体积，特别不改变任一点的两个近似极限。

对固定阈值 \(t\) 与半径 \(r\)，球体积积分作为中心的函数是 Borel 的。
对固定中心，它关于正半径连续，因为球壳的体积趋于零。
因此“\(\theta_r(\{u<t\},x)\to0\)”可以只用正有理半径写成可数个 Borel 条件。
例如对充分小的半径，使用
\(\inf_{n}\sup_{r\in\mathbb Q,\ 0<r<1/n}\theta_r\)；
球未必留在 \(\Omega\) 时可先在域外把阈值集补空，极限不受影响。
阈值的单调性又使定义中的上、下确界可只取有理 \(t\)。
所以两个近似极限都是 Borel 函数。

若 \(u^\wedge(x)>u^\vee(x)\)，可在二者之间取 \(a<b\)，使
\(\{u<b\}\) 与 \(\{u>a\}\) 在 \(x\) 都有密度零。
但这两个集合的并为 \(\Omega\)，不可能在内部一点有密度零，矛盾。
这个论证同样涵盖端点为无穷的情况。

若 \(u\le v\) 几乎处处，则
\(\{v<t\}\subseteq\{u<t\}\)、\(\{u>t\}\subseteq\{v>t\}\)，均按零测集忽略理解。
比较容许阈值区间即可得到两个单调关系。
最后，把 \(\{-u<t\}\) 写成 \(\{u>-t\}\)，便得到取负号的公式。
\end{proof}

\subsection{近似连续点}
\label{b1:subsec:280102}

\begin{definition}\label{b1:def:bv-finite-approximate-limit}
若存在有限实数 \(\ell\)，使对每个 \(\varepsilon>0\)，都有
\begin{equation}\label{b1:eq:approximate-limit-density}
 \lim_{r\downarrow0}
 \frac{m_d\bigl(\{y\in B(x,r):|u(y)-\ell|>\varepsilon\}\bigr)}
 {\omega_dr^d}=0,
\end{equation}
则按\cref{b1:def:approximate-limit}，\(u\) 在 \(x\) 有近似极限 \(\ell\)。
记该点集为 \(G_u\)，并在其上记 \(\widetilde u(x)=\ell\)。
\end{definition}

这里并不要求原代表已经满足 \(u(x)=\ell\)。
\chapref{b1:ch:14}的“近似连续”还要求这个等式成立，属于指定代表的性质；
本章的 \(G_u\) 只记录等价类有有限近似极限的点。
在 \(G_u\) 上重定义为 \(\widetilde u\) 后，该代表才在这些点近似连续。
后文为简便说等价类的近似连续点时，均指这项重定义后的性质，
不更改\chapref{b1:ch:14}对于指定函数值的定义。

\begin{proposition}\label{b1:prop:approximate-continuity-characterization}
对 \(x\in\Omega\) 及有限实数 \(\ell\)，下列条件等价：
\begin{equivalences}
 \item\label{b1:item:approximate-limit-density}
 \(u\) 在 \(x\) 的近似极限为 \(\ell\)。
 \item\label{b1:item:approximate-limits-equal}
 \(u^\wedge(x)=u^\vee(x)=\ell\)。
\end{equivalences}
近似极限若存在便唯一。若 \(u\in L^1_{\mathrm{loc}}(\Omega)\)，则几乎每点属于 \(G_u\)，
且 \(\widetilde u=u\) 几乎处处。
\end{proposition}
\begin{proof}
若密度条件成立，则对 \(a<\ell<b\)，集合 \(\{u<a\}\)、\(\{u>b\}\) 都有密度零。
故 \(u^\wedge\ge\ell\)、\(u^\vee\le\ell\)，结合前一命题得相等。
反过来，若两极限都等于 \(\ell\)，给定 \(\varepsilon>0\)，
可在 \(\ell-\varepsilon\) 与 \(\ell\) 之间取一个下极限的容许阈值，
在 \(\ell\) 与 \(\ell+\varepsilon\) 之间取一个上极限的容许阈值。
两项密度零条件的并包含 \(\{|u-\ell|>\varepsilon\}\)，所以该集合也有密度零。
唯一性随两极限的数值唯一性得到。

由\chapref{b1:ch:14}的 Lebesgue 微分定理，局部可积函数的几乎每点满足
\[
 \lim_{r\downarrow0}\frac1{\omega_dr^d}\int_{B(x,r)}|u(y)-u(x)|\,\mathrm dy=0.
\]
Chebyshev 不等式将左侧平均振荡控制转成%
\eqref{b1:eq:approximate-limit-density}，从而得到最后一项。
\end{proof}

Lebesgue 点的平均振荡趋零，比一般的密度意义近似极限强。
若 \(u\) 在 \(x\) 的某个邻域本性有界，二者确实等价：
密度意义收敛给出
\[
 \frac1{\omega_dr^d}\int_{B(x,r)}|u-\ell|
 \le\varepsilon+M\,\theta_r(\{|u-\ell|>\varepsilon\},x),
\]
其中 \(M\) 是邻域内 \(|u-\ell|\) 的本性上界。先令 \(r\downarrow0\)，再令
\(\varepsilon\downarrow0\) 即可。无界函数则可能在密度很小的部分上留下较大的积分。

\secref{b1:sec:1405}已构造可积尖峰的反例：非零部分在原点的密度趋零，
而球中归一化积分不趋零。因此这里不重复该构造，也不把近似极限存在自动当作一个
\(L^1\) 平均极限。有界变差函数在余维一尺度上还具有更强的性质，证明这点需要下一节的结构输入。

\subsection{不连续点集}
\label{b1:subsec:280103}

\begin{definition}\label{b1:def:bv-approximate-discontinuity-set}
定义 \(u\) 的\term[jinsi-bulianxu-ji]{近似不连续集}[approximate discontinuity set]为
\[
 S_u=\Omega\setminus G_u.
\]
\end{definition}
\symidx{\(S_u\)：没有有限近似极限的点集}
\symidx{\(\widetilde u\)：近似极限代表}

由前面的刻画，\(S_u\) 为 Borel 集。它包括 \(u^\wedge<u^\vee\) 的点，
也包括两极限都等于同一个无穷值的点；后者不能因上下极限相等而被算作有限值近似连续点。
在 \(S_u\) 上暂给 \(\widetilde u\) 补零，便得到一个全域 Borel 函数；
在尚未知道它与 \(u\) 几乎处处相等时，只把这项补值看作 Borel 规范化。
若 \(u\) 局部可积，前一命题进一步说明二者几乎处处相等；
此时才称 \(\widetilde u\) 为
\term[jinsi-jixian-daibiao]{近似极限代表}[approximate-limit representative]。
这个补零约定不声称零值具有特殊的边界意义；在跳跃点上，稍后还会定义两侧值的平均。

对任何 \(L^1_{\mathrm{loc}}\) 函数，目前只能由微分定理断言 \(m_d(S_u)=0\)。
对于有界变差函数，下一节将把它加强为一个余维一的结构结论。
加强的内容不是把 \(S_u\) 全部删掉，而是说明除了一个
\(\mathcal H^{d-1}\)-零集，近似不连续表现为沿界面的两侧跳跃。

例如 \(u=\mathbf1_{\{x_d>0\}}\) 在平面外近似连续，在平面上有
\(u^\wedge=0\)、\(u^\vee=1\)，因此 \(S_u=\{x_d=0\}\)。
它是局部有界变差函数，却在一个具有正余维一测度的集合上没有单一近似极限。
反之，\(\mathbf1_{\mathbb Q^d}\) 的任意点均有近似极限零，
所以 \(S_u=\varnothing\)，尽管这个指定代表在通常意义下处处不连续。
两个例子分别说明，跳跃不能通过给界面点重赋单一值消除，
而零测修改造成的普通不连续根本不应计入 \(S_u\)。

无穷近似值即使对有界变差函数也可能发生。设 \(d\ge2\)、\(0<\alpha<d-1\)，
在 \(B(0,1)\) 上令 \(u(x)=|x|^{-\alpha}\)，原点任意补值。
球壳积分给出 \(u,\,\alpha|x|^{-\alpha-1}\in L^1\)。
在删去半径 \(\varepsilon\) 小球的区域上分部积分，内边界项的绝对值至多为
\[
 C\varepsilon^{d-1-\alpha}\|\varphi\|_\infty\longrightarrow0.
\]
因此经典梯度在原点外的表达式也是全域弱梯度，\(u\in W^{1,1}\subset BV\)。
但原点的两个近似极限均为 \(+\infty\)，故 \(0\in S_u\)，而这不是两个有限值之间的跳跃。
单点的 \(\mathcal H^{d-1}\) 测度为零，所以这个现象与下一节的结构定理相容。

\begin{proposition}[近似振荡与层集边界]\label{b1:prop:approximate-oscillation-level-boundary}
设 \(T\subset\mathbb R\) 为任意可数稠密集。则
\[
 \{x:u^\wedge(x)<u^\vee(x)\}
 \subseteq\bigcup_{t\in T}\partial^e\{u>t\}.
\]
更具体地，若 \(u^\wedge(x)<t<u^\vee(x)\)，则 \(x\) 属于该层集的测度论边界；
若 \(t<u^\wedge(x)\)，层集在 \(x\) 的密度为 \(1\)；
若 \(t>u^\vee(x)\)，密度为 \(0\)。
\end{proposition}
\begin{proof}
若 \(t<u^\wedge(x)\)，在 \(t\) 与下极限之间选容许阈值 \(a\)，
则 \(\{u\le t\}\subseteq\{u<a\}\) 的密度为零，故 \(\{u>t\}\) 密度为一。
上极限情形直接由阈值单调性得到。

若 \(u^\wedge(x)<t<u^\vee(x)\)，而 \(\{u>t\}\) 密度为零，
则 \(t\) 是上极限的容许阈值，与 \(t<u^\vee(x)\) 矛盾。
若它密度为一，则 \(\{u<t\}\) 密度为零，推出 \(u^\wedge(x)\ge t\)，同样矛盾。
因而该点既不在测度论内部，也不在外部，属于 \(\partial^e\{u>t\}\)。
最后，在每段非空的扩展实数开区间 \((u^\wedge(x),u^\vee(x))\) 中取一个 \(T\) 的点，
就得到可数并包含。
\end{proof}

这条联系允许先选出一组周长有限的稠密好层值，再利用上一章的边界结构。
它还没有处理两极限同为无穷的点，也没有保证同一点上的不同层面法向一致；
后两项正是不能由这个简单包含替代的精细性质。

后续所有界面与代表仍由几乎处处类决定。特别是，
“导数由函数的精确代表给出”必须同时说明用了哪一种代表、
在哪一个测度的几乎处处意义下成立；Lebesgue 几乎处处相等的代表，
若被直接放进奇异导数测度的积分，结果未必相同。

% !TeX root = ../../Book1.tex
\section{跳跃集}
\label{b1:sec:2802}

近似不连续只说明一个球中不能用单一数值描述函数。
跳跃则有更确定的形状：把球沿一个平面分成两半后，两侧分别接近两个不同的有限值。
本节先把这个局部图像写成定义，确定两侧值和方向；
再用上一章的层集与周长结构，说明有界变差函数的跳跃点为何由可数个曲面承载。

以下均取实值函数。定义和最初几个命题只要求 \(u\in L^1_{\mathrm{loc}}(\Omega)\)；
涉及可整流性与导数测度时，另加 \(u\in BV(\Omega)\)。
文献有时允许任意选择法向并同时交换两侧值，本书则把较大值记作 \(u^+\)，固定取向。
这里的上标 \(u^\pm\) 专用于跳跃迹；一般函数的正部、负部已统一记为
\(u_+\)、\(u_-\)，二者不得混用。
这一选择不改变最终的向量测度。

% 实际阅读：AlbertiMantegazza1997SBV，作者公开正式副本，注1.2，
% 作者版内页2--3（刊印376--377页）。其中给出广义跳跃集上的两侧L1迹与导数公式，
% 并将精细结构证明指向Evans--Gariepy 1992第5.9节；未声称已实读该书证明。
% 下文唯一性、上下值关系、Borel性及Ju可整流性在本书工具内独立证明；
% 非有限近似值的余维一零测性、BV曲面两侧迹与界面导数公式明确作为深层输入。

\subsection{跳跃点}
\label{b1:subsec:280201}

对 \(x\in\Omega\)、单位向量 \(\nu\) 及足够小的 \(r>0\)，记
\[
 B_r^\pm(x,\nu)
 =\{y\in B(x,r):\ \pm(y-x)\cdot\nu>0\}.
\]
两半球的体积均为 \(\omega_dr^d/2\)；中间平面为体积零集，不影响平均。

\begin{definition}\label{b1:def:approximate-jump-point}
给定 \(u\in L^1_{\mathrm{loc}}(\Omega)\) 与 \(x\in\Omega\)。
若存在有限实数 \(a<b\) 和单位向量 \(\nu\)，使
\begin{equation}\label{b1:eq:approximate-jump-halves}
 \begin{aligned}
 \lim_{r\downarrow0}\frac2{\omega_dr^d}
       \int_{B_r^+(x,\nu)}|u(y)-b|\,\mathrm dy&=0,\\
 \lim_{r\downarrow0}\frac2{\omega_dr^d}
       \int_{B_r^-(x,\nu)}|u(y)-a|\,\mathrm dy&=0.
 \end{aligned}
\end{equation}
则称 \(x\) 为 \(u\) 的\term[jinsi-tiaoyue-dian]{近似跳跃点}[approximate jump point]。
全体这样的点组成\term[tiaoyue-ji]{跳跃集}[jump set]，记为 \(J_u\)。
\end{definition}
\symidx{\(J_u\)：函数的近似跳跃点集}

此处使用的是半球中的 \(L^1\) 平均，不只是某个稀疏点列上的极限，
也不只是半球内依测度趋近一个数。尤其对无界函数，
上一节已经说明后两种较弱信息不能自动给出平均绝对误差趋零。

把半球缩放到单位球更便于比较不同的跳跃描述。设
\[
 w_{a,b,\nu}(z)=
 \begin{cases}
  b,&z\cdot\nu>0,\\
  a,&z\cdot\nu<0.
 \end{cases}
\]
平面上任意补值。条件\eqref{b1:eq:approximate-jump-halves}恰好等价于
\begin{equation}\label{b1:eq:jump-step-blowup}
 u(x+r\,\cdot)\longrightarrow w_{a,b,\nu}
 \quad\text{于 }L^1\bigl(B(0,1)\bigr).
\end{equation}
这个缩放极限保留了跳跃的高度和方向，却不读取原代表在点 \(x\) 的指定值。

\begin{proposition}\label{b1:prop:approximate-jump-uniqueness}
在 \(a<b\) 的排序约定下，近似跳跃点的三元组 \((a,b,\nu)\) 唯一。
跳跃集及其三元组只依赖函数的几乎处处等价类。
\end{proposition}
\begin{proof}
若两个三元组满足定义，则由 \(L^1\) 极限的唯一性，所得两个阶跃函数几乎处处相等。
每个阶跃函数都恰好取两个不同的本性值，而且每个值占半个单位球。
因此它们的无序值集相同；排序以后，两侧的数值分别相同。

较大值所占的半球于是也几乎处处相同。若两个单位法向不同，
对应正半球的对称差便含有一个非空开锥与单位球的交，具有正体积，矛盾。
故法向也相同。若不排序，唯一的另一种写法是同时交换两个值并把法向反号。
最后，零测修改不改变任何半球积分，故不改变定义。
\end{proof}

\subsection{单侧近似极限}
\label{b1:subsec:280202}

现在可以在 \(J_u\) 上无歧义地记
\[
 u^-(x)\coloneqq a,\quad u^+(x)\coloneqq b,\quad \nu_u(x)\coloneqq\nu.
\]
称 \(u^\pm(x)\) 为相应两侧的\term[dance-jinsi-jixian]{单侧近似极限}[one-sided approximate limits]。
这里的正负号指向由 \(\nu_u\) 确定的半球；依本书约定，正侧同时是数值较大的一侧。
\symidx{\(u^+\)：跳跃点的较大单侧近似极限}
\symidx{\(u^-\)：跳跃点的较小单侧近似极限}

\begin{proposition}\label{b1:prop:jump-approximate-limits}
若 \(x\in J_u\)，则
\begin{equation}\label{b1:eq:jump-upper-lower-values}
 u^\wedge(x)=u^-(x)<u^+(x)=u^\vee(x).
\end{equation}
特别地，\(J_u\subseteq S_u\)。对每个 \(u^-(x)<t<u^+(x)\)，还有
\begin{equation}\label{b1:eq:jump-level-half-density}
 \lim_{r\downarrow0}
 \frac{m_d\bigl(\{u>t\}\cap B(x,r)\bigr)}{\omega_dr^d}
 =\frac12.
\end{equation}
\end{proposition}
\begin{proof}
记 \(a=u^-(x)\)、\(b=u^+(x)\)。在正半球中，\(\{u\le t\}\) 上的
\(|u-b|\) 至少为 \(b-t>0\)，故 Chebyshev 不等式说明这个例外集的相对体积趋零。
在负半球中同理，\(\{u>t\}\) 的相对体积趋零。
两项相加就得到\eqref{b1:eq:jump-level-half-density}。

若 \(s<a\)，两半球中的误差估计都给出 \(\{u<s\}\) 的密度为零；
若 \(s>a\)，在 \(a\) 与 \(\min(s,b)\) 之间取一个阈值，
负半球便给出 \(\{u<s\}\) 的正密度下界。
所以近似下极限恰为 \(a\)。对上极限作相同的论证，得到 \(b\)。
由于两者不同，有限近似极限不可能存在，故 \(x\in S_u\)。
\end{proof}

在跳跃点上还可取中点值
\[
 u^*(x)=\frac{u^+(x)+u^-(x)}2.
\]
由两个半球的平均收敛立即得到
\[
 \lim_{r\downarrow0}\frac1{\omega_dr^d}\int_{B(x,r)}u(y)\,\mathrm dy=u^*(x).
\]
但这不使跳跃点成为 Lebesgue 点：事实上
\[
 \lim_{r\downarrow0}\frac1{\omega_dr^d}\int_{B(x,r)}|u(y)-u^*(x)|\,\mathrm dy
 =\frac{u^+(x)-u^-(x)}2>0.
\]
因此球平均有极限、有限近似极限存在、平均绝对振荡趋零，是需要区别的三种陈述。
因此这里的 \(u^*\) 与\chapref{b1:ch:14}由球平均定义的精确代表在 \(J_u\) 上相同；
本书始终把 \(u^*\) 称为球平均精确代表，而把 \(\widetilde u\) 称为近似极限代表，
不另用这两个记号表示彼此冲突的全域点值约定。

\subsection{跳跃法向量}
\label{b1:subsec:280203}

向量 \(\nu_u\) 称为\term[tiaoyue-faxiangliang]{跳跃法向量}[jump normal]。
它从低值侧指向高值侧。这个取向首先由函数的两侧行为定义，
尚不需要预先知道 \(J_u\) 是一张光滑曲面。
\symidx{\(\nu_u\)：从低值侧指向高值侧的单位跳跃法向量}

取盒子 \(\Omega=(-1,1)^d\)，并令
\[
 u(x)=
 \begin{cases}
  a,&x_d<0,\\
  b,&x_d>0,
 \end{cases}
 \quad a<b.
\]
界面上任意补值。此时
\[
 J_u=\Omega\cap\{x_d=0\},\quad
 u^-=a,\quad u^+=b,\quad \nu_u=e_d.
\]
逐坐标应用 Fubini 定理和一维分部积分，可得
\[
 Du=(b-a)e_d\,\mathcal H^{d-1}\llcorner
       \bigl(\Omega\cap\{x_d=0\}\bigr).
\]
具体地，前 \(d-1\) 个方向的测试积分为零；最后一个方向则满足
\[
 -\int_\Omega u\,\partial_d\varphi\,\mathrm dx
 =(b-a)\int_{(-1,1)^{d-1}}\varphi(x',0)\,\mathrm dx'.
\]
这也直接检验了导数中法向因子的正负号。

对光滑集合的示性函数，数值 \(1\) 位于集合内部，数值 \(0\) 位于外部。
所以在其边界上，\(\nu_{\mathbf1_E}=-\nu_E\)，其中 \(\nu_E\) 是集合外法向。
跳跃公式因此写成
\[
 D\mathbf1_E=-\nu_E\,\mathcal H^{d-1}\llcorner\partial E,
\]
与\cref{b1:prop:smooth-domain-perimeter}一致。不能将集合外法向直接代入
排序后的函数跳跃公式而漏掉这个负号。

为了稍后把两侧值和法向放进测度积分，还需确认它们的可测性。
\begin{proposition}\label{b1:prop:jump-data-borel}
对 \(u\in L^1_{\mathrm{loc}}(\Omega)\)，集合 \(J_u\) 是 Borel 集，
而 \(u^\pm\) 与 \(\nu_u\) 是其上的 Borel 函数。
\end{proposition}
\begin{proof}
由\cref{b1:thm:lebesgue-borel-representative}，先选取 \(u\) 的 Borel 代表；
这不改变任何跳跃数据。
由\cref{b1:prop:approximate-limits-basic}，集合
\[
 A=\{x:-\infty<u^\wedge(x)<u^\vee(x)<+\infty\}
\]
是 Borel 的。在 \(A\) 上置 \(a=u^\wedge(x)\)、\(b=u^\vee(x)\)，并记
\[
 F_r(x,v)=\frac1{\omega_d}\int_{B(0,1)}
       |u(x+rz)-w_{a,b,v}(z)|\,\mathrm dz.
\]
先在距边界有正距离的可数个内部集合上考察足够小的 \(r\)。
对固定 \(r,v\)，这个函数关于 \(x\) 是 Borel 的；关于正半径 \(r\) 连续。
后一事实先对连续函数成立，再在共同紧集上用 \(L^1\) 逼近即可得到。
它关于 \(v\) 也连续，因为两半球的对称差体积随方向连续变化。

取单位球面的一个可数稠密子集 \(D\)。条件 \(x\in J_u\) 等价于：
对每个正整数 \(k\)，存在 \(v\in D\) 和正整数 \(m\)，使所有充分留在定义域内、
满足 \(0<r<1/m\) 的正有理半径都有 \(F_r(x,v)<1/k\)。
这是一组可数的 Borel 条件。
正向蕴含由阶跃极限及方向的连续性立即得到。
反过来，选取这些方向 \(v_k\)；在同时足够小的半径上用三角不等式，得到
\[
 \frac1{\omega_d}\int_{B(0,1)}
 |w_{a,b,v_k}-w_{a,b,v_\ell}|
 \le \frac1k+\frac1\ell.
\]
从单位球面到这些阶跃函数的映射连续且单射，球面又紧，
故 \(v_k\) 收敛到唯一方向 \(v\)。再次用三角不等式，并利用半径连续性，
便得 \(F_r(x,v)\to0\)。因此上述可数条件确实刻画 \(J_u\)。

每次按固定枚举选取第一个合条件的 \((v,m)\)，便使 \(v_k(x)\) 成为 Borel 函数；
其极限就是 \(\nu_u(x)\)。最后由\eqref{b1:eq:jump-upper-lower-values}，
两侧值也是 Borel 函数。
\end{proof}

\subsection{\texorpdfstring{\(J_u\)}{Ju} 的可整流性}
\label{b1:subsec:280204}

跳跃方向并不预设曲面，但有界变差条件将这些局部方向约束在可数个可整流集合上。
证明的关键是只选一组稠密的好层值，而不要求每一个层值都有有限周长。

\begin{proposition}\label{b1:prop:bv-jump-rectifiable}
若 \(u\in BV(\Omega)\)，则 \(J_u\) 可数 \((d-1)\)-可整流。
\end{proposition}
\begin{proof}
由有界变差函数的余面积公式，几乎每个 \(t\in\mathbb R\) 都使
\(E_t=\{u>t\}\) 在 \(\Omega\) 内有有限周长。
这个好层值集合中可以选取一个可数稠密集 \(D\)：在每个有理端点开区间中取一个好值即可。
不能未经选择便声称所有有理层值都好。

若 \(x\in J_u\)，在 \((u^-(x),u^+(x))\) 内取 \(t\in D\)。
由\eqref{b1:eq:jump-level-half-density}，\(E_t\) 在 \(x\) 的密度为 \(1/2\)，
所以 \(x\in\partial^eE_t\)。于是
\[
 J_u\subseteq\bigcup_{t\in D}(\partial^eE_t\cap\Omega).
\]
De Giorgi 结构\cref{b1:thm:de-giorgi-structure}说明，
每个 \(\partial^eE_t\cap\Omega\) 与可数可整流的
\(\partial^*E_t\cap\Omega\) 只差一个 \(\mathcal H^{d-1}\)-零集。
可数合并后就得到所需覆盖。这里没有声称每个跳跃点都属于某个指定层集的约化边界；
用到的是两类边界之间已经证明的模零关系。
当 \(d=1\) 时，各好层集的内部边界局部只有有限个点，
同一论证给出 \(J_u\) 至多可数。
\end{proof}

可整流性尚未回答两个问题：近似不连续能否在跳跃集之外占据正的余维一测度，
以及导数在跳跃面上究竟有多大。完整答案如下。
\begin{theorem}[有界变差函数的跳跃结构（深层结构定理）]\label{b1:thm:bv-jump-structure}
若 \(u\in BV(\Omega)\)，则
\begin{equation}\label{b1:eq:bv-discontinuity-jump-null}
 \mathcal H^{d-1}(S_u\setminus J_u)=0.
\end{equation}
而且其向量导数测度在 \(J_u\) 上满足
\begin{equation}\label{b1:eq:bv-jump-derivative}
 Du\llcorner J_u
 =(u^+-u^-)\nu_u\,\mathcal H^{d-1}\llcorner J_u.
\end{equation}
因而
\[
 |Du|\llcorner J_u
 =(u^+-u^-)\,\mathcal H^{d-1}\llcorner J_u,
 \quad
 \int_{J_u}(u^+-u^-)\,\mathrm d\mathcal H^{d-1}
 \le |Du|(\Omega).
\]
\end{theorem}
\begin{proofsketch}
上述可整流性已完整证明。余下的两项是有界变差函数精细结构的深层部分；
Alberti 与 Mantegazza 在《A Note on the Theory of SBV Functions》的注1.2中准确列出相应的两侧 \(L^1\) 迹及导数公式%
\cite[Remark 1.2, pp. 376--377]{AlbertiMantegazza1997SBV}。
该注把精细结构的证明指向 Evans 与 Gariepy 的相关论述，
并不在这几页中完整证明曲面迹定理。这里说明尚需的桥梁。

第一项输入是有界变差函数的非有限近似值集合
\[
 N_\infty
 =\{x:u^\wedge(x)\notin\mathbb R
          \text{ 或 }u^\vee(x)\notin\mathbb R\}
\]
具有零 \(\mathcal H^{d-1}\) 测度。
在 \(S_u\setminus N_\infty\) 上，两近似极限有限且不相等。
在它们之间选稠密好层值，由\cref{b1:prop:approximate-oscillation-level-boundary}，
便把这些点覆盖在好层集的测度论边界上。
由上一章，这些边界模零由可数个可整流曲面承载。

第二项输入是有界变差函数在可整流曲面上的两侧迹定理：
对曲面面积几乎每点，都存在有限的两个半球 \(L^1\) 迹；
在该曲面上的导数限制等于两侧迹之差乘有向单位法向与面积测度。
两迹相等时，合并半球平均便得到有限近似极限；
故在所覆盖的 \(S_u\) 上，两迹几乎处处不同，按大小排序即得到 \(J_u\)。
这证明\eqref{b1:eq:bv-discontinuity-jump-null}；
曲面上的导数限制公式再给出\eqref{b1:eq:bv-jump-derivative}。
最后，单位法向的长度为一，故向量测度的变差就是所写的非负密度积分。

本节不展开非有限近似值的余维一估计，以及有界变差函数在曲面两侧的迹及其分部积分公式。
它们需要进一步的有界变差函数切片和精细代表理论，不能把第23章的 Sobolev 迹定理
直接套在一般有界变差函数上来替代。
\end{proofsketch}

上一节的无穷近似值例子因此没有被漏掉：那些点仍属于 \(S_u\)，
却不符合两个有限值的跳跃定义；结构定理把它们连同其他非跳跃例外一起控制在
\(\mathcal H^{d-1}\)-零集中。
另一方面，公式控制的是跳跃高度与面积的乘积，
并不保证跳跃集本身有有限面积。
它只立即保证
\[
 \mathcal H^{d-1}
 \bigl(\{x\in J_u:u^+(x)-u^-(x)\ge 1/k\}\bigr)
 \le k\,|Du|(\Omega).
\]
下一节将把这部分导数与体积上的密度、以及剩余的 Cantor 部分分开，
从而得到三种互不混淆的变差来源。

% !TeX root = ../../Book1.tex
\section{导数测度的分解}
\label{b1:sec:2803}

分布导数把一个函数的全部一阶变化保存在同一向量测度中。
但这些变化可以分布在体积中，也可以集中在跳跃面上，还可能既不占体积、又不由跳跃面承载。
本节区分这三种情形。前两种分别连接 Sobolev 梯度与有限周长集合，
第三种则解释连续函数为何仍可能没有可积的弱导数。

以下 \(\Omega\subseteq\mathbb R^d\) 为开集，\(u\in BV(\Omega)\) 为实值。
跳跃值与法向沿用\secref{b1:sec:2802}的约定：
\(u^+>u^-\)，\(\nu_u\) 从低值侧指向高值侧。
深层结构的准确陈述可参看 Alberti、Mantegazza 的
《A Note on the Theory of SBV Functions》\cite[Remark 1.2]{AlbertiMantegazza1997SBV}；
系统学习可继续阅读 Ambrosio、Fusco、Pallara 的
《Functions of Bounded Variation and Free Discontinuity Problems》%
\cite[Chapter 3]{Ambrosio2000Functions}。
本节把测度分解的直接后果完整推出，而将近似微分的识别和链律中的深层步骤明确列为证明概要。
% 实读作者公开正式论文 Alberti--Mantegazza 1997：
% 作者版第2--3页（刊印376--377）Remark 1.2；
% 作者版第4--5页（刊印378--379）Remark 2.1；
% 其§2.3--2.5另供下一节紧性使用。AFP只核实出版与章节导航，无全文阅读声明。
% 原文Su表示无近似极限集；本书保留Su与真正两侧跳跃集Ju的区别。

\subsection{绝对连续部分}
\label{b1:subsec:280301}

先不使用函数的几何性质。对 \(Du\) 的每个分量应用 Lebesgue 分解，
得到唯一的表示
\begin{equation}\label{b1:eq:bv-lebesgue-decomposition}
 Du=g\,m_d+D^su,\quad g\in L^1(\Omega;\mathbb R^d),\quad
 |D^su|\perp m_d.
\end{equation}
记 \(D^au=g\,m_d\)，称它为导数测度的绝对连续部分。
具体说，各分量的奇异部分分别集中在某个体积零的 Borel 集上；
取这有限多个集合的并 \(N\)，便有
\[
 D^au=Du\mathbin{\llcorner}(\Omega\setminus N),
 \quad
 D^su=Du\mathbin{\llcorner}N.
\]
这里限制记号只表示在相应 Borel 集上取测度。
向量极分解给出 \(|D^au|=|g|m_d\)；
不同分量的分解由标量唯一性保证相容。因而这一步只需要第6章的测度论，
尚未证明 \(g\) 与逐点的微分有什么联系。
\symidx{\(D^au\)：\(Du\) 的绝对连续部分}
\symidx{\(D^su\)：\(Du\) 的奇异部分}

\begin{theorem}[有界变差函数的近似梯度]
\label{b1:thm:bv-approximate-gradient}
函数 \(u\) 在 \(m_d\)-几乎处处近似可微。
若在近似极限存在处用该极限选取中心值，则其近似微分可写成
\[
 \operatorname{ap}Du(x)\,h=\nabla u(x)\cdot h,
 \quad \nabla u(x)=g(x)
 \quad m_d\text{-几乎处处},
\]
其中 \(g\) 是\eqref{b1:eq:bv-lebesgue-decomposition}中的密度。
特别地，\(D^au=\nabla u\,m_d\)，且 \(\nabla u\in L^1(\Omega;\mathbb R^d)\)。
\end{theorem}
\begin{proofsketch}
这里不仅需要 Radon--Nikodym 定理，还要证明函数本身在几乎每点有一阶线性近似。
一种路线先用变差的局部平均控制函数振荡，把函数限制在可数个 Lipschitz 片上；
在这些片的密度点应用 Rademacher 定理，再用分布测试识别所得梯度与 \(g\)。
把振荡估计转成覆盖几乎所有点的 Lipschitz 分片，以及最后的密度识别，
是本节未展开的两步，不能由测度有密度直接代替。
上述结论在所引短文 Remark 1.2 和 equation (1.2) 之后明确陈述，
该处将完整论证指向 Evans--Gariepy 原版的 Section 6.1。
\end{proofsketch}

从此对一般有界变差函数写 \(\nabla u\) 时，指这个近似梯度的几乎处处等价类。
它不是 \(Du\) 的另一种写法，也不要求一个任意指定的代表具有经典偏导数。
若 \(u\in W^{1,1}(\Omega)\)，分布导数已经完全由弱梯度表示，
上述唯一性便说明近似梯度与弱梯度几乎处处一致。

\begin{proposition}\label{b1:prop:bv-sobolev-singular-criterion}
对于 \(u\in BV(\Omega)\)，有
\[
 u\in W^{1,1}(\Omega)\quad\Longleftrightarrow\quad D^su=0.
\]
\end{proposition}
\begin{proof}
正向由弱导数的定义和 Lebesgue 分解唯一性得到。
反向若 \(D^su=0\)，则每个分布偏导数都由 \(g_i\in L^1(\Omega)\) 表示。
再加上有界变差函数定义中的 \(u\in L^1(\Omega)\)，恰好满足 \(W^{1,1}\) 的定义。
\end{proof}

\subsection{跳跃部分}
\label{b1:subsec:280302}

\begin{definition}\label{b1:def:bv-jump-derivative}
将导数测度限制到跳跃集 \(J_u\)，所得向量测度记为
\[
 D^ju=Du\mathbin{\llcorner}J_u,
\]
称为导数测度的\term[tiaoyue-bufen]{跳跃部分}[jump part]。
\end{definition}
\symidx{\(D^ju\)：\(Du\) 的跳跃部分}

由\cref{b1:thm:bv-jump-structure}，它具有明确的表面密度：
\begin{equation}\label{b1:eq:bv-jump-density}
 D^ju=(u^+-u^-)\nu_u\,
       \mathcal H^{d-1}\mathbin{\llcorner}J_u.
\end{equation}
跳跃的高度、朝向与承载的面积在这个公式中各占一个位置。
如果改用不排序的两侧记号，同时交换 \(u^+\)、\(u^-\) 并把法向反号，
右侧向量不变。因此测度不依赖人为的定向选择。
由于 \(|\nu_u|=1\)，对每个 Borel 集 \(A\subseteq\Omega\)，
\[
 |D^ju|(A)=\int_{A\cap J_u}(u^+-u^-)\,\mathrm d\mathcal H^{d-1}.
\]
这给出跳跃高度的可积性，但不保证 \(\mathcal H^{d-1}(J_u)<\infty\)：
很多很小的跳跃，可以具有有限总变差和无限总面积。

例如在 \((0,1)\) 内令
\[
 v(x)=\sum_{n=1}^{\infty}2^{-n}\mathbf1_{(2^{-n},1)}(x).
\]
该级数一致绝对收敛，每项的分布导数为 \(2^{-n}\delta_{2^{-n}}\)。
对测试函数逐项积分合法，故
\[
 Dv=\sum_{n=1}^{\infty}2^{-n}\delta_{2^{-n}},
 \quad |Dv|\bigl((0,1)\bigr)=1.
\]
在每个 \(2^{-n}\) 处跳跃高度为 \(2^{-n}\)，其他内部点在一个邻域内函数恒定，
所以 \(J_v=\{\,2^{-n}:n\ge1\,\}\)。
一维的 \(\mathcal H^0\) 是计数测度，因而 \(\mathcal H^0(J_v)=\infty\)。

\subsection{Cantor 部分}
\label{b1:subsec:280303}

跳跃集可数 \((d-1)\)-可整流，因而是体积零集。
所以 \(D^ju\) 是 \(D^su\) 在 \(J_u\) 上的限制，余下的部分仍为奇异测度。

\begin{definition}\label{b1:def:bv-cantor-derivative}
规定
\[
 D^cu=D^su-D^ju
     =D^su\mathbin{\llcorner}(\Omega\setminus J_u),
\]
称 \(D^cu\) 为导数测度的\term[cantor-bufen]{Cantor 部分}[Cantor part]。
\end{definition}
\symidx{\(D^cu\)：\(Du\) 的 Cantor 部分}

\begin{theorem}[有界变差函数的导数测度三分解]
\label{b1:thm:bv-derivative-decomposition}
对于任意实值 \(u\in BV(\Omega)\)，有
\begin{equation}\label{b1:eq:bv-three-part-decomposition}
 Du=\nabla u\,m_d
    +(u^+-u^-)\nu_u\,
      \mathcal H^{d-1}\mathbin{\llcorner}J_u
    +D^cu.
\end{equation}
这三项两两互相奇异。更进一步，若 Borel 集 \(A\subseteq\Omega\)
关于 \(\mathcal H^{d-1}\) 为 \(\sigma\)-有限，则
\[
 |D^cu|(A)=0.
\]
有限近似极限 \(\widetilde u\) 在 \(|D^cu|\)-几乎处处存在。
\end{theorem}
\begin{proofsketch}
分解等式由前面两项的识别和 \(D^cu\) 的定义得到。
若 \(N\) 承载 \(D^su\) 且 \(m_d(N)=0\)，三项分别集中于
\(\Omega\setminus(N\cup J_u)\)、\(J_u\)、\(N\setminus J_u\)，故两两互相奇异。
这里尚不能只由 \(D^cu(J_u)=0\) 推出它不充任何有限面积集。

后一性质是有界变差函数的精细结构所给出的额外结论。
证明须把任意有限 \(\mathcal H^{d-1}\) 集与函数在不同方向上的一维截面比较，
识别其中能承载导数的部分为两侧迹的跳跃部分；
剩余部分在这样的集合上均为零，再由可数并处理 \(\sigma\)-有限集。
这一几何识别没有在本节证明，准确结论见前引 Remark 1.2 的式(1.1)--(1.4)。
结合\cref{b1:thm:bv-jump-structure}，
没有有限近似极限的集合除去 \(J_u\) 只有 \(\mathcal H^{d-1}\) 零测余项，
因此它也是 \(|D^cu|\) 的零测集。这给出最后一句。
\end{proofsketch}

\begin{proposition}\label{b1:prop:bv-decomposition-variation}
三部分分解唯一，并且对每个 Borel 集 \(A\subseteq\Omega\)，
\begin{equation}\label{b1:eq:bv-three-part-variation}
 |Du|(A)=\int_A|\nabla u|\,\mathrm dx
   +\int_{A\cap J_u}(u^+-u^-)\,\mathrm d\mathcal H^{d-1}
   +|D^cu|(A).
\end{equation}
\end{proposition}
\begin{proof}
Lebesgue 分解唯一确定第一项。
在 \(J_u\) 上限制剩余测度，唯一确定第二项；相减便唯一确定第三项。
取上述三个互不交的承载集，在每个承载集上 \(Du\) 等于相应分量，
所以其变差也等于该分量的变差。
由变差的可数可加性相加，再用向量密度的变差公式，即得所述等式。
\end{proof}

Cantor 部分的名称不是说它必须集中在三分 Cantor 集上，
而是说它代表不能由体积密度和跳跃面描述的奇异变化。
在一维，它恰好成为无原子的奇异部分，这一点可以不借高维结构定理直接证明。

\begin{proposition}\label{b1:prop:one-dimensional-bv-decomposition}
设 \(I=(a,b)\) 为有界区间，\(u\in BV(I)\)，\(\mu=Du\)。
则存在常数 \(c\)，使
\[
 \widehat u(x)=c+\mu\bigl((a,x]\bigr)
\]
是 \(u\) 的右连续代表。
其左右极限之差等于 \(\mu(\{x\})\)，且
\[
 D^ju=\sum_{x\in I}\mu(\{x\})\delta_x.
\]
这里非零项至多可数，绝对值之和有限。
测度 \(D^cu\) 则是 \(\mu\) 的奇异部分删去全部原子后的剩余测度；
它的累积分布函数连续。
\end{proposition}
\begin{proof}
函数 \(F(x)=\mu\bigl((a,x]\bigr)\) 有界。由符号测度积分的 Fubini 计算，
对 \(\varphi\in C_c^\infty(I)\)，
\[
 \int_a^b F(x)\varphi'(x)\,\mathrm dx
 =\int_I\left(\int_t^b\varphi'(x)\,\mathrm dx\right)\mathrm d\mu(t)
 =-\int_I\varphi\,\mathrm d\mu.
\]
因此 \(D(u-F)=0\)。
\secref{b1:sec:2004}证明的零弱导数判别说明 \(u-F=c\) 几乎处处。
有限测度的从上、从下连续性分别给出 \(F\) 的右连续性和左极限，
并且 \(F(x)-F(x-)=\mu(\{x\})\)。
非零原子至多可数，因为对每个正整数 \(n\)，质量绝对值至少 \(1/n\) 的原子只有有限个。
对有限原子集求和再取上确界，得到
\(\sum_x|\mu(\{x\})|\le|\mu|(I)\)。

累积分布代表在非原子点连续，在原子点具有不相同的两侧极限；
因此它的近似跳跃集恰为非零原子集。
每个原子的有向跳跃等于其有符号质量，给出 \(D^ju\) 的公式。
绝对连续测度没有原子，故原子全部属于 \(\mu\) 的奇异部分。
删去它们以后所得奇异测度无原子，其累积分布在每点左右连续，因而连续。
\end{proof}

例如，习题\ref{b1:ex:14-cantor-function}中的 \(F_C\) 满足
\[
 DF_C=\rho_C,\quad D^aF_C=0,\quad D^jF_C=0,\quad D^cF_C=\rho_C
 \quad\text{在 }(0,1)\text{ 上}.
\]
第一式由刚才的测试积分得出；后面三式来自 \(\rho_C\) 奇异且无原子。
于是函数
\[
 u(x)=x+\mathbf1_{(1/2,1)}(x)+F_C(x)
\]
在同一区间内同时具有三种非零变化：
\(D^au=m_1\)、\(D^ju=\delta_{1/2}\)、\(D^cu=\rho_C\)，总变差为 \(3\)。
这也表明连续性只能排除跳跃，不能排除 Cantor 部分。

最后记录复合函数所需的精细链律。
下一节会用它分别测试跳跃的存在与梯度的极限。

\begin{theorem}[Volpert 链律的 \(C^1\) 版本]
\label{b1:thm:volpert-chain-rule}
设 \(u\in BV(\Omega)\)，\(\Phi\in C^1(\mathbb R)\)，且 \(\Phi'\) 有界。
若 \(m_d(\Omega)<\infty\) 或 \(\Phi(0)=0\)，则 \(\Phi(u)\in BV(\Omega)\)，并且
\begin{equation}\label{b1:eq:volpert-chain-rule}
 D\bigl(\Phi(u)\bigr)
 =\Phi'(\widetilde u)(D^au+D^cu)
  +\bigl(\Phi(u^+)-\Phi(u^-)\bigr)\nu_u\,
    \mathcal H^{d-1}\mathbin{\llcorner}J_u.
\end{equation}
其中 \(\widetilde u\) 在有限近似极限不存在处可任意补值。
\end{theorem}
\term[volpert-lianlv]{Volpert 链律}[Volpert chain rule]在跳跃处使用两端值之差，
而不是把某个任选的点值代入 \(\Phi'\) 后乘上跳跃高度。
\begin{proofsketch}
一维时先用测度累积代表，把连续变化与原子变化分开：
连续部分满足通常的链式乘法，原子在复合后变成
\(\Phi(u(x+))-\Phi(u(x-))\)。
高维论证用有界变差函数的切片定理将各方向的分布导数还原为这些一维导数，
再重组为向量测度。这里省略的是切片定理及其与三个导数部分的相容性，
不是对公式作形式求导。
所用版本见 Alberti 与 Mantegazza 的《A Note on the Theory of SBV Functions》%
\cite[Remark 2.1, equation (2.1)]{AlbertiMantegazza1997SBV}；
本书按既定方向约定显式写出法向因子。
可积性由 \(|\Phi(t)|\le|\Phi(0)|+\|\Phi'\|_\infty\cdot|t|\) 及给定假设保证，
测度右侧的有限性则由三部分变差公式和 \(\Phi\) 的 Lipschitz 界保证。
\end{proofsketch}

% !TeX root = ../../Book1.tex
\section{特殊有界变差函数与紧性}
\label{b1:sec:2804}

在允许界面出现的变分问题中，变化的代价通常分成两类：
区域内部的梯度能量，以及界面本身的面积代价。
如果把所有有界变差函数都作为候选，而能量只计入这两项，
Cantor 部分便可能不付出任何代价。
因此需要一个排除该部分、同时保留跳跃的函数类。

\ProseParagraphBreak

本节先定义特殊有界变差函数，说明它对有界变差空间中的严格收敛并不闭合，
再记录带有值域、超线性梯度和跳跃面积控制的 Ambrosio 紧性定理。
定理用到\chapref{b1:ch:27}中有界变差函数的紧性，以及前节作为深层输入记录的 Volpert 链律。
后者如何用于特殊有界变差函数空间的闭合性，可参看 Alberti、Mantegazza 的
《A Note on the Theory of SBV Functions》%
\cite[Theorem 1.4, Propositions 2.3 and 2.5]{AlbertiMantegazza1997SBV}。
% 实读作者公开版第4--7页，对应刊印378--381页；
% Theorem1.4只陈述紧性及绝对连续/跳跃部分分开收敛，不包含面积下半连续性。
% 本节取其增长函数f(t)=1(t>0),f(0)=0，采用正弦/余弦测试展开闭合桥。

\subsection{特殊有界变差函数}
\label{b1:subsec:280401}

\begin{definition}\label{b1:def:sbv-function}
设 \(u\in BV(\Omega)\)。若 \(D^cu=0\)，则称 \(u\) 为%
\term[teshu-youjie-biancha-hanshu]{特殊有界变差函数}[special function of bounded variation]，
记 \(u\in SBV(\Omega)\)。\foreignidx{SBV}因此
\[
 Du=\nabla u\,m_d+(u^+-u^-)\nu_u\,
 \mathcal H^{d-1}\mathbin{\llcorner}J_u.
\]
\end{definition}
\symidx{\(SBV(\Omega)\)：特殊有界变差函数空间}

每个 \(W^{1,1}\) 函数都属于 \(SBV\)。
有界区域内有限周长集合的示性函数也属于 \(SBV\)：
它的导数由约化边界上的面积表示，故完全属于跳跃部分。
但特殊有界变差函数空间既不等于 Sobolev 空间，也不等于整个有界变差函数空间：
一个非零阶跃不属于 \(W^{1,1}\)，而前节的连续 Cantor 函数不属于 \(SBV\)。
定义本身也不要求跳跃集具有有限面积，更不要求它是闭集或由有限张光滑曲面组成。

\begin{proposition}\label{b1:prop:sbv-not-strictly-closed}
在区间 \((0,1)\) 上，特殊有界变差函数空间对有界变差空间中的严格收敛并不闭合。
甚至存在 \(0\le u_n\le1\) 的连续分段线性函数，
满足 \(u_n\to F_C\) 一致收敛、\(|Du_n|\bigl((0,1)\bigr)=1\)，
而 \(D^cF_C\ne0\)。
\end{proposition}
\begin{proof}
记 \(C_n\) 为三分 Cantor 集构造中的第 \(n\) 层闭区间之并。
在每个长度 \(3^{-n}\) 的基本区间上，用直线连接 Cantor 函数的两个端点值；
在其余间隙上保持 Cantor 函数的常值，所得函数记为 \(u_n\)。
两端值之差为 \(2^{-n}\)，故
\[
 u_n'=(3/2)^n\cdot\mathbf1_{C_n}\quad\text{几乎处处}.
\]
函数 \(u_n\) 连续且分段线性，因而属于 \(W^{1,1}\subseteq SBV\)。
它在每个基本区间内与 \(F_C\) 相差至多 \(2^{-n}\)，
在间隙上则相等，所以一致收敛成立。
又因为 \(m_1(C_n)=(2/3)^n\)，
\[
 \int_0^1|u_n'|\,\mathrm dx=1=|DF_C|\bigl((0,1)\bigr).
\]
这正是严格收敛，但 \(D^cF_C=\rho_C\ne0\)。
\end{proof}

这组函数没有跳跃，且梯度的 \(L^1\) 范数一致有界，
仍能把全部导数集中成奇异连续测度。
若 \(p>1\)，同一计算给出
\[
 \int_0^1|u_n'|^p\,\mathrm dx=(3/2)^{n(p-1)}\longrightarrow\infty.
\]
超线性的梯度控制正是后面紧性条件中的一个实质要求。

\subsection{紧性}
\label{b1:subsec:280403}

\begin{theorem}[Ambrosio 紧性：有界值域与超线性梯度（深层紧性定理）]
\label{b1:thm:ambrosio-sbv-compactness}
设 \(\Omega\) 为有界开集，\(1<p<\infty\)，\(u_j\in SBV(\Omega)\)，且
\[
 \sup_j\left\{\|u_j\|_{L^\infty(\Omega)}
       +\int_\Omega|\nabla u_j|^p\,\mathrm dx
       +\mathcal H^{d-1}(J_{u_j})\right\}<\infty.
\]
则存在子列与 \(u\in SBV(\Omega)\cap L^\infty(\Omega)\)，使
\[
 u_j\longrightarrow u\quad\text{在每个 }L^r(\Omega),\ 1\le r<\infty,
 \quad
 \nabla u_j\rightharpoonup\nabla u\quad\text{在 }L^p(\Omega;\mathbb R^d).
\]
此外，\(D^au_j\) 与 \(D^ju_j\) 分别\WeakStar{}收敛到 \(D^au\) 与 \(D^ju\)，
测度收敛均以 \(C_0(\Omega;\mathbb R^d)\) 为测试类。
\end{theorem}
\term[ambrosio-sbv-jinxing]{Ambrosio 紧性定理}[Ambrosio compactness theorem]
比一般有界变差函数的紧性多给出一个结论：极限仍没有 Cantor 部分。
\begin{proofsketch}
以下展开这个特殊版本的闭合步骤；所依赖的额外深层工具是
\cref{b1:thm:volpert-chain-rule}，不把它当作已经由普通 Sobolev 链律推出的结果。
记统一的值域界为 \(M\)，跳跃面积界为 \(S\)。
两侧近似值也在 \([-M,M]\) 中，故
\[
 |Du_j|(\Omega)
 \le m_d(\Omega)^{1-1/p}\|\nabla u_j\|_{L^p(\Omega)}
      +2M S.
\]
于是 \(BV\) 范数一致有界。由\cref{b1:thm:bv-local-compactness}取局部 \(L^1\) 收敛子列。
对于内部紧集穷竭 \(K_n\uparrow\Omega\)，
\(\int_{\Omega\setminus K_n}|u_j|\le M\,m_d(\Omega\setminus K_n)\) 一致趋零；
所以\cref{b1:cor:bv-global-compactness-with-tails}给出全域 \(L^1\) 收敛，
极限 \(u\in BV(\Omega)\) 且 \(|u|\le M\)。
再取子列，使 \(u_j\to u\) 几乎处处，且 \(\nabla u_j\rightharpoonup h\) 于 \(L^p\)。
尚需识别 \(h\) 并排除 \(D^cu\)。

取任意 \(\Phi\in C^1(\mathbb R)\)，满足 \(0\le\Phi\le1\) 且 \(\Phi'\) 有界。
由链律和 \(D^cu_j=0\)，测度
\[
 T_j^\Phi=D\bigl(\Phi(u_j)\bigr)-\Phi'(u_j)\nabla u_j\,m_d
\]
的总变差不超过 \(S\)。
函数 \(\Phi(u_j)\to\Phi(u)\) 于 \(L^1\)。
若 \(p'=p/(p-1)\)，对每个 \(q\in L^{p'}\)，
控制收敛给出
\(\Phi'(u_j)q\to\Phi'(u)q\) 于 \(L^{p'}\)，
因而 \(\Phi'(u_j)\nabla u_j\rightharpoonup\Phi'(u)h\)。
对紧支集光滑向量场取极限，再用变差的测试表述，得到
\[
 \left|D\bigl(\Phi(u)\bigr)-\Phi'(u)h\,m_d\right|(\Omega)\le S.
\]
再对 \(u\) 用链律，绝对连续、Cantor、跳跃三项互相奇异，因此
\[
 \int_\Omega|\Phi'(\widetilde u)|\,\mathrm d\sigma\le S,
 \quad
 \sigma=|\nabla u-h|\,m_d+|D^cu|.
\]
有限近似极限在 \(\sigma\)-几乎处处存在，故积分不依赖其余点的补值。
分别取
\[
 \Phi_N(t)=\frac{1+\sin(Nt)}2,\quad
 \Psi_N(t)=\frac{1+\cos(Nt)}2.
\]
因为 \(|\Phi_N'(t)|+|\Psi_N'(t)|\ge N/2\)，两次估计相加即得
\[
 \frac N2\,\sigma(\Omega)\le2S.
\]
令 \(N\to\infty\)，便有 \(\sigma=0\)，即 \(h=\nabla u\) 且 \(D^cu=0\)。
这个振荡测试正是所引论文 Propositions 2.3、2.5 在固定面积界下的核心机制。

最后，\(|u_j-u|\le2M\) 给出
\(\|u_j-u\|_r^r\le(2M)^{r-1}\|u_j-u\|_1\)，得到全部有限 \(r\) 的强收敛。
绝对连续部分的收敛来自梯度弱收敛，
跳跃部分的收敛来自 \(D^ju_j=Du_j-D^au_j\)
与\cref{b1:prop:bv-distributional-weak-star}。
\end{proofsketch}

这一定理为自由间断能量的直接法提供紧性，但尚未给出跳跃面积的下半连续性。
下一节先明确具体能量与容许类，再补上这项测度论接口。


\begin{chaptersummary}
BV 导数的三部分分别由体积、跳跃界面与不充余维一有限面积集的奇异变化承载。
这个区分不是把任意奇异测度机械拆成“有原子”和“无原子”：
只有在一维，原子恰与跳跃对应；高维的跳跃面测度本身也可以没有原子。

点值的恢复同样需要区分尺度。跳跃处没有单一有限近似极限，
但两个半球平均给出不同的有限值；它们的差与朝向共同决定跳跃导数。
给奇异测度乘一个函数时，应使用相应精细代表，
不能只凭 Lebesgue 几乎处处相等就任意替换。

SBV 排除了 Cantor 部分，却不自动对 BV 严格收敛闭合。
超线性梯度界与跳跃面积界分别防止不同形式的变化集中。
在有界值域内，这些界配合链律和下半连续性，给出 Mumford–Shah 型极小元。
存在性只产生函数等价类及其跳跃结构，没有同时给出唯一性、闭裂纹表示或界面正则性。
至此，本卷的积分、微分、覆盖、可整流性和紧性已经汇入同一个变分问题；
继续深入时，仍应逐项辨认哪些工具已经证明，哪些属于后续理论。
\end{chaptersummary}

\BookExercises

% 第二十八章习题临时稿；由根任务接在唯一的BookExercises之后，不另建入口。
% 八题均为自拟；所用BV结构、Volpert链律、SBV紧性来自本章明确陈述的版本，
% 不把其深层证明移入基础习题。逐题注释记录训练目标与相对正文的增量。

\begin{exercise}\label{b1:ex:28-radial-piecewise-decomposition}
在 \(\Omega=B(0,2)\subset\mathbb R^2\) 中，给定 \(a\in\mathbb R\)，令
\[
 u_a(x)=
 \begin{cases}
 |x|^2,&|x|<1,\\
 2|x|+a,&1<|x|<2.
 \end{cases}
\]
单位圆上的值任取。求 \(D^au_a,D^ju_a,D^cu_a\) 以及
\(|Du_a|(\Omega)\)，并判断 \(u_a\) 何时属于 \(W^{1,1}(\Omega)\)。
按照 \(u^+>u^-\) 的约定，分别写出可能跳跃处的两个值和法向。
\end{exercise}
% 来源：自拟；将两侧常数推广到两段径向光滑函数，训练体积分量、界面分量及排序反向。
\begin{solution}
两侧的经典梯度分别为 \(2x\) 和 \(2x/|x|\)。对内部光滑测试场逐区积分，
由\cref{b1:prop:smooth-domain-perimeter}，单位圆上的贡献等于外侧迹减内侧迹，
再乘从内向外的单位方向。因此
\[
 D^au_a=
 \left(2x\,\mathbf1_{\{|x|<1\}}
       +\frac{2x}{|x|}\mathbf1_{\{1<|x|<2\}}\right)m_2,
\]
\[
 D^ju_a=(a+1)x\,\mathcal H^1\llcorner\partial B(0,1),
 \quad D^cu_a=0.
\]
外圆是定义域的边界，不进入相对分布导数。两份非零分量互相奇异，所以
\[
 \begin{aligned}
 |Du_a|(\Omega)
 &=\int_{B(0,1)}2|x|\,\mathrm dx
   +\int_{B(0,2)\setminus B(0,1)}2\,\mathrm dx
   +2\pi|a+1|\\
 &=\frac{22\pi}{3}+2\pi|a+1|.
 \end{aligned}
\]
特别地，每个 \(u_a\) 都属于 \(SBV(\Omega)\)。

若 \(a>-1\)，圆周上的高值在外侧，
\[
 u_a^+=a+2,\quad u_a^-=1,\quad \nu_{u_a}=x.
\]
若 \(a<-1\)，高值改在内侧，
\[
 u_a^+=1,\quad u_a^-=a+2,\quad \nu_{u_a}=-x.
\]
两种情形中 \((u_a^+-u_a^-)\nu_{u_a}\) 都等于 \((a+1)x\)。
当 \(a=-1\) 时两侧迹相同，跳跃集为空；此时显式的有界梯度就是弱梯度，
甚至 \(u_{-1}\in W^{1,\infty}(\Omega)\)。
其余 \(a\) 有非零奇异导数，故不属于 \(W^{1,1}(\Omega)\)。
\end{solution}

\begin{exercise}\label{b1:ex:28-chain-jump-cancellation}
\begin{enumerate}
\item\label{b1:item:28-chain-square-cancellation}
在 \((-1,1)\) 上令 \(u(x)=x+\operatorname{sgn}x\)，原点任意补值。
取一个 \(C^1(\mathbb R)\) 函数 \(\Phi\)，使它在 \([-2,2]\) 上等于 \(t^2\)，且 \(\Phi'\) 有界。
计算 \(Du\) 与 \(D\bigl(\Phi(u)\bigr)\)，说明光滑的值域变换可以完全消除跳跃，却不使函数变成常数。
\item\label{b1:item:28-bilipschitz-value-change}
设 \(m_d(\Omega)<\infty\)、\(u\in BV(\Omega)\)，且
\[
 \Psi\in C^1(\mathbb R),\quad
 0<c\le|\Psi'(t)|\le C<\infty\quad(t\in\mathbb R).
\]
证明 \(J_{\Psi(u)}=J_u\)，写出排序后法向的变化，并证明
\[
 u\in SBV(\Omega)\quad\Longleftrightarrow\quad
 \Psi(u)\in SBV(\Omega).
\]
\end{enumerate}
\end{exercise}
% 来源：自拟；第一部分让两侧像值相合，第二部分辨明双Lipschitz值域变化不能产生同样消失。
\begin{solution}
第一部分中，原点两侧极限为 \(-1,1\)，所以
\[
 Du=m_1+2\delta_0.
\]
可具体选择 \(\Phi(t)=t^2\) 于 \(|t|\le2\)，在两侧分别以端点切线继续；
这样 \(\Phi\in C^1\) 且 \(|\Phi'|\le4\)。
复合函数为
\[
 \Phi(u(x))=(1+|x|)^2.
\]
它连续且 Lipschitz，其弱导数为
\[
 D\bigl(\Phi(u)\bigr)
 =2(1+|x|)\operatorname{sgn}x\,m_1.
\]
链律中的跳跃系数恰为 \(\Phi(1)-\Phi(-1)=0\)。
这不是把原点的代表值改掉，而是确实使两侧极限变为同一个值。

第二部分中，连续函数 \(\Psi'\) 从不为零，所以符号固定。
微积分基本定理给出
\[
 c|s-t|\le|\Psi(s)-\Psi(t)|\le C|s-t|.
\]
因此 \(\Psi\) 单调且映满 \(\mathbb R\)，其逆也为 Lipschitz 函数。
半球中的 \(L^1\) 逼近误差在复合后至多乘以 \(C\)，逆向至多乘以 \(1/c\)；
所以两侧不同有限值的跳跃点恰好保持不变。
若 \(\Psi'>0\)，高低次序与 \(\nu_u\) 不变；若 \(\Psi'<0\)，
高低次序交换，\(\nu_{\Psi(u)}=-\nu_u\)。
两种情形的跳跃高度均为 \(|\Psi(u^+)-\Psi(u^-)|\)。

由\cref{b1:thm:volpert-chain-rule}，
\[
 D^c\bigl(\Psi(u)\bigr)=\Psi'(\widetilde u)\,D^cu.
\]
近似代表在 \(|D^cu|\)-几乎处处有限，故
\[
 c|D^cu|\le |D^c\bigl(\Psi(u)\bigr)|\le C|D^cu|.
\]
Cantor 部分为零当且仅当复合后的 Cantor 部分为零，得到 \(SBV\) 的等价性。
有限体积条件保证即使 \(\Psi(0)\ne0\)，复合函数仍属于 \(L^1(\Omega)\)。
\end{solution}

\begin{exercise}\label{b1:ex:28-cantor-cylinder-and-atoms}
设 \(\mathcal C\subset[0,1]\) 为三分 Cantor 集，\(C\) 为其连续分布函数，
\(\mu_C=DC\) 为无原子的 Cantor 概率测度。在 \(\Omega=(0,1)^2\) 中令
\[
 u(x,y)=C(x),\quad v(x,y)=\mathbf1_{\{y>1/2\}}.
\]
\begin{enumerate}
\item\label{b1:item:28-cantor-cylinder-decomposition}
计算二者的三分解，并说明 \(|Du|\)、\(|Dv|\) 都没有原子，
但一份是纯 Cantor 部分，另一份是纯跳跃部分。
\item\label{b1:item:28-cantor-cylinder-diffuse}
不调用 Cantor 部分的结构定理，直接证明
\[
 (\mu_C\otimes m_1)(A)=0
\]
对每个 \(\mathcal H^1\)-\(\sigma\) 有限的 Borel 集 \(A\subset\Omega\) 成立。
据此说明 \((\mathcal C\times(0,1))\cap\Omega\) 不是
\(\mathcal H^1\)-\(\sigma\) 有限集。
\end{enumerate}
\end{exercise}
% 来源：自拟；从一维Cantor函数增至二维乘积，并用互不相交的竖截面直接区分无原子与无跳跃。
\begin{solution}
对测试函数先沿 \(x\) 积分，再沿 \(y\) 积分，得到
\[
 Du=e_1(\mu_C\otimes m_1),\quad
 Dv=e_2\mathcal H^1\llcorner\bigl((0,1)\times\{1/2\}\bigr).
\]
第一份测度集中在体积为零的 Cantor 柱上。函数 \(u\) 连续，没有跳跃点，
所以 \(D^cu=Du\)，而 \(D^au=D^ju=0\)。
函数 \(v\) 在水平线两侧有不同常值，因此 \(D^jv=Dv\)，其他两部分为零。
乘积测度和线上的长度测度都不给单点质量。
可见高维导数测度“无原子”并不等于“没有跳跃部分”。

先设 \(\mathcal H^1(A)<\infty\)。对每个 \(x\)，记
\(A_x=\{y:(x,y)\in A\}\)。
不同竖线上的集合互不相交，且由等距不变性，
\[
 \mathcal H^1\bigl(\{x\}\times A_x\bigr)=m_1(A_x).
\]
有限测度空间中，互不相交且各自具有正测度的集合至多可数：
对每个正整数 \(n\)，其中测度至少为 \(1/n\) 的集合只能有有限个。
所以 \(m_1(A_x)>0\) 的 \(x\) 至多可数。
由 Tonelli 定理与 \(\mu_C\) 无原子，
\[
 (\mu_C\otimes m_1)(A)=\int m_1(A_x)\,\mathrm d\mu_C(x)=0.
\]
对 \(\mathcal H^1\)-\(\sigma\) 有限集作可数分解，即得一般结论。
另一方面，\(\mu_C\otimes m_1\) 给 Cantor 柱的质量为 \(1\)，
所以该柱不可能具有所述 \(\sigma\) 有限性。
\end{solution}

\begin{exercise}\label{b1:ex:28-jump-cantor-approximation}
沿用习题\ref{b1:ex:28-cantor-cylinder-and-atoms}的一维 Cantor 分布函数。
令 \(I_{n,k}\)、\(1\le k\le2^n\)，为第 \(n\) 步留下的基本闭区间，
按从左到右排列，并记其中心为 \(c_{n,k}\)。定义
\[
 \mu_n=2^{-n}\sum_{k=1}^{2^n}\delta_{c_{n,k}},
 \quad
 u_n(x)=\mu_n((0,x])\quad(0<x<1).
\]
证明 \(u_n\in SBV(0,1)\)，且 \(u_n\to C\) 严格于 \(BV(0,1)\)，
但极限不属于 \(SBV(0,1)\)。
分别计算
\[
 \mathcal H^0(J_{u_n}),\quad
 \int_{J_{u_n}}(u_n^+-u_n^-)^q\,\mathrm d\mathcal H^0
 \quad(q>0),
\]
并指出\cref{b1:thm:ambrosio-sbv-compactness}中的哪项统一控制失效。
最后判断这是否也是通常 \(BV\) 范数下的收敛。
\end{exercise}
% 来源：自拟；与正文连续分段线性逼近相对照，令梯度恒为零而跳跃点数量发散，并量化振幅罚项。
\begin{solution}
每个 \(u_n\) 是有限阶跃函数，\(Du_n=\mu_n\)，所以它只有跳跃部分，
\[
 |Du_n|\bigl((0,1)\bigr)=1,\quad \nabla u_n=0
 \quad\text{几乎处处}.
\]
每个基本区间的 Cantor 质量与 \(\mu_n\) 质量都等于 \(2^{-n}\)。
在基本区间之间的间隙，两份分布函数取同一个值；在一个基本区间内，
二者都位于相同的两个相邻累计质量之间。因此
\[
 \sup_{0<x<1}|u_n-C|\le2^{-n}.
\]
于是有 \(L^1\) 收敛，同时 \(|Du_n|\bigl((0,1)\bigr)=|DC|\bigl((0,1)\bigr)=1\)，
这正是严格收敛。极限连续而具有非零奇异导数，故其导数全是 Cantor 部分，不属于 \(SBV\)。

每个中心点的跳跃高度为 \(2^{-n}\)，故
\[
 \mathcal H^0(J_{u_n})=2^n,\quad
 \int_{J_{u_n}}(u_n^+-u_n^-)^q\,\mathrm d\mathcal H^0
 =2^{n(1-q)}.
\]
这列函数的值域一致受控，所有 \(p>1\) 的梯度积分都为零；
失效的是跳跃集的面积，即此处的点数控制。
特别地，\(q\ge1\) 时振幅的 \(q\) 次罚项有界，仍不足以代替跳跃点数；
\(q>1\) 时它甚至趋零。

最后，有限原子测度 \(\mu_n\) 与无原子的 \(\mu_C\) 相互奇异，所以
\[
 |D(u_n-C)|\bigl((0,1)\bigr)=|\mu_n-\mu_C|\bigl((0,1)\bigr)=2.
\]
因此并没有通常 \(BV\) 范数的强收敛。严格收敛只要求各自总变差的数值收敛，
不要求导数测度之差的全变差趋零。
\end{solution}

\begin{exercise}\label{b1:ex:28-bv-removable-interface}
设 \(d\ge2\)，\(K\) 是开集 \(\Omega\) 中的相对闭集，并且
\(\mathcal H^{d-1}|_K\) 为 \(\sigma\) 有限。
给定 \(u\in BV(\Omega)\)，假设
\[
 u\in W_{\mathrm{loc}}^{1,1}(\Omega\setminus K),
 \quad \mathcal H^{d-1}(K\cap J_u)=0.
\]
证明 \(u\in W^{1,1}(\Omega)\)。
分别利用习题\ref{b1:ex:28-radial-piecewise-decomposition}和%
\ref{b1:ex:28-cantor-cylinder-and-atoms}说明：
若删除“没有界面跳跃”的条件，或删除 \(K\) 的 \(\sigma\) 有限条件，结论均可能失败。
\end{exercise}
% 来源：自拟；将三分解转为带BV先验的奇集可去性，两个反例分别孤立跳跃和Cantor障碍。
\begin{solution}
在开集 \(\Omega\setminus K\) 内，局部 Sobolev 条件说明分布导数由局部可积函数表示。
按可数内部开覆盖及分解唯一性，\(D^ju\)、\(D^cu\) 在该开集上的限制均为零，
所以两份奇异测度都集中在 \(K\) 上。

跳跃表示与 \(\mathcal H^{d-1}(K\cap J_u)=0\) 给出 \(|D^ju|(K)=0\)。
由\cref{b1:thm:bv-derivative-decomposition}中 Cantor 部分不充
\(\mathcal H^{d-1}\)-\(\sigma\) 有限集的结论，\(|D^cu|(K)=0\)。
因此 \(D^su=0\)。再由\cref{b1:prop:bv-sobolev-singular-criterion}，
得到 \(u\in W^{1,1}(\Omega)\)。
这里的全域可积性来自原先的 \(BV(\Omega)\) 假设，不是仅由域外局部 Sobolev 性推出。

若删除跳跃条件，取第一题中 \(a\ne-1\) 的函数，并令 \(K=\partial B(0,1)\)。
它在 \(K\) 外分片光滑，\(K\) 的长度有限，却有非零界面导数。
若删除 \(\sigma\) 有限条件，取第三题中的 \(u(x,y)=C(x)\) 及 Cantor 柱 \(K\)。
它连续，故没有跳跃；在 \(K\) 外局部常于 \(x\)，因而光滑，
却仍有非零 Cantor 导数。第三题已经证明这时 \(K\) 确实不是长度的 \(\sigma\) 有限集。
\end{solution}

\begin{exercise}\label{b1:ex:28-cantor-reparameterization}
设 \(C\) 与 \(\mu_C\) 如前，\(\Phi\in C^1(\mathbb R)\) 且 \(\Phi'\) 有界。
证明
\[
 C_\#\mu_C=m_1\llcorner[0,1],
 \quad
 |D^c\bigl(\Phi(C)\bigr)|((0,1))=\int_0^1|\Phi'(t)|\,\mathrm dt.
\]
据此证明：\(\Phi(C)\in SBV(0,1)\) 当且仅当 \(\Phi\) 在 \([0,1]\) 上为常数。
将这一结论与习题\ref{b1:ex:28-chain-jump-cancellation}中的消跳现象比较。
\end{exercise}
% 来源：自拟；通过分布函数的推前测度精确计算链律中的Cantor能量，不重复原始Cantor反例。
\begin{solution}
对 \(0<t<1\)，由 \(C\) 连续、非减且映满 \([0,1]\)，
存在最大的 \(b_t\in[0,1]\)，使 \(C(b_t)=t\)。
于是
\[
 \{x:C(x)\le t\}=[0,b_t],\quad
 \mu_C(\{C\le t\})=\mu_C([0,b_t])=C(b_t)=t.
\]
端点 \(t=0,1\) 同样成立。比较区间的分布函数并用测度唯一性，
得到 \(C_\#\mu_C=m_1\llcorner[0,1]\)。

函数 \(C\) 连续，近似代表就是它本身，导数为纯 Cantor 测度。
由链律，
\[
 D\bigl(\Phi(C)\bigr)=\Phi'(C)\,\mu_C.
\]
复合函数仍连续，因而没有跳跃；上式又集中在 Lebesgue 零集上，
所以整份导数都是 Cantor 部分。推前等式给出
\[
 |D^c\bigl(\Phi(C)\bigr)|((0,1))
 =\int|\Phi'(C(x))|\,\mathrm d\mu_C(x)
 =\int_0^1|\Phi'(t)|\,\mathrm dt.
\]
它为零当且仅当连续函数 \(\Phi'\) 在 \([0,1]\) 上恒为零，
也就是 \(\Phi\) 在该区间为常数。
跳跃只比较两个端点的像值，两点可以在非恒定变换下合并；
这里的奇异变化却经 \(C\) 遍历整个值域，不能只让两个值相合就消除。
\end{solution}

\begin{exercise}\label{b1:ex:28-energy-scaling}
设 \(\Omega\subset\mathbb R^d\) 有界，\(u\in BV(\Omega)\)。
给定 \(\rho>0\)、\(s\ne0\)、\(b\in\mathbb R\)，令
\[
 \Omega_\rho=\rho\Omega,\quad
 T_\rho(x)=\rho x,\quad
 v(x)=s\,u(x/\rho)+b.
\]
证明三份导数均满足
\[
 D^\ell v=s\rho^{d-1}(T_\rho)_\#D^\ell u,
 \quad \ell\in\{a,j,c\},
\]
并写出跳跃集与法向的变化。
进一步设 \(u\in SBV(\Omega)\)、\(g\in L^2(\Omega)\)、\(1<p<\infty\)，
给定 \(\beta,\lambda>0\)，记允许取值 \(+\infty\) 的能量为
\[
 \mathcal F_{\beta,\lambda}(u;\Omega,g)
 =\int_\Omega|\nabla u|^p\,\mathrm dx
 +\beta\mathcal H^{d-1}(J_u)
 +\lambda\int_\Omega|u-g|^2\,\mathrm dx.
\]
令 \(g_\rho(x)=s\,g(x/\rho)+b\)，求这三项各自的尺度。
再求 \(\beta_\rho,\lambda_\rho\)，使
\[
 \mathcal F_{\beta_\rho,\lambda_\rho}(v;\Omega_\rho,g_\rho)
 =|s|^p\rho^{d-p}\mathcal F_{\beta,\lambda}(u;\Omega,g).
\]
\end{exercise}
% 来源：自拟；同时追踪体积、余维一面积及振幅，要求权重随尺度改变而不是只比较导数范数。
\begin{solution}
在分布测试中作 \(x=\rho y\) 的变量代换，得到
\[
 Dv=s\rho^{d-1}(T_\rho)_\#Du.
\]
伸缩保持 Lebesgue 零集和正则分布密度的类型，也把半球近似跳跃条件逐项带到新点，
因此
\[
 J_v=\rho J_u,\quad
 \nu_v(\rho x)=\operatorname{sgn}(s)\nu_u(x).
\]
跳跃值是 \(su^\pm+b\) 重新排序后的两个值。绝对连续部分和跳跃部分分别变换为题述形式，
从总导数中减去这两部分，Cantor 部分也得到同一公式。
特别地，
\[
 \nabla v(x)=\frac{s}{\rho}\nabla u(x/\rho)
 \quad\text{几乎处处}.
\]

于是三项分别为
\[
 \int_{\Omega_\rho}|\nabla v|^p
 =|s|^p\rho^{d-p}\int_\Omega|\nabla u|^p,
\]
\[
 \mathcal H^{d-1}(J_v)=\rho^{d-1}\mathcal H^{d-1}(J_u),
 \quad
 \int_{\Omega_\rho}|v-g_\rho|^2
 =s^2\rho^d\int_\Omega|u-g|^2.
\]
比较三个系数，得到
\[
 \beta_\rho=\beta|s|^p\rho^{1-p},
 \quad
 \lambda_\rho=\lambda|s|^{p-2}\rho^{-p}.
\]
可见梯度、跳跃面积和拟合误差不具有相同的长度尺度。
若同时保留正文的值域约束，原来的 \(|u|\le M\) 应相应改为
\(|v|\le |s|M+|b|\)；不能在任意振幅变换后仍声称同一个固定值域类不变。
\end{solution}

\begin{exercise}\label{b1:ex:28-one-jump-mumford-shah}
取 \(L,a,\beta,\lambda>0\)，在 \(\Omega=(-L,L)\) 上令
\(g(x)=a\operatorname{sgn}x\)，并考虑 \(p=2\) 的能量
\[
 \mathcal F(u)=\int_{-L}^L|u'|^2\,\mathrm dx
 +\beta\mathcal H^0(J_u)
 +\lambda\int_{-L}^L|u-g|^2\,\mathrm dx.
\]
先把容许函数限制为 \([-a,a]\) 值的常数函数，或至多有一个跳跃的分段常值函数。
求此受限类中的最小值及全部极小元，说明何时出现两个不同的极小元。
然后用一个无跳跃的线性函数说明：这个阈值不能直接宣称为完整
\(SBV\) 容许类中的全局阈值。
\end{exercise}
% 来源：自拟；精确解一个有限参数模型，再用新竞争者主动检验“受限极小值”与真正全局极小值的差别。
\begin{solution}
常数函数的最优值为 \(0\)，因为
\[
 \lambda\int_{-L}^L|c-g|^2\,\mathrm dx
 =2\lambda L(c^2+a^2).
\]
这给出能量 \(2\lambda La^2\)。

若实际跳跃发生在 \(t\in[0,L)\)，左右两段分别以数据的均值作为最优常数：
\[
 c_1=a\frac{t-L}{t+L},\quad c_2=a.
\]
二者都在容许值域内且彼此不同。右段没有拟合误差，左段的最小平方误差为
\[
 a^2(t+L)-a^2\frac{(t-L)^2}{t+L}
 =\frac{4a^2Lt}{L+t}.
\]
若 \(t<0\)，相同计算给出 \(4a^2L|t|/(L+|t|)\)。
因此一个真实跳跃的最小能量是
\[
 \beta+\frac{4\lambda a^2L|t|}{L+|t|}\ge\beta,
\]
且等号只在 \(t=0\) 时成立，此时函数正是 \(g\)。
故受限类的最小值为
\[
 \min\{\beta,\,2\lambda La^2\}.
\]
若 \(\beta<2\lambda La^2\)，唯一极小元为 \(g\)；
若 \(\beta>2\lambda La^2\)，唯一极小元为零函数；
若二者相等，这两个函数恰为全部极小元。函数在单点上的重定义不计作新的极小元。

现在取 \(u_\varepsilon(x)=\varepsilon x\)，其中 \(0<\varepsilon L\le a\)。
它满足同一值域约束且没有跳跃，但不属于上面的分段常值受限类。直接积分得到
\[
 \mathcal F(u_\varepsilon)-\mathcal F(0)
 =-2\lambda aL^2\varepsilon
  +\left(2L+\frac{2\lambda L^3}{3}\right)\varepsilon^2.
\]
充分小的正 \(\varepsilon\) 使右边为负。
特别在 \(\beta=2\lambda La^2\) 时，它的能量同时低于上述两个受限极小元。
所以受限类中的非唯一性计算完全成立，却没有证明完整模型在该参数处具有同样的极小元或阈值。
\end{solution}

